\documentclass[11pt,a4paper]{book}
\usepackage{isabelle,isabellesym}
\usepackage{graphicx}
\graphicspath{document}

% further packages required for unusual symbols (see also
% isabellesym.sty), use only when needed

\usepackage{pifont}

\usepackage{latexsym}
\usepackage{amssymb}
  %for \<leadsto>, \<box>, \<diamond>, \<sqsupset>, \<mho>, \<Join>,
  %\<lhd>, \<lesssim>, \<greatersim>, \<lessapprox>, \<greaterapprox>,
  %\<triangleq>, \<yen>, \<lozenge>

%\usepackage[greek,english]{babel}
  %option greek for \<euro>
  %option english (default language) for \<guillemotleft>, \<guillemotright>

%\usepackage[latin1]{inputenc}
  %for \<onesuperior>, \<onequarter>, \<twosuperior>, \<onehalf>,
  %\<threesuperior>, \<threequarters>, \<degree>

%\usepackage[only,bigsqcap]{stmaryrd}
  %for \<Sqinter>

%\usepackage{eufrak}
  %for \<AA> ... \<ZZ>, \<aa> ... \<zz> (also included in amssymb)

%\usepackage{textcomp}
  %for \<cent>, \<currency>

% this should be the last package used
\usepackage{pdfsetup}

% urls in roman style, theory text in math-similar italics
\urlstyle{rm}
\isabellestyle{it}

\begin{document}

\title{$IMP_{concur}$ - A Semantics for a Language with Concurrency based on HOL-CSP}
\author{Burkhart Wolff}

\maketitle

\chapter*{Abstract}
The theory $IMP_{concur}$  provides a  
programming-language aspect to HOL-CSP. It extends the well-known 
imperative language \verb+IMP+ originated by Glenn Winskell by concurrent
primitives like \verb+lock+ and \verb+unlock+ for semaphores, thread-local and
system global program variables and the notion of threads that run in interleaving
semantics. Threads can be parametric and used to construct concurrent process
architectures.

Representing $IMP_{concur}$ semantics can be done
by a straight-forward translation into HOL-CSP; thus, $IMP_{concur}$ can be seen as
a thin layer on top of the theory of Concurrent Sequential Processes 
of Hoare, Brookes and Roscoe in its formalization HOL-CSP in Isabelle/HOL.

Notwithstanding its simplicity, $IMP_{concur}$ covers many aspects of concurrency, 
non-termination, divergence, synchronization, states and race-conditions.
We believe that $IMP_{concur}$ may be relevant for readers interested in the
link between programming languages and process algebras. Resulting new ways to 
concurrent program verification are an interesting object for future study.
We also present a language extension for dynamic reconfigurations of thread-
architectures. 

This session is complemented by a syntactic frontend for $IMP_{concur}$, called $CIL$,
which has been authored by Mathilde Needham and Zineddine Kermadj on their 
internship project in summer 2025.

\tableofcontents

\pagebreak

% sane default for proof documents
\parindent 0pt\parskip 0.5ex

% generated text of all theories
%
\begin{isabellebody}%
\setisabellecontext{IMP{\isacharunderscore}{\kern0pt}concur}%
%
\isadelimdocument
%
\endisadelimdocument
%
\isatagdocument
%
\isamarkupchapter{\isa{IMP\isactrlsub c\isactrlsub o\isactrlsub n\isactrlsub c\isactrlsub u\isactrlsub r}-A Shallow Embedding in HOL-CSP%
}
\isamarkuptrue%
%
\endisatagdocument
{\isafolddocument}%
%
\isadelimdocument
%
\endisadelimdocument
%
\isadelimtheory
%
\endisadelimtheory
%
\isatagtheory
\isakeywordONE{theory}\isamarkupfalse%
\ IMP{\isacharunderscore}{\kern0pt}concur\isanewline
\isakeywordTWO{imports}\ {\isachardoublequoteopen}HOL{\isacharminus}{\kern0pt}CSPM{\isachardoublequoteclose}\ \ \isanewline
\isanewline
\isakeywordTWO{begin}%
\endisatagtheory
{\isafoldtheory}%
%
\isadelimtheory
%
\endisadelimtheory
%
\isadelimdocument
%
\endisadelimdocument
%
\isatagdocument
%
\isamarkupsection{Introduction%
}
\isamarkuptrue%
%
\endisatagdocument
{\isafolddocument}%
%
\isadelimdocument
%
\endisadelimdocument
%
\begin{isamarkuptext}%
\isa{IMP\isactrlsub c\isactrlsub o\isactrlsub n\isactrlsub c\isactrlsub u\isactrlsub r} is intended to provide a more common (though still theoretic)
programming-language aspect to CSP, and could be an interesting object
of study in itself. This thin layer over HOL-CSP establishes a link 
between programming languages and process-algebras, a link that had already
been present in the influential OCCAM language.

\isa{IMP\isactrlsub c\isactrlsub o\isactrlsub n\isactrlsub c\isactrlsub u\isactrlsub r} comprises the standard elements of the IMP - language (\isa{SKIP}, \isa{assignment},
\isa{IF\ {\isacharunderscore}{\kern0pt}\ THEN\ {\isacharunderscore}{\kern0pt}\ ELSE}, \isa{WHILE\ {\isacharunderscore}{\kern0pt}\ DO\ {\isacharunderscore}{\kern0pt}} and the sequential composition \isa{{\isacharunderscore}{\kern0pt}\ {\isacharsemicolon}{\kern0pt}\ {\isacharunderscore}{\kern0pt}}), plus 
the new features:

%
\begin{enumerate}%
\item semaphores with \isa{lock} and \isa{unlock}, and

\item thread-global shared variables accessible 
via \isa{LOAD} and \isa{STORE} operations.%
\end{enumerate}

This file contains:

%
\begin{enumerate}%
\item the abstract and concrete syntax of \isa{IMP\isactrlsub c\isactrlsub o\isactrlsub n\isactrlsub c\isactrlsub u\isactrlsub r}

\item bricks-specifications for semaphores and global shared variables,

\item a denotational semantics of \isa{IMP\isactrlsub c\isactrlsub o\isactrlsub n\isactrlsub c\isactrlsub u\isactrlsub r} converting \isa{IMP\isactrlsub c\isactrlsub o\isactrlsub n\isactrlsub c\isactrlsub u\isactrlsub r}-program 
systems into HOL-CSPM, and

\item some examples and tests.%
\end{enumerate}

The general theory should be developed elsewhere; this sample file shows the 
global construction principle of a combined semantics.%
\end{isamarkuptext}\isamarkuptrue%
%
\isadelimdocument
%
\endisadelimdocument
%
\isatagdocument
%
\isamarkupsection{The Syntax of \isa{IMP\isactrlsub c\isactrlsub o\isactrlsub n\isactrlsub c\isactrlsub u\isactrlsub r}%
}
\isamarkuptrue%
%
\endisatagdocument
{\isafolddocument}%
%
\isadelimdocument
%
\endisadelimdocument
\isakeywordONE{type{\isacharunderscore}{\kern0pt}synonym}\isamarkupfalse%
\ SV\ {\isacharequal}{\kern0pt}\ string\ \ %
\isamarkupcmt{Shared Variables%
}\isanewline
\isakeywordONE{type{\isacharunderscore}{\kern0pt}synonym}\isamarkupfalse%
\ V\ \ {\isacharequal}{\kern0pt}\ string\ \ %
\isamarkupcmt{(Thread-)Local Variables%
}\isanewline
\isakeywordONE{type{\isacharunderscore}{\kern0pt}synonym}\isamarkupfalse%
\ MV\ {\isacharequal}{\kern0pt}\ int\ \ \ \ \ %
\isamarkupcmt{Mutual exclusion variables (MUTEX ids)%
}\isanewline
\isanewline
\isakeywordONE{type{\isacharunderscore}{\kern0pt}synonym}\isamarkupfalse%
\ D\ {\isacharequal}{\kern0pt}\ int\isanewline
\isanewline
\isakeywordONE{type{\isacharunderscore}{\kern0pt}synonym}\isamarkupfalse%
\ {\isachardoublequoteopen}{\isasymsigma}{\isachardoublequoteclose}\ {\isacharequal}{\kern0pt}\ {\isacartoucheopen}V\ {\isasymRightarrow}\ D{\isacartoucheclose}\isanewline
\isanewline
\isakeywordONE{type{\isacharunderscore}{\kern0pt}synonym}\isamarkupfalse%
\ {\isachardoublequoteopen}E\isactrlsub a\isactrlsub r\isactrlsub i\isactrlsub t\isactrlsub h{\isachardoublequoteclose}\ {\isacharequal}{\kern0pt}\ {\isacartoucheopen}{\isasymsigma}\ {\isasymRightarrow}\ D{\isacartoucheclose}\isanewline
\isakeywordONE{type{\isacharunderscore}{\kern0pt}synonym}\isamarkupfalse%
\ {\isachardoublequoteopen}E\isactrlsub b\isactrlsub o\isactrlsub o\isactrlsub l{\isachardoublequoteclose}\ \ {\isacharequal}{\kern0pt}\ {\isacartoucheopen}{\isasymsigma}\ {\isasymRightarrow}\ bool{\isacartoucheclose}\isanewline
\isakeywordONE{type{\isacharunderscore}{\kern0pt}synonym}\isamarkupfalse%
\ {\isachardoublequoteopen}F{\isachardoublequoteclose}\ \ \ \ \ {\isacharequal}{\kern0pt}\ {\isachardoublequoteopen}{\isasymsigma}\ {\isasymRightarrow}\ {\isasymsigma}{\isachardoublequoteclose}\isanewline
\isanewline
\isakeywordONE{datatype}\isamarkupfalse%
\ com\ {\isacharequal}{\kern0pt}\ SKIP\isanewline
\ \ \ \ \ \ \ \ \ \ \ \ \ {\isacharbar}{\kern0pt}\ assign\ \ V\ E\isactrlsub a\isactrlsub r\isactrlsub i\isactrlsub t\isactrlsub h\ \ \ \ \ \ \ {\isacharparenleft}{\kern0pt}{\isachardoublequoteopen}\ {\isacharunderscore}{\kern0pt}\ {\isacharcolon}{\kern0pt}{\isacharequal}{\kern0pt}\ {\isacharunderscore}{\kern0pt}{\isachardoublequoteclose}\ {\isacharbrackleft}{\kern0pt}{\isadigit{9}}{\isadigit{0}}{\isacharcomma}{\kern0pt}{\isadigit{9}}{\isadigit{0}}{\isacharbrackright}{\kern0pt}{\isadigit{8}}{\isadigit{0}}{\isacharparenright}{\kern0pt}\isanewline
\ \ \ \ \ \ \ \ \ \ \ \ \ {\isacharbar}{\kern0pt}\ seq\ \ \ \ \ com\ com\ \ \ \ \ \ \ {\isacharparenleft}{\kern0pt}\isakeywordTWO{infixr}\ {\isachardoublequoteopen}{\isacharsemicolon}{\kern0pt}{\isachardoublequoteclose}\ {\isadigit{7}}{\isadigit{8}}{\isacharparenright}{\kern0pt}\isanewline
\ \ \ \ \ \ \ \ \ \ \ \ \ {\isacharbar}{\kern0pt}\ cond\ \ \ {\isachardoublequoteopen}E\isactrlsub b\isactrlsub o\isactrlsub o\isactrlsub l{\isachardoublequoteclose}\ com\ com\ {\isacharparenleft}{\kern0pt}{\isachardoublequoteopen}IF\ {\isacharparenleft}{\kern0pt}{\isacharunderscore}{\kern0pt}{\isacharparenright}{\kern0pt}{\isacharslash}{\kern0pt}\ THEN\ {\isacharparenleft}{\kern0pt}{\isacharunderscore}{\kern0pt}{\isacharparenright}{\kern0pt}{\isacharslash}{\kern0pt}\ ELSE\ {\isacharparenleft}{\kern0pt}{\isacharunderscore}{\kern0pt}{\isacharparenright}{\kern0pt}{\isacharslash}{\kern0pt}{\isachardoublequoteclose}\ \isanewline
\ \ \ \ \ \ \ \ \ \ \ \ \ \ \ \ \ \ \ \ \ \ \ \ \ \ \ \ \ \ \ \ \ \ \ \ \ \ \ \ \ \ \ \ \ \ \ {\isacharbrackleft}{\kern0pt}{\isadigit{0}}{\isacharcomma}{\kern0pt}{\isadigit{0}}{\isacharcomma}{\kern0pt}{\isadigit{7}}{\isadigit{9}}{\isacharbrackright}{\kern0pt}{\isadigit{7}}{\isadigit{9}}{\isacharparenright}{\kern0pt}\isanewline
\ \ \ \ \ \ \ \ \ \ \ \ \ {\isacharbar}{\kern0pt}\ while\ \ {\isachardoublequoteopen}E\isactrlsub b\isactrlsub o\isactrlsub o\isactrlsub l{\isachardoublequoteclose}\ com\ \ \ \ \ {\isacharparenleft}{\kern0pt}{\isachardoublequoteopen}WHILE\ {\isacharparenleft}{\kern0pt}{\isacharunderscore}{\kern0pt}{\isacharparenright}{\kern0pt}{\isacharslash}{\kern0pt}\ DO\ {\isacharparenleft}{\kern0pt}{\isacharunderscore}{\kern0pt}{\isacharparenright}{\kern0pt}{\isachardoublequoteclose}\ {\isacharbrackleft}{\kern0pt}{\isadigit{0}}{\isacharcomma}{\kern0pt}{\isadigit{8}}{\isadigit{0}}{\isacharbrackright}{\kern0pt}{\isadigit{8}}{\isadigit{0}}{\isacharparenright}{\kern0pt}\isanewline
\ \ \ \ \ \ \ \ \ \ \ \ \ {\isacharbar}{\kern0pt}\ call\ \ \ \ F\isanewline
\isanewline
\ \ \ \ \ \ \ \ \ \ \ \ \ {\isacharbar}{\kern0pt}\ STOP\isanewline
\ \ \ \ \ \ \ \ \ \ \ \ \ {\isacharbar}{\kern0pt}\ lock\ \ \ \ MV\isanewline
\ \ \ \ \ \ \ \ \ \ \ \ \ {\isacharbar}{\kern0pt}\ unlock\ \ MV\ \ \ \ \ \ \ \ \ \ \ \ \isanewline
\ \ \ \ \ \ \ \ \ \ \ \ \ {\isacharbar}{\kern0pt}\ send\ \ \ \ E\isactrlsub a\isactrlsub r\isactrlsub i\isactrlsub t\isactrlsub h\ SV\ \ \ \ \ \ {\isacharparenleft}{\kern0pt}{\isachardoublequoteopen}STORE\ {\isacharunderscore}{\kern0pt}\ TO\ \ {\isacharunderscore}{\kern0pt}{\isachardoublequoteclose}\ {\isacharbrackleft}{\kern0pt}{\isadigit{9}}{\isadigit{0}}{\isacharcomma}{\kern0pt}{\isadigit{9}}{\isadigit{0}}{\isacharbrackright}{\kern0pt}{\isadigit{8}}{\isadigit{0}}{\isacharparenright}{\kern0pt}\ \ \ \ \ \ \ \ \ \ \ \ \ \ \isanewline
\ \ \ \ \ \ \ \ \ \ \ \ \ {\isacharbar}{\kern0pt}\ rec\ \ \ \ \ V\ SV\ \ \ \ \ \ \ \ \ \ {\isacharparenleft}{\kern0pt}{\isachardoublequoteopen}LOAD\ {\isacharunderscore}{\kern0pt}\ FROM\ {\isacharunderscore}{\kern0pt}{\isachardoublequoteclose}\ {\isacharbrackleft}{\kern0pt}{\isadigit{9}}{\isadigit{0}}{\isacharcomma}{\kern0pt}{\isadigit{9}}{\isadigit{0}}{\isacharbrackright}{\kern0pt}{\isadigit{8}}{\isadigit{0}}{\isacharparenright}{\kern0pt}%
\begin{isamarkuptext}%
Note that it is straight forward to formulate the assert primitive in \isa{IMP\isactrlsub c\isactrlsub o\isactrlsub n\isactrlsub c\isactrlsub u\isactrlsub r}:%
\end{isamarkuptext}\isamarkuptrue%
\isakeywordONE{definition}\isamarkupfalse%
\ assert\ \isakeywordTWO{where}\ {\isacartoucheopen}assert\ E\ {\isasymequiv}\ IF\ E\ THEN\ SKIP\ ELSE\ STOP\ {\isacartoucheclose}%
\begin{isamarkuptext}%
... which reduces the checking of assertions to the proof of deedlock-freeness.%
\end{isamarkuptext}\isamarkuptrue%
%
\begin{isamarkuptext}%
Two (non-sensical) example threads:%
\end{isamarkuptext}\isamarkuptrue%
\isakeywordONE{definition}\isamarkupfalse%
\ Thread{\isadigit{1}}\ \isakeywordTWO{where}\isanewline
\ \ \ \ \ \ {\isacartoucheopen}\ Thread{\isadigit{1}}\ {\isasymequiv}\ SKIP\ {\isacharsemicolon}{\kern0pt}\ \isanewline
\ \ \ \ \ \ \ \ \ \ \ \ \ \ \ \ \ \ {\isacharprime}{\kern0pt}{\isacharprime}{\kern0pt}a{\isacharprime}{\kern0pt}{\isacharprime}{\kern0pt}\ {\isacharcolon}{\kern0pt}{\isacharequal}{\kern0pt}\ \ {\isacharparenleft}{\kern0pt}{\isasymlambda}{\isasymsigma}{\isachardot}{\kern0pt}\ {\isadigit{1}}\ {\isacharplus}{\kern0pt}\ {\isadigit{2}}{\isacharparenright}{\kern0pt}{\isacharsemicolon}{\kern0pt}\ \isanewline
\ \ \ \ \ \ \ \ \ \ \ \ \ \ \ \ \ \ IF\ {\isacharparenleft}{\kern0pt}{\isasymlambda}{\isasymsigma}{\isachardot}{\kern0pt}\ {\isasymsigma}\ {\isacharprime}{\kern0pt}{\isacharprime}{\kern0pt}d{\isacharprime}{\kern0pt}{\isacharprime}{\kern0pt}\ {\isachargreater}{\kern0pt}\ {\isadigit{3}}{\isacharparenright}{\kern0pt}\ THEN\ lock\ {\isadigit{4}}\ ELSE\ lock\ {\isadigit{3}}{\isacharsemicolon}{\kern0pt}\ \ \isanewline
\ \ \ \ \ \ \ \ \ \ \ \ \ \ \ \ \ \ WHILE\ {\isacharparenleft}{\kern0pt}{\isasymlambda}{\isasymsigma}{\isachardot}{\kern0pt}\ {\isasymsigma}\ {\isacharprime}{\kern0pt}{\isacharprime}{\kern0pt}e{\isacharprime}{\kern0pt}{\isacharprime}{\kern0pt}\ {\isachargreater}{\kern0pt}\ {\isadigit{3}}{\isacharparenright}{\kern0pt}\ DO\ lock\ {\isadigit{4}}\ {\isacharsemicolon}{\kern0pt}\ \ \isanewline
\ \ \ \ \ \ \ \ \ \ \ \ \ \ \ \ \ \ STORE\ {\isacharparenleft}{\kern0pt}{\isasymlambda}{\isasymsigma}{\isachardot}{\kern0pt}\ {\isasymsigma}\ {\isacharprime}{\kern0pt}{\isacharprime}{\kern0pt}d{\isacharprime}{\kern0pt}{\isacharprime}{\kern0pt}\ {\isacharplus}{\kern0pt}\ {\isadigit{1}}{\isacharparenright}{\kern0pt}\ TO\ {\isacharprime}{\kern0pt}{\isacharprime}{\kern0pt}b{\isacharprime}{\kern0pt}{\isacharprime}{\kern0pt}{\isacharsemicolon}{\kern0pt}\ \isanewline
\ \ \ \ \ \ \ \ \ \ \ \ \ \ \ \ \ \ LOAD\ {\isacharprime}{\kern0pt}{\isacharprime}{\kern0pt}a{\isacharprime}{\kern0pt}{\isacharprime}{\kern0pt}\ FROM\ \ {\isacharprime}{\kern0pt}{\isacharprime}{\kern0pt}b{\isacharprime}{\kern0pt}{\isacharprime}{\kern0pt}{\isacharsemicolon}{\kern0pt}\ \isanewline
\ \ \ \ \ \ \ \ \ \ \ \ \ \ \ \ \ \ unlock\ {\isadigit{3}}\isanewline
\ \ \ \ \ \ {\isacartoucheclose}\isanewline
\isanewline
\isakeywordONE{definition}\isamarkupfalse%
\ Thread{\isadigit{2}}\ \isakeywordTWO{where}\isanewline
\ \ \ \ \ \ {\isacartoucheopen}\ Thread{\isadigit{2}}\ {\isasymequiv}\ SKIP\ {\isacharsemicolon}{\kern0pt}\ \isanewline
\ \ \ \ \ \ \ \ \ \ \ \ \ \ \ \ \ \ {\isacharprime}{\kern0pt}{\isacharprime}{\kern0pt}a{\isacharprime}{\kern0pt}{\isacharprime}{\kern0pt}\ {\isacharcolon}{\kern0pt}{\isacharequal}{\kern0pt}\ \ {\isacharparenleft}{\kern0pt}{\isasymlambda}{\isasymsigma}{\isachardot}{\kern0pt}\ {\isadigit{1}}\ {\isacharplus}{\kern0pt}\ {\isadigit{2}}{\isacharparenright}{\kern0pt}{\isacharsemicolon}{\kern0pt}\ \isanewline
\ \ \ \ \ \ \ \ \ \ \ \ \ \ \ \ \ \ IF\ {\isacharparenleft}{\kern0pt}{\isasymlambda}{\isasymsigma}{\isachardot}{\kern0pt}\ {\isasymsigma}\ {\isacharprime}{\kern0pt}{\isacharprime}{\kern0pt}d{\isacharprime}{\kern0pt}{\isacharprime}{\kern0pt}\ {\isachargreater}{\kern0pt}\ {\isadigit{3}}{\isacharparenright}{\kern0pt}\ THEN\ lock\ {\isadigit{4}}\ ELSE\ lock\ {\isadigit{3}}{\isacharsemicolon}{\kern0pt}\ \ \isanewline
\ \ \ \ \ \ \ \ \ \ \ \ \ \ \ \ \ \ WHILE\ {\isacharparenleft}{\kern0pt}{\isasymlambda}{\isasymsigma}{\isachardot}{\kern0pt}\ True{\isacharparenright}{\kern0pt}\ DO\ lock\ {\isadigit{4}}\ {\isacharsemicolon}{\kern0pt}\ \ \isanewline
\ \ \ \ \ \ \ \ \ \ \ \ \ \ \ \ \ \ STORE\ {\isacharparenleft}{\kern0pt}{\isasymlambda}{\isasymsigma}{\isachardot}{\kern0pt}\ {\isasymsigma}\ {\isacharprime}{\kern0pt}{\isacharprime}{\kern0pt}d{\isacharprime}{\kern0pt}{\isacharprime}{\kern0pt}\ {\isacharplus}{\kern0pt}\ {\isadigit{1}}{\isacharparenright}{\kern0pt}\ TO\ {\isacharprime}{\kern0pt}{\isacharprime}{\kern0pt}b{\isacharprime}{\kern0pt}{\isacharprime}{\kern0pt}{\isacharsemicolon}{\kern0pt}\ \isanewline
\ \ \ \ \ \ \ \ \ \ \ \ \ \ \ \ \ \ LOAD\ {\isacharprime}{\kern0pt}{\isacharprime}{\kern0pt}a{\isacharprime}{\kern0pt}{\isacharprime}{\kern0pt}\ FROM\ \ {\isacharprime}{\kern0pt}{\isacharprime}{\kern0pt}b{\isacharprime}{\kern0pt}{\isacharprime}{\kern0pt}{\isacharsemicolon}{\kern0pt}\ \isanewline
\ \ \ \ \ \ \ \ \ \ \ \ \ \ \ \ \ \ unlock\ {\isadigit{3}}\isanewline
\ \ \ \ \ \ {\isacartoucheclose}\isanewline
\isanewline
\isakeywordONE{thm}\isamarkupfalse%
\ Thread{\isadigit{2}}{\isacharunderscore}{\kern0pt}def\isanewline
\isakeywordONE{definition}\isamarkupfalse%
\ Thread{\isadigit{3}}\ \isakeywordTWO{where}\ \isanewline
\ \ {\isachardoublequoteopen}Thread{\isadigit{3}}\ {\isacharequal}{\kern0pt}\ SKIP\ {\isacharsemicolon}{\kern0pt}\isanewline
\ \ \ \ \ \ \ \ \ \ \ \ \ {\isacharparenleft}{\kern0pt}\ {\isacharprime}{\kern0pt}{\isacharprime}{\kern0pt}var{\isacharunderscore}{\kern0pt}local{\isacharprime}{\kern0pt}{\isacharprime}{\kern0pt}\ {\isacharcolon}{\kern0pt}{\isacharequal}{\kern0pt}\ {\isacharparenleft}{\kern0pt}{\isasymlambda}{\isasymsigma}{\isachardot}{\kern0pt}\ {\isadigit{4}}\ {\isacharplus}{\kern0pt}\ {\isadigit{4}}{\isadigit{2}}{\isacharparenright}{\kern0pt}\ {\isacharsemicolon}{\kern0pt}\isanewline
\ \ \ \ \ \ \ \ \ \ \ \ \ {\isacharparenleft}{\kern0pt}\ {\isacharprime}{\kern0pt}{\isacharprime}{\kern0pt}test{\isacharprime}{\kern0pt}{\isacharprime}{\kern0pt}\ {\isacharcolon}{\kern0pt}{\isacharequal}{\kern0pt}\ {\isacharparenleft}{\kern0pt}{\isasymlambda}{\isasymsigma}{\isachardot}{\kern0pt}\ {\isadigit{2}}\ {\isacharasterisk}{\kern0pt}\ {\isasymsigma}\ {\isacharprime}{\kern0pt}{\isacharprime}{\kern0pt}var{\isacharunderscore}{\kern0pt}local{\isacharprime}{\kern0pt}{\isacharprime}{\kern0pt}{\isacharparenright}{\kern0pt}\ {\isacharsemicolon}{\kern0pt}\isanewline
\ \ \ \ \ \ \ \ \ \ \ \ \ {\isacharparenleft}{\kern0pt}\ {\isacharprime}{\kern0pt}{\isacharprime}{\kern0pt}test{\isacharprime}{\kern0pt}{\isacharprime}{\kern0pt}\ {\isacharcolon}{\kern0pt}{\isacharequal}{\kern0pt}\ {\isacharparenleft}{\kern0pt}{\isasymlambda}{\isasymsigma}{\isachardot}{\kern0pt}\ {\isasymsigma}\ {\isacharprime}{\kern0pt}{\isacharprime}{\kern0pt}test{\isacharprime}{\kern0pt}{\isacharprime}{\kern0pt}\ {\isacharplus}{\kern0pt}\ {\isadigit{3}}{\isacharparenright}{\kern0pt}\ {\isacharsemicolon}{\kern0pt}\isanewline
\ \ \ \ \ \ \ \ \ \ \ \ \ \ IF\ {\isacharparenleft}{\kern0pt}{\isasymlambda}{\isasymsigma}{\isachardot}{\kern0pt}\ {\isasymsigma}\ {\isacharprime}{\kern0pt}{\isacharprime}{\kern0pt}x{\isacharprime}{\kern0pt}{\isacharprime}{\kern0pt}{\isachargreater}{\kern0pt}{\isadigit{5}}{\isacharparenright}{\kern0pt}\isanewline
\ \ \ \ \ \ \ \ \ \ \ \ \ \ THEN\ WHILE\ {\isacharparenleft}{\kern0pt}{\isasymlambda}{\isasymsigma}{\isachardot}{\kern0pt}\ {\isasymsigma}\ {\isacharprime}{\kern0pt}{\isacharprime}{\kern0pt}x{\isacharprime}{\kern0pt}{\isacharprime}{\kern0pt}\ {\isachargreater}{\kern0pt}{\isadigit{5}}{\isacharparenright}{\kern0pt}\ DO\ \ {\isacharprime}{\kern0pt}{\isacharprime}{\kern0pt}test{\isacharprime}{\kern0pt}{\isacharprime}{\kern0pt}\ {\isacharcolon}{\kern0pt}{\isacharequal}{\kern0pt}\ {\isacharparenleft}{\kern0pt}{\isasymlambda}{\isasymsigma}{\isachardot}{\kern0pt}\ {\isasymsigma}\ {\isacharprime}{\kern0pt}{\isacharprime}{\kern0pt}test{\isacharprime}{\kern0pt}{\isacharprime}{\kern0pt}\ {\isacharminus}{\kern0pt}\ {\isadigit{3}}{\isacharparenright}{\kern0pt}\isanewline
\ \ \ \ \ \ \ \ \ \ \ \ \ \ ELSE\ \ {\isacharprime}{\kern0pt}{\isacharprime}{\kern0pt}test{\isacharprime}{\kern0pt}{\isacharprime}{\kern0pt}\ {\isacharcolon}{\kern0pt}{\isacharequal}{\kern0pt}\ {\isacharparenleft}{\kern0pt}{\isasymlambda}{\isasymsigma}{\isachardot}{\kern0pt}\ {\isasymsigma}\ {\isacharprime}{\kern0pt}{\isacharprime}{\kern0pt}test{\isacharprime}{\kern0pt}{\isacharprime}{\kern0pt}\ {\isacharplus}{\kern0pt}\ {\isadigit{3}}{\isacharparenright}{\kern0pt}{\isacharparenright}{\kern0pt}{\isacharparenright}{\kern0pt}{\isacharparenright}{\kern0pt}{\isachardoublequoteclose}\isanewline
%
\isadelimML
\isanewline
%
\endisadelimML
%
\isatagML
\isakeywordONE{ML}\isamarkupfalse%
{\isacartoucheopen}val\ Thread{\isadigit{2}}{\isacharunderscore}{\kern0pt}term\ {\isacharequal}{\kern0pt}\ \isactrlterm {\isasymopen}Thread{\isadigit{2}}{\isasymclose}{\isacharsemicolon}{\kern0pt}\isanewline
\ \ \ val\ {\isacharunderscore}{\kern0pt}\ {\isachardollar}{\kern0pt}\ {\isacharunderscore}{\kern0pt}\ {\isachardollar}{\kern0pt}\ Thread{\isadigit{2}}{\isacharunderscore}{\kern0pt}rhs\ {\isacharequal}{\kern0pt}\ Thm{\isachardot}{\kern0pt}concl{\isacharunderscore}{\kern0pt}of\ {\isacharat}{\kern0pt}{\isacharbraceleft}{\kern0pt}thm\ Thread{\isadigit{2}}{\isacharunderscore}{\kern0pt}def{\isacharbraceright}{\kern0pt}{\isacartoucheclose}%
\endisatagML
{\isafoldML}%
%
\isadelimML
%
\endisadelimML
%
\isadelimdocument
%
\endisadelimdocument
%
\isatagdocument
%
\isamarkupsection{Events, States and Semaphores ("Bricks").%
}
\isamarkuptrue%
%
\endisatagdocument
{\isafolddocument}%
%
\isadelimdocument
%
\endisadelimdocument
%
\begin{isamarkuptext}%
The common set of events in a thread-system%
\end{isamarkuptext}\isamarkuptrue%
\isakeywordONE{datatype}\isamarkupfalse%
\ evs\ {\isacharequal}{\kern0pt}\ lock\ int\ {\isacharbar}{\kern0pt}\ unlock\ int\ {\isacharbar}{\kern0pt}\ read\isactrlsub g\isactrlsub l\isactrlsub o\isactrlsub b\isactrlsub a\isactrlsub l\ SV\ D\ {\isacharbar}{\kern0pt}\ upd\isactrlsub g\isactrlsub l\isactrlsub o\isactrlsub b\isactrlsub a\isactrlsub l\ \ SV\ D\ \isanewline
\ \ \ \ \ \ \ \ \ \ \ \ \ \ \ \ \ \ \ \ \ \ \ \ \ \ \ \ \ \ \ \ \ \ \ \ \ {\isacharbar}{\kern0pt}\ read\isactrlsub l\isactrlsub o\isactrlsub c\isactrlsub a\isactrlsub l\ SV\ D\ {\isacharbar}{\kern0pt}\ update\isactrlsub l\isactrlsub o\isactrlsub c\isactrlsub a\isactrlsub l\ SV\ D\ \isanewline
\isanewline
\isakeywordONE{definition}\isamarkupfalse%
\ LOCKS\ \isakeywordTWO{where}\ {\isacartoucheopen}LOCKS\ {\isasymequiv}\ range\ lock\ {\isasymunion}\ range\ unlock{\isacartoucheclose}\isanewline
\isakeywordONE{definition}\isamarkupfalse%
\ GVARS\ \isakeywordTWO{where}\ {\isacartoucheopen}GVARS\ {\isasymequiv}\ {\isacharbraceleft}{\kern0pt}e{\isachardot}{\kern0pt}\ {\isasymexists}x{\isachardot}{\kern0pt}\ {\isasymexists}y{\isachardot}{\kern0pt}\ e\ {\isacharequal}{\kern0pt}\ read\isactrlsub g\isactrlsub l\isactrlsub o\isactrlsub b\isactrlsub a\isactrlsub l\ x\ y\ {\isacharbraceright}{\kern0pt}\ {\isasymunion}\ \isanewline
\ \ \ \ \ \ \ \ \ \ \ \ \ \ \ \ \ \ \ \ \ \ \ \ \ \ \ \ \ \ \ \ \ {\isacharbraceleft}{\kern0pt}e{\isachardot}{\kern0pt}\ {\isasymexists}x{\isachardot}{\kern0pt}\ {\isasymexists}y{\isachardot}{\kern0pt}\ e\ {\isacharequal}{\kern0pt}\ upd\isactrlsub g\isactrlsub l\isactrlsub o\isactrlsub b\isactrlsub a\isactrlsub l\ x\ y\ {\isacharbraceright}{\kern0pt}{\isacartoucheclose}\isanewline
\isanewline
\isakeywordONE{lemma}\isamarkupfalse%
\ L{\isadigit{1}}\ {\isacharbrackleft}{\kern0pt}simp{\isacharbrackright}{\kern0pt}{\isacharcolon}{\kern0pt}\ {\isacartoucheopen}lock\ n\ {\isasymnotin}\ GVARS{\isacartoucheclose}%
\isadelimproof
\ \ \ %
\endisadelimproof
%
\isatagproof
\isakeywordONE{unfolding}\isamarkupfalse%
\ GVARS{\isacharunderscore}{\kern0pt}def\ \isakeywordONE{by}\isamarkupfalse%
\ auto%
\endisatagproof
{\isafoldproof}%
%
\isadelimproof
%
\endisadelimproof
\isanewline
\isakeywordONE{lemma}\isamarkupfalse%
\ L{\isadigit{2}}\ {\isacharbrackleft}{\kern0pt}simp{\isacharbrackright}{\kern0pt}{\isacharcolon}{\kern0pt}\ {\isacartoucheopen}lock\ n\ {\isasymin}\ LOCKS{\isacartoucheclose}%
\isadelimproof
\ \ \ %
\endisadelimproof
%
\isatagproof
\isakeywordONE{unfolding}\isamarkupfalse%
\ LOCKS{\isacharunderscore}{\kern0pt}def\ \isakeywordONE{by}\isamarkupfalse%
\ auto%
\endisatagproof
{\isafoldproof}%
%
\isadelimproof
%
\endisadelimproof
\isanewline
\isakeywordONE{lemma}\isamarkupfalse%
\ L{\isadigit{3}}\ {\isacharbrackleft}{\kern0pt}simp{\isacharbrackright}{\kern0pt}{\isacharcolon}{\kern0pt}\ {\isacartoucheopen}unlock\ n\ {\isasymnotin}\ GVARS{\isacartoucheclose}%
\isadelimproof
\ %
\endisadelimproof
%
\isatagproof
\isakeywordONE{unfolding}\isamarkupfalse%
\ GVARS{\isacharunderscore}{\kern0pt}def\ \isakeywordONE{by}\isamarkupfalse%
\ auto%
\endisatagproof
{\isafoldproof}%
%
\isadelimproof
%
\endisadelimproof
\isanewline
\isakeywordONE{lemma}\isamarkupfalse%
\ L{\isadigit{4}}\ {\isacharbrackleft}{\kern0pt}simp{\isacharbrackright}{\kern0pt}{\isacharcolon}{\kern0pt}\ {\isacartoucheopen}unlock\ n\ {\isasymin}\ LOCKS{\isacartoucheclose}%
\isadelimproof
\ %
\endisadelimproof
%
\isatagproof
\isakeywordONE{unfolding}\isamarkupfalse%
\ LOCKS{\isacharunderscore}{\kern0pt}def\ \isakeywordONE{by}\isamarkupfalse%
\ auto%
\endisatagproof
{\isafoldproof}%
%
\isadelimproof
%
\endisadelimproof
\isanewline
\isanewline
\isakeywordONE{lemma}\isamarkupfalse%
\ L{\isadigit{5}}\ {\isacharbrackleft}{\kern0pt}simp{\isacharbrackright}{\kern0pt}{\isacharcolon}{\kern0pt}\ {\isacartoucheopen}read\isactrlsub g\isactrlsub l\isactrlsub o\isactrlsub b\isactrlsub a\isactrlsub l\ n\ m\ {\isasymnotin}\ LOCKS{\isacartoucheclose}%
\isadelimproof
\ \ \ %
\endisadelimproof
%
\isatagproof
\isakeywordONE{unfolding}\isamarkupfalse%
\ LOCKS{\isacharunderscore}{\kern0pt}def\ GVARS{\isacharunderscore}{\kern0pt}def\ \isakeywordONE{by}\isamarkupfalse%
\ auto%
\endisatagproof
{\isafoldproof}%
%
\isadelimproof
%
\endisadelimproof
\isanewline
\isakeywordONE{lemma}\isamarkupfalse%
\ L{\isadigit{6}}\ {\isacharbrackleft}{\kern0pt}simp{\isacharbrackright}{\kern0pt}{\isacharcolon}{\kern0pt}\ {\isacartoucheopen}read\isactrlsub g\isactrlsub l\isactrlsub o\isactrlsub b\isactrlsub a\isactrlsub l\ n\ m\ {\isasymin}\ GVARS{\isacartoucheclose}%
\isadelimproof
\ \ \ %
\endisadelimproof
%
\isatagproof
\isakeywordONE{unfolding}\isamarkupfalse%
\ LOCKS{\isacharunderscore}{\kern0pt}def\ GVARS{\isacharunderscore}{\kern0pt}def\ \isakeywordONE{by}\isamarkupfalse%
\ auto%
\endisatagproof
{\isafoldproof}%
%
\isadelimproof
%
\endisadelimproof
\isanewline
\isakeywordONE{lemma}\isamarkupfalse%
\ L{\isadigit{7}}\ {\isacharbrackleft}{\kern0pt}simp{\isacharbrackright}{\kern0pt}{\isacharcolon}{\kern0pt}\ {\isacartoucheopen}upd\isactrlsub g\isactrlsub l\isactrlsub o\isactrlsub b\isactrlsub a\isactrlsub l\ n\ m\ {\isasymnotin}\ LOCKS{\isacartoucheclose}%
\isadelimproof
\ \ \ \ %
\endisadelimproof
%
\isatagproof
\isakeywordONE{unfolding}\isamarkupfalse%
\ LOCKS{\isacharunderscore}{\kern0pt}def\ GVARS{\isacharunderscore}{\kern0pt}def\ \isakeywordONE{by}\isamarkupfalse%
\ auto%
\endisatagproof
{\isafoldproof}%
%
\isadelimproof
%
\endisadelimproof
\isanewline
\isakeywordONE{lemma}\isamarkupfalse%
\ L{\isadigit{8}}\ {\isacharbrackleft}{\kern0pt}simp{\isacharbrackright}{\kern0pt}{\isacharcolon}{\kern0pt}\ {\isacartoucheopen}upd\isactrlsub g\isactrlsub l\isactrlsub o\isactrlsub b\isactrlsub a\isactrlsub l\ n\ m\ {\isasymin}\ GVARS{\isacartoucheclose}%
\isadelimproof
\ \ \ \ %
\endisadelimproof
%
\isatagproof
\isakeywordONE{unfolding}\isamarkupfalse%
\ LOCKS{\isacharunderscore}{\kern0pt}def\ GVARS{\isacharunderscore}{\kern0pt}def\ \isakeywordONE{by}\isamarkupfalse%
\ auto%
\endisatagproof
{\isafoldproof}%
%
\isadelimproof
%
\endisadelimproof
\isanewline
\isanewline
\isakeywordONE{lemma}\isamarkupfalse%
\ L{\isadigit{9}}\ {\isacharbrackleft}{\kern0pt}simp{\isacharbrackright}{\kern0pt}{\isacharcolon}{\kern0pt}\ {\isacartoucheopen}read\isactrlsub l\isactrlsub o\isactrlsub c\isactrlsub a\isactrlsub l\ n\ m\ {\isasymnotin}\ LOCKS{\isacartoucheclose}%
\isadelimproof
\ \ \ \ %
\endisadelimproof
%
\isatagproof
\isakeywordONE{unfolding}\isamarkupfalse%
\ LOCKS{\isacharunderscore}{\kern0pt}def\ GVARS{\isacharunderscore}{\kern0pt}def\ \isakeywordONE{by}\isamarkupfalse%
\ auto%
\endisatagproof
{\isafoldproof}%
%
\isadelimproof
%
\endisadelimproof
\isanewline
\isakeywordONE{lemma}\isamarkupfalse%
\ L{\isadigit{1}}{\isadigit{0}}{\isacharbrackleft}{\kern0pt}simp{\isacharbrackright}{\kern0pt}{\isacharcolon}{\kern0pt}\ {\isacartoucheopen}read\isactrlsub l\isactrlsub o\isactrlsub c\isactrlsub a\isactrlsub l\ n\ m\ {\isasymnotin}\ GVARS{\isacartoucheclose}%
\isadelimproof
\ \ \ \ %
\endisadelimproof
%
\isatagproof
\isakeywordONE{unfolding}\isamarkupfalse%
\ LOCKS{\isacharunderscore}{\kern0pt}def\ GVARS{\isacharunderscore}{\kern0pt}def\ \isakeywordONE{by}\isamarkupfalse%
\ auto%
\endisatagproof
{\isafoldproof}%
%
\isadelimproof
%
\endisadelimproof
\isanewline
\isakeywordONE{lemma}\isamarkupfalse%
\ L{\isadigit{1}}{\isadigit{1}}{\isacharbrackleft}{\kern0pt}simp{\isacharbrackright}{\kern0pt}{\isacharcolon}{\kern0pt}\ {\isacartoucheopen}update\isactrlsub l\isactrlsub o\isactrlsub c\isactrlsub a\isactrlsub l\ n\ m\ {\isasymnotin}\ LOCKS{\isacartoucheclose}%
\isadelimproof
\ \ %
\endisadelimproof
%
\isatagproof
\isakeywordONE{unfolding}\isamarkupfalse%
\ LOCKS{\isacharunderscore}{\kern0pt}def\ GVARS{\isacharunderscore}{\kern0pt}def\ \isakeywordONE{by}\isamarkupfalse%
\ auto%
\endisatagproof
{\isafoldproof}%
%
\isadelimproof
%
\endisadelimproof
\isanewline
\isakeywordONE{lemma}\isamarkupfalse%
\ L{\isadigit{1}}{\isadigit{2}}{\isacharbrackleft}{\kern0pt}simp{\isacharbrackright}{\kern0pt}{\isacharcolon}{\kern0pt}\ {\isacartoucheopen}update\isactrlsub l\isactrlsub o\isactrlsub c\isactrlsub a\isactrlsub l\ n\ m\ {\isasymnotin}\ GVARS{\isacartoucheclose}%
\isadelimproof
\ \ %
\endisadelimproof
%
\isatagproof
\isakeywordONE{unfolding}\isamarkupfalse%
\ LOCKS{\isacharunderscore}{\kern0pt}def\ GVARS{\isacharunderscore}{\kern0pt}def\ \isakeywordONE{by}\isamarkupfalse%
\ auto%
\endisatagproof
{\isafoldproof}%
%
\isadelimproof
%
\endisadelimproof
%
\isadelimdocument
%
\endisadelimdocument
%
\isatagdocument
%
\isamarkupsubsection{A simple Model of a Semaphore Family%
}
\isamarkuptrue%
%
\endisatagdocument
{\isafolddocument}%
%
\isadelimdocument
%
\endisadelimdocument
\isakeywordONE{definition}\isamarkupfalse%
\ semaphore\ {\isacharcolon}{\kern0pt}{\isacharcolon}{\kern0pt}\ {\isacartoucheopen}int\ {\isasymRightarrow}\ evs\ process{\isacartoucheclose}\isanewline
\ \ \isakeywordTWO{where}\ \ \ {\isacartoucheopen}semaphore\ {\isasymequiv}\ {\isasymmu}\ X{\isachardot}{\kern0pt}\ {\isacharparenleft}{\kern0pt}{\isasymlambda}n{\isachardot}{\kern0pt}\ lock\ n\ {\isasymrightarrow}\ unlock\ n\ {\isasymrightarrow}\ X\ n{\isacharparenright}{\kern0pt}{\isacartoucheclose}\isanewline
\isanewline
\isakeywordONE{lemma}\isamarkupfalse%
\ sema{\isacharunderscore}{\kern0pt}rec\ {\isacharcolon}{\kern0pt}\ {\isacartoucheopen}semaphore\ n\ {\isacharequal}{\kern0pt}\ lock\ n\ {\isasymrightarrow}\ unlock\ n\ {\isasymrightarrow}\ semaphore\ n{\isacartoucheclose}\isanewline
%
\isadelimproof
\ \ %
\endisadelimproof
%
\isatagproof
\isakeywordONE{by}\isamarkupfalse%
\ {\isacharparenleft}{\kern0pt}subst\ cont{\isacharunderscore}{\kern0pt}process{\isacharunderscore}{\kern0pt}rec{\isacharbrackleft}{\kern0pt}OF\ semaphore{\isacharunderscore}{\kern0pt}def{\isacharbrackleft}{\kern0pt}THEN\ meta{\isacharunderscore}{\kern0pt}eq{\isacharunderscore}{\kern0pt}to{\isacharunderscore}{\kern0pt}obj{\isacharunderscore}{\kern0pt}eq{\isacharbrackright}{\kern0pt}{\isacharbrackright}{\kern0pt}{\isacharparenright}{\kern0pt}\ simp{\isacharunderscore}{\kern0pt}all%
\endisatagproof
{\isafoldproof}%
%
\isadelimproof
%
\endisadelimproof
%
\isadelimdocument
%
\endisadelimdocument
%
\isatagdocument
%
\isamarkupsubsection{A simple Model of a Global Variable Family%
}
\isamarkuptrue%
%
\endisatagdocument
{\isafolddocument}%
%
\isadelimdocument
%
\endisadelimdocument
\isakeywordONE{definition}\isamarkupfalse%
\ global{\isacharunderscore}{\kern0pt}vars\ {\isacharcolon}{\kern0pt}{\isacharcolon}{\kern0pt}\ {\isacartoucheopen}{\isasymsigma}\ {\isasymRightarrow}\ SV\ {\isasymRightarrow}\ evs\ process{\isacartoucheclose}\ \isanewline
\ \ \isakeywordTWO{where}\ \ \ {\isacartoucheopen}global{\isacharunderscore}{\kern0pt}vars\ {\isasymequiv}\ {\isasymmu}\ X{\isachardot}{\kern0pt}\ \ \ {\isacharparenleft}{\kern0pt}{\isasymlambda}\ {\isasymsigma}{\isachardot}{\kern0pt}\ {\isacharparenleft}{\kern0pt}{\isasymlambda}\ id{\isachardot}{\kern0pt}\ \ {\isacharparenleft}{\kern0pt}{\isacharparenleft}{\kern0pt}read\isactrlsub g\isactrlsub l\isactrlsub o\isactrlsub b\isactrlsub a\isactrlsub l\ id{\isacharparenright}{\kern0pt}\isactrlbold {\isacharbang}{\kern0pt}{\isacharparenleft}{\kern0pt}{\isasymsigma}\ id{\isacharparenright}{\kern0pt}\ {\isasymrightarrow}\ X\ {\isasymsigma}\ id\ {\isacharparenright}{\kern0pt}\ \isanewline
\ \ \ \ \ \ \ \ \ \ \ \ \ \ \ \ \ \ \ \ \ \ \ \ \ \ \ \ \ \ \ \ \ \ \ \ \ \ \ \ \ \ \ \ {\isasymbox}\ {\isacharparenleft}{\kern0pt}{\isacharparenleft}{\kern0pt}upd\isactrlsub g\isactrlsub l\isactrlsub o\isactrlsub b\isactrlsub a\isactrlsub l\ id{\isacharparenright}{\kern0pt}\isactrlbold {\isacharquery}{\kern0pt}v\ {\isasymrightarrow}\ X\ {\isacharparenleft}{\kern0pt}{\isasymsigma}{\isacharparenleft}{\kern0pt}id\ {\isacharcolon}{\kern0pt}{\isacharequal}{\kern0pt}\ v{\isacharparenright}{\kern0pt}{\isacharparenright}{\kern0pt}\ id\ {\isacharparenright}{\kern0pt}{\isacharparenright}{\kern0pt}{\isacharparenright}{\kern0pt}{\isacartoucheclose}\isanewline
\isanewline
\isanewline
\isakeywordONE{lemma}\isamarkupfalse%
\ global{\isacharunderscore}{\kern0pt}vars{\isacharunderscore}{\kern0pt}rec\ {\isacharcolon}{\kern0pt}\ {\isacartoucheopen}global{\isacharunderscore}{\kern0pt}vars\ {\isasymsigma}\ n\ {\isacharequal}{\kern0pt}\ \ \ \ {\isacharparenleft}{\kern0pt}{\isacharparenleft}{\kern0pt}read\isactrlsub g\isactrlsub l\isactrlsub o\isactrlsub b\isactrlsub a\isactrlsub l\ n{\isacharparenright}{\kern0pt}\isactrlbold {\isacharbang}{\kern0pt}{\isacharparenleft}{\kern0pt}{\isasymsigma}\ n{\isacharparenright}{\kern0pt}\ {\isasymrightarrow}\ global{\isacharunderscore}{\kern0pt}vars\ {\isasymsigma}\ n{\isacharparenright}{\kern0pt}\ \isanewline
\ \ \ \ \ \ \ \ \ \ \ \ \ \ \ \ \ \ \ \ \ \ \ \ \ \ \ \ \ \ \ \ \ \ \ \ \ \ \ \ \ \ \ {\isasymbox}\ {\isacharparenleft}{\kern0pt}{\isacharparenleft}{\kern0pt}upd\isactrlsub g\isactrlsub l\isactrlsub o\isactrlsub b\isactrlsub a\isactrlsub l\ n{\isacharparenright}{\kern0pt}\isactrlbold {\isacharquery}{\kern0pt}v\ {\isasymrightarrow}\ global{\isacharunderscore}{\kern0pt}vars\ {\isacharparenleft}{\kern0pt}{\isasymsigma}{\isacharparenleft}{\kern0pt}n\ {\isacharcolon}{\kern0pt}{\isacharequal}{\kern0pt}\ v{\isacharparenright}{\kern0pt}{\isacharparenright}{\kern0pt}\ n{\isacharparenright}{\kern0pt}{\isacartoucheclose}\ \isanewline
%
\isadelimproof
\ \ %
\endisadelimproof
%
\isatagproof
\isakeywordONE{by}\isamarkupfalse%
\ {\isacharparenleft}{\kern0pt}subst\ cont{\isacharunderscore}{\kern0pt}process{\isacharunderscore}{\kern0pt}rec{\isacharbrackleft}{\kern0pt}OF\ global{\isacharunderscore}{\kern0pt}vars{\isacharunderscore}{\kern0pt}def{\isacharbrackleft}{\kern0pt}THEN\ meta{\isacharunderscore}{\kern0pt}eq{\isacharunderscore}{\kern0pt}to{\isacharunderscore}{\kern0pt}obj{\isacharunderscore}{\kern0pt}eq{\isacharbrackright}{\kern0pt}{\isacharbrackright}{\kern0pt}{\isacharparenright}{\kern0pt}\ simp{\isacharunderscore}{\kern0pt}all%
\endisatagproof
{\isafoldproof}%
%
\isadelimproof
%
\endisadelimproof
%
\isadelimdocument
%
\endisadelimdocument
%
\isatagdocument
%
\isamarkupsubsection{A simple Model of a Local Variable Family (Optional)%
}
\isamarkuptrue%
%
\endisatagdocument
{\isafolddocument}%
%
\isadelimdocument
%
\endisadelimdocument
\isakeywordONE{definition}\isamarkupfalse%
\ local{\isacharunderscore}{\kern0pt}vars\ {\isacharcolon}{\kern0pt}{\isacharcolon}{\kern0pt}\ {\isachardoublequoteopen}\ {\isasymsigma}\ {\isasymRightarrow}\ SV\ {\isasymRightarrow}\ evs\ process{\isachardoublequoteclose}\ \isanewline
\ \ \isakeywordTWO{where}\ \ \ {\isacartoucheopen}local{\isacharunderscore}{\kern0pt}vars\ {\isasymequiv}\ {\isasymmu}\ X{\isachardot}{\kern0pt}\ \ \ {\isacharparenleft}{\kern0pt}{\isasymlambda}\ {\isasymsigma}{\isachardot}{\kern0pt}\ {\isacharparenleft}{\kern0pt}{\isasymlambda}\ id{\isachardot}{\kern0pt}\ \ {\isacharparenleft}{\kern0pt}{\isacharparenleft}{\kern0pt}read\isactrlsub l\isactrlsub o\isactrlsub c\isactrlsub a\isactrlsub l\ id{\isacharparenright}{\kern0pt}\isactrlbold {\isacharbang}{\kern0pt}{\isacharparenleft}{\kern0pt}{\isasymsigma}\ id{\isacharparenright}{\kern0pt}\ {\isasymrightarrow}\ X\ {\isasymsigma}\ id{\isacharparenright}{\kern0pt}\ \isanewline
\ \ \ \ \ \ \ \ \ \ \ \ \ \ \ \ \ \ \ \ \ \ \ \ \ \ \ \ \ \ \ \ \ \ \ \ \ \ \ \ \ \ \ {\isasymbox}\ {\isacharparenleft}{\kern0pt}{\isacharparenleft}{\kern0pt}update\isactrlsub l\isactrlsub o\isactrlsub c\isactrlsub a\isactrlsub l\ id{\isacharparenright}{\kern0pt}\isactrlbold {\isacharquery}{\kern0pt}v\ {\isasymrightarrow}\ X\ {\isacharparenleft}{\kern0pt}{\isasymsigma}{\isacharparenleft}{\kern0pt}id\ {\isacharcolon}{\kern0pt}{\isacharequal}{\kern0pt}\ v{\isacharparenright}{\kern0pt}{\isacharparenright}{\kern0pt}\ id\ {\isacharparenright}{\kern0pt}{\isacharparenright}{\kern0pt}{\isacharparenright}{\kern0pt}{\isacartoucheclose}\isanewline
\isanewline
\isakeywordONE{lemma}\isamarkupfalse%
\ local{\isacharunderscore}{\kern0pt}vars{\isacharunderscore}{\kern0pt}rec\ {\isacharcolon}{\kern0pt}\ {\isacartoucheopen}local{\isacharunderscore}{\kern0pt}vars\ {\isasymsigma}\ n\ {\isacharequal}{\kern0pt}\ \ \ \ \ {\isacharparenleft}{\kern0pt}{\isacharparenleft}{\kern0pt}read\isactrlsub l\isactrlsub o\isactrlsub c\isactrlsub a\isactrlsub l\ n{\isacharparenright}{\kern0pt}\isactrlbold {\isacharbang}{\kern0pt}{\isacharparenleft}{\kern0pt}{\isasymsigma}\ n{\isacharparenright}{\kern0pt}\ {\isasymrightarrow}\ local{\isacharunderscore}{\kern0pt}vars\ {\isasymsigma}\ n{\isacharparenright}{\kern0pt}\ \isanewline
\ \ \ \ \ \ \ \ \ \ \ \ \ \ \ \ \ \ \ \ \ \ \ \ \ \ \ \ \ \ \ \ \ \ \ \ \ \ \ \ \ \ \ {\isasymbox}\ {\isacharparenleft}{\kern0pt}{\isacharparenleft}{\kern0pt}update\isactrlsub l\isactrlsub o\isactrlsub c\isactrlsub a\isactrlsub l\ n{\isacharparenright}{\kern0pt}\isactrlbold {\isacharquery}{\kern0pt}v\ {\isasymrightarrow}\ local{\isacharunderscore}{\kern0pt}vars\ {\isacharparenleft}{\kern0pt}{\isasymsigma}{\isacharparenleft}{\kern0pt}n\ {\isacharcolon}{\kern0pt}{\isacharequal}{\kern0pt}\ v{\isacharparenright}{\kern0pt}{\isacharparenright}{\kern0pt}\ n{\isacharparenright}{\kern0pt}{\isacartoucheclose}\ \isanewline
%
\isadelimproof
\ \ %
\endisadelimproof
%
\isatagproof
\isakeywordONE{by}\isamarkupfalse%
\ {\isacharparenleft}{\kern0pt}subst\ cont{\isacharunderscore}{\kern0pt}process{\isacharunderscore}{\kern0pt}rec{\isacharbrackleft}{\kern0pt}OF\ local{\isacharunderscore}{\kern0pt}vars{\isacharunderscore}{\kern0pt}def{\isacharbrackleft}{\kern0pt}THEN\ meta{\isacharunderscore}{\kern0pt}eq{\isacharunderscore}{\kern0pt}to{\isacharunderscore}{\kern0pt}obj{\isacharunderscore}{\kern0pt}eq{\isacharbrackright}{\kern0pt}{\isacharbrackright}{\kern0pt}{\isacharparenright}{\kern0pt}\ simp{\isacharunderscore}{\kern0pt}all%
\endisatagproof
{\isafoldproof}%
%
\isadelimproof
%
\endisadelimproof
%
\isadelimdocument
%
\endisadelimdocument
%
\isatagdocument
%
\isamarkupsection{Denotational Semantics of \isa{IMP\isactrlsub c\isactrlsub o\isactrlsub n\isactrlsub c\isactrlsub u\isactrlsub r}%
}
\isamarkuptrue%
%
\endisatagdocument
{\isafolddocument}%
%
\isadelimdocument
%
\endisadelimdocument
%
\begin{isamarkuptext}%
The denotational semantics of \isa{IMP\isactrlsub c\isactrlsub o\isactrlsub n\isactrlsub c\isactrlsub u\isactrlsub r} is a straight recursive interpretation 
in the domain of HOL-CSP processes. Thread local variables were represented in this 
presentation by updates on a thread-local environment, not by interactions with processes.
The recursive definition of the semantic function fits on a post-card:%
\end{isamarkuptext}\isamarkuptrue%
\isakeywordONE{fun}\isamarkupfalse%
\ Sem\isactrlsub {\isadigit{0}}\ {\isacharcolon}{\kern0pt}{\isacharcolon}{\kern0pt}\ {\isachardoublequoteopen}com\ {\isasymRightarrow}\ {\isacharparenleft}{\kern0pt}{\isasymsigma}\ {\isasymRightarrow}\ evs\ \ process{\isacharparenright}{\kern0pt}\ {\isasymRightarrow}\ {\isasymsigma}\ {\isasymRightarrow}\ evs\ process{\isachardoublequoteclose}\ \isanewline
\ \ \isakeywordTWO{where}\ {\isachardoublequoteopen}Sem\isactrlsub {\isadigit{0}}\ SKIP\ C\ \ \ \ \ \ \ \ \ \ \ \ \ \ \ \ \ \ {\isacharequal}{\kern0pt}\ C{\isachardoublequoteclose}\ \isanewline
\ \ \ \ \ \ \ {\isacharbar}{\kern0pt}{\isachardoublequoteopen}Sem\isactrlsub {\isadigit{0}}\ {\isacharparenleft}{\kern0pt}x\ {\isacharcolon}{\kern0pt}{\isacharequal}{\kern0pt}\ E{\isacharparenright}{\kern0pt}\ C\ \ \ \ \ \ \ \ \ \ \ \ \ \ {\isacharequal}{\kern0pt}\ {\isacharparenleft}{\kern0pt}{\isasymlambda}\ {\isasymsigma}{\isachardot}{\kern0pt}\ C\ {\isacharparenleft}{\kern0pt}{\isasymsigma}{\isacharparenleft}{\kern0pt}x\ {\isacharcolon}{\kern0pt}{\isacharequal}{\kern0pt}\ E\ {\isasymsigma}{\isacharparenright}{\kern0pt}{\isacharparenright}{\kern0pt}{\isacharparenright}{\kern0pt}{\isachardoublequoteclose}\isanewline
\ \ \ \ \ \ \ {\isacharbar}{\kern0pt}{\isachardoublequoteopen}Sem\isactrlsub {\isadigit{0}}\ {\isacharparenleft}{\kern0pt}P\ {\isacharsemicolon}{\kern0pt}\ Q{\isacharparenright}{\kern0pt}\ C\ \ \ \ \ \ \ \ \ \ \ \ \ \ \ {\isacharequal}{\kern0pt}\ {\isacharparenleft}{\kern0pt}Sem\isactrlsub {\isadigit{0}}\ P\ {\isacharparenleft}{\kern0pt}Sem\isactrlsub {\isadigit{0}}\ Q\ C{\isacharparenright}{\kern0pt}{\isacharparenright}{\kern0pt}{\isachardoublequoteclose}\ \isanewline
\ \ \ \ \ \ \ {\isacharbar}{\kern0pt}{\isachardoublequoteopen}Sem\isactrlsub {\isadigit{0}}\ {\isacharparenleft}{\kern0pt}IF\ E\ THEN\ C{\isadigit{1}}\ ELSE\ C{\isadigit{2}}{\isacharparenright}{\kern0pt}\ C{\isacharequal}{\kern0pt}\ {\isacharparenleft}{\kern0pt}{\isasymlambda}\ {\isasymsigma}{\isachardot}{\kern0pt}\ if\ E\ {\isasymsigma}\ \isanewline
\ \ \ \ \ \ \ \ \ \ \ \ \ \ \ \ \ \ \ \ \ \ \ \ \ \ \ \ \ \ \ \ \ \ \ \ \ \ \ \ \ \ \ \ \ \ then\ Sem\isactrlsub {\isadigit{0}}\ C{\isadigit{1}}\ C\ {\isasymsigma}\ \isanewline
\ \ \ \ \ \ \ \ \ \ \ \ \ \ \ \ \ \ \ \ \ \ \ \ \ \ \ \ \ \ \ \ \ \ \ \ \ \ \ \ \ \ \ \ \ \ else\ Sem\isactrlsub {\isadigit{0}}\ C{\isadigit{2}}\ C\ {\isasymsigma}{\isacharparenright}{\kern0pt}{\isachardoublequoteclose}\isanewline
\ \ \ \ \ \ \ {\isacharbar}{\kern0pt}{\isachardoublequoteopen}Sem\isactrlsub {\isadigit{0}}\ {\isacharparenleft}{\kern0pt}WHILE\ E\ DO\ B{\isacharparenright}{\kern0pt}\ C\ \ \ \ \ \ \ \ {\isacharequal}{\kern0pt}\ {\isacharparenleft}{\kern0pt}{\isasymmu}\ X{\isachardot}{\kern0pt}\ {\isacharparenleft}{\kern0pt}{\isasymlambda}\ {\isasymsigma}{\isachardot}{\kern0pt}\ if\ E\ {\isasymsigma}\ \isanewline
\ \ \ \ \ \ \ \ \ \ \ \ \ \ \ \ \ \ \ \ \ \ \ \ \ \ \ \ \ \ \ \ \ \ \ \ \ \ \ \ \ \ \ \ \ \ \ \ \ \ \ \ then\ Sem\isactrlsub {\isadigit{0}}\ B\ X\ {\isasymsigma}\ \isanewline
\ \ \ \ \ \ \ \ \ \ \ \ \ \ \ \ \ \ \ \ \ \ \ \ \ \ \ \ \ \ \ \ \ \ \ \ \ \ \ \ \ \ \ \ \ \ \ \ \ \ \ \ else\ C\ {\isasymsigma}{\isacharparenright}{\kern0pt}{\isacharparenright}{\kern0pt}{\isachardoublequoteclose}\isanewline
\ \ \ \ \ \ \ {\isacharbar}{\kern0pt}{\isachardoublequoteopen}Sem\isactrlsub {\isadigit{0}}\ {\isacharparenleft}{\kern0pt}call\ F{\isacharparenright}{\kern0pt}\ C\ \ \ \ \ \ \ \ \ \ \ \ \ \ {\isacharequal}{\kern0pt}\ {\isacharparenleft}{\kern0pt}{\isasymlambda}\ {\isasymsigma}{\isachardot}{\kern0pt}\ C\ {\isacharparenleft}{\kern0pt}F\ {\isasymsigma}{\isacharparenright}{\kern0pt}{\isacharparenright}{\kern0pt}{\isachardoublequoteclose}\isanewline
\ \ \ \ \ \ \ {\isacharbar}{\kern0pt}{\isachardoublequoteopen}Sem\isactrlsub {\isadigit{0}}\ {\isacharparenleft}{\kern0pt}com{\isachardot}{\kern0pt}STOP{\isacharparenright}{\kern0pt}\ C\ \ \ \ \ \ \ \ \ \ \ \ {\isacharequal}{\kern0pt}\ {\isacharparenleft}{\kern0pt}{\isasymlambda}\ {\isasymsigma}{\isachardot}{\kern0pt}\ Constant{\isacharunderscore}{\kern0pt}Processes{\isachardot}{\kern0pt}STOP{\isacharparenright}{\kern0pt}{\isachardoublequoteclose}\isanewline
\ \ \ \ \ \ \ {\isacharbar}{\kern0pt}{\isachardoublequoteopen}Sem\isactrlsub {\isadigit{0}}\ {\isacharparenleft}{\kern0pt}com{\isachardot}{\kern0pt}lock\ n{\isacharparenright}{\kern0pt}\ C\ \ \ \ \ \ \ \ \ \ {\isacharequal}{\kern0pt}\ {\isacharparenleft}{\kern0pt}{\isasymlambda}\ {\isasymsigma}{\isachardot}{\kern0pt}\ evs{\isachardot}{\kern0pt}lock\ n\ {\isasymrightarrow}\ C\ {\isasymsigma}{\isacharparenright}{\kern0pt}{\isachardoublequoteclose}\isanewline
\ \ \ \ \ \ \ {\isacharbar}{\kern0pt}{\isachardoublequoteopen}Sem\isactrlsub {\isadigit{0}}\ {\isacharparenleft}{\kern0pt}com{\isachardot}{\kern0pt}unlock\ n{\isacharparenright}{\kern0pt}\ C\ \ \ \ \ \ \ \ {\isacharequal}{\kern0pt}\ {\isacharparenleft}{\kern0pt}{\isasymlambda}\ {\isasymsigma}{\isachardot}{\kern0pt}\ evs{\isachardot}{\kern0pt}unlock\ n\ \ {\isasymrightarrow}\ C\ {\isasymsigma}{\isacharparenright}{\kern0pt}{\isachardoublequoteclose}\isanewline
\ \ \ \ \ \ \ {\isacharbar}{\kern0pt}{\isachardoublequoteopen}Sem\isactrlsub {\isadigit{0}}\ {\isacharparenleft}{\kern0pt}STORE\ E\ TO\ X\isactrlsub g\isactrlsub l\isactrlsub o{\isacharparenright}{\kern0pt}\ C\ \ \ \ \ \ {\isacharequal}{\kern0pt}\ {\isacharparenleft}{\kern0pt}{\isasymlambda}\ {\isasymsigma}{\isachardot}{\kern0pt}\ upd\isactrlsub g\isactrlsub l\isactrlsub o\isactrlsub b\isactrlsub a\isactrlsub l\ X\isactrlsub g\isactrlsub l\isactrlsub o\ {\isacharparenleft}{\kern0pt}E\ {\isasymsigma}{\isacharparenright}{\kern0pt}\ {\isasymrightarrow}\ C\ {\isasymsigma}{\isacharparenright}{\kern0pt}{\isachardoublequoteclose}\isanewline
\ \ \ \ \ \ \ {\isacharbar}{\kern0pt}{\isachardoublequoteopen}Sem\isactrlsub {\isadigit{0}}\ {\isacharparenleft}{\kern0pt}LOAD\ X\isactrlsub l\isactrlsub o\isactrlsub c\ FROM\ X\isactrlsub g\isactrlsub l\isactrlsub o{\isacharparenright}{\kern0pt}\ C\ \ \ {\isacharequal}{\kern0pt}\ {\isacharparenleft}{\kern0pt}{\isasymlambda}\ {\isasymsigma}{\isachardot}{\kern0pt}\ {\isacharparenleft}{\kern0pt}read\isactrlsub g\isactrlsub l\isactrlsub o\isactrlsub b\isactrlsub a\isactrlsub l\ X\isactrlsub g\isactrlsub l\isactrlsub o{\isacharparenright}{\kern0pt}\isactrlbold {\isacharquery}{\kern0pt}x\ \isanewline
\ \ \ \ \ \ \ \ \ \ \ \ \ \ \ \ \ \ \ \ \ \ \ \ \ \ \ \ \ \ \ \ \ \ \ \ \ \ \ \ \ \ \ \ \ \ \ \ \ \ {\isasymrightarrow}\ C{\isacharparenleft}{\kern0pt}{\isasymsigma}{\isacharparenleft}{\kern0pt}X\isactrlsub l\isactrlsub o\isactrlsub c{\isacharcolon}{\kern0pt}{\isacharequal}{\kern0pt}x{\isacharparenright}{\kern0pt}{\isacharparenright}{\kern0pt}{\isacharparenright}{\kern0pt}{\isachardoublequoteclose}\isanewline
\isanewline
\isanewline
\isakeywordONE{definition}\isamarkupfalse%
\ initial{\isacharunderscore}{\kern0pt}state\ {\isacharcolon}{\kern0pt}{\isacharcolon}{\kern0pt}\ {\isasymsigma}\ {\isacharparenleft}{\kern0pt}{\isacartoucheopen}{\isasymsigma}\isactrlsub {\isadigit{0}}{\isacartoucheclose}{\isacharparenright}{\kern0pt}\isanewline
\ \ \isakeywordTWO{where}\ \ \ {\isacartoucheopen}{\isasymsigma}\isactrlsub {\isadigit{0}}\ {\isasymequiv}\ {\isacharparenleft}{\kern0pt}{\isasymlambda}{\isacharunderscore}{\kern0pt}{\isachardot}{\kern0pt}\ {\isadigit{0}}{\isacharcolon}{\kern0pt}{\isacharcolon}{\kern0pt}int{\isacharparenright}{\kern0pt}{\isacartoucheclose}\isanewline
\isanewline
\isakeywordONE{definition}\isamarkupfalse%
\ some{\isacharunderscore}{\kern0pt}state\ {\isacharcolon}{\kern0pt}{\isacharcolon}{\kern0pt}\ {\isasymsigma}\ {\isacharparenleft}{\kern0pt}{\isacartoucheopen}{\isasymsigma}\isactrlsub e\isactrlsub x\isactrlsub p\isactrlsub l{\isacartoucheclose}{\isacharparenright}{\kern0pt}\isanewline
\ \ \isakeywordTWO{where}\ \ \ {\isacartoucheopen}{\isasymsigma}\isactrlsub e\isactrlsub x\isactrlsub p\isactrlsub l\ v\ {\isasymequiv}\ if\ v\ {\isacharequal}{\kern0pt}\ {\isacharprime}{\kern0pt}{\isacharprime}{\kern0pt}a{\isacharprime}{\kern0pt}{\isacharprime}{\kern0pt}\ then\ {\isadigit{3}}\ else\ {\isadigit{0}}{\isacartoucheclose}%
\begin{isamarkuptext}%
In case that an \isatt{a{\kern0pt}s{\kern0pt}s{\kern0pt}u{\kern0pt}m{\kern0pt}e{\kern0pt}}-primitive on initial states has to be modeled
(the usual contruct for modeling preconditions), this can be done by the 
Hilbert-Choice operator \isa{SOME\ \isabound{{\isasymsigma}}{\isachardot}{\kern0pt}\ \isafree{E}} on initial local and global states.%
\end{isamarkuptext}\isamarkuptrue%
\isakeywordONE{definition}\isamarkupfalse%
\ Sem\ \ \ \ \ \ \ \ \ \isakeywordTWO{where}\ {\isacartoucheopen}Sem\ Thread\ {\isasymequiv}\ Sem\isactrlsub {\isadigit{0}}\ Thread\ {\isacharparenleft}{\kern0pt}{\isasymlambda}{\isacharunderscore}{\kern0pt}{\isachardot}{\kern0pt}\ Skip{\isacharparenright}{\kern0pt}\ {\isasymsigma}\isactrlsub {\isadigit{0}}{\isacartoucheclose}\ \ \isanewline
\ \ \ %
\isamarkupcmt{final continuation: \isa{Skip}, execution starts with initialized local vars
       for simplicity.%
}%
\isadelimdocument
%
\endisadelimdocument
%
\isatagdocument
%
\isamarkupsection{A Thread-System as HOL-CSPM-Architecture%
}
\isamarkuptrue%
%
\endisatagdocument
{\isafolddocument}%
%
\isadelimdocument
%
\endisadelimdocument
%
\begin{isamarkuptext}%
An example of a global \isa{IMP\isactrlsub c\isactrlsub o\isactrlsub n\isactrlsub c\isactrlsub u\isactrlsub r}-System with 2 threads, 3 global variables and 4 semaphores:%
\end{isamarkuptext}\isamarkuptrue%
\isakeywordONE{definition}\isamarkupfalse%
{\isacartoucheopen}Thread{\isadigit{3}}{\isacharunderscore}{\kern0pt}sem\ {\isasymequiv}\ {\isacharparenleft}{\kern0pt}\ \isactrlbold {\isacharbar}{\kern0pt}\isactrlbold {\isacharbar}{\kern0pt}\isactrlbold {\isacharbar}{\kern0pt}\ idx\ {\isasymin}{\isacharhash}{\kern0pt}\ mset\ {\isacharbrackleft}{\kern0pt}{\isacharprime}{\kern0pt}{\isacharprime}{\kern0pt}a{\isacharprime}{\kern0pt}{\isacharprime}{\kern0pt}{\isacharcomma}{\kern0pt}{\isacharprime}{\kern0pt}{\isacharprime}{\kern0pt}b{\isacharprime}{\kern0pt}{\isacharprime}{\kern0pt}{\isacharcomma}{\kern0pt}{\isacharprime}{\kern0pt}{\isacharprime}{\kern0pt}c{\isacharprime}{\kern0pt}{\isacharprime}{\kern0pt}{\isacharbrackright}{\kern0pt}{\isachardot}{\kern0pt}\ \ global{\isacharunderscore}{\kern0pt}vars\ {\isasymsigma}\isactrlsub {\isadigit{0}}\ idx\ \ \isanewline
\ \ \ \ \ \ \ \ \ \ \ \ \ \ \ \ \ \ \ \ \ \ \ \ \ \ \ {\isacharbar}{\kern0pt}{\isacharbar}{\kern0pt}{\isacharbar}{\kern0pt}\isanewline
\ \ \ \ \ \ \ \ \ \ \ \ \ \ \ \ \ \ \ \ \ \ \ \ \ \ \ \isactrlbold {\isacharbar}{\kern0pt}\isactrlbold {\isacharbar}{\kern0pt}\isactrlbold {\isacharbar}{\kern0pt}\ idx\ {\isasymin}{\isacharhash}{\kern0pt}\ mset\ {\isacharbrackleft}{\kern0pt}{\isadigit{1}}{\isachardot}{\kern0pt}{\isachardot}{\kern0pt}{\isadigit{4}}{\isacharbrackright}{\kern0pt}{\isachardot}{\kern0pt}\ semaphore\ idx\ {\isacharparenright}{\kern0pt}\isanewline
\ \ \ \ \ \ \ \ \ \ \isanewline
\ \ \ \ \ \ \ \ \ \ \ \ \ \ \ \ \ \ \ \ \ \ \ \ \ \ \ {\isacharbar}{\kern0pt}{\isacharbar}{\kern0pt}\isanewline
\ \ \ \ \ \ \ \ \ \ \ \ \isanewline
\ \ \ \ \ \ \ \ \ \ \ \ \ \ \ \ \ \ \ \ \ \ \ \ \ \ \ {\isacharparenleft}{\kern0pt}Sem\ Thread{\isadigit{1}}\ {\isacharbar}{\kern0pt}{\isacharbar}{\kern0pt}{\isacharbar}{\kern0pt}\ Sem\ Thread{\isadigit{2}}{\isacharparenright}{\kern0pt}\isanewline
\ \ \ \ \ \ \ \ \ \ \ \ \ \ \ \ \ \ \ \ \ \ \ \ \ \ \ {\isacartoucheclose}\isanewline
\isanewline
%
\isadelimML
\isanewline
%
\endisadelimML
%
\isatagML
\isakeywordONE{ML}\isamarkupfalse%
{\isacartoucheopen}{\isacharparenleft}{\kern0pt}{\isacharasterisk}{\kern0pt}\ {\isachardot}{\kern0pt}{\isachardot}{\kern0pt}{\isachardot}{\kern0pt}\ and\ this\ reads\ at\ the\ term{\isacharminus}{\kern0pt}level\ as\ follows\ {\isacharcolon}{\kern0pt}\ {\isacharasterisk}{\kern0pt}{\isacharparenright}{\kern0pt}\isanewline
\ \ \ val\ {\isacharunderscore}{\kern0pt}\ {\isachardollar}{\kern0pt}\ {\isacharunderscore}{\kern0pt}\ {\isachardollar}{\kern0pt}\ temp{\isacharequal}{\kern0pt}\ Thm{\isachardot}{\kern0pt}concl{\isacharunderscore}{\kern0pt}of\ {\isacharat}{\kern0pt}{\isacharbraceleft}{\kern0pt}thm\ Thread{\isadigit{3}}{\isacharunderscore}{\kern0pt}sem{\isacharunderscore}{\kern0pt}def{\isacharbraceright}{\kern0pt}\ \ \isanewline
\ \ \ val\ temp{\isadigit{2}}{\isacharequal}{\kern0pt}\ \isactrlterm {\isasymopen}\isactrlbold {\isacharbar}{\kern0pt}\isactrlbold {\isacharbar}{\kern0pt}\isactrlbold {\isacharbar}{\kern0pt}\ idx\ {\isasymin}{\isacharhash}{\kern0pt}\ mset\ {\isacharbrackleft}{\kern0pt}{\isacharprime}{\kern0pt}{\isacharprime}{\kern0pt}a{\isacharprime}{\kern0pt}{\isacharprime}{\kern0pt}{\isacharcomma}{\kern0pt}{\isacharprime}{\kern0pt}{\isacharprime}{\kern0pt}b{\isacharprime}{\kern0pt}{\isacharprime}{\kern0pt}{\isacharcomma}{\kern0pt}{\isacharprime}{\kern0pt}{\isacharprime}{\kern0pt}c{\isacharprime}{\kern0pt}{\isacharprime}{\kern0pt}{\isacharbrackright}{\kern0pt}{\isachardot}{\kern0pt}\ \ global{\isacharunderscore}{\kern0pt}vars\ {\isasymsigma}\isactrlsub {\isadigit{0}}\ idx{\isasymclose}\isanewline
\ \ {\isacartoucheclose}%
\endisatagML
{\isafoldML}%
%
\isadelimML
%
\endisadelimML
%
\begin{isamarkuptext}%
Tests:%
\end{isamarkuptext}\isamarkuptrue%
\isakeywordONE{lemma}\isamarkupfalse%
\ Thread{\isadigit{2}}{\isacharunderscore}{\kern0pt}sem\ {\isacharcolon}{\kern0pt}\ \isanewline
\ \ {\isachardoublequoteopen}Sem{\isacharparenleft}{\kern0pt}Thread{\isadigit{2}}{\isacharparenright}{\kern0pt}\ {\isacharequal}{\kern0pt}\ \ evs{\isachardot}{\kern0pt}lock\ {\isadigit{3}}\ {\isasymrightarrow}\ {\isacharparenleft}{\kern0pt}{\isasymmu}\ x{\isachardot}{\kern0pt}\ {\isacharparenleft}{\kern0pt}{\isasymlambda}{\isasymsigma}{\isacharcolon}{\kern0pt}{\isacharcolon}{\kern0pt}{\isasymsigma}{\isachardot}{\kern0pt}\ evs{\isachardot}{\kern0pt}lock\ {\isadigit{4}}\ {\isasymrightarrow}\ x\ {\isasymsigma}{\isacharparenright}{\kern0pt}{\isacharparenright}{\kern0pt}\ {\isasymsigma}\isactrlsub e\isactrlsub x\isactrlsub p\isactrlsub l{\isachardoublequoteclose}\isanewline
%
\isadelimproof
\ \ %
\endisadelimproof
%
\isatagproof
\isakeywordONE{unfolding}\isamarkupfalse%
\ Thread{\isadigit{2}}{\isacharunderscore}{\kern0pt}def\ Sem{\isacharunderscore}{\kern0pt}def\ some{\isacharunderscore}{\kern0pt}state{\isacharunderscore}{\kern0pt}def\ initial{\isacharunderscore}{\kern0pt}state{\isacharunderscore}{\kern0pt}def\isanewline
\ \ \isakeywordONE{by}\isamarkupfalse%
\ {\isacharparenleft}{\kern0pt}simp\ add{\isacharcolon}{\kern0pt}\ fun{\isacharunderscore}{\kern0pt}upd{\isacharunderscore}{\kern0pt}def{\isacharparenright}{\kern0pt}%
\endisatagproof
{\isafoldproof}%
%
\isadelimproof
\isanewline
%
\endisadelimproof
\isanewline
\isanewline
\isakeywordONE{schematic{\isacharunderscore}{\kern0pt}goal}\isamarkupfalse%
\ K\ {\isacharcolon}{\kern0pt}\ {\isachardoublequoteopen}Thread{\isadigit{3}}{\isacharunderscore}{\kern0pt}sem\ {\isacharequal}{\kern0pt}\ {\isacharquery}{\kern0pt}X{\isachardoublequoteclose}\isanewline
%
\isadelimproof
\ \ %
\endisadelimproof
%
\isatagproof
\isakeywordONE{unfolding}\isamarkupfalse%
\ Thread{\isadigit{3}}{\isacharunderscore}{\kern0pt}sem{\isacharunderscore}{\kern0pt}def\ Sem{\isacharunderscore}{\kern0pt}def\isanewline
\ \ \isakeywordONE{unfolding}\isamarkupfalse%
\ Thread{\isadigit{1}}{\isacharunderscore}{\kern0pt}def\ Thread{\isadigit{2}}{\isacharunderscore}{\kern0pt}def\isanewline
\ \ \ \isakeywordONE{apply}\isamarkupfalse%
{\isacharparenleft}{\kern0pt}rule\ trans{\isacharparenright}{\kern0pt}\isanewline
\ \ \ \isakeywordONE{apply}\isamarkupfalse%
\ {\isacharparenleft}{\kern0pt}simp\ add{\isacharcolon}{\kern0pt}\ List{\isachardot}{\kern0pt}upto{\isachardot}{\kern0pt}simps\ cong{\isacharcolon}{\kern0pt}\ HOL{\isachardot}{\kern0pt}if{\isacharunderscore}{\kern0pt}cong{\isacharparenright}{\kern0pt}\isanewline
\ \ \isakeywordONE{by}\isamarkupfalse%
\ {\isacharparenleft}{\kern0pt}rule\ refl{\isacharparenright}{\kern0pt}%
\endisatagproof
{\isafoldproof}%
%
\isadelimproof
\isanewline
%
\endisadelimproof
\ \ \isanewline
\isanewline
\isanewline
\isakeywordONE{lemma}\isamarkupfalse%
\ Thread{\isadigit{3}}{\isacharunderscore}{\kern0pt}sem\ {\isacharcolon}{\kern0pt}\ \isanewline
\ \ \ \ {\isacartoucheopen}\ Thread{\isadigit{3}}{\isacharunderscore}{\kern0pt}sem\ {\isacharequal}{\kern0pt}\ {\isacharparenleft}{\kern0pt}{\isacharparenleft}{\kern0pt}global{\isacharunderscore}{\kern0pt}vars\ {\isasymsigma}\isactrlsub {\isadigit{0}}\ {\isacharprime}{\kern0pt}{\isacharprime}{\kern0pt}a{\isacharprime}{\kern0pt}{\isacharprime}{\kern0pt}\ {\isacharbar}{\kern0pt}{\isacharbar}{\kern0pt}{\isacharbar}{\kern0pt}\ {\isacharparenleft}{\kern0pt}global{\isacharunderscore}{\kern0pt}vars\ {\isasymsigma}\isactrlsub {\isadigit{0}}\ {\isacharprime}{\kern0pt}{\isacharprime}{\kern0pt}b{\isacharprime}{\kern0pt}{\isacharprime}{\kern0pt}\ {\isacharbar}{\kern0pt}{\isacharbar}{\kern0pt}{\isacharbar}{\kern0pt}\ global{\isacharunderscore}{\kern0pt}vars\ {\isasymsigma}\isactrlsub {\isadigit{0}}\ {\isacharprime}{\kern0pt}{\isacharprime}{\kern0pt}c{\isacharprime}{\kern0pt}{\isacharprime}{\kern0pt}{\isacharparenright}{\kern0pt}{\isacharparenright}{\kern0pt}\ \isanewline
\ \ \ \ \ \ \ \ \ \ \ \ \ \ \ \ \ \ \ \ \ {\isacharbar}{\kern0pt}{\isacharbar}{\kern0pt}{\isacharbar}{\kern0pt}\ {\isacharparenleft}{\kern0pt}semaphore\ {\isadigit{1}}\ {\isacharbar}{\kern0pt}{\isacharbar}{\kern0pt}{\isacharbar}{\kern0pt}\ {\isacharparenleft}{\kern0pt}semaphore\ {\isadigit{2}}\ {\isacharbar}{\kern0pt}{\isacharbar}{\kern0pt}{\isacharbar}{\kern0pt}\ {\isacharparenleft}{\kern0pt}semaphore\ {\isadigit{3}}\ {\isacharbar}{\kern0pt}{\isacharbar}{\kern0pt}{\isacharbar}{\kern0pt}\ semaphore\ {\isadigit{4}}{\isacharparenright}{\kern0pt}{\isacharparenright}{\kern0pt}{\isacharparenright}{\kern0pt}{\isacharparenright}{\kern0pt}\isanewline
\ \ \ \ \ \ \ \ \ \ \ \ \ \ \ \ \ \ \ \ {\isacharbar}{\kern0pt}{\isacharbar}{\kern0pt}\ \isanewline
\ \ \ \ \ \ \ \ \ \ \ \ \ \ \ \ \ \ \ \ {\isacharparenleft}{\kern0pt}{\isacharparenleft}{\kern0pt}evs{\isachardot}{\kern0pt}lock\ {\isadigit{3}}\ {\isasymrightarrow}\ {\isacharparenleft}{\kern0pt}{\isasymmu}\ x{\isachardot}{\kern0pt}\ {\isacharparenleft}{\kern0pt}{\isasymlambda}{\isasymsigma}{\isachardot}{\kern0pt}\ if\ {\isadigit{3}}\ {\isacharless}{\kern0pt}\ {\isasymsigma}\ {\isacharprime}{\kern0pt}{\isacharprime}{\kern0pt}e{\isacharprime}{\kern0pt}{\isacharprime}{\kern0pt}\ \isanewline
\ \ \ \ \ \ \ \ \ \ \ \ \ \ \ \ \ \ \ \ \ \ \ \ \ \ \ \ \ \ \ \ \ \ \ \ \ \ \ \ \ \ \ \ \ \ \ then\ evs{\isachardot}{\kern0pt}lock\ {\isadigit{4}}\ {\isasymrightarrow}\ x\ {\isasymsigma}\isanewline
\ \ \ \ \ \ \ \ \ \ \ \ \ \ \ \ \ \ \ \ \ \ \ \ \ \ \ \ \ \ \ \ \ \ \ \ \ \ \ \ \ \ \ \ \ \ \ else\ upd\isactrlsub g\isactrlsub l\isactrlsub o\isactrlsub b\isactrlsub a\isactrlsub l\ {\isacharprime}{\kern0pt}{\isacharprime}{\kern0pt}b{\isacharprime}{\kern0pt}{\isacharprime}{\kern0pt}\ {\isacharparenleft}{\kern0pt}{\isasymsigma}\ {\isacharprime}{\kern0pt}{\isacharprime}{\kern0pt}d{\isacharprime}{\kern0pt}{\isacharprime}{\kern0pt}\ {\isacharplus}{\kern0pt}\ {\isadigit{1}}{\isacharparenright}{\kern0pt}\isanewline
\ \ \ \ \ \ \ \ \ \ \ \ \ \ \ \ \ \ \ \ \ \ \ \ \ \ \ \ \ \ \ \ \ \ \ \ \ \ \ \ \ \ \ \ \ \ \ \ \ \ \ \ \ {\isasymrightarrow}\ read\isactrlsub g\isactrlsub l\isactrlsub o\isactrlsub b\isactrlsub a\isactrlsub l\ {\isacharprime}{\kern0pt}{\isacharprime}{\kern0pt}b{\isacharprime}{\kern0pt}{\isacharprime}{\kern0pt}\isactrlbold {\isacharquery}{\kern0pt}x\ \isanewline
\ \ \ \ \ \ \ \ \ \ \ \ \ \ \ \ \ \ \ \ \ \ \ \ \ \ \ \ \ \ \ \ \ \ \ \ \ \ \ \ \ \ \ \ \ \ \ \ \ \ \ \ \ {\isasymrightarrow}\ evs{\isachardot}{\kern0pt}unlock\ {\isadigit{3}}\ {\isasymrightarrow}\ Skip{\isacharparenright}{\kern0pt}{\isacharparenright}{\kern0pt}\isanewline
\ \ \ \ \ \ \ \ \ \ \ \ \ \ \ \ \ \ \ \ \ \ \ {\isacharparenleft}{\kern0pt}{\isacharparenleft}{\kern0pt}{\isasymlambda}{\isacharunderscore}{\kern0pt}{\isachardot}{\kern0pt}\ {\isadigit{0}}{\isacharcolon}{\kern0pt}{\isacharcolon}{\kern0pt}int{\isacharparenright}{\kern0pt}{\isacharparenleft}{\kern0pt}{\isacharprime}{\kern0pt}{\isacharprime}{\kern0pt}a{\isacharprime}{\kern0pt}{\isacharprime}{\kern0pt}\ {\isacharcolon}{\kern0pt}{\isacharequal}{\kern0pt}\ {\isadigit{3}}{\isacharcolon}{\kern0pt}{\isacharcolon}{\kern0pt}int{\isacharparenright}{\kern0pt}{\isacharparenright}{\kern0pt}{\isacharparenright}{\kern0pt}\ \isanewline
\ \ \ \ \ \ \ \ \ \ \ \ \ \ \ \ \ \ \ \ \ {\isacharbar}{\kern0pt}{\isacharbar}{\kern0pt}{\isacharbar}{\kern0pt}\ \isanewline
\ \ \ \ \ \ \ \ \ \ \ \ \ \ \ \ \ \ \ \ \ {\isacharparenleft}{\kern0pt}evs{\isachardot}{\kern0pt}lock\ {\isadigit{3}}\ {\isasymrightarrow}\ {\isacharparenleft}{\kern0pt}{\isasymmu}\ x{\isachardot}{\kern0pt}\ {\isacharparenleft}{\kern0pt}{\isasymlambda}{\isasymsigma}{\isachardot}{\kern0pt}\ evs{\isachardot}{\kern0pt}lock\ {\isadigit{4}}\ {\isasymrightarrow}\ x\ {\isasymsigma}{\isacharparenright}{\kern0pt}{\isacharparenright}{\kern0pt}\ {\isacharparenleft}{\kern0pt}{\isacharparenleft}{\kern0pt}{\isasymlambda}{\isacharunderscore}{\kern0pt}{\isachardot}{\kern0pt}\ {\isadigit{0}}{\isacharcolon}{\kern0pt}{\isacharcolon}{\kern0pt}int{\isacharparenright}{\kern0pt}{\isacharparenleft}{\kern0pt}{\isacharprime}{\kern0pt}{\isacharprime}{\kern0pt}a{\isacharprime}{\kern0pt}{\isacharprime}{\kern0pt}{\isacharcolon}{\kern0pt}{\isacharequal}{\kern0pt}\ {\isadigit{3}}{\isacharcolon}{\kern0pt}{\isacharcolon}{\kern0pt}int{\isacharparenright}{\kern0pt}{\isacharparenright}{\kern0pt}{\isacharparenright}{\kern0pt}{\isacharparenright}{\kern0pt}\ \isanewline
\ \ \ \ {\isacartoucheclose}\isanewline
%
\isadelimproof
\ \ %
\endisadelimproof
%
\isatagproof
\isakeywordONE{unfolding}\isamarkupfalse%
\ Thread{\isadigit{3}}{\isacharunderscore}{\kern0pt}sem{\isacharunderscore}{\kern0pt}def\ Sem{\isacharunderscore}{\kern0pt}def\ Thread{\isadigit{1}}{\isacharunderscore}{\kern0pt}def\ Thread{\isadigit{2}}{\isacharunderscore}{\kern0pt}def\ initial{\isacharunderscore}{\kern0pt}state{\isacharunderscore}{\kern0pt}def\isanewline
\ \ \isakeywordONE{by}\isamarkupfalse%
\ {\isacharparenleft}{\kern0pt}simp\ add{\isacharcolon}{\kern0pt}\ List{\isachardot}{\kern0pt}upto{\isachardot}{\kern0pt}simps\ cong{\isacharcolon}{\kern0pt}\ HOL{\isachardot}{\kern0pt}if{\isacharunderscore}{\kern0pt}cong{\isacharparenright}{\kern0pt}%
\endisatagproof
{\isafoldproof}%
%
\isadelimproof
\isanewline
%
\endisadelimproof
\ \ \isanewline
\isanewline
%
\isadelimML
\isanewline
%
\endisadelimML
%
\isatagML
\isakeywordONE{ML}\isamarkupfalse%
{\isacartoucheopen}\isanewline
val\ x\ {\isacharequal}{\kern0pt}HOLogic{\isachardot}{\kern0pt}dest{\isacharunderscore}{\kern0pt}eq\ {\isacharparenleft}{\kern0pt}HOLogic{\isachardot}{\kern0pt}dest{\isacharunderscore}{\kern0pt}Trueprop\ {\isacharparenleft}{\kern0pt}\ Thm{\isachardot}{\kern0pt}concl{\isacharunderscore}{\kern0pt}of\ \ {\isacharat}{\kern0pt}{\isacharbraceleft}{\kern0pt}thm\ K{\isacharbraceright}{\kern0pt}{\isacharparenright}{\kern0pt}{\isacharparenright}{\kern0pt}\ {\isacharsemicolon}{\kern0pt}\ \isanewline
val\ y\ {\isacharequal}{\kern0pt}\ {\isacharat}{\kern0pt}{\isacharbraceleft}{\kern0pt}thm\ Thread{\isadigit{3}}{\isacharunderscore}{\kern0pt}sem{\isacharbraceright}{\kern0pt}\ {\isacharbar}{\kern0pt}{\isachargreater}{\kern0pt}\ Thm{\isachardot}{\kern0pt}concl{\isacharunderscore}{\kern0pt}of\ \isanewline
\ \ \ \ \ \ \ \ \ \ \ \ \ \ \ \ \ \ \ \ \ \ \ \ \ \ \ {\isacharbar}{\kern0pt}{\isachargreater}{\kern0pt}\ HOLogic{\isachardot}{\kern0pt}dest{\isacharunderscore}{\kern0pt}Trueprop\ \isanewline
\ \ \ \ \ \ \ \ \ \ \ \ \ \ \ \ \ \ \ \ \ \ \ \ \ \ \ {\isacharbar}{\kern0pt}{\isachargreater}{\kern0pt}\ HOLogic{\isachardot}{\kern0pt}dest{\isacharunderscore}{\kern0pt}eq\isanewline
{\isacartoucheclose}%
\endisatagML
{\isafoldML}%
%
\isadelimML
%
\endisadelimML
%
\isadelimdocument
%
\endisadelimdocument
%
\isatagdocument
%
\isamarkupsection{A Semantic Corner-Case of \isa{IMP\isactrlsub c\isactrlsub o\isactrlsub n\isactrlsub c\isactrlsub u\isactrlsub r}%
}
\isamarkuptrue%
%
\endisatagdocument
{\isafolddocument}%
%
\isadelimdocument
%
\endisadelimdocument
%
\begin{isamarkuptext}%
We study the case of three global variables and 4 locks 
and a computation that spends its time in an infinite loop without
communicating. (The result would be the same if inside the loop,
only compuations on local variables were done.%
\end{isamarkuptext}\isamarkuptrue%
\isakeywordONE{definition}\isamarkupfalse%
\ Thread{\isadigit{4}}\ \isakeywordTWO{where}\ \isanewline
\ \ {\isachardoublequoteopen}Thread{\isadigit{4}}\ {\isacharequal}{\kern0pt}\ WHILE\ {\isacharparenleft}{\kern0pt}{\isasymlambda}{\isasymsigma}{\isachardot}{\kern0pt}\ True{\isacharparenright}{\kern0pt}\ DO\ \ SKIP{\isachardoublequoteclose}\isanewline
\isanewline
\isakeywordONE{definition}\isamarkupfalse%
\ Thread{\isadigit{4}}{\isacharunderscore}{\kern0pt}sem\ \isakeywordTWO{where}\ \isanewline
{\isachardoublequoteopen}Thread{\isadigit{4}}{\isacharunderscore}{\kern0pt}sem\ {\isasymequiv}\ {\isacharparenleft}{\kern0pt}\ \isactrlbold {\isacharbar}{\kern0pt}\isactrlbold {\isacharbar}{\kern0pt}\isactrlbold {\isacharbar}{\kern0pt}\ idx\ {\isasymin}{\isacharhash}{\kern0pt}\ mset\ {\isacharbrackleft}{\kern0pt}{\isacharprime}{\kern0pt}{\isacharprime}{\kern0pt}a{\isacharprime}{\kern0pt}{\isacharprime}{\kern0pt}{\isacharcomma}{\kern0pt}{\isacharprime}{\kern0pt}{\isacharprime}{\kern0pt}b{\isacharprime}{\kern0pt}{\isacharprime}{\kern0pt}{\isacharcomma}{\kern0pt}{\isacharprime}{\kern0pt}{\isacharprime}{\kern0pt}c{\isacharprime}{\kern0pt}{\isacharprime}{\kern0pt}{\isacharbrackright}{\kern0pt}{\isachardot}{\kern0pt}\ \ global{\isacharunderscore}{\kern0pt}vars\ {\isasymsigma}\isactrlsub {\isadigit{0}}\ idx\ \ \isanewline
\ \ \ \ \ \ \ \ \ \ \ \ \ \ \ \ \ \ {\isacharbar}{\kern0pt}{\isacharbar}{\kern0pt}{\isacharbar}{\kern0pt}\isanewline
\ \ \ \ \ \ \ \ \ \ \ \ \ \ \ \ \ \ \isactrlbold {\isacharbar}{\kern0pt}\isactrlbold {\isacharbar}{\kern0pt}\isactrlbold {\isacharbar}{\kern0pt}\ idx\ {\isasymin}{\isacharhash}{\kern0pt}\ mset\ {\isacharbrackleft}{\kern0pt}{\isadigit{1}}{\isachardot}{\kern0pt}{\isachardot}{\kern0pt}{\isadigit{4}}{\isacharbrackright}{\kern0pt}{\isachardot}{\kern0pt}\ semaphore\ idx\ {\isacharparenright}{\kern0pt}\isanewline
\ \ \ \ \ \ \ \ \ \ \ \ \ \ \ \ \ \isanewline
\ \ \ \ \ \ \ \ \ \ \ \ \ \ \ \ {\isacharbar}{\kern0pt}{\isacharbar}{\kern0pt}\isanewline
\ \ \ \isanewline
\ \ \ \ \ \ \ \ \ \ \ \ \ \ \ \ \ \ {\isacharparenleft}{\kern0pt}Sem\ Thread{\isadigit{4}}{\isacharparenright}{\kern0pt}{\isachardoublequoteclose}%
\begin{isamarkuptext}%
The following theory shows that this type of programs collapses to \isa{{\isasymbottom}}.%
\end{isamarkuptext}\isamarkuptrue%
\isakeywordONE{theorem}\isamarkupfalse%
\ bang\ {\isacharcolon}{\kern0pt}\ {\isacartoucheopen}Thread{\isadigit{4}}{\isacharunderscore}{\kern0pt}sem\ {\isacharequal}{\kern0pt}\ {\isasymbottom}{\isacartoucheclose}\isanewline
%
\isadelimproof
\ \ %
\endisadelimproof
%
\isatagproof
\isakeywordONE{unfolding}\isamarkupfalse%
\ Thread{\isadigit{4}}{\isacharunderscore}{\kern0pt}sem{\isacharunderscore}{\kern0pt}def\ Thread{\isadigit{4}}{\isacharunderscore}{\kern0pt}def\ Sem{\isacharunderscore}{\kern0pt}def\isanewline
\ \ \isakeywordONE{by}\isamarkupfalse%
\ {\isacharparenleft}{\kern0pt}simp\ add{\isacharcolon}{\kern0pt}\ List{\isachardot}{\kern0pt}upto{\isachardot}{\kern0pt}simps\ Fixrec{\isachardot}{\kern0pt}fix{\isacharunderscore}{\kern0pt}id\ cong{\isacharcolon}{\kern0pt}\ HOL{\isachardot}{\kern0pt}if{\isacharunderscore}{\kern0pt}cong{\isacharparenright}{\kern0pt}%
\endisatagproof
{\isafoldproof}%
%
\isadelimproof
%
\endisadelimproof
%
\isadelimdocument
%
\endisadelimdocument
%
\isatagdocument
%
\isamarkupsection{Another Semantic Corner-Case: Deadlock%
}
\isamarkuptrue%
%
\endisatagdocument
{\isafolddocument}%
%
\isadelimdocument
%
\endisadelimdocument
%
\begin{isamarkuptext}%
For simplicity of the subsequent proof, we model a system without global variables:%
\end{isamarkuptext}\isamarkuptrue%
\isakeywordONE{definition}\isamarkupfalse%
\ Thread{\isadigit{5}}\ \isakeywordTWO{where}\ \isanewline
\ \ {\isachardoublequoteopen}Thread{\isadigit{5}}\ {\isacharequal}{\kern0pt}\ WHILE\ {\isacharparenleft}{\kern0pt}{\isasymlambda}{\isasymsigma}{\isachardot}{\kern0pt}\ True{\isacharparenright}{\kern0pt}\ DO\ {\isacharparenleft}{\kern0pt}com{\isachardot}{\kern0pt}lock\ {\isadigit{1}}{\isacharparenright}{\kern0pt}{\isachardoublequoteclose}\isanewline
\isanewline
\isakeywordONE{definition}\isamarkupfalse%
\ Thread{\isadigit{5}}{\isacharunderscore}{\kern0pt}sem\ \isakeywordTWO{where}\ \isanewline
{\isachardoublequoteopen}Thread{\isadigit{5}}{\isacharunderscore}{\kern0pt}sem\ {\isasymequiv}\ \ {\isacharparenleft}{\kern0pt}\ \isactrlbold {\isacharbar}{\kern0pt}\isactrlbold {\isacharbar}{\kern0pt}\isactrlbold {\isacharbar}{\kern0pt}\ idx\ {\isasymin}{\isacharhash}{\kern0pt}\ mset\ {\isacharbrackleft}{\kern0pt}{\isadigit{1}}{\isachardot}{\kern0pt}{\isachardot}{\kern0pt}{\isadigit{4}}{\isacharbrackright}{\kern0pt}{\isachardot}{\kern0pt}\ semaphore\ idx\ {\isacharparenright}{\kern0pt}\ \ \isanewline
\ \ \ \ \ \ \ \ \ \ \ \ \ \ \ \ \ \isanewline
\ \ \ \ \ \ \ \ \ \ \ \ \ \ \ \ \ \ {\isacharbar}{\kern0pt}{\isacharbar}{\kern0pt}\isanewline
\ \isanewline
\ \ \ \ \ \ \ \ \ \ \ \ \ \ \ \ \ \ \ Sem\ Thread{\isadigit{5}}{\isachardoublequoteclose}\isanewline
\isanewline
\isakeywordONE{lemma}\isamarkupfalse%
\ Thread{\isadigit{5}}{\isacharunderscore}{\kern0pt}while{\isacharunderscore}{\kern0pt}rec\ {\isacharcolon}{\kern0pt}\ {\isachardoublequoteopen}{\isacharparenleft}{\kern0pt}{\isasymmu}\ x{\isacharcolon}{\kern0pt}{\isacharcolon}{\kern0pt}{\isacharprime}{\kern0pt}b\ {\isasymRightarrow}\ evs\ process{\isachardot}{\kern0pt}\ {\isacharparenleft}{\kern0pt}{\isasymlambda}{\isasymsigma}{\isachardot}{\kern0pt}\ evs{\isachardot}{\kern0pt}lock\ {\isacharparenleft}{\kern0pt}{\isadigit{1}}{\isacharcolon}{\kern0pt}{\isacharcolon}{\kern0pt}int{\isacharparenright}{\kern0pt}\ {\isasymrightarrow}\ x\ {\isasymsigma}{\isacharparenright}{\kern0pt}{\isacharparenright}{\kern0pt}\ A\ \isanewline
\ \ \ \ \ \ \ \ \ \ \ \ \ {\isacharequal}{\kern0pt}\ evs{\isachardot}{\kern0pt}lock\ {\isadigit{1}}\ {\isasymrightarrow}\ {\isacharparenleft}{\kern0pt}{\isasymmu}\ x{\isachardot}{\kern0pt}\ {\isacharparenleft}{\kern0pt}{\isasymlambda}{\isasymsigma}{\isachardot}{\kern0pt}\ evs{\isachardot}{\kern0pt}lock\ {\isacharparenleft}{\kern0pt}{\isadigit{1}}{\isacharcolon}{\kern0pt}{\isacharcolon}{\kern0pt}int{\isacharparenright}{\kern0pt}\ {\isasymrightarrow}\ x\ {\isasymsigma}{\isacharparenright}{\kern0pt}{\isacharparenright}{\kern0pt}\ A{\isachardoublequoteclose}\isanewline
%
\isadelimproof
\ \ %
\endisadelimproof
%
\isatagproof
\isakeywordONE{by}\isamarkupfalse%
{\isacharparenleft}{\kern0pt}subst\ cont{\isacharunderscore}{\kern0pt}process{\isacharunderscore}{\kern0pt}rec{\isacharbrackleft}{\kern0pt}\isakeywordTWO{where}\ P\ {\isacharequal}{\kern0pt}\ {\isacartoucheopen}{\isacharparenleft}{\kern0pt}{\isasymmu}\ x{\isachardot}{\kern0pt}\ {\isacharparenleft}{\kern0pt}{\isasymlambda}{\isasymsigma}{\isachardot}{\kern0pt}\ evs{\isachardot}{\kern0pt}lock\ {\isacharparenleft}{\kern0pt}{\isadigit{1}}{\isacharcolon}{\kern0pt}{\isacharcolon}{\kern0pt}int{\isacharparenright}{\kern0pt}\ {\isasymrightarrow}\ x\ {\isasymsigma}{\isacharparenright}{\kern0pt}{\isacharparenright}{\kern0pt}{\isacartoucheclose}{\isacharcomma}{\kern0pt}\ OF\ refl{\isacharbrackright}{\kern0pt}{\isacharparenright}{\kern0pt}\ simp{\isacharunderscore}{\kern0pt}all%
\endisatagproof
{\isafoldproof}%
%
\isadelimproof
%
\endisadelimproof
%
\begin{isamarkuptext}%
This process system will deadlock after the first attempt to get \isa{\isaconst{com{\isachardot}{\kern0pt}lock}\ {\isadigit{0}}}.
This can be formally proven in HOL-CSP by stating: 

\isa{{\isasymnot}\ deadlock{\isacharunderscore}{\kern0pt}free\ Thread{\isadigit{5}}{\isacharunderscore}{\kern0pt}sem}.

The proof is shown in the following subsection.%
\end{isamarkuptext}\isamarkuptrue%
%
\isadelimdocument
%
\endisadelimdocument
%
\isatagdocument
%
\isamarkupsubsection{A Technique to rewrite Interleaves to MPrefixes%
}
\isamarkuptrue%
%
\endisatagdocument
{\isafolddocument}%
%
\isadelimdocument
%
\endisadelimdocument
%
\begin{isamarkuptext}%
Note that the following description is not necessarily the most effective manner to establish
deadlock-freeness via fixed-point induction; it is  clear that it will not scale up to large
process systems. (For a more powerful and automated approach, we refer to the handling via Proc-Omata 
\cite{BallenghienW24} or the approach in \cite{ITP2026}). Our proceeding has, though, the advantage 
to use only very basic means and is therefore a self-contained demonstration.%
\end{isamarkuptext}\isamarkuptrue%
%
\begin{isamarkuptext}%
A key obstacle for the establishment of deadlock-properties is that we need to bring
both sides of the parallel composition operator \isa{{\isacharunderscore}{\kern0pt}\ {\isacharbar}{\kern0pt}{\isacharbar}{\kern0pt}\ {\isacharunderscore}{\kern0pt}} into the form \isa{\isaconst{Mprefix}\ \isafree{A}\ \isafree{P}} in order
to apply the rule:

 %
\begin{isabelle}%
\ \ \ \ \ \ \ \ \ \ \isaconst{Mprefix}\ \isavar{{\isacharquery}{\kern0pt}A}\ \isavar{{\isacharquery}{\kern0pt}P}\ {\isacharbar}{\kern0pt}{\isacharbar}{\kern0pt}\ \isaconst{Mprefix}\ \isavar{{\isacharquery}{\kern0pt}B}\ \isavar{{\isacharquery}{\kern0pt}Q}\ {\isacharequal}{\kern0pt}\ {\isasymbox}\isabound{x}{\isasymin}{\isacharparenleft}{\kern0pt}\isavar{{\isacharquery}{\kern0pt}A}\ {\isasyminter}\ \isavar{{\isacharquery}{\kern0pt}B}{\isacharparenright}{\kern0pt}\ {\isasymrightarrow}\ {\isacharparenleft}{\kern0pt}\isavar{{\isacharquery}{\kern0pt}P}\ \isabound{x}\ {\isacharbar}{\kern0pt}{\isacharbar}{\kern0pt}\ \isavar{{\isacharquery}{\kern0pt}Q}\ \isabound{x}{\isacharparenright}{\kern0pt}%
\end{isabelle} 

The  \isa{\isaconst{Mprefix}}-terms are typically
quite large if-then-else cascades representing some form of decision diagram that collapse to 
small terms after synchronization.%
\end{isamarkuptext}\isamarkuptrue%
\isakeywordONE{lemma}\isamarkupfalse%
\ inter{\isadigit{2}}Mprefix{\isadigit{1}}{\isacharcolon}{\kern0pt}\ \isanewline
\ \ \isakeywordTWO{assumes}\ \ {\isachardoublequoteopen}a\ {\isasymnoteq}\ b{\isachardoublequoteclose}\ \isanewline
\ \ \isakeywordTWO{shows}\ {\isachardoublequoteopen}{\isacharparenleft}{\kern0pt}a\ {\isasymrightarrow}\ P\ {\isacharbar}{\kern0pt}{\isacharbar}{\kern0pt}{\isacharbar}{\kern0pt}\ b\ {\isasymrightarrow}\ Q{\isacharparenright}{\kern0pt}\ \isanewline
\ \ \ \ \ \ \ \ \ \ {\isacharequal}{\kern0pt}\ Mprefix\ {\isacharbraceleft}{\kern0pt}a{\isacharcomma}{\kern0pt}b{\isacharbraceright}{\kern0pt}\ {\isacharparenleft}{\kern0pt}{\isasymlambda}x{\isachardot}{\kern0pt}\ if\ x\ {\isacharequal}{\kern0pt}\ a\ then\ {\isacharparenleft}{\kern0pt}P\ {\isacharbar}{\kern0pt}{\isacharbar}{\kern0pt}{\isacharbar}{\kern0pt}\ b\ {\isasymrightarrow}\ Q{\isacharparenright}{\kern0pt}\ \isanewline
\ \ \ \ \ \ \ \ \ \ \ \ \ \ \ \ \ \ \ \ \ \ \ \ \ \ \ \ \ \ \ \ \ \ \ \ \ \ \ \ else\ {\isacharparenleft}{\kern0pt}a\ {\isasymrightarrow}\ P{\isacharparenright}{\kern0pt}\ {\isacharbar}{\kern0pt}{\isacharbar}{\kern0pt}{\isacharbar}{\kern0pt}\ Q\ {\isacharparenright}{\kern0pt}{\isachardoublequoteclose}\isanewline
%
\isadelimproof
\ \ %
\endisadelimproof
%
\isatagproof
\isakeywordONE{apply}\isamarkupfalse%
{\isacharparenleft}{\kern0pt}simp\ add{\isacharcolon}{\kern0pt}write{\isadigit{0}}{\isacharunderscore}{\kern0pt}Inter{\isacharunderscore}{\kern0pt}write{\isadigit{0}}{\isacharparenright}{\kern0pt}\isanewline
\ \ \isakeywordONE{by}\isamarkupfalse%
{\isacharparenleft}{\kern0pt}subst\ {\isacharparenleft}{\kern0pt}{\isadigit{3}}{\isacharparenright}{\kern0pt}\ write{\isadigit{0}}{\isacharunderscore}{\kern0pt}def{\isacharcomma}{\kern0pt}subst\ write{\isadigit{0}}{\isacharunderscore}{\kern0pt}Det{\isacharunderscore}{\kern0pt}Mprefix{\isacharcomma}{\kern0pt}simp\ add{\isacharcolon}{\kern0pt}\ assms{\isacharparenright}{\kern0pt}%
\endisatagproof
{\isafoldproof}%
%
\isadelimproof
\isanewline
%
\endisadelimproof
\isanewline
\isanewline
\isakeywordONE{lemma}\isamarkupfalse%
\ inter{\isadigit{2}}Mprefix{\isadigit{2}}{\isacharcolon}{\kern0pt}\ \isanewline
\ \ \isakeywordTWO{assumes}\ {\isacharasterisk}{\kern0pt}\ {\isacharcolon}{\kern0pt}\ {\isachardoublequoteopen}a\ {\isasymnotin}\ A{\isachardoublequoteclose}\isanewline
\ \ \isakeywordTWO{shows}\ {\isachardoublequoteopen}{\isacharparenleft}{\kern0pt}{\isacharparenleft}{\kern0pt}a\ {\isasymrightarrow}\ P{\isacharparenright}{\kern0pt}\ {\isacharbar}{\kern0pt}{\isacharbar}{\kern0pt}{\isacharbar}{\kern0pt}\ Mprefix\ A\ Q{\isacharparenright}{\kern0pt}\ \isanewline
\ \ \ \ \ \ \ \ \ \ {\isacharequal}{\kern0pt}\ Mprefix\ {\isacharparenleft}{\kern0pt}insert\ a\ A{\isacharparenright}{\kern0pt}\ {\isacharparenleft}{\kern0pt}{\isasymlambda}x{\isachardot}{\kern0pt}\ if\ x\ {\isacharequal}{\kern0pt}\ a\ then\ {\isacharparenleft}{\kern0pt}P\ {\isacharbar}{\kern0pt}{\isacharbar}{\kern0pt}{\isacharbar}{\kern0pt}\ Mprefix\ A\ Q{\isacharparenright}{\kern0pt}\ \isanewline
\ \ \ \ \ \ \ \ \ \ \ \ \ \ \ \ \ \ \ \ \ \ \ \ \ \ \ \ \ \ \ \ \ \ \ \ \ \ \ \ \ \ \ \ \ \ \ else\ {\isacharparenleft}{\kern0pt}{\isacharparenleft}{\kern0pt}a\ {\isasymrightarrow}\ P{\isacharparenright}{\kern0pt}\ {\isacharbar}{\kern0pt}{\isacharbar}{\kern0pt}{\isacharbar}{\kern0pt}\ Q\ x{\isacharparenright}{\kern0pt}\ {\isacharparenright}{\kern0pt}{\isachardoublequoteclose}\isanewline
%
\isadelimproof
\ \ %
\endisadelimproof
%
\isatagproof
\isakeywordONE{apply}\isamarkupfalse%
{\isacharparenleft}{\kern0pt}subst\ write{\isadigit{0}}{\isacharunderscore}{\kern0pt}def{\isacharcomma}{\kern0pt}simp\ add{\isacharcolon}{\kern0pt}\ Mprefix{\isacharunderscore}{\kern0pt}Inter{\isacharunderscore}{\kern0pt}Mprefix{\isacharparenright}{\kern0pt}\isanewline
\ \ \isakeywordONE{apply}\isamarkupfalse%
{\isacharparenleft}{\kern0pt}subst\ Mprefix{\isacharunderscore}{\kern0pt}Det{\isacharunderscore}{\kern0pt}Mprefix{\isacharcomma}{\kern0pt}\ simp\ add{\isacharcolon}{\kern0pt}\ assms{\isacharparenright}{\kern0pt}\isanewline
\ \ \isakeywordONE{by}\isamarkupfalse%
{\isacharparenleft}{\kern0pt}fold\ write{\isadigit{0}}{\isacharunderscore}{\kern0pt}def{\isacharcomma}{\kern0pt}\ simp{\isacharparenright}{\kern0pt}%
\endisatagproof
{\isafoldproof}%
%
\isadelimproof
\isanewline
%
\endisadelimproof
\isanewline
\isakeywordONE{lemma}\isamarkupfalse%
\ inter{\isadigit{2}}Mprefix{\isadigit{3}}{\isacharcolon}{\kern0pt}\isanewline
\ \ {\isachardoublequoteopen}a\ {\isasymnoteq}\ b\ {\isasymLongrightarrow}\ {\isacharparenleft}{\kern0pt}{\isacharparenleft}{\kern0pt}a\ {\isasymrightarrow}\ P\ a{\isacharparenright}{\kern0pt}\ {\isasymbox}\ {\isacharparenleft}{\kern0pt}b\ {\isasymrightarrow}\ Q\ b{\isacharparenright}{\kern0pt}{\isacharparenright}{\kern0pt}\ \isanewline
\ \ \ \ \ \ \ \ \ \ \ \ \ {\isacharequal}{\kern0pt}\ Mprefix\ {\isacharbraceleft}{\kern0pt}a{\isacharcomma}{\kern0pt}b{\isacharbraceright}{\kern0pt}\ {\isacharparenleft}{\kern0pt}{\isasymlambda}x{\isachardot}{\kern0pt}\ if\ x{\isacharequal}{\kern0pt}\ a\ then\ P\ a\ else\ Q\ b{\isacharparenright}{\kern0pt}{\isachardoublequoteclose}\isanewline
%
\isadelimproof
\ \ %
\endisadelimproof
%
\isatagproof
\isakeywordONE{apply}\isamarkupfalse%
{\isacharparenleft}{\kern0pt}simp\ add{\isacharcolon}{\kern0pt}\ Mprefix{\isacharunderscore}{\kern0pt}singl{\isacharbrackleft}{\kern0pt}symmetric{\isacharbrackright}{\kern0pt}\ {\isacharparenright}{\kern0pt}\isanewline
\ \ \isakeywordONE{by}\isamarkupfalse%
\ {\isacharparenleft}{\kern0pt}smt\ {\isacharparenleft}{\kern0pt}verit{\isacharcomma}{\kern0pt}\ ccfv{\isacharunderscore}{\kern0pt}threshold{\isacharparenright}{\kern0pt}\ Mprefix{\isacharunderscore}{\kern0pt}Un{\isacharunderscore}{\kern0pt}distrib\ Mprefix{\isacharunderscore}{\kern0pt}singl\ \isanewline
\ \ \ \ \ \ \ \ \ \ Un{\isacharunderscore}{\kern0pt}insert{\isacharunderscore}{\kern0pt}right\ insert{\isacharunderscore}{\kern0pt}commute\ sup{\isacharunderscore}{\kern0pt}bot{\isachardot}{\kern0pt}right{\isacharunderscore}{\kern0pt}neutral{\isacharparenright}{\kern0pt}%
\endisatagproof
{\isafoldproof}%
%
\isadelimproof
%
\endisadelimproof
%
\begin{isamarkuptext}%
After these generalities, we turn to the core lemmas relevant to the subsequent proofs.
They construct compact normalforms for states (represented by CSPM-expressions) linked
by the events \isa{\isaconst{evs{\isachardot}{\kern0pt}lock}\ {\isadigit{1}}} and  \isa{\isaconst{evs{\isachardot}{\kern0pt}unlock}\ {\isadigit{1}}}. They proceed by blowing both sides
up into a lower-level normalform in which their equality can be demonstrated.%
\end{isamarkuptext}\isamarkuptrue%
\isakeywordONE{lemma}\isamarkupfalse%
\ contractMultiInter{\isadigit{1}}\ {\isacharcolon}{\kern0pt}\ \isanewline
\ \ {\isachardoublequoteopen}MultiInter\ {\isacharparenleft}{\kern0pt}mset\ {\isacharbrackleft}{\kern0pt}{\isadigit{1}}{\isachardot}{\kern0pt}{\isachardot}{\kern0pt}{\isadigit{4}}{\isacharbrackright}{\kern0pt}{\isacharparenright}{\kern0pt}\ semaphore\ {\isacharbar}{\kern0pt}{\isacharbar}{\kern0pt}\ evs{\isachardot}{\kern0pt}lock\ {\isadigit{1}}\ {\isasymrightarrow}\ P\ \isanewline
\ \ \ {\isacharequal}{\kern0pt}\ \ evs{\isachardot}{\kern0pt}lock\ {\isadigit{1}}\ {\isasymrightarrow}\ {\isacharparenleft}{\kern0pt}{\isacharparenleft}{\kern0pt}evs{\isachardot}{\kern0pt}unlock\ {\isadigit{1}}\ {\isasymrightarrow}\ semaphore\ {\isadigit{1}}\ \ {\isacharbar}{\kern0pt}{\isacharbar}{\kern0pt}{\isacharbar}{\kern0pt}\ MultiInter\ {\isacharparenleft}{\kern0pt}mset\ {\isacharbrackleft}{\kern0pt}{\isadigit{2}}{\isachardot}{\kern0pt}{\isachardot}{\kern0pt}{\isadigit{4}}{\isacharbrackright}{\kern0pt}{\isacharparenright}{\kern0pt}\ semaphore{\isacharparenright}{\kern0pt}\ {\isacharbar}{\kern0pt}{\isacharbar}{\kern0pt}\ P{\isacharparenright}{\kern0pt}{\isachardoublequoteclose}\isanewline
\ \ {\isacharparenleft}{\kern0pt}\isakeywordTWO{is}\ {\isachardoublequoteopen}{\isacharquery}{\kern0pt}lhs\ {\isacharequal}{\kern0pt}\ evs{\isachardot}{\kern0pt}lock\ {\isadigit{1}}\ {\isasymrightarrow}\ {\isacharquery}{\kern0pt}rhs{\isachardoublequoteclose}{\isacharparenright}{\kern0pt}\isanewline
%
\isadelimproof
%
\endisadelimproof
%
\isatagproof
\isakeywordONE{proof}\isamarkupfalse%
\ {\isacharminus}{\kern0pt}\isanewline
\ \ \isakeywordONE{have}\isamarkupfalse%
\ {\isacharasterisk}{\kern0pt}\ {\isacharcolon}{\kern0pt}\ {\isachardoublequoteopen}{\isacharquery}{\kern0pt}lhs\ {\isacharequal}{\kern0pt}\ evs{\isachardot}{\kern0pt}lock\ {\isadigit{1}}\ {\isasymrightarrow}\isanewline
\ \ \ \ \ \ \ \ \ \ \ \ \ \ \ \ \ \ \ \ {\isacharparenleft}{\kern0pt}{\isacharparenleft}{\kern0pt}evs{\isachardot}{\kern0pt}unlock\ {\isadigit{1}}\ {\isasymrightarrow}\ semaphore\ {\isadigit{1}}\ \isanewline
\ \ \ \ \ \ \ \ \ \ \ \ \ \ \ \ \ \ \ \ \ {\isacharbar}{\kern0pt}{\isacharbar}{\kern0pt}{\isacharbar}{\kern0pt}\ {\isasymbox}x{\isasymin}{\isacharbraceleft}{\kern0pt}evs{\isachardot}{\kern0pt}lock\ {\isadigit{2}}{\isacharcomma}{\kern0pt}\ evs{\isachardot}{\kern0pt}lock\ {\isadigit{3}}{\isacharcomma}{\kern0pt}\ evs{\isachardot}{\kern0pt}lock\ {\isadigit{4}}{\isacharbraceright}{\kern0pt}\isanewline
\ \ \ \ \ \ \ \ \ \ \ \ \ \ \ \ \ \ \ \ \ \ \ \ \ \ \ \ \ \ \ \ {\isasymrightarrow}\ {\isacharparenleft}{\kern0pt}if\ x\ {\isacharequal}{\kern0pt}\ evs{\isachardot}{\kern0pt}lock\ {\isadigit{2}}\isanewline
\ \ \ \ \ \ \ \ \ \ \ \ \ \ \ \ \ \ \ \ \ \ \ \ \ \ \ \ \ \ \ \ \ \ \ \ then\ evs{\isachardot}{\kern0pt}unlock\ {\isadigit{2}}\ {\isasymrightarrow}\ semaphore\ {\isadigit{2}}\ \isanewline
\ \ \ \ \ \ \ \ \ \ \ \ \ \ \ \ \ \ \ \ \ \ \ \ \ \ \ \ \ \ \ \ \ \ \ \ \ \ \ \ \ {\isacharbar}{\kern0pt}{\isacharbar}{\kern0pt}{\isacharbar}{\kern0pt}\ \isanewline
\ \ \ \ \ \ \ \ \ \ \ \ \ \ \ \ \ \ \ \ \ \ \ \ \ \ \ \ \ \ \ \ \ \ \ \ \ \ \ \ \ {\isasymbox}x{\isasymin}{\isacharbraceleft}{\kern0pt}evs{\isachardot}{\kern0pt}lock\ {\isadigit{3}}{\isacharcomma}{\kern0pt}\ evs{\isachardot}{\kern0pt}lock\ {\isadigit{4}}{\isacharbraceright}{\kern0pt}\isanewline
\ \ \ \ \ \ \ \ \ \ \ \ \ \ \ \ \ \ \ \ \ \ \ \ \ \ \ \ \ \ \ \ \ \ \ \ \ \ \ \ \ \ \ \ \ \ \ {\isasymrightarrow}\ {\isacharparenleft}{\kern0pt}if\ x\ {\isacharequal}{\kern0pt}\ evs{\isachardot}{\kern0pt}lock\ {\isadigit{3}}\ \isanewline
\ \ \ \ \ \ \ \ \ \ \ \ \ \ \ \ \ \ \ \ \ \ \ \ \ \ \ \ \ \ \ \ \ \ \ \ \ \ \ \ \ \ \ \ \ \ \ \ \ \ \ then\ evs{\isachardot}{\kern0pt}unlock\ {\isadigit{3}}\ {\isasymrightarrow}\ semaphore\ {\isadigit{3}}\ {\isacharbar}{\kern0pt}{\isacharbar}{\kern0pt}{\isacharbar}{\kern0pt}\ semaphore\ {\isadigit{4}}\isanewline
\ \ \ \ \ \ \ \ \ \ \ \ \ \ \ \ \ \ \ \ \ \ \ \ \ \ \ \ \ \ \ \ \ \ \ \ \ \ \ \ \ \ \ \ \ \ \ \ \ \ \ else\ semaphore\ {\isadigit{3}}\ {\isacharbar}{\kern0pt}{\isacharbar}{\kern0pt}{\isacharbar}{\kern0pt}\ evs{\isachardot}{\kern0pt}unlock\ {\isadigit{4}}\ {\isasymrightarrow}\ semaphore\ {\isadigit{4}}{\isacharparenright}{\kern0pt}\isanewline
\ \ \ \ \ \ \ \ \ \ \ \ \ \ \ \ \ \ \ \ \ \ \ \ \ \ \ \ \ \ \ \ \ \ \ \ else\ semaphore\ {\isadigit{2}}\ \isanewline
\ \ \ \ \ \ \ \ \ \ \ \ \ \ \ \ \ \ \ \ \ \ \ \ \ \ \ \ \ \ \ \ \ \ \ \ \ \ \ \ \ {\isacharbar}{\kern0pt}{\isacharbar}{\kern0pt}{\isacharbar}{\kern0pt}\ \isanewline
\ \ \ \ \ \ \ \ \ \ \ \ \ \ \ \ \ \ \ \ \ \ \ \ \ \ \ \ \ \ \ \ \ \ \ \ \ \ \ \ \ {\isacharparenleft}{\kern0pt}if\ x\ {\isacharequal}{\kern0pt}\ evs{\isachardot}{\kern0pt}lock\ {\isadigit{3}}\ \isanewline
\ \ \ \ \ \ \ \ \ \ \ \ \ \ \ \ \ \ \ \ \ \ \ \ \ \ \ \ \ \ \ \ \ \ \ \ \ \ \ \ \ \ then\ evs{\isachardot}{\kern0pt}unlock\ {\isadigit{3}}\ {\isasymrightarrow}\ semaphore\ {\isadigit{3}}\ {\isacharbar}{\kern0pt}{\isacharbar}{\kern0pt}{\isacharbar}{\kern0pt}\ semaphore\ {\isadigit{4}}\isanewline
\ \ \ \ \ \ \ \ \ \ \ \ \ \ \ \ \ \ \ \ \ \ \ \ \ \ \ \ \ \ \ \ \ \ \ \ \ \ \ \ \ \ else\ semaphore\ {\isadigit{3}}\ {\isacharbar}{\kern0pt}{\isacharbar}{\kern0pt}{\isacharbar}{\kern0pt}\ evs{\isachardot}{\kern0pt}unlock\ {\isadigit{4}}\ {\isasymrightarrow}\ semaphore\ {\isadigit{4}}{\isacharparenright}{\kern0pt}{\isacharparenright}{\kern0pt}{\isacharparenright}{\kern0pt}\ {\isacharbar}{\kern0pt}{\isacharbar}{\kern0pt}\ P{\isacharparenright}{\kern0pt}{\isachardoublequoteclose}\isanewline
\ \ \ \ {\isacharparenleft}{\kern0pt}\isakeywordTWO{is}\ {\isachardoublequoteopen}{\isacharunderscore}{\kern0pt}\ {\isacharequal}{\kern0pt}\ evs{\isachardot}{\kern0pt}lock\ {\isadigit{1}}\ {\isasymrightarrow}\ {\isacharquery}{\kern0pt}rhs{\isacharprime}{\kern0pt}{\isachardoublequoteclose}{\isacharparenright}{\kern0pt}\isanewline
\ \ \ \ \ \ \ \ \ \ \ \isakeywordONE{apply}\isamarkupfalse%
{\isacharparenleft}{\kern0pt}rule\ trans{\isacharparenright}{\kern0pt}\isanewline
\ \ \ \ \ \ \ \ \ \ \ \isakeywordONE{apply}\isamarkupfalse%
{\isacharparenleft}{\kern0pt}simp\ add{\isacharcolon}{\kern0pt}\ List{\isachardot}{\kern0pt}upto{\isachardot}{\kern0pt}simps{\isacharparenright}{\kern0pt}\isanewline
\ \ \ \ \ \ \ \ \ \ \ \isakeywordONE{apply}\isamarkupfalse%
{\isacharparenleft}{\kern0pt}subst\ sema{\isacharunderscore}{\kern0pt}rec{\isacharbrackleft}{\kern0pt}\isakeywordTWO{where}\ n{\isacharequal}{\kern0pt}{\isacartoucheopen}{\isadigit{1}}{\isacartoucheclose}{\isacharbrackright}{\kern0pt}{\isacharcomma}{\kern0pt}\isanewline
\ \ \ \ \ \ \ \ \ \ \ \ \ \ \ \ \ subst\ sema{\isacharunderscore}{\kern0pt}rec{\isacharbrackleft}{\kern0pt}\isakeywordTWO{where}\ n{\isacharequal}{\kern0pt}{\isacartoucheopen}{\isadigit{2}}{\isacartoucheclose}{\isacharbrackright}{\kern0pt}{\isacharcomma}{\kern0pt}\isanewline
\ \ \ \ \ \ \ \ \ \ \ \ \ \ \ \ \ subst\ sema{\isacharunderscore}{\kern0pt}rec{\isacharbrackleft}{\kern0pt}\isakeywordTWO{where}\ n{\isacharequal}{\kern0pt}{\isacartoucheopen}{\isadigit{3}}{\isacartoucheclose}{\isacharbrackright}{\kern0pt}{\isacharcomma}{\kern0pt}\isanewline
\ \ \ \ \ \ \ \ \ \ \ \ \ \ \ \ \ subst\ sema{\isacharunderscore}{\kern0pt}rec{\isacharbrackleft}{\kern0pt}\isakeywordTWO{where}\ n{\isacharequal}{\kern0pt}{\isacartoucheopen}{\isadigit{4}}{\isacartoucheclose}{\isacharbrackright}{\kern0pt}{\isacharparenright}{\kern0pt}\ \isanewline
\ \ \ \ \ \ \ \ \ \ \ \isakeywordONE{apply}\isamarkupfalse%
{\isacharparenleft}{\kern0pt}rule\ trans{\isacharparenright}{\kern0pt}\isanewline
\ \ \ \ \ \ \ \ \ \ \ \ \isakeywordONE{apply}\isamarkupfalse%
{\isacharparenleft}{\kern0pt}simp\ add{\isacharcolon}{\kern0pt}\ inter{\isadigit{2}}Mprefix{\isadigit{1}}\ inter{\isadigit{2}}Mprefix{\isadigit{2}}{\isacharparenright}{\kern0pt}\isanewline
\ \ \ \ \ \ \ \ \ \ \ \ \ \isakeywordONE{apply}\isamarkupfalse%
{\isacharparenleft}{\kern0pt}simp\ add{\isacharcolon}{\kern0pt}\ sema{\isacharunderscore}{\kern0pt}rec{\isacharbrackleft}{\kern0pt}\isakeywordTWO{where}\ n{\isacharequal}{\kern0pt}{\isacartoucheopen}{\isadigit{1}}{\isacartoucheclose}{\isacharcomma}{\kern0pt}symmetric{\isacharbrackright}{\kern0pt}\ \isanewline
\ \ \ \ \ \ \ \ \ \ \ \ \ \ \ \ \ \ \ \ \ \ \ \ \ \ \ \ \ sema{\isacharunderscore}{\kern0pt}rec{\isacharbrackleft}{\kern0pt}\isakeywordTWO{where}\ n{\isacharequal}{\kern0pt}{\isacartoucheopen}{\isadigit{2}}{\isacartoucheclose}{\isacharcomma}{\kern0pt}symmetric{\isacharbrackright}{\kern0pt}\ \isanewline
\ \ \ \ \ \ \ \ \ \ \ \ \ \ \ \ \ \ \ \ \ \ \ \ \ \ \ \ \ sema{\isacharunderscore}{\kern0pt}rec{\isacharbrackleft}{\kern0pt}\isakeywordTWO{where}\ n{\isacharequal}{\kern0pt}{\isacartoucheopen}{\isadigit{3}}{\isacartoucheclose}{\isacharcomma}{\kern0pt}symmetric{\isacharbrackright}{\kern0pt}\ \isanewline
\ \ \ \ \ \ \ \ \ \ \ \ \ \ \ \ \ \ \ \ \ \ \ \ \ \ \ \ \ sema{\isacharunderscore}{\kern0pt}rec{\isacharbrackleft}{\kern0pt}\isakeywordTWO{where}\ n{\isacharequal}{\kern0pt}{\isacartoucheopen}{\isadigit{4}}{\isacartoucheclose}{\isacharcomma}{\kern0pt}symmetric{\isacharbrackright}{\kern0pt}\ \isanewline
\ \ \ \ \ \ \ \ \ \ \ \ \ \ \ \ \ \ \ \ \ \ \ \ cong{\isacharcolon}{\kern0pt}\ HOL{\isachardot}{\kern0pt}if{\isacharunderscore}{\kern0pt}cong{\isacharparenright}{\kern0pt}\isanewline
\ \ \ \ \ \ \ \ \ \ \ \ \ \isakeywordONE{apply}\isamarkupfalse%
{\isacharparenleft}{\kern0pt}subst\ write{\isadigit{0}}{\isacharunderscore}{\kern0pt}def{\isacharbrackleft}{\kern0pt}THEN\ meta{\isacharunderscore}{\kern0pt}eq{\isacharunderscore}{\kern0pt}to{\isacharunderscore}{\kern0pt}obj{\isacharunderscore}{\kern0pt}eq{\isacharcomma}{\kern0pt}\ of{\isachardoublequoteopen}evs{\isachardot}{\kern0pt}lock\ {\isacharparenleft}{\kern0pt}{\isadigit{1}}{\isacharcolon}{\kern0pt}{\isacharcolon}{\kern0pt}int{\isacharparenright}{\kern0pt}{\isachardoublequoteclose}{\isacharbrackright}{\kern0pt}{\isacharparenright}{\kern0pt}\isanewline
\ \ \ \ \ \ \ \ \ \ \ \ \ \isakeywordONE{by}\isamarkupfalse%
{\isacharparenleft}{\kern0pt}subst\ Mprefix{\isacharunderscore}{\kern0pt}Par{\isacharunderscore}{\kern0pt}Mprefix{\isacharcomma}{\kern0pt}\ simp\ add{\isacharcolon}{\kern0pt}\ Mprefix{\isacharunderscore}{\kern0pt}singl{\isacharparenright}{\kern0pt}\ %
\isamarkupcmt{et bang!%
}\isanewline
\ \ \isakeywordTHREE{define}\isamarkupfalse%
\ X\ \isakeywordTWO{where}\ {\isachardoublequoteopen}X\ {\isacharequal}{\kern0pt}\ {\isacharquery}{\kern0pt}rhs{\isacharprime}{\kern0pt}{\isachardoublequoteclose}\isanewline
\ \ \isakeywordONE{have}\isamarkupfalse%
\ {\isacharasterisk}{\kern0pt}{\isacharasterisk}{\kern0pt}\ {\isacharcolon}{\kern0pt}\ {\isachardoublequoteopen}{\isacharquery}{\kern0pt}rhs\ {\isacharequal}{\kern0pt}\ X{\isachardoublequoteclose}\ \isanewline
\ \ \ \ \ \ \ \ \ \ \ \isakeywordONE{apply}\isamarkupfalse%
{\isacharparenleft}{\kern0pt}subst\ write{\isadigit{0}}{\isacharunderscore}{\kern0pt}def{\isacharbrackleft}{\kern0pt}THEN\ meta{\isacharunderscore}{\kern0pt}eq{\isacharunderscore}{\kern0pt}to{\isacharunderscore}{\kern0pt}obj{\isacharunderscore}{\kern0pt}eq{\isacharbrackright}{\kern0pt}{\isacharparenright}{\kern0pt}\isanewline
\ \ \ \ \ \ \ \ \ \ \ \isakeywordONE{apply}\isamarkupfalse%
{\isacharparenleft}{\kern0pt}simp\ add{\isacharcolon}{\kern0pt}\ List{\isachardot}{\kern0pt}upto{\isachardot}{\kern0pt}simps{\isacharparenright}{\kern0pt}\ \ \isanewline
\ \ \ \ \ \ \ \ \ \ \ \isakeywordONE{apply}\isamarkupfalse%
{\isacharparenleft}{\kern0pt}subst\ sema{\isacharunderscore}{\kern0pt}rec{\isacharbrackleft}{\kern0pt}\isakeywordTWO{where}\ n{\isacharequal}{\kern0pt}{\isacartoucheopen}{\isadigit{1}}{\isacartoucheclose}{\isacharbrackright}{\kern0pt}{\isacharcomma}{\kern0pt}subst\ sema{\isacharunderscore}{\kern0pt}rec{\isacharbrackleft}{\kern0pt}\isakeywordTWO{where}\ n{\isacharequal}{\kern0pt}{\isacartoucheopen}{\isadigit{2}}{\isacartoucheclose}{\isacharbrackright}{\kern0pt}{\isacharcomma}{\kern0pt}\isanewline
\ \ \ \ \ \ \ \ \ \ \ \ \ \ \ \ \ subst\ sema{\isacharunderscore}{\kern0pt}rec{\isacharbrackleft}{\kern0pt}\isakeywordTWO{where}\ n{\isacharequal}{\kern0pt}{\isacartoucheopen}{\isadigit{3}}{\isacartoucheclose}{\isacharbrackright}{\kern0pt}{\isacharcomma}{\kern0pt}subst\ sema{\isacharunderscore}{\kern0pt}rec{\isacharbrackleft}{\kern0pt}\isakeywordTWO{where}\ n{\isacharequal}{\kern0pt}{\isacartoucheopen}{\isadigit{4}}{\isacartoucheclose}{\isacharbrackright}{\kern0pt}{\isacharparenright}{\kern0pt}\ \isanewline
\ \ \ \ \ \ \ \ \ \ \ \isakeywordONE{apply}\isamarkupfalse%
{\isacharparenleft}{\kern0pt}simp\ add{\isacharcolon}{\kern0pt}\ inter{\isadigit{2}}Mprefix{\isadigit{1}}\ inter{\isadigit{2}}Mprefix{\isadigit{2}}{\isacharparenright}{\kern0pt}\ \isanewline
\ \ \ \ \ \ \ \ \ \ \ \isakeywordONE{by}\isamarkupfalse%
\ {\isacharparenleft}{\kern0pt}metis\ X{\isacharunderscore}{\kern0pt}def\ sema{\isacharunderscore}{\kern0pt}rec\ Mprefix{\isacharunderscore}{\kern0pt}singl{\isacharparenright}{\kern0pt}\isanewline
\ \ \isakeywordTHREE{show}\isamarkupfalse%
\ {\isacharquery}{\kern0pt}thesis\isanewline
\ \ \ \ \isakeywordONE{by}\isamarkupfalse%
\ {\isacharparenleft}{\kern0pt}metis\ {\isacharasterisk}{\kern0pt}\ {\isacharasterisk}{\kern0pt}{\isacharasterisk}{\kern0pt}\ X{\isacharunderscore}{\kern0pt}def{\isacharparenright}{\kern0pt}\isanewline
\isakeywordONE{qed}\isamarkupfalse%
%
\endisatagproof
{\isafoldproof}%
%
\isadelimproof
\isanewline
%
\endisadelimproof
\isanewline
\isanewline
\isakeywordONE{lemma}\isamarkupfalse%
\ contractMultiInter{\isadigit{2}}{\isacharcolon}{\kern0pt}\ \isanewline
\ \ \ \ \ \ {\isachardoublequoteopen}\ \ {\isacharparenleft}{\kern0pt}{\isacharparenleft}{\kern0pt}evs{\isachardot}{\kern0pt}unlock\ {\isadigit{1}}\ {\isasymrightarrow}\ semaphore\ {\isadigit{1}}\ \ {\isacharbar}{\kern0pt}{\isacharbar}{\kern0pt}{\isacharbar}{\kern0pt}\ MultiInter\ {\isacharparenleft}{\kern0pt}mset\ {\isacharbrackleft}{\kern0pt}{\isadigit{2}}{\isachardot}{\kern0pt}{\isachardot}{\kern0pt}{\isadigit{4}}{\isacharbrackright}{\kern0pt}{\isacharparenright}{\kern0pt}\ semaphore{\isacharparenright}{\kern0pt}\ \isanewline
\ \ \ \ \ \ \ \ \ {\isacharbar}{\kern0pt}{\isacharbar}{\kern0pt}\ \ evs{\isachardot}{\kern0pt}lock\ {\isadigit{1}}\ {\isasymrightarrow}\ P{\isacharparenright}{\kern0pt}\ \isanewline
\ \ \ \ \ \ \ {\isacharequal}{\kern0pt}\ Constant{\isacharunderscore}{\kern0pt}Processes{\isachardot}{\kern0pt}STOP{\isachardoublequoteclose}\isanewline
\ \ \ \ \ \ \ \ \ {\isacharparenleft}{\kern0pt}\isakeywordTWO{is}\ {\isachardoublequoteopen}{\isacharparenleft}{\kern0pt}{\isacharquery}{\kern0pt}R\ {\isacharbar}{\kern0pt}{\isacharbar}{\kern0pt}\ {\isacharquery}{\kern0pt}S{\isacharparenright}{\kern0pt}\ {\isacharequal}{\kern0pt}\ {\isacharunderscore}{\kern0pt}{\isachardoublequoteclose}{\isacharparenright}{\kern0pt}\isanewline
%
\isadelimproof
%
\endisadelimproof
%
\isatagproof
\isakeywordONE{proof}\isamarkupfalse%
\ {\isacharminus}{\kern0pt}\ \isanewline
\ \ \isakeywordONE{have}\isamarkupfalse%
\ {\isadigit{1}}\ {\isacharcolon}{\kern0pt}\ {\isachardoublequoteopen}{\isacharquery}{\kern0pt}R\ {\isacharequal}{\kern0pt}\ {\isasymbox}x{\isasymin}{\isacharbraceleft}{\kern0pt}evs{\isachardot}{\kern0pt}unlock\ {\isadigit{1}}{\isacharcomma}{\kern0pt}\ evs{\isachardot}{\kern0pt}lock\ {\isadigit{2}}{\isacharcomma}{\kern0pt}\ evs{\isachardot}{\kern0pt}lock\ {\isadigit{3}}{\isacharcomma}{\kern0pt}\ evs{\isachardot}{\kern0pt}lock\ {\isadigit{4}}{\isacharbraceright}{\kern0pt}\isanewline
\ \ \ \ \ \ \ {\isasymrightarrow}\ {\isacharparenleft}{\kern0pt}if\ x\ {\isacharequal}{\kern0pt}\ evs{\isachardot}{\kern0pt}unlock\ {\isadigit{1}}\isanewline
\ \ \ \ \ \ \ \ \ \ \ then\ semaphore\ {\isadigit{1}}\isanewline
\ \ \ \ \ \ \ \ \ \ \ \ \ \ \ \ {\isacharbar}{\kern0pt}{\isacharbar}{\kern0pt}{\isacharbar}{\kern0pt}\ {\isasymbox}x{\isasymin}{\isacharbraceleft}{\kern0pt}evs{\isachardot}{\kern0pt}lock\ {\isadigit{2}}{\isacharcomma}{\kern0pt}\ evs{\isachardot}{\kern0pt}lock\ {\isadigit{3}}{\isacharcomma}{\kern0pt}\ evs{\isachardot}{\kern0pt}lock\ {\isadigit{4}}{\isacharbraceright}{\kern0pt}\isanewline
\ \ \ \ \ \ \ \ \ \ \ \ \ \ \ \ \ \ \ \ \ \ \ \ \ \ {\isasymrightarrow}\ {\isacharparenleft}{\kern0pt}if\ x\ {\isacharequal}{\kern0pt}\ evs{\isachardot}{\kern0pt}lock\ {\isadigit{2}}\isanewline
\ \ \ \ \ \ \ \ \ \ \ \ \ \ \ \ \ \ \ \ \ \ \ \ \ \ \ \ \ \ then\ evs{\isachardot}{\kern0pt}unlock\ {\isadigit{2}}\ {\isasymrightarrow}\ semaphore\ {\isadigit{2}}\isanewline
\ \ \ \ \ \ \ \ \ \ \ \ \ \ \ \ \ \ \ \ \ \ \ \ \ \ \ \ \ \ {\isacharbar}{\kern0pt}{\isacharbar}{\kern0pt}{\isacharbar}{\kern0pt}\ {\isasymbox}x{\isasymin}{\isacharbraceleft}{\kern0pt}evs{\isachardot}{\kern0pt}lock\ {\isadigit{3}}{\isacharcomma}{\kern0pt}\ evs{\isachardot}{\kern0pt}lock\ {\isadigit{4}}{\isacharbraceright}{\kern0pt}\isanewline
\ \ \ \ \ \ \ \ \ \ \ \ \ \ \ \ \ \ \ \ \ \ \ \ \ \ \ \ \ \ \ \ \ \ \ \ \ \ {\isasymrightarrow}\ {\isacharparenleft}{\kern0pt}if\ x\ {\isacharequal}{\kern0pt}\ evs{\isachardot}{\kern0pt}lock\ {\isadigit{3}}\ \isanewline
\ \ \ \ \ \ \ \ \ \ \ \ \ \ \ \ \ \ \ \ \ \ \ \ \ \ \ \ \ \ \ \ \ \ \ \ \ \ \ \ \ \ then\ evs{\isachardot}{\kern0pt}unlock\ {\isadigit{3}}\ {\isasymrightarrow}\ semaphore\ {\isadigit{3}}\ {\isacharbar}{\kern0pt}{\isacharbar}{\kern0pt}{\isacharbar}{\kern0pt}\ semaphore\ {\isadigit{4}}\isanewline
\ \ \ \ \ \ \ \ \ \ \ \ \ \ \ \ \ \ \ \ \ \ \ \ \ \ \ \ \ \ \ \ \ \ \ \ \ \ \ \ \ \ else\ semaphore\ {\isadigit{3}}\ {\isacharbar}{\kern0pt}{\isacharbar}{\kern0pt}{\isacharbar}{\kern0pt}\ evs{\isachardot}{\kern0pt}unlock\ {\isadigit{4}}\ {\isasymrightarrow}\ semaphore\ {\isadigit{4}}{\isacharparenright}{\kern0pt}\isanewline
\ \ \ \ \ \ \ \ \ \ \ \ \ \ \ \ \ \ \ \ \ \ \ \ \ \ \ \ \ \ else\ semaphore\ \ {\isadigit{2}}\isanewline
\ \ \ \ \ \ \ \ \ \ \ \ \ \ \ \ \ \ \ \ \ \ \ \ \ \ \ \ \ \ \ \ \ \ \ {\isacharbar}{\kern0pt}{\isacharbar}{\kern0pt}{\isacharbar}{\kern0pt}\ {\isacharparenleft}{\kern0pt}if\ x\ {\isacharequal}{\kern0pt}\ evs{\isachardot}{\kern0pt}lock\ {\isadigit{3}}\ \isanewline
\ \ \ \ \ \ \ \ \ \ \ \ \ \ \ \ \ \ \ \ \ \ \ \ \ \ \ \ \ \ \ \ \ \ \ \ \ \ \ \ then\ evs{\isachardot}{\kern0pt}unlock\ {\isadigit{3}}\ {\isasymrightarrow}\ semaphore\ {\isadigit{3}}\ {\isacharbar}{\kern0pt}{\isacharbar}{\kern0pt}{\isacharbar}{\kern0pt}\ semaphore\ {\isadigit{4}}\isanewline
\ \ \ \ \ \ \ \ \ \ \ \ \ \ \ \ \ \ \ \ \ \ \ \ \ \ \ \ \ \ \ \ \ \ \ \ \ \ \ \ else\ semaphore\ {\isadigit{3}}\ {\isacharbar}{\kern0pt}{\isacharbar}{\kern0pt}{\isacharbar}{\kern0pt}\ evs{\isachardot}{\kern0pt}unlock\ {\isadigit{4}}\ {\isasymrightarrow}\ semaphore\ {\isadigit{4}}{\isacharparenright}{\kern0pt}{\isacharparenright}{\kern0pt}\isanewline
\ \ \ \ \ \ \ \ \ \ \ else\ evs{\isachardot}{\kern0pt}unlock\ {\isadigit{1}}\ {\isasymrightarrow}\ semaphore\ {\isadigit{1}}\isanewline
\ \ \ \ \ \ \ \ \ \ \ \ \ \ \ {\isacharbar}{\kern0pt}{\isacharbar}{\kern0pt}{\isacharbar}{\kern0pt}\ {\isacharparenleft}{\kern0pt}if\ x\ {\isacharequal}{\kern0pt}\ evs{\isachardot}{\kern0pt}lock\ {\isadigit{2}}\isanewline
\ \ \ \ \ \ \ \ \ \ \ \ \ \ \ \ \ \ \ \ then\ evs{\isachardot}{\kern0pt}unlock\ {\isadigit{2}}\ {\isasymrightarrow}\ semaphore\ {\isadigit{2}}\isanewline
\ \ \ \ \ \ \ \ \ \ \ \ \ \ \ \ \ \ \ \ \ \ \ \ \ {\isacharbar}{\kern0pt}{\isacharbar}{\kern0pt}{\isacharbar}{\kern0pt}\ {\isasymbox}x{\isasymin}{\isacharbraceleft}{\kern0pt}evs{\isachardot}{\kern0pt}lock\ {\isadigit{3}}{\isacharcomma}{\kern0pt}\ evs{\isachardot}{\kern0pt}lock\ {\isadigit{4}}{\isacharbraceright}{\kern0pt}\isanewline
\ \ \ \ \ \ \ \ \ \ \ \ \ \ \ \ \ \ \ \ \ \ \ \ \ \ \ \ \ {\isasymrightarrow}\ {\isacharparenleft}{\kern0pt}if\ x\ {\isacharequal}{\kern0pt}\ evs{\isachardot}{\kern0pt}lock\ {\isadigit{3}}\ then\ evs{\isachardot}{\kern0pt}unlock\ {\isadigit{3}}\ {\isasymrightarrow}\ semaphore\ {\isadigit{3}}\ {\isacharbar}{\kern0pt}{\isacharbar}{\kern0pt}{\isacharbar}{\kern0pt}\ semaphore\ {\isadigit{4}}\isanewline
\ \ \ \ \ \ \ \ \ \ \ \ \ \ \ \ \ \ \ \ \ \ \ \ \ \ \ \ \ \ \ \ else\ semaphore\ {\isadigit{3}}\ {\isacharbar}{\kern0pt}{\isacharbar}{\kern0pt}{\isacharbar}{\kern0pt}\ evs{\isachardot}{\kern0pt}unlock\ {\isadigit{4}}\ {\isasymrightarrow}\ semaphore\ {\isadigit{4}}{\isacharparenright}{\kern0pt}\isanewline
\ \ \ \ \ \ \ \ \ \ \ \ \ \ \ \ \ \ \ \ else\ semaphore\ {\isadigit{2}}\ \isanewline
\ \ \ \ \ \ \ \ \ \ \ \ \ \ \ \ \ \ \ \ \ \ \ \ \ {\isacharbar}{\kern0pt}{\isacharbar}{\kern0pt}{\isacharbar}{\kern0pt}\ {\isacharparenleft}{\kern0pt}if\ x\ {\isacharequal}{\kern0pt}\ evs{\isachardot}{\kern0pt}lock\ {\isadigit{3}}\isanewline
\ \ \ \ \ \ \ \ \ \ \ \ \ \ \ \ \ \ \ \ \ \ \ \ \ \ \ \ \ \ then\ evs{\isachardot}{\kern0pt}unlock\ {\isadigit{3}}\ {\isasymrightarrow}\ semaphore\ {\isadigit{3}}\ {\isacharbar}{\kern0pt}{\isacharbar}{\kern0pt}{\isacharbar}{\kern0pt}\ semaphore\ {\isadigit{4}}\isanewline
\ \ \ \ \ \ \ \ \ \ \ \ \ \ \ \ \ \ \ \ \ \ \ \ \ \ \ \ \ \ else\ semaphore\ {\isadigit{3}}\ {\isacharbar}{\kern0pt}{\isacharbar}{\kern0pt}{\isacharbar}{\kern0pt}\ evs{\isachardot}{\kern0pt}unlock\ {\isadigit{4}}\ {\isasymrightarrow}\ semaphore\ {\isadigit{4}}{\isacharparenright}{\kern0pt}{\isacharparenright}{\kern0pt}{\isacharparenright}{\kern0pt}{\isachardoublequoteclose}\isanewline
\ \ \ \ \ \ \ \ \ \isakeywordONE{apply}\isamarkupfalse%
{\isacharparenleft}{\kern0pt}simp\ add{\isacharcolon}{\kern0pt}\ List{\isachardot}{\kern0pt}upto{\isachardot}{\kern0pt}simps{\isacharparenright}{\kern0pt}\isanewline
\ \ \ \ \ \ \ \ \ \isakeywordONE{apply}\isamarkupfalse%
{\isacharparenleft}{\kern0pt}subst\ sema{\isacharunderscore}{\kern0pt}rec{\isacharbrackleft}{\kern0pt}\isakeywordTWO{where}\ n{\isacharequal}{\kern0pt}{\isacartoucheopen}{\isadigit{1}}{\isacartoucheclose}{\isacharbrackright}{\kern0pt}{\isacharcomma}{\kern0pt}subst\ sema{\isacharunderscore}{\kern0pt}rec{\isacharbrackleft}{\kern0pt}\isakeywordTWO{where}\ n{\isacharequal}{\kern0pt}{\isacartoucheopen}{\isadigit{2}}{\isacartoucheclose}{\isacharbrackright}{\kern0pt}{\isacharcomma}{\kern0pt}\isanewline
\ \ \ \ \ \ \ \ \ \ \ \ \ \ \ subst\ sema{\isacharunderscore}{\kern0pt}rec{\isacharbrackleft}{\kern0pt}\isakeywordTWO{where}\ n{\isacharequal}{\kern0pt}{\isacartoucheopen}{\isadigit{3}}{\isacartoucheclose}{\isacharbrackright}{\kern0pt}{\isacharcomma}{\kern0pt}subst\ sema{\isacharunderscore}{\kern0pt}rec{\isacharbrackleft}{\kern0pt}\isakeywordTWO{where}\ n{\isacharequal}{\kern0pt}{\isacartoucheopen}{\isadigit{4}}{\isacartoucheclose}{\isacharbrackright}{\kern0pt}{\isacharparenright}{\kern0pt}\ \isanewline
\ \ \ \ \ \ \ \ \ \isakeywordONE{apply}\isamarkupfalse%
{\isacharparenleft}{\kern0pt}simp\ add{\isacharcolon}{\kern0pt}\ inter{\isadigit{2}}Mprefix{\isadigit{1}}\ inter{\isadigit{2}}Mprefix{\isadigit{2}}{\isacharparenright}{\kern0pt}\isanewline
\ \ \ \ \ \ \ \ \ \isakeywordONE{by}\isamarkupfalse%
{\isacharparenleft}{\kern0pt}simp\ add{\isacharcolon}{\kern0pt}\ sema{\isacharunderscore}{\kern0pt}rec{\isacharbrackleft}{\kern0pt}\isakeywordTWO{where}\ n{\isacharequal}{\kern0pt}{\isacartoucheopen}{\isadigit{1}}{\isacartoucheclose}{\isacharcomma}{\kern0pt}symmetric{\isacharbrackright}{\kern0pt}\ \isanewline
\ \ \ \ \ \ \ \ \ \ \ \ \ \ \ \ \ \ \ \ \ \ sema{\isacharunderscore}{\kern0pt}rec{\isacharbrackleft}{\kern0pt}\isakeywordTWO{where}\ n{\isacharequal}{\kern0pt}{\isacartoucheopen}{\isadigit{2}}{\isacartoucheclose}{\isacharcomma}{\kern0pt}symmetric{\isacharbrackright}{\kern0pt}\ \isanewline
\ \ \ \ \ \ \ \ \ \ \ \ \ \ \ \ \ \ \ \ \ \ sema{\isacharunderscore}{\kern0pt}rec{\isacharbrackleft}{\kern0pt}\isakeywordTWO{where}\ n{\isacharequal}{\kern0pt}{\isacartoucheopen}{\isadigit{3}}{\isacartoucheclose}{\isacharcomma}{\kern0pt}symmetric{\isacharbrackright}{\kern0pt}\ \isanewline
\ \ \ \ \ \ \ \ \ \ \ \ \ \ \ \ \ \ \ \ \ \ sema{\isacharunderscore}{\kern0pt}rec{\isacharbrackleft}{\kern0pt}\isakeywordTWO{where}\ n{\isacharequal}{\kern0pt}{\isacartoucheopen}{\isadigit{4}}{\isacartoucheclose}{\isacharcomma}{\kern0pt}symmetric{\isacharbrackright}{\kern0pt}\ \isanewline
\ \ \ \ \ \ \ \ \ \ \ \ \ \ \ \ \ cong{\isacharcolon}{\kern0pt}\ HOL{\isachardot}{\kern0pt}if{\isacharunderscore}{\kern0pt}cong{\isacharparenright}{\kern0pt}\isanewline
\ \ \isakeywordTHREE{show}\isamarkupfalse%
\ {\isacharquery}{\kern0pt}thesis\isanewline
\ \ \ \ \isakeywordONE{apply}\isamarkupfalse%
{\isacharparenleft}{\kern0pt}subst{\isacharparenleft}{\kern0pt}{\isadigit{2}}{\isacharparenright}{\kern0pt}\ write{\isadigit{0}}{\isacharunderscore}{\kern0pt}def{\isacharbrackleft}{\kern0pt}THEN\ meta{\isacharunderscore}{\kern0pt}eq{\isacharunderscore}{\kern0pt}to{\isacharunderscore}{\kern0pt}obj{\isacharunderscore}{\kern0pt}eq{\isacharbrackright}{\kern0pt}{\isacharparenright}{\kern0pt}\isanewline
\ \ \ \ \isakeywordONE{apply}\isamarkupfalse%
{\isacharparenleft}{\kern0pt}subst\ {\isadigit{1}}{\isacharparenright}{\kern0pt}\isanewline
\ \ \ \ \isakeywordONE{by}\isamarkupfalse%
\ {\isacharparenleft}{\kern0pt}simp\ add{\isacharcolon}{\kern0pt}\ Mprefix{\isacharunderscore}{\kern0pt}Par{\isacharunderscore}{\kern0pt}Mprefix{\isacharparenright}{\kern0pt}\isanewline
\isakeywordONE{qed}\isamarkupfalse%
%
\endisatagproof
{\isafoldproof}%
%
\isadelimproof
%
\endisadelimproof
%
\begin{isamarkuptext}%
After these preparatory lemmas, the path to the deadlock is straight-forward:
it proceeds by unroling the while-loop two times; unfolding
\isa{\isaconst{semaphore}\ {\isadigit{1}}} and its fixpoint two times; simplifying to \isa{\isaconst{com{\isachardot}{\kern0pt}STOP}}
since \isa{\isaconst{evs{\isachardot}{\kern0pt}lock}\ {\isadigit{0}}} and \isa{\isaconst{evs{\isachardot}{\kern0pt}unlock}\ {\isadigit{0}}} are forced  to be synchronized, 
propagating this result to the refinement level and reduce it to a contradiction
of the refinement statement with the \isa{\isaconst{DF}\ \isaconst{UNIV}} process.
We skip the formal proof here since it is out of scope of this presentation.%
\end{isamarkuptext}\isamarkuptrue%
\isakeywordONE{theorem}\isamarkupfalse%
\ bullocks{\isacharcolon}{\kern0pt}\ {\isacartoucheopen}{\isasymnot}\ deadlock{\isacharunderscore}{\kern0pt}free\ Thread{\isadigit{5}}{\isacharunderscore}{\kern0pt}sem{\isacartoucheclose}\isanewline
%
\isadelimproof
\ \ %
\endisadelimproof
%
\isatagproof
\isakeywordONE{unfolding}\isamarkupfalse%
\ Thread{\isadigit{5}}{\isacharunderscore}{\kern0pt}sem{\isacharunderscore}{\kern0pt}def\ Sem{\isacharunderscore}{\kern0pt}def\ Thread{\isadigit{5}}{\isacharunderscore}{\kern0pt}def\isanewline
\isakeywordONE{proof}\isamarkupfalse%
\ {\isacharminus}{\kern0pt}\ \isanewline
\ \ \isakeywordTHREE{show}\isamarkupfalse%
\ {\isacartoucheopen}{\isasymnot}\ deadlock{\isacharunderscore}{\kern0pt}free\isanewline
\ \ \ \ \ \ \ \ \ \ \ \ \ \ \ {\isacharparenleft}{\kern0pt}\ MultiInter\ {\isacharparenleft}{\kern0pt}mset\ {\isacharbrackleft}{\kern0pt}{\isadigit{1}}{\isachardot}{\kern0pt}{\isachardot}{\kern0pt}{\isadigit{4}}{\isacharbrackright}{\kern0pt}{\isacharparenright}{\kern0pt}\ semaphore\isanewline
\ \ \ \ \ \ \ \ \ \ \ \ \ \ \ \ {\isacharbar}{\kern0pt}{\isacharbar}{\kern0pt}\ Sem\isactrlsub {\isadigit{0}}\ {\isacharparenleft}{\kern0pt}WHILE\ {\isasymlambda}{\isasymsigma}{\isachardot}{\kern0pt}\ True\ DO\ com{\isachardot}{\kern0pt}lock\ {\isadigit{1}}{\isacharparenright}{\kern0pt}\ {\isacharparenleft}{\kern0pt}{\isasymlambda}{\isacharunderscore}{\kern0pt}{\isachardot}{\kern0pt}\ Skip{\isacharparenright}{\kern0pt}\ {\isasymsigma}\isactrlsub {\isadigit{0}}{\isacharparenright}{\kern0pt}{\isacartoucheclose}\isanewline
\ \ \ \ \ \ \ \ \ \ \ \ \ \ \ \ {\isacharparenleft}{\kern0pt}\isakeywordTWO{is}\ {\isacartoucheopen}{\isasymnot}\ deadlock{\isacharunderscore}{\kern0pt}free\ {\isacharparenleft}{\kern0pt}{\isacharquery}{\kern0pt}SEMA\ {\isacharbar}{\kern0pt}{\isacharbar}{\kern0pt}\ {\isacharquery}{\kern0pt}P\ {\isacharparenright}{\kern0pt}{\isacartoucheclose}{\isacharparenright}{\kern0pt}\isanewline
\ \ \isakeywordONE{proof}\isamarkupfalse%
\ {\isacharminus}{\kern0pt}\ \isanewline
\ \ \ \ \isakeywordONE{have}\isamarkupfalse%
\ {\isadigit{1}}\ {\isacharcolon}{\kern0pt}\ {\isacartoucheopen}{\isacharquery}{\kern0pt}P\ {\isacharequal}{\kern0pt}\ evs{\isachardot}{\kern0pt}lock\ {\isadigit{1}}\ {\isasymrightarrow}\ {\isacharquery}{\kern0pt}P{\isacartoucheclose}\ \isakeywordONE{using}\isamarkupfalse%
\ Thread{\isadigit{5}}{\isacharunderscore}{\kern0pt}while{\isacharunderscore}{\kern0pt}rec\ \isakeywordONE{by}\isamarkupfalse%
\ auto\isanewline
\ \ \ \ \isakeywordONE{have}\isamarkupfalse%
\ {\isadigit{2}}\ {\isacharcolon}{\kern0pt}\ {\isacartoucheopen}{\isacharparenleft}{\kern0pt}{\isacharquery}{\kern0pt}SEMA\ {\isacharbar}{\kern0pt}{\isacharbar}{\kern0pt}\ evs{\isachardot}{\kern0pt}lock\ {\isadigit{1}}\ {\isasymrightarrow}\ {\isacharquery}{\kern0pt}P{\isacharparenright}{\kern0pt}\ \isanewline
\ \ \ \ \ \ \ \ \ \ \ \ \ \ {\isacharequal}{\kern0pt}\ \ evs{\isachardot}{\kern0pt}lock\ {\isadigit{1}}\ {\isasymrightarrow}\ {\isacharparenleft}{\kern0pt}{\isacharparenleft}{\kern0pt}evs{\isachardot}{\kern0pt}unlock\ {\isadigit{1}}\ {\isasymrightarrow}\ semaphore\ {\isadigit{1}}\ \ \isanewline
\ \ \ \ \ \ \ \ \ \ \ \ \ \ \ \ \ \ \ \ \ \ \ \ \ \ \ \ \ \ \ \ {\isacharbar}{\kern0pt}{\isacharbar}{\kern0pt}{\isacharbar}{\kern0pt}\ MultiInter\ {\isacharparenleft}{\kern0pt}mset\ {\isacharbrackleft}{\kern0pt}{\isadigit{2}}{\isachardot}{\kern0pt}{\isachardot}{\kern0pt}{\isadigit{4}}{\isacharbrackright}{\kern0pt}{\isacharparenright}{\kern0pt}\ semaphore{\isacharparenright}{\kern0pt}\ {\isacharbar}{\kern0pt}{\isacharbar}{\kern0pt}\ {\isacharquery}{\kern0pt}P{\isacharparenright}{\kern0pt}{\isacartoucheclose}\isanewline
\ \ \ \ \ \ \ \ \ \ \ \ \ {\isacharparenleft}{\kern0pt}\isakeywordTWO{is}\ {\isachardoublequoteopen}{\isacharquery}{\kern0pt}lhs\ {\isacharequal}{\kern0pt}\ evs{\isachardot}{\kern0pt}lock\ {\isadigit{1}}\ {\isasymrightarrow}\ {\isacharquery}{\kern0pt}R{\isacharprime}{\kern0pt}{\isachardoublequoteclose}{\isacharparenright}{\kern0pt}\isanewline
\ \ \ \ \ \ \isakeywordONE{by}\isamarkupfalse%
\ {\isacharparenleft}{\kern0pt}metis\ contractMultiInter{\isadigit{1}}{\isacharparenright}{\kern0pt}\isanewline
\ \ \ \ \isakeywordONE{have}\isamarkupfalse%
\ {\isadigit{3}}\ {\isacharcolon}{\kern0pt}\ {\isacartoucheopen}{\isacharparenleft}{\kern0pt}{\isasymnot}\ deadlock{\isacharunderscore}{\kern0pt}free\ Thread{\isadigit{5}}{\isacharunderscore}{\kern0pt}sem{\isacharparenright}{\kern0pt}\ {\isacharequal}{\kern0pt}\ {\isacharparenleft}{\kern0pt}{\isasymnot}\ deadlock{\isacharunderscore}{\kern0pt}free\ {\isacharquery}{\kern0pt}R{\isacharprime}{\kern0pt}{\isacharparenright}{\kern0pt}{\isacartoucheclose}\isanewline
\ \ \ \ \ \ \isakeywordONE{by}\isamarkupfalse%
\ {\isacharparenleft}{\kern0pt}metis\ {\isadigit{1}}\ {\isadigit{2}}\ Sem{\isacharunderscore}{\kern0pt}def\ Thread{\isadigit{5}}{\isacharunderscore}{\kern0pt}def\ Thread{\isadigit{5}}{\isacharunderscore}{\kern0pt}sem{\isacharunderscore}{\kern0pt}def\ deadlock{\isacharunderscore}{\kern0pt}free{\isacharunderscore}{\kern0pt}write{\isadigit{0}}{\isacharunderscore}{\kern0pt}iff{\isacharparenright}{\kern0pt}\isanewline
\ \ \ \ \isakeywordONE{have}\isamarkupfalse%
\ {\isadigit{4}}{\isacharcolon}{\kern0pt}\ {\isacartoucheopen}{\isacharquery}{\kern0pt}R{\isacharprime}{\kern0pt}\ {\isacharequal}{\kern0pt}\ Constant{\isacharunderscore}{\kern0pt}Processes{\isachardot}{\kern0pt}STOP{\isacartoucheclose}\ \isanewline
\ \ \ \ \ \ \isakeywordONE{by}\isamarkupfalse%
\ {\isacharparenleft}{\kern0pt}metis\ {\isadigit{1}}\ contractMultiInter{\isadigit{2}}{\isacharparenright}{\kern0pt}\isanewline
\ \ \ \ \isakeywordTHREE{show}\isamarkupfalse%
\ {\isacharquery}{\kern0pt}thesis\isanewline
\ \ \ \ \ \ \isakeywordONE{using}\isamarkupfalse%
\ {\isadigit{3}}\ {\isadigit{4}}\ Sem{\isacharunderscore}{\kern0pt}def\ Thread{\isadigit{5}}{\isacharunderscore}{\kern0pt}def\ Thread{\isadigit{5}}{\isacharunderscore}{\kern0pt}sem{\isacharunderscore}{\kern0pt}def\ non{\isacharunderscore}{\kern0pt}deadlock{\isacharunderscore}{\kern0pt}free{\isacharunderscore}{\kern0pt}STOP\ \isakeywordONE{by}\isamarkupfalse%
\ auto\isanewline
\ \ \isakeywordONE{qed}\isamarkupfalse%
\isanewline
\isakeywordONE{qed}\isamarkupfalse%
%
\endisatagproof
{\isafoldproof}%
%
\isadelimproof
%
\endisadelimproof
%
\begin{isamarkuptext}%
This proof can be drastically shortened / automatized if we would use a more advanced 
background theory of HOL-CSP, for example:%
\end{isamarkuptext}\isamarkuptrue%
\isakeywordONE{term}\isamarkupfalse%
{\isacartoucheopen}Initials\ Q\ {\isasymsubseteq}\ ev\ {\isacharbackquote}{\kern0pt}\ S\ \isanewline
\ \ \ \ \ {\isasymLongrightarrow}\ Initials\ Q\ {\isasyminter}\ ev\ {\isacharbackquote}{\kern0pt}\ T\ {\isacharequal}{\kern0pt}\ {\isacharbraceleft}{\kern0pt}{\isacharbraceright}{\kern0pt}\ \isanewline
\ \ \ \ \ {\isasymLongrightarrow}\ \ {\isacharparenleft}{\kern0pt}{\isacharparenleft}{\kern0pt}P\ {\isasymbox}\ Q{\isacharparenright}{\kern0pt}\ {\isasymlbrakk}S{\isasymrbrakk}\ {\isacharparenleft}{\kern0pt}Mprefix\ T\ R{\isacharparenright}{\kern0pt}{\isacharparenright}{\kern0pt}\ {\isacharequal}{\kern0pt}\ P\ {\isasymlbrakk}S{\isasymrbrakk}\ {\isacharparenleft}{\kern0pt}Mprefix\ T\ R{\isacharparenright}{\kern0pt}{\isacartoucheclose}\isanewline
\isanewline
\isakeywordONE{term}\isamarkupfalse%
{\isacartoucheopen}{\isasymalpha}{\isacharparenleft}{\kern0pt}Q{\isacharparenright}{\kern0pt}\ {\isasymsubseteq}\ S\ \isanewline
\ \ \ \ \ {\isasymLongrightarrow}\ {\isasymalpha}{\isacharparenleft}{\kern0pt}Q{\isacharparenright}{\kern0pt}\ {\isasyminter}\ T\ {\isacharequal}{\kern0pt}\ {\isacharbraceleft}{\kern0pt}{\isacharbraceright}{\kern0pt}\ \isanewline
\ \ \ \ \ {\isasymLongrightarrow}\ {\isasymalpha}{\isacharparenleft}{\kern0pt}Q{\isacharparenright}{\kern0pt}\ {\isasyminter}\ {\isasymalpha}{\isacharparenleft}{\kern0pt}R{\isacharparenright}{\kern0pt}\ {\isacharequal}{\kern0pt}\ {\isacharbraceleft}{\kern0pt}{\isacharbraceright}{\kern0pt}\ \isanewline
\ \ \ \ \ {\isasymLongrightarrow}\ {\isacharparenleft}{\kern0pt}{\isacharparenleft}{\kern0pt}P\ {\isacharbar}{\kern0pt}{\isacharbar}{\kern0pt}{\isacharbar}{\kern0pt}\ Q{\isacharparenright}{\kern0pt}\ {\isasymlbrakk}S{\isasymrbrakk}\ R{\isacharparenright}{\kern0pt}\ {\isacharequal}{\kern0pt}\ P\ {\isasymlbrakk}S{\isasymrbrakk}\ R{\isacartoucheclose}\isanewline
\isanewline
\isakeywordONE{term}\isamarkupfalse%
{\isacartoucheopen}{\isasymalpha}{\isacharparenleft}{\kern0pt}Q{\isacharparenright}{\kern0pt}\ {\isasymsubseteq}\ S\ \isanewline
\ \ \ \ \ {\isasymLongrightarrow}\ {\isasymalpha}{\isacharparenleft}{\kern0pt}Q{\isacharparenright}{\kern0pt}\ {\isasyminter}\ T\ {\isacharequal}{\kern0pt}\ {\isacharbraceleft}{\kern0pt}{\isacharbraceright}{\kern0pt}\ \isanewline
\ \ \ \ \ {\isasymLongrightarrow}\ Initials\ Q\ {\isasyminter}\ {\isasymalpha}{\isacharparenleft}{\kern0pt}{\isacharparenleft}{\kern0pt}Mprefix\ T\ R{\isacharparenright}{\kern0pt}{\isacharparenright}{\kern0pt}\ {\isacharequal}{\kern0pt}\ {\isacharbraceleft}{\kern0pt}{\isacharbraceright}{\kern0pt}\ \isanewline
\ \ \ \ \ {\isasymLongrightarrow}\ {\isacharparenleft}{\kern0pt}{\isacharparenleft}{\kern0pt}P\ {\isacharbar}{\kern0pt}{\isacharbar}{\kern0pt}{\isacharbar}{\kern0pt}\ Q{\isacharparenright}{\kern0pt}\ {\isasymlbrakk}S{\isasymrbrakk}\ {\isacharparenleft}{\kern0pt}Mprefix\ T\ R{\isacharparenright}{\kern0pt}{\isacharparenright}{\kern0pt}\ {\isacharequal}{\kern0pt}\ P\ {\isasymlbrakk}S{\isasymrbrakk}\ {\isacharparenleft}{\kern0pt}Mprefix\ T\ R{\isacharparenright}{\kern0pt}{\isacartoucheclose}%
\isadelimdocument
%
\endisadelimdocument
%
\isatagdocument
%
\isamarkupsection{Yet another Semantic Corner-Case: A Deadlock-Freeness Proof%
}
\isamarkuptrue%
%
\endisatagdocument
{\isafolddocument}%
%
\isadelimdocument
%
\endisadelimdocument
%
\begin{isamarkuptext}%
We slightly mofify the previous example by adding the following thread into the
architecture:%
\end{isamarkuptext}\isamarkuptrue%
\isakeywordONE{definition}\isamarkupfalse%
\ Thread{\isadigit{6}}\ \isakeywordTWO{where}\ \isanewline
\ \ {\isachardoublequoteopen}Thread{\isadigit{6}}\ {\isacharequal}{\kern0pt}\ WHILE\ {\isacharparenleft}{\kern0pt}{\isasymlambda}{\isasymsigma}{\isachardot}{\kern0pt}\ True{\isacharparenright}{\kern0pt}\ DO\ {\isacharparenleft}{\kern0pt}com{\isachardot}{\kern0pt}unlock\ {\isadigit{1}}{\isacharparenright}{\kern0pt}{\isachardoublequoteclose}\isanewline
\isanewline
\isakeywordONE{lemma}\isamarkupfalse%
\ Thread{\isadigit{6}}{\isacharunderscore}{\kern0pt}while{\isacharunderscore}{\kern0pt}rec\ {\isacharcolon}{\kern0pt}\ \isanewline
\ \ \ \ \ \ \ \ \ \ {\isachardoublequoteopen}{\isacharparenleft}{\kern0pt}{\isasymmu}\ x{\isacharcolon}{\kern0pt}{\isacharcolon}{\kern0pt}{\isacharprime}{\kern0pt}b\ {\isasymRightarrow}\ evs\ process{\isachardot}{\kern0pt}\ {\isacharparenleft}{\kern0pt}{\isasymlambda}{\isasymsigma}{\isachardot}{\kern0pt}\ evs{\isachardot}{\kern0pt}unlock\ {\isacharparenleft}{\kern0pt}{\isadigit{1}}{\isacharcolon}{\kern0pt}{\isacharcolon}{\kern0pt}int{\isacharparenright}{\kern0pt}\ {\isasymrightarrow}\ x\ {\isasymsigma}{\isacharparenright}{\kern0pt}{\isacharparenright}{\kern0pt}\ A\ \isanewline
\ \ \ \ \ \ \ \ \ \ \ \ {\isacharequal}{\kern0pt}\ evs{\isachardot}{\kern0pt}unlock\ {\isadigit{1}}\ {\isasymrightarrow}\ {\isacharparenleft}{\kern0pt}{\isasymmu}\ x{\isachardot}{\kern0pt}\ {\isacharparenleft}{\kern0pt}{\isasymlambda}{\isasymsigma}{\isachardot}{\kern0pt}\ evs{\isachardot}{\kern0pt}unlock\ {\isacharparenleft}{\kern0pt}{\isadigit{1}}{\isacharcolon}{\kern0pt}{\isacharcolon}{\kern0pt}int{\isacharparenright}{\kern0pt}\ {\isasymrightarrow}\ x\ {\isasymsigma}{\isacharparenright}{\kern0pt}{\isacharparenright}{\kern0pt}\ A{\isachardoublequoteclose}\isanewline
%
\isadelimproof
\ \ %
\endisadelimproof
%
\isatagproof
\isakeywordONE{by}\isamarkupfalse%
{\isacharparenleft}{\kern0pt}subst\ cont{\isacharunderscore}{\kern0pt}process{\isacharunderscore}{\kern0pt}rec{\isacharbrackleft}{\kern0pt}\isakeywordTWO{where}\ P\ {\isacharequal}{\kern0pt}\ {\isacartoucheopen}{\isacharparenleft}{\kern0pt}{\isasymmu}\ x{\isachardot}{\kern0pt}\ {\isacharparenleft}{\kern0pt}{\isasymlambda}{\isasymsigma}{\isachardot}{\kern0pt}\ evs{\isachardot}{\kern0pt}unlock\ {\isacharparenleft}{\kern0pt}{\isadigit{1}}{\isacharcolon}{\kern0pt}{\isacharcolon}{\kern0pt}int{\isacharparenright}{\kern0pt}\ {\isasymrightarrow}\ x\ {\isasymsigma}{\isacharparenright}{\kern0pt}{\isacharparenright}{\kern0pt}{\isacartoucheclose}{\isacharcomma}{\kern0pt}\ \isanewline
\ \ \ \ \ \ \ \ \ \ \ \ \ \ \ \ \ \ \ \ \ \ \ \ \ \ \ \ OF\ refl{\isacharbrackright}{\kern0pt}{\isacharparenright}{\kern0pt}\ simp{\isacharunderscore}{\kern0pt}all%
\endisatagproof
{\isafoldproof}%
%
\isadelimproof
\isanewline
%
\endisadelimproof
\isanewline
\isanewline
\isakeywordONE{definition}\isamarkupfalse%
\ Thread{\isadigit{6}}{\isacharunderscore}{\kern0pt}sem\ \isakeywordTWO{where}\ \isanewline
{\isachardoublequoteopen}Thread{\isadigit{6}}{\isacharunderscore}{\kern0pt}sem\ {\isasymequiv}\ {\isacharparenleft}{\kern0pt}\ \isactrlbold {\isacharbar}{\kern0pt}\isactrlbold {\isacharbar}{\kern0pt}\isactrlbold {\isacharbar}{\kern0pt}\ idx\ {\isasymin}{\isacharhash}{\kern0pt}\ mset\ {\isacharbrackleft}{\kern0pt}{\isadigit{1}}{\isachardot}{\kern0pt}{\isachardot}{\kern0pt}{\isadigit{4}}{\isacharbrackright}{\kern0pt}{\isachardot}{\kern0pt}\ semaphore\ idx\ {\isacharparenright}{\kern0pt}\ \ \isanewline
\ \ \ \ \ \ \ \ \ \ \ \ \ \ \ \ \ \isanewline
\ \ \ \ \ \ \ \ \ \ \ \ \ \ \ \ \ {\isacharbar}{\kern0pt}{\isacharbar}{\kern0pt}\isanewline
\ \isanewline
\ \ \ \ \ \ \ \ \ \ \ \ \ \ \ \ \ {\isacharparenleft}{\kern0pt}Sem\ Thread{\isadigit{5}}\ {\isacharbar}{\kern0pt}{\isacharbar}{\kern0pt}{\isacharbar}{\kern0pt}\ Sem\ Thread{\isadigit{6}}{\isacharparenright}{\kern0pt}{\isachardoublequoteclose}%
\begin{isamarkuptext}%
The architecture \isa{\isaconst{Thread{\isadigit{6}}}}, however, will not deadlock, although each single thread will.
This shows that in \isa{IMP\isactrlsub c\isactrlsub o\isactrlsub n\isactrlsub c\isactrlsub u\isactrlsub r} locks can mutually deblock each other and have ---
together with the global variables --- global  visibility inside the thread system.%
\end{isamarkuptext}\isamarkuptrue%
%
\begin{isamarkuptext}%
The statement for this claim looks as follows:%
\end{isamarkuptext}\isamarkuptrue%
\isakeywordONE{term}\isamarkupfalse%
\ \ {\isacartoucheopen}deadlock{\isacharunderscore}{\kern0pt}free\ Thread{\isadigit{6}}{\isacharunderscore}{\kern0pt}sem{\isacartoucheclose}%
\begin{isamarkuptext}%
We will present in this theory a straight-forward proof approach (HOL-CSP has more
sophisticated and also more automatic methods than the suggested one). We need a little
background theory concerning \isa{\isaconst{deadlock{\isacharunderscore}{\kern0pt}free}}'ness here; in principle, it can be 
established via a consequence of the fixpoint induction. A complication to be tackled
is that we need a format of this induction that takes two events --- instead of just one ---
to arrive at the anchor of the induction, namely \isa{\isaconst{evs{\isachardot}{\kern0pt}lock}} and \isa{\isaconst{evs{\isachardot}{\kern0pt}unlock}}.%
\end{isamarkuptext}\isamarkuptrue%
%
\isadelimdocument
%
\endisadelimdocument
%
\isatagdocument
%
\isamarkupsubsection{On Coinduction for deadlock-freeness.%
}
\isamarkuptrue%
%
\endisatagdocument
{\isafolddocument}%
%
\isadelimdocument
%
\endisadelimdocument
%
\begin{isamarkuptext}%
In the following, we derive the necessary deadlock-coinduction:%
\end{isamarkuptext}\isamarkuptrue%
\isakeywordONE{lemma}\isamarkupfalse%
\ lasso{\isacharunderscore}{\kern0pt}equiv{\isadigit{1}}{\isacharcolon}{\kern0pt}\ {\isachardoublequoteopen}DF\ UNIV\ {\isacharequal}{\kern0pt}\ {\isacharparenleft}{\kern0pt}{\isasymmu}\ x{\isachardot}{\kern0pt}\ {\isasymsqinter}a{\isasymin}UNIV\ {\isasymrightarrow}\ \ {\isasymsqinter}a{\isasymin}UNIV\ {\isasymrightarrow}\ x{\isacharparenright}{\kern0pt}{\isachardoublequoteclose}\isanewline
%
\isadelimproof
%
\endisadelimproof
%
\isatagproof
\isakeywordONE{unfolding}\isamarkupfalse%
\ DF{\isacharunderscore}{\kern0pt}def\isanewline
\isakeywordONE{proof}\isamarkupfalse%
{\isacharparenleft}{\kern0pt}rule\ FD{\isacharunderscore}{\kern0pt}antisym{\isacharparenright}{\kern0pt}\isanewline
\ \ \isakeywordONE{have}\isamarkupfalse%
\ {\isacharasterisk}{\kern0pt}\ {\isacharcolon}{\kern0pt}\ {\isacartoucheopen}X\ {\isacharequal}{\kern0pt}\ {\isacharparenleft}{\kern0pt}X\ {\isasymand}\ X{\isacharparenright}{\kern0pt}{\isacartoucheclose}\ \isakeywordTWO{for}\ X\ \isakeywordONE{by}\isamarkupfalse%
\ simp\isanewline
\ \ \isakeywordTHREE{show}\isamarkupfalse%
\ {\isacartoucheopen}{\isacharparenleft}{\kern0pt}{\isasymmu}\ x{\isachardot}{\kern0pt}\ {\isasymsqinter}a{\isasymin}UNIV\ {\isasymrightarrow}\ x{\isacharparenright}{\kern0pt}\ {\isasymsqsubseteq}\isactrlsub F\isactrlsub D\ {\isacharparenleft}{\kern0pt}{\isasymmu}\ x{\isachardot}{\kern0pt}\ {\isasymsqinter}a{\isasymin}UNIV\ {\isasymrightarrow}\ {\isasymsqinter}a{\isasymin}UNIV\ {\isasymrightarrow}\ x{\isacharparenright}{\kern0pt}{\isacartoucheclose}\isanewline
\ \ \ \ \isakeywordONE{apply}\isamarkupfalse%
{\isacharparenleft}{\kern0pt}subst\ {\isacharasterisk}{\kern0pt}{\isacharparenright}{\kern0pt}\isanewline
\ \ \ \ \isakeywordONE{apply}\isamarkupfalse%
{\isacharparenleft}{\kern0pt}subst\ cont{\isacharunderscore}{\kern0pt}process{\isacharunderscore}{\kern0pt}rec{\isacharbrackleft}{\kern0pt}\isakeywordTWO{where}\ P\ {\isacharequal}{\kern0pt}\ {\isacartoucheopen}{\isacharparenleft}{\kern0pt}{\isasymmu}\ x{\isachardot}{\kern0pt}\ {\isasymsqinter}a{\isasymin}UNIV\ {\isasymrightarrow}\ x{\isacharparenright}{\kern0pt}{\isacartoucheclose}\ \isanewline
\ \ \ \ \ \ \ \ \ \ \ \ \ \ \ \ \ \ \ \ \ \ \ \ \ \ \ \ \ \ \ \ \ \ \ \isakeywordTWO{and}\ f\ {\isacharequal}{\kern0pt}\ {\isacartoucheopen}{\isasymlambda}x{\isachardot}{\kern0pt}\ {\isasymsqinter}a{\isasymin}UNIV\ {\isasymrightarrow}\ x{\isacartoucheclose}{\isacharbrackright}{\kern0pt}{\isacharcomma}{\kern0pt}simp{\isacharunderscore}{\kern0pt}all{\isacharparenright}{\kern0pt}\isanewline
\ \ \ \ \isakeywordONE{apply}\isamarkupfalse%
{\isacharparenleft}{\kern0pt}rule\ fix{\isacharunderscore}{\kern0pt}ind{\isacharbrackleft}{\kern0pt}\isakeywordTWO{where}\ F\ {\isacharequal}{\kern0pt}\ {\isacartoucheopen}{\isacharparenleft}{\kern0pt}{\isasymLambda}\ x{\isachardot}{\kern0pt}\ {\isasymsqinter}a{\isasymin}UNIV\ {\isasymrightarrow}\ x{\isacharparenright}{\kern0pt}{\isacartoucheclose}{\isacharbrackright}{\kern0pt}{\isacharcomma}{\kern0pt}\ simp{\isacharcomma}{\kern0pt}simp{\isacharparenright}{\kern0pt}\ \isanewline
\ \ \ \ \ \isakeywordONE{apply}\isamarkupfalse%
{\isacharparenleft}{\kern0pt}subst\ cont{\isacharunderscore}{\kern0pt}process{\isacharunderscore}{\kern0pt}rec{\isacharbrackleft}{\kern0pt}\isakeywordTWO{where}\ P\ {\isacharequal}{\kern0pt}\ {\isacartoucheopen}{\isacharparenleft}{\kern0pt}{\isasymmu}\ x{\isachardot}{\kern0pt}\ {\isasymsqinter}a{\isasymin}UNIV\ {\isasymrightarrow}\ {\isasymsqinter}a{\isasymin}UNIV\ {\isasymrightarrow}\ x{\isacharparenright}{\kern0pt}{\isacartoucheclose}\ \isanewline
\ \ \ \ \ \ \ \ \ \ \ \ \ \ \ \ \ \ \ \ \ \ \ \ \ \ \ \ \ \ \ \ \ \ \ \ \isakeywordTWO{and}\ f\ {\isacharequal}{\kern0pt}\ {\isacartoucheopen}{\isasymlambda}x{\isachardot}{\kern0pt}\ {\isasymsqinter}a{\isasymin}UNIV\ {\isasymrightarrow}\ {\isasymsqinter}a{\isasymin}UNIV\ {\isasymrightarrow}\ x{\isacartoucheclose}{\isacharbrackright}{\kern0pt}{\isacharcomma}{\kern0pt}\ simp{\isacharunderscore}{\kern0pt}all{\isacharparenright}{\kern0pt}\isanewline
\ \ \ \ \ \isakeywordONE{apply}\isamarkupfalse%
{\isacharparenleft}{\kern0pt}rule\ mono{\isacharunderscore}{\kern0pt}Mndetprefix{\isacharunderscore}{\kern0pt}FD{\isacharcomma}{\kern0pt}\ simp{\isacharparenright}{\kern0pt}\ \isanewline
\ \ \ \ \isakeywordONE{apply}\isamarkupfalse%
{\isacharparenleft}{\kern0pt}subst\ cont{\isacharunderscore}{\kern0pt}process{\isacharunderscore}{\kern0pt}rec{\isacharbrackleft}{\kern0pt}\isakeywordTWO{where}\ P\ {\isacharequal}{\kern0pt}\ {\isacartoucheopen}{\isacharparenleft}{\kern0pt}{\isasymmu}\ x{\isachardot}{\kern0pt}\ {\isasymsqinter}a{\isasymin}UNIV\ {\isasymrightarrow}\ {\isasymsqinter}a{\isasymin}UNIV\ {\isasymrightarrow}\ x{\isacharparenright}{\kern0pt}{\isacartoucheclose}\ \isanewline
\ \ \ \ \ \ \ \ \ \ \ \ \ \ \ \ \ \ \ \ \ \ \ \ \ \ \ \ \ \ \ \ \ \ \ \ \ \ \ \isakeywordTWO{and}\ f\ {\isacharequal}{\kern0pt}\ {\isacartoucheopen}{\isasymlambda}x{\isachardot}{\kern0pt}\ {\isasymsqinter}a{\isasymin}UNIV\ {\isasymrightarrow}\ {\isasymsqinter}a{\isasymin}UNIV\ {\isasymrightarrow}\ x{\isacartoucheclose}{\isacharbrackright}{\kern0pt}{\isacharcomma}{\kern0pt}\ simp{\isacharunderscore}{\kern0pt}all{\isacharparenright}{\kern0pt}\isanewline
\ \ \ \ \isakeywordONE{by}\isamarkupfalse%
\ {\isacharparenleft}{\kern0pt}simp\ add{\isacharcolon}{\kern0pt}\ mono{\isacharunderscore}{\kern0pt}Mndetprefix{\isacharunderscore}{\kern0pt}FD{\isacharparenright}{\kern0pt}\isanewline
\isakeywordONE{next}\isamarkupfalse%
\ \isanewline
\ \ \isakeywordTHREE{show}\isamarkupfalse%
\ {\isacartoucheopen}{\isacharparenleft}{\kern0pt}{\isasymmu}\ x{\isachardot}{\kern0pt}\ {\isasymsqinter}a{\isasymin}UNIV\ {\isasymrightarrow}\ {\isasymsqinter}a{\isasymin}UNIV\ {\isasymrightarrow}\ x{\isacharparenright}{\kern0pt}\ {\isasymsqsubseteq}\isactrlsub F\isactrlsub D\ {\isacharparenleft}{\kern0pt}{\isasymmu}\ x{\isachardot}{\kern0pt}\ {\isasymsqinter}a{\isasymin}UNIV\ {\isasymrightarrow}\ x{\isacharparenright}{\kern0pt}{\isacartoucheclose}\ \isanewline
\ \ \ \ \isakeywordONE{apply}\isamarkupfalse%
{\isacharparenleft}{\kern0pt}rule\ fix{\isacharunderscore}{\kern0pt}ind{\isacharbrackleft}{\kern0pt}\isakeywordTWO{where}\ F\ {\isacharequal}{\kern0pt}\ {\isacartoucheopen}{\isacharparenleft}{\kern0pt}{\isasymLambda}\ x{\isachardot}{\kern0pt}\ {\isasymsqinter}a{\isasymin}UNIV\ {\isasymrightarrow}\ {\isasymsqinter}a{\isasymin}UNIV\ {\isasymrightarrow}\ x{\isacharparenright}{\kern0pt}{\isacartoucheclose}{\isacharbrackright}{\kern0pt}\ {\isacharparenright}{\kern0pt}\ \ \isanewline
\ \ \ \ \ \ \isakeywordONE{apply}\isamarkupfalse%
\ {\isacharparenleft}{\kern0pt}simp{\isacharunderscore}{\kern0pt}all{\isacharparenright}{\kern0pt}\isanewline
\ \ \ \ \isakeywordONE{apply}\isamarkupfalse%
{\isacharparenleft}{\kern0pt}subst\ cont{\isacharunderscore}{\kern0pt}process{\isacharunderscore}{\kern0pt}rec{\isacharbrackleft}{\kern0pt}\isakeywordTWO{where}\ P\ {\isacharequal}{\kern0pt}\ {\isacartoucheopen}{\isacharparenleft}{\kern0pt}{\isasymmu}\ x{\isachardot}{\kern0pt}\ {\isasymsqinter}a{\isasymin}UNIV\ {\isasymrightarrow}\ x{\isacharparenright}{\kern0pt}{\isacartoucheclose}\ \isanewline
\ \ \ \ \ \ \ \ \ \ \ \ \ \ \ \ \ \ \ \ \ \ \ \ \ \ \ \ \ \ \ \ \ \ \ \ \ \ \ \isakeywordTWO{and}\ f\ {\isacharequal}{\kern0pt}\ {\isacartoucheopen}{\isasymlambda}x{\isachardot}{\kern0pt}\ {\isasymsqinter}a{\isasymin}UNIV\ {\isasymrightarrow}\ x{\isacartoucheclose}{\isacharbrackright}{\kern0pt}{\isacharcomma}{\kern0pt}\ simp{\isacharunderscore}{\kern0pt}all{\isacharparenright}{\kern0pt}\isanewline
\ \ \ \ \isakeywordONE{apply}\isamarkupfalse%
{\isacharparenleft}{\kern0pt}subst\ cont{\isacharunderscore}{\kern0pt}process{\isacharunderscore}{\kern0pt}rec{\isacharbrackleft}{\kern0pt}\isakeywordTWO{where}\ P\ {\isacharequal}{\kern0pt}\ {\isacartoucheopen}{\isacharparenleft}{\kern0pt}{\isasymmu}\ x{\isachardot}{\kern0pt}\ {\isasymsqinter}a{\isasymin}UNIV\ {\isasymrightarrow}\ x{\isacharparenright}{\kern0pt}{\isacartoucheclose}\ \isanewline
\ \ \ \ \ \ \ \ \ \ \ \ \ \ \ \ \ \ \ \ \ \ \ \ \ \ \ \ \ \ \ \ \ \ \ \ \ \ \ \isakeywordTWO{and}\ f\ {\isacharequal}{\kern0pt}\ {\isacartoucheopen}{\isasymlambda}x{\isachardot}{\kern0pt}\ {\isasymsqinter}a{\isasymin}UNIV\ {\isasymrightarrow}\ x{\isacartoucheclose}{\isacharbrackright}{\kern0pt}{\isacharcomma}{\kern0pt}\ simp{\isacharunderscore}{\kern0pt}all{\isacharparenright}{\kern0pt}\isanewline
\ \ \ \ \isakeywordONE{apply}\isamarkupfalse%
{\isacharparenleft}{\kern0pt}rule\ mono{\isacharunderscore}{\kern0pt}Mndetprefix{\isacharunderscore}{\kern0pt}FD{\isacharcomma}{\kern0pt}\ simp{\isacharparenright}{\kern0pt}\ \isanewline
\ \ \ \ \isakeywordONE{by}\isamarkupfalse%
{\isacharparenleft}{\kern0pt}rule\ mono{\isacharunderscore}{\kern0pt}Mndetprefix{\isacharunderscore}{\kern0pt}FD{\isacharcomma}{\kern0pt}\ simp{\isacharparenright}{\kern0pt}\ \isanewline
\isakeywordONE{qed}\isamarkupfalse%
%
\endisatagproof
{\isafoldproof}%
%
\isadelimproof
%
\endisadelimproof
%
\begin{isamarkuptext}%
Here is the neccessary versions of \isa{\isaconst{deadlock{\isacharunderscore}{\kern0pt}free}}-ness induction for
cases with a lasso-length 1 and 2:%
\end{isamarkuptext}\isamarkuptrue%
\isakeywordONE{lemma}\isamarkupfalse%
\ deadlock{\isacharunderscore}{\kern0pt}free{\isacharunderscore}{\kern0pt}{\isadigit{1}}{\isacharunderscore}{\kern0pt}coinduct\ {\isacharcolon}{\kern0pt}\isanewline
\ \ \ \ \ \ \ \ \ \ {\isacartoucheopen}\ \ \ \ {\isacharparenleft}{\kern0pt}{\isasymAnd}x{\isachardot}{\kern0pt}\ x\ {\isasymsqsubseteq}\isactrlsub F\isactrlsub D\ P\ {\isasymLongrightarrow}\ {\isasymsqinter}a{\isasymin}UNIV\ {\isasymrightarrow}\ x\ {\isasymsqsubseteq}\isactrlsub F\isactrlsub D\ P{\isacharparenright}{\kern0pt}\isanewline
\ \ \ \ \ \ \ \ \ \ \ {\isasymLongrightarrow}\ deadlock{\isacharunderscore}{\kern0pt}free\ {\isacharparenleft}{\kern0pt}P{\isacharcolon}{\kern0pt}{\isacharcolon}{\kern0pt}{\isacharparenleft}{\kern0pt}{\isacharprime}{\kern0pt}{\isasymalpha}{\isacharcomma}{\kern0pt}\ {\isacharprime}{\kern0pt}{\isasymdelta}{\isacharparenright}{\kern0pt}\ process\isactrlsub p\isactrlsub t\isactrlsub i\isactrlsub c\isactrlsub k{\isacharparenright}{\kern0pt}{\isacartoucheclose}\isanewline
%
\isadelimproof
\ \ %
\endisadelimproof
%
\isatagproof
\isakeywordONE{apply}\isamarkupfalse%
{\isacharparenleft}{\kern0pt}rule\ DF{\isacharunderscore}{\kern0pt}Univ{\isacharunderscore}{\kern0pt}freeness{\isacharbrackleft}{\kern0pt}of\ UNIV{\isacharbrackright}{\kern0pt}{\isacharcomma}{\kern0pt}\ simp{\isacharcomma}{\kern0pt}unfold\ DF{\isacharunderscore}{\kern0pt}def{\isacharparenright}{\kern0pt}\isanewline
\ \ \isakeywordONE{by}\isamarkupfalse%
{\isacharparenleft}{\kern0pt}rule\ fix{\isacharunderscore}{\kern0pt}ind{\isacharparenright}{\kern0pt}{\isacharparenleft}{\kern0pt}simp{\isacharunderscore}{\kern0pt}all{\isacharparenright}{\kern0pt}%
\endisatagproof
{\isafoldproof}%
%
\isadelimproof
\isanewline
%
\endisadelimproof
\ \isanewline
\isakeywordONE{lemma}\isamarkupfalse%
\ deadlock{\isacharunderscore}{\kern0pt}free{\isacharunderscore}{\kern0pt}{\isadigit{2}}{\isacharunderscore}{\kern0pt}coinduct\ {\isacharcolon}{\kern0pt}\isanewline
\ \ \ \ \ \ \ {\isacartoucheopen}\ {\isacharparenleft}{\kern0pt}{\isasymAnd}x{\isachardot}{\kern0pt}\ x\ {\isasymsqsubseteq}\isactrlsub F\isactrlsub D\ P\ {\isasymLongrightarrow}\ {\isasymsqinter}a{\isasymin}UNIV\ {\isasymrightarrow}\ {\isasymsqinter}a{\isasymin}UNIV\ {\isasymrightarrow}\ x\ {\isasymsqsubseteq}\isactrlsub F\isactrlsub D\ P{\isacharparenright}{\kern0pt}\isanewline
\ \ \ \ \ \ \ \ \ {\isasymLongrightarrow}\ deadlock{\isacharunderscore}{\kern0pt}free\ {\isacharparenleft}{\kern0pt}P{\isacharcolon}{\kern0pt}{\isacharcolon}{\kern0pt}{\isacharparenleft}{\kern0pt}{\isacharprime}{\kern0pt}{\isasymalpha}{\isacharcomma}{\kern0pt}\ {\isacharprime}{\kern0pt}{\isasymdelta}{\isacharparenright}{\kern0pt}\ process\isactrlsub p\isactrlsub t\isactrlsub i\isactrlsub c\isactrlsub k{\isacharparenright}{\kern0pt}{\isacartoucheclose}\isanewline
%
\isadelimproof
\ \ %
\endisadelimproof
%
\isatagproof
\isakeywordONE{apply}\isamarkupfalse%
{\isacharparenleft}{\kern0pt}rule\ DF{\isacharunderscore}{\kern0pt}Univ{\isacharunderscore}{\kern0pt}freeness{\isacharbrackleft}{\kern0pt}of\ UNIV{\isacharbrackright}{\kern0pt}{\isacharcomma}{\kern0pt}simp{\isacharunderscore}{\kern0pt}all\ add{\isacharcolon}{\kern0pt}\ lasso{\isacharunderscore}{\kern0pt}equiv{\isadigit{1}}{\isacharparenright}{\kern0pt}\isanewline
\ \ \isakeywordONE{by}\isamarkupfalse%
{\isacharparenleft}{\kern0pt}rule\ fix{\isacharunderscore}{\kern0pt}ind{\isacharparenright}{\kern0pt}{\isacharparenleft}{\kern0pt}simp{\isacharunderscore}{\kern0pt}all{\isacharparenright}{\kern0pt}%
\endisatagproof
{\isafoldproof}%
%
\isadelimproof
%
\endisadelimproof
%
\begin{isamarkuptext}%
And we need two more technical lemmas (analogously to the proofs in the previous session)
that enable us to represent the intermediate process states by compact process expressions 
in \isatt{H{\kern0pt}O{\kern0pt}L{\kern0pt}{\char`\-}{\kern0pt}C{\kern0pt}S{\kern0pt}P{\kern0pt}M{\kern0pt}}.%
\end{isamarkuptext}\isamarkuptrue%
\isakeywordONE{lemma}\isamarkupfalse%
\ contractMultiInter{\isadigit{1}}{\isacharprime}{\kern0pt}{\isacharcolon}{\kern0pt}\isanewline
\ \ \isakeywordTWO{assumes}\ {\isacartoucheopen}evs{\isachardot}{\kern0pt}unlock\ {\isadigit{1}}\ {\isasymrightarrow}\ P{\isadigit{2}}\ {\isacharequal}{\kern0pt}\ P{\isadigit{2}}{\isacartoucheclose}\isanewline
\ \ \isakeywordTWO{shows}\ \ \ {\isacartoucheopen}MultiInter\ {\isacharparenleft}{\kern0pt}mset\ {\isacharbrackleft}{\kern0pt}{\isadigit{1}}{\isachardot}{\kern0pt}{\isachardot}{\kern0pt}{\isadigit{4}}{\isacharbrackright}{\kern0pt}{\isacharparenright}{\kern0pt}\ semaphore\isanewline
\ \ \ \ \ \ \ \ \ \ \ {\isacharbar}{\kern0pt}{\isacharbar}{\kern0pt}\ \isanewline
\ \ \ \ \ \ \ \ \ \ \ {\isasymbox}x{\isasymin}{\isacharbraceleft}{\kern0pt}evs{\isachardot}{\kern0pt}lock\ {\isadigit{1}}{\isacharcomma}{\kern0pt}\ evs{\isachardot}{\kern0pt}unlock\ {\isadigit{1}}{\isacharbraceright}{\kern0pt}\isanewline
\ \ \ \ \ \ \ \ \ \ \ \ \ \ {\isasymrightarrow}\ {\isacharparenleft}{\kern0pt}if\ x\ {\isacharequal}{\kern0pt}\ evs{\isachardot}{\kern0pt}lock\ {\isadigit{1}}\ then\ P{\isadigit{1}}\ {\isacharbar}{\kern0pt}{\isacharbar}{\kern0pt}{\isacharbar}{\kern0pt}\ evs{\isachardot}{\kern0pt}unlock\ {\isadigit{1}}\ {\isasymrightarrow}\ P{\isadigit{2}}\ else\ evs{\isachardot}{\kern0pt}lock\ {\isadigit{1}}\ {\isasymrightarrow}\ P{\isadigit{1}}\ {\isacharbar}{\kern0pt}{\isacharbar}{\kern0pt}{\isacharbar}{\kern0pt}\ P{\isadigit{2}}{\isacharparenright}{\kern0pt}\ {\isacharequal}{\kern0pt}\isanewline
\ \ \ \ \ \ \ \ \ \ \ evs{\isachardot}{\kern0pt}lock\ {\isadigit{1}}\ {\isasymrightarrow}\isanewline
\ \ \ \ \ \ \ \ \ \ \ {\isacharparenleft}{\kern0pt}\ \ \ {\isacharparenleft}{\kern0pt}evs{\isachardot}{\kern0pt}unlock\ {\isadigit{1}}\ {\isasymrightarrow}\ semaphore\ {\isadigit{1}}\ {\isacharbar}{\kern0pt}{\isacharbar}{\kern0pt}{\isacharbar}{\kern0pt}\ MultiInter\ {\isacharparenleft}{\kern0pt}mset\ {\isacharbrackleft}{\kern0pt}{\isadigit{2}}{\isachardot}{\kern0pt}{\isachardot}{\kern0pt}{\isadigit{4}}{\isacharbrackright}{\kern0pt}{\isacharparenright}{\kern0pt}\ semaphore{\isacharparenright}{\kern0pt}\ \isanewline
\ \ \ \ \ \ \ \ \ \ \ {\isacharbar}{\kern0pt}{\isacharbar}{\kern0pt}\ {\isacharparenleft}{\kern0pt}P{\isadigit{1}}\ {\isacharbar}{\kern0pt}{\isacharbar}{\kern0pt}{\isacharbar}{\kern0pt}\ P{\isadigit{2}}{\isacharparenright}{\kern0pt}{\isacharparenright}{\kern0pt}{\isacartoucheclose}\isanewline
\ \ {\isacharparenleft}{\kern0pt}\isakeywordTWO{is}\ {\isachardoublequoteopen}{\isacharquery}{\kern0pt}lhs\ {\isacharequal}{\kern0pt}\ evs{\isachardot}{\kern0pt}lock\ {\isadigit{1}}\ {\isasymrightarrow}\ {\isacharquery}{\kern0pt}rhs{\isachardoublequoteclose}{\isacharparenright}{\kern0pt}\isanewline
%
\isadelimproof
%
\endisadelimproof
%
\isatagproof
\isakeywordONE{proof}\isamarkupfalse%
\ {\isacharminus}{\kern0pt}\ \isanewline
\ \ \ \isakeywordONE{have}\isamarkupfalse%
\ {\isacharasterisk}{\kern0pt}\ {\isacharcolon}{\kern0pt}\ {\isachardoublequoteopen}{\isacharquery}{\kern0pt}lhs\ {\isacharequal}{\kern0pt}\ evs{\isachardot}{\kern0pt}lock\ {\isadigit{1}}\ {\isasymrightarrow}\isanewline
\ \ \ \ \ \ \ \ \ \ \ \ \ \ \ \ \ \ \ \ {\isacharparenleft}{\kern0pt}{\isacharparenleft}{\kern0pt}evs{\isachardot}{\kern0pt}unlock\ {\isadigit{1}}\ {\isasymrightarrow}\ semaphore\ {\isadigit{1}}\ \isanewline
\ \ \ \ \ \ \ \ \ \ \ \ \ \ \ \ \ \ \ \ \ {\isacharbar}{\kern0pt}{\isacharbar}{\kern0pt}{\isacharbar}{\kern0pt}\ {\isasymbox}x{\isasymin}{\isacharbraceleft}{\kern0pt}evs{\isachardot}{\kern0pt}lock\ {\isadigit{2}}{\isacharcomma}{\kern0pt}\ evs{\isachardot}{\kern0pt}lock\ {\isadigit{3}}{\isacharcomma}{\kern0pt}\ evs{\isachardot}{\kern0pt}lock\ {\isadigit{4}}{\isacharbraceright}{\kern0pt}\isanewline
\ \ \ \ \ \ \ \ \ \ \ \ \ \ \ \ \ \ \ \ \ \ \ \ \ \ \ \ \ \ \ \ {\isasymrightarrow}\ {\isacharparenleft}{\kern0pt}if\ x\ {\isacharequal}{\kern0pt}\ evs{\isachardot}{\kern0pt}lock\ {\isadigit{2}}\isanewline
\ \ \ \ \ \ \ \ \ \ \ \ \ \ \ \ \ \ \ \ \ \ \ \ \ \ \ \ \ \ \ \ \ \ \ \ then\ evs{\isachardot}{\kern0pt}unlock\ {\isadigit{2}}\ {\isasymrightarrow}\ semaphore\ {\isadigit{2}}\ \isanewline
\ \ \ \ \ \ \ \ \ \ \ \ \ \ \ \ \ \ \ \ \ \ \ \ \ \ \ \ \ \ \ \ \ \ \ \ \ \ \ \ \ {\isacharbar}{\kern0pt}{\isacharbar}{\kern0pt}{\isacharbar}{\kern0pt}\ \isanewline
\ \ \ \ \ \ \ \ \ \ \ \ \ \ \ \ \ \ \ \ \ \ \ \ \ \ \ \ \ \ \ \ \ \ \ \ \ \ \ \ \ {\isasymbox}x{\isasymin}{\isacharbraceleft}{\kern0pt}evs{\isachardot}{\kern0pt}lock\ {\isadigit{3}}{\isacharcomma}{\kern0pt}\ evs{\isachardot}{\kern0pt}lock\ {\isadigit{4}}{\isacharbraceright}{\kern0pt}\isanewline
\ \ \ \ \ \ \ \ \ \ \ \ \ \ \ \ \ \ \ \ \ \ \ \ \ \ \ \ \ \ \ \ \ \ \ \ \ \ \ \ \ \ \ \ \ \ \ {\isasymrightarrow}\ {\isacharparenleft}{\kern0pt}if\ x\ {\isacharequal}{\kern0pt}\ evs{\isachardot}{\kern0pt}lock\ {\isadigit{3}}\ \isanewline
\ \ \ \ \ \ \ \ \ \ \ \ \ \ \ \ \ \ \ \ \ \ \ \ \ \ \ \ \ \ \ \ \ \ \ \ \ \ \ \ \ \ \ \ \ \ \ \ \ \ \ then\ evs{\isachardot}{\kern0pt}unlock\ {\isadigit{3}}\ {\isasymrightarrow}\ semaphore\ {\isadigit{3}}\ {\isacharbar}{\kern0pt}{\isacharbar}{\kern0pt}{\isacharbar}{\kern0pt}\ semaphore\ {\isadigit{4}}\isanewline
\ \ \ \ \ \ \ \ \ \ \ \ \ \ \ \ \ \ \ \ \ \ \ \ \ \ \ \ \ \ \ \ \ \ \ \ \ \ \ \ \ \ \ \ \ \ \ \ \ \ \ else\ semaphore\ {\isadigit{3}}\ {\isacharbar}{\kern0pt}{\isacharbar}{\kern0pt}{\isacharbar}{\kern0pt}\ evs{\isachardot}{\kern0pt}unlock\ {\isadigit{4}}\ {\isasymrightarrow}\ semaphore\ {\isadigit{4}}{\isacharparenright}{\kern0pt}\isanewline
\ \ \ \ \ \ \ \ \ \ \ \ \ \ \ \ \ \ \ \ \ \ \ \ \ \ \ \ \ \ \ \ \ \ \ \ else\ semaphore\ {\isadigit{2}}\ \isanewline
\ \ \ \ \ \ \ \ \ \ \ \ \ \ \ \ \ \ \ \ \ \ \ \ \ \ \ \ \ \ \ \ \ \ \ \ \ \ \ \ \ {\isacharbar}{\kern0pt}{\isacharbar}{\kern0pt}{\isacharbar}{\kern0pt}\ \isanewline
\ \ \ \ \ \ \ \ \ \ \ \ \ \ \ \ \ \ \ \ \ \ \ \ \ \ \ \ \ \ \ \ \ \ \ \ \ \ \ \ \ {\isacharparenleft}{\kern0pt}if\ x\ {\isacharequal}{\kern0pt}\ evs{\isachardot}{\kern0pt}lock\ {\isadigit{3}}\ \isanewline
\ \ \ \ \ \ \ \ \ \ \ \ \ \ \ \ \ \ \ \ \ \ \ \ \ \ \ \ \ \ \ \ \ \ \ \ \ \ \ \ \ \ then\ evs{\isachardot}{\kern0pt}unlock\ {\isadigit{3}}\ {\isasymrightarrow}\ semaphore\ {\isadigit{3}}\ {\isacharbar}{\kern0pt}{\isacharbar}{\kern0pt}{\isacharbar}{\kern0pt}\ semaphore\ {\isadigit{4}}\isanewline
\ \ \ \ \ \ \ \ \ \ \ \ \ \ \ \ \ \ \ \ \ \ \ \ \ \ \ \ \ \ \ \ \ \ \ \ \ \ \ \ \ \ else\ semaphore\ {\isadigit{3}}\ {\isacharbar}{\kern0pt}{\isacharbar}{\kern0pt}{\isacharbar}{\kern0pt}\ evs{\isachardot}{\kern0pt}unlock\ {\isadigit{4}}\ {\isasymrightarrow}\ semaphore\ {\isadigit{4}}{\isacharparenright}{\kern0pt}{\isacharparenright}{\kern0pt}{\isacharparenright}{\kern0pt}\ \isanewline
\ \ \ \ \ \ \ \ \ \ \ \ \ \ \ \ \ \ \ \ {\isacharbar}{\kern0pt}{\isacharbar}{\kern0pt}\ {\isacharparenleft}{\kern0pt}P{\isadigit{1}}\ {\isacharbar}{\kern0pt}{\isacharbar}{\kern0pt}{\isacharbar}{\kern0pt}\ evs{\isachardot}{\kern0pt}unlock\ {\isadigit{1}}\ {\isasymrightarrow}\ P{\isadigit{2}}{\isacharparenright}{\kern0pt}{\isacharparenright}{\kern0pt}{\isachardoublequoteclose}\isanewline
\ \ \ \ {\isacharparenleft}{\kern0pt}\isakeywordTWO{is}\ {\isachardoublequoteopen}{\isacharunderscore}{\kern0pt}\ {\isacharequal}{\kern0pt}\ evs{\isachardot}{\kern0pt}lock\ {\isadigit{1}}\ {\isasymrightarrow}\ {\isacharquery}{\kern0pt}rhs{\isacharprime}{\kern0pt}{\isachardoublequoteclose}{\isacharparenright}{\kern0pt}\isanewline
\ \ \ \ \ \ \ \ \ \ \isakeywordONE{apply}\isamarkupfalse%
{\isacharparenleft}{\kern0pt}rule\ trans{\isacharparenright}{\kern0pt}\isanewline
\ \ \ \ \ \ \ \ \ \ \ \isakeywordONE{apply}\isamarkupfalse%
{\isacharparenleft}{\kern0pt}simp\ add{\isacharcolon}{\kern0pt}\ List{\isachardot}{\kern0pt}upto{\isachardot}{\kern0pt}simps{\isacharparenright}{\kern0pt}\isanewline
\ \ \ \ \ \ \ \ \ \ \isakeywordONE{apply}\isamarkupfalse%
{\isacharparenleft}{\kern0pt}subst\ sema{\isacharunderscore}{\kern0pt}rec{\isacharbrackleft}{\kern0pt}\isakeywordTWO{where}\ n{\isacharequal}{\kern0pt}{\isacartoucheopen}{\isadigit{1}}{\isacartoucheclose}{\isacharbrackright}{\kern0pt}{\isacharcomma}{\kern0pt}subst\ sema{\isacharunderscore}{\kern0pt}rec{\isacharbrackleft}{\kern0pt}\isakeywordTWO{where}\ n{\isacharequal}{\kern0pt}{\isacartoucheopen}{\isadigit{2}}{\isacartoucheclose}{\isacharbrackright}{\kern0pt}{\isacharcomma}{\kern0pt}\isanewline
\ \ \ \ \ \ \ \ \ \ \ \ \ \ \ \ subst\ sema{\isacharunderscore}{\kern0pt}rec{\isacharbrackleft}{\kern0pt}\isakeywordTWO{where}\ n{\isacharequal}{\kern0pt}{\isacartoucheopen}{\isadigit{3}}{\isacartoucheclose}{\isacharbrackright}{\kern0pt}{\isacharcomma}{\kern0pt}subst\ sema{\isacharunderscore}{\kern0pt}rec{\isacharbrackleft}{\kern0pt}\isakeywordTWO{where}\ n{\isacharequal}{\kern0pt}{\isacartoucheopen}{\isadigit{4}}{\isacartoucheclose}{\isacharbrackright}{\kern0pt}{\isacharparenright}{\kern0pt}\ \isanewline
\ \ \ \ \ \ \ \ \ \ \isakeywordONE{apply}\isamarkupfalse%
{\isacharparenleft}{\kern0pt}rule\ trans{\isacharparenright}{\kern0pt}\isanewline
\ \ \ \ \ \ \ \ \ \ \isakeywordONE{apply}\isamarkupfalse%
{\isacharparenleft}{\kern0pt}simp\ add{\isacharcolon}{\kern0pt}\ inter{\isadigit{2}}Mprefix{\isadigit{1}}\ inter{\isadigit{2}}Mprefix{\isadigit{2}}{\isacharparenright}{\kern0pt}\isanewline
\ \ \ \ \ \ \ \ \ \ \isakeywordONE{apply}\isamarkupfalse%
{\isacharparenleft}{\kern0pt}simp\ add{\isacharcolon}{\kern0pt}\ sema{\isacharunderscore}{\kern0pt}rec{\isacharbrackleft}{\kern0pt}\isakeywordTWO{where}\ n{\isacharequal}{\kern0pt}{\isacartoucheopen}{\isadigit{1}}{\isacartoucheclose}{\isacharcomma}{\kern0pt}symmetric{\isacharbrackright}{\kern0pt}\ \isanewline
\ \ \ \ \ \ \ \ \ \ \ \ \ \ \ \ \ \ \ \ \ \ \ \ \ \ sema{\isacharunderscore}{\kern0pt}rec{\isacharbrackleft}{\kern0pt}\isakeywordTWO{where}\ n{\isacharequal}{\kern0pt}{\isacartoucheopen}{\isadigit{2}}{\isacartoucheclose}{\isacharcomma}{\kern0pt}symmetric{\isacharbrackright}{\kern0pt}\ \isanewline
\ \ \ \ \ \ \ \ \ \ \ \ \ \ \ \ \ \ \ \ \ \ \ \ \ \ sema{\isacharunderscore}{\kern0pt}rec{\isacharbrackleft}{\kern0pt}\isakeywordTWO{where}\ n{\isacharequal}{\kern0pt}{\isacartoucheopen}{\isadigit{3}}{\isacartoucheclose}{\isacharcomma}{\kern0pt}symmetric{\isacharbrackright}{\kern0pt}\ \isanewline
\ \ \ \ \ \ \ \ \ \ \ \ \ \ \ \ \ \ \ \ \ \ \ \ \ \ sema{\isacharunderscore}{\kern0pt}rec{\isacharbrackleft}{\kern0pt}\isakeywordTWO{where}\ n{\isacharequal}{\kern0pt}{\isacartoucheopen}{\isadigit{4}}{\isacartoucheclose}{\isacharcomma}{\kern0pt}symmetric{\isacharbrackright}{\kern0pt}\isanewline
\ \ \ \ \ \ \ \ \ \ \ \ \ \ \ \ \ \ \ \ \ cong{\isacharcolon}{\kern0pt}\ HOL{\isachardot}{\kern0pt}if{\isacharunderscore}{\kern0pt}cong{\isacharparenright}{\kern0pt}\isanewline
\ \ \ \ \ \ \ \ \ \ \isakeywordONE{apply}\isamarkupfalse%
{\isacharparenleft}{\kern0pt}subst\ write{\isadigit{0}}{\isacharunderscore}{\kern0pt}def{\isacharbrackleft}{\kern0pt}THEN\ meta{\isacharunderscore}{\kern0pt}eq{\isacharunderscore}{\kern0pt}to{\isacharunderscore}{\kern0pt}obj{\isacharunderscore}{\kern0pt}eq{\isacharcomma}{\kern0pt}\ of{\isachardoublequoteopen}evs{\isachardot}{\kern0pt}lock\ {\isacharparenleft}{\kern0pt}{\isadigit{1}}{\isacharcolon}{\kern0pt}{\isacharcolon}{\kern0pt}int{\isacharparenright}{\kern0pt}{\isachardoublequoteclose}{\isacharbrackright}{\kern0pt}{\isacharparenright}{\kern0pt}\isanewline
\ \ \ \ \ \ \ \ \ \ \isakeywordONE{by}\isamarkupfalse%
{\isacharparenleft}{\kern0pt}subst\ Mprefix{\isacharunderscore}{\kern0pt}Par{\isacharunderscore}{\kern0pt}Mprefix{\isacharcomma}{\kern0pt}\ simp\ add{\isacharcolon}{\kern0pt}\ Mprefix{\isacharunderscore}{\kern0pt}singl{\isacharparenright}{\kern0pt}\isanewline
\ \ \isakeywordTHREE{define}\isamarkupfalse%
\ X\ \isakeywordTWO{where}\ {\isachardoublequoteopen}X\ {\isacharequal}{\kern0pt}\ {\isacharquery}{\kern0pt}rhs{\isacharprime}{\kern0pt}{\isachardoublequoteclose}\isanewline
\ \ \isakeywordONE{have}\isamarkupfalse%
\ {\isacharasterisk}{\kern0pt}{\isacharasterisk}{\kern0pt}\ {\isacharcolon}{\kern0pt}\ {\isachardoublequoteopen}{\isacharquery}{\kern0pt}rhs\ {\isacharequal}{\kern0pt}\ X{\isachardoublequoteclose}\ \isanewline
\ \ \ \ \ \ \ \ \ \ \isakeywordONE{apply}\isamarkupfalse%
{\isacharparenleft}{\kern0pt}subst\ write{\isadigit{0}}{\isacharunderscore}{\kern0pt}def{\isacharbrackleft}{\kern0pt}THEN\ meta{\isacharunderscore}{\kern0pt}eq{\isacharunderscore}{\kern0pt}to{\isacharunderscore}{\kern0pt}obj{\isacharunderscore}{\kern0pt}eq{\isacharbrackright}{\kern0pt}{\isacharparenright}{\kern0pt}\isanewline
\ \ \ \ \ \ \ \ \ \ \isakeywordONE{apply}\isamarkupfalse%
{\isacharparenleft}{\kern0pt}simp\ add{\isacharcolon}{\kern0pt}\ List{\isachardot}{\kern0pt}upto{\isachardot}{\kern0pt}simps{\isacharparenright}{\kern0pt}\ \ \isanewline
\ \ \ \ \ \ \ \ \ \ \isakeywordONE{apply}\isamarkupfalse%
{\isacharparenleft}{\kern0pt}subst\ sema{\isacharunderscore}{\kern0pt}rec{\isacharbrackleft}{\kern0pt}\isakeywordTWO{where}\ n{\isacharequal}{\kern0pt}{\isacartoucheopen}{\isadigit{1}}{\isacartoucheclose}{\isacharbrackright}{\kern0pt}{\isacharcomma}{\kern0pt}subst\ sema{\isacharunderscore}{\kern0pt}rec{\isacharbrackleft}{\kern0pt}\isakeywordTWO{where}\ n{\isacharequal}{\kern0pt}{\isacartoucheopen}{\isadigit{2}}{\isacartoucheclose}{\isacharbrackright}{\kern0pt}{\isacharcomma}{\kern0pt}\isanewline
\ \ \ \ \ \ \ \ \ \ \ \ \ \ \ \ subst\ sema{\isacharunderscore}{\kern0pt}rec{\isacharbrackleft}{\kern0pt}\isakeywordTWO{where}\ n{\isacharequal}{\kern0pt}{\isacartoucheopen}{\isadigit{3}}{\isacartoucheclose}{\isacharbrackright}{\kern0pt}{\isacharcomma}{\kern0pt}subst\ sema{\isacharunderscore}{\kern0pt}rec{\isacharbrackleft}{\kern0pt}\isakeywordTWO{where}\ n{\isacharequal}{\kern0pt}{\isacartoucheopen}{\isadigit{4}}{\isacartoucheclose}{\isacharbrackright}{\kern0pt}{\isacharparenright}{\kern0pt}\ \isanewline
\ \ \ \ \ \ \ \ \ \ \isakeywordONE{apply}\isamarkupfalse%
{\isacharparenleft}{\kern0pt}simp\ add{\isacharcolon}{\kern0pt}\ inter{\isadigit{2}}Mprefix{\isadigit{1}}\ inter{\isadigit{2}}Mprefix{\isadigit{2}}{\isacharparenright}{\kern0pt}\ \isanewline
\ \ \ \ \ \ \ \ \ \ \isakeywordONE{apply}\isamarkupfalse%
{\isacharparenleft}{\kern0pt}simp\ only{\isacharcolon}{\kern0pt}\ X{\isacharunderscore}{\kern0pt}def{\isacharparenright}{\kern0pt}\ \isanewline
\ \ \ \ \ \ \ \ \ \ \isakeywordONE{by}\isamarkupfalse%
\ {\isacharparenleft}{\kern0pt}metis\ sema{\isacharunderscore}{\kern0pt}rec\ assms\ Mprefix{\isacharunderscore}{\kern0pt}singl{\isacharparenright}{\kern0pt}\isanewline
\ \ \isakeywordTHREE{show}\isamarkupfalse%
\ {\isacharquery}{\kern0pt}thesis\isanewline
\ \ \ \ \isakeywordONE{by}\isamarkupfalse%
\ {\isacharparenleft}{\kern0pt}metis\ {\isacharasterisk}{\kern0pt}\ {\isacharasterisk}{\kern0pt}{\isacharasterisk}{\kern0pt}\ X{\isacharunderscore}{\kern0pt}def{\isacharparenright}{\kern0pt}\isanewline
\isakeywordONE{qed}\isamarkupfalse%
%
\endisatagproof
{\isafoldproof}%
%
\isadelimproof
\isanewline
%
\endisadelimproof
\isanewline
\isanewline
\isakeywordONE{lemma}\isamarkupfalse%
\ contractMultiInter{\isadigit{1}}{\isacharprime}{\kern0pt}{\isacharprime}{\kern0pt}\ {\isacharcolon}{\kern0pt}\ \isanewline
\ \ \isakeywordTWO{assumes}\ {\isachardoublequoteopen}evs{\isachardot}{\kern0pt}lock\ {\isadigit{1}}\ {\isasymrightarrow}\ P{\isadigit{1}}\ {\isacharequal}{\kern0pt}\ P{\isadigit{1}}{\isachardoublequoteclose}\isanewline
\ \ \isakeywordTWO{shows}\ {\isacartoucheopen}\isanewline
\ \ \ \ \ \ \ \ \ \ {\isacharparenleft}{\kern0pt}evs{\isachardot}{\kern0pt}unlock\ {\isadigit{1}}\ {\isasymrightarrow}\ semaphore\ {\isadigit{1}}\ \ {\isacharbar}{\kern0pt}{\isacharbar}{\kern0pt}{\isacharbar}{\kern0pt}\ MultiInter\ {\isacharparenleft}{\kern0pt}mset\ {\isacharbrackleft}{\kern0pt}{\isadigit{2}}{\isachardot}{\kern0pt}{\isachardot}{\kern0pt}{\isadigit{4}}{\isacharbrackright}{\kern0pt}{\isacharparenright}{\kern0pt}\ semaphore{\isacharparenright}{\kern0pt}\ \isanewline
\ \ \ \ \ \ \ \ {\isacharbar}{\kern0pt}{\isacharbar}{\kern0pt}\ {\isasymbox}x{\isasymin}{\isacharbraceleft}{\kern0pt}evs{\isachardot}{\kern0pt}lock\ {\isadigit{1}}{\isacharcomma}{\kern0pt}\ evs{\isachardot}{\kern0pt}unlock\ {\isadigit{1}}{\isacharbraceright}{\kern0pt}\isanewline
\ \ \ \ \ \ \ \ \ \ \ \ \ \ {\isasymrightarrow}\ {\isacharparenleft}{\kern0pt}if\ x\ {\isacharequal}{\kern0pt}\ evs{\isachardot}{\kern0pt}lock\ {\isadigit{1}}\ \isanewline
\ \ \ \ \ \ \ \ \ \ \ \ \ \ \ \ \ \ then\ P{\isadigit{1}}\ {\isacharbar}{\kern0pt}{\isacharbar}{\kern0pt}{\isacharbar}{\kern0pt}\ evs{\isachardot}{\kern0pt}unlock\ {\isadigit{1}}\ {\isasymrightarrow}\ P{\isadigit{2}}\ \isanewline
\ \ \ \ \ \ \ \ \ \ \ \ \ \ \ \ \ \ else\ evs{\isachardot}{\kern0pt}lock\ {\isadigit{1}}\ {\isasymrightarrow}\ P{\isadigit{1}}\ {\isacharbar}{\kern0pt}{\isacharbar}{\kern0pt}{\isacharbar}{\kern0pt}\ P{\isadigit{2}}{\isacharparenright}{\kern0pt}\ \isanewline
\ \ \ \ \ \ \ \ {\isacharequal}{\kern0pt}\ \isanewline
\ \ \ \ \ \ \ \ \ \ {\isacharparenleft}{\kern0pt}evs{\isachardot}{\kern0pt}unlock\ {\isadigit{1}}\ {\isasymrightarrow}\ {\isacharparenleft}{\kern0pt}{\isacharparenleft}{\kern0pt}MultiInter\ {\isacharparenleft}{\kern0pt}mset\ {\isacharbrackleft}{\kern0pt}{\isadigit{1}}{\isachardot}{\kern0pt}{\isachardot}{\kern0pt}{\isadigit{4}}{\isacharbrackright}{\kern0pt}{\isacharparenright}{\kern0pt}\ semaphore{\isacharparenright}{\kern0pt}\ \isanewline
\ \ \ \ \ \ \ \ \ \ \ \ \ \ \ \ \ \ \ \ \ \ \ \ \ \ \ \ {\isacharbar}{\kern0pt}{\isacharbar}{\kern0pt}\ {\isacharparenleft}{\kern0pt}P{\isadigit{1}}\ \ {\isacharbar}{\kern0pt}{\isacharbar}{\kern0pt}{\isacharbar}{\kern0pt}\ P{\isadigit{2}}{\isacharparenright}{\kern0pt}{\isacharparenright}{\kern0pt}{\isacharparenright}{\kern0pt}\isanewline
\ \ \ \ \ \ \ \ {\isacartoucheclose}\isanewline
\ \ \ \ \ \ \ {\isacharparenleft}{\kern0pt}\isakeywordTWO{is}\ {\isachardoublequoteopen}{\isacharparenleft}{\kern0pt}{\isacharquery}{\kern0pt}R\ {\isacharbar}{\kern0pt}{\isacharbar}{\kern0pt}\ {\isacharquery}{\kern0pt}S{\isacharparenright}{\kern0pt}\ {\isacharequal}{\kern0pt}\ evs{\isachardot}{\kern0pt}unlock\ {\isadigit{1}}\ {\isasymrightarrow}\ {\isacharquery}{\kern0pt}rhs{\isachardoublequoteclose}{\isacharparenright}{\kern0pt}\ \isanewline
%
\isadelimproof
%
\endisadelimproof
%
\isatagproof
\isakeywordONE{proof}\isamarkupfalse%
\ {\isacharminus}{\kern0pt}\isanewline
\ \ \isakeywordONE{have}\isamarkupfalse%
\ {\isadigit{1}}\ {\isacharcolon}{\kern0pt}\ {\isachardoublequoteopen}{\isacharquery}{\kern0pt}R\ {\isacharequal}{\kern0pt}\ {\isasymbox}x{\isasymin}{\isacharbraceleft}{\kern0pt}evs{\isachardot}{\kern0pt}unlock\ {\isadigit{1}}{\isacharcomma}{\kern0pt}\ evs{\isachardot}{\kern0pt}lock\ {\isadigit{2}}{\isacharcomma}{\kern0pt}\ evs{\isachardot}{\kern0pt}lock\ {\isadigit{3}}{\isacharcomma}{\kern0pt}\ evs{\isachardot}{\kern0pt}lock\ {\isadigit{4}}{\isacharbraceright}{\kern0pt}\isanewline
\ \ \ \ \ \ \ \ \ \ \ \ \ \ \ \ \ {\isasymrightarrow}\ {\isacharparenleft}{\kern0pt}if\ x\ {\isacharequal}{\kern0pt}\ evs{\isachardot}{\kern0pt}unlock\ {\isadigit{1}}\isanewline
\ \ \ \ \ \ \ \ \ \ \ \ \ \ \ \ \ \ \ \ \ then\ semaphore\ {\isadigit{1}}\isanewline
\ \ \ \ \ \ \ \ \ \ \ \ \ \ \ \ \ \ \ \ \ \ \ \ \ \ {\isacharbar}{\kern0pt}{\isacharbar}{\kern0pt}{\isacharbar}{\kern0pt}\ {\isasymbox}x{\isasymin}{\isacharbraceleft}{\kern0pt}evs{\isachardot}{\kern0pt}lock\ {\isadigit{2}}{\isacharcomma}{\kern0pt}\ evs{\isachardot}{\kern0pt}lock\ {\isadigit{3}}{\isacharcomma}{\kern0pt}\ evs{\isachardot}{\kern0pt}lock\ {\isadigit{4}}{\isacharbraceright}{\kern0pt}\isanewline
\ \ \ \ \ \ \ \ \ \ \ \ \ \ \ \ \ \ \ \ \ \ \ \ \ \ \ \ \ \ {\isasymrightarrow}\ {\isacharparenleft}{\kern0pt}if\ x\ {\isacharequal}{\kern0pt}\ evs{\isachardot}{\kern0pt}lock\ {\isadigit{2}}\isanewline
\ \ \ \ \ \ \ \ \ \ \ \ \ \ \ \ \ \ \ \ \ \ \ \ \ \ \ \ \ \ \ \ \ \ then\ evs{\isachardot}{\kern0pt}unlock\ {\isadigit{2}}\ {\isasymrightarrow}\ semaphore\ {\isadigit{2}}\isanewline
\ \ \ \ \ \ \ \ \ \ \ \ \ \ \ \ \ \ \ \ \ \ \ \ \ \ \ \ \ \ \ \ \ \ \ \ \ \ \ {\isacharbar}{\kern0pt}{\isacharbar}{\kern0pt}{\isacharbar}{\kern0pt}\ {\isasymbox}x{\isasymin}{\isacharbraceleft}{\kern0pt}evs{\isachardot}{\kern0pt}lock\ {\isadigit{3}}{\isacharcomma}{\kern0pt}\ evs{\isachardot}{\kern0pt}lock\ {\isadigit{4}}{\isacharbraceright}{\kern0pt}\isanewline
\ \ \ \ \ \ \ \ \ \ \ \ \ \ \ \ \ \ \ \ \ \ \ \ \ \ \ \ \ \ \ \ \ \ \ \ \ \ \ \ \ \ \ {\isasymrightarrow}\ {\isacharparenleft}{\kern0pt}if\ x\ {\isacharequal}{\kern0pt}\ evs{\isachardot}{\kern0pt}lock\ {\isadigit{3}}\ \isanewline
\ \ \ \ \ \ \ \ \ \ \ \ \ \ \ \ \ \ \ \ \ \ \ \ \ \ \ \ \ \ \ \ \ \ \ \ \ \ \ \ \ \ \ \ \ \ \ then\ evs{\isachardot}{\kern0pt}unlock\ {\isadigit{3}}\ {\isasymrightarrow}\ semaphore\ {\isadigit{3}}\ {\isacharbar}{\kern0pt}{\isacharbar}{\kern0pt}{\isacharbar}{\kern0pt}\ semaphore\ {\isadigit{4}}\isanewline
\ \ \ \ \ \ \ \ \ \ \ \ \ \ \ \ \ \ \ \ \ \ \ \ \ \ \ \ \ \ \ \ \ \ \ \ \ \ \ \ \ \ \ \ \ \ \ else\ semaphore\ {\isadigit{3}}\ {\isacharbar}{\kern0pt}{\isacharbar}{\kern0pt}{\isacharbar}{\kern0pt}\ evs{\isachardot}{\kern0pt}unlock\ {\isadigit{4}}\ {\isasymrightarrow}\ semaphore\ {\isadigit{4}}{\isacharparenright}{\kern0pt}\isanewline
\ \ \ \ \ \ \ \ \ \ \ \ \ \ \ \ \ \ \ \ \ \ \ \ \ \ \ \ \ \ \ \ \ \ else\ semaphore\ \ {\isadigit{2}}\isanewline
\ \ \ \ \ \ \ \ \ \ \ \ \ \ \ \ \ \ \ \ \ \ \ \ \ \ \ \ \ \ \ \ \ \ \ \ \ \ \ {\isacharbar}{\kern0pt}{\isacharbar}{\kern0pt}{\isacharbar}{\kern0pt}\ {\isacharparenleft}{\kern0pt}if\ x\ {\isacharequal}{\kern0pt}\ evs{\isachardot}{\kern0pt}lock\ {\isadigit{3}}\ \isanewline
\ \ \ \ \ \ \ \ \ \ \ \ \ \ \ \ \ \ \ \ \ \ \ \ \ \ \ \ \ \ \ \ \ \ \ \ \ \ \ \ \ \ \ \ then\ evs{\isachardot}{\kern0pt}unlock\ {\isadigit{3}}\ {\isasymrightarrow}\ semaphore\ {\isadigit{3}}\ {\isacharbar}{\kern0pt}{\isacharbar}{\kern0pt}{\isacharbar}{\kern0pt}\ semaphore\ {\isadigit{4}}\isanewline
\ \ \ \ \ \ \ \ \ \ \ \ \ \ \ \ \ \ \ \ \ \ \ \ \ \ \ \ \ \ \ \ \ \ \ \ \ \ \ \ \ \ \ \ else\ semaphore\ {\isadigit{3}}\ {\isacharbar}{\kern0pt}{\isacharbar}{\kern0pt}{\isacharbar}{\kern0pt}\ evs{\isachardot}{\kern0pt}unlock\ {\isadigit{4}}\ {\isasymrightarrow}\ semaphore\ {\isadigit{4}}{\isacharparenright}{\kern0pt}{\isacharparenright}{\kern0pt}\isanewline
\ \ \ \ \ \ \ \ \ \ \ \ \ \ \ \ \ \ \ \ \ else\ evs{\isachardot}{\kern0pt}unlock\ {\isadigit{1}}\ {\isasymrightarrow}\ semaphore\ {\isadigit{1}}\isanewline
\ \ \ \ \ \ \ \ \ \ \ \ \ \ \ \ \ \ \ \ \ \ \ \ \ \ {\isacharbar}{\kern0pt}{\isacharbar}{\kern0pt}{\isacharbar}{\kern0pt}\ {\isacharparenleft}{\kern0pt}if\ x\ {\isacharequal}{\kern0pt}\ evs{\isachardot}{\kern0pt}lock\ {\isadigit{2}}\isanewline
\ \ \ \ \ \ \ \ \ \ \ \ \ \ \ \ \ \ \ \ \ \ \ \ \ \ \ \ \ \ \ then\ evs{\isachardot}{\kern0pt}unlock\ {\isadigit{2}}\ {\isasymrightarrow}\ semaphore\ {\isadigit{2}}\isanewline
\ \ \ \ \ \ \ \ \ \ \ \ \ \ \ \ \ \ \ \ \ \ \ \ \ \ \ \ \ \ \ \ \ \ \ \ {\isacharbar}{\kern0pt}{\isacharbar}{\kern0pt}{\isacharbar}{\kern0pt}\ {\isasymbox}x{\isasymin}{\isacharbraceleft}{\kern0pt}evs{\isachardot}{\kern0pt}lock\ {\isadigit{3}}{\isacharcomma}{\kern0pt}\ evs{\isachardot}{\kern0pt}lock\ {\isadigit{4}}{\isacharbraceright}{\kern0pt}\isanewline
\ \ \ \ \ \ \ \ \ \ \ \ \ \ \ \ \ \ \ \ \ \ \ \ \ \ \ \ \ \ \ \ \ \ \ \ \ \ \ \ {\isasymrightarrow}\ {\isacharparenleft}{\kern0pt}if\ x\ {\isacharequal}{\kern0pt}\ evs{\isachardot}{\kern0pt}lock\ {\isadigit{3}}\ \isanewline
\ \ \ \ \ \ \ \ \ \ \ \ \ \ \ \ \ \ \ \ \ \ \ \ \ \ \ \ \ \ \ \ \ \ \ \ \ \ \ \ \ \ \ \ then\ evs{\isachardot}{\kern0pt}unlock\ {\isadigit{3}}\ {\isasymrightarrow}\ semaphore\ {\isadigit{3}}\ {\isacharbar}{\kern0pt}{\isacharbar}{\kern0pt}{\isacharbar}{\kern0pt}\ semaphore\ {\isadigit{4}}\isanewline
\ \ \ \ \ \ \ \ \ \ \ \ \ \ \ \ \ \ \ \ \ \ \ \ \ \ \ \ \ \ \ \ \ \ \ \ \ \ \ \ \ \ \ \ else\ semaphore\ {\isadigit{3}}\ {\isacharbar}{\kern0pt}{\isacharbar}{\kern0pt}{\isacharbar}{\kern0pt}\ evs{\isachardot}{\kern0pt}unlock\ {\isadigit{4}}\ {\isasymrightarrow}\ semaphore\ {\isadigit{4}}{\isacharparenright}{\kern0pt}\isanewline
\ \ \ \ \ \ \ \ \ \ \ \ \ \ \ \ \ \ \ \ \ \ \ \ \ \ \ \ \ \ \ else\ semaphore\ {\isadigit{2}}\ \isanewline
\ \ \ \ \ \ \ \ \ \ \ \ \ \ \ \ \ \ \ \ \ \ \ \ \ \ \ \ \ \ \ \ \ \ \ \ {\isacharbar}{\kern0pt}{\isacharbar}{\kern0pt}{\isacharbar}{\kern0pt}\ {\isacharparenleft}{\kern0pt}if\ x\ {\isacharequal}{\kern0pt}\ evs{\isachardot}{\kern0pt}lock\ {\isadigit{3}}\isanewline
\ \ \ \ \ \ \ \ \ \ \ \ \ \ \ \ \ \ \ \ \ \ \ \ \ \ \ \ \ \ \ \ \ \ \ \ \ \ \ \ \ then\ evs{\isachardot}{\kern0pt}unlock\ {\isadigit{3}}\ {\isasymrightarrow}\ semaphore\ {\isadigit{3}}\ {\isacharbar}{\kern0pt}{\isacharbar}{\kern0pt}{\isacharbar}{\kern0pt}\ semaphore\ {\isadigit{4}}\isanewline
\ \ \ \ \ \ \ \ \ \ \ \ \ \ \ \ \ \ \ \ \ \ \ \ \ \ \ \ \ \ \ \ \ \ \ \ \ \ \ \ \ else\ semaphore\ {\isadigit{3}}\ {\isacharbar}{\kern0pt}{\isacharbar}{\kern0pt}{\isacharbar}{\kern0pt}\ evs{\isachardot}{\kern0pt}unlock\ {\isadigit{4}}\ {\isasymrightarrow}\ semaphore\ {\isadigit{4}}{\isacharparenright}{\kern0pt}{\isacharparenright}{\kern0pt}{\isacharparenright}{\kern0pt}{\isachardoublequoteclose}\isanewline
\ \ \ \ \ \ \ \ \ \ \ \isakeywordONE{apply}\isamarkupfalse%
{\isacharparenleft}{\kern0pt}simp\ add{\isacharcolon}{\kern0pt}\ List{\isachardot}{\kern0pt}upto{\isachardot}{\kern0pt}simps{\isacharparenright}{\kern0pt}\isanewline
\ \ \ \ \ \ \ \ \ \ \ \isakeywordONE{apply}\isamarkupfalse%
{\isacharparenleft}{\kern0pt}subst\ sema{\isacharunderscore}{\kern0pt}rec{\isacharbrackleft}{\kern0pt}\isakeywordTWO{where}\ n{\isacharequal}{\kern0pt}{\isacartoucheopen}{\isadigit{1}}{\isacartoucheclose}{\isacharbrackright}{\kern0pt}{\isacharcomma}{\kern0pt}subst\ sema{\isacharunderscore}{\kern0pt}rec{\isacharbrackleft}{\kern0pt}\isakeywordTWO{where}\ n{\isacharequal}{\kern0pt}{\isacartoucheopen}{\isadigit{2}}{\isacartoucheclose}{\isacharbrackright}{\kern0pt}{\isacharcomma}{\kern0pt}\isanewline
\ \ \ \ \ \ \ \ \ \ \ \ \ \ \ \ \ subst\ sema{\isacharunderscore}{\kern0pt}rec{\isacharbrackleft}{\kern0pt}\isakeywordTWO{where}\ n{\isacharequal}{\kern0pt}{\isacartoucheopen}{\isadigit{3}}{\isacartoucheclose}{\isacharbrackright}{\kern0pt}{\isacharcomma}{\kern0pt}subst\ sema{\isacharunderscore}{\kern0pt}rec{\isacharbrackleft}{\kern0pt}\isakeywordTWO{where}\ n{\isacharequal}{\kern0pt}{\isacartoucheopen}{\isadigit{4}}{\isacartoucheclose}{\isacharbrackright}{\kern0pt}{\isacharparenright}{\kern0pt}\ \isanewline
\ \ \ \ \ \ \ \ \ \ \ \isakeywordONE{apply}\isamarkupfalse%
{\isacharparenleft}{\kern0pt}simp\ add{\isacharcolon}{\kern0pt}\ inter{\isadigit{2}}Mprefix{\isadigit{1}}\ inter{\isadigit{2}}Mprefix{\isadigit{2}}{\isacharparenright}{\kern0pt}\isanewline
\ \ \ \ \ \ \ \ \ \ \ \isakeywordONE{by}\isamarkupfalse%
{\isacharparenleft}{\kern0pt}simp\ add{\isacharcolon}{\kern0pt}\ sema{\isacharunderscore}{\kern0pt}rec{\isacharbrackleft}{\kern0pt}\isakeywordTWO{where}\ n{\isacharequal}{\kern0pt}{\isacartoucheopen}{\isadigit{1}}{\isacartoucheclose}{\isacharcomma}{\kern0pt}symmetric{\isacharbrackright}{\kern0pt}\ \isanewline
\ \ \ \ \ \ \ \ \ \ \ \ \ \ \ \ \ \ \ \ \ \ \ \ sema{\isacharunderscore}{\kern0pt}rec{\isacharbrackleft}{\kern0pt}\isakeywordTWO{where}\ n{\isacharequal}{\kern0pt}{\isacartoucheopen}{\isadigit{2}}{\isacartoucheclose}{\isacharcomma}{\kern0pt}symmetric{\isacharbrackright}{\kern0pt}\ \isanewline
\ \ \ \ \ \ \ \ \ \ \ \ \ \ \ \ \ \ \ \ \ \ \ \ sema{\isacharunderscore}{\kern0pt}rec{\isacharbrackleft}{\kern0pt}\isakeywordTWO{where}\ n{\isacharequal}{\kern0pt}{\isacartoucheopen}{\isadigit{3}}{\isacartoucheclose}{\isacharcomma}{\kern0pt}symmetric{\isacharbrackright}{\kern0pt}\ \isanewline
\ \ \ \ \ \ \ \ \ \ \ \ \ \ \ \ \ \ \ \ \ \ \ \ sema{\isacharunderscore}{\kern0pt}rec{\isacharbrackleft}{\kern0pt}\isakeywordTWO{where}\ n{\isacharequal}{\kern0pt}{\isacartoucheopen}{\isadigit{4}}{\isacartoucheclose}{\isacharcomma}{\kern0pt}symmetric{\isacharbrackright}{\kern0pt}\ \isanewline
\ \ \ \ \ \ \ \ \ \ \ \ \ \ \ \ \ \ \ cong{\isacharcolon}{\kern0pt}\ HOL{\isachardot}{\kern0pt}if{\isacharunderscore}{\kern0pt}cong{\isacharparenright}{\kern0pt}\isanewline
\ \ \isakeywordONE{have}\isamarkupfalse%
\ {\isadigit{2}}{\isacharcolon}{\kern0pt}\ {\isachardoublequoteopen}evs{\isachardot}{\kern0pt}unlock\ {\isadigit{1}}\ {\isasymrightarrow}\ {\isacharquery}{\kern0pt}rhs\ {\isacharequal}{\kern0pt}\ {\isasymbox}x{\isasymin}{\isacharbraceleft}{\kern0pt}evs{\isachardot}{\kern0pt}unlock\ {\isadigit{1}}{\isacharbraceright}{\kern0pt}\ {\isasymrightarrow}\ {\isacharquery}{\kern0pt}rhs{\isachardoublequoteclose}\ \isakeywordONE{by}\isamarkupfalse%
\ {\isacharparenleft}{\kern0pt}metis\ write{\isadigit{0}}{\isacharunderscore}{\kern0pt}def{\isacharparenright}{\kern0pt}\isanewline
\ \ \isakeywordONE{have}\isamarkupfalse%
\ {\isadigit{3}}{\isacharcolon}{\kern0pt}\ {\isachardoublequoteopen}{\isacharquery}{\kern0pt}rhs\ {\isacharequal}{\kern0pt}\ {\isasymbox}x{\isasymin}{\isacharbraceleft}{\kern0pt}evs{\isachardot}{\kern0pt}lock\ {\isadigit{1}}{\isacharcomma}{\kern0pt}\ evs{\isachardot}{\kern0pt}lock\ {\isadigit{2}}{\isacharcomma}{\kern0pt}\ evs{\isachardot}{\kern0pt}lock\ {\isadigit{3}}{\isacharcomma}{\kern0pt}\ evs{\isachardot}{\kern0pt}lock\ {\isadigit{4}}{\isacharbraceright}{\kern0pt}\isanewline
\ \ \ \ \ \ \ \ \ \ \ \ \ \ \ \ \ \ \ \ {\isasymrightarrow}\ {\isacharparenleft}{\kern0pt}if\ x\ {\isacharequal}{\kern0pt}\ evs{\isachardot}{\kern0pt}lock\ {\isadigit{1}}\isanewline
\ \ \ \ \ \ \ \ \ \ \ \ \ \ \ \ \ \ \ \ \ \ \ \ then\ evs{\isachardot}{\kern0pt}unlock\ {\isadigit{1}}\ {\isasymrightarrow}\ semaphore\ \ {\isadigit{1}}\ \isanewline
\ \ \ \ \ \ \ \ \ \ \ \ \ \ \ \ \ \ \ \ \ \ \ \ \ \ \ \ \ {\isacharbar}{\kern0pt}{\isacharbar}{\kern0pt}{\isacharbar}{\kern0pt}\ {\isasymbox}x{\isasymin}{\isacharbraceleft}{\kern0pt}evs{\isachardot}{\kern0pt}lock\ {\isadigit{2}}{\isacharcomma}{\kern0pt}\ evs{\isachardot}{\kern0pt}lock\ {\isadigit{3}}{\isacharcomma}{\kern0pt}\ evs{\isachardot}{\kern0pt}lock\ {\isadigit{4}}{\isacharbraceright}{\kern0pt}\isanewline
\ \ \ \ \ \ \ \ \ \ \ \ \ \ \ \ \ \ \ \ \ \ \ \ \ \ \ \ \ \ \ \ \ {\isasymrightarrow}\ {\isacharparenleft}{\kern0pt}if\ x\ {\isacharequal}{\kern0pt}\ evs{\isachardot}{\kern0pt}lock\ {\isadigit{2}}\isanewline
\ \ \ \ \ \ \ \ \ \ \ \ \ \ \ \ \ \ \ \ \ \ \ \ \ \ \ \ \ \ \ \ \ \ \ \ \ then\ evs{\isachardot}{\kern0pt}unlock\ {\isadigit{2}}\ {\isasymrightarrow}\ semaphore\ {\isadigit{2}}\ \isanewline
\ \ \ \ \ \ \ \ \ \ \ \ \ \ \ \ \ \ \ \ \ \ \ \ \ \ \ \ \ \ \ \ \ \ \ \ \ \ \ \ \ \ {\isacharbar}{\kern0pt}{\isacharbar}{\kern0pt}{\isacharbar}{\kern0pt}\ {\isasymbox}x{\isasymin}{\isacharbraceleft}{\kern0pt}evs{\isachardot}{\kern0pt}lock\ {\isadigit{3}}{\isacharcomma}{\kern0pt}\ evs{\isachardot}{\kern0pt}lock\ {\isadigit{4}}{\isacharbraceright}{\kern0pt}\isanewline
\ \ \ \ \ \ \ \ \ \ \ \ \ \ \ \ \ \ \ \ \ \ \ \ \ \ \ \ \ \ \ \ \ \ \ \ \ \ \ \ \ \ \ \ \ \ {\isasymrightarrow}\ {\isacharparenleft}{\kern0pt}if\ x\ {\isacharequal}{\kern0pt}\ evs{\isachardot}{\kern0pt}lock\ {\isadigit{3}}\ \isanewline
\ \ \ \ \ \ \ \ \ \ \ \ \ \ \ \ \ \ \ \ \ \ \ \ \ \ \ \ \ \ \ \ \ \ \ \ \ \ \ \ \ \ \ \ \ \ \ \ \ \ then\ evs{\isachardot}{\kern0pt}unlock\ {\isadigit{3}}\ {\isasymrightarrow}\ semaphore\ {\isadigit{3}}\ {\isacharbar}{\kern0pt}{\isacharbar}{\kern0pt}{\isacharbar}{\kern0pt}\ semaphore\ {\isadigit{4}}\isanewline
\ \ \ \ \ \ \ \ \ \ \ \ \ \ \ \ \ \ \ \ \ \ \ \ \ \ \ \ \ \ \ \ \ \ \ \ \ \ \ \ \ \ \ \ \ \ \ \ \ \ else\ semaphore\ {\isadigit{3}}\ {\isacharbar}{\kern0pt}{\isacharbar}{\kern0pt}{\isacharbar}{\kern0pt}\ evs{\isachardot}{\kern0pt}unlock\ {\isadigit{4}}\ {\isasymrightarrow}\ semaphore\ {\isadigit{4}}{\isacharparenright}{\kern0pt}\isanewline
\ \ \ \ \ \ \ \ \ \ \ \ \ \ \ \ \ \ \ \ \ \ \ \ \ \ \ \ \ \ \ \ \ \ \ \ \ else\ semaphore\ {\isadigit{2}}\ \isanewline
\ \ \ \ \ \ \ \ \ \ \ \ \ \ \ \ \ \ \ \ \ \ \ \ \ \ \ \ \ \ \ \ \ \ \ \ \ \ \ \ \ \ {\isacharbar}{\kern0pt}{\isacharbar}{\kern0pt}{\isacharbar}{\kern0pt}\ {\isacharparenleft}{\kern0pt}if\ x\ {\isacharequal}{\kern0pt}\ evs{\isachardot}{\kern0pt}lock\ {\isadigit{3}}\ \isanewline
\ \ \ \ \ \ \ \ \ \ \ \ \ \ \ \ \ \ \ \ \ \ \ \ \ \ \ \ \ \ \ \ \ \ \ \ \ \ \ \ \ \ \ \ \ \ \ then\ evs{\isachardot}{\kern0pt}unlock\ {\isadigit{3}}\ {\isasymrightarrow}\ semaphore\ {\isadigit{3}}\ {\isacharbar}{\kern0pt}{\isacharbar}{\kern0pt}{\isacharbar}{\kern0pt}\ semaphore\ {\isadigit{4}}\isanewline
\ \ \ \ \ \ \ \ \ \ \ \ \ \ \ \ \ \ \ \ \ \ \ \ \ \ \ \ \ \ \ \ \ \ \ \ \ \ \ \ \ \ \ \ \ \ \ else\ semaphore\ {\isadigit{3}}\ {\isacharbar}{\kern0pt}{\isacharbar}{\kern0pt}{\isacharbar}{\kern0pt}\ evs{\isachardot}{\kern0pt}unlock\ {\isadigit{4}}\ {\isasymrightarrow}\ semaphore\ {\isadigit{4}}{\isacharparenright}{\kern0pt}{\isacharparenright}{\kern0pt}\isanewline
\ \ \ \ \ \ \ \ \ \ \ \ \ \ \ \ \ \ \ \ \ \ else\ semaphore\ {\isadigit{1}}\ \isanewline
\ \ \ \ \ \ \ \ \ \ \ \ \ \ \ \ \ \ \ \ \ \ \ \ \ \ \ {\isacharbar}{\kern0pt}{\isacharbar}{\kern0pt}{\isacharbar}{\kern0pt}\ {\isacharparenleft}{\kern0pt}if\ x\ {\isacharequal}{\kern0pt}\ evs{\isachardot}{\kern0pt}lock\ {\isadigit{2}}\isanewline
\ \ \ \ \ \ \ \ \ \ \ \ \ \ \ \ \ \ \ \ \ \ \ \ \ \ \ \ \ \ \ \ then\ evs{\isachardot}{\kern0pt}unlock\ {\isadigit{2}}\ {\isasymrightarrow}\ \ semaphore\ {\isadigit{2}}\ \isanewline
\ \ \ \ \ \ \ \ \ \ \ \ \ \ \ \ \ \ \ \ \ \ \ \ \ \ \ \ \ \ \ \ \ \ \ \ \ {\isacharbar}{\kern0pt}{\isacharbar}{\kern0pt}{\isacharbar}{\kern0pt}\ {\isasymbox}x{\isasymin}{\isacharbraceleft}{\kern0pt}evs{\isachardot}{\kern0pt}lock\ {\isadigit{3}}{\isacharcomma}{\kern0pt}\ evs{\isachardot}{\kern0pt}lock\ {\isadigit{4}}{\isacharbraceright}{\kern0pt}\isanewline
\ \ \ \ \ \ \ \ \ \ \ \ \ \ \ \ \ \ \ \ \ \ \ \ \ \ \ \ \ \ \ \ \ \ \ \ \ \ \ \ \ {\isasymrightarrow}\ {\isacharparenleft}{\kern0pt}if\ x\ {\isacharequal}{\kern0pt}\ evs{\isachardot}{\kern0pt}lock\ {\isadigit{3}}\ \isanewline
\ \ \ \ \ \ \ \ \ \ \ \ \ \ \ \ \ \ \ \ \ \ \ \ \ \ \ \ \ \ \ \ \ \ \ \ \ \ \ \ \ \ \ \ \ then\ evs{\isachardot}{\kern0pt}unlock\ {\isadigit{3}}\ {\isasymrightarrow}\ semaphore\ {\isadigit{3}}\ {\isacharbar}{\kern0pt}{\isacharbar}{\kern0pt}{\isacharbar}{\kern0pt}\ semaphore\ {\isadigit{4}}\isanewline
\ \ \ \ \ \ \ \ \ \ \ \ \ \ \ \ \ \ \ \ \ \ \ \ \ \ \ \ \ \ \ \ \ \ \ \ \ \ \ \ \ \ \ \ \ else\ semaphore\ {\isadigit{3}}\ {\isacharbar}{\kern0pt}{\isacharbar}{\kern0pt}{\isacharbar}{\kern0pt}\ evs{\isachardot}{\kern0pt}unlock\ {\isadigit{4}}\ {\isasymrightarrow}\ semaphore\ {\isadigit{4}}{\isacharparenright}{\kern0pt}\isanewline
\ \ \ \ \ \ \ \ \ \ \ \ \ \ \ \ \ \ \ \ \ \ \ \ \ \ \ \ \ \ \ \ else\ semaphore\ {\isadigit{2}}\ \isanewline
\ \ \ \ \ \ \ \ \ \ \ \ \ \ \ \ \ \ \ \ \ \ \ \ \ \ \ \ \ \ \ \ \ \ \ \ \ {\isacharbar}{\kern0pt}{\isacharbar}{\kern0pt}{\isacharbar}{\kern0pt}\ {\isacharparenleft}{\kern0pt}if\ x\ {\isacharequal}{\kern0pt}\ evs{\isachardot}{\kern0pt}lock\ {\isadigit{3}}\ \isanewline
\ \ \ \ \ \ \ \ \ \ \ \ \ \ \ \ \ \ \ \ \ \ \ \ \ \ \ \ \ \ \ \ \ \ \ \ \ \ \ \ \ \ then\ evs{\isachardot}{\kern0pt}unlock\ {\isadigit{3}}\ {\isasymrightarrow}\ semaphore\ {\isadigit{3}}\ {\isacharbar}{\kern0pt}{\isacharbar}{\kern0pt}{\isacharbar}{\kern0pt}\ semaphore\ {\isadigit{4}}\isanewline
\ \ \ \ \ \ \ \ \ \ \ \ \ \ \ \ \ \ \ \ \ \ \ \ \ \ \ \ \ \ \ \ \ \ \ \ \ \ \ \ \ \ else\ semaphore\ {\isadigit{3}}\ {\isacharbar}{\kern0pt}{\isacharbar}{\kern0pt}{\isacharbar}{\kern0pt}\ evs{\isachardot}{\kern0pt}unlock\ {\isadigit{4}}\ {\isasymrightarrow}\ semaphore\ {\isadigit{4}}{\isacharparenright}{\kern0pt}{\isacharparenright}{\kern0pt}{\isacharparenright}{\kern0pt}\ \isanewline
\ \ \ \ \ \ \ \ \ \ \ \ \ \ \ \ \ \ \ \ \ \ \ \ {\isacharbar}{\kern0pt}{\isacharbar}{\kern0pt}\ {\isacharparenleft}{\kern0pt}P{\isadigit{1}}\ {\isacharbar}{\kern0pt}{\isacharbar}{\kern0pt}{\isacharbar}{\kern0pt}\ P{\isadigit{2}}{\isacharparenright}{\kern0pt}{\isachardoublequoteclose}\isanewline
\ \ \ \ \ \ \ \ \ \isakeywordONE{apply}\isamarkupfalse%
{\isacharparenleft}{\kern0pt}simp\ add{\isacharcolon}{\kern0pt}\ List{\isachardot}{\kern0pt}upto{\isachardot}{\kern0pt}simps{\isacharparenright}{\kern0pt}\isanewline
\ \ \ \ \ \ \ \ \ \isakeywordONE{apply}\isamarkupfalse%
{\isacharparenleft}{\kern0pt}subst\ sema{\isacharunderscore}{\kern0pt}rec{\isacharbrackleft}{\kern0pt}\isakeywordTWO{where}\ n{\isacharequal}{\kern0pt}{\isacartoucheopen}{\isadigit{1}}{\isacartoucheclose}{\isacharbrackright}{\kern0pt}{\isacharcomma}{\kern0pt}subst\ sema{\isacharunderscore}{\kern0pt}rec{\isacharbrackleft}{\kern0pt}\isakeywordTWO{where}\ n{\isacharequal}{\kern0pt}{\isacartoucheopen}{\isadigit{2}}{\isacartoucheclose}{\isacharbrackright}{\kern0pt}{\isacharcomma}{\kern0pt}\isanewline
\ \ \ \ \ \ \ \ \ \ \ \ \ \ \ subst\ sema{\isacharunderscore}{\kern0pt}rec{\isacharbrackleft}{\kern0pt}\isakeywordTWO{where}\ n{\isacharequal}{\kern0pt}{\isacartoucheopen}{\isadigit{3}}{\isacartoucheclose}{\isacharbrackright}{\kern0pt}{\isacharcomma}{\kern0pt}subst\ sema{\isacharunderscore}{\kern0pt}rec{\isacharbrackleft}{\kern0pt}\isakeywordTWO{where}\ n{\isacharequal}{\kern0pt}{\isacartoucheopen}{\isadigit{4}}{\isacartoucheclose}{\isacharbrackright}{\kern0pt}{\isacharparenright}{\kern0pt}\ \isanewline
\ \ \ \ \ \ \ \ \ \isakeywordONE{apply}\isamarkupfalse%
{\isacharparenleft}{\kern0pt}simp\ add{\isacharcolon}{\kern0pt}\ inter{\isadigit{2}}Mprefix{\isadigit{1}}\ inter{\isadigit{2}}Mprefix{\isadigit{2}}{\isacharparenright}{\kern0pt}\isanewline
\ \ \ \ \ \ \ \ \ \isakeywordONE{by}\isamarkupfalse%
{\isacharparenleft}{\kern0pt}simp\ add{\isacharcolon}{\kern0pt}\ sema{\isacharunderscore}{\kern0pt}rec{\isacharbrackleft}{\kern0pt}\isakeywordTWO{where}\ n{\isacharequal}{\kern0pt}{\isacartoucheopen}{\isadigit{1}}{\isacartoucheclose}{\isacharcomma}{\kern0pt}symmetric{\isacharbrackright}{\kern0pt}\ \isanewline
\ \ \ \ \ \ \ \ \ \ \ \ \ \ \ \ \ \ \ \ \ \ sema{\isacharunderscore}{\kern0pt}rec{\isacharbrackleft}{\kern0pt}\isakeywordTWO{where}\ n{\isacharequal}{\kern0pt}{\isacartoucheopen}{\isadigit{2}}{\isacartoucheclose}{\isacharcomma}{\kern0pt}symmetric{\isacharbrackright}{\kern0pt}\ \isanewline
\ \ \ \ \ \ \ \ \ \ \ \ \ \ \ \ \ \ \ \ \ \ sema{\isacharunderscore}{\kern0pt}rec{\isacharbrackleft}{\kern0pt}\isakeywordTWO{where}\ n{\isacharequal}{\kern0pt}{\isacartoucheopen}{\isadigit{3}}{\isacartoucheclose}{\isacharcomma}{\kern0pt}symmetric{\isacharbrackright}{\kern0pt}\ \isanewline
\ \ \ \ \ \ \ \ \ \ \ \ \ \ \ \ \ \ \ \ \ \ sema{\isacharunderscore}{\kern0pt}rec{\isacharbrackleft}{\kern0pt}\isakeywordTWO{where}\ n{\isacharequal}{\kern0pt}{\isacartoucheopen}{\isadigit{4}}{\isacartoucheclose}{\isacharcomma}{\kern0pt}symmetric{\isacharbrackright}{\kern0pt}\ \isanewline
\ \ \ \ \ \ \ \ \ \ \ \ \ \ \ \ \ cong{\isacharcolon}{\kern0pt}\ HOL{\isachardot}{\kern0pt}if{\isacharunderscore}{\kern0pt}cong{\isacharparenright}{\kern0pt}\isanewline
\ \ \isakeywordTHREE{show}\isamarkupfalse%
\ {\isacharquery}{\kern0pt}thesis\isanewline
\ \ \ \ \isakeywordONE{apply}\isamarkupfalse%
{\isacharparenleft}{\kern0pt}subst\ {\isadigit{1}}{\isacharcomma}{\kern0pt}\ subst\ {\isadigit{2}}{\isacharparenright}{\kern0pt}\isanewline
\ \ \ \ \isakeywordONE{apply}\isamarkupfalse%
{\isacharparenleft}{\kern0pt}subst\ Mprefix{\isacharunderscore}{\kern0pt}Par{\isacharunderscore}{\kern0pt}Mprefix{\isacharcomma}{\kern0pt}\ simp{\isacharunderscore}{\kern0pt}all{\isacharparenright}{\kern0pt}\isanewline
\ \ \ \ \isakeywordONE{apply}\isamarkupfalse%
{\isacharparenleft}{\kern0pt}simp{\isacharunderscore}{\kern0pt}all\ add{\isacharcolon}{\kern0pt}\ Mprefix{\isacharunderscore}{\kern0pt}singl\ \ cong{\isacharcolon}{\kern0pt}\ HOL{\isachardot}{\kern0pt}if{\isacharunderscore}{\kern0pt}cong{\isacharparenright}{\kern0pt}\isanewline
\ \ \ \ \isakeywordONE{apply}\isamarkupfalse%
{\isacharparenleft}{\kern0pt}subst\ sema{\isacharunderscore}{\kern0pt}rec{\isacharbrackleft}{\kern0pt}\isakeywordTWO{where}\ n{\isacharequal}{\kern0pt}{\isacartoucheopen}{\isadigit{1}}{\isacartoucheclose}{\isacharbrackright}{\kern0pt}{\isacharparenright}{\kern0pt}\isanewline
\ \ \ \ \isakeywordONE{apply}\isamarkupfalse%
{\isacharparenleft}{\kern0pt}subst\ inter{\isadigit{2}}Mprefix{\isadigit{2}}{\isacharcomma}{\kern0pt}\ simp{\isacharparenright}{\kern0pt}\isanewline
\ \ \ \ \isakeywordONE{apply}\isamarkupfalse%
{\isacharparenleft}{\kern0pt}subst\ inter{\isadigit{2}}Mprefix{\isadigit{2}}{\isacharcomma}{\kern0pt}\ simp{\isacharparenright}{\kern0pt}\isanewline
\ \ \ \ \isakeywordONE{apply}\isamarkupfalse%
{\isacharparenleft}{\kern0pt}subst\ assms{\isacharparenright}{\kern0pt}\isanewline
\ \ \ \ \isakeywordONE{apply}\isamarkupfalse%
{\isacharparenleft}{\kern0pt}subst\ sema{\isacharunderscore}{\kern0pt}rec{\isacharbrackleft}{\kern0pt}\isakeywordTWO{where}\ n{\isacharequal}{\kern0pt}{\isacartoucheopen}{\isadigit{1}}{\isacartoucheclose}{\isacharcomma}{\kern0pt}\ symmetric{\isacharbrackright}{\kern0pt}{\isacharparenright}{\kern0pt}\isanewline
\ \ \ \ \isakeywordONE{apply}\isamarkupfalse%
{\isacharparenleft}{\kern0pt}subst\ {\isadigit{3}}{\isacharparenright}{\kern0pt}\isanewline
\ \ \ \ \isakeywordONE{by}\isamarkupfalse%
\ {\isacharparenleft}{\kern0pt}metis\ {\isacharparenleft}{\kern0pt}no{\isacharunderscore}{\kern0pt}types{\isacharcomma}{\kern0pt}lifting{\isacharparenright}{\kern0pt}\ empty{\isacharunderscore}{\kern0pt}iff\ evs{\isachardot}{\kern0pt}distinct{\isacharparenleft}{\kern0pt}{\isadigit{1}}{\isacharparenright}{\kern0pt}\ insert{\isacharunderscore}{\kern0pt}iff\ \isanewline
\ \ \ \ \ \ \ \ \ \ \ \ \ \ inter{\isadigit{2}}Mprefix{\isadigit{2}}\ mono{\isacharunderscore}{\kern0pt}Mprefix{\isacharunderscore}{\kern0pt}eq{\isacharparenright}{\kern0pt}\isanewline
\ \ \isakeywordONE{qed}\isamarkupfalse%
%
\endisatagproof
{\isafoldproof}%
%
\isadelimproof
%
\endisadelimproof
%
\begin{isamarkuptext}%
And now we attack the main theorem:%
\end{isamarkuptext}\isamarkuptrue%
\isakeywordONE{theorem}\isamarkupfalse%
\ bullocks{\isadigit{2}}\ {\isacharcolon}{\kern0pt}\ \ {\isacartoucheopen}deadlock{\isacharunderscore}{\kern0pt}free\ Thread{\isadigit{6}}{\isacharunderscore}{\kern0pt}sem{\isacartoucheclose}\ \ \isanewline
%
\isadelimproof
\ \ %
\endisadelimproof
%
\isatagproof
\isakeywordONE{unfolding}\isamarkupfalse%
\ Thread{\isadigit{6}}{\isacharunderscore}{\kern0pt}sem{\isacharunderscore}{\kern0pt}def\ Sem{\isacharunderscore}{\kern0pt}def\ Thread{\isadigit{5}}{\isacharunderscore}{\kern0pt}def\ Thread{\isadigit{6}}{\isacharunderscore}{\kern0pt}def\isanewline
\isakeywordONE{proof}\isamarkupfalse%
\ {\isacharminus}{\kern0pt}\ \isanewline
\ \ \isakeywordTHREE{show}\isamarkupfalse%
\ {\isacartoucheopen}deadlock{\isacharunderscore}{\kern0pt}free\isanewline
\ \ \ \ \ \ \ \ \ \ \ \ \ \ \ {\isacharparenleft}{\kern0pt}\ MultiInter\ {\isacharparenleft}{\kern0pt}mset\ {\isacharbrackleft}{\kern0pt}{\isadigit{1}}{\isachardot}{\kern0pt}{\isachardot}{\kern0pt}{\isadigit{4}}{\isacharbrackright}{\kern0pt}{\isacharparenright}{\kern0pt}\ semaphore\isanewline
\ \ \ \ \ \ \ \ \ \ \ \ \ \ \ \ {\isacharbar}{\kern0pt}{\isacharbar}{\kern0pt}\ {\isacharparenleft}{\kern0pt}\ \ \ \ Sem\isactrlsub {\isadigit{0}}\ {\isacharparenleft}{\kern0pt}WHILE\ {\isasymlambda}{\isasymsigma}{\isachardot}{\kern0pt}\ True\ DO\ com{\isachardot}{\kern0pt}lock\ {\isadigit{1}}{\isacharparenright}{\kern0pt}\ {\isacharparenleft}{\kern0pt}{\isasymlambda}{\isacharunderscore}{\kern0pt}{\isachardot}{\kern0pt}\ Skip{\isacharparenright}{\kern0pt}\ {\isasymsigma}\isactrlsub {\isadigit{0}}\ \isanewline
\ \ \ \ \ \ \ \ \ \ \ \ \ \ \ \ \ \ \ \ {\isacharbar}{\kern0pt}{\isacharbar}{\kern0pt}{\isacharbar}{\kern0pt}\ Sem\isactrlsub {\isadigit{0}}\ {\isacharparenleft}{\kern0pt}WHILE\ {\isasymlambda}{\isasymsigma}{\isachardot}{\kern0pt}\ True\ DO\ com{\isachardot}{\kern0pt}unlock\ {\isadigit{1}}{\isacharparenright}{\kern0pt}\ {\isacharparenleft}{\kern0pt}{\isasymlambda}{\isacharunderscore}{\kern0pt}{\isachardot}{\kern0pt}\ Skip{\isacharparenright}{\kern0pt}\ {\isasymsigma}\isactrlsub {\isadigit{0}}{\isacharparenright}{\kern0pt}{\isacharparenright}{\kern0pt}{\isacartoucheclose}\isanewline
\ \ \ \ \ \ \ \ \ \ \ \ \ \ \ \ {\isacharparenleft}{\kern0pt}\isakeywordTWO{is}\ {\isacartoucheopen}deadlock{\isacharunderscore}{\kern0pt}free\ {\isacharparenleft}{\kern0pt}{\isacharquery}{\kern0pt}SEMA\ {\isacharbar}{\kern0pt}{\isacharbar}{\kern0pt}\ {\isacharparenleft}{\kern0pt}{\isacharquery}{\kern0pt}P{\isadigit{1}}\ \ {\isacharbar}{\kern0pt}{\isacharbar}{\kern0pt}{\isacharbar}{\kern0pt}\ {\isacharquery}{\kern0pt}P{\isadigit{2}}{\isacharparenright}{\kern0pt}{\isacharparenright}{\kern0pt}{\isacartoucheclose}{\isacharparenright}{\kern0pt}\isanewline
\ \ \isakeywordONE{proof}\isamarkupfalse%
\ {\isacharminus}{\kern0pt}\isanewline
\ \ \ \ \isakeywordONE{have}\isamarkupfalse%
\ {\isadigit{1}}\ {\isacharcolon}{\kern0pt}\ {\isacartoucheopen}{\isacharquery}{\kern0pt}P{\isadigit{1}}\ {\isacharequal}{\kern0pt}\ evs{\isachardot}{\kern0pt}lock\ {\isadigit{1}}\ {\isasymrightarrow}\ {\isacharquery}{\kern0pt}P{\isadigit{1}}{\isacartoucheclose}\ \isakeywordONE{using}\isamarkupfalse%
\ Thread{\isadigit{5}}{\isacharunderscore}{\kern0pt}while{\isacharunderscore}{\kern0pt}rec\ \isakeywordONE{by}\isamarkupfalse%
\ auto\isanewline
\ \ \ \ \isakeywordONE{have}\isamarkupfalse%
\ {\isadigit{2}}\ {\isacharcolon}{\kern0pt}\ {\isacartoucheopen}{\isacharquery}{\kern0pt}P{\isadigit{2}}\ {\isacharequal}{\kern0pt}\ evs{\isachardot}{\kern0pt}unlock\ {\isadigit{1}}\ {\isasymrightarrow}\ {\isacharquery}{\kern0pt}P{\isadigit{2}}{\isacartoucheclose}\ \isakeywordONE{using}\isamarkupfalse%
\ Thread{\isadigit{6}}{\isacharunderscore}{\kern0pt}while{\isacharunderscore}{\kern0pt}rec\ \isakeywordONE{by}\isamarkupfalse%
\ auto\isanewline
\ \ \ \ \isakeywordONE{have}\isamarkupfalse%
\ {\isadigit{3}}\ {\isacharcolon}{\kern0pt}\ {\isacartoucheopen}{\isacharquery}{\kern0pt}P{\isadigit{1}}\ \ {\isacharbar}{\kern0pt}{\isacharbar}{\kern0pt}{\isacharbar}{\kern0pt}\ {\isacharquery}{\kern0pt}P{\isadigit{2}}\ {\isacharequal}{\kern0pt}\ {\isasymbox}x{\isasymin}{\isacharbraceleft}{\kern0pt}evs{\isachardot}{\kern0pt}lock\ {\isadigit{1}}{\isacharcomma}{\kern0pt}\ evs{\isachardot}{\kern0pt}unlock\ {\isadigit{1}}{\isacharbraceright}{\kern0pt}\isanewline
\ \ \ \ \ \ \ \ \ \ \ \ \ \ \ \ \ \ \ \ \ \ \ \ \ \ \ \ \ {\isasymrightarrow}\ {\isacharparenleft}{\kern0pt}if\ x\ {\isacharequal}{\kern0pt}\ evs{\isachardot}{\kern0pt}lock\ {\isadigit{1}}\ then\ {\isacharquery}{\kern0pt}P{\isadigit{1}}\ {\isacharbar}{\kern0pt}{\isacharbar}{\kern0pt}{\isacharbar}{\kern0pt}\ evs{\isachardot}{\kern0pt}unlock\ {\isadigit{1}}\ {\isasymrightarrow}\ {\isacharquery}{\kern0pt}P{\isadigit{2}}\isanewline
\ \ \ \ \ \ \ \ \ \ \ \ \ \ \ \ \ \ \ \ \ \ \ \ \ \ \ \ \ \ \ \ \ \ \ \ \ \ \ \ \ \ \ \ \ \ \ \ \ \ \ else\ evs{\isachardot}{\kern0pt}lock\ {\isadigit{1}}\ {\isasymrightarrow}\ {\isacharquery}{\kern0pt}P{\isadigit{1}}\ {\isacharbar}{\kern0pt}{\isacharbar}{\kern0pt}{\isacharbar}{\kern0pt}\ {\isacharquery}{\kern0pt}P{\isadigit{2}}{\isacharparenright}{\kern0pt}{\isacartoucheclose}\isanewline
\ \ \ \ \ \ \ \ \ \ \ \ \ \isakeywordONE{apply}\isamarkupfalse%
{\isacharparenleft}{\kern0pt}subst\ {\isadigit{1}}{\isacharcomma}{\kern0pt}\ subst\ {\isadigit{2}}{\isacharparenright}{\kern0pt}\ \isakeywordONE{by}\isamarkupfalse%
{\isacharparenleft}{\kern0pt}subst\ inter{\isadigit{2}}Mprefix{\isadigit{1}}{\isacharcomma}{\kern0pt}simp{\isacharunderscore}{\kern0pt}all{\isacharparenright}{\kern0pt}\isanewline
\ \ \ \ \isakeywordONE{have}\isamarkupfalse%
\ {\isadigit{4}}\ {\isacharcolon}{\kern0pt}\ {\isacartoucheopen}{\isacharparenleft}{\kern0pt}{\isacharquery}{\kern0pt}SEMA\ {\isacharbar}{\kern0pt}{\isacharbar}{\kern0pt}\ {\isachardot}{\kern0pt}{\isachardot}{\kern0pt}{\isachardot}{\kern0pt}{\isacharparenright}{\kern0pt}\ \isanewline
\ \ \ \ \ \ \ \ \ \ \ \ \ \ {\isacharequal}{\kern0pt}\ evs{\isachardot}{\kern0pt}lock\ {\isadigit{1}}\ {\isasymrightarrow}\ {\isacharparenleft}{\kern0pt}{\isacharparenleft}{\kern0pt}evs{\isachardot}{\kern0pt}unlock\ {\isadigit{1}}\ {\isasymrightarrow}\ semaphore\ {\isadigit{1}}\ \ {\isacharbar}{\kern0pt}{\isacharbar}{\kern0pt}{\isacharbar}{\kern0pt}\ MultiInter\ {\isacharparenleft}{\kern0pt}mset\ {\isacharbrackleft}{\kern0pt}{\isadigit{2}}{\isachardot}{\kern0pt}{\isachardot}{\kern0pt}{\isadigit{4}}{\isacharbrackright}{\kern0pt}{\isacharparenright}{\kern0pt}\ semaphore{\isacharparenright}{\kern0pt}\ \isanewline
\ \ \ \ \ \ \ \ \ \ \ \ \ \ \ \ {\isacharbar}{\kern0pt}{\isacharbar}{\kern0pt}\ {\isacharparenleft}{\kern0pt}{\isacharquery}{\kern0pt}P{\isadigit{1}}\ {\isacharbar}{\kern0pt}{\isacharbar}{\kern0pt}{\isacharbar}{\kern0pt}\ {\isacharquery}{\kern0pt}P{\isadigit{2}}{\isacharparenright}{\kern0pt}{\isacharparenright}{\kern0pt}{\isacartoucheclose}\isanewline
\ \ \ \ \ \ \ \ \ \ \ \ \ {\isacharparenleft}{\kern0pt}\isakeywordTWO{is}\ {\isacartoucheopen}{\isacharunderscore}{\kern0pt}\ {\isacharequal}{\kern0pt}\ evs{\isachardot}{\kern0pt}lock\ {\isadigit{1}}\ {\isasymrightarrow}\ {\isacharquery}{\kern0pt}P{\isadigit{3}}{\isacartoucheclose}{\isacharparenright}{\kern0pt}\isanewline
\ \ \ \ \ \ \ \ \ \ \ \ \ \isakeywordONE{using}\isamarkupfalse%
\ {\isadigit{2}}\ contractMultiInter{\isadigit{1}}{\isacharprime}{\kern0pt}\ \isakeywordONE{by}\isamarkupfalse%
\ auto\isanewline
\ \ \ \ \isakeywordONE{have}\isamarkupfalse%
\ {\isadigit{5}}\ {\isacharcolon}{\kern0pt}\ {\isacartoucheopen}\ {\isacharparenleft}{\kern0pt}evs{\isachardot}{\kern0pt}unlock\ {\isadigit{1}}\ {\isasymrightarrow}\ semaphore\ {\isadigit{1}}\ \ {\isacharbar}{\kern0pt}{\isacharbar}{\kern0pt}{\isacharbar}{\kern0pt}\ MultiInter\ {\isacharparenleft}{\kern0pt}mset\ {\isacharbrackleft}{\kern0pt}{\isadigit{2}}{\isachardot}{\kern0pt}{\isachardot}{\kern0pt}{\isadigit{4}}{\isacharbrackright}{\kern0pt}{\isacharparenright}{\kern0pt}\ semaphore{\isacharparenright}{\kern0pt}\ {\isacharbar}{\kern0pt}{\isacharbar}{\kern0pt}\ {\isacharparenleft}{\kern0pt}{\isacharquery}{\kern0pt}P{\isadigit{1}}{\isacharbar}{\kern0pt}{\isacharbar}{\kern0pt}{\isacharbar}{\kern0pt}{\isacharquery}{\kern0pt}P{\isadigit{2}}{\isacharparenright}{\kern0pt}\ \isanewline
\ \ \ \ \ \ \ \ \ \ \ \ \ \ {\isacharequal}{\kern0pt}\ \isanewline
\ \ \ \ \ \ \ \ \ \ \ \ \ \ \ {\isacharparenleft}{\kern0pt}evs{\isachardot}{\kern0pt}unlock\ {\isadigit{1}}\ {\isasymrightarrow}\ {\isacharparenleft}{\kern0pt}{\isacharparenleft}{\kern0pt}MultiInter\ {\isacharparenleft}{\kern0pt}mset\ {\isacharbrackleft}{\kern0pt}{\isadigit{1}}{\isachardot}{\kern0pt}{\isachardot}{\kern0pt}{\isadigit{4}}{\isacharbrackright}{\kern0pt}{\isacharparenright}{\kern0pt}\ semaphore{\isacharparenright}{\kern0pt}\ {\isacharbar}{\kern0pt}{\isacharbar}{\kern0pt}\ {\isacharparenleft}{\kern0pt}{\isacharquery}{\kern0pt}P{\isadigit{1}}{\isacharbar}{\kern0pt}{\isacharbar}{\kern0pt}{\isacharbar}{\kern0pt}{\isacharquery}{\kern0pt}P{\isadigit{2}}{\isacharparenright}{\kern0pt}{\isacharparenright}{\kern0pt}{\isacharparenright}{\kern0pt}{\isacartoucheclose}\isanewline
\ \ \ \ \ \ \ \ \ \ \ \ \ \isakeywordONE{using}\isamarkupfalse%
\ {\isadigit{1}}\ {\isadigit{3}}\ contractMultiInter{\isadigit{1}}{\isacharprime}{\kern0pt}{\isacharprime}{\kern0pt}\ \isakeywordONE{by}\isamarkupfalse%
\ fastforce\isanewline
\ \ \ \ \isakeywordTHREE{show}\isamarkupfalse%
\ {\isacharquery}{\kern0pt}thesis\isanewline
\ \ \ \ \ \ \isakeywordONE{apply}\isamarkupfalse%
{\isacharparenleft}{\kern0pt}rule\ deadlock{\isacharunderscore}{\kern0pt}free{\isacharunderscore}{\kern0pt}{\isadigit{2}}{\isacharunderscore}{\kern0pt}coinduct{\isacharparenright}{\kern0pt}\ \isanewline
\ \ \ \ \ \ \isakeywordONE{apply}\isamarkupfalse%
{\isacharparenleft}{\kern0pt}subst\ {\isadigit{3}}{\isacharcomma}{\kern0pt}\ subst\ {\isadigit{4}}{\isacharparenright}{\kern0pt}\isanewline
\ \ \ \ \ \ \isakeywordONE{apply}\isamarkupfalse%
{\isacharparenleft}{\kern0pt}rule\ Mndetprefix{\isacharunderscore}{\kern0pt}FD{\isacharunderscore}{\kern0pt}write{\isadigit{0}}{\isacharbrackleft}{\kern0pt}of\ {\isacartoucheopen}evs{\isachardot}{\kern0pt}lock\ {\isadigit{1}}{\isacartoucheclose}{\isacharcomma}{\kern0pt}\ THEN\ trans{\isacharunderscore}{\kern0pt}FD{\isacharbrackright}{\kern0pt}{\isacharcomma}{\kern0pt}simp{\isacharparenright}{\kern0pt}\isanewline
\ \ \ \ \ \ \isakeywordONE{apply}\isamarkupfalse%
{\isacharparenleft}{\kern0pt}rule\ mono{\isacharunderscore}{\kern0pt}write{\isadigit{0}}{\isacharunderscore}{\kern0pt}FD{\isacharparenright}{\kern0pt}\isanewline
\ \ \ \ \ \ \isakeywordONE{apply}\isamarkupfalse%
{\isacharparenleft}{\kern0pt}subst\ {\isadigit{5}}{\isacharparenright}{\kern0pt}\ \isanewline
\ \ \ \ \ \ \isakeywordONE{by}\isamarkupfalse%
\ {\isacharparenleft}{\kern0pt}metis\ Mndetprefix{\isacharunderscore}{\kern0pt}FD\ UNIV{\isacharunderscore}{\kern0pt}I\ mono{\isacharunderscore}{\kern0pt}write{\isadigit{0}}{\isacharunderscore}{\kern0pt}FD{\isacharparenright}{\kern0pt}\isanewline
\ \ \isakeywordONE{qed}\isamarkupfalse%
\isanewline
\isakeywordONE{qed}\isamarkupfalse%
%
\endisatagproof
{\isafoldproof}%
%
\isadelimproof
%
\endisadelimproof
%
\isadelimdocument
%
\endisadelimdocument
%
\isatagdocument
%
\isamarkupsection{A more serious Example: An MPI Program%
}
\isamarkuptrue%
%
\endisatagdocument
{\isafolddocument}%
%
\isadelimdocument
%
\endisadelimdocument
%
\begin{isamarkuptext}%
The example is drawn from Steven F. Siegels paper:
``Parameterized Verification of Deterministic MPI Programs''
\url{https://arxiv.org/pdf/2607.18049}, 2026. It addresses 
the verification of message passing programs (MPI) 
\footnote{Message-Passing Interface Forum. 2025. MPI: A Message-Passing
Interface Standard, Version 5.0. 
\url{https://www.mpi-forum.org/docs/mpi-5.0/mpi50-report.pdf}.}

The informal spec of the running example reads as follows:
The program \isatt{c{\kern0pt}y{\kern0pt}c{\kern0pt}s{\kern0pt}u{\kern0pt}m{\kern0pt}} is a cyclic sum reduction to all processes. Each
process stores a int number in \isa{x}, then repeatedly sends to
its left, receives from its right, and adds the received value to
sum. At the end, each process holds the sum of the original
array \isa{A} of values.

A bit more formally, this reads as follows:
%
\begin{isabelle}%
\ \ \ \ \ \ \ \ \ \ {\isadigit{1}}\ \ param\ A{\isacharsemicolon}{\kern0pt}\isanewline
\ \ \ \ \ \ \ \ \ \ {\isadigit{2}}\ \ var\ i{\isacharcomma}{\kern0pt}\ x{\isacharcomma}{\kern0pt}\ sum{\isacharcomma}{\kern0pt}\ left{\isacharcomma}{\kern0pt}\ right{\isacharsemicolon}{\kern0pt}\isanewline
\ \ \ \ \ \ \ \ \ \ {\isadigit{3}}\ \ assume\ isRealArray{\isacharparenleft}{\kern0pt}A{\isacharcomma}{\kern0pt}NP{\isacharparenright}{\kern0pt}\ {\isasymand}\ x\ {\isacharequal}{\kern0pt}\ A{\isacharbrackleft}{\kern0pt}pid{\isacharbrackright}{\kern0pt}{\isacharsemicolon}{\kern0pt}\isanewline
\ \ \ \ \ \ \ \ \ \ {\isadigit{4}}\ \ left\ {\isacharcolon}{\kern0pt}{\isacharequal}{\kern0pt}\ {\isacharparenleft}{\kern0pt}pid\ {\isacharplus}{\kern0pt}\ NP−{\isadigit{1}}{\isacharparenright}{\kern0pt}{\isacharpercent}{\kern0pt}NP{\isacharsemicolon}{\kern0pt}\isanewline
\ \ \ \ \ \ \ \ \ \ {\isadigit{5}}\ \ right\ {\isacharcolon}{\kern0pt}{\isacharequal}{\kern0pt}\ {\isacharparenleft}{\kern0pt}pid\ {\isacharplus}{\kern0pt}\ {\isadigit{1}}{\isacharparenright}{\kern0pt}{\isacharpercent}{\kern0pt}NP{\isacharsemicolon}{\kern0pt}\isanewline
\ \ \ \ \ \ \ \ \ \ {\isadigit{6}}\ \ i\ {\isacharcolon}{\kern0pt}{\isacharequal}{\kern0pt}\ {\isadigit{1}}{\isacharsemicolon}{\kern0pt}\isanewline
\ \ \ \ \ \ \ \ \ \ {\isadigit{7}}\ \ sum\ {\isacharcolon}{\kern0pt}{\isacharequal}{\kern0pt}\ x{\isacharsemicolon}{\kern0pt}\isanewline
\ \ \ \ \ \ \ \ \ \ {\isadigit{8}}\ \ while\ {\isacharparenleft}{\kern0pt}\ i\ {\isacharless}{\kern0pt}\ NP\ {\isacharparenright}{\kern0pt}\ {\isacharbraceleft}{\kern0pt}\isanewline
\ \ \ \ \ \ \ \ \ \ {\isadigit{9}}\ \ \ \ \ send\ x\ to\ left{\isacharsemicolon}{\kern0pt}\isanewline
\ \ \ \ \ \ \ \ \ \ {\isadigit{1}}{\isadigit{0}}\ \ \ \ recv\ x\ from\ right{\isacharsemicolon}{\kern0pt}\isanewline
\ \ \ \ \ \ \ \ \ \ {\isadigit{1}}{\isadigit{1}}\ \ \ \ sum\ {\isacharcolon}{\kern0pt}{\isacharequal}{\kern0pt}\ sum\ {\isacharplus}{\kern0pt}\ x{\isacharsemicolon}{\kern0pt}\isanewline
\ \ \ \ \ \ \ \ \ \ {\isadigit{1}}{\isadigit{2}}\ \ \ \ x\ {\isacharcolon}{\kern0pt}{\isacharequal}{\kern0pt}\ x{\isacharplus}{\kern0pt}{\isadigit{1}}{\isacharsemicolon}{\kern0pt}\isanewline
\ \ \ \ \ \ \ \ \ \ {\isadigit{1}}{\isadigit{3}}\ {\isacharbraceright}{\kern0pt}\isanewline
\ \ \ \ \ \ \ \ \ \ {\isadigit{1}}{\isadigit{4}}\ assert\ sum\ {\isacharequal}{\kern0pt}\ {\isasymSigma}\isactrlsup N\isactrlsup P\isactrlsup −\isactrlsup {\isadigit{1}}\isactrlsub i\isactrlsub {\isacharequal}{\kern0pt}\isactrlsub {\isadigit{0}}\ A{\isacharbrackleft}{\kern0pt}i{\isacharbrackright}{\kern0pt}{\isacharsemicolon}{\kern0pt}%
\end{isabelle}

Note that this is a conceptual notation. The full fledged proof in Frama-C
involves several transformations and an unverfied VCG in the tool-chain; however, it
offers automated proofs of the  program code annotated by assertions, which 
comprises 17 pages of dense Frama-C code. 

Our modeling in \isa{IMP\isactrlsub c\isactrlsub o\isactrlsub n\isactrlsub c\isactrlsub u\isactrlsub r} reads as follows:%
\end{isamarkuptext}\isamarkuptrue%
\isakeywordONE{consts}\isamarkupfalse%
\ NP{\isacharcolon}{\kern0pt}{\isacharcolon}{\kern0pt}\ int\ %
\isamarkupcmt{Parameter of spec as unspecified constant%
}\isanewline
\isanewline
\isakeywordONE{consts}\isamarkupfalse%
\ array\ {\isacharcolon}{\kern0pt}{\isacharcolon}{\kern0pt}\ {\isacartoucheopen}string\ {\isasymtimes}\ int\ {\isasymRightarrow}\ string{\isacartoucheclose}\ \isanewline
\ \ \ \ \ \ \ \ \ \ \ \ \ \ \ \ %
\isamarkupcmt{array-name + offset :
                    the only thing we need to know: array is injective.%
}\isanewline
\isakeywordONE{consts}\isamarkupfalse%
\ array{\isacharunderscore}{\kern0pt}offset\ {\isacharcolon}{\kern0pt}{\isacharcolon}{\kern0pt}\ {\isacartoucheopen}string\ {\isasymtimes}\ int\ {\isasymRightarrow}\ int{\isacartoucheclose}\isanewline
\ \ \ \ \ \ \ \ \ \ \ \ \ \ \ \ %
\isamarkupcmt{the only thing we need to know: \isa{array{\isacharunderscore}{\kern0pt}offset{\isacharparenleft}{\kern0pt}array{\isacharparenleft}{\kern0pt}X{\isacharcomma}{\kern0pt}i{\isacharparenright}{\kern0pt}{\isacharparenright}{\kern0pt}\ {\isacharequal}{\kern0pt}\ i}.%
}\isanewline
\isanewline
\isakeywordONE{term}\isamarkupfalse%
{\isacartoucheopen}{\isacharbrackleft}{\kern0pt}{\isadigit{1}}{\isachardot}{\kern0pt}{\isachardot}{\kern0pt}N{\isacharbrackright}{\kern0pt}{\isacartoucheclose}\isanewline
\isakeywordONE{definition}\isamarkupfalse%
\ sums\ \isakeywordTWO{where}\ \ {\isacartoucheopen}sums\ A\ {\isacharequal}{\kern0pt}\ fold\ {\isacharparenleft}{\kern0pt}{\isacharplus}{\kern0pt}{\isacharparenright}{\kern0pt}\ A\ {\isacharparenleft}{\kern0pt}{\isadigit{0}}{\isacharcolon}{\kern0pt}{\isacharcolon}{\kern0pt}int{\isacharparenright}{\kern0pt}\ {\isacartoucheclose}\isanewline
\isanewline
\isakeywordONE{definition}\isamarkupfalse%
\ Thread{\isadigit{7}}\ \isakeywordTWO{where}\isanewline
{\isacartoucheopen}Thread{\isadigit{7}}\ pid\ A\ {\isacharequal}{\kern0pt}\ {\isacharparenleft}{\kern0pt}\ {\isacharprime}{\kern0pt}{\isacharprime}{\kern0pt}left{\isacharprime}{\kern0pt}{\isacharprime}{\kern0pt}\ \ {\isacharcolon}{\kern0pt}{\isacharequal}{\kern0pt}\ {\isacharparenleft}{\kern0pt}{\isasymlambda}{\isasymsigma}{\isachardot}{\kern0pt}\ {\isacharparenleft}{\kern0pt}pid\ {\isacharplus}{\kern0pt}\ NP\ {\isacharminus}{\kern0pt}\ {\isadigit{1}}{\isacharparenright}{\kern0pt}\ mod\ NP{\isacharparenright}{\kern0pt}{\isacharparenright}{\kern0pt}\ {\isacharsemicolon}{\kern0pt}\isanewline
\ \ \ \ \ \ \ \ \ \ \ \ \ \ \ \ \ {\isacharparenleft}{\kern0pt}\ {\isacharprime}{\kern0pt}{\isacharprime}{\kern0pt}right{\isacharprime}{\kern0pt}{\isacharprime}{\kern0pt}\ {\isacharcolon}{\kern0pt}{\isacharequal}{\kern0pt}\ {\isacharparenleft}{\kern0pt}{\isasymlambda}{\isasymsigma}{\isachardot}{\kern0pt}\ {\isacharparenleft}{\kern0pt}pid\ {\isacharplus}{\kern0pt}\ {\isadigit{1}}{\isacharparenright}{\kern0pt}\ mod\ NP{\isacharparenright}{\kern0pt}{\isacharparenright}{\kern0pt}\ {\isacharsemicolon}{\kern0pt}\isanewline
\ \ \ \ \ \ \ \ \ \ \ \ \ \ \ \ \ {\isacharparenleft}{\kern0pt}\ {\isacharprime}{\kern0pt}{\isacharprime}{\kern0pt}i{\isacharprime}{\kern0pt}{\isacharprime}{\kern0pt}\ \ \ \ \ {\isacharcolon}{\kern0pt}{\isacharequal}{\kern0pt}\ {\isacharparenleft}{\kern0pt}{\isasymlambda}{\isasymsigma}{\isachardot}{\kern0pt}\ {\isadigit{1}}{\isacharparenright}{\kern0pt}{\isacharparenright}{\kern0pt}\ {\isacharsemicolon}{\kern0pt}\isanewline
\ \ \ \ \ \ \ \ \ \ \ \ \ \ \ \ \ com{\isachardot}{\kern0pt}lock\ pid{\isacharsemicolon}{\kern0pt}\isanewline
\ \ \ \ \ \ \ \ \ \ \ \ \ \ \ \ \ LOAD\ {\isacharprime}{\kern0pt}{\isacharprime}{\kern0pt}x{\isacharprime}{\kern0pt}{\isacharprime}{\kern0pt}\ FROM\ array{\isacharparenleft}{\kern0pt}{\isacharprime}{\kern0pt}{\isacharprime}{\kern0pt}A{\isacharprime}{\kern0pt}{\isacharprime}{\kern0pt}{\isacharcomma}{\kern0pt}pid{\isacharparenright}{\kern0pt}{\isacharsemicolon}{\kern0pt}\isanewline
\ \ \ \ \ \ \ \ \ \ \ \ \ \ \ \ \ com{\isachardot}{\kern0pt}unlock\ pid\ {\isacharsemicolon}{\kern0pt}\isanewline
\ \ \ \ \ \ \ \ \ \ \ \ \ \ \ \ \ {\isacharparenleft}{\kern0pt}\ {\isacharprime}{\kern0pt}{\isacharprime}{\kern0pt}sum{\isacharprime}{\kern0pt}{\isacharprime}{\kern0pt}\ \ \ \ \ {\isacharcolon}{\kern0pt}{\isacharequal}{\kern0pt}\ {\isacharparenleft}{\kern0pt}{\isasymlambda}{\isasymsigma}{\isachardot}{\kern0pt}\ {\isasymsigma}\ {\isacharprime}{\kern0pt}{\isacharprime}{\kern0pt}x{\isacharprime}{\kern0pt}{\isacharprime}{\kern0pt}{\isacharparenright}{\kern0pt}{\isacharparenright}{\kern0pt}\ {\isacharsemicolon}{\kern0pt}\isanewline
\ \ \ \ \ \ \ \ \ \ \ \ \ \ \ \ \ WHILE\ {\isacharparenleft}{\kern0pt}{\isasymlambda}{\isasymsigma}{\isachardot}{\kern0pt}\ {\isasymsigma}\ {\isacharprime}{\kern0pt}{\isacharprime}{\kern0pt}i{\isacharprime}{\kern0pt}{\isacharprime}{\kern0pt}\ {\isacharless}{\kern0pt}\ NP{\isacharparenright}{\kern0pt}\ DO\isanewline
\ \ \ \ \ \ \ \ \ \ \ \ \ \ \ \ \ \ \ \ \ \ {\isacharparenleft}{\kern0pt}com{\isachardot}{\kern0pt}lock\ {\isacharparenleft}{\kern0pt}{\isacharparenleft}{\kern0pt}pid\ {\isacharminus}{\kern0pt}\ {\isadigit{1}}\ {\isacharplus}{\kern0pt}\ NP{\isacharparenright}{\kern0pt}\ mod\ NP{\isacharparenright}{\kern0pt}{\isacharsemicolon}{\kern0pt}\isanewline
\ \ \ \ \ \ \ \ \ \ \ \ \ \ \ \ \ \ \ \ \ \ \ STORE\ {\isacharparenleft}{\kern0pt}{\isasymlambda}{\isasymsigma}{\isachardot}{\kern0pt}\ {\isasymsigma}\ {\isacharprime}{\kern0pt}{\isacharprime}{\kern0pt}x{\isacharprime}{\kern0pt}{\isacharprime}{\kern0pt}{\isacharparenright}{\kern0pt}\ TO\ {\isacharparenleft}{\kern0pt}array{\isacharparenleft}{\kern0pt}{\isacharprime}{\kern0pt}{\isacharprime}{\kern0pt}A{\isacharprime}{\kern0pt}{\isacharprime}{\kern0pt}{\isacharcomma}{\kern0pt}{\isacharparenleft}{\kern0pt}pid\ {\isacharminus}{\kern0pt}\ {\isadigit{1}}\ {\isacharplus}{\kern0pt}\ NP{\isacharparenright}{\kern0pt}\ mod\ NP{\isacharparenright}{\kern0pt}{\isacharparenright}{\kern0pt}{\isacharsemicolon}{\kern0pt}\isanewline
\ \ \ \ \ \ \ \ \ \ \ \ \ \ \ \ \ \ \ \ \ \ \ com{\isachardot}{\kern0pt}unlock\ {\isacharparenleft}{\kern0pt}{\isacharparenleft}{\kern0pt}pid\ {\isacharminus}{\kern0pt}\ {\isadigit{1}}\ {\isacharplus}{\kern0pt}\ NP{\isacharparenright}{\kern0pt}\ mod\ NP{\isacharparenright}{\kern0pt}\ {\isacharsemicolon}{\kern0pt}\isanewline
\ \ \ \ \ \ \ \ \ \ \ \ \ \ \ \ \ \ \ \ \ \ \ com{\isachardot}{\kern0pt}lock\ {\isacharparenleft}{\kern0pt}{\isacharparenleft}{\kern0pt}pid\ {\isacharplus}{\kern0pt}\ {\isadigit{1}}{\isacharparenright}{\kern0pt}\ mod\ NP{\isacharparenright}{\kern0pt}{\isacharsemicolon}{\kern0pt}\isanewline
\ \ \ \ \ \ \ \ \ \ \ \ \ \ \ \ \ \ \ \ \ \ \ LOAD\ {\isacharprime}{\kern0pt}{\isacharprime}{\kern0pt}x{\isacharprime}{\kern0pt}{\isacharprime}{\kern0pt}\ FROM\ array{\isacharparenleft}{\kern0pt}{\isacharprime}{\kern0pt}{\isacharprime}{\kern0pt}A{\isacharprime}{\kern0pt}{\isacharprime}{\kern0pt}{\isacharcomma}{\kern0pt}{\isacharparenleft}{\kern0pt}pid\ {\isacharplus}{\kern0pt}\ {\isadigit{1}}{\isacharparenright}{\kern0pt}\ mod\ NP{\isacharparenright}{\kern0pt}{\isacharsemicolon}{\kern0pt}\isanewline
\ \ \ \ \ \ \ \ \ \ \ \ \ \ \ \ \ \ \ \ \ \ \ com{\isachardot}{\kern0pt}unlock\ {\isacharparenleft}{\kern0pt}{\isacharparenleft}{\kern0pt}pid\ {\isacharplus}{\kern0pt}\ {\isadigit{1}}{\isacharparenright}{\kern0pt}\ mod\ NP{\isacharparenright}{\kern0pt}{\isacharsemicolon}{\kern0pt}\isanewline
\ \ \ \ \ \ \ \ \ \ \ \ \ \ \ \ \ \ \ \ \ \ \ {\isacharparenleft}{\kern0pt}{\isacharprime}{\kern0pt}{\isacharprime}{\kern0pt}sum{\isacharprime}{\kern0pt}{\isacharprime}{\kern0pt}\ {\isacharcolon}{\kern0pt}{\isacharequal}{\kern0pt}\ {\isacharparenleft}{\kern0pt}{\isasymlambda}{\isasymsigma}{\isachardot}{\kern0pt}\ {\isasymsigma}\ {\isacharprime}{\kern0pt}{\isacharprime}{\kern0pt}sum{\isacharprime}{\kern0pt}{\isacharprime}{\kern0pt}\ {\isacharplus}{\kern0pt}\ {\isasymsigma}\ {\isacharprime}{\kern0pt}{\isacharprime}{\kern0pt}x{\isacharprime}{\kern0pt}{\isacharprime}{\kern0pt}{\isacharparenright}{\kern0pt}{\isacharparenright}{\kern0pt}{\isacharsemicolon}{\kern0pt}\isanewline
\ \ \ \ \ \ \ \ \ \ \ \ \ \ \ \ \ \ \ \ \ \ \ {\isacharparenleft}{\kern0pt}{\isacharprime}{\kern0pt}{\isacharprime}{\kern0pt}i{\isacharprime}{\kern0pt}{\isacharprime}{\kern0pt}\ {\isacharcolon}{\kern0pt}{\isacharequal}{\kern0pt}\ {\isacharparenleft}{\kern0pt}{\isasymlambda}{\isasymsigma}{\isachardot}{\kern0pt}\ {\isasymsigma}\ {\isacharprime}{\kern0pt}{\isacharprime}{\kern0pt}i{\isacharprime}{\kern0pt}{\isacharprime}{\kern0pt}\ {\isacharplus}{\kern0pt}\ {\isadigit{1}}{\isacharparenright}{\kern0pt}{\isacharparenright}{\kern0pt}\isanewline
\ \ \ \ \ \ \ \ \ \ \ \ \ \ \ \ \ \ \ \ \ \ {\isacharparenright}{\kern0pt}{\isacharsemicolon}{\kern0pt}\isanewline
\ \ \ \ \ \ \ \ \ \ \ \ \ \ \ \ \ \ assert\ {\isacharparenleft}{\kern0pt}{\isasymlambda}{\isasymsigma}{\isachardot}{\kern0pt}\ {\isasymsigma}\ {\isacharprime}{\kern0pt}{\isacharprime}{\kern0pt}sum{\isacharprime}{\kern0pt}{\isacharprime}{\kern0pt}\ {\isacharequal}{\kern0pt}\ sums\ A{\isacharparenright}{\kern0pt}\isanewline
{\isacartoucheclose}%
\begin{isamarkuptext}%
Some remarks: we assume that the \isa{send} and \isa{recv} semantics of the
above informal notation implies an implicit lock-unlock mechanism  of these
global variables. Moreover, the notation is not quite consequent: while MPI programs
are explicitely designed for fixed architectures (in our case: a ring), the notation
implies access to the local variables which could be modified. In our example \isatt{c{\kern0pt}y{\kern0pt}c{\kern0pt}s{\kern0pt}u{\kern0pt}m{\kern0pt}},
\isa{left} and \isa{right} are constants which have obviously been introduced just to increase
readability. The assertion syntax is, as common in PL-annotation languages, a mixture
of logical variables and program variables which are syntactically identified: the array \isa{A} is
referring to the \emph{content} of the state of the array \isa{A}. This has to be made explicit
in the correctness statement.%
\end{isamarkuptext}\isamarkuptrue%
\isakeywordONE{definition}\isamarkupfalse%
\ Thread{\isadigit{7}}{\isacharunderscore}{\kern0pt}sem\ \isakeywordTWO{where}\ \isanewline
{\isachardoublequoteopen}Thread{\isadigit{7}}{\isacharunderscore}{\kern0pt}sem\ A\ A{\isacharprime}{\kern0pt}\ {\isasymequiv}\ \isanewline
\ \ \ \ \ \ \ \ \ \ {\isacharparenleft}{\kern0pt}\ \isactrlbold {\isacharbar}{\kern0pt}\isactrlbold {\isacharbar}{\kern0pt}\isactrlbold {\isacharbar}{\kern0pt}\ idx\ {\isasymin}{\isacharhash}{\kern0pt}\ mset\ {\isacharbrackleft}{\kern0pt}{\isacharprime}{\kern0pt}{\isacharprime}{\kern0pt}a{\isacharprime}{\kern0pt}{\isacharprime}{\kern0pt}{\isacharcomma}{\kern0pt}{\isacharprime}{\kern0pt}{\isacharprime}{\kern0pt}b{\isacharprime}{\kern0pt}{\isacharprime}{\kern0pt}{\isacharcomma}{\kern0pt}{\isacharprime}{\kern0pt}{\isacharprime}{\kern0pt}c{\isacharprime}{\kern0pt}{\isacharprime}{\kern0pt}{\isacharbrackright}{\kern0pt}{\isachardot}{\kern0pt}\ \ global{\isacharunderscore}{\kern0pt}vars\ {\isasymsigma}\isactrlsub {\isadigit{0}}\ idx\ \ \isanewline
\ \ \ \ \ \ \ \ \ \ \ \ {\isacharbar}{\kern0pt}{\isacharbar}{\kern0pt}{\isacharbar}{\kern0pt}\isanewline
\ \ \ \ \ \ \ \ \ \ \ \ \isactrlbold {\isacharbar}{\kern0pt}\isactrlbold {\isacharbar}{\kern0pt}\isactrlbold {\isacharbar}{\kern0pt}\ idx\ {\isasymin}{\isacharhash}{\kern0pt}\ mset\ {\isacharbrackleft}{\kern0pt}{\isadigit{1}}{\isachardot}{\kern0pt}{\isachardot}{\kern0pt}{\isadigit{4}}{\isacharbrackright}{\kern0pt}{\isachardot}{\kern0pt}\ semaphore\ idx\ {\isacharparenright}{\kern0pt}\ \ \isanewline
\ \ \ \ \ \ \ \ \ \ \ \ \ \ \ \ \ \isanewline
\ \ \ \ \ \ \ \ \ \ {\isacharbar}{\kern0pt}{\isacharbar}{\kern0pt}\isanewline
\isanewline
\ \ \ \ \ \ \ \ \ \ {\isacharparenleft}{\kern0pt}\ \isactrlbold {\isacharbar}{\kern0pt}\isactrlbold {\isacharbar}{\kern0pt}\isactrlbold {\isacharbar}{\kern0pt}\ idx\ {\isasymin}{\isacharhash}{\kern0pt}\ mset\ {\isacharbrackleft}{\kern0pt}{\isadigit{1}}{\isachardot}{\kern0pt}{\isachardot}{\kern0pt}NP{\isacharbrackright}{\kern0pt}{\isachardot}{\kern0pt}\ semaphore\ idx\ {\isacharparenright}{\kern0pt}{\isachardoublequoteclose}%
\begin{isamarkuptext}%
... where \isa{A{\isacharprime}{\kern0pt}}, the array-cell-environment, should be set to its content \isa{A}: 
     \isa{A{\isacharprime}{\kern0pt}{\isacharparenleft}{\kern0pt}array{\isacharparenleft}{\kern0pt}{\isacharprime}{\kern0pt}{\isacharprime}{\kern0pt}A{\isacharprime}{\kern0pt}{\isacharprime}{\kern0pt}{\isacharcomma}{\kern0pt}idx{\isacharparenright}{\kern0pt}{\isacharparenright}{\kern0pt}\ {\isacharequal}{\kern0pt}\ A{\isacharbang}{\kern0pt}{\isacharparenleft}{\kern0pt}idx{\isacharminus}{\kern0pt}{\isadigit{1}}{\isacharparenright}{\kern0pt}} for all idx in \isa{{\isacharbrackleft}{\kern0pt}{\isadigit{1}}{\isachardot}{\kern0pt}{\isachardot}{\kern0pt}NP{\isacharbrackright}{\kern0pt}}. 
     See remark on the difference between array and its content above.%
\end{isamarkuptext}\isamarkuptrue%
%
\begin{isamarkuptext}%
Again, the functional correctness proof boils down to a deadlock-freeness proof:%
\end{isamarkuptext}\isamarkuptrue%
\isakeywordONE{theorem}\isamarkupfalse%
\ correct{\isacharcolon}{\kern0pt}\ \isanewline
\ \ \isakeywordTWO{assumes}\ {\isacartoucheopen}{\isasymforall}\ i\ {\isasymin}\ set{\isacharbrackleft}{\kern0pt}{\isadigit{1}}{\isachardot}{\kern0pt}{\isachardot}{\kern0pt}NP{\isacharbrackright}{\kern0pt}{\isachardot}{\kern0pt}\ A{\isacharprime}{\kern0pt}{\isacharparenleft}{\kern0pt}array{\isacharparenleft}{\kern0pt}{\isacharprime}{\kern0pt}{\isacharprime}{\kern0pt}A{\isacharprime}{\kern0pt}{\isacharprime}{\kern0pt}{\isacharcomma}{\kern0pt}idx{\isacharparenright}{\kern0pt}{\isacharparenright}{\kern0pt}\ {\isacharequal}{\kern0pt}\ A{\isacharbang}{\kern0pt}{\isacharparenleft}{\kern0pt}nat{\isacharparenleft}{\kern0pt}idx{\isacharminus}{\kern0pt}{\isadigit{1}}{\isacharparenright}{\kern0pt}{\isacharparenright}{\kern0pt}{\isacartoucheclose}\ \isanewline
\ \ \ \isakeywordTWO{and}\ \ \ \ {\isacartoucheopen}length\ A\ {\isacharequal}{\kern0pt}\ nat\ NP{\isacartoucheclose}\isanewline
\ \ \isakeywordTWO{shows}\ \ \ {\isacartoucheopen}deadlock{\isacharunderscore}{\kern0pt}free\ {\isacharparenleft}{\kern0pt}Thread{\isadigit{7}}{\isacharunderscore}{\kern0pt}sem\ A\ A{\isacharprime}{\kern0pt}{\isacharparenright}{\kern0pt}{\isacartoucheclose}\isanewline
%
\isadelimproof
\ \ %
\endisadelimproof
%
\isatagproof
\isakeywordONE{unfolding}\isamarkupfalse%
\ deadlock{\isacharunderscore}{\kern0pt}free{\isacharunderscore}{\kern0pt}def\ Thread{\isadigit{7}}{\isacharunderscore}{\kern0pt}sem{\isacharunderscore}{\kern0pt}def\ Sem{\isacharunderscore}{\kern0pt}def\ Thread{\isadigit{7}}{\isacharunderscore}{\kern0pt}def\isanewline
\ \ \isakeywordONE{oops}\isamarkupfalse%
%
\endisatagproof
{\isafoldproof}%
%
\isadelimproof
%
\endisadelimproof
%
\begin{isamarkuptext}%
We consider a formal proof here as out of scope.%
\end{isamarkuptext}\isamarkuptrue%
%
\isadelimdocument
%
\endisadelimdocument
%
\isatagdocument
%
\isamarkupsection{An Architectural Extension: Parameterized Threads%
}
\isamarkuptrue%
%
\endisatagdocument
{\isafolddocument}%
%
\isadelimdocument
%
\endisadelimdocument
%
\begin{isamarkuptext}%
With a tiny bit of tinkering, we can extend the language \isa{IMP\isactrlsub c\isactrlsub o\isactrlsub n\isactrlsub c\isactrlsub u\isactrlsub r} to a language
allowing to express dynamic calculations of architectural connections:%
\end{isamarkuptext}\isamarkuptrue%
\isakeywordONE{type{\isacharunderscore}{\kern0pt}synonym}\isamarkupfalse%
\ {\isachardoublequoteopen}LV\isactrlsub e\isactrlsub x\isactrlsub p\isactrlsub r{\isachardoublequoteclose}\ \ {\isacharequal}{\kern0pt}\ {\isacartoucheopen}{\isasymsigma}\ {\isasymRightarrow}\ V{\isacartoucheclose}\ \ \ \ \ \ \ %
\isamarkupcmt{Local Variable Expr%
}\isanewline
\isakeywordONE{type{\isacharunderscore}{\kern0pt}synonym}\isamarkupfalse%
\ {\isachardoublequoteopen}SV\isactrlsub e\isactrlsub x\isactrlsub p\isactrlsub r{\isachardoublequoteclose}\ \ {\isacharequal}{\kern0pt}\ {\isacartoucheopen}{\isasymsigma}\ {\isasymRightarrow}\ MV{\isacartoucheclose}\ \ \ \ \ \ %
\isamarkupcmt{Semaphore Expr%
}\isanewline
\isakeywordONE{type{\isacharunderscore}{\kern0pt}synonym}\isamarkupfalse%
\ \ GV\isactrlsub e\isactrlsub x\isactrlsub p\isactrlsub r\ \ \ {\isacharequal}{\kern0pt}\ {\isacartoucheopen}{\isasymsigma}\ {\isasymRightarrow}\ string{\isacartoucheclose}\ \ %
\isamarkupcmt{Global Variable Expr%
}\isanewline
\isanewline
\isakeywordONE{datatype}\isamarkupfalse%
\ com\isactrlsub a\ {\isacharequal}{\kern0pt}\ SKIP\isanewline
\ \ \ \ \ \ \ \ \ \ \ \ \ {\isacharbar}{\kern0pt}\ assign\ \ LV\isactrlsub e\isactrlsub x\isactrlsub p\isactrlsub r\ E\isactrlsub a\isactrlsub r\isactrlsub i\isactrlsub t\isactrlsub h\ \ \ \ \ \ \ \ {\isacharparenleft}{\kern0pt}{\isachardoublequoteopen}\ {\isacharunderscore}{\kern0pt}\ {\isacharcolon}{\kern0pt}{\isacharequal}{\kern0pt}\isactrlsub a\ {\isacharunderscore}{\kern0pt}{\isachardoublequoteclose}\ {\isacharbrackleft}{\kern0pt}{\isadigit{9}}{\isadigit{0}}{\isacharcomma}{\kern0pt}{\isadigit{9}}{\isadigit{0}}{\isacharbrackright}{\kern0pt}{\isadigit{8}}{\isadigit{0}}{\isacharparenright}{\kern0pt}\isanewline
\ \ \ \ \ \ \ \ \ \ \ \ \ {\isacharbar}{\kern0pt}\ seq\ \ \ \ \ com\isactrlsub a\ com\isactrlsub a\ \ \ \ \ \ \ \ \ \ {\isacharparenleft}{\kern0pt}\isakeywordTWO{infixl}\ {\isachardoublequoteopen}{\isacharsemicolon}{\kern0pt}{\isacharsemicolon}{\kern0pt}{\isachardoublequoteclose}\ {\isadigit{7}}{\isadigit{8}}{\isacharparenright}{\kern0pt}\isanewline
\ \ \ \ \ \ \ \ \ \ \ \ \ {\isacharbar}{\kern0pt}\ cond\ \ \ {\isachardoublequoteopen}E\isactrlsub b\isactrlsub o\isactrlsub o\isactrlsub l{\isachardoublequoteclose}\ com\isactrlsub a\ com\isactrlsub a\ \ \ \ {\isacharparenleft}{\kern0pt}{\isachardoublequoteopen}IF\isactrlsub a\ {\isacharparenleft}{\kern0pt}{\isacharunderscore}{\kern0pt}{\isacharparenright}{\kern0pt}{\isacharslash}{\kern0pt}\ THEN\ {\isacharparenleft}{\kern0pt}{\isacharunderscore}{\kern0pt}{\isacharparenright}{\kern0pt}{\isacharslash}{\kern0pt}\ ELSE\ {\isacharparenleft}{\kern0pt}{\isacharunderscore}{\kern0pt}{\isacharparenright}{\kern0pt}{\isacharslash}{\kern0pt}{\isachardoublequoteclose}\ {\isacharbrackleft}{\kern0pt}{\isadigit{0}}{\isacharcomma}{\kern0pt}{\isadigit{0}}{\isacharcomma}{\kern0pt}{\isadigit{7}}{\isadigit{9}}{\isacharbrackright}{\kern0pt}{\isadigit{7}}{\isadigit{9}}{\isacharparenright}{\kern0pt}\isanewline
\ \ \ \ \ \ \ \ \ \ \ \ \ {\isacharbar}{\kern0pt}\ while\ \ {\isachardoublequoteopen}E\isactrlsub b\isactrlsub o\isactrlsub o\isactrlsub l{\isachardoublequoteclose}\ com\isactrlsub a\ \ \ \ \ \ \ \ \ {\isacharparenleft}{\kern0pt}{\isachardoublequoteopen}WHILE\isactrlsub a\ {\isacharparenleft}{\kern0pt}{\isacharunderscore}{\kern0pt}{\isacharparenright}{\kern0pt}{\isacharslash}{\kern0pt}\ DO\ {\isacharparenleft}{\kern0pt}{\isacharunderscore}{\kern0pt}{\isacharparenright}{\kern0pt}{\isachardoublequoteclose}\ {\isacharbrackleft}{\kern0pt}{\isadigit{0}}{\isacharcomma}{\kern0pt}{\isadigit{8}}{\isadigit{0}}{\isacharbrackright}{\kern0pt}{\isadigit{8}}{\isadigit{0}}{\isacharparenright}{\kern0pt}\isanewline
\ \ \ \ \ \ \ \ \ \ \ \ \ {\isacharbar}{\kern0pt}\ call\ \ \ \ F\isanewline
\isanewline
\ \ \ \ \ \ \ \ \ \ \ \ \ {\isacharbar}{\kern0pt}\ lock\ \ \ \ SV\isactrlsub e\isactrlsub x\isactrlsub p\isactrlsub r\isanewline
\ \ \ \ \ \ \ \ \ \ \ \ \ {\isacharbar}{\kern0pt}\ unlock\ \ SV\isactrlsub e\isactrlsub x\isactrlsub p\isactrlsub r\ \ \ \ \ \ \ \ \ \ \ \ \isanewline
\ \ \ \ \ \ \ \ \ \ \ \ \ {\isacharbar}{\kern0pt}\ send\ \ \ \ E\isactrlsub a\isactrlsub r\isactrlsub i\isactrlsub t\isactrlsub h\ GV\isactrlsub e\isactrlsub x\isactrlsub p\isactrlsub r\ \ \ \ \ \ \ \ {\isacharparenleft}{\kern0pt}{\isachardoublequoteopen}STORE\isactrlsub a\ {\isacharunderscore}{\kern0pt}\ TO\ \ {\isacharunderscore}{\kern0pt}{\isachardoublequoteclose}\ {\isacharbrackleft}{\kern0pt}{\isadigit{9}}{\isadigit{0}}{\isacharcomma}{\kern0pt}{\isadigit{9}}{\isadigit{0}}{\isacharbrackright}{\kern0pt}{\isadigit{8}}{\isadigit{0}}{\isacharparenright}{\kern0pt}\ \ \ \ \ \ \ \ \ \ \ \ \ \ \isanewline
\ \ \ \ \ \ \ \ \ \ \ \ \ {\isacharbar}{\kern0pt}\ rec\ \ \ \ \ E\isactrlsub a\isactrlsub r\isactrlsub i\isactrlsub t\isactrlsub h\ GV\isactrlsub e\isactrlsub x\isactrlsub p\isactrlsub r\ \ \ \ \ \ \ \ {\isacharparenleft}{\kern0pt}{\isachardoublequoteopen}LOAD\isactrlsub a\ {\isacharunderscore}{\kern0pt}\ FROM\ {\isacharunderscore}{\kern0pt}{\isachardoublequoteclose}\ {\isacharbrackleft}{\kern0pt}{\isadigit{9}}{\isadigit{0}}{\isacharcomma}{\kern0pt}{\isadigit{9}}{\isadigit{0}}{\isacharbrackright}{\kern0pt}{\isadigit{8}}{\isadigit{0}}{\isacharparenright}{\kern0pt}%
\isadelimdocument
%
\endisadelimdocument
%
\isatagdocument
%
\isamarkupsection{Denotational Semantics of \isa{IMP\isactrlsub c\isactrlsub o\isactrlsub n\isactrlsub c\isactrlsub u\isactrlsub r}-P, a parametric version of \isa{IMP\isactrlsub c\isactrlsub o\isactrlsub n\isactrlsub c\isactrlsub u\isactrlsub r}%
}
\isamarkuptrue%
%
\endisatagdocument
{\isafolddocument}%
%
\isadelimdocument
%
\endisadelimdocument
\isakeywordONE{consts}\isamarkupfalse%
\ VarName\ {\isacharcolon}{\kern0pt}{\isacharcolon}{\kern0pt}\ {\isacartoucheopen}int\ {\isasymRightarrow}\ string{\isacartoucheclose}\ %
\isamarkupcmt{converts a reference to a global var name%
}\isanewline
\isanewline
\isakeywordONE{fun}\isamarkupfalse%
\ Sem\isactrlsub a\isactrlsub {\isadigit{0}}\ {\isacharcolon}{\kern0pt}{\isacharcolon}{\kern0pt}\ {\isachardoublequoteopen}com\isactrlsub a\ {\isasymRightarrow}\ {\isacharparenleft}{\kern0pt}{\isasymsigma}\ {\isasymRightarrow}\ evs\ \ process{\isacharparenright}{\kern0pt}\ {\isasymRightarrow}\ {\isasymsigma}\ {\isasymRightarrow}\ evs\ process{\isachardoublequoteclose}\ \isanewline
\ \ \isakeywordTWO{where}\ {\isacartoucheopen}Sem\isactrlsub a\isactrlsub {\isadigit{0}}\ SKIP\ C\ \ \ \ \ \ \ \ \ \ \ \ \ \ \ \ \ \ \ \ {\isacharequal}{\kern0pt}\ C{\isacartoucheclose}\ \isanewline
\ \ \ \ \ \ \ {\isacharbar}{\kern0pt}{\isacartoucheopen}Sem\isactrlsub a\isactrlsub {\isadigit{0}}\ {\isacharparenleft}{\kern0pt}x\ {\isacharcolon}{\kern0pt}{\isacharequal}{\kern0pt}\isactrlsub a\ E{\isacharparenright}{\kern0pt}\ C\ \ \ \ \ \ \ \ \ \ \ \ \ \ \ {\isacharequal}{\kern0pt}\ {\isacharparenleft}{\kern0pt}{\isasymlambda}\ {\isasymsigma}{\isachardot}{\kern0pt}\ C\ {\isacharparenleft}{\kern0pt}{\isasymsigma}{\isacharparenleft}{\kern0pt}{\isacharparenleft}{\kern0pt}x\ {\isasymsigma}{\isacharparenright}{\kern0pt}\ {\isacharcolon}{\kern0pt}{\isacharequal}{\kern0pt}\ E\ {\isasymsigma}{\isacharparenright}{\kern0pt}{\isacharparenright}{\kern0pt}{\isacharparenright}{\kern0pt}{\isacartoucheclose}\isanewline
\ \ \ \ \ \ \ {\isacharbar}{\kern0pt}{\isacartoucheopen}Sem\isactrlsub a\isactrlsub {\isadigit{0}}\ {\isacharparenleft}{\kern0pt}P\ {\isacharsemicolon}{\kern0pt}{\isacharsemicolon}{\kern0pt}\ Q{\isacharparenright}{\kern0pt}\ C\ \ \ \ \ \ \ \ \ \ \ \ \ \ \ \ {\isacharequal}{\kern0pt}\ {\isacharparenleft}{\kern0pt}Sem\isactrlsub a\isactrlsub {\isadigit{0}}\ P\ {\isacharparenleft}{\kern0pt}Sem\isactrlsub a\isactrlsub {\isadigit{0}}\ Q\ C{\isacharparenright}{\kern0pt}{\isacharparenright}{\kern0pt}{\isacartoucheclose}\ \isanewline
\ \ \ \ \ \ \ {\isacharbar}{\kern0pt}{\isacartoucheopen}Sem\isactrlsub a\isactrlsub {\isadigit{0}}\ {\isacharparenleft}{\kern0pt}IF\isactrlsub a\ E\ THEN\ C{\isadigit{1}}\ ELSE\ C{\isadigit{2}}{\isacharparenright}{\kern0pt}\ C\ {\isacharequal}{\kern0pt}\ {\isacharparenleft}{\kern0pt}{\isasymlambda}\ {\isasymsigma}{\isachardot}{\kern0pt}\ if\ E\ {\isasymsigma}\ \isanewline
\ \ \ \ \ \ \ \ \ \ \ \ \ \ \ \ \ \ \ \ \ \ \ \ \ \ \ \ \ \ \ \ \ \ \ \ \ \ \ \ \ \ \ \ \ \ \ \ then\ Sem\isactrlsub a\isactrlsub {\isadigit{0}}\ C{\isadigit{1}}\ C\ {\isasymsigma}\ \isanewline
\ \ \ \ \ \ \ \ \ \ \ \ \ \ \ \ \ \ \ \ \ \ \ \ \ \ \ \ \ \ \ \ \ \ \ \ \ \ \ \ \ \ \ \ \ \ \ \ else\ Sem\isactrlsub a\isactrlsub {\isadigit{0}}\ C{\isadigit{2}}\ C\ {\isasymsigma}{\isacharparenright}{\kern0pt}{\isacartoucheclose}\isanewline
\ \ \ \ \ \ \ {\isacharbar}{\kern0pt}{\isacartoucheopen}Sem\isactrlsub a\isactrlsub {\isadigit{0}}\ {\isacharparenleft}{\kern0pt}WHILE\isactrlsub a\ E\ DO\ B{\isacharparenright}{\kern0pt}\ C\ \ \ \ \ \ \ \ \ {\isacharequal}{\kern0pt}\ {\isacharparenleft}{\kern0pt}{\isasymmu}\ X{\isachardot}{\kern0pt}\ {\isacharparenleft}{\kern0pt}{\isasymlambda}\ {\isasymsigma}{\isachardot}{\kern0pt}\ if\ E\ {\isasymsigma}\ \isanewline
\ \ \ \ \ \ \ \ \ \ \ \ \ \ \ \ \ \ \ \ \ \ \ \ \ \ \ \ \ \ \ \ \ \ \ \ \ \ \ \ \ \ \ \ \ \ \ \ \ \ \ \ \ \ then\ Sem\isactrlsub a\isactrlsub {\isadigit{0}}\ B\ X\ {\isasymsigma}\ \isanewline
\ \ \ \ \ \ \ \ \ \ \ \ \ \ \ \ \ \ \ \ \ \ \ \ \ \ \ \ \ \ \ \ \ \ \ \ \ \ \ \ \ \ \ \ \ \ \ \ \ \ \ \ \ \ else\ C\ {\isasymsigma}{\isacharparenright}{\kern0pt}{\isacharparenright}{\kern0pt}{\isacartoucheclose}\isanewline
\ \ \ \ \ \ \ {\isacharbar}{\kern0pt}{\isacartoucheopen}Sem\isactrlsub a\isactrlsub {\isadigit{0}}\ {\isacharparenleft}{\kern0pt}com\isactrlsub a{\isachardot}{\kern0pt}call\ F{\isacharparenright}{\kern0pt}\ C\ \ \ \ \ \ \ \ \ \ \ {\isacharequal}{\kern0pt}\ {\isacharparenleft}{\kern0pt}{\isasymlambda}\ {\isasymsigma}{\isachardot}{\kern0pt}\ C\ {\isacharparenleft}{\kern0pt}F\ {\isasymsigma}{\isacharparenright}{\kern0pt}{\isacharparenright}{\kern0pt}{\isacartoucheclose}\isanewline
\ \ \ \ \ \ \ {\isacharbar}{\kern0pt}{\isacartoucheopen}Sem\isactrlsub a\isactrlsub {\isadigit{0}}\ {\isacharparenleft}{\kern0pt}com\isactrlsub a{\isachardot}{\kern0pt}lock\ n{\isacharparenright}{\kern0pt}\ C\ \ \ \ \ \ \ \ \ \ \ {\isacharequal}{\kern0pt}\ {\isacharparenleft}{\kern0pt}{\isasymlambda}\ {\isasymsigma}{\isachardot}{\kern0pt}\ evs{\isachardot}{\kern0pt}lock\ {\isacharparenleft}{\kern0pt}n\ {\isasymsigma}{\isacharparenright}{\kern0pt}\ {\isasymrightarrow}\ C\ {\isasymsigma}{\isacharparenright}{\kern0pt}{\isacartoucheclose}\isanewline
\ \ \ \ \ \ \ {\isacharbar}{\kern0pt}{\isacartoucheopen}Sem\isactrlsub a\isactrlsub {\isadigit{0}}\ {\isacharparenleft}{\kern0pt}com\isactrlsub a{\isachardot}{\kern0pt}unlock\ n{\isacharparenright}{\kern0pt}\ C\ \ \ \ \ \ \ \ \ {\isacharequal}{\kern0pt}\ {\isacharparenleft}{\kern0pt}{\isasymlambda}\ {\isasymsigma}{\isachardot}{\kern0pt}\ evs{\isachardot}{\kern0pt}unlock\ {\isacharparenleft}{\kern0pt}n\ {\isasymsigma}{\isacharparenright}{\kern0pt}\ \ {\isasymrightarrow}\ C\ {\isasymsigma}{\isacharparenright}{\kern0pt}{\isacartoucheclose}\isanewline
\ \ \ \ \ \ \ {\isacharbar}{\kern0pt}{\isacartoucheopen}Sem\isactrlsub a\isactrlsub {\isadigit{0}}\ {\isacharparenleft}{\kern0pt}STORE\isactrlsub a\ E\ TO\ X\isactrlsub g\isactrlsub l\isactrlsub o{\isacharparenright}{\kern0pt}\ C\ \ \ \ \ \ \ {\isacharequal}{\kern0pt}\ {\isacharparenleft}{\kern0pt}{\isasymlambda}\ {\isasymsigma}{\isachardot}{\kern0pt}\ upd\isactrlsub g\isactrlsub l\isactrlsub o\isactrlsub b\isactrlsub a\isactrlsub l\ {\isacharparenleft}{\kern0pt}X\isactrlsub g\isactrlsub l\isactrlsub o\ {\isasymsigma}{\isacharparenright}{\kern0pt}\ {\isacharparenleft}{\kern0pt}E\ {\isasymsigma}{\isacharparenright}{\kern0pt}\ {\isasymrightarrow}\ C\ {\isasymsigma}{\isacharparenright}{\kern0pt}{\isacartoucheclose}\isanewline
\ \ \ \ \ \ \ {\isacharbar}{\kern0pt}{\isacartoucheopen}Sem\isactrlsub a\isactrlsub {\isadigit{0}}\ {\isacharparenleft}{\kern0pt}LOAD\isactrlsub a\ X\isactrlsub l\isactrlsub o\isactrlsub c\ FROM\ \ X\isactrlsub g\isactrlsub l\isactrlsub o{\isacharparenright}{\kern0pt}\ C\ \ \ {\isacharequal}{\kern0pt}\ {\isacharparenleft}{\kern0pt}{\isasymlambda}\ {\isasymsigma}{\isachardot}{\kern0pt}\ {\isacharparenleft}{\kern0pt}read\isactrlsub g\isactrlsub l\isactrlsub o\isactrlsub b\isactrlsub a\isactrlsub l\ {\isacharparenleft}{\kern0pt}X\isactrlsub g\isactrlsub l\isactrlsub o\ {\isasymsigma}{\isacharparenright}{\kern0pt}{\isacharparenright}{\kern0pt}\isactrlbold {\isacharquery}{\kern0pt}x\isanewline
\ \ \ \ \ \ \ \ \ \ \ \ \ \ \ \ \ \ \ \ \ \ \ \ \ \ \ \ \ \ \ \ \ \ \ \ \ \ \ \ \ \ \ \ \ \ \ \ \ \ \ \ \ {\isasymrightarrow}\ C{\isacharparenleft}{\kern0pt}{\isasymsigma}{\isacharparenleft}{\kern0pt}{\isacharparenleft}{\kern0pt}VarName{\isacharparenleft}{\kern0pt}X\isactrlsub l\isactrlsub o\isactrlsub c\ {\isasymsigma}{\isacharparenright}{\kern0pt}{\isacharparenright}{\kern0pt}{\isacharcolon}{\kern0pt}{\isacharequal}{\kern0pt}x{\isacharparenright}{\kern0pt}{\isacharparenright}{\kern0pt}{\isacharparenright}{\kern0pt}{\isacartoucheclose}%
\isadelimdocument
%
\endisadelimdocument
%
\isatagdocument
%
\isamarkupsection{Conclusion and Future Work%
}
\isamarkuptrue%
%
\endisatagdocument
{\isafolddocument}%
%
\isadelimdocument
%
\endisadelimdocument
%
\begin{isamarkuptext}%
We have shown the language and semantics of \isa{IMP\isactrlsub c\isactrlsub o\isactrlsub n\isactrlsub c\isactrlsub u\isactrlsub r}, a thin layer over HOL-CSP
giving this process-algebra a more programming language flavor. It is conceived to 
show the close link between both worlds considered quite distinct by many.

The embedding provides:

%
\begin{enumerate}%
\item a clear semantic foundation via HOL-CSP to HOL,
comprising  computation and concurrency, 

\item an end-to-end verification of a verification method inside 
a highly trustable interactive proof assistant, 

\item a path to theorem proving of functional properties,
via deadlock-freeness proofs,

\item a path to model-checking of functional properties,
via deadlock-freeness inside FDR4 (unpublished so far).%
\end{enumerate}%
\end{isamarkuptext}\isamarkuptrue%
%
\begin{isamarkuptext}%
An interesting direction for future work is, for example, how  a Hoare Calculus for \isa{IMP\isactrlsub c\isactrlsub o\isactrlsub n\isactrlsub c\isactrlsub u\isactrlsub r} 
would look like, or how a conversion to FDR4 could be used in a SPIN-like manner 
for finitized programs as model-checking-based backend for practical purposes.%
\end{isamarkuptext}\isamarkuptrue%
%
\isadelimtheory
%
\endisadelimtheory
%
\isatagtheory
%
\endisatagtheory
{\isafoldtheory}%
%
\isadelimtheory
%
\endisadelimtheory
%
\end{isabellebody}%
\endinput
%:%file=$AFP/IMP_concur/IMP_concur.thy%:%
%:%11=1%:%
%:%27=3%:%
%:%28=3%:%
%:%29=4%:%
%:%30=5%:%
%:%31=6%:%
%:%45=8%:%
%:%57=11%:%
%:%58=12%:%
%:%59=13%:%
%:%60=14%:%
%:%61=15%:%
%:%63=17%:%
%:%64=18%:%
%:%65=19%:%
%:%69=20%:%
%:%71=21%:%
%:%72=22%:%
%:%75=24%:%
%:%79=25%:%
%:%81=26%:%
%:%83=27%:%
%:%84=28%:%
%:%86=29%:%
%:%89=31%:%
%:%90=32%:%
%:%99=35%:%
%:%109=37%:%
%:%110=37%:%
%:%111=37%:%
%:%112=37%:%
%:%113=38%:%
%:%114=38%:%
%:%115=38%:%
%:%116=38%:%
%:%117=39%:%
%:%118=39%:%
%:%119=39%:%
%:%120=39%:%
%:%121=40%:%
%:%122=41%:%
%:%123=41%:%
%:%124=42%:%
%:%125=43%:%
%:%126=43%:%
%:%127=44%:%
%:%128=45%:%
%:%129=45%:%
%:%130=46%:%
%:%131=46%:%
%:%132=47%:%
%:%133=47%:%
%:%134=48%:%
%:%135=49%:%
%:%136=49%:%
%:%137=50%:%
%:%138=51%:%
%:%139=52%:%
%:%140=53%:%
%:%141=54%:%
%:%142=55%:%
%:%143=56%:%
%:%144=57%:%
%:%145=58%:%
%:%146=59%:%
%:%147=60%:%
%:%148=61%:%
%:%150=64%:%
%:%152=66%:%
%:%153=66%:%
%:%155=68%:%
%:%159=70%:%
%:%161=72%:%
%:%162=72%:%
%:%163=73%:%
%:%170=80%:%
%:%171=81%:%
%:%172=82%:%
%:%173=82%:%
%:%174=83%:%
%:%181=90%:%
%:%182=91%:%
%:%183=92%:%
%:%184=92%:%
%:%185=93%:%
%:%186=93%:%
%:%187=94%:%
%:%193=100%:%
%:%196=101%:%
%:%201=102%:%
%:%202=102%:%
%:%217=106%:%
%:%229=108%:%
%:%231=110%:%
%:%232=110%:%
%:%233=111%:%
%:%234=112%:%
%:%235=113%:%
%:%236=113%:%
%:%237=114%:%
%:%238=114%:%
%:%239=115%:%
%:%240=116%:%
%:%241=117%:%
%:%242=117%:%
%:%244=117%:%
%:%248=117%:%
%:%249=117%:%
%:%250=117%:%
%:%257=117%:%
%:%258=118%:%
%:%259=118%:%
%:%261=118%:%
%:%265=118%:%
%:%266=118%:%
%:%267=118%:%
%:%274=118%:%
%:%275=119%:%
%:%276=119%:%
%:%278=119%:%
%:%282=119%:%
%:%283=119%:%
%:%284=119%:%
%:%291=119%:%
%:%292=120%:%
%:%293=120%:%
%:%295=120%:%
%:%299=120%:%
%:%300=120%:%
%:%301=120%:%
%:%308=120%:%
%:%309=121%:%
%:%310=122%:%
%:%311=122%:%
%:%313=122%:%
%:%317=122%:%
%:%318=122%:%
%:%319=122%:%
%:%326=122%:%
%:%327=123%:%
%:%328=123%:%
%:%330=123%:%
%:%334=123%:%
%:%335=123%:%
%:%336=123%:%
%:%343=123%:%
%:%344=124%:%
%:%345=124%:%
%:%347=124%:%
%:%351=124%:%
%:%352=124%:%
%:%353=124%:%
%:%360=124%:%
%:%361=125%:%
%:%362=125%:%
%:%364=125%:%
%:%368=125%:%
%:%369=125%:%
%:%370=125%:%
%:%377=125%:%
%:%378=126%:%
%:%379=127%:%
%:%380=127%:%
%:%382=127%:%
%:%386=127%:%
%:%387=127%:%
%:%388=127%:%
%:%395=127%:%
%:%396=128%:%
%:%397=128%:%
%:%399=128%:%
%:%403=128%:%
%:%404=128%:%
%:%405=128%:%
%:%412=128%:%
%:%413=129%:%
%:%414=129%:%
%:%416=129%:%
%:%420=129%:%
%:%421=129%:%
%:%422=129%:%
%:%429=129%:%
%:%430=130%:%
%:%431=130%:%
%:%433=130%:%
%:%437=130%:%
%:%438=130%:%
%:%439=130%:%
%:%453=134%:%
%:%463=139%:%
%:%464=139%:%
%:%465=140%:%
%:%466=141%:%
%:%467=142%:%
%:%468=142%:%
%:%471=143%:%
%:%475=143%:%
%:%476=143%:%
%:%490=145%:%
%:%500=147%:%
%:%501=147%:%
%:%502=148%:%
%:%503=149%:%
%:%504=150%:%
%:%505=151%:%
%:%506=152%:%
%:%507=152%:%
%:%508=153%:%
%:%511=154%:%
%:%515=154%:%
%:%516=154%:%
%:%530=157%:%
%:%540=159%:%
%:%541=159%:%
%:%542=160%:%
%:%543=161%:%
%:%544=162%:%
%:%545=163%:%
%:%546=163%:%
%:%547=164%:%
%:%550=165%:%
%:%554=165%:%
%:%555=165%:%
%:%569=168%:%
%:%581=170%:%
%:%582=171%:%
%:%583=172%:%
%:%584=173%:%
%:%586=175%:%
%:%587=175%:%
%:%588=176%:%
%:%589=177%:%
%:%590=178%:%
%:%591=179%:%
%:%593=181%:%
%:%594=182%:%
%:%596=184%:%
%:%597=185%:%
%:%598=186%:%
%:%599=187%:%
%:%600=188%:%
%:%601=189%:%
%:%602=190%:%
%:%603=191%:%
%:%604=193%:%
%:%605=194%:%
%:%606=195%:%
%:%607=195%:%
%:%608=196%:%
%:%609=197%:%
%:%610=198%:%
%:%611=198%:%
%:%612=199%:%
%:%614=201%:%
%:%615=202%:%
%:%616=203%:%
%:%618=206%:%
%:%619=206%:%
%:%620=207%:%
%:%621=207%:%
%:%622=208%:%
%:%630=210%:%
%:%642=212%:%
%:%644=214%:%
%:%645=214%:%
%:%652=221%:%
%:%653=222%:%
%:%656=223%:%
%:%661=224%:%
%:%662=224%:%
%:%674=229%:%
%:%676=231%:%
%:%677=231%:%
%:%678=232%:%
%:%681=233%:%
%:%685=233%:%
%:%686=233%:%
%:%687=234%:%
%:%688=234%:%
%:%693=234%:%
%:%696=235%:%
%:%697=236%:%
%:%698=237%:%
%:%699=237%:%
%:%702=238%:%
%:%706=238%:%
%:%707=238%:%
%:%708=239%:%
%:%709=239%:%
%:%710=240%:%
%:%711=240%:%
%:%712=241%:%
%:%713=241%:%
%:%714=242%:%
%:%715=242%:%
%:%720=242%:%
%:%723=243%:%
%:%724=244%:%
%:%725=245%:%
%:%726=246%:%
%:%727=246%:%
%:%728=247%:%
%:%739=258%:%
%:%742=259%:%
%:%746=259%:%
%:%747=259%:%
%:%748=260%:%
%:%749=260%:%
%:%754=260%:%
%:%757=261%:%
%:%758=262%:%
%:%761=263%:%
%:%766=264%:%
%:%767=264%:%
%:%786=271%:%
%:%798=273%:%
%:%799=274%:%
%:%800=275%:%
%:%801=276%:%
%:%803=278%:%
%:%804=278%:%
%:%805=279%:%
%:%806=280%:%
%:%807=281%:%
%:%808=281%:%
%:%809=282%:%
%:%817=290%:%
%:%819=291%:%
%:%820=291%:%
%:%823=292%:%
%:%827=292%:%
%:%828=292%:%
%:%829=293%:%
%:%830=293%:%
%:%844=296%:%
%:%856=298%:%
%:%858=300%:%
%:%859=300%:%
%:%860=301%:%
%:%861=302%:%
%:%862=303%:%
%:%863=303%:%
%:%864=304%:%
%:%868=308%:%
%:%869=309%:%
%:%870=310%:%
%:%871=310%:%
%:%872=311%:%
%:%875=312%:%
%:%879=312%:%
%:%880=312%:%
%:%889=314%:%
%:%890=315%:%
%:%891=316%:%
%:%892=317%:%
%:%893=318%:%
%:%894=319%:%
%:%903=322%:%
%:%915=324%:%
%:%916=325%:%
%:%917=326%:%
%:%918=327%:%
%:%919=328%:%
%:%923=331%:%
%:%924=332%:%
%:%925=333%:%
%:%926=334%:%
%:%927=335%:%
%:%930=335%:%
%:%931=336%:%
%:%932=337%:%
%:%933=338%:%
%:%934=339%:%
%:%936=341%:%
%:%937=341%:%
%:%938=342%:%
%:%939=343%:%
%:%941=345%:%
%:%944=346%:%
%:%948=346%:%
%:%949=346%:%
%:%950=347%:%
%:%951=347%:%
%:%956=347%:%
%:%959=348%:%
%:%960=349%:%
%:%961=350%:%
%:%962=350%:%
%:%963=351%:%
%:%964=352%:%
%:%966=354%:%
%:%969=355%:%
%:%973=355%:%
%:%974=355%:%
%:%975=356%:%
%:%976=356%:%
%:%977=357%:%
%:%978=357%:%
%:%983=357%:%
%:%986=358%:%
%:%987=359%:%
%:%988=359%:%
%:%989=360%:%
%:%990=361%:%
%:%993=362%:%
%:%997=362%:%
%:%998=362%:%
%:%999=363%:%
%:%1000=363%:%
%:%1001=364%:%
%:%1010=366%:%
%:%1011=367%:%
%:%1012=368%:%
%:%1013=369%:%
%:%1015=371%:%
%:%1016=371%:%
%:%1017=372%:%
%:%1018=373%:%
%:%1019=374%:%
%:%1026=375%:%
%:%1027=375%:%
%:%1028=376%:%
%:%1029=376%:%
%:%1043=390%:%
%:%1044=391%:%
%:%1045=392%:%
%:%1046=392%:%
%:%1047=393%:%
%:%1048=393%:%
%:%1049=394%:%
%:%1050=394%:%
%:%1051=395%:%
%:%1052=396%:%
%:%1053=397%:%
%:%1054=398%:%
%:%1055=398%:%
%:%1056=399%:%
%:%1057=399%:%
%:%1058=400%:%
%:%1059=400%:%
%:%1060=401%:%
%:%1061=402%:%
%:%1062=403%:%
%:%1063=404%:%
%:%1064=405%:%
%:%1065=405%:%
%:%1066=406%:%
%:%1067=406%:%
%:%1068=406%:%
%:%1069=406%:%
%:%1070=407%:%
%:%1071=407%:%
%:%1072=408%:%
%:%1073=408%:%
%:%1074=409%:%
%:%1075=409%:%
%:%1076=410%:%
%:%1077=410%:%
%:%1078=411%:%
%:%1079=411%:%
%:%1080=412%:%
%:%1081=413%:%
%:%1082=413%:%
%:%1083=414%:%
%:%1084=414%:%
%:%1085=415%:%
%:%1086=415%:%
%:%1087=416%:%
%:%1088=416%:%
%:%1089=417%:%
%:%1095=417%:%
%:%1098=418%:%
%:%1099=419%:%
%:%1100=420%:%
%:%1101=420%:%
%:%1102=421%:%
%:%1104=423%:%
%:%1105=424%:%
%:%1112=425%:%
%:%1113=425%:%
%:%1114=426%:%
%:%1115=426%:%
%:%1138=449%:%
%:%1139=450%:%
%:%1140=450%:%
%:%1141=451%:%
%:%1142=451%:%
%:%1143=452%:%
%:%1144=453%:%
%:%1145=453%:%
%:%1146=454%:%
%:%1147=454%:%
%:%1148=455%:%
%:%1149=456%:%
%:%1150=457%:%
%:%1151=458%:%
%:%1152=459%:%
%:%1153=459%:%
%:%1154=460%:%
%:%1155=460%:%
%:%1156=461%:%
%:%1157=461%:%
%:%1158=462%:%
%:%1159=462%:%
%:%1160=463%:%
%:%1170=466%:%
%:%1171=467%:%
%:%1172=468%:%
%:%1173=469%:%
%:%1174=470%:%
%:%1175=471%:%
%:%1176=472%:%
%:%1178=475%:%
%:%1179=475%:%
%:%1182=476%:%
%:%1186=476%:%
%:%1187=476%:%
%:%1188=477%:%
%:%1189=477%:%
%:%1190=478%:%
%:%1191=478%:%
%:%1193=480%:%
%:%1194=481%:%
%:%1195=482%:%
%:%1196=482%:%
%:%1197=483%:%
%:%1198=483%:%
%:%1199=483%:%
%:%1200=483%:%
%:%1201=484%:%
%:%1202=484%:%
%:%1204=486%:%
%:%1205=487%:%
%:%1206=488%:%
%:%1207=488%:%
%:%1208=489%:%
%:%1209=489%:%
%:%1210=490%:%
%:%1211=490%:%
%:%1212=491%:%
%:%1213=491%:%
%:%1214=492%:%
%:%1215=492%:%
%:%1216=493%:%
%:%1217=493%:%
%:%1218=494%:%
%:%1219=494%:%
%:%1220=494%:%
%:%1221=495%:%
%:%1222=495%:%
%:%1223=496%:%
%:%1233=499%:%
%:%1234=500%:%
%:%1236=502%:%
%:%1237=502%:%
%:%1239=504%:%
%:%1240=505%:%
%:%1241=506%:%
%:%1242=506%:%
%:%1245=509%:%
%:%1246=510%:%
%:%1247=511%:%
%:%1248=511%:%
%:%1258=516%:%
%:%1270=518%:%
%:%1271=519%:%
%:%1273=521%:%
%:%1274=521%:%
%:%1275=522%:%
%:%1276=523%:%
%:%1277=524%:%
%:%1278=524%:%
%:%1279=525%:%
%:%1280=526%:%
%:%1283=527%:%
%:%1287=527%:%
%:%1288=527%:%
%:%1289=528%:%
%:%1294=528%:%
%:%1297=529%:%
%:%1298=530%:%
%:%1299=531%:%
%:%1300=531%:%
%:%1301=532%:%
%:%1307=538%:%
%:%1308=539%:%
%:%1309=540%:%
%:%1313=543%:%
%:%1315=545%:%
%:%1316=545%:%
%:%1318=547%:%
%:%1319=548%:%
%:%1320=549%:%
%:%1321=550%:%
%:%1322=551%:%
%:%1323=552%:%
%:%1332=554%:%
%:%1344=556%:%
%:%1346=558%:%
%:%1347=558%:%
%:%1354=559%:%
%:%1355=559%:%
%:%1356=560%:%
%:%1357=560%:%
%:%1358=561%:%
%:%1359=561%:%
%:%1360=561%:%
%:%1361=562%:%
%:%1362=562%:%
%:%1363=563%:%
%:%1364=563%:%
%:%1365=564%:%
%:%1366=564%:%
%:%1367=565%:%
%:%1368=566%:%
%:%1369=566%:%
%:%1370=567%:%
%:%1371=567%:%
%:%1372=568%:%
%:%1373=569%:%
%:%1374=569%:%
%:%1375=570%:%
%:%1376=570%:%
%:%1377=571%:%
%:%1378=572%:%
%:%1379=572%:%
%:%1380=573%:%
%:%1381=573%:%
%:%1382=574%:%
%:%1383=574%:%
%:%1384=575%:%
%:%1385=575%:%
%:%1386=576%:%
%:%1387=576%:%
%:%1388=577%:%
%:%1389=577%:%
%:%1390=578%:%
%:%1391=579%:%
%:%1392=579%:%
%:%1393=580%:%
%:%1394=581%:%
%:%1395=581%:%
%:%1396=582%:%
%:%1397=582%:%
%:%1398=583%:%
%:%1408=587%:%
%:%1409=588%:%
%:%1411=589%:%
%:%1412=589%:%
%:%1413=590%:%
%:%1414=591%:%
%:%1417=592%:%
%:%1421=592%:%
%:%1422=592%:%
%:%1423=593%:%
%:%1424=593%:%
%:%1429=593%:%
%:%1432=594%:%
%:%1433=595%:%
%:%1434=595%:%
%:%1435=596%:%
%:%1436=597%:%
%:%1439=598%:%
%:%1443=598%:%
%:%1444=598%:%
%:%1445=599%:%
%:%1446=599%:%
%:%1455=601%:%
%:%1456=602%:%
%:%1457=603%:%
%:%1459=606%:%
%:%1460=606%:%
%:%1461=607%:%
%:%1462=608%:%
%:%1468=614%:%
%:%1469=615%:%
%:%1476=616%:%
%:%1477=616%:%
%:%1478=617%:%
%:%1479=617%:%
%:%1494=632%:%
%:%1495=633%:%
%:%1496=634%:%
%:%1497=634%:%
%:%1498=635%:%
%:%1499=635%:%
%:%1500=636%:%
%:%1501=636%:%
%:%1502=637%:%
%:%1503=638%:%
%:%1504=638%:%
%:%1505=639%:%
%:%1506=639%:%
%:%1507=640%:%
%:%1508=640%:%
%:%1509=641%:%
%:%1510=642%:%
%:%1511=643%:%
%:%1512=644%:%
%:%1513=645%:%
%:%1514=645%:%
%:%1515=646%:%
%:%1516=646%:%
%:%1517=647%:%
%:%1518=647%:%
%:%1519=648%:%
%:%1520=648%:%
%:%1521=649%:%
%:%1522=649%:%
%:%1523=650%:%
%:%1524=650%:%
%:%1525=651%:%
%:%1526=651%:%
%:%1527=652%:%
%:%1528=653%:%
%:%1529=653%:%
%:%1530=654%:%
%:%1531=654%:%
%:%1532=655%:%
%:%1533=655%:%
%:%1534=656%:%
%:%1535=656%:%
%:%1536=657%:%
%:%1537=657%:%
%:%1538=658%:%
%:%1544=658%:%
%:%1547=659%:%
%:%1548=660%:%
%:%1549=661%:%
%:%1550=661%:%
%:%1551=662%:%
%:%1552=663%:%
%:%1561=672%:%
%:%1562=673%:%
%:%1569=674%:%
%:%1570=674%:%
%:%1571=675%:%
%:%1572=675%:%
%:%1596=699%:%
%:%1597=700%:%
%:%1598=700%:%
%:%1599=701%:%
%:%1600=701%:%
%:%1601=702%:%
%:%1602=703%:%
%:%1603=703%:%
%:%1604=704%:%
%:%1605=704%:%
%:%1606=705%:%
%:%1607=706%:%
%:%1608=707%:%
%:%1609=708%:%
%:%1610=709%:%
%:%1611=709%:%
%:%1612=709%:%
%:%1613=710%:%
%:%1614=710%:%
%:%1639=735%:%
%:%1640=736%:%
%:%1641=736%:%
%:%1642=737%:%
%:%1643=737%:%
%:%1644=738%:%
%:%1645=739%:%
%:%1646=739%:%
%:%1647=740%:%
%:%1648=740%:%
%:%1649=741%:%
%:%1650=742%:%
%:%1651=743%:%
%:%1652=744%:%
%:%1653=745%:%
%:%1654=745%:%
%:%1655=746%:%
%:%1656=746%:%
%:%1657=747%:%
%:%1658=747%:%
%:%1659=748%:%
%:%1660=748%:%
%:%1661=749%:%
%:%1662=749%:%
%:%1663=750%:%
%:%1664=750%:%
%:%1665=751%:%
%:%1666=751%:%
%:%1667=752%:%
%:%1668=752%:%
%:%1669=753%:%
%:%1670=753%:%
%:%1671=754%:%
%:%1672=754%:%
%:%1673=755%:%
%:%1674=755%:%
%:%1675=756%:%
%:%1676=757%:%
%:%1686=760%:%
%:%1688=762%:%
%:%1689=762%:%
%:%1692=763%:%
%:%1696=763%:%
%:%1697=763%:%
%:%1698=764%:%
%:%1699=764%:%
%:%1700=765%:%
%:%1701=765%:%
%:%1704=768%:%
%:%1705=769%:%
%:%1706=770%:%
%:%1707=770%:%
%:%1708=771%:%
%:%1709=771%:%
%:%1710=771%:%
%:%1711=771%:%
%:%1712=772%:%
%:%1713=772%:%
%:%1714=772%:%
%:%1715=772%:%
%:%1716=773%:%
%:%1717=773%:%
%:%1719=775%:%
%:%1720=776%:%
%:%1721=776%:%
%:%1722=776%:%
%:%1723=777%:%
%:%1724=777%:%
%:%1726=779%:%
%:%1727=780%:%
%:%1728=781%:%
%:%1729=781%:%
%:%1730=781%:%
%:%1731=782%:%
%:%1732=782%:%
%:%1734=784%:%
%:%1735=785%:%
%:%1736=785%:%
%:%1737=785%:%
%:%1738=786%:%
%:%1739=786%:%
%:%1740=787%:%
%:%1741=787%:%
%:%1742=788%:%
%:%1743=788%:%
%:%1744=789%:%
%:%1745=789%:%
%:%1746=790%:%
%:%1747=790%:%
%:%1748=791%:%
%:%1749=791%:%
%:%1750=792%:%
%:%1751=792%:%
%:%1752=793%:%
%:%1753=793%:%
%:%1754=794%:%
%:%1769=797%:%
%:%1781=799%:%
%:%1782=800%:%
%:%1783=801%:%
%:%1784=802%:%
%:%1785=803%:%
%:%1786=804%:%
%:%1787=805%:%
%:%1788=806%:%
%:%1789=807%:%
%:%1790=808%:%
%:%1791=809%:%
%:%1792=810%:%
%:%1793=811%:%
%:%1794=812%:%
%:%1795=813%:%
%:%1796=814%:%
%:%1813=829%:%
%:%1814=830%:%
%:%1815=831%:%
%:%1816=832%:%
%:%1817=833%:%
%:%1818=834%:%
%:%1819=835%:%
%:%1820=836%:%
%:%1822=838%:%
%:%1823=838%:%
%:%1824=838%:%
%:%1825=838%:%
%:%1826=839%:%
%:%1827=840%:%
%:%1828=840%:%
%:%1829=841%:%
%:%1830=841%:%
%:%1831=842%:%
%:%1832=842%:%
%:%1833=843%:%
%:%1834=843%:%
%:%1835=844%:%
%:%1836=844%:%
%:%1837=844%:%
%:%1838=845%:%
%:%1839=846%:%
%:%1840=846%:%
%:%1841=847%:%
%:%1842=847%:%
%:%1843=848%:%
%:%1844=849%:%
%:%1845=849%:%
%:%1846=850%:%
%:%1866=870%:%
%:%1867=871%:%
%:%1868=872%:%
%:%1869=873%:%
%:%1870=874%:%
%:%1871=875%:%
%:%1872=876%:%
%:%1873=877%:%
%:%1874=878%:%
%:%1875=879%:%
%:%1877=883%:%
%:%1878=883%:%
%:%1879=884%:%
%:%1888=893%:%
%:%1889=894%:%
%:%1890=895%:%
%:%1894=897%:%
%:%1896=899%:%
%:%1897=899%:%
%:%1898=900%:%
%:%1899=901%:%
%:%1900=902%:%
%:%1903=903%:%
%:%1907=903%:%
%:%1908=903%:%
%:%1909=904%:%
%:%1919=906%:%
%:%1928=909%:%
%:%1940=911%:%
%:%1941=912%:%
%:%1943=914%:%
%:%1944=914%:%
%:%1945=914%:%
%:%1946=914%:%
%:%1947=915%:%
%:%1948=915%:%
%:%1949=915%:%
%:%1950=915%:%
%:%1951=916%:%
%:%1952=916%:%
%:%1953=916%:%
%:%1954=916%:%
%:%1955=917%:%
%:%1956=918%:%
%:%1957=918%:%
%:%1958=919%:%
%:%1959=920%:%
%:%1960=921%:%
%:%1961=922%:%
%:%1962=923%:%
%:%1963=924%:%
%:%1964=925%:%
%:%1965=926%:%
%:%1966=927%:%
%:%1967=928%:%
%:%1974=931%:%
%:%1984=933%:%
%:%1985=933%:%
%:%1986=933%:%
%:%1987=933%:%
%:%1988=934%:%
%:%1989=935%:%
%:%1990=935%:%
%:%1991=936%:%
%:%1992=937%:%
%:%1993=938%:%
%:%1994=939%:%
%:%1996=941%:%
%:%1997=942%:%
%:%1999=944%:%
%:%2000=945%:%
%:%2001=946%:%
%:%2002=947%:%
%:%2003=948%:%
%:%2004=949%:%
%:%2012=952%:%
%:%2024=954%:%
%:%2025=955%:%
%:%2026=956%:%
%:%2028=958%:%
%:%2032=959%:%
%:%2033=960%:%
%:%2035=961%:%
%:%2036=962%:%
%:%2038=963%:%
%:%2039=964%:%
%:%2041=965%:%
%:%2042=966%:%
%:%2047=970%:%
%:%2048=971%:%
%:%2049=972%:%

%
\begin{isabellebody}%
\setisabellecontext{CIL}%
%
\isadelimdocument
%
\endisadelimdocument
%
\isatagdocument
%
\isamarkupchapter{Appendix: A Syntactic Frontend for \isa{IMP\isactrlsub c\isactrlsub o\isactrlsub n\isactrlsub c\isactrlsub u\isactrlsub r}%
}
\isamarkuptrue%
%
\endisatagdocument
{\isafolddocument}%
%
\isadelimdocument
%
\endisadelimdocument
%
\isadelimtheory
%
\endisadelimtheory
%
\isatagtheory
\isakeywordONE{theory}\isamarkupfalse%
\ CIL\isanewline
\ \ \isakeywordTWO{imports}\ \ \ {\isachardoublequoteopen}IMP{\isacharunderscore}{\kern0pt}concur{\isachardoublequoteclose}\ \isanewline
\ \ \isakeywordTWO{keywords}\ \ {\isachardoublequoteopen}SYSTEM{\isachardoublequoteclose}\ {\isacharcolon}{\kern0pt}{\isacharcolon}{\kern0pt}\ thy{\isacharunderscore}{\kern0pt}decl\isanewline
\ \ \ \ \ \ \isakeywordTWO{and}\ \ \ {\isachardoublequoteopen}globals{\isachardoublequoteclose}\ {\isachardoublequoteopen}locks{\isachardoublequoteclose}\ {\isachardoublequoteopen}LOCK{\isachardoublequoteclose}\ {\isachardoublequoteopen}UNLOCK{\isachardoublequoteclose}\ {\isachardoublequoteopen}thread{\isachardoublequoteclose}\ {\isachardoublequoteopen}WHILE{\isachardoublequoteclose}\ {\isachardoublequoteopen}IF{\isachardoublequoteclose}\ {\isachardoublequoteopen}ELSE{\isachardoublequoteclose}\ {\isachardoublequoteopen}THEN{\isachardoublequoteclose}\isanewline
\ \ \ \ \ \ \ \ \ \ \ \ {\isachardoublequoteopen}end{\isacharsemicolon}{\kern0pt}{\isachardoublequoteclose}\ {\isachardoublequoteopen}any{\isachardoublequoteclose}\ {\isachardoublequoteopen}actions{\isachardoublequoteclose}\ {\isachardoublequoteopen}DO{\isachardoublequoteclose}\ {\isachardoublequoteopen}{\isacharless}{\kern0pt}{\isacharminus}{\kern0pt}{\isachardoublequoteclose}\ {\isachardoublequoteopen}{\isacharminus}{\kern0pt}{\isachargreater}{\kern0pt}{\isachardoublequoteclose}\ {\isachardoublequoteopen}SKIP{\isachardoublequoteclose}\ {\isachardoublequoteopen}DONE{\isachardoublequoteclose}\isanewline
\ \isanewline
\isakeywordTWO{begin}%
\endisatagtheory
{\isafoldtheory}%
%
\isadelimtheory
%
\endisadelimtheory
%
\isadelimdocument
%
\endisadelimdocument
%
\isatagdocument
%
\isamarkupsection{Introduction%
}
\isamarkuptrue%
%
\endisatagdocument
{\isafolddocument}%
%
\isadelimdocument
%
\endisadelimdocument
%
\begin{isamarkuptext}%
This theory defines syntactic support for a Concurrent Imperative Language (CIL) on top of
\isa{IMP\isactrlsub c\isactrlsub o\isactrlsub n\isactrlsub c\isactrlsub u\isactrlsub r}, a semantic language with threads, thread-local and thread-global states on which 
the threads may synchronize via semaphores (locks and unlocks). \isa{IMP\isactrlsub c\isactrlsub o\isactrlsub n\isactrlsub c\isactrlsub u\isactrlsub r} itself is based on 
HOL-CSPM; the semantics is constructed in a way that a conversion to Proc-Omata is possible.%
\end{isamarkuptext}\isamarkuptrue%
%
\begin{isamarkuptext}%
This appendix has been authored by Mathilde Needham and Zineddine Kermadj on their 
internship in summer 2025.%
\end{isamarkuptext}\isamarkuptrue%
%
\isadelimdocument
%
\endisadelimdocument
%
\isatagdocument
%
\isamarkupsection{Interface%
}
\isamarkuptrue%
%
\endisatagdocument
{\isafolddocument}%
%
\isadelimdocument
%
\endisadelimdocument
\isakeywordONE{definition}\isamarkupfalse%
\ SEMAPHORES\ \isanewline
\ \ \isakeywordTWO{where}\ \ \ {\isacartoucheopen}SEMAPHORES\ idx\ {\isasymequiv}\ Renaming\ {\isacharparenleft}{\kern0pt}semaphore\ idx{\isacharparenright}{\kern0pt}\ id\ id{\isacartoucheclose}\isanewline
\ \ \ \ \ \ \ \ \ %
\isamarkupcmt{lifting semaphores to the global alphabet%
}\ \ \ \isanewline
\ \isanewline
\isakeywordONE{definition}\isamarkupfalse%
\ GLOBALVARS\ {\isacharcolon}{\kern0pt}{\isacharcolon}{\kern0pt}\ \ {\isachardoublequoteopen}{\isasymsigma}\ {\isasymRightarrow}\ string\ {\isasymRightarrow}\ evs\ process{\isachardoublequoteclose}\isanewline
\ \ \isakeywordTWO{where}\ \ \ {\isacartoucheopen}GLOBALVARS\ {\isasymsigma}\isactrlsub i\isactrlsub n\isactrlsub i\isactrlsub t\ idx\ {\isasymequiv}\ Renaming\ {\isacharparenleft}{\kern0pt}global{\isacharunderscore}{\kern0pt}vars\ {\isasymsigma}\isactrlsub i\isactrlsub n\isactrlsub i\isactrlsub t\ idx{\isacharparenright}{\kern0pt}\ {\isacharparenleft}{\kern0pt}{\isasymlambda}x{\isachardot}{\kern0pt}\ x{\isacharparenright}{\kern0pt}\ {\isacharparenleft}{\kern0pt}{\isasymlambda}y{\isachardot}{\kern0pt}\ y{\isacharparenright}{\kern0pt}{\isacartoucheclose}\isanewline
\isanewline
\isakeywordONE{definition}\isamarkupfalse%
\ LOCALVARS\ \isanewline
\ \ \isakeywordTWO{where}\ \ \ {\isacartoucheopen}LOCALVARS\ \ {\isasymsigma}\isactrlsub i\isactrlsub n\isactrlsub i\isactrlsub t\ idx\ {\isasymequiv}\ Renaming\ {\isacharparenleft}{\kern0pt}local{\isacharunderscore}{\kern0pt}vars\ \ {\isasymsigma}\isactrlsub i\isactrlsub n\isactrlsub i\isactrlsub t\ idx{\isacharparenright}{\kern0pt}{\isacharparenleft}{\kern0pt}{\isasymlambda}x{\isachardot}{\kern0pt}\ x{\isacharparenright}{\kern0pt}\ {\isacharparenleft}{\kern0pt}{\isasymlambda}y{\isachardot}{\kern0pt}\ y{\isacharparenright}{\kern0pt}{\isacartoucheclose}%
\isadelimdocument
%
\endisadelimdocument
%
\isatagdocument
%
\isamarkupsection{Parsing and Term-Reading%
}
\isamarkuptrue%
%
\endisatagdocument
{\isafolddocument}%
%
\isadelimdocument
%
\endisadelimdocument
%
\isadelimML
%
\endisadelimML
%
\isatagML
\isakeywordONE{ML}\isamarkupfalse%
{\isacartoucheopen}\ \isanewline
\isanewline
val\ quiet{\isacharunderscore}{\kern0pt}mode\ {\isacharequal}{\kern0pt}\ false\isanewline
\isanewline
fun\ message\ quiet{\isacharunderscore}{\kern0pt}mode\ s\ {\isacharequal}{\kern0pt}\ if\ quiet{\isacharunderscore}{\kern0pt}mode\ then\ {\isacharparenleft}{\kern0pt}{\isacharparenright}{\kern0pt}\ else\ writeln\ s{\isacharsemicolon}{\kern0pt}\isanewline
\isanewline
datatype\ {\isacharparenleft}{\kern0pt}{\isacharprime}{\kern0pt}a{\isacharcomma}{\kern0pt}\ {\isacharprime}{\kern0pt}b{\isacharparenright}{\kern0pt}\ sum\ {\isacharequal}{\kern0pt}\ inl\ of\ {\isacharprime}{\kern0pt}a\ {\isacharbar}{\kern0pt}\ inr\ of\ {\isacharprime}{\kern0pt}b{\isacharsemicolon}{\kern0pt}\isanewline
\isanewline
exception\ Undeclared{\isacharunderscore}{\kern0pt}Variable\ of\ string\isanewline
\isanewline
local\ open\ HOLogic\ in\isanewline
\isanewline
fun\ map{\isacharunderscore}{\kern0pt}opt\ f\ {\isacharunderscore}{\kern0pt}\ {\isacharparenleft}{\kern0pt}SOME\ a{\isacharparenright}{\kern0pt}\ {\isacharequal}{\kern0pt}\ f\ a\isanewline
\ \ \ {\isacharbar}{\kern0pt}map{\isacharunderscore}{\kern0pt}opt\ {\isacharunderscore}{\kern0pt}\ a\ NONE\ {\isacharequal}{\kern0pt}\ a\isanewline
\isanewline
fun\ map{\isacharunderscore}{\kern0pt}option\ f\ \ {\isacharequal}{\kern0pt}\ map{\isacharunderscore}{\kern0pt}opt\ {\isacharparenleft}{\kern0pt}SOME\ o\ f{\isacharparenright}{\kern0pt}\ NONE\isanewline
\isanewline
{\isacharparenleft}{\kern0pt}{\isacharasterisk}{\kern0pt}\ remove\ this\ {\isacharasterisk}{\kern0pt}{\isacharparenright}{\kern0pt}\isanewline
fun\ map{\isacharunderscore}{\kern0pt}option\ f\ {\isacharparenleft}{\kern0pt}SOME\ a{\isacharparenright}{\kern0pt}\ {\isacharequal}{\kern0pt}\ SOME{\isacharparenleft}{\kern0pt}f\ a{\isacharparenright}{\kern0pt}\isanewline
\ \ \ {\isacharbar}{\kern0pt}map{\isacharunderscore}{\kern0pt}option\ {\isacharunderscore}{\kern0pt}\ NONE\ {\isacharequal}{\kern0pt}\ NONE\isanewline
\isanewline
\isanewline
\isanewline
val\ parse{\isacharunderscore}{\kern0pt}var{\isacharunderscore}{\kern0pt}decl\ {\isacharequal}{\kern0pt}\ {\isacharparenleft}{\kern0pt}Parse{\isachardot}{\kern0pt}binding\ {\isacharminus}{\kern0pt}{\isacharminus}{\kern0pt}{\isacharbar}{\kern0pt}\ Parse{\isachardot}{\kern0pt}{\isachardollar}{\kern0pt}{\isachardollar}{\kern0pt}{\isachardollar}{\kern0pt}\ {\isachardoublequote}{\kern0pt}{\isacharcolon}{\kern0pt}{\isachardoublequote}{\kern0pt}\ {\isacharminus}{\kern0pt}{\isacharminus}{\kern0pt}\ {\isacharparenleft}{\kern0pt}Parse{\isachardot}{\kern0pt}position\ {\isacharparenleft}{\kern0pt}Parse{\isachardot}{\kern0pt}term{\isacharparenright}{\kern0pt}{\isacharparenright}{\kern0pt}{\isacharparenright}{\kern0pt}\isanewline
\ \ \ \ \ \ \ \ \ \ \ \ \ \ \ \ \ \ \ {\isacharminus}{\kern0pt}{\isacharminus}{\kern0pt}\ {\isacharparenleft}{\kern0pt}Scan{\isachardot}{\kern0pt}option\ {\isacharparenleft}{\kern0pt}\ Parse{\isachardot}{\kern0pt}{\isachardollar}{\kern0pt}{\isachardollar}{\kern0pt}{\isachardollar}{\kern0pt}\ {\isachardoublequote}{\kern0pt}{\isacharequal}{\kern0pt}{\isachardoublequote}{\kern0pt}\ {\isacharbar}{\kern0pt}{\isacharminus}{\kern0pt}{\isacharminus}{\kern0pt}\ Parse{\isachardot}{\kern0pt}term{\isacharparenright}{\kern0pt}{\isacharparenright}{\kern0pt}\isanewline
\ \ \ \ \ \ \ \ \ \ \ \ \ \ \ \ \ \ \ \ {\isacharcolon}{\kern0pt}\ {\isacharparenleft}{\kern0pt}{\isacharparenleft}{\kern0pt}binding{\isacharasterisk}{\kern0pt}{\isacharparenleft}{\kern0pt}string{\isacharasterisk}{\kern0pt}Position{\isachardot}{\kern0pt}T{\isacharparenright}{\kern0pt}{\isacharparenright}{\kern0pt}{\isacharasterisk}{\kern0pt}string\ option{\isacharparenright}{\kern0pt}\ parser\isanewline
\isanewline
{\isacharparenleft}{\kern0pt}{\isacharasterisk}{\kern0pt}declarations\ de\ locks{\isacharasterisk}{\kern0pt}{\isacharparenright}{\kern0pt}\isanewline
val\ parse{\isacharunderscore}{\kern0pt}locks\ {\isacharequal}{\kern0pt}\ Parse{\isachardot}{\kern0pt}binding\ {\isacharminus}{\kern0pt}{\isacharminus}{\kern0pt}{\isacharbar}{\kern0pt}\ Parse{\isachardot}{\kern0pt}{\isachardollar}{\kern0pt}{\isachardollar}{\kern0pt}{\isachardollar}{\kern0pt}\ {\isachardoublequote}{\kern0pt}{\isacharcolon}{\kern0pt}{\isachardoublequote}{\kern0pt}\ {\isacharminus}{\kern0pt}{\isacharminus}{\kern0pt}\ {\isacharparenleft}{\kern0pt}Parse{\isachardot}{\kern0pt}position\ {\isacharparenleft}{\kern0pt}Parse{\isachardot}{\kern0pt}term{\isacharparenright}{\kern0pt}{\isacharparenright}{\kern0pt}\ \isanewline
\ \ \ \ \ \ \ \ \ \ \ \ \ \ \ \ \ \ {\isacharcolon}{\kern0pt}\ {\isacharparenleft}{\kern0pt}binding{\isacharasterisk}{\kern0pt}{\isacharparenleft}{\kern0pt}string{\isacharasterisk}{\kern0pt}Position{\isachardot}{\kern0pt}T{\isacharparenright}{\kern0pt}{\isacharparenright}{\kern0pt}\ parser\isanewline
\isanewline
\isanewline
datatype\ a\ {\isacharequal}{\kern0pt}\ Skip\isanewline
\ \ \ \ \ \ \ \ \ \ \ \ {\isacharbar}{\kern0pt}\ Assign\ of\ binding{\isacharasterisk}{\kern0pt}string\ \isanewline
\ \ \ \ \ \ \ \ \ \ \ \ {\isacharbar}{\kern0pt}\ Unlock\ of\ binding\ \isanewline
\ \ \ \ \ \ \ \ \ \ \ \ {\isacharbar}{\kern0pt}\ Lock\ of\ binding\isanewline
\ \ \ \ \ \ \ \ \ \ \ \ {\isacharbar}{\kern0pt}\ Send\ of\ binding{\isacharasterisk}{\kern0pt}string\ \isanewline
\ \ \ \ \ \ \ \ \ \ \ \ {\isacharbar}{\kern0pt}\ Receive\ of\ binding{\isacharasterisk}{\kern0pt}binding\isanewline
\ \ \ \ \ \ \ \ \ \ \ \ {\isacharbar}{\kern0pt}\ Ifelse\ of\ {\isacharparenleft}{\kern0pt}string{\isacharasterisk}{\kern0pt}{\isacharparenleft}{\kern0pt}a\ list{\isacharparenright}{\kern0pt}{\isacharparenright}{\kern0pt}{\isacharasterisk}{\kern0pt}{\isacharparenleft}{\kern0pt}a\ list{\isacharparenright}{\kern0pt}\ \isanewline
\ \ \ \ \ \ \ \ \ \ \ \ {\isacharbar}{\kern0pt}\ While\ of\ string{\isacharasterisk}{\kern0pt}{\isacharparenleft}{\kern0pt}a\ list{\isacharparenright}{\kern0pt}\isanewline
\isanewline
val\ parse{\isacharunderscore}{\kern0pt}lock\ {\isacharequal}{\kern0pt}\ \isactrlkeyword {\isasymopen}LOCK{\isasymclose}\ {\isacharbar}{\kern0pt}{\isacharminus}{\kern0pt}{\isacharminus}{\kern0pt}\ Parse{\isachardot}{\kern0pt}binding\ {\isacharminus}{\kern0pt}{\isacharminus}{\kern0pt}{\isacharbar}{\kern0pt}\ Parse{\isachardot}{\kern0pt}{\isachardollar}{\kern0pt}{\isachardollar}{\kern0pt}{\isachardollar}{\kern0pt}\ {\isachardoublequote}{\kern0pt}{\isacharsemicolon}{\kern0pt}{\isachardoublequote}{\kern0pt}\isanewline
\ \ \ \ \ \ \ \ \ \ \ \ \ \ \ \ \ {\isachargreater}{\kern0pt}{\isachargreater}{\kern0pt}\ {\isacharparenleft}{\kern0pt}fn\ {\isacharparenleft}{\kern0pt}x{\isacharparenright}{\kern0pt}\ {\isacharequal}{\kern0pt}{\isachargreater}{\kern0pt}\ Lock{\isacharparenleft}{\kern0pt}x{\isacharparenright}{\kern0pt}{\isacharparenright}{\kern0pt}\isanewline
\ \ \ \ \ \ \ \ \ \ \ \ \ \ \ \ \ \ {\isacharcolon}{\kern0pt}\ a\ parser\isanewline
val\ parse{\isacharunderscore}{\kern0pt}unlock\ {\isacharequal}{\kern0pt}\ \ \isactrlkeyword {\isasymopen}UNLOCK{\isasymclose}\ {\isacharbar}{\kern0pt}{\isacharminus}{\kern0pt}{\isacharminus}{\kern0pt}\ Parse{\isachardot}{\kern0pt}binding\ {\isacharminus}{\kern0pt}{\isacharminus}{\kern0pt}{\isacharbar}{\kern0pt}\ Parse{\isachardot}{\kern0pt}{\isachardollar}{\kern0pt}{\isachardollar}{\kern0pt}{\isachardollar}{\kern0pt}\ {\isachardoublequote}{\kern0pt}{\isacharsemicolon}{\kern0pt}{\isachardoublequote}{\kern0pt}\isanewline
\ \ \ \ \ \ \ \ \ \ \ \ \ \ \ \ \ {\isachargreater}{\kern0pt}{\isachargreater}{\kern0pt}\ \ {\isacharparenleft}{\kern0pt}fn\ {\isacharparenleft}{\kern0pt}x{\isacharparenright}{\kern0pt}\ {\isacharequal}{\kern0pt}{\isachargreater}{\kern0pt}\ Unlock{\isacharparenleft}{\kern0pt}x{\isacharparenright}{\kern0pt}{\isacharparenright}{\kern0pt}\isanewline
\ \ \ \ \ \ \ \ \ \ \ \ \ \ \ \ \ \ {\isacharcolon}{\kern0pt}\ a\ parser\isanewline
val\ parse{\isacharunderscore}{\kern0pt}assign\ {\isacharequal}{\kern0pt}\ Parse{\isachardot}{\kern0pt}binding\ {\isacharminus}{\kern0pt}{\isacharminus}{\kern0pt}{\isacharbar}{\kern0pt}\ Parse{\isachardot}{\kern0pt}{\isachardollar}{\kern0pt}{\isachardollar}{\kern0pt}{\isachardollar}{\kern0pt}\ {\isachardoublequote}{\kern0pt}{\isacharequal}{\kern0pt}{\isachardoublequote}{\kern0pt}\ {\isacharminus}{\kern0pt}{\isacharminus}{\kern0pt}\ Parse{\isachardot}{\kern0pt}term\ {\isacharminus}{\kern0pt}{\isacharminus}{\kern0pt}{\isacharbar}{\kern0pt}\ Parse{\isachardot}{\kern0pt}{\isachardollar}{\kern0pt}{\isachardollar}{\kern0pt}{\isachardollar}{\kern0pt}\ {\isachardoublequote}{\kern0pt}{\isacharsemicolon}{\kern0pt}{\isachardoublequote}{\kern0pt}\isanewline
\ \ \ \ \ \ \ \ \ \ \ \ \ \ \ \ {\isachargreater}{\kern0pt}{\isachargreater}{\kern0pt}\ {\isacharparenleft}{\kern0pt}fn\ {\isacharparenleft}{\kern0pt}x{\isacharparenright}{\kern0pt}\ {\isacharequal}{\kern0pt}{\isachargreater}{\kern0pt}\ Assign{\isacharparenleft}{\kern0pt}x{\isacharparenright}{\kern0pt}{\isacharparenright}{\kern0pt}\isanewline
\ \ \ \ \ \ \ \ \ \ \ \ \ \ \ \ \ \ {\isacharcolon}{\kern0pt}\ a\ parser\isanewline
val\ parse{\isacharunderscore}{\kern0pt}send\ {\isacharequal}{\kern0pt}\ Parse{\isachardot}{\kern0pt}binding\ {\isacharminus}{\kern0pt}{\isacharminus}{\kern0pt}{\isacharbar}{\kern0pt}\ \isactrlkeyword {\isasymopen}{\isacharminus}{\kern0pt}{\isachargreater}{\kern0pt}{\isasymclose}\ {\isacharminus}{\kern0pt}{\isacharminus}{\kern0pt}\ Parse{\isachardot}{\kern0pt}term\ {\isacharminus}{\kern0pt}{\isacharminus}{\kern0pt}{\isacharbar}{\kern0pt}\ Parse{\isachardot}{\kern0pt}{\isachardollar}{\kern0pt}{\isachardollar}{\kern0pt}{\isachardollar}{\kern0pt}\ {\isachardoublequote}{\kern0pt}{\isacharsemicolon}{\kern0pt}{\isachardoublequote}{\kern0pt}\isanewline
\ \ \ \ \ \ \ \ \ \ \ \ \ \ \ \ \ {\isachargreater}{\kern0pt}{\isachargreater}{\kern0pt}\ {\isacharparenleft}{\kern0pt}fn\ {\isacharparenleft}{\kern0pt}x{\isacharparenright}{\kern0pt}\ {\isacharequal}{\kern0pt}{\isachargreater}{\kern0pt}\ Send{\isacharparenleft}{\kern0pt}x{\isacharparenright}{\kern0pt}{\isacharparenright}{\kern0pt}\isanewline
\ \ \ \ \ \ \ \ \ \ \ \ \ \ \ \ \ \ {\isacharcolon}{\kern0pt}\ a\ parser\isanewline
val\ parse{\isacharunderscore}{\kern0pt}receive\ {\isacharequal}{\kern0pt}\ Parse{\isachardot}{\kern0pt}binding\ {\isacharminus}{\kern0pt}{\isacharminus}{\kern0pt}{\isacharbar}{\kern0pt}\ \isactrlkeyword {\isasymopen}{\isacharless}{\kern0pt}{\isacharminus}{\kern0pt}{\isasymclose}\ {\isacharminus}{\kern0pt}{\isacharminus}{\kern0pt}\ Parse{\isachardot}{\kern0pt}binding\ {\isacharminus}{\kern0pt}{\isacharminus}{\kern0pt}{\isacharbar}{\kern0pt}\ Parse{\isachardot}{\kern0pt}{\isachardollar}{\kern0pt}{\isachardollar}{\kern0pt}{\isachardollar}{\kern0pt}\ {\isachardoublequote}{\kern0pt}{\isacharsemicolon}{\kern0pt}{\isachardoublequote}{\kern0pt}\isanewline
\ \ \ \ \ \ \ \ \ \ \ \ \ \ \ \ {\isachargreater}{\kern0pt}{\isachargreater}{\kern0pt}\ {\isacharparenleft}{\kern0pt}fn\ {\isacharparenleft}{\kern0pt}x{\isacharparenright}{\kern0pt}\ {\isacharequal}{\kern0pt}{\isachargreater}{\kern0pt}\ Receive{\isacharparenleft}{\kern0pt}x{\isacharparenright}{\kern0pt}{\isacharparenright}{\kern0pt}\isanewline
\ \ \ \ \ \ \ \ \ \ \ \ \ \ \ \ \ \ {\isacharcolon}{\kern0pt}\ a\ parser\isanewline
val\ parse{\isacharunderscore}{\kern0pt}skip\ {\isacharequal}{\kern0pt}\ \isactrlkeyword {\isasymopen}SKIP{\isasymclose}\ {\isacharbar}{\kern0pt}{\isacharminus}{\kern0pt}{\isacharminus}{\kern0pt}\ Parse{\isachardot}{\kern0pt}{\isachardollar}{\kern0pt}{\isachardollar}{\kern0pt}{\isachardollar}{\kern0pt}\ {\isachardoublequote}{\kern0pt}{\isacharsemicolon}{\kern0pt}{\isachardoublequote}{\kern0pt}\isanewline
\ \ \ \ \ \ \ \ \ \ \ \ \ \ \ \ {\isachargreater}{\kern0pt}{\isachargreater}{\kern0pt}\ {\isacharparenleft}{\kern0pt}fn\ {\isacharunderscore}{\kern0pt}\ \ {\isacharequal}{\kern0pt}{\isachargreater}{\kern0pt}\ Skip{\isacharparenright}{\kern0pt}\ {\isacharcolon}{\kern0pt}\ a\ parser\isanewline
\isanewline
fun\ parse{\isacharunderscore}{\kern0pt}actions\ TL\ \ \ {\isacharequal}{\kern0pt}{\isacharparenleft}{\kern0pt}\ Scan{\isachardot}{\kern0pt}repeat\ {\isacharparenleft}{\kern0pt}\ parse{\isacharunderscore}{\kern0pt}lock\ \isanewline
\ \ \ \ \ \ \ \ \ \ \ \ \ \ \ \ \ \ \ \ \ \ \ \ \ \ \ \ \ \ \ \ \ \ {\isacharbar}{\kern0pt}{\isacharbar}{\kern0pt}\ parse{\isacharunderscore}{\kern0pt}unlock\ \isanewline
\ \ \ \ \ \ \ \ \ \ \ \ \ \ \ \ \ \ \ \ \ \ \ \ \ \ \ \ \ \ \ \ \ \ {\isacharbar}{\kern0pt}{\isacharbar}{\kern0pt}\ parse{\isacharunderscore}{\kern0pt}assign\isanewline
\ \ \ \ \ \ \ \ \ \ \ \ \ \ \ \ \ \ \ \ \ \ \ \ \ \ \ \ \ \ \ \ \ \ {\isacharbar}{\kern0pt}{\isacharbar}{\kern0pt}\ parse{\isacharunderscore}{\kern0pt}send\ \isanewline
\ \ \ \ \ \ \ \ \ \ \ \ \ \ \ \ \ \ \ \ \ \ \ \ \ \ \ \ \ \ \ \ \ \ {\isacharbar}{\kern0pt}{\isacharbar}{\kern0pt}\ parse{\isacharunderscore}{\kern0pt}receive\isanewline
\ \ \ \ \ \ \ \ \ \ \ \ \ \ \ \ \ \ \ \ \ \ \ \ \ \ \ \ \ \ \ \ \ \ {\isacharbar}{\kern0pt}{\isacharbar}{\kern0pt}\ parse{\isacharunderscore}{\kern0pt}skip\isanewline
\ \ \ \ \ \ \ \ \ \ \ \ \ \ \ \ \ \ \ \ \ \ \ \ \ \ \ \ \ \ \ \ \ \ {\isacharbar}{\kern0pt}{\isacharbar}{\kern0pt}\ parse{\isacharunderscore}{\kern0pt}if{\isacharunderscore}{\kern0pt}else\ \isanewline
\ \ \ \ \ \ \ \ \ \ \ \ \ \ \ \ \ \ \ \ \ \ \ \ \ \ \ \ \ \ \ \ \ \ {\isacharbar}{\kern0pt}{\isacharbar}{\kern0pt}\ parse{\isacharunderscore}{\kern0pt}while\ \isanewline
\ \ \ \ \ \ \ \ \ \ \ \ \ \ \ \ \ \ \ \ \ \ \ \ \ \ \ \ \ \ \ \ \ \ {\isacharparenright}{\kern0pt}{\isacharparenright}{\kern0pt}\ TL\isanewline
and\ \ parse{\isacharunderscore}{\kern0pt}if{\isacharunderscore}{\kern0pt}else\ TL\ {\isacharequal}{\kern0pt}\ {\isacharparenleft}{\kern0pt}{\isacharparenleft}{\kern0pt}\isactrlkeyword {\isasymopen}IF{\isasymclose}\ {\isacharbar}{\kern0pt}{\isacharminus}{\kern0pt}{\isacharminus}{\kern0pt}\ Parse{\isachardot}{\kern0pt}term\ {\isacharminus}{\kern0pt}{\isacharminus}{\kern0pt}{\isacharbar}{\kern0pt}\ \isactrlkeyword {\isasymopen}THEN{\isasymclose}\isanewline
\ \ \ \ \ \ \ \ \ \ \ \ \ \ \ \ \ \ \ \ {\isacharminus}{\kern0pt}{\isacharminus}{\kern0pt}\ parse{\isacharunderscore}{\kern0pt}actions\ {\isacharminus}{\kern0pt}{\isacharminus}{\kern0pt}{\isacharbar}{\kern0pt}\ \isactrlkeyword {\isasymopen}ELSE{\isasymclose}\ {\isacharminus}{\kern0pt}{\isacharminus}{\kern0pt}\ parse{\isacharunderscore}{\kern0pt}actions{\isacharminus}{\kern0pt}{\isacharminus}{\kern0pt}{\isacharbar}{\kern0pt}\ \isactrlkeyword {\isasymopen}DONE{\isasymclose}{\isacharparenright}{\kern0pt}\isanewline
\ \ \ \ \ \ \ \ \ \ \ \ \ \ \ \ \ \ \ \ {\isachargreater}{\kern0pt}{\isachargreater}{\kern0pt}{\isacharparenleft}{\kern0pt}fn\ {\isacharparenleft}{\kern0pt}x{\isacharparenright}{\kern0pt}\ {\isacharequal}{\kern0pt}{\isachargreater}{\kern0pt}\ Ifelse{\isacharparenleft}{\kern0pt}x{\isacharparenright}{\kern0pt}{\isacharparenright}{\kern0pt}{\isacharparenright}{\kern0pt}\ TL\isanewline
and\ parse{\isacharunderscore}{\kern0pt}while\ TL\ {\isacharequal}{\kern0pt}{\isacharparenleft}{\kern0pt}{\isacharparenleft}{\kern0pt}\ \isactrlkeyword {\isasymopen}WHILE{\isasymclose}\ {\isacharbar}{\kern0pt}{\isacharminus}{\kern0pt}{\isacharminus}{\kern0pt}\ Parse{\isachardot}{\kern0pt}term\ {\isacharminus}{\kern0pt}{\isacharminus}{\kern0pt}{\isacharbar}{\kern0pt}\ \isactrlkeyword {\isasymopen}DO{\isasymclose}\ {\isacharminus}{\kern0pt}{\isacharminus}{\kern0pt}\ parse{\isacharunderscore}{\kern0pt}actions\ {\isacharminus}{\kern0pt}{\isacharminus}{\kern0pt}{\isacharbar}{\kern0pt}\isactrlkeyword {\isasymopen}DONE{\isasymclose}{\isacharparenright}{\kern0pt}\isanewline
\ \ \ \ \ \ \ \ \ \ \ \ \ \ \ \ \ \ \ \ {\isachargreater}{\kern0pt}{\isachargreater}{\kern0pt}\ {\isacharparenleft}{\kern0pt}fn\ {\isacharparenleft}{\kern0pt}x{\isacharparenright}{\kern0pt}\ {\isacharequal}{\kern0pt}{\isachargreater}{\kern0pt}\ While{\isacharparenleft}{\kern0pt}x{\isacharparenright}{\kern0pt}{\isacharparenright}{\kern0pt}{\isacharparenright}{\kern0pt}\ TL\isanewline
\isanewline
\isanewline
\ \ \ \ \ \ \ \ \ \ \ \ \ \ \ \ \ \ \ \ \ \ \ \ \ \ \ \ \ \ \ \ \isanewline
\isanewline
{\isacharparenleft}{\kern0pt}{\isacharasterisk}{\kern0pt}thread\ T{\isadigit{1}}\ {\isacharcolon}{\kern0pt}\ where\ var{\isacharunderscore}{\kern0pt}local{\isacharunderscore}{\kern0pt}decl\isanewline
\ \ \ \ \ \ \ \ \ \ \ \ \ actions\ \ action{\isadigit{1}}{\isacharsemicolon}{\kern0pt}\isanewline
\ \ \ \ \ \ \ \ \ \ \ \ \ \ \ \ \ \ \ \ \ \ \ {\isachardot}{\kern0pt}{\isachardot}{\kern0pt}{\isachardot}{\kern0pt}\ {\isacharsemicolon}{\kern0pt}\isanewline
\ \ \ \ \ \ \ \ \ \ \ \ \ \ \ \ \ \ \ \ \ \ actionN{\isacharsemicolon}{\kern0pt}\isanewline
\ \ \ thread\ T{\isadigit{2}}{\isacharcolon}{\kern0pt}\ actions\ \ action{\isadigit{1}}{\isacharsemicolon}{\kern0pt}\isanewline
\ \ \ \ \ \ \ \ \ \ \ \ \ \ \ \ \ \ \ \ \ \ \ {\isachardot}{\kern0pt}{\isachardot}{\kern0pt}{\isachardot}{\kern0pt}\ {\isacharsemicolon}{\kern0pt}\isanewline
\ \ \ \ \ \ \ \ \ \ \ \ \ \ \ \ \ \ \ \ \ \ actionN{\isacharsemicolon}{\kern0pt}\isanewline
\ {\isacharasterisk}{\kern0pt}{\isacharparenright}{\kern0pt}\isanewline
val\ parse{\isacharunderscore}{\kern0pt}threads\ {\isacharequal}{\kern0pt}{\isacharparenleft}{\kern0pt}\ \isactrlkeyword {\isasymopen}thread{\isasymclose}\ {\isacharbar}{\kern0pt}{\isacharminus}{\kern0pt}{\isacharminus}{\kern0pt}\ {\isacharparenleft}{\kern0pt}Scan{\isachardot}{\kern0pt}option\ {\isacharparenleft}{\kern0pt}Parse{\isachardot}{\kern0pt}binding{\isacharparenright}{\kern0pt}{\isacharparenright}{\kern0pt}\ {\isacharminus}{\kern0pt}{\isacharminus}{\kern0pt}{\isacharbar}{\kern0pt}\ Parse{\isachardot}{\kern0pt}{\isachardollar}{\kern0pt}{\isachardollar}{\kern0pt}{\isachardollar}{\kern0pt}\ {\isachardoublequote}{\kern0pt}{\isacharcolon}{\kern0pt}{\isachardoublequote}{\kern0pt}\isanewline
\ \ \ \ \ \ \ \ \ \ \ \ \ \ \ \ \ \ \ \ \ \ \ \ \ \ \ \ \ \ \ \ \ \ \ \ \ \ {\isacharminus}{\kern0pt}{\isacharminus}{\kern0pt}\ {\isacharparenleft}{\kern0pt}Scan{\isachardot}{\kern0pt}option\ {\isacharparenleft}{\kern0pt}\isactrlkeyword {\isasymopen}any{\isasymclose}\ {\isacharbar}{\kern0pt}{\isacharminus}{\kern0pt}{\isacharminus}{\kern0pt}\ \ Scan{\isachardot}{\kern0pt}repeat{\isadigit{1}}\ {\isacharparenleft}{\kern0pt}parse{\isacharunderscore}{\kern0pt}var{\isacharunderscore}{\kern0pt}decl{\isacharparenright}{\kern0pt}\ {\isacharparenright}{\kern0pt}{\isacharparenright}{\kern0pt}\isanewline
\ \ \ \ \ \ \ \ \ \ \ \ \ \ \ \ \ \ \ \ \ \ \ \ \ \ \ \ \ \ \ \ \ \ \ \ \ {\isacharminus}{\kern0pt}{\isacharminus}{\kern0pt}\ {\isacharparenleft}{\kern0pt}\isactrlkeyword {\isasymopen}actions{\isasymclose}\ {\isacharbar}{\kern0pt}{\isacharminus}{\kern0pt}{\isacharminus}{\kern0pt}\ parse{\isacharunderscore}{\kern0pt}actions{\isacharparenright}{\kern0pt}{\isacharparenright}{\kern0pt}\isanewline
\isanewline
type\ raw{\isacharunderscore}{\kern0pt}absy\ {\isacharequal}{\kern0pt}\ \ {\isacharparenleft}{\kern0pt}{\isacharparenleft}{\kern0pt}{\isacharparenleft}{\kern0pt}binding\ {\isacharparenleft}{\kern0pt}{\isacharasterisk}{\kern0pt}system\ name{\isacharasterisk}{\kern0pt}{\isacharparenright}{\kern0pt}\isanewline
\ \ \ \ \ \ \ \ \ \ \ \ {\isacharasterisk}{\kern0pt}\isanewline
\ \ \ \ \ \ \ \ \ \ \ \ \ {\isacharparenleft}{\kern0pt}{\isacharparenleft}{\kern0pt}binding\ {\isacharasterisk}{\kern0pt}\ {\isacharparenleft}{\kern0pt}string\ {\isacharasterisk}{\kern0pt}\ Position{\isachardot}{\kern0pt}T{\isacharparenright}{\kern0pt}{\isacharparenright}{\kern0pt}\ {\isacharasterisk}{\kern0pt}\ string\ option{\isacharparenright}{\kern0pt}list{\isacharparenright}{\kern0pt}\ {\isacharparenleft}{\kern0pt}{\isacharasterisk}{\kern0pt}globals{\isacharasterisk}{\kern0pt}{\isacharparenright}{\kern0pt}\isanewline
\ \ \ \ \ \ \ \ \ \ \ \ {\isacharasterisk}{\kern0pt}\isanewline
\ \ \ \ \ \ \ \ \ \ \ \ {\isacharparenleft}{\kern0pt}binding\ {\isacharasterisk}{\kern0pt}\ {\isacharparenleft}{\kern0pt}string\ {\isacharasterisk}{\kern0pt}\ Position{\isachardot}{\kern0pt}T{\isacharparenright}{\kern0pt}{\isacharparenright}{\kern0pt}\ list{\isacharparenright}{\kern0pt}\ {\isacharparenleft}{\kern0pt}{\isacharasterisk}{\kern0pt}locks{\isacharasterisk}{\kern0pt}{\isacharparenright}{\kern0pt}\isanewline
\ \ \ \ \ \ \ \ \ \ \ {\isacharasterisk}{\kern0pt}\isanewline
\ \ \ \ \ \ \ \ \ \ \ {\isacharparenleft}{\kern0pt}{\isacharparenleft}{\kern0pt}binding\ option\ {\isacharasterisk}{\kern0pt}\ {\isacharparenleft}{\kern0pt}{\isacharasterisk}{\kern0pt}a\ thread{\isacharprime}{\kern0pt}s\ name{\isacharasterisk}{\kern0pt}{\isacharparenright}{\kern0pt}\isanewline
\ \ \ \ \ \ \ \ \ \ \ \ \ {\isacharparenleft}{\kern0pt}{\isacharparenleft}{\kern0pt}binding\ {\isacharasterisk}{\kern0pt}\ {\isacharparenleft}{\kern0pt}string\ {\isacharasterisk}{\kern0pt}\ Position{\isachardot}{\kern0pt}T{\isacharparenright}{\kern0pt}{\isacharparenright}{\kern0pt}\ {\isacharasterisk}{\kern0pt}string\ option{\isacharparenright}{\kern0pt}list\ option{\isacharparenright}{\kern0pt}\ {\isacharparenleft}{\kern0pt}{\isacharasterisk}{\kern0pt}a\ thread{\isacharprime}{\kern0pt}s\ locals{\isacharasterisk}{\kern0pt}{\isacharparenright}{\kern0pt}\isanewline
\ \ \ \ \ \ \ \ \ \ \ \ {\isacharasterisk}{\kern0pt}\isanewline
\ \ \ \ \ \ \ \ \ \ \ \ \ \ a\ list\ {\isacharparenleft}{\kern0pt}{\isacharasterisk}{\kern0pt}a\ thread{\isacharprime}{\kern0pt}s\ actions{\isacharasterisk}{\kern0pt}{\isacharparenright}{\kern0pt}\isanewline
\ \ \ \ \ \ \ \ \ \ \ {\isacharparenright}{\kern0pt}\ list{\isacharparenright}{\kern0pt}\ {\isacharparenleft}{\kern0pt}{\isacharasterisk}{\kern0pt}threads\ lists{\isacharasterisk}{\kern0pt}{\isacharparenright}{\kern0pt}\isanewline
\isanewline
type\ raw{\isacharunderscore}{\kern0pt}thread{\isacharunderscore}{\kern0pt}absy\ {\isacharequal}{\kern0pt}\ {\isacharparenleft}{\kern0pt}binding\ option\ {\isacharasterisk}{\kern0pt}\ {\isacharparenleft}{\kern0pt}{\isacharasterisk}{\kern0pt}a\ thread{\isacharprime}{\kern0pt}s\ name{\isacharasterisk}{\kern0pt}{\isacharparenright}{\kern0pt}\isanewline
\ \ \ \ \ \ \ \ \ \ \ \ \ \ \ \ \ \ \ \ \ \ \ {\isacharparenleft}{\kern0pt}{\isacharparenleft}{\kern0pt}binding\ {\isacharasterisk}{\kern0pt}\ {\isacharparenleft}{\kern0pt}string\ {\isacharasterisk}{\kern0pt}\ Position{\isachardot}{\kern0pt}T{\isacharparenright}{\kern0pt}{\isacharparenright}{\kern0pt}\ {\isacharasterisk}{\kern0pt}string\ option{\isacharparenright}{\kern0pt}list\ option{\isacharparenright}{\kern0pt}{\isacharasterisk}{\kern0pt}\ {\isacharparenleft}{\kern0pt}{\isacharasterisk}{\kern0pt}a\ thread{\isacharprime}{\kern0pt}s\ locals{\isacharasterisk}{\kern0pt}{\isacharparenright}{\kern0pt}\isanewline
\ \ \ \ \ \ \ \ \ \ \ \ \ \ \ \ \ \ \ \ \ \ \ \ a\ list\ {\isacharparenleft}{\kern0pt}{\isacharasterisk}{\kern0pt}a\ thread{\isacharprime}{\kern0pt}s\ actions{\isacharasterisk}{\kern0pt}{\isacharparenright}{\kern0pt}\isanewline
\ \ \ \ \ \ \ \ \ \ \ \ \ \ \ \ \ \ \ \ \ \isanewline
\isanewline
val\ parse{\isacharunderscore}{\kern0pt}system{\isacharunderscore}{\kern0pt}spec\ {\isacharequal}{\kern0pt}\ {\isacharparenleft}{\kern0pt}\isanewline
\ \ \ \ \ \ \ \ \ \ Parse{\isachardot}{\kern0pt}binding\ \isanewline
\ \ \ \ \ \ \ {\isacharminus}{\kern0pt}{\isacharminus}{\kern0pt}{\isacharbar}{\kern0pt}\ \isactrlkeyword {\isasymopen}globals{\isasymclose}\ \ \ \ {\isacharminus}{\kern0pt}{\isacharminus}{\kern0pt}\ {\isacharparenleft}{\kern0pt}Scan{\isachardot}{\kern0pt}repeat\ parse{\isacharunderscore}{\kern0pt}var{\isacharunderscore}{\kern0pt}decl{\isacharparenright}{\kern0pt}\isanewline
\ \ \ \ \ \ \ {\isacharminus}{\kern0pt}{\isacharminus}{\kern0pt}{\isacharbar}{\kern0pt}\ \isactrlkeyword {\isasymopen}locks{\isasymclose}\ \ \ \ \ \ {\isacharminus}{\kern0pt}{\isacharminus}{\kern0pt}\ {\isacharparenleft}{\kern0pt}Scan{\isachardot}{\kern0pt}repeat\ parse{\isacharunderscore}{\kern0pt}locks{\isacharparenright}{\kern0pt}\isanewline
\ \ \ \ \ \ \ {\isacharminus}{\kern0pt}{\isacharminus}{\kern0pt}\ {\isacharparenleft}{\kern0pt}Scan{\isachardot}{\kern0pt}repeat{\isadigit{1}}\ parse{\isacharunderscore}{\kern0pt}threads{\isacharparenright}{\kern0pt}\isanewline
\ \ \ \ \ \ \ {\isacharminus}{\kern0pt}{\isacharminus}{\kern0pt}{\isacharbar}{\kern0pt}\ \isactrlkeyword {\isasymopen}end{\isacharsemicolon}{\kern0pt}{\isasymclose}\isanewline
\ \ \ \ \ \ {\isacharparenright}{\kern0pt}\ {\isacharcolon}{\kern0pt}\ raw{\isacharunderscore}{\kern0pt}absy\ parser\isanewline
\isanewline
end\isanewline
{\isacartoucheclose}%
\endisatagML
{\isafoldML}%
%
\isadelimML
%
\endisadelimML
\isanewline
\isanewline
\isanewline
\isakeywordONE{typedecl}\isamarkupfalse%
\ pre{\isacharunderscore}{\kern0pt}event{\isacharunderscore}{\kern0pt}type%
\isadelimdocument
%
\endisadelimdocument
%
\isatagdocument
%
\isamarkupsection{Context Checking%
}
\isamarkuptrue%
%
\endisatagdocument
{\isafolddocument}%
%
\isadelimdocument
%
\endisadelimdocument
%
\isadelimML
%
\endisadelimML
%
\isatagML
\isakeywordONE{ML}\isamarkupfalse%
{\isacartoucheopen}\isanewline
val\ pre{\isacharunderscore}{\kern0pt}event{\isacharunderscore}{\kern0pt}type\ {\isacharequal}{\kern0pt}\ {\isacharat}{\kern0pt}{\isacharbraceleft}{\kern0pt}typ\ {\isasymopen}pre{\isacharunderscore}{\kern0pt}event{\isacharunderscore}{\kern0pt}type{\isasymclose}{\isacharbraceright}{\kern0pt}\isanewline
\isanewline
datatype\ a{\isacharunderscore}{\kern0pt}term\ {\isacharequal}{\kern0pt}\ SkipA\isanewline
\ \ \ \ \ \ \ \ \ \ \ \ {\isacharbar}{\kern0pt}\ AssignA\ of\ binding{\isacharasterisk}{\kern0pt}term\ \isanewline
\ \ \ \ \ \ \ \ \ \ \ \ {\isacharbar}{\kern0pt}\ UnlockA\ of\ binding\ \isanewline
\ \ \ \ \ \ \ \ \ \ \ \ {\isacharbar}{\kern0pt}\ LockA\ of\ binding\isanewline
\ \ \ \ \ \ \ \ \ \ \ \ {\isacharbar}{\kern0pt}\ SendA\ of\ binding{\isacharasterisk}{\kern0pt}term\ \isanewline
\ \ \ \ \ \ \ \ \ \ \ \ {\isacharbar}{\kern0pt}\ ReceiveA\ of\ binding{\isacharasterisk}{\kern0pt}binding\isanewline
\ \ \ \ \ \ \ \ \ \ \ \ {\isacharbar}{\kern0pt}\ IfelseA\ of\ {\isacharparenleft}{\kern0pt}term{\isacharasterisk}{\kern0pt}{\isacharparenleft}{\kern0pt}a{\isacharunderscore}{\kern0pt}term\ list{\isacharparenright}{\kern0pt}{\isacharparenright}{\kern0pt}{\isacharasterisk}{\kern0pt}{\isacharparenleft}{\kern0pt}a{\isacharunderscore}{\kern0pt}term\ list{\isacharparenright}{\kern0pt}\ \isanewline
\ \ \ \ \ \ \ \ \ \ \ \ {\isacharbar}{\kern0pt}\ WhileA\ of\ term{\isacharasterisk}{\kern0pt}{\isacharparenleft}{\kern0pt}a{\isacharunderscore}{\kern0pt}term\ list{\isacharparenright}{\kern0pt}\isanewline
\ \ \ \ \ \ \ \ \ \ \ \ {\isacharbar}{\kern0pt}\ SeqA\ of\ a{\isacharunderscore}{\kern0pt}term\ {\isacharasterisk}{\kern0pt}\ a{\isacharunderscore}{\kern0pt}term\isanewline
\isanewline
type\ thread{\isacharunderscore}{\kern0pt}absy\ {\isacharequal}{\kern0pt}\ {\isacharbraceleft}{\kern0pt}nom{\isacharunderscore}{\kern0pt}thread\ \ {\isacharcolon}{\kern0pt}\ binding\ {\isacharcomma}{\kern0pt}\isanewline
\ \ \ \ \ \ \ \ \ \ \ \ \ \ \ \ \ \ \ \ locals{\isacharunderscore}{\kern0pt}decl\ {\isacharcolon}{\kern0pt}\ {\isacharparenleft}{\kern0pt}{\isacharparenleft}{\kern0pt}binding\ {\isacharasterisk}{\kern0pt}\ typ{\isacharparenright}{\kern0pt}\ {\isacharasterisk}{\kern0pt}\ term\ option{\isacharparenright}{\kern0pt}\ list\ {\isacharcomma}{\kern0pt}\isanewline
\ \ \ \ \ \ \ \ \ \ \ \ \ \ \ \ \ \ \ \ actions\ \ \ \ \ {\isacharcolon}{\kern0pt}\ a{\isacharunderscore}{\kern0pt}term\ list{\isacharcomma}{\kern0pt}\isanewline
\ \ \ \ \ \ \ \ \ \ \ \ \ \ \ \ \ \ \ \ locstab\ \ \ \ \ {\isacharcolon}{\kern0pt}\ {\isacharparenleft}{\kern0pt}term\ option{\isacharparenright}{\kern0pt}\ Symtab{\isachardot}{\kern0pt}table\isanewline
\ \ \ \ \ \ \ \ \ \ \ \ \ \ \ \ \ \ \ {\isacharbraceright}{\kern0pt}\isanewline
type\ absy\ {\isacharequal}{\kern0pt}\ {\isacharbraceleft}{\kern0pt}system\ \ \ \ \ \ \ \ \ \ {\isacharcolon}{\kern0pt}\ binding{\isacharcomma}{\kern0pt}\isanewline
\ \ \ \ \ \ \ \ \ \ \ \ \ globals{\isacharunderscore}{\kern0pt}decls\ \ \ {\isacharcolon}{\kern0pt}\ {\isacharparenleft}{\kern0pt}binding\ {\isacharasterisk}{\kern0pt}\ typ\ {\isacharasterisk}{\kern0pt}\ Position{\isachardot}{\kern0pt}T\ {\isacharasterisk}{\kern0pt}\ term\ option{\isacharparenright}{\kern0pt}\ list{\isacharcomma}{\kern0pt}\isanewline
\ \ \ \ \ \ \ \ \ \ \ \ \ varstab\ \ \ \ \ \ \ \ \ {\isacharcolon}{\kern0pt}\ {\isacharparenleft}{\kern0pt}term\ option{\isacharparenright}{\kern0pt}\ Symtab{\isachardot}{\kern0pt}table{\isacharcomma}{\kern0pt}\isanewline
\ \ \ \ \ \ \ \ \ \ \ \ \ locks{\isacharunderscore}{\kern0pt}decls\ \ \ \ \ {\isacharcolon}{\kern0pt}\ {\isacharparenleft}{\kern0pt}binding\ {\isacharasterisk}{\kern0pt}\ term\ {\isacharasterisk}{\kern0pt}\ Position{\isachardot}{\kern0pt}T{\isacharparenright}{\kern0pt}\ list{\isacharcomma}{\kern0pt}\isanewline
\ \ \ \ \ \ \ \ \ \ \ \ \ threads{\isacharunderscore}{\kern0pt}decls\ \ \ {\isacharcolon}{\kern0pt}\ thread{\isacharunderscore}{\kern0pt}absy\ list\isanewline
\ \ \ \ \ \ \ \ \ \ \ \ {\isacharbraceright}{\kern0pt}\isanewline
\isanewline
\isanewline
val\ SPY\ {\isacharequal}{\kern0pt}\ Unsynchronized{\isachardot}{\kern0pt}ref{\isacharparenleft}{\kern0pt}NONE{\isacharcolon}{\kern0pt}absy\ option{\isacharparenright}{\kern0pt}\isanewline
val\ SPYG{\isacharequal}{\kern0pt}\ Unsynchronized{\isachardot}{\kern0pt}ref{\isacharparenleft}{\kern0pt}{\isacharbrackleft}{\kern0pt}Bound\ {\isadigit{0}}{\isacharbrackright}{\kern0pt}{\isacharparenright}{\kern0pt}\isanewline
\isanewline
fun\ pre{\isacharunderscore}{\kern0pt}const\ {\isacharparenleft}{\kern0pt}{\isacharparenleft}{\kern0pt}bdg{\isacharcomma}{\kern0pt}\ typ{\isacharparenright}{\kern0pt}{\isacharcomma}{\kern0pt}\ {\isacharunderscore}{\kern0pt}{\isacharparenright}{\kern0pt}\ {\isacharequal}{\kern0pt}\ {\isacharparenleft}{\kern0pt}bdg{\isacharcomma}{\kern0pt}\ typ{\isacharcomma}{\kern0pt}\ NoSyn{\isacharparenright}{\kern0pt}\isanewline
fun\ pre{\isacharunderscore}{\kern0pt}const{\isacharprime}{\kern0pt}\ {\isacharparenleft}{\kern0pt}b{\isacharcomma}{\kern0pt}\ trm{\isacharparenright}{\kern0pt}\ {\isacharequal}{\kern0pt}\ {\isacharparenleft}{\kern0pt}b{\isacharcomma}{\kern0pt}\ \ \ fastype{\isacharunderscore}{\kern0pt}of\ trm{\isacharcomma}{\kern0pt}\ Mixfix{\isachardot}{\kern0pt}NoSyn{\isacharparenright}{\kern0pt}\ \isanewline
fun\ pre{\isacharunderscore}{\kern0pt}const{\isacharunderscore}{\kern0pt}bar\ {\isacharparenleft}{\kern0pt}b{\isacharcomma}{\kern0pt}\ trm{\isacharcomma}{\kern0pt}\ {\isacharunderscore}{\kern0pt}{\isacharparenright}{\kern0pt}\ {\isacharequal}{\kern0pt}\ {\isacharparenleft}{\kern0pt}Binding{\isachardot}{\kern0pt}suffix{\isacharunderscore}{\kern0pt}name\ {\isachardoublequote}{\kern0pt}{\isacharprime}{\kern0pt}{\isachardoublequote}{\kern0pt}\ b{\isacharcomma}{\kern0pt}\ \ \isanewline
\ \ \ \ \ \ \ \ \ \ \ \ \ \ \ \ \ \ \ \ \ \ \ \ \ \ \ \ \ \ \ \ \ \ fastype{\isacharunderscore}{\kern0pt}of\ trm{\isacharcomma}{\kern0pt}\ Mixfix{\isachardot}{\kern0pt}NoSyn{\isacharparenright}{\kern0pt}\ \isanewline
fun\ pre{\isacharunderscore}{\kern0pt}const{\isadigit{2}}\ {\isacharparenleft}{\kern0pt}b{\isacharcomma}{\kern0pt}\ typ{\isacharcomma}{\kern0pt}\ {\isacharunderscore}{\kern0pt}{\isacharcomma}{\kern0pt}\ {\isacharunderscore}{\kern0pt}{\isacharparenright}{\kern0pt}\ {\isacharequal}{\kern0pt}\ {\isacharparenleft}{\kern0pt}b{\isacharcomma}{\kern0pt}\ \ typ{\isacharcomma}{\kern0pt}\ Mixfix{\isachardot}{\kern0pt}NoSyn{\isacharparenright}{\kern0pt}\ \isanewline
fun\ pre{\isacharunderscore}{\kern0pt}const{\isadigit{3}}\ {\isacharparenleft}{\kern0pt}b{\isacharcomma}{\kern0pt}\ trm{\isacharcomma}{\kern0pt}\ {\isacharunderscore}{\kern0pt}{\isacharparenright}{\kern0pt}\ {\isacharequal}{\kern0pt}\ {\isacharparenleft}{\kern0pt}b{\isacharcomma}{\kern0pt}\ fastype{\isacharunderscore}{\kern0pt}of\ trm{\isacharcomma}{\kern0pt}\ Mixfix{\isachardot}{\kern0pt}NoSyn{\isacharparenright}{\kern0pt}\ \isanewline
\isanewline
fun\ pre{\isacharunderscore}{\kern0pt}event\ {\isacharparenleft}{\kern0pt}b{\isacharcomma}{\kern0pt}\ trm{\isacharcomma}{\kern0pt}\ {\isacharunderscore}{\kern0pt}{\isacharparenright}{\kern0pt}\ {\isacharequal}{\kern0pt}\ {\isacharparenleft}{\kern0pt}b{\isacharcomma}{\kern0pt}\ fastype{\isacharunderscore}{\kern0pt}of\ trm\ {\isacharminus}{\kern0pt}{\isacharminus}{\kern0pt}{\isachargreater}{\kern0pt}\ pre{\isacharunderscore}{\kern0pt}event{\isacharunderscore}{\kern0pt}type{\isacharcomma}{\kern0pt}\ \isanewline
\ \ \ \ \ \ \ \ \ \ \ \ \ \ \ \ \ \ \ \ \ \ \ \ \ \ \ \ \ \ \ \ \ \ \ Mixfix{\isachardot}{\kern0pt}NoSyn{\isacharparenright}{\kern0pt}\ \ \isanewline
\isanewline
\isanewline
{\isacharparenleft}{\kern0pt}{\isacharasterisk}{\kern0pt}\ Got\ replaced\ but\ would\ be\ cool\ to\ know\ why\ isn{\isacharprime}{\kern0pt}t\ it\ working\isanewline
\isanewline
fun\ check{\isacharunderscore}{\kern0pt}vars\ term\ varstab\ locstab\ {\isacharequal}{\kern0pt}\isanewline
\ \ let\isanewline
\ \ \ \ val\ vars\ {\isacharequal}{\kern0pt}\ Term{\isachardot}{\kern0pt}add{\isacharunderscore}{\kern0pt}free{\isacharunderscore}{\kern0pt}names\ term\ {\isacharbrackleft}{\kern0pt}{\isacharbrackright}{\kern0pt}\isanewline
\ \ \ \ fun\ is{\isacharunderscore}{\kern0pt}declared\ v\ {\isacharequal}{\kern0pt}\ Symtab{\isachardot}{\kern0pt}defined\ varstab\ v\ orelse\ Symtab{\isachardot}{\kern0pt}defined\ locstab\ v\isanewline
\ \ \ \ val\ undeclared\ {\isacharequal}{\kern0pt}\ List{\isachardot}{\kern0pt}filter\ {\isacharparenleft}{\kern0pt}not\ o\ is{\isacharunderscore}{\kern0pt}declared{\isacharparenright}{\kern0pt}\ vars\isanewline
\ \ in\isanewline
\ \ \ \ if\ null\ undeclared\ then\ {\isacharparenleft}{\kern0pt}{\isacharparenright}{\kern0pt}\isanewline
\ \ \ \ else\ error\ {\isacharparenleft}{\kern0pt}{\isachardoublequote}{\kern0pt}Undeclared\ variable{\isacharparenleft}{\kern0pt}s{\isacharparenright}{\kern0pt}\ used{\isacharcolon}{\kern0pt}\ {\isachardoublequote}{\kern0pt}\ {\isacharcircum}{\kern0pt}\ String{\isachardot}{\kern0pt}concatWith\ {\isachardoublequote}{\kern0pt}{\isacharcomma}{\kern0pt}\ {\isachardoublequote}{\kern0pt}\ undeclared{\isacharparenright}{\kern0pt}\isanewline
\ \ end{\isacharasterisk}{\kern0pt}{\isacharparenright}{\kern0pt}\isanewline
\isanewline
\isanewline
fun\ read{\isacharunderscore}{\kern0pt}term{\isacharunderscore}{\kern0pt}err\ thy\ msg\ bind\ str\ {\isacharequal}{\kern0pt}\ \isanewline
\ \ \ \ \ \ \ \ let\ val\ pos\ {\isacharequal}{\kern0pt}\ Position{\isachardot}{\kern0pt}here{\isacharparenleft}{\kern0pt}Binding{\isachardot}{\kern0pt}pos{\isacharunderscore}{\kern0pt}of\ bind{\isacharparenright}{\kern0pt}\isanewline
\ \ \ \ \ \ \ \ in\ \ {\isacharparenleft}{\kern0pt}Syntax{\isachardot}{\kern0pt}read{\isacharunderscore}{\kern0pt}term{\isacharunderscore}{\kern0pt}global\ thy\ str\isanewline
\ \ \ \ \ \ \ \ \ \ \ \ \ handle\ ERROR\ s\ {\isacharequal}{\kern0pt}{\isachargreater}{\kern0pt}\ error{\isacharparenleft}{\kern0pt}{\isachardoublequote}{\kern0pt}error\ in\ {\isachardoublequote}{\kern0pt}{\isacharcircum}{\kern0pt}msg{\isacharcircum}{\kern0pt}{\isachardoublequote}{\kern0pt}\ {\isacharcolon}{\kern0pt}{\isachardoublequote}{\kern0pt}{\isacharcircum}{\kern0pt}\ pos\ {\isacharcircum}{\kern0pt}{\isachardoublequote}{\kern0pt}{\isacharbackslash}{\kern0pt}n{\isachardoublequote}{\kern0pt}{\isacharcircum}{\kern0pt}\ s\ {\isacharparenright}{\kern0pt}\ {\isacharparenright}{\kern0pt}\ \isanewline
\ \ \ \ \ \ \ \ end\isanewline
\isanewline
fun\ read{\isacharunderscore}{\kern0pt}term{\isacharunderscore}{\kern0pt}err{\isacharunderscore}{\kern0pt}pos\ thy\ msg\ {\isacharunderscore}{\kern0pt}\ str\ pos\ {\isacharequal}{\kern0pt}\ \isanewline
\ \ \ \ \ \ \ \ {\isacharparenleft}{\kern0pt}Syntax{\isachardot}{\kern0pt}read{\isacharunderscore}{\kern0pt}term{\isacharunderscore}{\kern0pt}global\ thy\ str\isanewline
\ \ \ \ \ \ \ \ \ \ \ \ \ handle\ ERROR\ s\ {\isacharequal}{\kern0pt}{\isachargreater}{\kern0pt}\ error{\isacharparenleft}{\kern0pt}{\isachardoublequote}{\kern0pt}error\ in\ {\isachardoublequote}{\kern0pt}{\isacharcircum}{\kern0pt}msg{\isacharcircum}{\kern0pt}{\isachardoublequote}{\kern0pt}\ {\isacharcolon}{\kern0pt}{\isachardoublequote}{\kern0pt}{\isacharcircum}{\kern0pt}\ {\isacharparenleft}{\kern0pt}Position{\isachardot}{\kern0pt}here\ pos{\isacharparenright}{\kern0pt}\ {\isacharcircum}{\kern0pt}{\isachardoublequote}{\kern0pt}{\isacharbackslash}{\kern0pt}n{\isachardoublequote}{\kern0pt}{\isacharcircum}{\kern0pt}\ s\ {\isacharparenright}{\kern0pt}\ {\isacharparenright}{\kern0pt}\ \isanewline
\isanewline
{\isacharparenleft}{\kern0pt}{\isacharasterisk}{\kern0pt}\ Utilitaire\ pour\ lire\ un\ terme\ à\ partir\ d’un\ string\ {\isacharasterisk}{\kern0pt}{\isacharparenright}{\kern0pt}\isanewline
fun\ get{\isacharunderscore}{\kern0pt}term{\isacharunderscore}{\kern0pt}and{\isacharunderscore}{\kern0pt}vars\ thy\ str\ {\isacharequal}{\kern0pt}\isanewline
\ \ let\isanewline
\ \ \ \ val\ ctxt\ {\isacharequal}{\kern0pt}\ Proof{\isacharunderscore}{\kern0pt}Context{\isachardot}{\kern0pt}init{\isacharunderscore}{\kern0pt}global\ thy\isanewline
\ \ \ \ val\ term\ {\isacharequal}{\kern0pt}\ Syntax{\isachardot}{\kern0pt}read{\isacharunderscore}{\kern0pt}term\ ctxt\ str\isanewline
\ \ \ \ val\ vars\ {\isacharequal}{\kern0pt}\ Term{\isachardot}{\kern0pt}add{\isacharunderscore}{\kern0pt}free{\isacharunderscore}{\kern0pt}names\ term\ {\isacharbrackleft}{\kern0pt}{\isacharbrackright}{\kern0pt}\isanewline
\ \ in\isanewline
\ \ \ \ {\isacharparenleft}{\kern0pt}term{\isacharcomma}{\kern0pt}\ vars{\isacharparenright}{\kern0pt}\isanewline
\ \ end\isanewline
\isanewline
{\isacharparenleft}{\kern0pt}{\isacharasterisk}{\kern0pt}\ Résout\ une\ variable\ selon\ priorité\ locale\ ou\ globale\ {\isacharasterisk}{\kern0pt}{\isacharparenright}{\kern0pt}\isanewline
fun\ resolve{\isacharunderscore}{\kern0pt}var\ prefer{\isacharunderscore}{\kern0pt}global\ var\ locstab\ varstab\ {\isacharequal}{\kern0pt}\isanewline
\ \ let\isanewline
\ \ \ \ val\ loc\ {\isacharequal}{\kern0pt}\ Symtab{\isachardot}{\kern0pt}lookup\ locstab\ var\isanewline
\ \ \ \ val\ global\ {\isacharequal}{\kern0pt}\ Symtab{\isachardot}{\kern0pt}lookup\ varstab\ var\isanewline
\ \ in\isanewline
\ \ \ \ if\ prefer{\isacharunderscore}{\kern0pt}global\ then\isanewline
\ \ \ \ \ \ {\isacharparenleft}{\kern0pt}case\ global\ of\isanewline
\ \ \ \ \ \ \ \ \ SOME\ {\isacharparenleft}{\kern0pt}SOME\ t{\isacharparenright}{\kern0pt}\ {\isacharequal}{\kern0pt}{\isachargreater}{\kern0pt}\ SOME\ t\isanewline
\ \ \ \ \ \ \ {\isacharbar}{\kern0pt}\ {\isacharunderscore}{\kern0pt}\ {\isacharequal}{\kern0pt}{\isachargreater}{\kern0pt}\ {\isacharparenleft}{\kern0pt}case\ loc\ of\isanewline
\ \ \ \ \ \ \ \ \ \ \ \ \ \ \ \ \ SOME\ {\isacharparenleft}{\kern0pt}SOME\ t{\isacharparenright}{\kern0pt}\ {\isacharequal}{\kern0pt}{\isachargreater}{\kern0pt}\ SOME\ t\isanewline
\ \ \ \ \ \ \ \ \ \ \ \ \ \ \ {\isacharbar}{\kern0pt}\ {\isacharunderscore}{\kern0pt}\ {\isacharequal}{\kern0pt}{\isachargreater}{\kern0pt}\ NONE{\isacharparenright}{\kern0pt}{\isacharparenright}{\kern0pt}\isanewline
\ \ \ \ else\isanewline
\ \ \ \ \ \ {\isacharparenleft}{\kern0pt}case\ loc\ of\isanewline
\ \ \ \ \ \ \ \ \ SOME\ {\isacharparenleft}{\kern0pt}SOME\ t{\isacharparenright}{\kern0pt}\ {\isacharequal}{\kern0pt}{\isachargreater}{\kern0pt}\ SOME\ t\isanewline
\ \ \ \ \ \ \ {\isacharbar}{\kern0pt}\ {\isacharunderscore}{\kern0pt}\ {\isacharequal}{\kern0pt}{\isachargreater}{\kern0pt}\ {\isacharparenleft}{\kern0pt}case\ global\ of\isanewline
\ \ \ \ \ \ \ \ \ \ \ \ \ \ \ \ \ SOME\ {\isacharparenleft}{\kern0pt}SOME\ t{\isacharparenright}{\kern0pt}\ {\isacharequal}{\kern0pt}{\isachargreater}{\kern0pt}\ SOME\ t\isanewline
\ \ \ \ \ \ \ \ \ \ \ \ \ \ \ {\isacharbar}{\kern0pt}\ {\isacharunderscore}{\kern0pt}\ {\isacharequal}{\kern0pt}{\isachargreater}{\kern0pt}\ NONE{\isacharparenright}{\kern0pt}{\isacharparenright}{\kern0pt}\isanewline
\ \ end\isanewline
\isanewline
{\isacharparenleft}{\kern0pt}{\isacharasterisk}{\kern0pt}\ Variante\ stricte\ qui\ renvoie\ une\ erreur\ si\ non\ trouvé\ {\isacharasterisk}{\kern0pt}{\isacharparenright}{\kern0pt}\isanewline
fun\ resolve{\isacharunderscore}{\kern0pt}var{\isacharunderscore}{\kern0pt}strict\ prefer{\isacharunderscore}{\kern0pt}global\ var\ locstab\ varstab\ {\isacharequal}{\kern0pt}\isanewline
\ \ case\ resolve{\isacharunderscore}{\kern0pt}var\ prefer{\isacharunderscore}{\kern0pt}global\ var\ locstab\ varstab\ of\isanewline
\ \ \ \ SOME\ t\ {\isacharequal}{\kern0pt}{\isachargreater}{\kern0pt}\ t\isanewline
\ \ {\isacharbar}{\kern0pt}\ NONE\ {\isacharequal}{\kern0pt}{\isachargreater}{\kern0pt}\ error\ {\isacharparenleft}{\kern0pt}{\isachardoublequote}{\kern0pt}Variable\ {\isachardoublequote}{\kern0pt}{\isacharcircum}{\kern0pt}var{\isacharcircum}{\kern0pt}{\isachardoublequote}{\kern0pt}\ non\ déclarée{\isachardoublequote}{\kern0pt}{\isacharparenright}{\kern0pt}\isanewline
\isanewline
fun\ check{\isacharunderscore}{\kern0pt}thread\ thy\ thy{\isadigit{3}}\ vars{\isacharunderscore}{\kern0pt}tab\ {\isacharparenleft}{\kern0pt}{\isacharparenleft}{\kern0pt}{\isacharparenleft}{\kern0pt}SOME\ thread{\isacharunderscore}{\kern0pt}name{\isacharcomma}{\kern0pt}\ locals{\isacharunderscore}{\kern0pt}opt{\isacharparenright}{\kern0pt}{\isacharcomma}{\kern0pt}\ actions{\isacharunderscore}{\kern0pt}list{\isacharparenright}{\kern0pt}{\isacharcolon}{\kern0pt}\ raw{\isacharunderscore}{\kern0pt}thread{\isacharunderscore}{\kern0pt}absy{\isacharparenright}{\kern0pt}\ {\isacharequal}{\kern0pt}\isanewline
\ \ let\isanewline
\ \ \ \ fun\ loc{\isacharunderscore}{\kern0pt}decls{\isacharunderscore}{\kern0pt}conv\ {\isacharparenleft}{\kern0pt}{\isacharparenleft}{\kern0pt}bdg{\isacharcomma}{\kern0pt}\ {\isacharparenleft}{\kern0pt}s{\isacharcomma}{\kern0pt}\ pos{\isacharparenright}{\kern0pt}{\isacharparenright}{\kern0pt}{\isacharcomma}{\kern0pt}\ SOME\ init{\isacharunderscore}{\kern0pt}term{\isacharparenright}{\kern0pt}\ {\isacharequal}{\kern0pt}\isanewline
\ \ \ \ \ \ \ \ \ \ let\ \isanewline
\ \ \ \ \ \ \ \ \ \ \ \ val\ typ\ {\isacharequal}{\kern0pt}\ Syntax{\isachardot}{\kern0pt}read{\isacharunderscore}{\kern0pt}typ{\isacharunderscore}{\kern0pt}global\ thy\ s\isanewline
\ \ \ \ \ \ \ \ \ \ \ \ val\ inferred{\isacharunderscore}{\kern0pt}type\ {\isacharequal}{\kern0pt}\ fastype{\isacharunderscore}{\kern0pt}of\ init{\isacharunderscore}{\kern0pt}term\isanewline
\ \ \ \ \ \ \ \ \ \ \ \ val\ ctxt\ {\isacharequal}{\kern0pt}\ Proof{\isacharunderscore}{\kern0pt}Context{\isachardot}{\kern0pt}init{\isacharunderscore}{\kern0pt}global\ thy\isanewline
\ \ \ \ \ \ \ \ \ \ \ \ val\ inferred{\isacharunderscore}{\kern0pt}str\ {\isacharequal}{\kern0pt}\ Syntax{\isachardot}{\kern0pt}string{\isacharunderscore}{\kern0pt}of{\isacharunderscore}{\kern0pt}typ\ ctxt\ inferred{\isacharunderscore}{\kern0pt}type\isanewline
\ \ \ \ \ \ \ \ \ \ \ \ val\ declared{\isacharunderscore}{\kern0pt}str\ {\isacharequal}{\kern0pt}\ Syntax{\isachardot}{\kern0pt}string{\isacharunderscore}{\kern0pt}of{\isacharunderscore}{\kern0pt}typ\ ctxt\ typ\isanewline
\ \ \ \ \ \ \ \ \ \ in\isanewline
\ \ \ \ \ \ \ \ \ \ \ \ if\ inferred{\isacharunderscore}{\kern0pt}type\ {\isacharless}{\kern0pt}{\isachargreater}{\kern0pt}\ typ\ then\isanewline
\ \ \ \ \ \ \ \ \ \ \ \ \ \ error\ {\isacharparenleft}{\kern0pt}{\isachardoublequote}{\kern0pt}Type\ mismatch\ in\ variable\ {\isachardoublequote}{\kern0pt}\ {\isacharcircum}{\kern0pt}\ Binding{\isachardot}{\kern0pt}name{\isacharunderscore}{\kern0pt}of\ bdg\ {\isacharcircum}{\kern0pt}\isanewline
\ \ \ \ \ \ \ \ \ \ \ \ \ \ \ \ \ \ \ \ \ {\isachardoublequote}{\kern0pt}{\isacharcolon}{\kern0pt}\ expected\ {\isachardoublequote}{\kern0pt}\ {\isacharcircum}{\kern0pt}\ declared{\isacharunderscore}{\kern0pt}str\ {\isacharcircum}{\kern0pt}\ {\isachardoublequote}{\kern0pt}{\isacharcomma}{\kern0pt}\ but\ got\ {\isachardoublequote}{\kern0pt}\ {\isacharcircum}{\kern0pt}\ inferred{\isacharunderscore}{\kern0pt}str\ {\isacharcircum}{\kern0pt}\ {\isachardoublequote}{\kern0pt}at\ {\isachardoublequote}{\kern0pt}\ {\isacharcircum}{\kern0pt}\ Position{\isachardot}{\kern0pt}here\ pos\ {\isacharparenright}{\kern0pt}\isanewline
\ \ \ \ \ \ \ \ \ \ \ \ else\isanewline
\ \ \ \ \ \ \ \ \ \ \ \ \ \ {\isacharparenleft}{\kern0pt}{\isacharparenleft}{\kern0pt}bdg{\isacharcomma}{\kern0pt}\ typ{\isacharparenright}{\kern0pt}{\isacharcomma}{\kern0pt}\ SOME\ init{\isacharunderscore}{\kern0pt}term{\isacharparenright}{\kern0pt}\isanewline
\ \ \ \ \ \ \ \ \ \ end\isanewline
\ \ \ \ \ \ \ \ {\isacharbar}{\kern0pt}\ loc{\isacharunderscore}{\kern0pt}decls{\isacharunderscore}{\kern0pt}conv\ {\isacharparenleft}{\kern0pt}{\isacharparenleft}{\kern0pt}bdg{\isacharcomma}{\kern0pt}\ {\isacharparenleft}{\kern0pt}s{\isacharcomma}{\kern0pt}\ pos{\isacharparenright}{\kern0pt}{\isacharparenright}{\kern0pt}{\isacharcomma}{\kern0pt}\ NONE{\isacharparenright}{\kern0pt}\ {\isacharequal}{\kern0pt}\isanewline
\ \ \ \ \ \ \ \ \ \ {\isacharparenleft}{\kern0pt}{\isacharparenleft}{\kern0pt}bdg{\isacharcomma}{\kern0pt}\ Syntax{\isachardot}{\kern0pt}read{\isacharunderscore}{\kern0pt}typ{\isacharunderscore}{\kern0pt}global\ thy\ s{\isacharparenright}{\kern0pt}{\isacharcomma}{\kern0pt}\ NONE{\isacharparenright}{\kern0pt}\isanewline
\isanewline
\ \ \ \ fun\ action{\isacharunderscore}{\kern0pt}conv\ {\isacharunderscore}{\kern0pt}\ {\isacharunderscore}{\kern0pt}\ Skip\ {\isacharequal}{\kern0pt}\ SkipA\isanewline
\ \ \ \ \ \ {\isacharbar}{\kern0pt}\ action{\isacharunderscore}{\kern0pt}conv\ {\isacharunderscore}{\kern0pt}\ {\isacharunderscore}{\kern0pt}\ {\isacharparenleft}{\kern0pt}Lock\ b{\isacharparenright}{\kern0pt}\ {\isacharequal}{\kern0pt}\ LockA\ b\isanewline
\ \ \ \ \ \ {\isacharbar}{\kern0pt}\ action{\isacharunderscore}{\kern0pt}conv\ {\isacharunderscore}{\kern0pt}\ {\isacharunderscore}{\kern0pt}\ {\isacharparenleft}{\kern0pt}Unlock\ b{\isacharparenright}{\kern0pt}\ {\isacharequal}{\kern0pt}\ UnlockA\ b\isanewline
\ \ \ \ \ \ {\isacharbar}{\kern0pt}\ action{\isacharunderscore}{\kern0pt}conv\ varstab\ locstab\ {\isacharparenleft}{\kern0pt}Assign\ {\isacharparenleft}{\kern0pt}bdg{\isacharcomma}{\kern0pt}\ str{\isacharparenright}{\kern0pt}{\isacharparenright}{\kern0pt}\ {\isacharequal}{\kern0pt}\isanewline
\ \ \ \ \ \ \ \ \ \ \ let\isanewline
\ \ \ \ \ \ \ \ \ \ \ \ \ \ val\ {\isacharparenleft}{\kern0pt}term{\isacharcomma}{\kern0pt}\ vars{\isacharparenright}{\kern0pt}\ {\isacharequal}{\kern0pt}\ get{\isacharunderscore}{\kern0pt}term{\isacharunderscore}{\kern0pt}and{\isacharunderscore}{\kern0pt}vars\ thy\ str\isanewline
\ \ \ \ \ \ \ \ \ \ \ \ \ \ val\ assign{\isacharunderscore}{\kern0pt}term\ {\isacharequal}{\kern0pt}\ term\isanewline
\ \ \ \ \ \ \ \ \ \ \ \ \ val\ {\isacharunderscore}{\kern0pt}\ {\isacharequal}{\kern0pt}\isanewline
\ \ \ \ \ \ \ \ \ \ \ \ \ \ if\ Symtab{\isachardot}{\kern0pt}defined\ locstab\ {\isacharparenleft}{\kern0pt}Binding{\isachardot}{\kern0pt}name{\isacharunderscore}{\kern0pt}of\ bdg{\isacharparenright}{\kern0pt}\ \isanewline
\ \ \ \ \ \ \ \ \ \ \ \ \ \ \ \ \ orelse\ Symtab{\isachardot}{\kern0pt}defined\ varstab\ {\isacharparenleft}{\kern0pt}Binding{\isachardot}{\kern0pt}name{\isacharunderscore}{\kern0pt}of\ bdg{\isacharparenright}{\kern0pt}\isanewline
\ \ \ \ \ \ \ \ \ \ \ \ \ \ then\ {\isacharparenleft}{\kern0pt}{\isacharparenright}{\kern0pt}\isanewline
\ \ \ \ \ \ \ \ \ \ \ \ \ \ else\ error\ {\isacharparenleft}{\kern0pt}{\isachardoublequote}{\kern0pt}Variable\ {\isachardoublequote}{\kern0pt}\ {\isacharcircum}{\kern0pt}\ {\isacharparenleft}{\kern0pt}Binding{\isachardot}{\kern0pt}name{\isacharunderscore}{\kern0pt}of\ bdg{\isacharparenright}{\kern0pt}\ {\isacharcircum}{\kern0pt}\ {\isachardoublequote}{\kern0pt}\ is\ not\ declared\ {\isacharparenleft}{\kern0pt}in\ assign\ binding{\isacharparenright}{\kern0pt}{\isachardoublequote}{\kern0pt}{\isacharparenright}{\kern0pt}\isanewline
\ \ \ \ \ \ \ \ \ \ \ \ in\isanewline
\ \ \ \ \ \ \ \ \ \ \ \ \ \ AssignA\ {\isacharparenleft}{\kern0pt}bdg{\isacharcomma}{\kern0pt}\ assign{\isacharunderscore}{\kern0pt}term{\isacharparenright}{\kern0pt}\isanewline
\ \ \ \ \ \ \ \ \ \ \ \ end\isanewline
\isanewline
\isanewline
\ \ \ \ \ \ {\isacharbar}{\kern0pt}\ action{\isacharunderscore}{\kern0pt}conv\ varstab\ locstab\ {\isacharparenleft}{\kern0pt}Send\ {\isacharparenleft}{\kern0pt}b{\isacharcomma}{\kern0pt}\ str{\isacharparenright}{\kern0pt}{\isacharparenright}{\kern0pt}\ {\isacharequal}{\kern0pt}\isanewline
\ \ \ \ \ \ \ \ \ \ let\isanewline
\ \ \ \ \ \ \ \ \ \ \ \ val\ {\isacharparenleft}{\kern0pt}term{\isacharcomma}{\kern0pt}\ {\isacharunderscore}{\kern0pt}{\isacharparenright}{\kern0pt}\ {\isacharequal}{\kern0pt}\ get{\isacharunderscore}{\kern0pt}term{\isacharunderscore}{\kern0pt}and{\isacharunderscore}{\kern0pt}vars\ thy\ str\isanewline
\ \ \ \ \ \ \ \ \ \ \ \ val\ {\isacharunderscore}{\kern0pt}\ {\isacharequal}{\kern0pt}\isanewline
\ \ \ \ \ \ \ \ \ \ \ \ \ \ if\ Symtab{\isachardot}{\kern0pt}defined\ locstab\ {\isacharparenleft}{\kern0pt}Binding{\isachardot}{\kern0pt}name{\isacharunderscore}{\kern0pt}of\ b{\isacharparenright}{\kern0pt}\ \isanewline
\ \ \ \ \ \ \ \ \ \ \ \ \ \ \ \ \ orelse\ Symtab{\isachardot}{\kern0pt}defined\ varstab\ {\isacharparenleft}{\kern0pt}Binding{\isachardot}{\kern0pt}name{\isacharunderscore}{\kern0pt}of\ b{\isacharparenright}{\kern0pt}\isanewline
\ \ \ \ \ \ \ \ \ \ \ \ \ \ then\ {\isacharparenleft}{\kern0pt}{\isacharparenright}{\kern0pt}\isanewline
\ \ \ \ \ \ \ \ \ \ \ \ \ \ else\ error\ {\isacharparenleft}{\kern0pt}{\isachardoublequote}{\kern0pt}Variable\ {\isachardoublequote}{\kern0pt}\ {\isacharcircum}{\kern0pt}\ {\isacharparenleft}{\kern0pt}Binding{\isachardot}{\kern0pt}name{\isacharunderscore}{\kern0pt}of\ b{\isacharparenright}{\kern0pt}\ {\isacharcircum}{\kern0pt}\ {\isachardoublequote}{\kern0pt}\ is\ not\ declared\ {\isacharparenleft}{\kern0pt}in\ send{\isacharparenright}{\kern0pt}{\isachardoublequote}{\kern0pt}{\isacharparenright}{\kern0pt}\isanewline
\ \ \ \ \ \ \ \ \ \ in\isanewline
\ \ \ \ \ \ \ \ \ \ \ \ SendA\ {\isacharparenleft}{\kern0pt}b{\isacharcomma}{\kern0pt}\ term{\isacharparenright}{\kern0pt}\isanewline
\ \ \ \ \ \ \ \ \ \ end\isanewline
\isanewline
\ \ \ \ \ \ {\isacharbar}{\kern0pt}\ action{\isacharunderscore}{\kern0pt}conv\ varstab\ locstab\ {\isacharparenleft}{\kern0pt}Receive\ {\isacharparenleft}{\kern0pt}blv{\isacharcomma}{\kern0pt}\ bsv{\isacharparenright}{\kern0pt}{\isacharparenright}{\kern0pt}\ {\isacharequal}{\kern0pt}\isanewline
\ \ \ \ \ \ \ \ let\isanewline
\ \ \ \ \ \ \ \ \ \ \ val\ {\isacharunderscore}{\kern0pt}\ {\isacharequal}{\kern0pt}\isanewline
\ \ \ \ \ \ \ \ \ \ \ \ \ \ if\ Symtab{\isachardot}{\kern0pt}defined\ locstab\ {\isacharparenleft}{\kern0pt}Binding{\isachardot}{\kern0pt}name{\isacharunderscore}{\kern0pt}of\ blv{\isacharparenright}{\kern0pt}\isanewline
\ \ \ \ \ \ \ \ \ \ \ \ \ \ \ \ \ then\ \ \ if\ {\isacharparenleft}{\kern0pt}Symtab{\isachardot}{\kern0pt}defined\ varstab\ {\isacharparenleft}{\kern0pt}Binding{\isachardot}{\kern0pt}name{\isacharunderscore}{\kern0pt}of\ bsv{\isacharparenright}{\kern0pt}{\isacharparenright}{\kern0pt}\ then\ {\isacharparenleft}{\kern0pt}{\isacharparenright}{\kern0pt}\isanewline
\ \ \ \ \ \ \ \ \ \ \ \ \ \ \ \ \ \ \ \ \ \ \ \ else\ error\ {\isacharparenleft}{\kern0pt}{\isachardoublequote}{\kern0pt}Variable\ {\isachardoublequote}{\kern0pt}\ {\isacharcircum}{\kern0pt}\ {\isacharparenleft}{\kern0pt}Binding{\isachardot}{\kern0pt}name{\isacharunderscore}{\kern0pt}of\ bsv{\isacharparenright}{\kern0pt}\ {\isacharcircum}{\kern0pt}\ {\isachardoublequote}{\kern0pt}\ is\ not\ declared\ {\isacharparenleft}{\kern0pt}in\ receive\ binding{\isacharparenright}{\kern0pt}{\isachardoublequote}{\kern0pt}{\isacharparenright}{\kern0pt}\isanewline
\ \ \ \ \ \ \ \ \ \ \ \ \ \ else\ error\ {\isacharparenleft}{\kern0pt}{\isachardoublequote}{\kern0pt}Variable\ {\isachardoublequote}{\kern0pt}\ {\isacharcircum}{\kern0pt}\ {\isacharparenleft}{\kern0pt}Binding{\isachardot}{\kern0pt}name{\isacharunderscore}{\kern0pt}of\ blv{\isacharparenright}{\kern0pt}\ {\isacharcircum}{\kern0pt}\ {\isachardoublequote}{\kern0pt}\ is\ not\ declared\ {\isacharparenleft}{\kern0pt}in\ receive\ binding{\isacharparenright}{\kern0pt}{\isachardoublequote}{\kern0pt}{\isacharparenright}{\kern0pt}\isanewline
\ \ \ \ \ \ \ \ in\isanewline
\ \ \ \ \ \ \ \ \ \ ReceiveA\ {\isacharparenleft}{\kern0pt}blv{\isacharcomma}{\kern0pt}\ bsv{\isacharparenright}{\kern0pt}\isanewline
\ \ \ \ \ \ \ \ end\isanewline
\isanewline
\ \ \ \ \ \ {\isacharbar}{\kern0pt}\ action{\isacharunderscore}{\kern0pt}conv\ varstab\ locstab\ {\isacharparenleft}{\kern0pt}Ifelse\ {\isacharparenleft}{\kern0pt}{\isacharparenleft}{\kern0pt}cond{\isacharunderscore}{\kern0pt}str{\isacharcomma}{\kern0pt}\ then{\isacharunderscore}{\kern0pt}branch{\isacharparenright}{\kern0pt}{\isacharcomma}{\kern0pt}\ else{\isacharunderscore}{\kern0pt}branch{\isacharparenright}{\kern0pt}{\isacharparenright}{\kern0pt}\ {\isacharequal}{\kern0pt}\isanewline
\ \ \ \ \ \ \ \ \ \ let\isanewline
\ \ \ \ \ \ \ \ \ \ val\ {\isacharparenleft}{\kern0pt}term{\isacharcomma}{\kern0pt}\ vars{\isacharparenright}{\kern0pt}\ {\isacharequal}{\kern0pt}\ get{\isacharunderscore}{\kern0pt}term{\isacharunderscore}{\kern0pt}and{\isacharunderscore}{\kern0pt}vars\ thy\ cond{\isacharunderscore}{\kern0pt}str\isanewline
\ \ \ \ \ \ \ \ \ \ val\ cond{\isacharunderscore}{\kern0pt}term\ {\isacharequal}{\kern0pt}\ term\isanewline
\ \ \ \ \ \ \ \ \ \ val\ then{\isacharunderscore}{\kern0pt}branch{\isacharprime}{\kern0pt}\ {\isacharequal}{\kern0pt}\ map\ {\isacharparenleft}{\kern0pt}action{\isacharunderscore}{\kern0pt}conv\ varstab\ locstab{\isacharparenright}{\kern0pt}\ then{\isacharunderscore}{\kern0pt}branch\isanewline
\ \ \ \ \ \ \ \ \ \ val\ else{\isacharunderscore}{\kern0pt}branch{\isacharprime}{\kern0pt}\ {\isacharequal}{\kern0pt}\ map\ {\isacharparenleft}{\kern0pt}action{\isacharunderscore}{\kern0pt}conv\ varstab\ locstab{\isacharparenright}{\kern0pt}\ else{\isacharunderscore}{\kern0pt}branch\isanewline
\ \ \ \ \ \ \ \ in\isanewline
\ \ \ \ \ \ \ \ \ \ IfelseA\ {\isacharparenleft}{\kern0pt}{\isacharparenleft}{\kern0pt}cond{\isacharunderscore}{\kern0pt}term{\isacharcomma}{\kern0pt}\ then{\isacharunderscore}{\kern0pt}branch{\isacharprime}{\kern0pt}{\isacharparenright}{\kern0pt}{\isacharcomma}{\kern0pt}\ else{\isacharunderscore}{\kern0pt}branch{\isacharprime}{\kern0pt}{\isacharparenright}{\kern0pt}\isanewline
\ \ \ \ \ \ \ \ end\isanewline
\isanewline
\ \ \ \ \ \ {\isacharbar}{\kern0pt}\ action{\isacharunderscore}{\kern0pt}conv\ varstab\ locstab\ {\isacharparenleft}{\kern0pt}While\ {\isacharparenleft}{\kern0pt}cond{\isacharunderscore}{\kern0pt}str{\isacharcomma}{\kern0pt}\ body{\isacharparenright}{\kern0pt}{\isacharparenright}{\kern0pt}\ {\isacharequal}{\kern0pt}\isanewline
\ \ \ \ \ \ \ \ \ \ let\isanewline
\ \ \ \ \ \ \ \ \ \ \ \ val\ {\isacharparenleft}{\kern0pt}term{\isacharcomma}{\kern0pt}\ vars{\isacharparenright}{\kern0pt}\ {\isacharequal}{\kern0pt}\ get{\isacharunderscore}{\kern0pt}term{\isacharunderscore}{\kern0pt}and{\isacharunderscore}{\kern0pt}vars\ thy\ cond{\isacharunderscore}{\kern0pt}str\isanewline
\ \ \ \ \ \ \ \ \ \ \ {\isacharparenleft}{\kern0pt}{\isacharasterisk}{\kern0pt}\ val\ cond{\isacharunderscore}{\kern0pt}term\ {\isacharequal}{\kern0pt}\isanewline
\ \ \ \ \ \ \ \ \ \ \ \ \ \ case\ vars\ of\isanewline
\ \ \ \ \ \ \ \ \ \ \ \ \ \ \ \ {\isacharbrackleft}{\kern0pt}x{\isacharbrackright}{\kern0pt}\ {\isacharequal}{\kern0pt}{\isachargreater}{\kern0pt}\ resolve{\isacharunderscore}{\kern0pt}var{\isacharunderscore}{\kern0pt}strict\ false\ x\ locstab\ varstab\isanewline
\ \ \ \ \ \ \ \ \ \ \ \ \ \ {\isacharbar}{\kern0pt}\ {\isacharunderscore}{\kern0pt}\ {\isacharequal}{\kern0pt}{\isachargreater}{\kern0pt}\ term{\isacharasterisk}{\kern0pt}{\isacharparenright}{\kern0pt}\isanewline
\ \ \ \ \ \ \ \ \ \ \ \ val\ cond{\isacharunderscore}{\kern0pt}term\ {\isacharequal}{\kern0pt}\ \ term\ \ \isanewline
\ \ \ \ \ \ \ \ \ \ \ \ val\ actions\ {\isacharequal}{\kern0pt}\ map\ {\isacharparenleft}{\kern0pt}action{\isacharunderscore}{\kern0pt}conv\ varstab\ locstab{\isacharparenright}{\kern0pt}\ body\isanewline
\ \ \ \ \ \ \ \ \ \ in\isanewline
\ \ \ \ \ \ \ \ \ \ \ \ WhileA\ {\isacharparenleft}{\kern0pt}cond{\isacharunderscore}{\kern0pt}term{\isacharcomma}{\kern0pt}\ actions{\isacharparenright}{\kern0pt}\isanewline
\ \ \ \ \ \ \ \ \ \ end\isanewline
\isanewline
\ \ \ \ val\ {\isacharparenleft}{\kern0pt}loc{\isacharunderscore}{\kern0pt}decls{\isacharprime}{\kern0pt}{\isacharcomma}{\kern0pt}\ thy{\isadigit{3}}{\isacharprime}{\kern0pt}{\isacharparenright}{\kern0pt}\ {\isacharequal}{\kern0pt}\ case\ locals{\isacharunderscore}{\kern0pt}opt\ of\isanewline
\ \ \ \ \ \ \ \ NONE\ {\isacharequal}{\kern0pt}{\isachargreater}{\kern0pt}\ {\isacharparenleft}{\kern0pt}{\isacharbrackleft}{\kern0pt}{\isacharbrackright}{\kern0pt}{\isacharcomma}{\kern0pt}\ thy{\isadigit{3}}{\isacharparenright}{\kern0pt}\isanewline
\ \ \ \ \ \ {\isacharbar}{\kern0pt}\ SOME\ locals\ {\isacharequal}{\kern0pt}{\isachargreater}{\kern0pt}\isanewline
\ \ \ \ \ \ \ \ \ \ let\isanewline
\ \ \ \ \ \ \ \ \ \ \ \ val\ locals{\isacharunderscore}{\kern0pt}terms\ {\isacharequal}{\kern0pt}\ map\ {\isacharparenleft}{\kern0pt}fn\ {\isacharparenleft}{\kern0pt}{\isacharparenleft}{\kern0pt}bdg{\isacharcomma}{\kern0pt}{\isacharparenleft}{\kern0pt}s{\isacharcomma}{\kern0pt}pos{\isacharparenright}{\kern0pt}{\isacharparenright}{\kern0pt}{\isacharcomma}{\kern0pt}\ init{\isacharunderscore}{\kern0pt}opt{\isacharunderscore}{\kern0pt}str{\isacharparenright}{\kern0pt}\ {\isacharequal}{\kern0pt}{\isachargreater}{\kern0pt}\isanewline
\ \ \ \ \ \ \ \ \ \ \ \ \ \ \ let\isanewline
\ \ \ \ \ \ \ \ \ \ \ \ \ \ \ \ \ val\ init{\isacharunderscore}{\kern0pt}opt{\isacharunderscore}{\kern0pt}term\ {\isacharequal}{\kern0pt}\ case\ init{\isacharunderscore}{\kern0pt}opt{\isacharunderscore}{\kern0pt}str\ of\isanewline
\ \ \ \ \ \ \ \ \ \ \ \ \ \ \ \ \ \ \ \ \ NONE\ {\isacharequal}{\kern0pt}{\isachargreater}{\kern0pt}\ NONE\isanewline
\ \ \ \ \ \ \ \ \ \ \ \ \ \ \ \ \ \ \ {\isacharbar}{\kern0pt}\ SOME\ str\ {\isacharequal}{\kern0pt}{\isachargreater}{\kern0pt}\ SOME\ {\isacharparenleft}{\kern0pt}read{\isacharunderscore}{\kern0pt}term{\isacharunderscore}{\kern0pt}err\ thy\ {\isachardoublequote}{\kern0pt}local\ init{\isachardoublequote}{\kern0pt}\ bdg\ str{\isacharparenright}{\kern0pt}\isanewline
\ \ \ \ \ \ \ \ \ \ \ \ \ \ \ in\isanewline
\ \ \ \ \ \ \ \ \ \ \ \ \ \ \ \ \ {\isacharparenleft}{\kern0pt}{\isacharparenleft}{\kern0pt}bdg{\isacharcomma}{\kern0pt}\ {\isacharparenleft}{\kern0pt}s{\isacharcomma}{\kern0pt}pos{\isacharparenright}{\kern0pt}{\isacharparenright}{\kern0pt}{\isacharcomma}{\kern0pt}\ init{\isacharunderscore}{\kern0pt}opt{\isacharunderscore}{\kern0pt}term{\isacharparenright}{\kern0pt}\isanewline
\ \ \ \ \ \ \ \ \ \ \ \ \ \ \ end{\isacharparenright}{\kern0pt}\ locals\isanewline
\ \ \ \ \isanewline
\ \ \ \ \ \ \ \ \ \ \ \ val\ decls\ {\isacharequal}{\kern0pt}\ map\ loc{\isacharunderscore}{\kern0pt}decls{\isacharunderscore}{\kern0pt}conv\ locals{\isacharunderscore}{\kern0pt}terms\isanewline
\ \ \ \ \ \ \ \ \ \ \ \ val\ thy{\isadigit{4}}\ {\isacharequal}{\kern0pt}\ Sign{\isachardot}{\kern0pt}add{\isacharunderscore}{\kern0pt}consts\ {\isacharparenleft}{\kern0pt}map\ pre{\isacharunderscore}{\kern0pt}const\ decls{\isacharparenright}{\kern0pt}\ thy{\isadigit{3}}\isanewline
\ \ \ \ \ \ \ \ \ \ in\isanewline
\ \ \ \ \ \ \ \ \ \ \ \ {\isacharparenleft}{\kern0pt}decls{\isacharcomma}{\kern0pt}\ thy{\isadigit{4}}{\isacharparenright}{\kern0pt}\isanewline
\ \ \ \ \ \ \ \ \ \ end\isanewline
\isanewline
\ \ \ \ fun\ conv{\isacharunderscore}{\kern0pt}decl{\isacharunderscore}{\kern0pt}init\ {\isacharparenleft}{\kern0pt}{\isacharparenleft}{\kern0pt}n{\isacharcomma}{\kern0pt}\ {\isacharunderscore}{\kern0pt}{\isacharparenright}{\kern0pt}{\isacharcomma}{\kern0pt}\ init{\isacharunderscore}{\kern0pt}opt{\isacharparenright}{\kern0pt}\ {\isacharequal}{\kern0pt}\ {\isacharparenleft}{\kern0pt}Binding{\isachardot}{\kern0pt}name{\isacharunderscore}{\kern0pt}of\ n{\isacharcomma}{\kern0pt}\ init{\isacharunderscore}{\kern0pt}opt{\isacharparenright}{\kern0pt}\isanewline
\ \ \ \ val\ locstab{\isacharprime}{\kern0pt}\ \ \ \ \ {\isacharequal}{\kern0pt}\ Symtab{\isachardot}{\kern0pt}make\ {\isacharparenleft}{\kern0pt}map\ conv{\isacharunderscore}{\kern0pt}decl{\isacharunderscore}{\kern0pt}init\ loc{\isacharunderscore}{\kern0pt}decls{\isacharprime}{\kern0pt}{\isacharparenright}{\kern0pt}\isanewline
\ \ \ \ val\ actions{\isacharprime}{\kern0pt}\ {\isacharequal}{\kern0pt}\ map\ {\isacharparenleft}{\kern0pt}action{\isacharunderscore}{\kern0pt}conv\ vars{\isacharunderscore}{\kern0pt}tab\ locstab{\isacharprime}{\kern0pt}{\isacharparenright}{\kern0pt}\ actions{\isacharunderscore}{\kern0pt}list\ \ in\isanewline
\ \ \ \ {\isacharbraceleft}{\kern0pt}\ nom{\isacharunderscore}{\kern0pt}thread\ {\isacharequal}{\kern0pt}\ thread{\isacharunderscore}{\kern0pt}name{\isacharcomma}{\kern0pt}\ locals{\isacharunderscore}{\kern0pt}decl\ {\isacharequal}{\kern0pt}\ loc{\isacharunderscore}{\kern0pt}decls{\isacharprime}{\kern0pt}{\isacharcomma}{\kern0pt}\ actions\ {\isacharequal}{\kern0pt}\ actions{\isacharprime}{\kern0pt}{\isacharcomma}{\kern0pt}\ locstab\ {\isacharequal}{\kern0pt}\ locstab{\isacharprime}{\kern0pt}\ {\isacharbraceright}{\kern0pt}\isanewline
\ \ end\isanewline
\ \ {\isacharbar}{\kern0pt}\ check{\isacharunderscore}{\kern0pt}thread\ {\isacharunderscore}{\kern0pt}\ {\isacharunderscore}{\kern0pt}\ {\isacharunderscore}{\kern0pt}\ {\isacharparenleft}{\kern0pt}{\isacharparenleft}{\kern0pt}NONE{\isacharcomma}{\kern0pt}\ {\isacharunderscore}{\kern0pt}{\isacharparenright}{\kern0pt}{\isacharcomma}{\kern0pt}\ {\isacharunderscore}{\kern0pt}{\isacharparenright}{\kern0pt}\ {\isacharequal}{\kern0pt}\isanewline
\ \ \ \ \ \ error\ {\isachardoublequote}{\kern0pt}Current\ syntax\ restriction{\isacharcolon}{\kern0pt}\ threads\ must\ have\ label{\isachardot}{\kern0pt}{\isachardoublequote}{\kern0pt}\isanewline
\isanewline
\isanewline
\isanewline
\isanewline
val\ context{\isacharunderscore}{\kern0pt}check\ {\isacharequal}{\kern0pt}\ fn\ {\isacharparenleft}{\kern0pt}{\isacharparenleft}{\kern0pt}{\isacharparenleft}{\kern0pt}system{\isacharcolon}{\kern0pt}\ binding\ {\isacharcomma}{\kern0pt}\ \isanewline
\ \ \ \ \ \ \ \ \ \ \ \ \ \ \ \ \ \ \ \ \ \ \ \ \ \ globals\ {\isacharcolon}{\kern0pt}\ \ {\isacharparenleft}{\kern0pt}{\isacharparenleft}{\kern0pt}binding\ {\isacharasterisk}{\kern0pt}\ {\isacharparenleft}{\kern0pt}string\ {\isacharasterisk}{\kern0pt}\ Position{\isachardot}{\kern0pt}T{\isacharparenright}{\kern0pt}{\isacharparenright}{\kern0pt}\ {\isacharasterisk}{\kern0pt}\ string\ option{\isacharparenright}{\kern0pt}list{\isacharparenright}{\kern0pt}{\isacharcomma}{\kern0pt}\isanewline
\ \ \ \ \ \ \ \ \ \ \ \ \ \ \ \ \ \ \ \ \ \ \ \ \ \ locks\ {\isacharcolon}{\kern0pt}\ \ {\isacharparenleft}{\kern0pt}binding\ {\isacharasterisk}{\kern0pt}\ {\isacharparenleft}{\kern0pt}string\ {\isacharasterisk}{\kern0pt}\ Position{\isachardot}{\kern0pt}T{\isacharparenright}{\kern0pt}{\isacharparenright}{\kern0pt}\ list{\isacharparenright}{\kern0pt}{\isacharcomma}{\kern0pt}\isanewline
\ \ \ \ \ \ \ \ \ \ \ \ \ \ \ \ \ \ \ \ \ \ \ \ \ \ threads\ {\isacharcolon}{\kern0pt}\ raw{\isacharunderscore}{\kern0pt}thread{\isacharunderscore}{\kern0pt}absy\ list{\isacharparenright}{\kern0pt}\isanewline
\ \ \ \ \ \ \ \ \ \ \ \ \ \ \ \ \ \ \ \ \ \ \ \ \ \ \ {\isacharequal}{\kern0pt}{\isachargreater}{\kern0pt}\isanewline
\ \ \ \ \ \ \ \ \ \ \ \ \ \ \ \ \ \ \ \ \ fn\ thy\ {\isacharcolon}{\kern0pt}\ theory\ {\isacharequal}{\kern0pt}{\isachargreater}{\kern0pt}\ \isanewline
\ \ \ \ \ \ \ \ \ \ \ \ \ \ \ \ \ \ \ \ \ let\ val\ read\ {\isacharequal}{\kern0pt}\ read{\isacharunderscore}{\kern0pt}term{\isacharunderscore}{\kern0pt}err\ thy\isanewline
\ \ \ \ \ \ \ \ \ \ \ \ \ \ \ \ \ \ \ \ \ \ \ \ \ fun\ global{\isacharunderscore}{\kern0pt}decls{\isacharunderscore}{\kern0pt}conv\ {\isacharparenleft}{\kern0pt}{\isacharparenleft}{\kern0pt}bdg{\isacharcomma}{\kern0pt}\ {\isacharparenleft}{\kern0pt}s{\isacharcomma}{\kern0pt}\ pos{\isacharparenright}{\kern0pt}{\isacharparenright}{\kern0pt}{\isacharcomma}{\kern0pt}\ SOME\ init{\isacharunderscore}{\kern0pt}term{\isacharparenright}{\kern0pt}\ {\isacharequal}{\kern0pt}\isanewline
\ \ \ \ \ \ \ \ \ \ \ \ \ \ \ \ \ \ \ \ \ \ \ \ \ \ \ \ let\ \isanewline
\ \ \ \ \ \ \ \ \ \ \ \ \ \ \ \ \ \ \ \ \ \ \ \ \ \ \ \ \ \ val\ typ\ {\isacharequal}{\kern0pt}\ Syntax{\isachardot}{\kern0pt}read{\isacharunderscore}{\kern0pt}typ{\isacharunderscore}{\kern0pt}global\ thy\ s\isanewline
\ \ \ \ \ \ \ \ \ \ \ \ \ \ \ \ \ \ \ \ \ \ \ \ \ \ \ \ \ \ val\ inferred{\isacharunderscore}{\kern0pt}type\ {\isacharequal}{\kern0pt}\ fastype{\isacharunderscore}{\kern0pt}of\ init{\isacharunderscore}{\kern0pt}term\isanewline
\ \ \ \ \ \ \ \ \ \ \ \ \ \ \ \ \ \ \ \ \ \ \ \ \ \ \ \ \ \ val\ ctxt\ {\isacharequal}{\kern0pt}\ Proof{\isacharunderscore}{\kern0pt}Context{\isachardot}{\kern0pt}init{\isacharunderscore}{\kern0pt}global\ thy\isanewline
\ \ \ \ \ \ \ \ \ \ \ \ \ \ \ \ \ \ \ \ \ \ \ \ \ \ \ \ \ \ val\ inferred{\isacharunderscore}{\kern0pt}str\ {\isacharequal}{\kern0pt}\ Syntax{\isachardot}{\kern0pt}string{\isacharunderscore}{\kern0pt}of{\isacharunderscore}{\kern0pt}typ\ ctxt\ inferred{\isacharunderscore}{\kern0pt}type\isanewline
\ \ \ \ \ \ \ \ \ \ \ \ \ \ \ \ \ \ \ \ \ \ \ \ \ \ \ \ \ \ val\ declared{\isacharunderscore}{\kern0pt}str\ {\isacharequal}{\kern0pt}\ Syntax{\isachardot}{\kern0pt}string{\isacharunderscore}{\kern0pt}of{\isacharunderscore}{\kern0pt}typ\ ctxt\ typ\isanewline
\ \ \ \ \ \ \ \ \ \ \ \ \ \ \ \ \ \ \ \ \ \ \ \ \ \ \ \ in\isanewline
\ \ \ \ \ \ \ \ \ \ \ \ \ \ \ \ \ \ \ \ \ \ \ \ \ \ \ \ \ \ if\ inferred{\isacharunderscore}{\kern0pt}type\ {\isacharless}{\kern0pt}{\isachargreater}{\kern0pt}\ typ\ then\isanewline
\ \ \ \ \ \ \ \ \ \ \ \ \ \ \ \ \ \ \ \ \ \ \ \ \ \ \ \ \ \ \ \ error\ {\isacharparenleft}{\kern0pt}{\isachardoublequote}{\kern0pt}Type\ mismatch\ in\ variable\ {\isachardoublequote}{\kern0pt}\ {\isacharcircum}{\kern0pt}\ Binding{\isachardot}{\kern0pt}name{\isacharunderscore}{\kern0pt}of\ bdg\ {\isacharcircum}{\kern0pt}\isanewline
\ \ \ \ \ \ \ \ \ \ \ \ \ \ \ \ \ \ \ \ \ \ \ \ \ \ \ \ \ \ \ \ \ \ \ \ \ \ \ {\isachardoublequote}{\kern0pt}{\isacharcolon}{\kern0pt}\ expected\ {\isachardoublequote}{\kern0pt}\ {\isacharcircum}{\kern0pt}\ declared{\isacharunderscore}{\kern0pt}str\ {\isacharcircum}{\kern0pt}\ {\isachardoublequote}{\kern0pt}{\isacharcomma}{\kern0pt}\ but\ got\ {\isachardoublequote}{\kern0pt}\ {\isacharcircum}{\kern0pt}\ inferred{\isacharunderscore}{\kern0pt}str\ {\isacharcircum}{\kern0pt}\ \ \isanewline
\ \ \ \ \ \ \ \ \ \ \ \ \ \ \ \ \ \ \ \ \ \ \ \ \ \ \ \ \ \ \ \ \ \ \ \ \ \ \ {\isachardoublequote}{\kern0pt}at\ {\isachardoublequote}{\kern0pt}\ {\isacharcircum}{\kern0pt}Position{\isachardot}{\kern0pt}here\ pos{\isacharparenright}{\kern0pt}\isanewline
\ \ \ \ \ \ \ \ \ \ \ \ \ \ \ \ \ \ \ \ \ \ \ \ \ \ \ \ \ \ else\isanewline
\ \ \ \ \ \ \ \ \ \ \ \ \ \ \ \ \ \ \ \ \ \ \ \ \ \ \ \ \ \ \ \ {\isacharparenleft}{\kern0pt}bdg{\isacharcomma}{\kern0pt}\ typ{\isacharcomma}{\kern0pt}\ pos{\isacharcomma}{\kern0pt}\ SOME\ init{\isacharunderscore}{\kern0pt}term{\isacharparenright}{\kern0pt}\isanewline
\ \ \ \ \ \ \ \ \ \ \ \ \ \ \ \ \ \ \ \ \ \ \ \ \ \ \ \ end\isanewline
\ \ \ \ \ \ \ \ \ \ \ \ \ \ \ \ \ \ \ \ \ \ \ \ \ \ {\isacharbar}{\kern0pt}\ global{\isacharunderscore}{\kern0pt}decls{\isacharunderscore}{\kern0pt}conv\ {\isacharparenleft}{\kern0pt}{\isacharparenleft}{\kern0pt}bdg{\isacharcomma}{\kern0pt}\ {\isacharparenleft}{\kern0pt}s{\isacharcomma}{\kern0pt}\ pos{\isacharparenright}{\kern0pt}{\isacharparenright}{\kern0pt}{\isacharcomma}{\kern0pt}\ NONE{\isacharparenright}{\kern0pt}\ {\isacharequal}{\kern0pt}\isanewline
\ \ \ \ \ \ \ \ \ \ \ \ \ \ \ \ \ \ \ \ \ \ \ \ \ \ \ \ {\isacharparenleft}{\kern0pt}bdg{\isacharcomma}{\kern0pt}\ Syntax{\isachardot}{\kern0pt}read{\isacharunderscore}{\kern0pt}typ{\isacharunderscore}{\kern0pt}global\ thy\ s{\isacharcomma}{\kern0pt}\ pos{\isacharcomma}{\kern0pt}\ NONE{\isacharparenright}{\kern0pt}\isanewline
\ \ \ \ \ \ \ \ \ \ \ \ \ \ \ \ \ \ \ \ \ val\ {\isacharparenleft}{\kern0pt}global{\isacharunderscore}{\kern0pt}decls{\isacharprime}{\kern0pt}{\isacharcomma}{\kern0pt}\ thy{\isadigit{3}}{\isacharprime}{\kern0pt}{\isacharparenright}{\kern0pt}\ {\isacharequal}{\kern0pt}\ \isanewline
\ \ \ \ \ \ \ \ \ \ \ \ \ \ \ \ \ \ \ \ \ \ \ \ \ \ \ \ \ \ \ \ \ \ let\isanewline
\ \ \ \ \ \ \ \ \ \ \ \ \ \ \ \ \ \ \ \ \ \ \ \ \ \ \ \ \ \ \ \ \ \ \ \ val\ global{\isacharunderscore}{\kern0pt}terms\ {\isacharequal}{\kern0pt}\ map\ {\isacharparenleft}{\kern0pt}fn\ {\isacharparenleft}{\kern0pt}{\isacharparenleft}{\kern0pt}bdg{\isacharcomma}{\kern0pt}{\isacharparenleft}{\kern0pt}s{\isacharcomma}{\kern0pt}pos{\isacharparenright}{\kern0pt}{\isacharparenright}{\kern0pt}{\isacharcomma}{\kern0pt}\ init{\isacharunderscore}{\kern0pt}opt{\isacharunderscore}{\kern0pt}str{\isacharparenright}{\kern0pt}\ {\isacharequal}{\kern0pt}{\isachargreater}{\kern0pt}\isanewline
\ \ \ \ \ \ \ \ \ \ \ \ \ \ \ \ \ \ \ \ \ \ \ \ \ \ \ \ \ \ \ \ \ \ \ \ \ \ \ let\isanewline
\ \ \ \ \ \ \ \ \ \ \ \ \ \ \ \ \ \ \ \ \ \ \ \ \ \ \ \ \ \ \ \ \ \ \ \ \ \ \ \ \ val\ init{\isacharunderscore}{\kern0pt}opt{\isacharunderscore}{\kern0pt}term\ {\isacharequal}{\kern0pt}\ case\ init{\isacharunderscore}{\kern0pt}opt{\isacharunderscore}{\kern0pt}str\ of\isanewline
\ \ \ \ \ \ \ \ \ \ \ \ \ \ \ \ \ \ \ \ \ \ \ \ \ \ \ \ \ \ \ \ \ \ \ \ \ \ \ \ \ \ \ \ \ NONE\ {\isacharequal}{\kern0pt}{\isachargreater}{\kern0pt}\ NONE\isanewline
\ \ \ \ \ \ \ \ \ \ \ \ \ \ \ \ \ \ \ \ \ \ \ \ \ \ \ \ \ \ \ \ \ \ \ \ \ \ \ \ \ \ \ {\isacharbar}{\kern0pt}\ SOME\ str\ {\isacharequal}{\kern0pt}{\isachargreater}{\kern0pt}\ SOME\ {\isacharparenleft}{\kern0pt}read{\isacharunderscore}{\kern0pt}term{\isacharunderscore}{\kern0pt}err\ thy\ {\isachardoublequote}{\kern0pt}local\ init{\isachardoublequote}{\kern0pt}\ bdg\ str{\isacharparenright}{\kern0pt}\isanewline
\ \ \ \ \ \ \ \ \ \ \ \ \ \ \ \ \ \ \ \ \ \ \ \ \ \ \ \ \ \ \ \ \ \ \ \ \ \ \ in\isanewline
\ \ \ \ \ \ \ \ \ \ \ \ \ \ \ \ \ \ \ \ \ \ \ \ \ \ \ \ \ \ \ \ \ \ \ \ \ \ \ \ \ {\isacharparenleft}{\kern0pt}{\isacharparenleft}{\kern0pt}bdg{\isacharcomma}{\kern0pt}\ {\isacharparenleft}{\kern0pt}s{\isacharcomma}{\kern0pt}pos{\isacharparenright}{\kern0pt}{\isacharparenright}{\kern0pt}{\isacharcomma}{\kern0pt}\ init{\isacharunderscore}{\kern0pt}opt{\isacharunderscore}{\kern0pt}term{\isacharparenright}{\kern0pt}\isanewline
\ \ \ \ \ \ \ \ \ \ \ \ \ \ \ \ \ \ \ \ \ \ \ \ \ \ \ \ \ \ \ \ \ \ \ \ \ \ \ end{\isacharparenright}{\kern0pt}\ globals\isanewline
\ \ \ \ \ \ \ \ \ \ \ \ \ \ \ \ \ \ \ \ \ \ \ \ \ \ \ \ \isanewline
\ \ \ \ \ \ \ \ \ \ \ \ \ \ \ \ \ \ \ \ \ \ \ \ \ \ \ \ \ \ \ \ \ \ \ \ val\ decls\ {\isacharequal}{\kern0pt}\ map\ global{\isacharunderscore}{\kern0pt}decls{\isacharunderscore}{\kern0pt}conv\ global{\isacharunderscore}{\kern0pt}terms\isanewline
\ \ \ \ \ \ \ \ \ \ \ \ \ \ \ \ \ \ \ \ \ \ \ \ \ \ \ \ \ \ \ \ \ \ \ \ val\ thy{\isadigit{4}}\ {\isacharequal}{\kern0pt}\ Sign{\isachardot}{\kern0pt}add{\isacharunderscore}{\kern0pt}consts\ {\isacharparenleft}{\kern0pt}map\ pre{\isacharunderscore}{\kern0pt}const{\isadigit{2}}\ decls{\isacharparenright}{\kern0pt}\ thy\isanewline
\ \ \ \ \ \ \ \ \ \ \ \ \ \ \ \ \ \ \ \ \ \ \ \ \ \ \ \ \ \ \ \ \ \ in\isanewline
\ \ \ \ \ \ \ \ \ \ \ \ \ \ \ \ \ \ \ \ \ \ \ \ \ \ \ \ \ \ \ \ \ \ \ \ {\isacharparenleft}{\kern0pt}decls{\isacharcomma}{\kern0pt}\ thy{\isadigit{4}}{\isacharparenright}{\kern0pt}\isanewline
\ \ \ \ \ \ \ \ \ \ \ \ \ \ \ \ \ \ \ \ \ \ \ \ \ \ \ \ \ \ \ \ \ \ end\isanewline
\isanewline
\ \ \ \ \ \ \ \ \ \ \ \ \ \ \ \ \ \ \ \ \ \ \ \ \ fun\ conv{\isacharunderscore}{\kern0pt}decl{\isacharunderscore}{\kern0pt}init\ {\isacharparenleft}{\kern0pt}n{\isacharcomma}{\kern0pt}{\isacharunderscore}{\kern0pt}{\isacharcomma}{\kern0pt}\ {\isacharunderscore}{\kern0pt}\ \ {\isacharcomma}{\kern0pt}\ init{\isacharunderscore}{\kern0pt}opt{\isacharparenright}{\kern0pt}\ {\isacharequal}{\kern0pt}\ {\isacharparenleft}{\kern0pt}Binding{\isachardot}{\kern0pt}name{\isacharunderscore}{\kern0pt}of\ n{\isacharcomma}{\kern0pt}\ init{\isacharunderscore}{\kern0pt}opt{\isacharparenright}{\kern0pt}\isanewline
\ \ \ \ \ \ \ \ \ \ \ \ \ \ \ \ \ \ \ \ \ \ \ \ \ val\ varstab\ \ \ \ \ {\isacharequal}{\kern0pt}\ Symtab{\isachardot}{\kern0pt}make\ {\isacharparenleft}{\kern0pt}map\ conv{\isacharunderscore}{\kern0pt}decl{\isacharunderscore}{\kern0pt}init\ global{\isacharunderscore}{\kern0pt}decls{\isacharprime}{\kern0pt}{\isacharparenright}{\kern0pt}\isanewline
\ \ \ \ \ \ \ \ \ \ \ \ \ \ \ \ \ \ \ \ \ \ \ \ \ val\ lock{\isacharunderscore}{\kern0pt}decls{\isacharprime}{\kern0pt}{\isacharequal}{\kern0pt}\ map\ {\isacharparenleft}{\kern0pt}fn\ {\isacharparenleft}{\kern0pt}b{\isacharcomma}{\kern0pt}\ {\isacharparenleft}{\kern0pt}str{\isacharcomma}{\kern0pt}pos{\isacharparenright}{\kern0pt}{\isacharparenright}{\kern0pt}\ {\isacharequal}{\kern0pt}{\isachargreater}{\kern0pt}\ {\isacharparenleft}{\kern0pt}b{\isacharcomma}{\kern0pt}\ read\ {\isachardoublequote}{\kern0pt}lock{\isacharunderscore}{\kern0pt}decl{\isachardoublequote}{\kern0pt}\ b\ str{\isacharcomma}{\kern0pt}\ pos{\isacharparenright}{\kern0pt}{\isacharparenright}{\kern0pt}\ locks\isanewline
\ \ \ \ \ \ \ \ \ \ \ \ \ \ \ \ \ \ \ \ \ \ \ \ \ val\ thy{\isacharprime}{\kern0pt}\ \ \ \ \ \ \ \ {\isacharequal}{\kern0pt}\ Sign{\isachardot}{\kern0pt}add{\isacharunderscore}{\kern0pt}consts\ {\isacharparenleft}{\kern0pt}map\ pre{\isacharunderscore}{\kern0pt}const{\isadigit{2}}\ global{\isacharunderscore}{\kern0pt}decls{\isacharprime}{\kern0pt}{\isacharparenright}{\kern0pt}\ thy\isanewline
\ \ \ \ \ \ \ \ \ \ \ \ \ \ \ \ \ \ \ \ \ \ \ \ \ val\ thy{\isacharprime}{\kern0pt}{\isacharprime}{\kern0pt}\ \ \ \ \ \ \ {\isacharequal}{\kern0pt}\ thy{\isacharprime}{\kern0pt}\ {\isacharbar}{\kern0pt}{\isachargreater}{\kern0pt}\ Sign{\isachardot}{\kern0pt}add{\isacharunderscore}{\kern0pt}consts\ {\isacharparenleft}{\kern0pt}map\ pre{\isacharunderscore}{\kern0pt}event\ lock{\isacharunderscore}{\kern0pt}decls{\isacharprime}{\kern0pt}{\isacharparenright}{\kern0pt}\isanewline
\isanewline
\ \ \ \ \ \ \ \ \ \ \ \ \ \ \ \ \ \ \ \ \ \ \ \ \ val\ threads{\isacharprime}{\kern0pt}\ {\isacharequal}{\kern0pt}\ map\ {\isacharparenleft}{\kern0pt}check{\isacharunderscore}{\kern0pt}thread\ thy\ thy{\isacharprime}{\kern0pt}{\isacharprime}{\kern0pt}\ varstab{\isacharparenright}{\kern0pt}\ threads\ \isanewline
\ \ \ \ \ \ \ \ \ \ \ \ \ \ \ \ \ \ \ \ \ \ \ \ val\ res\ {\isacharequal}{\kern0pt}\ {\isacharbraceleft}{\kern0pt}system\ {\isacharequal}{\kern0pt}\ system{\isacharcomma}{\kern0pt}\isanewline
\ \ \ \ \ \ \ \ \ \ \ \ \ \ \ \ \ \ \ \ \ \ \ \ \ \ \ \ \ \ \ \ \ \ \ globals{\isacharunderscore}{\kern0pt}decls\ {\isacharequal}{\kern0pt}\ global{\isacharunderscore}{\kern0pt}decls{\isacharprime}{\kern0pt}{\isacharcomma}{\kern0pt}\isanewline
\ \ \ \ \ \ \ \ \ \ \ \ \ \ \ \ \ \ \ \ \ \ \ \ \ \ \ \ \ \ \ \ \ \ \ varstab\ {\isacharequal}{\kern0pt}\ varstab{\isacharcomma}{\kern0pt}\isanewline
\ \ \ \ \ \ \ \ \ \ \ \ \ \ \ \ \ \ \ \ \ \ \ \ \ \ \ \ \ \ \ \ \ \ \ locks{\isacharunderscore}{\kern0pt}decls\ {\isacharequal}{\kern0pt}\ lock{\isacharunderscore}{\kern0pt}decls{\isacharprime}{\kern0pt}{\isacharcomma}{\kern0pt}\isanewline
\ \ \ \ \ \ \ \ \ \ \ \ \ \ \ \ \ \ \ \ \ \ \ \ \ \ \ \ \ \ \ \ \ \ \ threads{\isacharunderscore}{\kern0pt}decls\ {\isacharequal}{\kern0pt}\ threads{\isacharprime}{\kern0pt}{\isacharbraceright}{\kern0pt}\isanewline
\ \ \ \ \ \ \ \ \ \ \ \ \ \ \ \ \ \ \ \ \ in\ \ SPY{\isacharcolon}{\kern0pt}{\isacharequal}{\kern0pt}{\isacharparenleft}{\kern0pt}SOME\ {\isacharparenleft}{\kern0pt}res{\isacharparenright}{\kern0pt}{\isacharparenright}{\kern0pt}{\isacharsemicolon}{\kern0pt}\ res\ end\isanewline
\isanewline
\isanewline
{\isacartoucheclose}%
\endisatagML
{\isafoldML}%
%
\isadelimML
%
\endisadelimML
%
\isadelimdocument
%
\endisadelimdocument
%
\isatagdocument
%
\isamarkupsection{Semantic Check%
}
\isamarkuptrue%
%
\endisatagdocument
{\isafolddocument}%
%
\isadelimdocument
%
\endisadelimdocument
%
\isadelimML
%
\endisadelimML
%
\isatagML
\isakeywordONE{ML}\isamarkupfalse%
{\isacartoucheopen}\isanewline
{\isacharparenleft}{\kern0pt}{\isacharasterisk}{\kern0pt}\ Redéfinir\ les\ types\ de\ base\ de\ la\ sémantique\ {\isacharasterisk}{\kern0pt}{\isacharparenright}{\kern0pt}\isanewline
type\ SV\ {\isacharequal}{\kern0pt}\ string\ \ {\isacharparenleft}{\kern0pt}{\isacharasterisk}{\kern0pt}\ Shared\ Variables\ {\isacharasterisk}{\kern0pt}{\isacharparenright}{\kern0pt}\isanewline
type\ V\ {\isacharequal}{\kern0pt}\ string\ \ \ {\isacharparenleft}{\kern0pt}{\isacharasterisk}{\kern0pt}\ Thread{\isacharminus}{\kern0pt}Local\ Variables\ {\isacharasterisk}{\kern0pt}{\isacharparenright}{\kern0pt}\ \ \isanewline
type\ MV\ {\isacharequal}{\kern0pt}\ int\ \ \ \ \ {\isacharparenleft}{\kern0pt}{\isacharasterisk}{\kern0pt}\ Mutex\ IDs\ {\isacharasterisk}{\kern0pt}{\isacharparenright}{\kern0pt}\isanewline
type\ D\ {\isacharequal}{\kern0pt}\ int\ \ \ \ \ \ {\isacharparenleft}{\kern0pt}{\isacharasterisk}{\kern0pt}\ Data\ type\ {\isacharminus}{\kern0pt}\ simplifié\ à\ int\ {\isacharasterisk}{\kern0pt}{\isacharparenright}{\kern0pt}\isanewline
\isanewline
type\ state\ {\isacharequal}{\kern0pt}\ string\ {\isacharminus}{\kern0pt}{\isachargreater}{\kern0pt}\ int\isanewline
\isanewline
fun\ safe{\isacharunderscore}{\kern0pt}lookup\ var\ state\ {\isacharequal}{\kern0pt}\isanewline
\ \ case\ state\ var\ of\isanewline
\ \ \ \ {\isachartilde}{\kern0pt}{\isadigit{1}}\ {\isacharequal}{\kern0pt}{\isachargreater}{\kern0pt}\ raise\ Undeclared{\isacharunderscore}{\kern0pt}Variable\ var\isanewline
\ \ {\isacharbar}{\kern0pt}\ v\ {\isacharequal}{\kern0pt}{\isachargreater}{\kern0pt}\ v\isanewline
\isanewline
type\ arith{\isacharunderscore}{\kern0pt}expr\ {\isacharequal}{\kern0pt}\ state\ {\isacharminus}{\kern0pt}{\isachargreater}{\kern0pt}\ int\isanewline
type\ bool{\isacharunderscore}{\kern0pt}expr\ {\isacharequal}{\kern0pt}\ state\ {\isacharminus}{\kern0pt}{\isachargreater}{\kern0pt}\ bool\isanewline
\isanewline
datatype\ com\ {\isacharequal}{\kern0pt}\ \isanewline
\ \ \ \ SKIP\isanewline
\ \ {\isacharbar}{\kern0pt}\ Assign\ of\ string\ {\isacharasterisk}{\kern0pt}\ arith{\isacharunderscore}{\kern0pt}expr\isanewline
\ \ {\isacharbar}{\kern0pt}\ Seq\ of\ com\ {\isacharasterisk}{\kern0pt}\ com\ \ \isanewline
\ \ {\isacharbar}{\kern0pt}\ Cond\ of\ bool{\isacharunderscore}{\kern0pt}expr\ {\isacharasterisk}{\kern0pt}\ com\ {\isacharasterisk}{\kern0pt}\ com\isanewline
\ \ {\isacharbar}{\kern0pt}\ While\ of\ bool{\isacharunderscore}{\kern0pt}expr\ {\isacharasterisk}{\kern0pt}\ com\isanewline
\ \ {\isacharbar}{\kern0pt}\ Lock\ of\ string\isanewline
\ \ {\isacharbar}{\kern0pt}\ Unlock\ of\ string\isanewline
\ \ {\isacharbar}{\kern0pt}\ Store\ of\ arith{\isacharunderscore}{\kern0pt}expr\ {\isacharasterisk}{\kern0pt}\ string\isanewline
\ \ {\isacharbar}{\kern0pt}\ Load\ of\ string\ {\isacharasterisk}{\kern0pt}\ string\isanewline
\isanewline
{\isacharparenleft}{\kern0pt}{\isacharasterisk}{\kern0pt}\ Table\ de\ mapping\ pour\ les\ locks\ {\isacharasterisk}{\kern0pt}{\isacharparenright}{\kern0pt}\isanewline
val\ lock{\isacharunderscore}{\kern0pt}table\ {\isacharequal}{\kern0pt}\ Unsynchronized{\isachardot}{\kern0pt}ref\ {\isacharparenleft}{\kern0pt}Symtab{\isachardot}{\kern0pt}empty\ {\isacharcolon}{\kern0pt}\ int\ Symtab{\isachardot}{\kern0pt}table{\isacharparenright}{\kern0pt}\isanewline
val\ lock{\isacharunderscore}{\kern0pt}counter\ {\isacharequal}{\kern0pt}\ Unsynchronized{\isachardot}{\kern0pt}ref\ {\isadigit{1}}\isanewline
{\isacharparenleft}{\kern0pt}{\isacharasterisk}{\kern0pt}\ \ \ \ \ \ \ \ \ \ \ \ \ \ \ \ \ \isanewline
type\ absy\ {\isacharequal}{\kern0pt}\ {\isacharbraceleft}{\kern0pt}system\ \ \ \ \ \ \ \ \ \ {\isacharcolon}{\kern0pt}\ binding{\isacharcomma}{\kern0pt}\isanewline
\ \ \ \ \ \ \ \ \ \ \ \ \ globals{\isacharunderscore}{\kern0pt}decls\ \ \ {\isacharcolon}{\kern0pt}\ {\isacharparenleft}{\kern0pt}binding\ {\isacharasterisk}{\kern0pt}\ typ\ {\isacharasterisk}{\kern0pt}\ Position{\isachardot}{\kern0pt}T\ {\isacharasterisk}{\kern0pt}\ term\ option{\isacharparenright}{\kern0pt}\ list{\isacharcomma}{\kern0pt}\isanewline
\ \ \ \ \ \ \ \ \ \ \ \ \ varstab\ \ \ \ \ \ \ \ \ {\isacharcolon}{\kern0pt}\ {\isacharparenleft}{\kern0pt}term\ option{\isacharparenright}{\kern0pt}\ Symtab{\isachardot}{\kern0pt}table{\isacharcomma}{\kern0pt}\isanewline
\ \ \ \ \ \ \ \ \ \ \ \ \ locks{\isacharunderscore}{\kern0pt}decls\ \ \ \ \ {\isacharcolon}{\kern0pt}\ {\isacharparenleft}{\kern0pt}binding\ {\isacharasterisk}{\kern0pt}\ term\ {\isacharasterisk}{\kern0pt}\ Position{\isachardot}{\kern0pt}T{\isacharparenright}{\kern0pt}\ list{\isacharcomma}{\kern0pt}\isanewline
\ \ \ \ \ \ \ \ \ \ \ \ \ threads{\isacharunderscore}{\kern0pt}decls\ \ \ {\isacharcolon}{\kern0pt}\ thread{\isacharunderscore}{\kern0pt}absy\ list\isanewline
\ \ \ \ \ \ \ \ \ \ \ \ {\isacharbraceright}{\kern0pt}\isanewline
\isanewline
fun\ get{\isacharunderscore}{\kern0pt}lock{\isacharunderscore}{\kern0pt}id\ {\isacharparenleft}{\kern0pt}lock{\isacharunderscore}{\kern0pt}name\ {\isacharcolon}{\kern0pt}\ string{\isacharparenright}{\kern0pt}\ {\isacharcolon}{\kern0pt}\ int\ {\isacharequal}{\kern0pt}\isanewline
\ \ case\ Symtab{\isachardot}{\kern0pt}lookup\ {\isacharparenleft}{\kern0pt}{\isacharbang}{\kern0pt}lock{\isacharunderscore}{\kern0pt}table{\isacharparenright}{\kern0pt}\ lock{\isacharunderscore}{\kern0pt}name\ of\isanewline
\ \ \ \ SOME\ id\ {\isacharequal}{\kern0pt}{\isachargreater}{\kern0pt}\ id\isanewline
\ \ {\isacharbar}{\kern0pt}\ NONE\ {\isacharequal}{\kern0pt}{\isachargreater}{\kern0pt}\ \isanewline
\ \ \ \ \ \ let\ \isanewline
\ \ \ \ \ \ \ \ val\ new{\isacharunderscore}{\kern0pt}id\ {\isacharequal}{\kern0pt}\ {\isacharbang}{\kern0pt}lock{\isacharunderscore}{\kern0pt}counter\isanewline
\ \ \ \ \ \ \ \ val\ {\isacharunderscore}{\kern0pt}\ {\isacharequal}{\kern0pt}\ lock{\isacharunderscore}{\kern0pt}counter\ {\isacharcolon}{\kern0pt}{\isacharequal}{\kern0pt}\ new{\isacharunderscore}{\kern0pt}id\ {\isacharplus}{\kern0pt}\ {\isadigit{1}}\isanewline
\ \ \ \ \ \ \ \ val\ {\isacharunderscore}{\kern0pt}\ {\isacharequal}{\kern0pt}\ lock{\isacharunderscore}{\kern0pt}table\ {\isacharcolon}{\kern0pt}{\isacharequal}{\kern0pt}\ Symtab{\isachardot}{\kern0pt}update\ {\isacharparenleft}{\kern0pt}lock{\isacharunderscore}{\kern0pt}name{\isacharcomma}{\kern0pt}\ new{\isacharunderscore}{\kern0pt}id{\isacharparenright}{\kern0pt}\ {\isacharparenleft}{\kern0pt}{\isacharbang}{\kern0pt}lock{\isacharunderscore}{\kern0pt}table{\isacharparenright}{\kern0pt}\isanewline
\ \ \ \ \ \ in\ \isanewline
\ \ \ \ \ \ \ \ new{\isacharunderscore}{\kern0pt}id\ \isanewline
\ \ \ \ \ \ end\isanewline
{\isacharasterisk}{\kern0pt}{\isacharparenright}{\kern0pt}\isanewline
\isanewline
fun\ declared{\isacharunderscore}{\kern0pt}lock\ {\isacharparenleft}{\kern0pt}name{\isacharcolon}{\kern0pt}\ string{\isacharparenright}{\kern0pt}\ {\isacharparenleft}{\kern0pt}machine{\isacharcolon}{\kern0pt}\ absy{\isacharparenright}{\kern0pt}\ {\isacharcolon}{\kern0pt}\ bool\ {\isacharequal}{\kern0pt}\isanewline
\ \ List{\isachardot}{\kern0pt}exists\ {\isacharparenleft}{\kern0pt}fn\ {\isacharparenleft}{\kern0pt}b{\isacharcomma}{\kern0pt}\ {\isacharunderscore}{\kern0pt}{\isacharcomma}{\kern0pt}\ {\isacharunderscore}{\kern0pt}{\isacharparenright}{\kern0pt}\ {\isacharequal}{\kern0pt}{\isachargreater}{\kern0pt}\ Binding{\isachardot}{\kern0pt}name{\isacharunderscore}{\kern0pt}of\ b\ {\isacharequal}{\kern0pt}\ name{\isacharparenright}{\kern0pt}\ {\isacharparenleft}{\kern0pt}{\isacharhash}{\kern0pt}locks{\isacharunderscore}{\kern0pt}decls\ machine{\isacharparenright}{\kern0pt}\isanewline
\isanewline
fun\ get{\isacharunderscore}{\kern0pt}lock{\isacharunderscore}{\kern0pt}id\ {\isacharparenleft}{\kern0pt}lock{\isacharunderscore}{\kern0pt}name{\isacharcolon}{\kern0pt}\ string{\isacharparenright}{\kern0pt}\ {\isacharparenleft}{\kern0pt}machine{\isacharcolon}{\kern0pt}\ absy{\isacharparenright}{\kern0pt}\ {\isacharcolon}{\kern0pt}\ int\ {\isacharequal}{\kern0pt}\isanewline
\ \ if\ not\ {\isacharparenleft}{\kern0pt}declared{\isacharunderscore}{\kern0pt}lock\ lock{\isacharunderscore}{\kern0pt}name\ machine{\isacharparenright}{\kern0pt}\ then\isanewline
\ \ \ \ error\ {\isacharparenleft}{\kern0pt}{\isachardoublequote}{\kern0pt}Lock\ {\isachardoublequote}{\kern0pt}\ {\isacharcircum}{\kern0pt}\ lock{\isacharunderscore}{\kern0pt}name\ {\isacharcircum}{\kern0pt}\ {\isachardoublequote}{\kern0pt}\ not\ declared{\isachardoublequote}{\kern0pt}{\isacharparenright}{\kern0pt}\isanewline
\ \ else\isanewline
\ \ \ \ case\ Symtab{\isachardot}{\kern0pt}lookup\ {\isacharparenleft}{\kern0pt}{\isacharbang}{\kern0pt}lock{\isacharunderscore}{\kern0pt}table{\isacharparenright}{\kern0pt}\ lock{\isacharunderscore}{\kern0pt}name\ of\isanewline
\ \ \ \ \ \ \ \ SOME\ id\ {\isacharequal}{\kern0pt}{\isachargreater}{\kern0pt}\ id\isanewline
\ \ \ \ \ \ {\isacharbar}{\kern0pt}\ NONE\ {\isacharequal}{\kern0pt}{\isachargreater}{\kern0pt}\isanewline
\ \ \ \ \ \ \ \ \ \ let\isanewline
\ \ \ \ \ \ \ \ \ \ \ \ val\ new{\isacharunderscore}{\kern0pt}id\ {\isacharequal}{\kern0pt}\ {\isacharbang}{\kern0pt}lock{\isacharunderscore}{\kern0pt}counter\isanewline
\ \ \ \ \ \ \ \ \ \ \ \ val\ {\isacharunderscore}{\kern0pt}\ {\isacharequal}{\kern0pt}\ lock{\isacharunderscore}{\kern0pt}counter\ {\isacharcolon}{\kern0pt}{\isacharequal}{\kern0pt}\ new{\isacharunderscore}{\kern0pt}id\ {\isacharplus}{\kern0pt}\ {\isadigit{1}}\isanewline
\ \ \ \ \ \ \ \ \ \ \ \ val\ {\isacharunderscore}{\kern0pt}\ {\isacharequal}{\kern0pt}\ lock{\isacharunderscore}{\kern0pt}table\ {\isacharcolon}{\kern0pt}{\isacharequal}{\kern0pt}\ Symtab{\isachardot}{\kern0pt}update\ {\isacharparenleft}{\kern0pt}lock{\isacharunderscore}{\kern0pt}name{\isacharcomma}{\kern0pt}\ new{\isacharunderscore}{\kern0pt}id{\isacharparenright}{\kern0pt}\ {\isacharparenleft}{\kern0pt}{\isacharbang}{\kern0pt}lock{\isacharunderscore}{\kern0pt}table{\isacharparenright}{\kern0pt}\isanewline
\ \ \ \ \ \ \ \ \ \ in\isanewline
\ \ \ \ \ \ \ \ \ \ \ \ new{\isacharunderscore}{\kern0pt}id\isanewline
\ \ \ \ \ \ \ \ \ \ end{\isacharsemicolon}{\kern0pt}\isanewline
{\isacharparenleft}{\kern0pt}{\isacharasterisk}{\kern0pt}\ Évaluation\ des\ termes\ Isabelle{\isacharslash}{\kern0pt}HOL\ {\isacharasterisk}{\kern0pt}{\isacharparenright}{\kern0pt}\isanewline
fun\ eval{\isacharunderscore}{\kern0pt}term{\isacharunderscore}{\kern0pt}to{\isacharunderscore}{\kern0pt}int\ {\isacharparenleft}{\kern0pt}thy\ {\isacharcolon}{\kern0pt}\ theory{\isacharparenright}{\kern0pt}\ {\isacharparenleft}{\kern0pt}term\ {\isacharcolon}{\kern0pt}\ term{\isacharparenright}{\kern0pt}\ {\isacharparenleft}{\kern0pt}state\ {\isacharcolon}{\kern0pt}\ state{\isacharparenright}{\kern0pt}\ {\isacharcolon}{\kern0pt}\ int\ {\isacharequal}{\kern0pt}\isanewline
\ \ let\isanewline
\ \ \ \ {\isacharparenleft}{\kern0pt}{\isacharasterisk}{\kern0pt}\ Fonction\ récursive\ pour\ évaluer\ les\ termes\ {\isacharasterisk}{\kern0pt}{\isacharparenright}{\kern0pt}\isanewline
\ \ \ \ fun\ eval\ t\ {\isacharequal}{\kern0pt}\isanewline
\ \ \ \ \ \ case\ t\ of\isanewline
\ \ \ \ \ \ \ \ Const\ {\isacharparenleft}{\kern0pt}{\isacharat}{\kern0pt}{\isacharbraceleft}{\kern0pt}const{\isacharunderscore}{\kern0pt}name\ {\isachardoublequote}{\kern0pt}Groups{\isachardot}{\kern0pt}plus{\isacharunderscore}{\kern0pt}class{\isachardot}{\kern0pt}plus{\isachardoublequote}{\kern0pt}{\isacharbraceright}{\kern0pt}{\isacharcomma}{\kern0pt}\ {\isacharunderscore}{\kern0pt}{\isacharparenright}{\kern0pt}\ {\isachardollar}{\kern0pt}\ t{\isadigit{1}}\ {\isachardollar}{\kern0pt}\ t{\isadigit{2}}\ {\isacharequal}{\kern0pt}{\isachargreater}{\kern0pt}\isanewline
\ \ \ \ \ \ \ \ \ \ {\isacharparenleft}{\kern0pt}eval\ t{\isadigit{1}}{\isacharparenright}{\kern0pt}\ {\isacharplus}{\kern0pt}\ {\isacharparenleft}{\kern0pt}eval\ t{\isadigit{2}}{\isacharparenright}{\kern0pt}\isanewline
\ \ \ \ \ \ {\isacharbar}{\kern0pt}\ Const\ {\isacharparenleft}{\kern0pt}{\isacharat}{\kern0pt}{\isacharbraceleft}{\kern0pt}const{\isacharunderscore}{\kern0pt}name\ {\isachardoublequote}{\kern0pt}Groups{\isachardot}{\kern0pt}minus{\isacharunderscore}{\kern0pt}class{\isachardot}{\kern0pt}minus{\isachardoublequote}{\kern0pt}{\isacharbraceright}{\kern0pt}{\isacharcomma}{\kern0pt}\ {\isacharunderscore}{\kern0pt}{\isacharparenright}{\kern0pt}\ {\isachardollar}{\kern0pt}\ t{\isadigit{1}}\ {\isachardollar}{\kern0pt}\ t{\isadigit{2}}\ {\isacharequal}{\kern0pt}{\isachargreater}{\kern0pt}\isanewline
\ \ \ \ \ \ \ \ \ \ {\isacharparenleft}{\kern0pt}eval\ t{\isadigit{1}}{\isacharparenright}{\kern0pt}\ {\isacharminus}{\kern0pt}\ {\isacharparenleft}{\kern0pt}eval\ t{\isadigit{2}}{\isacharparenright}{\kern0pt}\isanewline
\ \ \ \ \ \ {\isacharbar}{\kern0pt}\ Const\ {\isacharparenleft}{\kern0pt}{\isacharat}{\kern0pt}{\isacharbraceleft}{\kern0pt}const{\isacharunderscore}{\kern0pt}name\ {\isachardoublequote}{\kern0pt}Groups{\isachardot}{\kern0pt}times{\isacharunderscore}{\kern0pt}class{\isachardot}{\kern0pt}times{\isachardoublequote}{\kern0pt}{\isacharbraceright}{\kern0pt}{\isacharcomma}{\kern0pt}\ {\isacharunderscore}{\kern0pt}{\isacharparenright}{\kern0pt}\ {\isachardollar}{\kern0pt}\ t{\isadigit{1}}\ {\isachardollar}{\kern0pt}\ t{\isadigit{2}}\ {\isacharequal}{\kern0pt}{\isachargreater}{\kern0pt}\isanewline
\ \ \ \ \ \ \ \ \ \ {\isacharparenleft}{\kern0pt}eval\ t{\isadigit{1}}{\isacharparenright}{\kern0pt}\ {\isacharasterisk}{\kern0pt}\ {\isacharparenleft}{\kern0pt}eval\ t{\isadigit{2}}{\isacharparenright}{\kern0pt}\isanewline
\ \ \ \ \ \ {\isacharbar}{\kern0pt}\ Const\ {\isacharparenleft}{\kern0pt}{\isacharat}{\kern0pt}{\isacharbraceleft}{\kern0pt}const{\isacharunderscore}{\kern0pt}name\ {\isachardoublequote}{\kern0pt}HOL{\isachardot}{\kern0pt}Trueprop{\isachardoublequote}{\kern0pt}{\isacharbraceright}{\kern0pt}{\isacharcomma}{\kern0pt}\ {\isacharunderscore}{\kern0pt}{\isacharparenright}{\kern0pt}\ {\isachardollar}{\kern0pt}\ t{\isadigit{1}}\ {\isacharequal}{\kern0pt}{\isachargreater}{\kern0pt}\ eval\ t{\isadigit{1}}\isanewline
\ \ \ \ \ \ {\isacharbar}{\kern0pt}\ {\isacharat}{\kern0pt}{\isacharbraceleft}{\kern0pt}term\ {\isachardoublequote}{\kern0pt}{\isadigit{0}}{\isacharcolon}{\kern0pt}{\isacharcolon}{\kern0pt}nat{\isachardoublequote}{\kern0pt}{\isacharbraceright}{\kern0pt}\ {\isacharequal}{\kern0pt}{\isachargreater}{\kern0pt}\ {\isadigit{0}}\isanewline
\ \ \ \ \ \ {\isacharbar}{\kern0pt}\ {\isacharat}{\kern0pt}{\isacharbraceleft}{\kern0pt}term\ {\isachardoublequote}{\kern0pt}{\isadigit{1}}{\isacharcolon}{\kern0pt}{\isacharcolon}{\kern0pt}nat{\isachardoublequote}{\kern0pt}{\isacharbraceright}{\kern0pt}\ {\isacharequal}{\kern0pt}{\isachargreater}{\kern0pt}\ {\isadigit{1}}\isanewline
\ \ \ \ \ \ {\isacharbar}{\kern0pt}\ Free\ {\isacharparenleft}{\kern0pt}name{\isacharcomma}{\kern0pt}\ {\isacharunderscore}{\kern0pt}{\isacharparenright}{\kern0pt}\ {\isacharequal}{\kern0pt}{\isachargreater}{\kern0pt}\ safe{\isacharunderscore}{\kern0pt}lookup\ name\ state\isanewline
\ \ \ \ \ \ {\isacharbar}{\kern0pt}\ Bound\ i\ {\isacharequal}{\kern0pt}{\isachargreater}{\kern0pt}\ \ i\ {\isacharparenleft}{\kern0pt}{\isacharasterisk}{\kern0pt}\ Placeholder\ pour\ les\ variables\ liées\ {\isacharasterisk}{\kern0pt}{\isacharparenright}{\kern0pt}\isanewline
\ \ \ \ \ \ {\isacharbar}{\kern0pt}\ {\isacharunderscore}{\kern0pt}\ {\isacharequal}{\kern0pt}{\isachargreater}{\kern0pt}\ \isanewline
\ \ \ \ \ \ \ \ \ \ {\isacharparenleft}{\kern0pt}{\isacharasterisk}{\kern0pt}\ Essayer\ d{\isacharprime}{\kern0pt}extraire\ les\ constantes\ numériques\ {\isacharasterisk}{\kern0pt}{\isacharparenright}{\kern0pt}\isanewline
\ \ \ \ \ \ \ \ \ \ {\isacharparenleft}{\kern0pt}case\ try\ {\isacharparenleft}{\kern0pt}HOLogic{\isachardot}{\kern0pt}dest{\isacharunderscore}{\kern0pt}number\ {\isacharhash}{\kern0pt}{\isachargreater}{\kern0pt}\ snd{\isacharparenright}{\kern0pt}\ t\ of\isanewline
\ \ \ \ \ \ \ \ \ \ \ \ SOME\ n\ {\isacharequal}{\kern0pt}{\isachargreater}{\kern0pt}\ n\isanewline
\ \ \ \ \ \ \ \ \ \ {\isacharbar}{\kern0pt}\ NONE\ {\isacharequal}{\kern0pt}{\isachargreater}{\kern0pt}\ {\isadigit{0}}{\isacharparenright}{\kern0pt}\ {\isacharparenleft}{\kern0pt}{\isacharasterisk}{\kern0pt}\ Valeur\ par\ défaut\ {\isacharasterisk}{\kern0pt}{\isacharparenright}{\kern0pt}\isanewline
\ \ in\isanewline
\ \ \ \ eval\ term\isanewline
\ \ end\isanewline
\isanewline
fun\ eval{\isacharunderscore}{\kern0pt}term{\isacharunderscore}{\kern0pt}to{\isacharunderscore}{\kern0pt}bool\ {\isacharparenleft}{\kern0pt}thy\ {\isacharcolon}{\kern0pt}\ theory{\isacharparenright}{\kern0pt}\ {\isacharparenleft}{\kern0pt}term\ {\isacharcolon}{\kern0pt}\ term{\isacharparenright}{\kern0pt}\ {\isacharparenleft}{\kern0pt}state\ {\isacharcolon}{\kern0pt}\ state{\isacharparenright}{\kern0pt}\ {\isacharcolon}{\kern0pt}\ bool\ {\isacharequal}{\kern0pt}\isanewline
\ \ let\isanewline
\ \ \ \ fun\ eval\ t\ {\isacharequal}{\kern0pt}\isanewline
\ \ \ \ \ \ case\ t\ of\isanewline
\ \ \ \ \ \ \ \ {\isacharat}{\kern0pt}{\isacharbraceleft}{\kern0pt}term\ {\isachardoublequote}{\kern0pt}True{\isachardoublequote}{\kern0pt}{\isacharbraceright}{\kern0pt}\ {\isacharequal}{\kern0pt}{\isachargreater}{\kern0pt}\ true\isanewline
\ \ \ \ \ \ {\isacharbar}{\kern0pt}\ {\isacharat}{\kern0pt}{\isacharbraceleft}{\kern0pt}term\ {\isachardoublequote}{\kern0pt}False{\isachardoublequote}{\kern0pt}{\isacharbraceright}{\kern0pt}\ {\isacharequal}{\kern0pt}{\isachargreater}{\kern0pt}\ false\isanewline
\ \ \ \ \ \ {\isacharbar}{\kern0pt}\ Const\ {\isacharparenleft}{\kern0pt}{\isacharat}{\kern0pt}{\isacharbraceleft}{\kern0pt}const{\isacharunderscore}{\kern0pt}name\ {\isachardoublequote}{\kern0pt}HOL{\isachardot}{\kern0pt}conj{\isachardoublequote}{\kern0pt}{\isacharbraceright}{\kern0pt}{\isacharcomma}{\kern0pt}\ {\isacharunderscore}{\kern0pt}{\isacharparenright}{\kern0pt}\ {\isachardollar}{\kern0pt}\ t{\isadigit{1}}\ {\isachardollar}{\kern0pt}\ t{\isadigit{2}}\ {\isacharequal}{\kern0pt}{\isachargreater}{\kern0pt}\isanewline
\ \ \ \ \ \ \ \ \ \ {\isacharparenleft}{\kern0pt}eval\ t{\isadigit{1}}{\isacharparenright}{\kern0pt}\ andalso\ {\isacharparenleft}{\kern0pt}eval\ t{\isadigit{2}}{\isacharparenright}{\kern0pt}\isanewline
\ \ \ \ \ \ {\isacharbar}{\kern0pt}\ Const\ {\isacharparenleft}{\kern0pt}{\isacharat}{\kern0pt}{\isacharbraceleft}{\kern0pt}const{\isacharunderscore}{\kern0pt}name\ {\isachardoublequote}{\kern0pt}HOL{\isachardot}{\kern0pt}disj{\isachardoublequote}{\kern0pt}{\isacharbraceright}{\kern0pt}{\isacharcomma}{\kern0pt}\ {\isacharunderscore}{\kern0pt}{\isacharparenright}{\kern0pt}\ {\isachardollar}{\kern0pt}\ t{\isadigit{1}}\ {\isachardollar}{\kern0pt}\ t{\isadigit{2}}\ {\isacharequal}{\kern0pt}{\isachargreater}{\kern0pt}\isanewline
\ \ \ \ \ \ \ \ \ \ {\isacharparenleft}{\kern0pt}eval\ t{\isadigit{1}}{\isacharparenright}{\kern0pt}\ orelse\ {\isacharparenleft}{\kern0pt}eval\ t{\isadigit{2}}{\isacharparenright}{\kern0pt}\isanewline
\ \ \ \ \ \ {\isacharbar}{\kern0pt}\ Const\ {\isacharparenleft}{\kern0pt}{\isacharat}{\kern0pt}{\isacharbraceleft}{\kern0pt}const{\isacharunderscore}{\kern0pt}name\ {\isachardoublequote}{\kern0pt}HOL{\isachardot}{\kern0pt}Not{\isachardoublequote}{\kern0pt}{\isacharbraceright}{\kern0pt}{\isacharcomma}{\kern0pt}\ {\isacharunderscore}{\kern0pt}{\isacharparenright}{\kern0pt}\ {\isachardollar}{\kern0pt}\ t{\isadigit{1}}\ {\isacharequal}{\kern0pt}{\isachargreater}{\kern0pt}\isanewline
\ \ \ \ \ \ \ \ \ \ not\ {\isacharparenleft}{\kern0pt}eval\ t{\isadigit{1}}{\isacharparenright}{\kern0pt}\isanewline
\ \ \ \ \ \ {\isacharbar}{\kern0pt}\ Const\ {\isacharparenleft}{\kern0pt}{\isacharat}{\kern0pt}{\isacharbraceleft}{\kern0pt}const{\isacharunderscore}{\kern0pt}name\ {\isachardoublequote}{\kern0pt}HOL{\isachardot}{\kern0pt}eq{\isachardoublequote}{\kern0pt}{\isacharbraceright}{\kern0pt}{\isacharcomma}{\kern0pt}\ {\isacharunderscore}{\kern0pt}{\isacharparenright}{\kern0pt}\ {\isachardollar}{\kern0pt}\ t{\isadigit{1}}\ {\isachardollar}{\kern0pt}\ t{\isadigit{2}}\ {\isacharequal}{\kern0pt}{\isachargreater}{\kern0pt}\isanewline
\ \ \ \ \ \ \ \ \ \ {\isacharparenleft}{\kern0pt}eval{\isacharunderscore}{\kern0pt}term{\isacharunderscore}{\kern0pt}to{\isacharunderscore}{\kern0pt}int\ thy\ t{\isadigit{1}}\ state{\isacharparenright}{\kern0pt}\ {\isacharequal}{\kern0pt}\ {\isacharparenleft}{\kern0pt}eval{\isacharunderscore}{\kern0pt}term{\isacharunderscore}{\kern0pt}to{\isacharunderscore}{\kern0pt}int\ thy\ t{\isadigit{2}}\ state{\isacharparenright}{\kern0pt}\isanewline
\ \ \ \ \ \ {\isacharbar}{\kern0pt}\ Const\ {\isacharparenleft}{\kern0pt}{\isacharat}{\kern0pt}{\isacharbraceleft}{\kern0pt}const{\isacharunderscore}{\kern0pt}name\ {\isachardoublequote}{\kern0pt}Orderings{\isachardot}{\kern0pt}ord{\isacharunderscore}{\kern0pt}class{\isachardot}{\kern0pt}less{\isachardoublequote}{\kern0pt}{\isacharbraceright}{\kern0pt}{\isacharcomma}{\kern0pt}\ {\isacharunderscore}{\kern0pt}{\isacharparenright}{\kern0pt}\ {\isachardollar}{\kern0pt}\ t{\isadigit{1}}\ {\isachardollar}{\kern0pt}\ t{\isadigit{2}}\ {\isacharequal}{\kern0pt}{\isachargreater}{\kern0pt}\isanewline
\ \ \ \ \ \ \ \ \ \ {\isacharparenleft}{\kern0pt}eval{\isacharunderscore}{\kern0pt}term{\isacharunderscore}{\kern0pt}to{\isacharunderscore}{\kern0pt}int\ thy\ t{\isadigit{1}}\ state{\isacharparenright}{\kern0pt}\ {\isacharless}{\kern0pt}\ {\isacharparenleft}{\kern0pt}eval{\isacharunderscore}{\kern0pt}term{\isacharunderscore}{\kern0pt}to{\isacharunderscore}{\kern0pt}int\ thy\ t{\isadigit{2}}\ state{\isacharparenright}{\kern0pt}\isanewline
\ \ \ \ \ \ {\isacharbar}{\kern0pt}\ Const\ {\isacharparenleft}{\kern0pt}{\isacharat}{\kern0pt}{\isacharbraceleft}{\kern0pt}const{\isacharunderscore}{\kern0pt}name\ {\isachardoublequote}{\kern0pt}Orderings{\isachardot}{\kern0pt}ord{\isacharunderscore}{\kern0pt}class{\isachardot}{\kern0pt}less{\isacharunderscore}{\kern0pt}eq{\isachardoublequote}{\kern0pt}{\isacharbraceright}{\kern0pt}{\isacharcomma}{\kern0pt}\ {\isacharunderscore}{\kern0pt}{\isacharparenright}{\kern0pt}\ {\isachardollar}{\kern0pt}\ t{\isadigit{1}}\ {\isachardollar}{\kern0pt}\ t{\isadigit{2}}\ {\isacharequal}{\kern0pt}{\isachargreater}{\kern0pt}\isanewline
\ \ \ \ \ \ \ \ \ \ {\isacharparenleft}{\kern0pt}eval{\isacharunderscore}{\kern0pt}term{\isacharunderscore}{\kern0pt}to{\isacharunderscore}{\kern0pt}int\ thy\ t{\isadigit{1}}\ state{\isacharparenright}{\kern0pt}\ {\isacharless}{\kern0pt}{\isacharequal}{\kern0pt}\ {\isacharparenleft}{\kern0pt}eval{\isacharunderscore}{\kern0pt}term{\isacharunderscore}{\kern0pt}to{\isacharunderscore}{\kern0pt}int\ thy\ t{\isadigit{2}}\ state{\isacharparenright}{\kern0pt}\isanewline
\ \ \ \ \ \ {\isacharbar}{\kern0pt}\ Free\ {\isacharparenleft}{\kern0pt}name{\isacharcomma}{\kern0pt}\ {\isacharunderscore}{\kern0pt}{\isacharparenright}{\kern0pt}\ {\isacharequal}{\kern0pt}{\isachargreater}{\kern0pt}\ \isanewline
\ \ \ \ \ \ \ \ \ \ {\isacharparenleft}{\kern0pt}case\ try\ {\isacharparenleft}{\kern0pt}fn\ {\isacharparenleft}{\kern0pt}{\isacharparenright}{\kern0pt}\ {\isacharequal}{\kern0pt}{\isachargreater}{\kern0pt}\ safe{\isacharunderscore}{\kern0pt}lookup\ name\ state{\isacharparenright}{\kern0pt}\ {\isacharparenleft}{\kern0pt}{\isacharparenright}{\kern0pt}\ of\isanewline
\ \ \ \ \ \ \ \ \ \ \ \ SOME\ {\isadigit{0}}\ {\isacharequal}{\kern0pt}{\isachargreater}{\kern0pt}\ false\isanewline
\ \ \ \ \ \ \ \ \ \ {\isacharbar}{\kern0pt}\ SOME\ {\isacharunderscore}{\kern0pt}\ {\isacharequal}{\kern0pt}{\isachargreater}{\kern0pt}\ true\isanewline
\ \ \ \ \ \ \ \ \ \ {\isacharbar}{\kern0pt}\ NONE\ {\isacharequal}{\kern0pt}{\isachargreater}{\kern0pt}\ false{\isacharparenright}{\kern0pt}\isanewline
\ \ \ \ \ \ {\isacharbar}{\kern0pt}\ {\isacharunderscore}{\kern0pt}\ {\isacharequal}{\kern0pt}{\isachargreater}{\kern0pt}\ error\ {\isacharparenleft}{\kern0pt}{\isachardoublequote}{\kern0pt}Unsupported\ boolean\ expression{\isacharcolon}{\kern0pt}\ {\isachardoublequote}{\kern0pt}\ {\isacharcircum}{\kern0pt}\ Pretty{\isachardot}{\kern0pt}string{\isacharunderscore}{\kern0pt}of\ {\isacharparenleft}{\kern0pt}Syntax{\isachardot}{\kern0pt}pretty{\isacharunderscore}{\kern0pt}term{\isacharunderscore}{\kern0pt}global\ thy\ t{\isacharparenright}{\kern0pt}\ {\isacharparenright}{\kern0pt}\isanewline
\ \ \ \ \ \ \ \ \ \ \ \ \ \ \ \ \ \ \ \ \ \ \ \ \ \ \ \ \ \ \ \ \ \ \ \ \ \ \ \ \ \ \ \ \ \ \ \ \ \ \ \ \ \ \ \ \ {\isacharparenleft}{\kern0pt}{\isacharasterisk}{\kern0pt}Syntax{\isachardot}{\kern0pt}string{\isacharunderscore}{\kern0pt}of{\isacharunderscore}{\kern0pt}term{\isacharunderscore}{\kern0pt}global\ thy\ t{\isacharparenright}{\kern0pt}{\isacharasterisk}{\kern0pt}{\isacharparenright}{\kern0pt}\isanewline
\ \ in\isanewline
\ \ \ \ eval\ term\isanewline
\ \ end\isanewline
{\isacharparenleft}{\kern0pt}{\isacharasterisk}{\kern0pt}\ Fixed\ record\ pattern\ matching\ {\isacharasterisk}{\kern0pt}{\isacharparenright}{\kern0pt}\isanewline
{\isacharparenleft}{\kern0pt}{\isacharasterisk}{\kern0pt}fun\ process{\isacharunderscore}{\kern0pt}absy\ absy{\isacharunderscore}{\kern0pt}data\ {\isacharequal}{\kern0pt}\isanewline
\ \ let\isanewline
\ \ \ \ val\ {\isacharbraceleft}{\kern0pt}system{\isacharcomma}{\kern0pt}\ globals{\isacharunderscore}{\kern0pt}decls{\isacharcomma}{\kern0pt}\ locks{\isacharunderscore}{\kern0pt}decls{\isacharcomma}{\kern0pt}\ threads{\isacharunderscore}{\kern0pt}decls{\isacharcomma}{\kern0pt}\ varstab{\isacharbraceright}{\kern0pt}\ {\isacharequal}{\kern0pt}\ absy{\isacharunderscore}{\kern0pt}data\isanewline
\ \ in\isanewline
\ \ \ \ {\isacharparenleft}{\kern0pt}{\isacharasterisk}{\kern0pt}\ Process\ the\ fields\ {\isacharasterisk}{\kern0pt}{\isacharparenright}{\kern0pt}\isanewline
\ \ \ \ {\isacharparenleft}{\kern0pt}system{\isacharcomma}{\kern0pt}\ globals{\isacharunderscore}{\kern0pt}decls{\isacharcomma}{\kern0pt}\ locks{\isacharunderscore}{\kern0pt}decls{\isacharcomma}{\kern0pt}\ threads{\isacharunderscore}{\kern0pt}decls{\isacharparenright}{\kern0pt}\isanewline
\ \ end\isanewline
{\isacharasterisk}{\kern0pt}{\isacharparenright}{\kern0pt}\isanewline
{\isacharparenleft}{\kern0pt}{\isacharasterisk}{\kern0pt}\ Fonction\ de\ conversion\ de\ l{\isacharprime}{\kern0pt}AST\ vers\ com\ {\isacharasterisk}{\kern0pt}{\isacharparenright}{\kern0pt}\isanewline
fun\ convert{\isacharunderscore}{\kern0pt}to{\isacharunderscore}{\kern0pt}com\ {\isacharparenleft}{\kern0pt}thy\ {\isacharcolon}{\kern0pt}\ theory{\isacharparenright}{\kern0pt}\ {\isacharparenleft}{\kern0pt}actions\ {\isacharcolon}{\kern0pt}\ a{\isacharunderscore}{\kern0pt}term\ list{\isacharparenright}{\kern0pt}\ \isanewline
\ \ \ \ \ \ \ \ \ \ \ \ \ \ \ \ \ \ \ {\isacharparenleft}{\kern0pt}varstab\ {\isacharcolon}{\kern0pt}\ {\isacharparenleft}{\kern0pt}term\ option{\isacharparenright}{\kern0pt}\ Symtab{\isachardot}{\kern0pt}table{\isacharparenright}{\kern0pt}\ \isanewline
\ \ \ \ \ \ \ \ \ \ \ \ \ \ \ \ \ \ \ {\isacharparenleft}{\kern0pt}locstab\ {\isacharcolon}{\kern0pt}\ {\isacharparenleft}{\kern0pt}term\ option{\isacharparenright}{\kern0pt}\ Symtab{\isachardot}{\kern0pt}table{\isacharparenright}{\kern0pt}\ {\isacharcolon}{\kern0pt}\ com\ {\isacharequal}{\kern0pt}\isanewline
\ \ let\isanewline
\ \ \ \ fun\ term{\isacharunderscore}{\kern0pt}to{\isacharunderscore}{\kern0pt}arith{\isacharunderscore}{\kern0pt}expr\ {\isacharparenleft}{\kern0pt}t\ {\isacharcolon}{\kern0pt}\ term{\isacharparenright}{\kern0pt}\ {\isacharcolon}{\kern0pt}\ arith{\isacharunderscore}{\kern0pt}expr\ {\isacharequal}{\kern0pt}\ \isanewline
\ \ \ \ \ \ fn\ state\ {\isacharequal}{\kern0pt}{\isachargreater}{\kern0pt}\ eval{\isacharunderscore}{\kern0pt}term{\isacharunderscore}{\kern0pt}to{\isacharunderscore}{\kern0pt}int\ thy\ t\ state\isanewline
\ \ \ \ \isanewline
\ \ \ \ fun\ term{\isacharunderscore}{\kern0pt}to{\isacharunderscore}{\kern0pt}bool{\isacharunderscore}{\kern0pt}expr\ {\isacharparenleft}{\kern0pt}t\ {\isacharcolon}{\kern0pt}\ term{\isacharparenright}{\kern0pt}\ {\isacharcolon}{\kern0pt}\ bool{\isacharunderscore}{\kern0pt}expr\ {\isacharequal}{\kern0pt}\isanewline
\ \ \ \ \ \ fn\ state\ {\isacharequal}{\kern0pt}{\isachargreater}{\kern0pt}\ eval{\isacharunderscore}{\kern0pt}term{\isacharunderscore}{\kern0pt}to{\isacharunderscore}{\kern0pt}bool\ thy\ t\ state\isanewline
\ \ \ \ \isanewline
\ \ \ \ fun\ convert{\isacharunderscore}{\kern0pt}action\ {\isacharparenleft}{\kern0pt}action\ {\isacharcolon}{\kern0pt}\ a{\isacharunderscore}{\kern0pt}term{\isacharparenright}{\kern0pt}\ {\isacharcolon}{\kern0pt}\ com\ {\isacharequal}{\kern0pt}\isanewline
\ \ \ \ \ \ case\ action\ of\isanewline
\ \ \ \ \ \ \ \ SkipA\ {\isacharequal}{\kern0pt}{\isachargreater}{\kern0pt}\ SKIP\isanewline
\ \ \ \ \ \ {\isacharbar}{\kern0pt}\ AssignA\ {\isacharparenleft}{\kern0pt}bdg{\isacharcomma}{\kern0pt}\ term{\isacharparenright}{\kern0pt}\ {\isacharequal}{\kern0pt}{\isachargreater}{\kern0pt}\ \isanewline
\ \ \ \ \ \ \ \ let\isanewline
\ \ \ \ \ \ \ \ \ \ val\ name\ {\isacharequal}{\kern0pt}\ Binding{\isachardot}{\kern0pt}name{\isacharunderscore}{\kern0pt}of\ bdg\isanewline
\ \ \ \ \ \ \ \ \ \ val\ is{\isacharunderscore}{\kern0pt}local\ {\isacharequal}{\kern0pt}\ Symtab{\isachardot}{\kern0pt}defined\ locstab\ name\isanewline
\ \ \ \ \ \ \ \ in\isanewline
\ \ \ \ \ \ \ \ \ \ if\ is{\isacharunderscore}{\kern0pt}local\isanewline
\ \ \ \ \ \ \ \ \ \ then\ Assign\ {\isacharparenleft}{\kern0pt}name{\isacharcomma}{\kern0pt}\ term{\isacharunderscore}{\kern0pt}to{\isacharunderscore}{\kern0pt}arith{\isacharunderscore}{\kern0pt}expr\ term{\isacharparenright}{\kern0pt}\ \ {\isacharparenleft}{\kern0pt}{\isacharasterisk}{\kern0pt}\ local\ {\isacharasterisk}{\kern0pt}{\isacharparenright}{\kern0pt}\isanewline
\ \ \ \ \ \ \ \ \ \ else\ Store\ {\isacharparenleft}{\kern0pt}term{\isacharunderscore}{\kern0pt}to{\isacharunderscore}{\kern0pt}arith{\isacharunderscore}{\kern0pt}expr\ term{\isacharcomma}{\kern0pt}\ name{\isacharparenright}{\kern0pt}\ \ \ {\isacharparenleft}{\kern0pt}{\isacharasterisk}{\kern0pt}\ global\ {\isacharasterisk}{\kern0pt}{\isacharparenright}{\kern0pt}\isanewline
\ \ \ \ \ \ \ \ end\isanewline
\ \ \ \ \ \ {\isacharbar}{\kern0pt}\ LockA\ bdg\ {\isacharequal}{\kern0pt}{\isachargreater}{\kern0pt}\ \isanewline
\ \ \ \ \ \ \ \ \ \ Lock\ {\isacharparenleft}{\kern0pt}Binding{\isachardot}{\kern0pt}name{\isacharunderscore}{\kern0pt}of\ bdg{\isacharparenright}{\kern0pt}\isanewline
\ \ \ \ \ \ {\isacharbar}{\kern0pt}\ UnlockA\ bdg\ {\isacharequal}{\kern0pt}{\isachargreater}{\kern0pt}\ \isanewline
\ \ \ \ \ \ \ \ \ \ Unlock\ {\isacharparenleft}{\kern0pt}Binding{\isachardot}{\kern0pt}name{\isacharunderscore}{\kern0pt}of\ bdg{\isacharparenright}{\kern0pt}\isanewline
\ \ \ \ \ \ {\isacharbar}{\kern0pt}\ SendA\ {\isacharparenleft}{\kern0pt}bdg{\isacharcomma}{\kern0pt}\ term{\isacharparenright}{\kern0pt}\ {\isacharequal}{\kern0pt}{\isachargreater}{\kern0pt}\ \isanewline
\ \ \ \ \ \ \ \ \ \ Store\ {\isacharparenleft}{\kern0pt}term{\isacharunderscore}{\kern0pt}to{\isacharunderscore}{\kern0pt}arith{\isacharunderscore}{\kern0pt}expr\ term{\isacharcomma}{\kern0pt}\ Binding{\isachardot}{\kern0pt}name{\isacharunderscore}{\kern0pt}of\ bdg{\isacharparenright}{\kern0pt}\isanewline
\ \ \ \ \ \ {\isacharbar}{\kern0pt}\ ReceiveA\ {\isacharparenleft}{\kern0pt}local{\isacharunderscore}{\kern0pt}bdg{\isacharcomma}{\kern0pt}\ shared{\isacharunderscore}{\kern0pt}bdg{\isacharparenright}{\kern0pt}\ {\isacharequal}{\kern0pt}{\isachargreater}{\kern0pt}\ \isanewline
\ \ \ \ \ \ \ \ \ \ Load\ {\isacharparenleft}{\kern0pt}Binding{\isachardot}{\kern0pt}name{\isacharunderscore}{\kern0pt}of\ local{\isacharunderscore}{\kern0pt}bdg{\isacharcomma}{\kern0pt}\ Binding{\isachardot}{\kern0pt}name{\isacharunderscore}{\kern0pt}of\ shared{\isacharunderscore}{\kern0pt}bdg{\isacharparenright}{\kern0pt}\isanewline
\ \ \ \ \ \ {\isacharbar}{\kern0pt}\ IfelseA\ {\isacharparenleft}{\kern0pt}{\isacharparenleft}{\kern0pt}cond{\isacharunderscore}{\kern0pt}term{\isacharcomma}{\kern0pt}\ then{\isacharunderscore}{\kern0pt}actions{\isacharparenright}{\kern0pt}{\isacharcomma}{\kern0pt}\ else{\isacharunderscore}{\kern0pt}actions{\isacharparenright}{\kern0pt}\ {\isacharequal}{\kern0pt}{\isachargreater}{\kern0pt}\isanewline
\ \ \ \ \ \ \ \ \ \ Cond\ {\isacharparenleft}{\kern0pt}term{\isacharunderscore}{\kern0pt}to{\isacharunderscore}{\kern0pt}bool{\isacharunderscore}{\kern0pt}expr\ cond{\isacharunderscore}{\kern0pt}term{\isacharcomma}{\kern0pt}\isanewline
\ \ \ \ \ \ \ \ \ \ \ \ \ \ \ \ convert{\isacharunderscore}{\kern0pt}actions{\isacharunderscore}{\kern0pt}to{\isacharunderscore}{\kern0pt}seq\ then{\isacharunderscore}{\kern0pt}actions{\isacharcomma}{\kern0pt}\isanewline
\ \ \ \ \ \ \ \ \ \ \ \ \ \ \ \ convert{\isacharunderscore}{\kern0pt}actions{\isacharunderscore}{\kern0pt}to{\isacharunderscore}{\kern0pt}seq\ else{\isacharunderscore}{\kern0pt}actions{\isacharparenright}{\kern0pt}\isanewline
\ \ \ \ \ \ {\isacharbar}{\kern0pt}\ WhileA\ {\isacharparenleft}{\kern0pt}cond{\isacharunderscore}{\kern0pt}term{\isacharcomma}{\kern0pt}\ body{\isacharunderscore}{\kern0pt}actions{\isacharparenright}{\kern0pt}\ {\isacharequal}{\kern0pt}{\isachargreater}{\kern0pt}\isanewline
\ \ \ \ \ \ \ \ \ \ While\ {\isacharparenleft}{\kern0pt}term{\isacharunderscore}{\kern0pt}to{\isacharunderscore}{\kern0pt}bool{\isacharunderscore}{\kern0pt}expr\ cond{\isacharunderscore}{\kern0pt}term{\isacharcomma}{\kern0pt}\ convert{\isacharunderscore}{\kern0pt}actions{\isacharunderscore}{\kern0pt}to{\isacharunderscore}{\kern0pt}seq\ body{\isacharunderscore}{\kern0pt}actions{\isacharparenright}{\kern0pt}\isanewline
\ \ \ \ \ \ \ \ \isanewline
\ \ \ \ \isanewline
\ \ \ \ and\ convert{\isacharunderscore}{\kern0pt}actions{\isacharunderscore}{\kern0pt}to{\isacharunderscore}{\kern0pt}seq\ {\isacharparenleft}{\kern0pt}actions\ {\isacharcolon}{\kern0pt}\ a{\isacharunderscore}{\kern0pt}term\ list{\isacharparenright}{\kern0pt}\ {\isacharcolon}{\kern0pt}\ com\ {\isacharequal}{\kern0pt}\isanewline
\ \ \ \ \ \ case\ actions\ of\isanewline
\ \ \ \ \ \ \ \ {\isacharbrackleft}{\kern0pt}{\isacharbrackright}{\kern0pt}\ {\isacharequal}{\kern0pt}{\isachargreater}{\kern0pt}\ SKIP\isanewline
\ \ \ \ \ \ {\isacharbar}{\kern0pt}\ {\isacharbrackleft}{\kern0pt}action{\isacharbrackright}{\kern0pt}\ {\isacharequal}{\kern0pt}{\isachargreater}{\kern0pt}\ convert{\isacharunderscore}{\kern0pt}action\ action\isanewline
\ \ \ \ \ \ {\isacharbar}{\kern0pt}\ action\ {\isacharcolon}{\kern0pt}{\isacharcolon}{\kern0pt}\ rest\ {\isacharequal}{\kern0pt}{\isachargreater}{\kern0pt}\ Seq\ {\isacharparenleft}{\kern0pt}convert{\isacharunderscore}{\kern0pt}action\ action{\isacharcomma}{\kern0pt}\ convert{\isacharunderscore}{\kern0pt}actions{\isacharunderscore}{\kern0pt}to{\isacharunderscore}{\kern0pt}seq\ rest{\isacharparenright}{\kern0pt}\isanewline
\ \ in\isanewline
\ \ \ \ convert{\isacharunderscore}{\kern0pt}actions{\isacharunderscore}{\kern0pt}to{\isacharunderscore}{\kern0pt}seq\ actions\isanewline
\ \ end\isanewline
\isanewline
fun\ check{\isacharunderscore}{\kern0pt}lock{\isacharunderscore}{\kern0pt}ownership\ {\isacharparenleft}{\kern0pt}thread{\isacharunderscore}{\kern0pt}name\ {\isacharcolon}{\kern0pt}\ string{\isacharparenright}{\kern0pt}\ {\isacharparenleft}{\kern0pt}actions\ {\isacharcolon}{\kern0pt}\ a{\isacharunderscore}{\kern0pt}term\ list{\isacharparenright}{\kern0pt}\ {\isacharequal}{\kern0pt}\isanewline
\ \ let\isanewline
\ \ \ \ fun\ check\ {\isacharbrackleft}{\kern0pt}{\isacharbrackright}{\kern0pt}\ {\isacharunderscore}{\kern0pt}\ {\isacharequal}{\kern0pt}\ NONE\isanewline
\ \ \ \ \ \ {\isacharbar}{\kern0pt}\ check\ {\isacharparenleft}{\kern0pt}LockA\ b\ {\isacharcolon}{\kern0pt}{\isacharcolon}{\kern0pt}\ rest{\isacharparenright}{\kern0pt}\ held\ {\isacharequal}{\kern0pt}\isanewline
\ \ \ \ \ \ \ \ \ \ let\ val\ l\ {\isacharequal}{\kern0pt}\ Binding{\isachardot}{\kern0pt}name{\isacharunderscore}{\kern0pt}of\ b\ in\isanewline
\ \ \ \ \ \ \ \ \ \ \ \ if\ List{\isachardot}{\kern0pt}exists\ {\isacharparenleft}{\kern0pt}fn\ x\ {\isacharequal}{\kern0pt}{\isachargreater}{\kern0pt}\ x\ {\isacharequal}{\kern0pt}\ l{\isacharparenright}{\kern0pt}\ held\ then\ \isanewline
\ \ \ \ \ \ \ \ \ \ \ \ \ \ SOME\ {\isacharparenleft}{\kern0pt}{\isachardoublequote}{\kern0pt}Thread\ {\isachardoublequote}{\kern0pt}\ {\isacharcircum}{\kern0pt}\ thread{\isacharunderscore}{\kern0pt}name\ {\isacharcircum}{\kern0pt}\ {\isachardoublequote}{\kern0pt}{\isacharcolon}{\kern0pt}\ LOCK\ {\isachardoublequote}{\kern0pt}\ {\isacharcircum}{\kern0pt}\ l\ {\isacharcircum}{\kern0pt}\ {\isachardoublequote}{\kern0pt}\ already\ held{\isachardoublequote}{\kern0pt}{\isacharparenright}{\kern0pt}\isanewline
\ \ \ \ \ \ \ \ \ \ \ \ else\isanewline
\ \ \ \ \ \ \ \ \ \ \ \ \ \ check\ rest\ {\isacharparenleft}{\kern0pt}l\ {\isacharcolon}{\kern0pt}{\isacharcolon}{\kern0pt}\ held{\isacharparenright}{\kern0pt}\isanewline
\ \ \ \ \ \ \ \ \ \ end\isanewline
\ \ \ \ \ \ {\isacharbar}{\kern0pt}\ check\ {\isacharparenleft}{\kern0pt}UnlockA\ b\ {\isacharcolon}{\kern0pt}{\isacharcolon}{\kern0pt}\ rest{\isacharparenright}{\kern0pt}\ held\ {\isacharequal}{\kern0pt}\isanewline
\ \ \ \ \ \ \ \ \ \ let\ val\ l\ {\isacharequal}{\kern0pt}\ Binding{\isachardot}{\kern0pt}name{\isacharunderscore}{\kern0pt}of\ b\ in\isanewline
\ \ \ \ \ \ \ \ \ \ \ \ if\ List{\isachardot}{\kern0pt}exists\ {\isacharparenleft}{\kern0pt}fn\ x\ {\isacharequal}{\kern0pt}{\isachargreater}{\kern0pt}\ x\ {\isacharequal}{\kern0pt}\ l{\isacharparenright}{\kern0pt}\ held\ then\isanewline
\ \ \ \ \ \ \ \ \ \ \ \ \ \ check\ rest\ {\isacharparenleft}{\kern0pt}remove\ {\isacharparenleft}{\kern0pt}op\ {\isacharequal}{\kern0pt}{\isacharparenright}{\kern0pt}\ l\ held{\isacharparenright}{\kern0pt}\isanewline
\ \ \ \ \ \ \ \ \ \ \ \ else\isanewline
\ \ \ \ \ \ \ \ \ \ \ \ \ \ SOME\ {\isacharparenleft}{\kern0pt}{\isachardoublequote}{\kern0pt}Thread\ {\isachardoublequote}{\kern0pt}\ {\isacharcircum}{\kern0pt}\ thread{\isacharunderscore}{\kern0pt}name\ {\isacharcircum}{\kern0pt}\ {\isachardoublequote}{\kern0pt}{\isacharcolon}{\kern0pt}\ tries\ to\ UNLOCK\ {\isachardoublequote}{\kern0pt}\ {\isacharcircum}{\kern0pt}\ l\ {\isacharcircum}{\kern0pt}\ {\isachardoublequote}{\kern0pt}\ without\ holding\ it{\isachardoublequote}{\kern0pt}{\isacharparenright}{\kern0pt}\isanewline
\ \ \ \ \ \ \ \ \ \ end\isanewline
\ \ \ \ \ \ {\isacharbar}{\kern0pt}\ check\ {\isacharparenleft}{\kern0pt}IfelseA\ {\isacharparenleft}{\kern0pt}{\isacharparenleft}{\kern0pt}{\isacharunderscore}{\kern0pt}{\isacharcomma}{\kern0pt}\ ifBody{\isacharparenright}{\kern0pt}{\isacharcomma}{\kern0pt}\ elseBody{\isacharparenright}{\kern0pt}\ {\isacharcolon}{\kern0pt}{\isacharcolon}{\kern0pt}\ rest{\isacharparenright}{\kern0pt}\ held\ {\isacharequal}{\kern0pt}\isanewline
\ \ \ \ \ \ \ \ \ \ {\isacharparenleft}{\kern0pt}case\ check\ ifBody\ held\ of\isanewline
\ \ \ \ \ \ \ \ \ \ \ \ \ \ SOME\ err\ {\isacharequal}{\kern0pt}{\isachargreater}{\kern0pt}\ SOME\ err\isanewline
\ \ \ \ \ \ \ \ \ \ \ \ {\isacharbar}{\kern0pt}\ NONE\ {\isacharequal}{\kern0pt}{\isachargreater}{\kern0pt}\isanewline
\ \ \ \ \ \ \ \ \ \ \ \ \ \ {\isacharparenleft}{\kern0pt}case\ check\ elseBody\ held\ of\isanewline
\ \ \ \ \ \ \ \ \ \ \ \ \ \ \ \ \ \ SOME\ err\ {\isacharequal}{\kern0pt}{\isachargreater}{\kern0pt}\ SOME\ err\isanewline
\ \ \ \ \ \ \ \ \ \ \ \ \ \ \ \ {\isacharbar}{\kern0pt}\ NONE\ {\isacharequal}{\kern0pt}{\isachargreater}{\kern0pt}\ check\ rest\ held{\isacharparenright}{\kern0pt}{\isacharparenright}{\kern0pt}\isanewline
\ \ \ \ \ \ {\isacharbar}{\kern0pt}\ check\ {\isacharparenleft}{\kern0pt}WhileA\ {\isacharparenleft}{\kern0pt}{\isacharunderscore}{\kern0pt}{\isacharcomma}{\kern0pt}\ body{\isacharparenright}{\kern0pt}\ {\isacharcolon}{\kern0pt}{\isacharcolon}{\kern0pt}\ rest{\isacharparenright}{\kern0pt}\ held\ {\isacharequal}{\kern0pt}\ \isanewline
\ \ \ \ \ \ \ \ \ \ {\isacharparenleft}{\kern0pt}case\ check\ body\ held\ of\isanewline
\ \ \ \ \ \ \ \ \ \ \ \ \ \ SOME\ err\ {\isacharequal}{\kern0pt}{\isachargreater}{\kern0pt}\ SOME\ err\isanewline
\ \ \ \ \ \ \ \ \ \ \ \ {\isacharbar}{\kern0pt}\ NONE\ {\isacharequal}{\kern0pt}{\isachargreater}{\kern0pt}\ check\ rest\ held{\isacharparenright}{\kern0pt}\isanewline
\ \ \ \ \ \ {\isacharbar}{\kern0pt}\ check\ {\isacharparenleft}{\kern0pt}{\isacharunderscore}{\kern0pt}\ {\isacharcolon}{\kern0pt}{\isacharcolon}{\kern0pt}\ rest{\isacharparenright}{\kern0pt}\ held\ {\isacharequal}{\kern0pt}\ check\ rest\ held\isanewline
\ \ in\isanewline
\ \ \ \ check\ actions\ {\isacharbrackleft}{\kern0pt}{\isacharbrackright}{\kern0pt}\isanewline
\ \ end\isanewline
\isanewline
{\isacharparenleft}{\kern0pt}{\isacharasterisk}{\kern0pt}\ Alternative\ CPS{\isacharminus}{\kern0pt}style\ semantics\ {\isacharasterisk}{\kern0pt}{\isacharparenright}{\kern0pt}\isanewline
datatype\ csp{\isacharunderscore}{\kern0pt}event\ {\isacharequal}{\kern0pt}\ \isanewline
\ \ \ \ LockEvent\ of\ string\ {\isacharasterisk}{\kern0pt}\ int\isanewline
\ \ {\isacharbar}{\kern0pt}\ UnlockEvent\ of\ string\ {\isacharasterisk}{\kern0pt}\ int\ \ \isanewline
\ \ {\isacharbar}{\kern0pt}\ ReadEvent\ of\ string\ {\isacharasterisk}{\kern0pt}\ int\isanewline
\ \ {\isacharbar}{\kern0pt}\ UpdateEvent\ of\ string\ {\isacharasterisk}{\kern0pt}\ int\isanewline
\ \ {\isacharbar}{\kern0pt}\ AssignEvent\ of\ string\ {\isacharasterisk}{\kern0pt}\ int\isanewline
\ \ {\isacharbar}{\kern0pt}\ SkipEvent\isanewline
{\isacharparenleft}{\kern0pt}{\isacharasterisk}{\kern0pt}\ Mise\ à\ jour\ de\ l’état\ fonctionnel\ {\isacharasterisk}{\kern0pt}{\isacharparenright}{\kern0pt}\isanewline
fun\ update{\isacharunderscore}{\kern0pt}state\ state\ x\ v\ {\isacharequal}{\kern0pt}\ fn\ y\ {\isacharequal}{\kern0pt}{\isachargreater}{\kern0pt}\ if\ y\ {\isacharequal}{\kern0pt}\ x\ then\ v\ else\ state\ y\isanewline
\isanewline
{\isacharparenleft}{\kern0pt}{\isacharasterisk}{\kern0pt}\ Semantique\ CPS\ complète\ {\isacharasterisk}{\kern0pt}{\isacharparenright}{\kern0pt}\isanewline
fun\ sem{\isacharunderscore}{\kern0pt}cps\ {\isacharunderscore}{\kern0pt}\ SKIP\ cont\ {\isacharequal}{\kern0pt}\isanewline
\ \ \ \ \ \ {\isacharparenleft}{\kern0pt}fn\ state\ {\isacharequal}{\kern0pt}{\isachargreater}{\kern0pt}\ SkipEvent\ {\isacharcolon}{\kern0pt}{\isacharcolon}{\kern0pt}\ cont\ state{\isacharparenright}{\kern0pt}\isanewline
\isanewline
\ \ {\isacharbar}{\kern0pt}\ sem{\isacharunderscore}{\kern0pt}cps\ {\isacharunderscore}{\kern0pt}\ {\isacharparenleft}{\kern0pt}Assign\ {\isacharparenleft}{\kern0pt}x{\isacharcomma}{\kern0pt}\ e{\isacharparenright}{\kern0pt}{\isacharparenright}{\kern0pt}\ cont\ {\isacharequal}{\kern0pt}\ \isanewline
\ \ \ \ \ \ {\isacharparenleft}{\kern0pt}fn\ state\ {\isacharequal}{\kern0pt}{\isachargreater}{\kern0pt}\isanewline
\ \ \ \ \ \ \ \ let\isanewline
\ \ \ \ \ \ \ \ \ \ val\ v\ {\isacharequal}{\kern0pt}\ e\ state\isanewline
\ \ \ \ \ \ \ \ in\isanewline
\ \ \ \ \ \ \ \ \ \ AssignEvent\ {\isacharparenleft}{\kern0pt}x{\isacharcomma}{\kern0pt}\ v{\isacharparenright}{\kern0pt}\ {\isacharcolon}{\kern0pt}{\isacharcolon}{\kern0pt}\ cont\ {\isacharparenleft}{\kern0pt}update{\isacharunderscore}{\kern0pt}state\ state\ x\ v{\isacharparenright}{\kern0pt}\isanewline
\ \ \ \ \ \ \ \ end{\isacharparenright}{\kern0pt}\isanewline
\isanewline
\ \ {\isacharbar}{\kern0pt}\ sem{\isacharunderscore}{\kern0pt}cps\ machine\ {\isacharparenleft}{\kern0pt}Seq\ {\isacharparenleft}{\kern0pt}P{\isacharcomma}{\kern0pt}\ Q{\isacharparenright}{\kern0pt}{\isacharparenright}{\kern0pt}\ cont\ {\isacharequal}{\kern0pt}\ \isanewline
\ \ \ \ \ \ sem{\isacharunderscore}{\kern0pt}cps\ machine\ P\ {\isacharparenleft}{\kern0pt}sem{\isacharunderscore}{\kern0pt}cps\ machine\ Q\ cont{\isacharparenright}{\kern0pt}\isanewline
\isanewline
\ \ {\isacharbar}{\kern0pt}\ sem{\isacharunderscore}{\kern0pt}cps\ machine\ {\isacharparenleft}{\kern0pt}Cond\ {\isacharparenleft}{\kern0pt}cond{\isacharcomma}{\kern0pt}\ t{\isadigit{1}}{\isacharcomma}{\kern0pt}\ t{\isadigit{2}}{\isacharparenright}{\kern0pt}{\isacharparenright}{\kern0pt}\ cont\ {\isacharequal}{\kern0pt}\isanewline
\ \ \ \ \ \ {\isacharparenleft}{\kern0pt}fn\ state\ {\isacharequal}{\kern0pt}{\isachargreater}{\kern0pt}\isanewline
\ \ \ \ \ \ \ \ if\ cond\ state\ then\ sem{\isacharunderscore}{\kern0pt}cps\ machine\ t{\isadigit{1}}\ cont\ state\isanewline
\ \ \ \ \ \ \ \ else\ sem{\isacharunderscore}{\kern0pt}cps\ machine\ t{\isadigit{2}}\ cont\ state{\isacharparenright}{\kern0pt}\isanewline
\isanewline
\ \ {\isacharbar}{\kern0pt}\ sem{\isacharunderscore}{\kern0pt}cps\ machine\ {\isacharparenleft}{\kern0pt}While\ {\isacharparenleft}{\kern0pt}cond{\isacharcomma}{\kern0pt}\ body{\isacharparenright}{\kern0pt}{\isacharparenright}{\kern0pt}\ cont\ {\isacharequal}{\kern0pt}\isanewline
\ \ \ \ \ \ {\isacharparenleft}{\kern0pt}let\isanewline
\ \ \ \ \ \ \ \ fun\ loop\ state\ {\isacharequal}{\kern0pt}\isanewline
\ \ \ \ \ \ \ \ \ \ if\ cond\ state\ then\isanewline
\ \ \ \ \ \ \ \ \ \ \ \ sem{\isacharunderscore}{\kern0pt}cps\ \ machine\ body\ loop\ state\isanewline
\ \ \ \ \ \ \ \ \ \ else\ cont\ state\isanewline
\ \ \ \ \ \ in\ loop\ end{\isacharparenright}{\kern0pt}\isanewline
\isanewline
\ \ {\isacharbar}{\kern0pt}\ sem{\isacharunderscore}{\kern0pt}cps\ machine\ {\isacharparenleft}{\kern0pt}Lock\ id{\isacharparenright}{\kern0pt}\ cont\ {\isacharequal}{\kern0pt}\isanewline
\ \ \ \ \ \ {\isacharparenleft}{\kern0pt}fn\ state\ {\isacharequal}{\kern0pt}{\isachargreater}{\kern0pt}\ LockEvent\ {\isacharparenleft}{\kern0pt}id{\isacharcomma}{\kern0pt}\ get{\isacharunderscore}{\kern0pt}lock{\isacharunderscore}{\kern0pt}id\ id\ machine{\isacharparenright}{\kern0pt}\ {\isacharcolon}{\kern0pt}{\isacharcolon}{\kern0pt}\ cont\ state{\isacharparenright}{\kern0pt}\isanewline
\isanewline
\ \ {\isacharbar}{\kern0pt}\ sem{\isacharunderscore}{\kern0pt}cps\ machine\ {\isacharparenleft}{\kern0pt}Unlock\ id{\isacharparenright}{\kern0pt}\ cont\ {\isacharequal}{\kern0pt}\isanewline
\ \ \ \ \ \ {\isacharparenleft}{\kern0pt}fn\ state\ {\isacharequal}{\kern0pt}{\isachargreater}{\kern0pt}\ UnlockEvent\ {\isacharparenleft}{\kern0pt}id{\isacharcomma}{\kern0pt}\ get{\isacharunderscore}{\kern0pt}lock{\isacharunderscore}{\kern0pt}id\ id\ machine{\isacharparenright}{\kern0pt}\ {\isacharcolon}{\kern0pt}{\isacharcolon}{\kern0pt}\ cont\ state{\isacharparenright}{\kern0pt}\isanewline
\isanewline
\ \ {\isacharbar}{\kern0pt}\ sem{\isacharunderscore}{\kern0pt}cps\ {\isacharunderscore}{\kern0pt}\ {\isacharparenleft}{\kern0pt}Store\ {\isacharparenleft}{\kern0pt}e{\isacharcomma}{\kern0pt}\ x{\isacharparenright}{\kern0pt}{\isacharparenright}{\kern0pt}\ cont\ {\isacharequal}{\kern0pt}\isanewline
\ \ \ \ {\isacharparenleft}{\kern0pt}fn\ state\ {\isacharequal}{\kern0pt}{\isachargreater}{\kern0pt}\isanewline
\ \ \ \ \ \ let\isanewline
\ \ \ \ \ \ \ \ val\ v\ {\isacharequal}{\kern0pt}\ e\ state\isanewline
\ \ \ \ \ \ \ \ val\ new{\isacharunderscore}{\kern0pt}state\ {\isacharequal}{\kern0pt}\ update{\isacharunderscore}{\kern0pt}state\ state\ x\ v\isanewline
\ \ \ \ \ \ in\isanewline
\ \ \ \ \ \ \ \ UpdateEvent\ {\isacharparenleft}{\kern0pt}x{\isacharcomma}{\kern0pt}\ v{\isacharparenright}{\kern0pt}\ {\isacharcolon}{\kern0pt}{\isacharcolon}{\kern0pt}\ cont\ new{\isacharunderscore}{\kern0pt}state\isanewline
\ \ \ \ \ \ end{\isacharparenright}{\kern0pt}\isanewline
\ \ {\isacharbar}{\kern0pt}\ sem{\isacharunderscore}{\kern0pt}cps\ {\isacharunderscore}{\kern0pt}\ {\isacharparenleft}{\kern0pt}Load\ {\isacharparenleft}{\kern0pt}x{\isacharcomma}{\kern0pt}\ y{\isacharparenright}{\kern0pt}{\isacharparenright}{\kern0pt}\ cont\ {\isacharequal}{\kern0pt}\isanewline
\ \ \ \ {\isacharparenleft}{\kern0pt}fn\ state\ {\isacharequal}{\kern0pt}{\isachargreater}{\kern0pt}\isanewline
\ \ \ \ \ \ let\ \ \ \ \ \ \ \ \ \isanewline
\ \ \ \ \ \ \ \ val\ v\ {\isacharequal}{\kern0pt}\ safe{\isacharunderscore}{\kern0pt}lookup\ y\ state\isanewline
\ \ \ \ \ \ \ \ val\ new{\isacharunderscore}{\kern0pt}state\ {\isacharequal}{\kern0pt}\ update{\isacharunderscore}{\kern0pt}state\ state\ x\ v\isanewline
\ \ \ \ \ \ in\isanewline
\ \ \ \ \ \ \ \ ReadEvent\ {\isacharparenleft}{\kern0pt}y{\isacharcomma}{\kern0pt}\ v{\isacharparenright}{\kern0pt}\ {\isacharcolon}{\kern0pt}{\isacharcolon}{\kern0pt}\ cont\ new{\isacharunderscore}{\kern0pt}state\isanewline
\ \ \ \ \ \ end{\isacharparenright}{\kern0pt}\isanewline
type\ csp{\isacharunderscore}{\kern0pt}trace\ {\isacharequal}{\kern0pt}\ csp{\isacharunderscore}{\kern0pt}event\ list\isanewline
\isanewline
{\isacharparenleft}{\kern0pt}{\isacharasterisk}{\kern0pt}\ Fonction\ Sem\ adaptée\ avec\ gestion\ des\ événements\ CSP\ {\isacharasterisk}{\kern0pt}{\isacharparenright}{\kern0pt}\isanewline
fun\ sem{\isacharunderscore}{\kern0pt}step\ {\isacharparenleft}{\kern0pt}thy\ {\isacharcolon}{\kern0pt}\ theory{\isacharparenright}{\kern0pt}\ {\isacharparenleft}{\kern0pt}machine{\isacharcolon}{\kern0pt}\ absy{\isacharparenright}{\kern0pt}\ {\isacharparenleft}{\kern0pt}cmd\ {\isacharcolon}{\kern0pt}\ com{\isacharparenright}{\kern0pt}\ {\isacharparenleft}{\kern0pt}cont\ {\isacharcolon}{\kern0pt}\ state\ {\isacharasterisk}{\kern0pt}\ csp{\isacharunderscore}{\kern0pt}trace\ {\isacharminus}{\kern0pt}{\isachargreater}{\kern0pt}\ state\ {\isacharasterisk}{\kern0pt}\ csp{\isacharunderscore}{\kern0pt}trace{\isacharparenright}{\kern0pt}\ \isanewline
\ \ \ \ \ \ \ \ \ \ \ \ \ {\isacharparenleft}{\kern0pt}s\ {\isacharcolon}{\kern0pt}\ state{\isacharcomma}{\kern0pt}\ trace\ {\isacharcolon}{\kern0pt}\ csp{\isacharunderscore}{\kern0pt}trace{\isacharparenright}{\kern0pt}\ {\isacharcolon}{\kern0pt}\ state\ {\isacharasterisk}{\kern0pt}\ csp{\isacharunderscore}{\kern0pt}trace\ {\isacharequal}{\kern0pt}\isanewline
\ \ case\ cmd\ of\isanewline
\ \ \ \ SKIP\ {\isacharequal}{\kern0pt}{\isachargreater}{\kern0pt}\ cont\ {\isacharparenleft}{\kern0pt}s{\isacharcomma}{\kern0pt}\ trace{\isacharparenright}{\kern0pt}\isanewline
\ \ {\isacharbar}{\kern0pt}\ Assign\ {\isacharparenleft}{\kern0pt}x{\isacharcomma}{\kern0pt}\ expr{\isacharparenright}{\kern0pt}\ {\isacharequal}{\kern0pt}{\isachargreater}{\kern0pt}\ \isanewline
\ \ \ \ \ \ let\ \isanewline
\ \ \ \ \ \ \ \ val\ new{\isacharunderscore}{\kern0pt}value\ {\isacharequal}{\kern0pt}\ expr\ s\isanewline
\ \ \ \ \ \ \ \ val\ new{\isacharunderscore}{\kern0pt}state\ {\isacharequal}{\kern0pt}\ fn\ v\ {\isacharequal}{\kern0pt}{\isachargreater}{\kern0pt}\ if\ v\ {\isacharequal}{\kern0pt}\ x\ then\ new{\isacharunderscore}{\kern0pt}value\ else\ s\ v\isanewline
\ \ \ \ \ \ in\ \isanewline
\ \ \ \ \ \ \ \ cont\ {\isacharparenleft}{\kern0pt}new{\isacharunderscore}{\kern0pt}state{\isacharcomma}{\kern0pt}\ trace{\isacharparenright}{\kern0pt}\ \isanewline
\ \ \ \ \ \ end\isanewline
\ \ {\isacharbar}{\kern0pt}\ Seq\ {\isacharparenleft}{\kern0pt}p{\isacharcomma}{\kern0pt}\ q{\isacharparenright}{\kern0pt}\ {\isacharequal}{\kern0pt}{\isachargreater}{\kern0pt}\ \isanewline
\ \ \ \ \ \ sem{\isacharunderscore}{\kern0pt}step\ thy\ machine\ p\ {\isacharparenleft}{\kern0pt}fn\ {\isacharparenleft}{\kern0pt}s{\isacharprime}{\kern0pt}{\isacharcomma}{\kern0pt}\ trace{\isacharprime}{\kern0pt}{\isacharparenright}{\kern0pt}\ {\isacharequal}{\kern0pt}{\isachargreater}{\kern0pt}\ sem{\isacharunderscore}{\kern0pt}step\ thy\ machine\ \ q\ cont\ {\isacharparenleft}{\kern0pt}s{\isacharprime}{\kern0pt}{\isacharcomma}{\kern0pt}\ trace{\isacharprime}{\kern0pt}{\isacharparenright}{\kern0pt}{\isacharparenright}{\kern0pt}\ {\isacharparenleft}{\kern0pt}s{\isacharcomma}{\kern0pt}\ trace{\isacharparenright}{\kern0pt}\isanewline
\ \ {\isacharbar}{\kern0pt}\ Cond\ {\isacharparenleft}{\kern0pt}guard{\isacharcomma}{\kern0pt}\ c{\isadigit{1}}{\isacharcomma}{\kern0pt}\ c{\isadigit{2}}{\isacharparenright}{\kern0pt}\ {\isacharequal}{\kern0pt}{\isachargreater}{\kern0pt}\isanewline
\ \ \ \ \ \ if\ guard\ s\ \isanewline
\ \ \ \ \ \ then\ sem{\isacharunderscore}{\kern0pt}step\ thy\ \ machine\ \ c{\isadigit{1}}\ cont\ {\isacharparenleft}{\kern0pt}s{\isacharcomma}{\kern0pt}\ trace{\isacharparenright}{\kern0pt}\isanewline
\ \ \ \ \ \ else\ sem{\isacharunderscore}{\kern0pt}step\ thy\ \ machine\ c{\isadigit{2}}\ cont\ {\isacharparenleft}{\kern0pt}s{\isacharcomma}{\kern0pt}\ trace{\isacharparenright}{\kern0pt}\isanewline
\ \ {\isacharbar}{\kern0pt}\ While\ {\isacharparenleft}{\kern0pt}guard{\isacharcomma}{\kern0pt}\ body{\isacharparenright}{\kern0pt}\ {\isacharequal}{\kern0pt}{\isachargreater}{\kern0pt}\isanewline
\ \ \ \ \ \ if\ guard\ s\ \isanewline
\ \ \ \ \ \ then\ sem{\isacharunderscore}{\kern0pt}step\ thy\ \ machine\ body\ {\isacharparenleft}{\kern0pt}fn\ {\isacharparenleft}{\kern0pt}s{\isacharprime}{\kern0pt}{\isacharcomma}{\kern0pt}\ trace{\isacharprime}{\kern0pt}{\isacharparenright}{\kern0pt}\ {\isacharequal}{\kern0pt}{\isachargreater}{\kern0pt}\ \isanewline
\ \ \ \ \ \ \ \ \ \ \ \ \ sem{\isacharunderscore}{\kern0pt}step\ thy\ \ machine\ {\isacharparenleft}{\kern0pt}While\ {\isacharparenleft}{\kern0pt}guard{\isacharcomma}{\kern0pt}\ body{\isacharparenright}{\kern0pt}{\isacharparenright}{\kern0pt}\ cont\ {\isacharparenleft}{\kern0pt}s{\isacharprime}{\kern0pt}{\isacharcomma}{\kern0pt}\ trace{\isacharprime}{\kern0pt}{\isacharparenright}{\kern0pt}{\isacharparenright}{\kern0pt}\ {\isacharparenleft}{\kern0pt}s{\isacharcomma}{\kern0pt}\ trace{\isacharparenright}{\kern0pt}\isanewline
\ \ \ \ \ \ else\ cont\ {\isacharparenleft}{\kern0pt}s{\isacharcomma}{\kern0pt}\ trace{\isacharparenright}{\kern0pt}\isanewline
\ \ {\isacharbar}{\kern0pt}\ Lock\ lock{\isacharunderscore}{\kern0pt}name\ {\isacharequal}{\kern0pt}{\isachargreater}{\kern0pt}\ \isanewline
\ \ \ \ \ \ let\isanewline
\ \ \ \ \ \ \ \ val\ lock{\isacharunderscore}{\kern0pt}id\ {\isacharequal}{\kern0pt}\ get{\isacharunderscore}{\kern0pt}lock{\isacharunderscore}{\kern0pt}id\ lock{\isacharunderscore}{\kern0pt}name\ machine\isanewline
\ \ \ \ \ \ \ \ val\ new{\isacharunderscore}{\kern0pt}trace\ {\isacharequal}{\kern0pt}\ LockEvent\ {\isacharparenleft}{\kern0pt}lock{\isacharunderscore}{\kern0pt}name{\isacharcomma}{\kern0pt}\ lock{\isacharunderscore}{\kern0pt}id{\isacharparenright}{\kern0pt}\ {\isacharcolon}{\kern0pt}{\isacharcolon}{\kern0pt}\ trace\isanewline
\ \ \ \ \ \ in\isanewline
\ \ \ \ \ \ \ \ cont\ {\isacharparenleft}{\kern0pt}s{\isacharcomma}{\kern0pt}\ new{\isacharunderscore}{\kern0pt}trace{\isacharparenright}{\kern0pt}\isanewline
\ \ \ \ \ \ end\isanewline
\ \ {\isacharbar}{\kern0pt}\ Unlock\ lock{\isacharunderscore}{\kern0pt}name\ {\isacharequal}{\kern0pt}{\isachargreater}{\kern0pt}\isanewline
\ \ \ \ \ \ let\isanewline
\ \ \ \ \ \ \ \ val\ lock{\isacharunderscore}{\kern0pt}id\ {\isacharequal}{\kern0pt}\ get{\isacharunderscore}{\kern0pt}lock{\isacharunderscore}{\kern0pt}id\ lock{\isacharunderscore}{\kern0pt}name\ machine\isanewline
\ \ \ \ \ \ \ \ val\ new{\isacharunderscore}{\kern0pt}trace\ {\isacharequal}{\kern0pt}\ UnlockEvent\ {\isacharparenleft}{\kern0pt}lock{\isacharunderscore}{\kern0pt}name{\isacharcomma}{\kern0pt}\ lock{\isacharunderscore}{\kern0pt}id{\isacharparenright}{\kern0pt}\ {\isacharcolon}{\kern0pt}{\isacharcolon}{\kern0pt}\ trace\isanewline
\ \ \ \ \ \ in\isanewline
\ \ \ \ \ \ \ \ cont\ {\isacharparenleft}{\kern0pt}s{\isacharcomma}{\kern0pt}\ new{\isacharunderscore}{\kern0pt}trace{\isacharparenright}{\kern0pt}\isanewline
\ \ \ \ \ \ end\isanewline
\ \ {\isacharbar}{\kern0pt}\ Store\ {\isacharparenleft}{\kern0pt}expr{\isacharcomma}{\kern0pt}\ var{\isacharparenright}{\kern0pt}\ {\isacharequal}{\kern0pt}{\isachargreater}{\kern0pt}\isanewline
\ \ \ \ let\isanewline
\ \ \ \ \ \ val\ value\ {\isacharequal}{\kern0pt}\ expr\ s\isanewline
\ \ \ \ \ \ val\ new{\isacharunderscore}{\kern0pt}state\ {\isacharequal}{\kern0pt}\ update{\isacharunderscore}{\kern0pt}state\ s\ var\ value\isanewline
\ \ \ \ \ \ val\ new{\isacharunderscore}{\kern0pt}trace\ {\isacharequal}{\kern0pt}\ UpdateEvent\ {\isacharparenleft}{\kern0pt}var{\isacharcomma}{\kern0pt}\ value{\isacharparenright}{\kern0pt}\ {\isacharcolon}{\kern0pt}{\isacharcolon}{\kern0pt}\ trace\isanewline
\ \ \ \ in\isanewline
\ \ \ \ \ \ cont\ {\isacharparenleft}{\kern0pt}new{\isacharunderscore}{\kern0pt}state{\isacharcomma}{\kern0pt}\ new{\isacharunderscore}{\kern0pt}trace{\isacharparenright}{\kern0pt}\isanewline
\ \ \ \ end\isanewline
\ \ {\isacharbar}{\kern0pt}\ Load\ {\isacharparenleft}{\kern0pt}local{\isacharunderscore}{\kern0pt}var{\isacharcomma}{\kern0pt}\ global{\isacharunderscore}{\kern0pt}var{\isacharparenright}{\kern0pt}\ {\isacharequal}{\kern0pt}{\isachargreater}{\kern0pt}\isanewline
\ \ \ \ \ \ let\isanewline
\ \ \ \ \ \ \ \ {\isacharparenleft}{\kern0pt}{\isacharasterisk}{\kern0pt}TODO{\isacharasterisk}{\kern0pt}{\isacharparenright}{\kern0pt}\isanewline
\ \ \ \ \ \ \ \ val\ read{\isacharunderscore}{\kern0pt}value\ {\isacharequal}{\kern0pt}\ s\ global{\isacharunderscore}{\kern0pt}var\isanewline
\ \ \ \ \ \ \ \ val\ new{\isacharunderscore}{\kern0pt}state\ {\isacharequal}{\kern0pt}\ fn\ v\ {\isacharequal}{\kern0pt}{\isachargreater}{\kern0pt}\ if\ v\ {\isacharequal}{\kern0pt}\ local{\isacharunderscore}{\kern0pt}var\ then\ read{\isacharunderscore}{\kern0pt}value\ else\ s\ v\isanewline
\ \ \ \ \ \ \ \ val\ new{\isacharunderscore}{\kern0pt}trace\ {\isacharequal}{\kern0pt}\ ReadEvent\ {\isacharparenleft}{\kern0pt}global{\isacharunderscore}{\kern0pt}var{\isacharcomma}{\kern0pt}\ read{\isacharunderscore}{\kern0pt}value{\isacharparenright}{\kern0pt}\ {\isacharcolon}{\kern0pt}{\isacharcolon}{\kern0pt}\ trace\isanewline
\ \ \ \ \ \ in\isanewline
\ \ \ \ \ \ \ \ cont\ {\isacharparenleft}{\kern0pt}new{\isacharunderscore}{\kern0pt}state{\isacharcomma}{\kern0pt}\ new{\isacharunderscore}{\kern0pt}trace{\isacharparenright}{\kern0pt}\isanewline
\ \ \ \ \ \ end\isanewline
\isanewline
{\isacharparenleft}{\kern0pt}{\isacharasterisk}{\kern0pt}\ Analyses\ sémantiques\ avancées\ {\isacharasterisk}{\kern0pt}{\isacharparenright}{\kern0pt}\isanewline
\isanewline
{\isacharparenleft}{\kern0pt}{\isacharasterisk}{\kern0pt}\ Détection\ de\ deadlocks\ simples\ {\isacharasterisk}{\kern0pt}{\isacharparenright}{\kern0pt}\isanewline
fun\ detect{\isacharunderscore}{\kern0pt}potential{\isacharunderscore}{\kern0pt}deadlock\ {\isacharparenleft}{\kern0pt}threads\ {\isacharcolon}{\kern0pt}\ thread{\isacharunderscore}{\kern0pt}absy\ list{\isacharparenright}{\kern0pt}\ {\isacharcolon}{\kern0pt}\ string\ list\ {\isacharequal}{\kern0pt}\isanewline
\ \ let\isanewline
\ \ \ \ fun\ extract{\isacharunderscore}{\kern0pt}locks{\isacharunderscore}{\kern0pt}from{\isacharunderscore}{\kern0pt}actions\ {\isacharparenleft}{\kern0pt}actions\ {\isacharcolon}{\kern0pt}\ a{\isacharunderscore}{\kern0pt}term\ list{\isacharparenright}{\kern0pt}\ {\isacharcolon}{\kern0pt}\ string\ list\ {\isacharequal}{\kern0pt}\isanewline
\ \ \ \ \ \ let\isanewline
\ \ \ \ \ \ \ \ fun\ extract{\isacharunderscore}{\kern0pt}from{\isacharunderscore}{\kern0pt}action\ action\ {\isacharequal}{\kern0pt}\isanewline
\ \ \ \ \ \ \ \ \ \ case\ action\ of\isanewline
\ \ \ \ \ \ \ \ \ \ \ \ LockA\ bdg\ {\isacharequal}{\kern0pt}{\isachargreater}{\kern0pt}\ {\isacharbrackleft}{\kern0pt}Binding{\isachardot}{\kern0pt}name{\isacharunderscore}{\kern0pt}of\ bdg{\isacharbrackright}{\kern0pt}\isanewline
\ \ \ \ \ \ \ \ \ \ {\isacharbar}{\kern0pt}\ UnlockA\ bdg\ {\isacharequal}{\kern0pt}{\isachargreater}{\kern0pt}\ {\isacharbrackleft}{\kern0pt}{\isacharbrackright}{\kern0pt}\isanewline
\ \ \ \ \ \ \ \ \ \ {\isacharbar}{\kern0pt}\ IfelseA\ {\isacharparenleft}{\kern0pt}{\isacharparenleft}{\kern0pt}{\isacharunderscore}{\kern0pt}{\isacharcomma}{\kern0pt}\ then{\isacharunderscore}{\kern0pt}actions{\isacharparenright}{\kern0pt}{\isacharcomma}{\kern0pt}\ else{\isacharunderscore}{\kern0pt}actions{\isacharparenright}{\kern0pt}\ {\isacharequal}{\kern0pt}{\isachargreater}{\kern0pt}\isanewline
\ \ \ \ \ \ \ \ \ \ \ \ \ \ {\isacharparenleft}{\kern0pt}extract{\isacharunderscore}{\kern0pt}locks{\isacharunderscore}{\kern0pt}from{\isacharunderscore}{\kern0pt}actions\ then{\isacharunderscore}{\kern0pt}actions{\isacharparenright}{\kern0pt}\ {\isacharat}{\kern0pt}\ \isanewline
\ \ \ \ \ \ \ \ \ \ \ \ \ \ {\isacharparenleft}{\kern0pt}extract{\isacharunderscore}{\kern0pt}locks{\isacharunderscore}{\kern0pt}from{\isacharunderscore}{\kern0pt}actions\ else{\isacharunderscore}{\kern0pt}actions{\isacharparenright}{\kern0pt}\isanewline
\ \ \ \ \ \ \ \ \ \ {\isacharbar}{\kern0pt}\ WhileA\ {\isacharparenleft}{\kern0pt}{\isacharunderscore}{\kern0pt}{\isacharcomma}{\kern0pt}\ body{\isacharunderscore}{\kern0pt}actions{\isacharparenright}{\kern0pt}\ {\isacharequal}{\kern0pt}{\isachargreater}{\kern0pt}\isanewline
\ \ \ \ \ \ \ \ \ \ \ \ \ \ extract{\isacharunderscore}{\kern0pt}locks{\isacharunderscore}{\kern0pt}from{\isacharunderscore}{\kern0pt}actions\ body{\isacharunderscore}{\kern0pt}actions\isanewline
\ \ \ \ \ \ \ \ \ \ {\isacharbar}{\kern0pt}\ {\isacharunderscore}{\kern0pt}\ {\isacharequal}{\kern0pt}{\isachargreater}{\kern0pt}\ {\isacharbrackleft}{\kern0pt}{\isacharbrackright}{\kern0pt}\ \ \ \ \ \ \ \ \ \ \ \ \ \ \ \ \ \ \ \ \ \ \ \ \ \ \ \ \ \ \ \ \ \ \ \ \ \ \ \ \ \ \ \ \ \ \ \ \ \ \ \ \ \ \ \ \ \ \ \ \ \ \ \ \ \ \ \ \ \ \ \ \ \ \ \ \ \ \ \ \ \ \ \ \ \ \ \ \ \ \ \ \ \ \ \ \ \ \ \ \ \ \ \isanewline
\ \ \ \ \ \ in\isanewline
\ \ \ \ \ \ \ \ List{\isachardot}{\kern0pt}concat\ {\isacharparenleft}{\kern0pt}map\ extract{\isacharunderscore}{\kern0pt}from{\isacharunderscore}{\kern0pt}action\ actions{\isacharparenright}{\kern0pt}\isanewline
\ \ \ \ \ \ end\isanewline
\ \ \ \ \isanewline
\ \ \ \ fun\ get{\isacharunderscore}{\kern0pt}thread{\isacharunderscore}{\kern0pt}locks\ {\isacharparenleft}{\kern0pt}thread\ {\isacharcolon}{\kern0pt}\ thread{\isacharunderscore}{\kern0pt}absy{\isacharparenright}{\kern0pt}\ {\isacharcolon}{\kern0pt}\ string\ {\isacharasterisk}{\kern0pt}\ string\ list\ {\isacharequal}{\kern0pt}\isanewline
\ \ \ \ \ \ let\isanewline
\ \ \ \ \ \ \ \ val\ {\isacharbraceleft}{\kern0pt}nom{\isacharunderscore}{\kern0pt}thread{\isacharcomma}{\kern0pt}\ actions{\isacharcomma}{\kern0pt}\ {\isachardot}{\kern0pt}{\isachardot}{\kern0pt}{\isachardot}{\kern0pt}{\isacharbraceright}{\kern0pt}\ {\isacharequal}{\kern0pt}\ thread\isanewline
\ \ \ \ \ \ \ \ val\ locks\ {\isacharequal}{\kern0pt}\ extract{\isacharunderscore}{\kern0pt}locks{\isacharunderscore}{\kern0pt}from{\isacharunderscore}{\kern0pt}actions\ actions\isanewline
\ \ \ \ \ \ in\isanewline
\ \ \ \ \ \ \ \ {\isacharparenleft}{\kern0pt}Binding{\isachardot}{\kern0pt}name{\isacharunderscore}{\kern0pt}of\ nom{\isacharunderscore}{\kern0pt}thread{\isacharcomma}{\kern0pt}\ locks{\isacharparenright}{\kern0pt}\isanewline
\ \ \ \ \ \ end\isanewline
\ \ \ \ \isanewline
\ \ \ \ val\ thread{\isacharunderscore}{\kern0pt}locks\ {\isacharequal}{\kern0pt}\ map\ get{\isacharunderscore}{\kern0pt}thread{\isacharunderscore}{\kern0pt}locks\ threads\isanewline
\ \ \ \ \isanewline
\ \ \ \ {\isacharparenleft}{\kern0pt}{\isacharasterisk}{\kern0pt}\ Détection\ simple\ {\isacharcolon}{\kern0pt}\ si\ deux\ threads\ utilisent\ les\ mêmes\ verrous\ {\isacharasterisk}{\kern0pt}{\isacharparenright}{\kern0pt}\isanewline
\ \ \ \ fun\ find{\isacharunderscore}{\kern0pt}shared{\isacharunderscore}{\kern0pt}locks\ {\isacharbrackleft}{\kern0pt}{\isacharbrackright}{\kern0pt}\ {\isacharequal}{\kern0pt}\ {\isacharbrackleft}{\kern0pt}{\isacharbrackright}{\kern0pt}\isanewline
\ \ \ \ \ \ {\isacharbar}{\kern0pt}\ find{\isacharunderscore}{\kern0pt}shared{\isacharunderscore}{\kern0pt}locks\ {\isacharparenleft}{\kern0pt}{\isacharparenleft}{\kern0pt}t{\isadigit{1}}{\isacharcomma}{\kern0pt}\ locks{\isadigit{1}}{\isacharparenright}{\kern0pt}\ {\isacharcolon}{\kern0pt}{\isacharcolon}{\kern0pt}\ rest{\isacharparenright}{\kern0pt}\ {\isacharequal}{\kern0pt}\isanewline
\ \ \ \ \ \ \ \ \ \ let\isanewline
\ \ \ \ \ \ \ \ \ \ \ \ fun\ check{\isacharunderscore}{\kern0pt}overlap\ {\isacharparenleft}{\kern0pt}t{\isadigit{2}}{\isacharcomma}{\kern0pt}\ locks{\isadigit{2}}{\isacharparenright}{\kern0pt}\ {\isacharequal}{\kern0pt}\isanewline
\ \ \ \ \ \ \ \ \ \ \ \ \ \ let\isanewline
\ \ \ \ \ \ \ \ \ \ \ \ \ \ \ \ val\ shared\ {\isacharequal}{\kern0pt}\ List{\isachardot}{\kern0pt}filter\ {\isacharparenleft}{\kern0pt}fn\ l\ {\isacharequal}{\kern0pt}{\isachargreater}{\kern0pt}\ List{\isachardot}{\kern0pt}exists\ {\isacharparenleft}{\kern0pt}fn\ l{\isadigit{2}}\ {\isacharequal}{\kern0pt}{\isachargreater}{\kern0pt}\ l\ {\isacharequal}{\kern0pt}\ l{\isadigit{2}}{\isacharparenright}{\kern0pt}\ locks{\isadigit{2}}{\isacharparenright}{\kern0pt}\ locks{\isadigit{1}}\isanewline
\ \ \ \ \ \ \ \ \ \ \ \ \ \ in\isanewline
\ \ \ \ \ \ \ \ \ \ \ \ \ \ \ \ if\ null\ shared\ then\ NONE\isanewline
\ \ \ \ \ \ \ \ \ \ \ \ \ \ \ \ else\ SOME\ {\isacharparenleft}{\kern0pt}{\isachardoublequote}{\kern0pt}Potential\ deadlock\ between\ {\isachardoublequote}{\kern0pt}\ {\isacharcircum}{\kern0pt}\ t{\isadigit{1}}\ {\isacharcircum}{\kern0pt}\ {\isachardoublequote}{\kern0pt}\ and\ {\isachardoublequote}{\kern0pt}\ {\isacharcircum}{\kern0pt}\ t{\isadigit{2}}\ {\isacharcircum}{\kern0pt}\ \isanewline
\ \ \ \ \ \ \ \ \ \ \ \ \ \ \ \ \ \ \ \ \ \ \ \ \ \ {\isachardoublequote}{\kern0pt}\ on\ locks{\isacharcolon}{\kern0pt}\ {\isachardoublequote}{\kern0pt}\ {\isacharcircum}{\kern0pt}\ String{\isachardot}{\kern0pt}concatWith\ {\isachardoublequote}{\kern0pt}{\isacharcomma}{\kern0pt}\ {\isachardoublequote}{\kern0pt}\ shared{\isacharparenright}{\kern0pt}\isanewline
\ \ \ \ \ \ \ \ \ \ \ \ \ \ end\isanewline
\ \ \ \ \ \ \ \ \ \ \ \ val\ warnings\ {\isacharequal}{\kern0pt}\ List{\isachardot}{\kern0pt}mapPartial\ check{\isacharunderscore}{\kern0pt}overlap\ rest\isanewline
\ \ \ \ \ \ \ \ \ \ in\isanewline
\ \ \ \ \ \ \ \ \ \ \ \ warnings\ {\isacharat}{\kern0pt}\ find{\isacharunderscore}{\kern0pt}shared{\isacharunderscore}{\kern0pt}locks\ rest\isanewline
\ \ \ \ \ \ \ \ \ \ end\isanewline
\ \ in\isanewline
\ \ \ \ find{\isacharunderscore}{\kern0pt}shared{\isacharunderscore}{\kern0pt}locks\ thread{\isacharunderscore}{\kern0pt}locks\isanewline
\ \ end\isanewline
\isanewline
{\isacharparenleft}{\kern0pt}{\isacharasterisk}{\kern0pt}\ Détection\ de\ race\ conditions\ {\isacharasterisk}{\kern0pt}{\isacharparenright}{\kern0pt}\isanewline
fun\ detect{\isacharunderscore}{\kern0pt}race{\isacharunderscore}{\kern0pt}conditions\ {\isacharparenleft}{\kern0pt}threads\ {\isacharcolon}{\kern0pt}\ thread{\isacharunderscore}{\kern0pt}absy\ list{\isacharparenright}{\kern0pt}\ {\isacharcolon}{\kern0pt}\ string\ list\ {\isacharequal}{\kern0pt}\isanewline
\ \ let\isanewline
\ \ \ \ fun\ extract{\isacharunderscore}{\kern0pt}global{\isacharunderscore}{\kern0pt}vars{\isacharunderscore}{\kern0pt}from{\isacharunderscore}{\kern0pt}actions\ {\isacharparenleft}{\kern0pt}actions\ {\isacharcolon}{\kern0pt}\ a{\isacharunderscore}{\kern0pt}term\ list{\isacharparenright}{\kern0pt}\ {\isacharcolon}{\kern0pt}\ {\isacharparenleft}{\kern0pt}string\ {\isacharasterisk}{\kern0pt}\ bool{\isacharparenright}{\kern0pt}\ list\ {\isacharequal}{\kern0pt}\isanewline
\ \ \ \ \ \ let\isanewline
\ \ \ \ \ \ \ \ fun\ extract{\isacharunderscore}{\kern0pt}from{\isacharunderscore}{\kern0pt}action\ action\ {\isacharequal}{\kern0pt}\isanewline
\ \ \ \ \ \ \ \ \ \ case\ action\ of\isanewline
\ \ \ \ \ \ \ \ \ \ \ \ SendA\ {\isacharparenleft}{\kern0pt}bdg{\isacharcomma}{\kern0pt}\ {\isacharunderscore}{\kern0pt}{\isacharparenright}{\kern0pt}\ {\isacharequal}{\kern0pt}{\isachargreater}{\kern0pt}\ {\isacharbrackleft}{\kern0pt}{\isacharparenleft}{\kern0pt}Binding{\isachardot}{\kern0pt}name{\isacharunderscore}{\kern0pt}of\ bdg{\isacharcomma}{\kern0pt}\ true{\isacharparenright}{\kern0pt}{\isacharbrackright}{\kern0pt}\ \ {\isacharparenleft}{\kern0pt}{\isacharasterisk}{\kern0pt}\ write\ {\isacharasterisk}{\kern0pt}{\isacharparenright}{\kern0pt}\isanewline
\ \ \ \ \ \ \ \ \ \ {\isacharbar}{\kern0pt}\ ReceiveA\ {\isacharparenleft}{\kern0pt}{\isacharunderscore}{\kern0pt}{\isacharcomma}{\kern0pt}\ bdg{\isacharparenright}{\kern0pt}\ {\isacharequal}{\kern0pt}{\isachargreater}{\kern0pt}\ {\isacharbrackleft}{\kern0pt}{\isacharparenleft}{\kern0pt}Binding{\isachardot}{\kern0pt}name{\isacharunderscore}{\kern0pt}of\ bdg{\isacharcomma}{\kern0pt}\ false{\isacharparenright}{\kern0pt}{\isacharbrackright}{\kern0pt}\ {\isacharparenleft}{\kern0pt}{\isacharasterisk}{\kern0pt}\ read\ {\isacharasterisk}{\kern0pt}{\isacharparenright}{\kern0pt}\isanewline
\ \ \ \ \ \ \ \ \ \ {\isacharbar}{\kern0pt}\ IfelseA\ {\isacharparenleft}{\kern0pt}{\isacharparenleft}{\kern0pt}{\isacharunderscore}{\kern0pt}{\isacharcomma}{\kern0pt}\ then{\isacharunderscore}{\kern0pt}actions{\isacharparenright}{\kern0pt}{\isacharcomma}{\kern0pt}\ else{\isacharunderscore}{\kern0pt}actions{\isacharparenright}{\kern0pt}\ {\isacharequal}{\kern0pt}{\isachargreater}{\kern0pt}\isanewline
\ \ \ \ \ \ \ \ \ \ \ \ \ \ {\isacharparenleft}{\kern0pt}extract{\isacharunderscore}{\kern0pt}global{\isacharunderscore}{\kern0pt}vars{\isacharunderscore}{\kern0pt}from{\isacharunderscore}{\kern0pt}actions\ then{\isacharunderscore}{\kern0pt}actions{\isacharparenright}{\kern0pt}\ {\isacharat}{\kern0pt}\ \isanewline
\ \ \ \ \ \ \ \ \ \ \ \ \ \ {\isacharparenleft}{\kern0pt}extract{\isacharunderscore}{\kern0pt}global{\isacharunderscore}{\kern0pt}vars{\isacharunderscore}{\kern0pt}from{\isacharunderscore}{\kern0pt}actions\ else{\isacharunderscore}{\kern0pt}actions{\isacharparenright}{\kern0pt}\isanewline
\ \ \ \ \ \ \ \ \ \ {\isacharbar}{\kern0pt}\ WhileA\ {\isacharparenleft}{\kern0pt}{\isacharunderscore}{\kern0pt}{\isacharcomma}{\kern0pt}\ body{\isacharunderscore}{\kern0pt}actions{\isacharparenright}{\kern0pt}\ {\isacharequal}{\kern0pt}{\isachargreater}{\kern0pt}\isanewline
\ \ \ \ \ \ \ \ \ \ \ \ \ \ extract{\isacharunderscore}{\kern0pt}global{\isacharunderscore}{\kern0pt}vars{\isacharunderscore}{\kern0pt}from{\isacharunderscore}{\kern0pt}actions\ body{\isacharunderscore}{\kern0pt}actions\isanewline
\ \ \ \ \ \ \ \ \ \ {\isacharbar}{\kern0pt}\ {\isacharunderscore}{\kern0pt}\ {\isacharequal}{\kern0pt}{\isachargreater}{\kern0pt}\ {\isacharbrackleft}{\kern0pt}{\isacharbrackright}{\kern0pt}\isanewline
\ \ \ \ \ \ in\isanewline
\ \ \ \ \ \ \ \ List{\isachardot}{\kern0pt}concat\ {\isacharparenleft}{\kern0pt}map\ extract{\isacharunderscore}{\kern0pt}from{\isacharunderscore}{\kern0pt}action\ actions{\isacharparenright}{\kern0pt}\isanewline
\ \ \ \ \ \ end\isanewline
\ \ \ \ \isanewline
\ \ \ \ fun\ get{\isacharunderscore}{\kern0pt}thread{\isacharunderscore}{\kern0pt}global{\isacharunderscore}{\kern0pt}access\ {\isacharparenleft}{\kern0pt}thread\ {\isacharcolon}{\kern0pt}\ thread{\isacharunderscore}{\kern0pt}absy{\isacharparenright}{\kern0pt}\ {\isacharcolon}{\kern0pt}\ string\ {\isacharasterisk}{\kern0pt}\ {\isacharparenleft}{\kern0pt}string\ {\isacharasterisk}{\kern0pt}\ bool{\isacharparenright}{\kern0pt}\ list\ {\isacharequal}{\kern0pt}\isanewline
\ \ \ \ \ \ let\isanewline
\ \ \ \ \ \ \ \ val\ {\isacharbraceleft}{\kern0pt}nom{\isacharunderscore}{\kern0pt}thread{\isacharcomma}{\kern0pt}\ actions{\isacharcomma}{\kern0pt}\ {\isachardot}{\kern0pt}{\isachardot}{\kern0pt}{\isachardot}{\kern0pt}{\isacharbraceright}{\kern0pt}\ {\isacharequal}{\kern0pt}\ thread\isanewline
\ \ \ \ \ \ \ \ val\ accesses\ {\isacharequal}{\kern0pt}\ extract{\isacharunderscore}{\kern0pt}global{\isacharunderscore}{\kern0pt}vars{\isacharunderscore}{\kern0pt}from{\isacharunderscore}{\kern0pt}actions\ actions\isanewline
\ \ \ \ \ \ in\isanewline
\ \ \ \ \ \ \ \ {\isacharparenleft}{\kern0pt}Binding{\isachardot}{\kern0pt}name{\isacharunderscore}{\kern0pt}of\ nom{\isacharunderscore}{\kern0pt}thread{\isacharcomma}{\kern0pt}\ accesses{\isacharparenright}{\kern0pt}\isanewline
\ \ \ \ \ \ end\isanewline
\ \ \ \ \isanewline
\ \ \ \ val\ thread{\isacharunderscore}{\kern0pt}accesses\ {\isacharequal}{\kern0pt}\ map\ get{\isacharunderscore}{\kern0pt}thread{\isacharunderscore}{\kern0pt}global{\isacharunderscore}{\kern0pt}access\ threads\isanewline
\ \ \ \ \isanewline
\ \ \ \ {\isacharparenleft}{\kern0pt}{\isacharasterisk}{\kern0pt}\ Détection\ {\isacharcolon}{\kern0pt}\ si\ deux\ threads\ accèdent\ à\ la\ même\ variable\ et\ au\ moins\ un\ écrit\ {\isacharasterisk}{\kern0pt}{\isacharparenright}{\kern0pt}\isanewline
\ \ \ \ fun\ find{\isacharunderscore}{\kern0pt}races\ {\isacharbrackleft}{\kern0pt}{\isacharbrackright}{\kern0pt}\ {\isacharequal}{\kern0pt}\ {\isacharbrackleft}{\kern0pt}{\isacharbrackright}{\kern0pt}\isanewline
\ \ \ \ \ \ {\isacharbar}{\kern0pt}\ find{\isacharunderscore}{\kern0pt}races\ {\isacharparenleft}{\kern0pt}{\isacharparenleft}{\kern0pt}t{\isadigit{1}}{\isacharcomma}{\kern0pt}\ accesses{\isadigit{1}}{\isacharparenright}{\kern0pt}\ {\isacharcolon}{\kern0pt}{\isacharcolon}{\kern0pt}\ rest{\isacharparenright}{\kern0pt}\ {\isacharequal}{\kern0pt}\isanewline
\ \ \ \ \ \ \ \ \ \ let\isanewline
\ \ \ \ \ \ \ \ \ \ \ \ fun\ check{\isacharunderscore}{\kern0pt}race\ {\isacharparenleft}{\kern0pt}t{\isadigit{2}}{\isacharcomma}{\kern0pt}\ accesses{\isadigit{2}}{\isacharparenright}{\kern0pt}\ {\isacharequal}{\kern0pt}\isanewline
\ \ \ \ \ \ \ \ \ \ \ \ \ \ let\isanewline
\ \ \ \ \ \ \ \ \ \ \ \ \ \ \ \ fun\ has{\isacharunderscore}{\kern0pt}conflict\ var\ {\isacharequal}{\kern0pt}\isanewline
\ \ \ \ \ \ \ \ \ \ \ \ \ \ \ \ \ \ let\isanewline
\ \ \ \ \ \ \ \ \ \ \ \ \ \ \ \ \ \ \ \ val\ t{\isadigit{1}}{\isacharunderscore}{\kern0pt}writes\ {\isacharequal}{\kern0pt}\ List{\isachardot}{\kern0pt}exists\ {\isacharparenleft}{\kern0pt}fn\ {\isacharparenleft}{\kern0pt}v{\isacharcomma}{\kern0pt}\ is{\isacharunderscore}{\kern0pt}write{\isacharparenright}{\kern0pt}\ {\isacharequal}{\kern0pt}{\isachargreater}{\kern0pt}\ v\ {\isacharequal}{\kern0pt}\ var\ andalso\ is{\isacharunderscore}{\kern0pt}write{\isacharparenright}{\kern0pt}\ accesses{\isadigit{1}}\isanewline
\ \ \ \ \ \ \ \ \ \ \ \ \ \ \ \ \ \ \ \ val\ t{\isadigit{1}}{\isacharunderscore}{\kern0pt}reads\ {\isacharequal}{\kern0pt}\ List{\isachardot}{\kern0pt}exists\ {\isacharparenleft}{\kern0pt}fn\ {\isacharparenleft}{\kern0pt}v{\isacharcomma}{\kern0pt}\ is{\isacharunderscore}{\kern0pt}write{\isacharparenright}{\kern0pt}\ {\isacharequal}{\kern0pt}{\isachargreater}{\kern0pt}\ v\ {\isacharequal}{\kern0pt}\ var\ andalso\ not\ is{\isacharunderscore}{\kern0pt}write{\isacharparenright}{\kern0pt}\ accesses{\isadigit{1}}\isanewline
\ \ \ \ \ \ \ \ \ \ \ \ \ \ \ \ \ \ \ \ val\ t{\isadigit{2}}{\isacharunderscore}{\kern0pt}writes\ {\isacharequal}{\kern0pt}\ List{\isachardot}{\kern0pt}exists\ {\isacharparenleft}{\kern0pt}fn\ {\isacharparenleft}{\kern0pt}v{\isacharcomma}{\kern0pt}\ is{\isacharunderscore}{\kern0pt}write{\isacharparenright}{\kern0pt}\ {\isacharequal}{\kern0pt}{\isachargreater}{\kern0pt}\ v\ {\isacharequal}{\kern0pt}\ var\ andalso\ is{\isacharunderscore}{\kern0pt}write{\isacharparenright}{\kern0pt}\ accesses{\isadigit{2}}\isanewline
\ \ \ \ \ \ \ \ \ \ \ \ \ \ \ \ \ \ \ \ val\ t{\isadigit{2}}{\isacharunderscore}{\kern0pt}reads\ {\isacharequal}{\kern0pt}\ List{\isachardot}{\kern0pt}exists\ {\isacharparenleft}{\kern0pt}fn\ {\isacharparenleft}{\kern0pt}v{\isacharcomma}{\kern0pt}\ is{\isacharunderscore}{\kern0pt}write{\isacharparenright}{\kern0pt}\ {\isacharequal}{\kern0pt}{\isachargreater}{\kern0pt}\ v\ {\isacharequal}{\kern0pt}\ var\ andalso\ not\ is{\isacharunderscore}{\kern0pt}write{\isacharparenright}{\kern0pt}\ accesses{\isadigit{2}}\isanewline
\ \ \ \ \ \ \ \ \ \ \ \ \ \ \ \ \ \ in\isanewline
\ \ \ \ \ \ \ \ \ \ \ \ \ \ \ \ \ \ \ \ {\isacharparenleft}{\kern0pt}t{\isadigit{1}}{\isacharunderscore}{\kern0pt}writes\ orelse\ t{\isadigit{1}}{\isacharunderscore}{\kern0pt}reads{\isacharparenright}{\kern0pt}\ andalso\ {\isacharparenleft}{\kern0pt}t{\isadigit{2}}{\isacharunderscore}{\kern0pt}writes\ orelse\ t{\isadigit{2}}{\isacharunderscore}{\kern0pt}reads{\isacharparenright}{\kern0pt}\ andalso\isanewline
\ \ \ \ \ \ \ \ \ \ \ \ \ \ \ \ \ \ \ \ {\isacharparenleft}{\kern0pt}t{\isadigit{1}}{\isacharunderscore}{\kern0pt}writes\ orelse\ t{\isadigit{2}}{\isacharunderscore}{\kern0pt}writes{\isacharparenright}{\kern0pt}\isanewline
\ \ \ \ \ \ \ \ \ \ \ \ \ \ \ \ \ \ end\isanewline
\ \ \ \ \ \ \ \ \ \ \ \ \ \ \ \ \isanewline
\ \ \ \ \ \ \ \ \ \ \ \ \ \ \ \ val\ all{\isacharunderscore}{\kern0pt}vars\ {\isacharequal}{\kern0pt}\ map\ fst\ {\isacharparenleft}{\kern0pt}accesses{\isadigit{1}}\ {\isacharat}{\kern0pt}\ accesses{\isadigit{2}}{\isacharparenright}{\kern0pt}\isanewline
\ \ \ \ \ \ \ \ \ \ \ \ \ \ \ \ val\ conflicting{\isacharunderscore}{\kern0pt}vars\ {\isacharequal}{\kern0pt}\ List{\isachardot}{\kern0pt}filter\ has{\isacharunderscore}{\kern0pt}conflict\ all{\isacharunderscore}{\kern0pt}vars\isanewline
\ \ \ \ \ \ \ \ \ \ \ \ \ \ in\isanewline
\ \ \ \ \ \ \ \ \ \ \ \ \ \ \ \ if\ null\ conflicting{\isacharunderscore}{\kern0pt}vars\ then\ {\isacharbrackleft}{\kern0pt}{\isacharbrackright}{\kern0pt}\isanewline
\ \ \ \ \ \ \ \ \ \ \ \ \ \ \ \ else\ {\isacharbrackleft}{\kern0pt}{\isachardoublequote}{\kern0pt}Potential\ race\ condition\ between\ {\isachardoublequote}{\kern0pt}\ {\isacharcircum}{\kern0pt}\ t{\isadigit{1}}\ {\isacharcircum}{\kern0pt}\ {\isachardoublequote}{\kern0pt}\ and\ {\isachardoublequote}{\kern0pt}\ {\isacharcircum}{\kern0pt}\ t{\isadigit{2}}\ {\isacharcircum}{\kern0pt}\ \isanewline
\ \ \ \ \ \ \ \ \ \ \ \ \ \ \ \ \ \ \ \ \ {\isachardoublequote}{\kern0pt}\ on\ variables{\isacharcolon}{\kern0pt}\ {\isachardoublequote}{\kern0pt}\ {\isacharcircum}{\kern0pt}\ String{\isachardot}{\kern0pt}concatWith\ {\isachardoublequote}{\kern0pt}{\isacharcomma}{\kern0pt}\ {\isachardoublequote}{\kern0pt}\ conflicting{\isacharunderscore}{\kern0pt}vars{\isacharbrackright}{\kern0pt}\isanewline
\ \ \ \ \ \ \ \ \ \ \ \ \ \ end\isanewline
\ \ \ \ \ \ \ \ \ \ \ \ val\ warnings\ {\isacharequal}{\kern0pt}\ List{\isachardot}{\kern0pt}concat\ {\isacharparenleft}{\kern0pt}map\ check{\isacharunderscore}{\kern0pt}race\ rest{\isacharparenright}{\kern0pt}\isanewline
\ \ \ \ \ \ \ \ \ \ in\isanewline
\ \ \ \ \ \ \ \ \ \ \ \ warnings\ {\isacharat}{\kern0pt}\ find{\isacharunderscore}{\kern0pt}races\ rest\isanewline
\ \ \ \ \ \ \ \ \ \ end\isanewline
\ \ in\isanewline
\ \ \ \ find{\isacharunderscore}{\kern0pt}races\ thread{\isacharunderscore}{\kern0pt}accesses\isanewline
\ \ end\isanewline
{\isacharparenleft}{\kern0pt}{\isacharasterisk}{\kern0pt}\ Vérification\ de\ cohérence\ des\ types\ {\isacharasterisk}{\kern0pt}{\isacharparenright}{\kern0pt}\isanewline
fun\ type{\isacharunderscore}{\kern0pt}check{\isacharunderscore}{\kern0pt}thread\ {\isacharparenleft}{\kern0pt}thy\ {\isacharcolon}{\kern0pt}\ theory{\isacharparenright}{\kern0pt}\ {\isacharparenleft}{\kern0pt}thread\ {\isacharcolon}{\kern0pt}\ thread{\isacharunderscore}{\kern0pt}absy{\isacharparenright}{\kern0pt}\ {\isacharparenleft}{\kern0pt}vars{\isacharunderscore}{\kern0pt}decl{\isacharparenright}{\kern0pt}\ {\isacharcolon}{\kern0pt}\ string\ list\ {\isacharequal}{\kern0pt}\isanewline
\ \ let\isanewline
\ \ \ \ val\ {\isacharbraceleft}{\kern0pt}nom{\isacharunderscore}{\kern0pt}thread{\isacharcomma}{\kern0pt}\ locals{\isacharunderscore}{\kern0pt}decl{\isacharcomma}{\kern0pt}\ actions{\isacharcomma}{\kern0pt}\ locstab{\isacharbraceright}{\kern0pt}\ {\isacharequal}{\kern0pt}\ thread\isanewline
\ \ \ \ fun\ check{\isacharunderscore}{\kern0pt}action{\isacharunderscore}{\kern0pt}types\ {\isacharparenleft}{\kern0pt}action{\isacharcomma}{\kern0pt}\ errors{\isacharparenright}{\kern0pt}\ {\isacharequal}{\kern0pt}\ \ \ \ \ \ case\ action\ of\isanewline
\ \ \ \ \ \ \ \ AssignA\ {\isacharparenleft}{\kern0pt}bdg{\isacharcomma}{\kern0pt}\ term{\isacharparenright}{\kern0pt}\ {\isacharequal}{\kern0pt}{\isachargreater}{\kern0pt}\isanewline
\ \ \ \ \ \ \ \ \ \ let\isanewline
\ \ \ \ \ \ \ \ \ \ \ \ val\ var{\isacharunderscore}{\kern0pt}name\ {\isacharequal}{\kern0pt}\ Binding{\isachardot}{\kern0pt}name{\isacharunderscore}{\kern0pt}of\ bdg\isanewline
\ \ \ \ \ \ \ \ \ \ \ \ val\ term{\isacharunderscore}{\kern0pt}type\ {\isacharequal}{\kern0pt}\ fastype{\isacharunderscore}{\kern0pt}of\ term\isanewline
\ \ \ \ \ \ \ \ \ \ \ \ val\ expected{\isacharunderscore}{\kern0pt}type{\isacharunderscore}{\kern0pt}opt\ {\isacharequal}{\kern0pt}\ \isanewline
\ \ \ \ \ \ \ \ \ \ \ \ \ \ \ \ \ \ \ \ \ \ \ \ \ \ case\ find{\isacharunderscore}{\kern0pt}first\ {\isacharparenleft}{\kern0pt}fn\ {\isacharparenleft}{\kern0pt}{\isacharparenleft}{\kern0pt}bdg{\isacharprime}{\kern0pt}{\isacharcomma}{\kern0pt}\ {\isacharunderscore}{\kern0pt}{\isacharparenright}{\kern0pt}{\isacharcomma}{\kern0pt}\ {\isacharunderscore}{\kern0pt}{\isacharparenright}{\kern0pt}\ {\isacharequal}{\kern0pt}{\isachargreater}{\kern0pt}\ Binding{\isachardot}{\kern0pt}name{\isacharunderscore}{\kern0pt}of\ bdg{\isacharprime}{\kern0pt}\ {\isacharequal}{\kern0pt}\ var{\isacharunderscore}{\kern0pt}name{\isacharparenright}{\kern0pt}\ locals{\isacharunderscore}{\kern0pt}decl\ of\isanewline
\ \ \ \ \ \ \ \ \ \ \ \ \ \ \ \ \ \ \ \ \ \ \ \ \ \ \ \ SOME\ {\isacharparenleft}{\kern0pt}{\isacharparenleft}{\kern0pt}{\isacharunderscore}{\kern0pt}{\isacharcomma}{\kern0pt}\ typ{\isacharparenright}{\kern0pt}{\isacharcomma}{\kern0pt}\ {\isacharunderscore}{\kern0pt}{\isacharparenright}{\kern0pt}\ {\isacharequal}{\kern0pt}{\isachargreater}{\kern0pt}\ SOME\ typ\isanewline
\ \ \ \ \ \ \ \ \ \ \ \ \ \ \ \ \ \ \ \ \ \ \ \ \ \ {\isacharbar}{\kern0pt}\ NONE\ {\isacharequal}{\kern0pt}{\isachargreater}{\kern0pt}\ NONE\isanewline
\ \ \ \ \ \ \ \ \ \ in\isanewline
\ \ \ \ \ \ \ \ \ \ \ \ case\ expected{\isacharunderscore}{\kern0pt}type{\isacharunderscore}{\kern0pt}opt\ of\isanewline
\ \ \ \ \ \ \ \ \ \ \ \ \ \ SOME\ expected{\isacharunderscore}{\kern0pt}type\ {\isacharequal}{\kern0pt}{\isachargreater}{\kern0pt}\isanewline
\ \ \ \ \ \ \ \ \ \ \ \ \ \ \ \ if\ term{\isacharunderscore}{\kern0pt}type\ {\isacharequal}{\kern0pt}\ expected{\isacharunderscore}{\kern0pt}type\ then\isanewline
\ \ \ \ \ \ \ \ \ \ \ \ \ \ \ \ \ \ errors\isanewline
\ \ \ \ \ \ \ \ \ \ \ \ \ \ \ \ else\isanewline
\ \ \ \ \ \ \ \ \ \ \ \ \ \ \ \ \ \ {\isacharparenleft}{\kern0pt}{\isachardoublequote}{\kern0pt}Type\ mismatch\ in\ assignment\ to\ {\isachardoublequote}{\kern0pt}\ {\isacharcircum}{\kern0pt}\ var{\isacharunderscore}{\kern0pt}name\ {\isacharcircum}{\kern0pt}\isanewline
\ \ \ \ \ \ \ \ \ \ \ \ \ \ \ \ \ \ \ {\isachardoublequote}{\kern0pt}{\isacharcolon}{\kern0pt}\ expected\ {\isachardoublequote}{\kern0pt}\ {\isacharcircum}{\kern0pt}\ Syntax{\isachardot}{\kern0pt}string{\isacharunderscore}{\kern0pt}of{\isacharunderscore}{\kern0pt}typ{\isacharunderscore}{\kern0pt}global\ thy\ expected{\isacharunderscore}{\kern0pt}type\ {\isacharcircum}{\kern0pt}\isanewline
\ \ \ \ \ \ \ \ \ \ \ \ \ \ \ \ \ \ \ {\isachardoublequote}{\kern0pt}{\isacharcomma}{\kern0pt}\ got\ {\isachardoublequote}{\kern0pt}\ {\isacharcircum}{\kern0pt}\ Syntax{\isachardot}{\kern0pt}string{\isacharunderscore}{\kern0pt}of{\isacharunderscore}{\kern0pt}typ{\isacharunderscore}{\kern0pt}global\ thy\ term{\isacharunderscore}{\kern0pt}type{\isacharparenright}{\kern0pt}\ {\isacharcolon}{\kern0pt}{\isacharcolon}{\kern0pt}\ errors\isanewline
\ \ \ \ \ \ \ \ \ \ \ \ {\isacharbar}{\kern0pt}\ NONE\ {\isacharequal}{\kern0pt}{\isachargreater}{\kern0pt}\isanewline
\ \ \ \ \ \ \ \ \ \ \ \ \ \ \ \ {\isacharparenleft}{\kern0pt}{\isachardoublequote}{\kern0pt}Variable\ {\isachardoublequote}{\kern0pt}\ {\isacharcircum}{\kern0pt}\ var{\isacharunderscore}{\kern0pt}name\ {\isacharcircum}{\kern0pt}\ {\isachardoublequote}{\kern0pt}\ is\ not\ locally\ declared{\isachardoublequote}{\kern0pt}{\isacharparenright}{\kern0pt}\ {\isacharcolon}{\kern0pt}{\isacharcolon}{\kern0pt}\ errors\isanewline
\ \ \ \ \ \ \ \ \ \ end\isanewline
\isanewline
\ \ \ \ \ \ {\isacharbar}{\kern0pt}\ SendA\ {\isacharparenleft}{\kern0pt}bdg{\isacharcomma}{\kern0pt}\ term{\isacharparenright}{\kern0pt}\ {\isacharequal}{\kern0pt}{\isachargreater}{\kern0pt}\isanewline
\ \ \ \ \ \ \ \ \ \ let\isanewline
\ \ \ \ \ \ \ \ \ \ \ \ val\ var{\isacharunderscore}{\kern0pt}name\ {\isacharequal}{\kern0pt}\ Binding{\isachardot}{\kern0pt}name{\isacharunderscore}{\kern0pt}of\ bdg\isanewline
\ \ \ \ \ \ \ \ \ \ \ \ val\ term{\isacharunderscore}{\kern0pt}type\ {\isacharequal}{\kern0pt}\ fastype{\isacharunderscore}{\kern0pt}of\ term\isanewline
\ \ \ \ \ \ \ \ \ \ \ \ val\ expected{\isacharunderscore}{\kern0pt}type{\isacharunderscore}{\kern0pt}opt\ {\isacharequal}{\kern0pt}\ \isanewline
\ \ \ \ \ \ \ \ \ \ \ \ \ \ \ \ \ \ \ \ \ \ \ \ \ \ case\ find{\isacharunderscore}{\kern0pt}first\ {\isacharparenleft}{\kern0pt}fn\ {\isacharparenleft}{\kern0pt}\ bdg{\isacharprime}{\kern0pt}{\isacharcomma}{\kern0pt}\ {\isacharunderscore}{\kern0pt}{\isacharcomma}{\kern0pt}\ {\isacharunderscore}{\kern0pt}{\isacharcomma}{\kern0pt}\ {\isacharunderscore}{\kern0pt}{\isacharparenright}{\kern0pt}\ {\isacharequal}{\kern0pt}{\isachargreater}{\kern0pt}\ Binding{\isachardot}{\kern0pt}name{\isacharunderscore}{\kern0pt}of\ bdg{\isacharprime}{\kern0pt}\ {\isacharequal}{\kern0pt}\ var{\isacharunderscore}{\kern0pt}name{\isacharparenright}{\kern0pt}\ vars{\isacharunderscore}{\kern0pt}decl\ of\isanewline
\ \ \ \ \ \ \ \ \ \ \ \ \ \ \ \ \ \ \ \ \ \ \ \ \ \ \ \ SOME\ {\isacharparenleft}{\kern0pt}\ {\isacharunderscore}{\kern0pt}{\isacharcomma}{\kern0pt}\ typ{\isacharcomma}{\kern0pt}\ {\isacharunderscore}{\kern0pt}{\isacharcomma}{\kern0pt}\ {\isacharunderscore}{\kern0pt}{\isacharparenright}{\kern0pt}\ {\isacharequal}{\kern0pt}{\isachargreater}{\kern0pt}\ SOME\ typ\isanewline
\ \ \ \ \ \ \ \ \ \ \ \ \ \ \ \ \ \ \ \ \ \ \ \ \ \ {\isacharbar}{\kern0pt}\ NONE\ {\isacharequal}{\kern0pt}{\isachargreater}{\kern0pt}\ NONE\isanewline
\ \ \ \ \ \ \ \ \ \ in\isanewline
\ \ \ \ \ \ \ \ \ \ \ \ case\ expected{\isacharunderscore}{\kern0pt}type{\isacharunderscore}{\kern0pt}opt\ of\isanewline
\ \ \ \ \ \ \ \ \ \ \ \ \ \ SOME\ expected{\isacharunderscore}{\kern0pt}type\ {\isacharequal}{\kern0pt}{\isachargreater}{\kern0pt}\isanewline
\ \ \ \ \ \ \ \ \ \ \ \ \ \ \ \ if\ term{\isacharunderscore}{\kern0pt}type\ {\isacharequal}{\kern0pt}\ expected{\isacharunderscore}{\kern0pt}type\ then\isanewline
\ \ \ \ \ \ \ \ \ \ \ \ \ \ \ \ \ \ errors\isanewline
\ \ \ \ \ \ \ \ \ \ \ \ \ \ \ \ else\isanewline
\ \ \ \ \ \ \ \ \ \ \ \ \ \ \ \ \ \ {\isacharparenleft}{\kern0pt}{\isachardoublequote}{\kern0pt}Type\ mismatch\ in\ assignment\ to\ {\isachardoublequote}{\kern0pt}\ {\isacharcircum}{\kern0pt}\ var{\isacharunderscore}{\kern0pt}name\ {\isacharcircum}{\kern0pt}\isanewline
\ \ \ \ \ \ \ \ \ \ \ \ \ \ \ \ \ \ \ {\isachardoublequote}{\kern0pt}{\isacharcolon}{\kern0pt}\ expected\ {\isachardoublequote}{\kern0pt}\ {\isacharcircum}{\kern0pt}\ Syntax{\isachardot}{\kern0pt}string{\isacharunderscore}{\kern0pt}of{\isacharunderscore}{\kern0pt}typ{\isacharunderscore}{\kern0pt}global\ thy\ expected{\isacharunderscore}{\kern0pt}type\ {\isacharcircum}{\kern0pt}\isanewline
\ \ \ \ \ \ \ \ \ \ \ \ \ \ \ \ \ \ \ {\isachardoublequote}{\kern0pt}{\isacharcomma}{\kern0pt}\ got\ {\isachardoublequote}{\kern0pt}\ {\isacharcircum}{\kern0pt}\ Syntax{\isachardot}{\kern0pt}string{\isacharunderscore}{\kern0pt}of{\isacharunderscore}{\kern0pt}typ{\isacharunderscore}{\kern0pt}global\ thy\ term{\isacharunderscore}{\kern0pt}type{\isacharparenright}{\kern0pt}\ {\isacharcolon}{\kern0pt}{\isacharcolon}{\kern0pt}\ errors\isanewline
\ \ \ \ \ \ \ \ \ \ \ \ {\isacharbar}{\kern0pt}\ NONE\ {\isacharequal}{\kern0pt}{\isachargreater}{\kern0pt}\isanewline
\ \ \ \ \ \ \ \ \ \ \ \ \ \ \ \ {\isacharparenleft}{\kern0pt}{\isachardoublequote}{\kern0pt}Variable\ {\isachardoublequote}{\kern0pt}\ {\isacharcircum}{\kern0pt}\ var{\isacharunderscore}{\kern0pt}name\ {\isacharcircum}{\kern0pt}\ {\isachardoublequote}{\kern0pt}\ is\ not\ globally\ declared{\isachardoublequote}{\kern0pt}{\isacharparenright}{\kern0pt}\ {\isacharcolon}{\kern0pt}{\isacharcolon}{\kern0pt}\ errors\isanewline
\ \ \ \ \ \ \ \ \ \ end\isanewline
\isanewline
\ \ \ \ \ \ {\isacharbar}{\kern0pt}\ ReceiveA\ {\isacharparenleft}{\kern0pt}local{\isacharunderscore}{\kern0pt}bdg{\isacharcomma}{\kern0pt}\ shared{\isacharunderscore}{\kern0pt}bdg{\isacharparenright}{\kern0pt}\ {\isacharequal}{\kern0pt}{\isachargreater}{\kern0pt}\isanewline
\ \ \ \ \ \ \ \ \ \ let\isanewline
\ \ \ \ \ \ \ \ \ \ \ \ val\ local{\isacharunderscore}{\kern0pt}var\ {\isacharequal}{\kern0pt}\ Binding{\isachardot}{\kern0pt}name{\isacharunderscore}{\kern0pt}of\ local{\isacharunderscore}{\kern0pt}bdg\isanewline
\ \ \ \ \ \ \ \ \ \ \ \ val\ shared{\isacharunderscore}{\kern0pt}var\ {\isacharequal}{\kern0pt}\ Binding{\isachardot}{\kern0pt}name{\isacharunderscore}{\kern0pt}of\ shared{\isacharunderscore}{\kern0pt}bdg\isanewline
\ \ \ \ \ \ \ \ \ \ \ \ val\ local{\isacharunderscore}{\kern0pt}type{\isacharunderscore}{\kern0pt}opt\ {\isacharequal}{\kern0pt}\isanewline
\ \ \ \ \ \ \ \ \ \ \ \ \ \ \ \ \ \ \ \ \ \ \ \ \ \ case\ find{\isacharunderscore}{\kern0pt}first\ {\isacharparenleft}{\kern0pt}fn\ {\isacharparenleft}{\kern0pt}{\isacharparenleft}{\kern0pt}bdg{\isacharprime}{\kern0pt}{\isacharcomma}{\kern0pt}\ {\isacharunderscore}{\kern0pt}{\isacharparenright}{\kern0pt}{\isacharcomma}{\kern0pt}\ {\isacharunderscore}{\kern0pt}{\isacharparenright}{\kern0pt}\ {\isacharequal}{\kern0pt}{\isachargreater}{\kern0pt}\ Binding{\isachardot}{\kern0pt}name{\isacharunderscore}{\kern0pt}of\ bdg{\isacharprime}{\kern0pt}\ {\isacharequal}{\kern0pt}\ local{\isacharunderscore}{\kern0pt}var{\isacharparenright}{\kern0pt}\ locals{\isacharunderscore}{\kern0pt}decl\ of\isanewline
\ \ \ \ \ \ \ \ \ \ \ \ \ \ \ \ \ \ \ \ \ \ \ \ \ \ \ \ SOME\ {\isacharparenleft}{\kern0pt}{\isacharparenleft}{\kern0pt}{\isacharunderscore}{\kern0pt}{\isacharcomma}{\kern0pt}\ typ{\isacharparenright}{\kern0pt}{\isacharcomma}{\kern0pt}\ {\isacharunderscore}{\kern0pt}{\isacharparenright}{\kern0pt}\ {\isacharequal}{\kern0pt}{\isachargreater}{\kern0pt}\ SOME\ typ\isanewline
\ \ \ \ \ \ \ \ \ \ \ \ \ \ \ \ \ \ \ \ \ \ \ \ \ \ {\isacharbar}{\kern0pt}\ NONE\ {\isacharequal}{\kern0pt}{\isachargreater}{\kern0pt}\ NONE\isanewline
\isanewline
\ \ \ \ \ \ \ \ \ \ \ \ val\ shared{\isacharunderscore}{\kern0pt}type{\isacharunderscore}{\kern0pt}opt\ {\isacharequal}{\kern0pt}\isanewline
\ \ \ \ \ \ \ \ \ \ \ \ \ \ {\isacharparenleft}{\kern0pt}{\isacharasterisk}{\kern0pt}\ DOUBLE\ CHECK\ IF\ THIS\ CAN\ SAFELY\ BE\ REMOVED\isanewline
\isanewline
\ \ \ \ \ \ \ \ \ \ \ \ \ \ case\ Symtab{\isachardot}{\kern0pt}lookup\ varstab\ shared{\isacharunderscore}{\kern0pt}var\ of\isanewline
\ \ \ \ \ \ \ \ \ \ \ \ \ \ \ \ SOME\ {\isacharparenleft}{\kern0pt}SOME\ t{\isacharparenright}{\kern0pt}\ {\isacharequal}{\kern0pt}{\isachargreater}{\kern0pt}\ fastype{\isacharunderscore}{\kern0pt}of\ t\isanewline
\ \ \ \ \ \ \ \ \ \ \ \ \ \ {\isacharbar}{\kern0pt}\ NONE\ {\isacharequal}{\kern0pt}{\isachargreater}{\kern0pt}\isanewline
\ \ \ \ \ \ \ \ \ \ \ \ \ \ \ \ case\ Symtab{\isachardot}{\kern0pt}lookup\ {\isacharparenleft}{\kern0pt}{\isacharbang}{\kern0pt}SPY\ {\isacharbar}{\kern0pt}{\isachargreater}{\kern0pt}\ Option{\isachardot}{\kern0pt}valOf\ {\isacharbar}{\kern0pt}{\isachargreater}{\kern0pt}\ {\isacharhash}{\kern0pt}varstab{\isacharparenright}{\kern0pt}\ shared{\isacharunderscore}{\kern0pt}var\ of\isanewline
\ \ \ \ \ \ \ \ \ \ \ \ \ \ \ \ \ \ SOME\ {\isacharparenleft}{\kern0pt}SOME\ t{\isacharparenright}{\kern0pt}\ {\isacharequal}{\kern0pt}{\isachargreater}{\kern0pt}\ fastype{\isacharunderscore}{\kern0pt}of\ t\isanewline
\ \ \ \ \ \ \ \ \ \ \ \ \ \ \ \ {\isacharbar}{\kern0pt}\ {\isacharunderscore}{\kern0pt}\ {\isacharequal}{\kern0pt}{\isachargreater}{\kern0pt}\ error\ {\isacharparenleft}{\kern0pt}{\isachardoublequote}{\kern0pt}Shared\ variable\ {\isachardoublequote}{\kern0pt}\ {\isacharcircum}{\kern0pt}\ shared{\isacharunderscore}{\kern0pt}var\ {\isacharcircum}{\kern0pt}\ {\isachardoublequote}{\kern0pt}\ not\ declared{\isachardoublequote}{\kern0pt}{\isacharparenright}{\kern0pt}\isanewline
\ \ \ \ \ \ \ \ \ \ \ \ \ \ {\isacharbar}{\kern0pt}\ SOME\ NONE\ {\isacharequal}{\kern0pt}{\isachargreater}{\kern0pt}\ error\ {\isacharparenleft}{\kern0pt}{\isachardoublequote}{\kern0pt}Variable\ {\isachardoublequote}{\kern0pt}\ {\isacharcircum}{\kern0pt}\ shared{\isacharunderscore}{\kern0pt}var\ {\isacharcircum}{\kern0pt}\ {\isachardoublequote}{\kern0pt}\ declared\ but\ has\ no\ initial\ value{\isachardoublequote}{\kern0pt}{\isacharparenright}{\kern0pt}{\isacharasterisk}{\kern0pt}{\isacharparenright}{\kern0pt}\isanewline
\ \ \ \ \ \ \ \ \ \ \ \ \ \ \ \ case\ find{\isacharunderscore}{\kern0pt}first\ {\isacharparenleft}{\kern0pt}fn\ {\isacharparenleft}{\kern0pt}\ bdg{\isacharprime}{\kern0pt}{\isacharcomma}{\kern0pt}\ {\isacharunderscore}{\kern0pt}{\isacharcomma}{\kern0pt}\ {\isacharunderscore}{\kern0pt}{\isacharcomma}{\kern0pt}\ {\isacharunderscore}{\kern0pt}{\isacharparenright}{\kern0pt}\ {\isacharequal}{\kern0pt}{\isachargreater}{\kern0pt}\ Binding{\isachardot}{\kern0pt}name{\isacharunderscore}{\kern0pt}of\ bdg{\isacharprime}{\kern0pt}\ {\isacharequal}{\kern0pt}\ shared{\isacharunderscore}{\kern0pt}var{\isacharparenright}{\kern0pt}\ vars{\isacharunderscore}{\kern0pt}decl\ of\isanewline
\ \ \ \ \ \ \ \ \ \ \ \ \ \ \ \ \ \ \ \ \ \ \ \ \ \ \ \ \ \ \ \ \ \ \ \ \ \ \ \ \ \ \ \ SOME\ {\isacharparenleft}{\kern0pt}\ {\isacharunderscore}{\kern0pt}{\isacharcomma}{\kern0pt}\ typ{\isacharcomma}{\kern0pt}\ {\isacharunderscore}{\kern0pt}{\isacharcomma}{\kern0pt}\ {\isacharunderscore}{\kern0pt}{\isacharparenright}{\kern0pt}\ {\isacharequal}{\kern0pt}{\isachargreater}{\kern0pt}\ SOME\ typ\isanewline
\ \ \ \ \ \ \ \ \ \ \ \ \ \ \ \ \ \ \ \ \ \ \ \ \ \ \ \ \ \ \ \ \ \ \ \ \ \ \ \ \ \ {\isacharbar}{\kern0pt}\ NONE\ {\isacharequal}{\kern0pt}{\isachargreater}{\kern0pt}\ NONE\isanewline
\isanewline
\ \ \ \ \ \ \ \ \ \ \ in\isanewline
\ \ \ \ \ \ \ \ \ \ \ \ case\ {\isacharparenleft}{\kern0pt}local{\isacharunderscore}{\kern0pt}type{\isacharunderscore}{\kern0pt}opt{\isacharcomma}{\kern0pt}\ shared{\isacharunderscore}{\kern0pt}type{\isacharunderscore}{\kern0pt}opt{\isacharparenright}{\kern0pt}\ of\isanewline
\ \ \ \ \ \ \ \ \ \ \ \ \ \ {\isacharparenleft}{\kern0pt}SOME\ local{\isacharunderscore}{\kern0pt}type{\isacharcomma}{\kern0pt}\ SOME\ shared{\isacharunderscore}{\kern0pt}type{\isacharparenright}{\kern0pt}\ {\isacharequal}{\kern0pt}{\isachargreater}{\kern0pt}\isanewline
\ \ \ \ \ \ \ \ \ \ \ \ \ \ \ \ if\ local{\isacharunderscore}{\kern0pt}type\ {\isacharequal}{\kern0pt}\ shared{\isacharunderscore}{\kern0pt}type\ then\ \isanewline
\ \ \ \ \ \ \ \ \ \ \ \ \ \ \ \ \ \ errors\isanewline
\ \ \ \ \ \ \ \ \ \ \ \ \ \ \ \ else\isanewline
\ \ \ \ \ \ \ \ \ \ \ \ \ \ \ \ \ \ {\isacharparenleft}{\kern0pt}{\isachardoublequote}{\kern0pt}Type\ mismatch\ in\ receive{\isacharcolon}{\kern0pt}\ cannot\ assign\ {\isachardoublequote}{\kern0pt}\ {\isacharcircum}{\kern0pt}\ shared{\isacharunderscore}{\kern0pt}var\ {\isacharcircum}{\kern0pt}\ {\isachardoublequote}{\kern0pt}\ to\ {\isachardoublequote}{\kern0pt}\ {\isacharcircum}{\kern0pt}\ local{\isacharunderscore}{\kern0pt}var\ {\isacharcircum}{\kern0pt}\isanewline
\ \ \ \ \ \ \ \ \ \ \ \ \ \ \ \ \ \ \ {\isachardoublequote}{\kern0pt}\ {\isacharparenleft}{\kern0pt}expected\ {\isachardoublequote}{\kern0pt}\ {\isacharcircum}{\kern0pt}\ Syntax{\isachardot}{\kern0pt}string{\isacharunderscore}{\kern0pt}of{\isacharunderscore}{\kern0pt}typ{\isacharunderscore}{\kern0pt}global\ thy\ local{\isacharunderscore}{\kern0pt}type\ {\isacharcircum}{\kern0pt}\isanewline
\ \ \ \ \ \ \ \ \ \ \ \ \ \ \ \ \ \ \ {\isachardoublequote}{\kern0pt}{\isacharcomma}{\kern0pt}\ got\ {\isachardoublequote}{\kern0pt}\ {\isacharcircum}{\kern0pt}\ Syntax{\isachardot}{\kern0pt}string{\isacharunderscore}{\kern0pt}of{\isacharunderscore}{\kern0pt}typ{\isacharunderscore}{\kern0pt}global\ thy\ shared{\isacharunderscore}{\kern0pt}type\ {\isacharcircum}{\kern0pt}\ {\isachardoublequote}{\kern0pt}{\isacharparenright}{\kern0pt}{\isachardoublequote}{\kern0pt}{\isacharparenright}{\kern0pt}\ {\isacharcolon}{\kern0pt}{\isacharcolon}{\kern0pt}\ errors\isanewline
\ \ \ \ \ \ \ \ \ \ \ \ {\isacharbar}{\kern0pt}\ {\isacharparenleft}{\kern0pt}NONE{\isacharcomma}{\kern0pt}\ {\isacharunderscore}{\kern0pt}{\isacharparenright}{\kern0pt}\ {\isacharequal}{\kern0pt}{\isachargreater}{\kern0pt}\isanewline
\ \ \ \ \ \ \ \ \ \ \ \ \ \ \ \ {\isacharparenleft}{\kern0pt}{\isachardoublequote}{\kern0pt}Local\ variable\ {\isachardoublequote}{\kern0pt}\ {\isacharcircum}{\kern0pt}\ local{\isacharunderscore}{\kern0pt}var\ {\isacharcircum}{\kern0pt}\ {\isachardoublequote}{\kern0pt}\ not\ declared{\isachardoublequote}{\kern0pt}{\isacharparenright}{\kern0pt}\ {\isacharcolon}{\kern0pt}{\isacharcolon}{\kern0pt}\ errors\isanewline
\ \ \ \ \ \ \ \ \ \ \ \ {\isacharbar}{\kern0pt}\ {\isacharparenleft}{\kern0pt}{\isacharunderscore}{\kern0pt}{\isacharcomma}{\kern0pt}\ NONE{\isacharparenright}{\kern0pt}\ {\isacharequal}{\kern0pt}{\isachargreater}{\kern0pt}\isanewline
\ \ \ \ \ \ \ \ \ \ \ \ \ \ \ \ {\isacharparenleft}{\kern0pt}{\isachardoublequote}{\kern0pt}Shared\ variable\ {\isachardoublequote}{\kern0pt}\ {\isacharcircum}{\kern0pt}\ shared{\isacharunderscore}{\kern0pt}var\ {\isacharcircum}{\kern0pt}\ {\isachardoublequote}{\kern0pt}\ not\ declared{\isachardoublequote}{\kern0pt}{\isacharparenright}{\kern0pt}\ {\isacharcolon}{\kern0pt}{\isacharcolon}{\kern0pt}\ errors\isanewline
\ \ \ \ \ \ \ \ \ \ end\isanewline
\ \ \ \ \ \ {\isacharbar}{\kern0pt}\ IfelseA\ {\isacharparenleft}{\kern0pt}{\isacharparenleft}{\kern0pt}cond{\isacharcomma}{\kern0pt}\ then{\isacharunderscore}{\kern0pt}branch{\isacharparenright}{\kern0pt}{\isacharcomma}{\kern0pt}\ else{\isacharunderscore}{\kern0pt}branch{\isacharparenright}{\kern0pt}\ {\isacharequal}{\kern0pt}{\isachargreater}{\kern0pt}\isanewline
\ \ \ \ \ \ \ \ let\isanewline
\ \ \ \ \ \ \ \ \ \ val\ cond{\isacharunderscore}{\kern0pt}type\ {\isacharequal}{\kern0pt}\isanewline
\ \ \ \ \ \ \ \ \ \ \ \ let\isanewline
\ \ \ \ \ \ \ \ \ \ \ \ \ \ val\ t\ {\isacharequal}{\kern0pt}\ cond\isanewline
\ \ \ \ \ \ \ \ \ \ \ \ in\isanewline
\ \ \ \ \ \ \ \ \ \ \ \ \ \ case\ cond\ of\isanewline
\ \ \ \ \ \ \ \ \ \ \ \ \ \ \ \ Free\ {\isacharparenleft}{\kern0pt}x{\isacharcomma}{\kern0pt}\ {\isacharunderscore}{\kern0pt}{\isacharparenright}{\kern0pt}\ {\isacharequal}{\kern0pt}{\isachargreater}{\kern0pt}\isanewline
\ \ \ \ \ \ \ \ \ \ \ \ \ \ \ \ \ \ {\isacharparenleft}{\kern0pt}case\ Symtab{\isachardot}{\kern0pt}lookup\ locstab\ x\ of\isanewline
\ \ \ \ \ \ \ \ \ \ \ \ \ \ \ \ \ \ \ \ \ SOME\ {\isacharparenleft}{\kern0pt}SOME\ t{\isacharprime}{\kern0pt}{\isacharparenright}{\kern0pt}\ {\isacharequal}{\kern0pt}{\isachargreater}{\kern0pt}\ fastype{\isacharunderscore}{\kern0pt}of\ t{\isacharprime}{\kern0pt}\isanewline
\ \ \ \ \ \ \ \ \ \ \ \ \ \ \ \ \ \ \ {\isacharbar}{\kern0pt}\ SOME\ NONE\ {\isacharequal}{\kern0pt}{\isachargreater}{\kern0pt}\ error\ {\isacharparenleft}{\kern0pt}{\isachardoublequote}{\kern0pt}Variable\ {\isachardoublequote}{\kern0pt}\ {\isacharcircum}{\kern0pt}\ x\ {\isacharcircum}{\kern0pt}\ {\isachardoublequote}{\kern0pt}\ declared\ but\ has\ no\ initial\ value{\isachardoublequote}{\kern0pt}{\isacharparenright}{\kern0pt}\isanewline
\ \ \ \ \ \ \ \ \ \ \ \ \ \ \ \ \ \ \ {\isacharbar}{\kern0pt}\ NONE\ {\isacharequal}{\kern0pt}{\isachargreater}{\kern0pt}\isanewline
\ \ \ \ \ \ \ \ \ \ \ \ \ \ \ \ \ \ \ \ \ \ \ {\isacharparenleft}{\kern0pt}case\ Symtab{\isachardot}{\kern0pt}lookup\ {\isacharparenleft}{\kern0pt}{\isacharbang}{\kern0pt}SPY\ {\isacharbar}{\kern0pt}{\isachargreater}{\kern0pt}\ Option{\isachardot}{\kern0pt}valOf\ {\isacharbar}{\kern0pt}{\isachargreater}{\kern0pt}\ {\isacharhash}{\kern0pt}varstab{\isacharparenright}{\kern0pt}\ x\ of\isanewline
\ \ \ \ \ \ \ \ \ \ \ \ \ \ \ \ \ \ \ \ \ \ \ \ \ \ SOME\ {\isacharparenleft}{\kern0pt}SOME\ t{\isacharprime}{\kern0pt}{\isacharparenright}{\kern0pt}\ {\isacharequal}{\kern0pt}{\isachargreater}{\kern0pt}\ fastype{\isacharunderscore}{\kern0pt}of\ t{\isacharprime}{\kern0pt}\isanewline
\ \ \ \ \ \ \ \ \ \ \ \ \ \ \ \ \ \ \ \ \ \ \ \ {\isacharbar}{\kern0pt}\ SOME\ NONE\ {\isacharequal}{\kern0pt}{\isachargreater}{\kern0pt}\ error\ {\isacharparenleft}{\kern0pt}{\isachardoublequote}{\kern0pt}Global\ variable\ {\isachardoublequote}{\kern0pt}\ {\isacharcircum}{\kern0pt}\ x\ {\isacharcircum}{\kern0pt}\ {\isachardoublequote}{\kern0pt}\ declared\ but\ has\ no\ initial\ value{\isachardoublequote}{\kern0pt}{\isacharparenright}{\kern0pt}\isanewline
\ \ \ \ \ \ \ \ \ \ \ \ \ \ \ \ \ \ \ \ \ \ \ \ {\isacharbar}{\kern0pt}\ NONE\ {\isacharequal}{\kern0pt}{\isachargreater}{\kern0pt}\ fastype{\isacharunderscore}{\kern0pt}of\ t{\isacharparenright}{\kern0pt}{\isacharparenright}{\kern0pt}\ \ {\isacharparenleft}{\kern0pt}{\isacharasterisk}{\kern0pt}\ fallback\ to\ actual\ term\ type\ {\isacharasterisk}{\kern0pt}{\isacharparenright}{\kern0pt}\isanewline
\ \ \ \ \ \ \ \ \ \ \ \ \ \ {\isacharbar}{\kern0pt}\ {\isacharunderscore}{\kern0pt}\ {\isacharequal}{\kern0pt}{\isachargreater}{\kern0pt}\ fastype{\isacharunderscore}{\kern0pt}of\ t\isanewline
\ \ \ \ \ \ \ \ \ \ \ \ end\isanewline
\ \ \ \ \isanewline
\ \ \ \ \ \ \ \ \ \ val\ cond{\isacharunderscore}{\kern0pt}error\ {\isacharequal}{\kern0pt}\isanewline
\ \ \ \ \ \ \ \ \ \ \ \ if\ cond{\isacharunderscore}{\kern0pt}type\ {\isacharequal}{\kern0pt}\ {\isacharat}{\kern0pt}{\isacharbraceleft}{\kern0pt}typ\ bool{\isacharbraceright}{\kern0pt}\ then\ {\isacharbrackleft}{\kern0pt}{\isacharbrackright}{\kern0pt}\isanewline
\ \ \ \ \ \ \ \ \ \ \ \ else\ {\isacharbrackleft}{\kern0pt}{\isachardoublequote}{\kern0pt}Condition\ must\ be\ boolean{\isacharcomma}{\kern0pt}\ got{\isacharcolon}{\kern0pt}\ {\isachardoublequote}{\kern0pt}\ {\isacharcircum}{\kern0pt}\ Syntax{\isachardot}{\kern0pt}string{\isacharunderscore}{\kern0pt}of{\isacharunderscore}{\kern0pt}typ{\isacharunderscore}{\kern0pt}global\ thy\ cond{\isacharunderscore}{\kern0pt}type{\isacharbrackright}{\kern0pt}\isanewline
\ \ \ \ \isanewline
\ \ \ \ \ \ \ \ \ \ val\ then{\isacharunderscore}{\kern0pt}errors\ {\isacharequal}{\kern0pt}\ List{\isachardot}{\kern0pt}foldl\ check{\isacharunderscore}{\kern0pt}action{\isacharunderscore}{\kern0pt}types\ {\isacharbrackleft}{\kern0pt}{\isacharbrackright}{\kern0pt}\ then{\isacharunderscore}{\kern0pt}branch\isanewline
\ \ \ \ \ \ \ \ \ \ val\ else{\isacharunderscore}{\kern0pt}errors\ {\isacharequal}{\kern0pt}\ List{\isachardot}{\kern0pt}foldl\ check{\isacharunderscore}{\kern0pt}action{\isacharunderscore}{\kern0pt}types\ {\isacharbrackleft}{\kern0pt}{\isacharbrackright}{\kern0pt}\ else{\isacharunderscore}{\kern0pt}branch\isanewline
\ \ \ \ \ \ \ \ in\isanewline
\ \ \ \ \ \ \ \ \ \ cond{\isacharunderscore}{\kern0pt}error\ {\isacharat}{\kern0pt}\ then{\isacharunderscore}{\kern0pt}errors\ {\isacharat}{\kern0pt}\ else{\isacharunderscore}{\kern0pt}errors\ {\isacharat}{\kern0pt}\ errors\isanewline
\ \ \ \ \ \ \ \ end\isanewline
\isanewline
\ \ \ \ \ {\isacharbar}{\kern0pt}\ WhileA\ {\isacharparenleft}{\kern0pt}cond{\isacharcomma}{\kern0pt}\ body{\isacharparenright}{\kern0pt}\ {\isacharequal}{\kern0pt}{\isachargreater}{\kern0pt}\isanewline
\ \ \ \ \ \ \ \ \ \ \ \ let\isanewline
\ \ \ \ \ \ \ \ \ \ \ \ \ \ val\ cond{\isacharunderscore}{\kern0pt}type\ {\isacharequal}{\kern0pt}\isanewline
\ \ \ \ \ \ \ \ \ \ \ \ \ \ \ \ let\isanewline
\ \ \ \ \ \ \ \ \ \ \ \ \ \ \ \ \ \ val\ t\ {\isacharequal}{\kern0pt}\ cond\isanewline
\ \ \ \ \ \ \ \ \ \ \ \ \ \ \ \ in\isanewline
\ \ \ \ \ \ \ \ \ \ \ \ \ \ \ \ \ \ case\ cond\ of\isanewline
\ \ \ \ \ \ \ \ \ \ \ \ \ \ \ \ \ \ \ \ Free\ {\isacharparenleft}{\kern0pt}x{\isacharcomma}{\kern0pt}\ {\isacharunderscore}{\kern0pt}{\isacharparenright}{\kern0pt}\ {\isacharequal}{\kern0pt}{\isachargreater}{\kern0pt}\isanewline
\ \ \ \ \ \ \ \ \ \ \ \ \ \ \ \ \ \ \ \ \ \ {\isacharparenleft}{\kern0pt}case\ Symtab{\isachardot}{\kern0pt}lookup\ locstab\ x\ of\isanewline
\ \ \ \ \ \ \ \ \ \ \ \ \ \ \ \ \ \ \ \ \ \ \ \ \ SOME\ {\isacharparenleft}{\kern0pt}SOME\ t{\isacharprime}{\kern0pt}{\isacharparenright}{\kern0pt}\ {\isacharequal}{\kern0pt}{\isachargreater}{\kern0pt}\ fastype{\isacharunderscore}{\kern0pt}of\ t{\isacharprime}{\kern0pt}\isanewline
\ \ \ \ \ \ \ \ \ \ \ \ \ \ \ \ \ \ \ \ \ \ \ {\isacharbar}{\kern0pt}\ SOME\ NONE\ {\isacharequal}{\kern0pt}{\isachargreater}{\kern0pt}\ error\ {\isacharparenleft}{\kern0pt}{\isachardoublequote}{\kern0pt}Variable\ {\isachardoublequote}{\kern0pt}\ {\isacharcircum}{\kern0pt}\ x\ {\isacharcircum}{\kern0pt}\ {\isachardoublequote}{\kern0pt}\ declared\ but\ has\ no\ initial\ value{\isachardoublequote}{\kern0pt}{\isacharparenright}{\kern0pt}\isanewline
\ \ \ \ \ \ \ \ \ \ \ \ \ \ \ \ \ \ \ \ \ \ \ {\isacharbar}{\kern0pt}\ NONE\ {\isacharequal}{\kern0pt}{\isachargreater}{\kern0pt}\isanewline
\ \ \ \ \ \ \ \ \ \ \ \ \ \ \ \ \ \ \ \ \ \ \ \ \ \ \ {\isacharparenleft}{\kern0pt}case\ Symtab{\isachardot}{\kern0pt}lookup\ {\isacharparenleft}{\kern0pt}{\isacharbang}{\kern0pt}SPY\ {\isacharbar}{\kern0pt}{\isachargreater}{\kern0pt}\ Option{\isachardot}{\kern0pt}valOf\ {\isacharbar}{\kern0pt}{\isachargreater}{\kern0pt}\ {\isacharhash}{\kern0pt}varstab{\isacharparenright}{\kern0pt}\ x\ of\isanewline
\ \ \ \ \ \ \ \ \ \ \ \ \ \ \ \ \ \ \ \ \ \ \ \ \ \ \ \ \ \ SOME\ {\isacharparenleft}{\kern0pt}SOME\ t{\isacharprime}{\kern0pt}{\isacharparenright}{\kern0pt}\ {\isacharequal}{\kern0pt}{\isachargreater}{\kern0pt}\ fastype{\isacharunderscore}{\kern0pt}of\ t{\isacharprime}{\kern0pt}\isanewline
\ \ \ \ \ \ \ \ \ \ \ \ \ \ \ \ \ \ \ \ \ \ \ \ \ \ \ \ {\isacharbar}{\kern0pt}\ SOME\ NONE\ {\isacharequal}{\kern0pt}{\isachargreater}{\kern0pt}\ error\ {\isacharparenleft}{\kern0pt}{\isachardoublequote}{\kern0pt}Global\ variable\ {\isachardoublequote}{\kern0pt}\ {\isacharcircum}{\kern0pt}\ x\ {\isacharcircum}{\kern0pt}\ {\isachardoublequote}{\kern0pt}\ declared\ but\ has\ no\ initial\ value{\isachardoublequote}{\kern0pt}{\isacharparenright}{\kern0pt}\isanewline
\ \ \ \ \ \ \ \ \ \ \ \ \ \ \ \ \ \ \ \ \ \ \ \ \ \ \ \ {\isacharbar}{\kern0pt}\ NONE\ {\isacharequal}{\kern0pt}{\isachargreater}{\kern0pt}\ fastype{\isacharunderscore}{\kern0pt}of\ t{\isacharparenright}{\kern0pt}{\isacharparenright}{\kern0pt}\ \ {\isacharparenleft}{\kern0pt}{\isacharasterisk}{\kern0pt}\ fallback\ {\isacharasterisk}{\kern0pt}{\isacharparenright}{\kern0pt}\isanewline
\ \ \ \ \ \ \ \ \ \ \ \ \ \ \ \ \ \ {\isacharbar}{\kern0pt}\ {\isacharunderscore}{\kern0pt}\ {\isacharequal}{\kern0pt}{\isachargreater}{\kern0pt}\ fastype{\isacharunderscore}{\kern0pt}of\ t\isanewline
\ \ \ \ \ \ \ \ \ \ \ \ \ \ \ \ end\isanewline
\ \ \ \ \ \ \ \ \isanewline
\ \ \ \ \ \ \ \ \ \ \ \ \ \ val\ cond{\isacharunderscore}{\kern0pt}error\ {\isacharequal}{\kern0pt}\isanewline
\ \ \ \ \ \ \ \ \ \ \ \ \ \ \ \ if\ cond{\isacharunderscore}{\kern0pt}type\ {\isacharequal}{\kern0pt}\ {\isacharat}{\kern0pt}{\isacharbraceleft}{\kern0pt}typ\ bool{\isacharbraceright}{\kern0pt}\ then\ {\isacharbrackleft}{\kern0pt}{\isacharbrackright}{\kern0pt}\isanewline
\ \ \ \ \ \ \ \ \ \ \ \ \ \ \ \ else\ {\isacharbrackleft}{\kern0pt}{\isachardoublequote}{\kern0pt}While\ condition\ must\ be\ boolean{\isacharcomma}{\kern0pt}\ got{\isacharcolon}{\kern0pt}\ {\isachardoublequote}{\kern0pt}\ {\isacharcircum}{\kern0pt}\ Syntax{\isachardot}{\kern0pt}string{\isacharunderscore}{\kern0pt}of{\isacharunderscore}{\kern0pt}typ{\isacharunderscore}{\kern0pt}global\ thy\ cond{\isacharunderscore}{\kern0pt}type{\isacharbrackright}{\kern0pt}\isanewline
\ \ \ \ \ \ \ \ \isanewline
\ \ \ \ \ \ \ \ \ \ \ \ \ \ val\ body{\isacharunderscore}{\kern0pt}errors\ {\isacharequal}{\kern0pt}\ List{\isachardot}{\kern0pt}foldl\ check{\isacharunderscore}{\kern0pt}action{\isacharunderscore}{\kern0pt}types\ {\isacharbrackleft}{\kern0pt}{\isacharbrackright}{\kern0pt}\ body\isanewline
\ \ \ \ \ \ \ \ \ \ \ \ in\isanewline
\ \ \ \ \ \ \ \ \ \ \ \ \ \ cond{\isacharunderscore}{\kern0pt}error\ {\isacharat}{\kern0pt}\ body{\isacharunderscore}{\kern0pt}errors\ {\isacharat}{\kern0pt}\ errors\isanewline
\ \ \ \ \ \ \ \ \ \ \ \ end\isanewline
\isanewline
\ \ \ \ \ \ {\isacharbar}{\kern0pt}\ {\isacharunderscore}{\kern0pt}\ {\isacharequal}{\kern0pt}{\isachargreater}{\kern0pt}\ errors\isanewline
\isanewline
\ \ in\isanewline
\ \ \ \ List{\isachardot}{\kern0pt}foldl\ check{\isacharunderscore}{\kern0pt}action{\isacharunderscore}{\kern0pt}types\ {\isacharbrackleft}{\kern0pt}{\isacharbrackright}{\kern0pt}\ actions\isanewline
\ \ end\isanewline
\isanewline
fun\ build{\isacharunderscore}{\kern0pt}initial{\isacharunderscore}{\kern0pt}state{\isacharunderscore}{\kern0pt}for{\isacharunderscore}{\kern0pt}thread\ \isanewline
\ \ \ \ \ \ {\isacharparenleft}{\kern0pt}globals\ {\isacharcolon}{\kern0pt}\ {\isacharparenleft}{\kern0pt}binding\ {\isacharasterisk}{\kern0pt}\ typ\ {\isacharasterisk}{\kern0pt}\ Position{\isachardot}{\kern0pt}T\ {\isacharasterisk}{\kern0pt}\ term\ option{\isacharparenright}{\kern0pt}\ list{\isacharparenright}{\kern0pt}\ \isanewline
\ \ \ \ \ \ {\isacharparenleft}{\kern0pt}locals\ {\isacharcolon}{\kern0pt}\ {\isacharparenleft}{\kern0pt}{\isacharparenleft}{\kern0pt}binding\ {\isacharasterisk}{\kern0pt}\ typ{\isacharparenright}{\kern0pt}\ {\isacharasterisk}{\kern0pt}\ term\ option{\isacharparenright}{\kern0pt}\ list{\isacharparenright}{\kern0pt}\ {\isacharparenleft}{\kern0pt}thy{\isacharcolon}{\kern0pt}\ theory{\isacharparenright}{\kern0pt}\ {\isacharcolon}{\kern0pt}\ state\ {\isacharequal}{\kern0pt}\isanewline
\ \ let\isanewline
\ \ \ \ fun\ term{\isacharunderscore}{\kern0pt}to{\isacharunderscore}{\kern0pt}int\ t{\isacharunderscore}{\kern0pt}opt\ {\isacharequal}{\kern0pt}\isanewline
\ \ \ \ \ \ \ \ case\ t{\isacharunderscore}{\kern0pt}opt\ of\isanewline
\ \ \ \ \ \ \ \ \ \ \ \ NONE\ {\isacharequal}{\kern0pt}{\isachargreater}{\kern0pt}\ {\isadigit{0}}\isanewline
\ \ \ \ \ \ \ \ \ \ {\isacharbar}{\kern0pt}\ SOME\ t\ {\isacharequal}{\kern0pt}{\isachargreater}{\kern0pt}\ eval{\isacharunderscore}{\kern0pt}term{\isacharunderscore}{\kern0pt}to{\isacharunderscore}{\kern0pt}int\ thy\ t\ {\isacharparenleft}{\kern0pt}fn\ {\isacharunderscore}{\kern0pt}\ {\isacharequal}{\kern0pt}{\isachargreater}{\kern0pt}\ {\isadigit{0}}{\isacharparenright}{\kern0pt}\isanewline
\isanewline
\ \ \ \ val\ global{\isacharunderscore}{\kern0pt}vals\ {\isacharequal}{\kern0pt}\ map\ {\isacharparenleft}{\kern0pt}fn\ {\isacharparenleft}{\kern0pt}bdg{\isacharcomma}{\kern0pt}\ {\isacharunderscore}{\kern0pt}{\isacharcomma}{\kern0pt}\ {\isacharunderscore}{\kern0pt}{\isacharcomma}{\kern0pt}\ t{\isacharunderscore}{\kern0pt}opt{\isacharparenright}{\kern0pt}\ {\isacharequal}{\kern0pt}{\isachargreater}{\kern0pt}\isanewline
\ \ \ \ \ \ \ \ \ \ \ \ \ \ \ \ \ \ \ \ \ \ \ \ \ \ \ \ {\isacharparenleft}{\kern0pt}Binding{\isachardot}{\kern0pt}name{\isacharunderscore}{\kern0pt}of\ bdg{\isacharcomma}{\kern0pt}\ term{\isacharunderscore}{\kern0pt}to{\isacharunderscore}{\kern0pt}int\ t{\isacharunderscore}{\kern0pt}opt{\isacharparenright}{\kern0pt}{\isacharparenright}{\kern0pt}\ globals\isanewline
\isanewline
\ \ \ \ val\ local{\isacharunderscore}{\kern0pt}vals\ \ {\isacharequal}{\kern0pt}\ map\ {\isacharparenleft}{\kern0pt}fn\ {\isacharparenleft}{\kern0pt}{\isacharparenleft}{\kern0pt}bdg{\isacharcomma}{\kern0pt}\ {\isacharunderscore}{\kern0pt}{\isacharparenright}{\kern0pt}{\isacharcomma}{\kern0pt}\ t{\isacharunderscore}{\kern0pt}opt{\isacharparenright}{\kern0pt}\ {\isacharequal}{\kern0pt}{\isachargreater}{\kern0pt}\isanewline
\ \ \ \ \ \ \ \ \ \ \ \ \ \ \ \ \ \ \ \ \ \ \ \ \ \ \ \ {\isacharparenleft}{\kern0pt}Binding{\isachardot}{\kern0pt}name{\isacharunderscore}{\kern0pt}of\ bdg{\isacharcomma}{\kern0pt}\ term{\isacharunderscore}{\kern0pt}to{\isacharunderscore}{\kern0pt}int\ t{\isacharunderscore}{\kern0pt}opt{\isacharparenright}{\kern0pt}{\isacharparenright}{\kern0pt}\ locals\isanewline
\isanewline
\ \ \ \ val\ all{\isacharunderscore}{\kern0pt}vals\ {\isacharequal}{\kern0pt}\ local{\isacharunderscore}{\kern0pt}vals\ {\isacharat}{\kern0pt}\ global{\isacharunderscore}{\kern0pt}vals\isanewline
\ \ in\isanewline
\ \ \ \ fn\ var\ {\isacharequal}{\kern0pt}{\isachargreater}{\kern0pt}\ case\ AList{\isachardot}{\kern0pt}lookup\ {\isacharparenleft}{\kern0pt}op\ {\isacharequal}{\kern0pt}{\isacharparenright}{\kern0pt}\ all{\isacharunderscore}{\kern0pt}vals\ var\ of\isanewline
\ \ \ \ \ \ \ \ \ \ \ \ \ \ \ \ SOME\ v\ {\isacharequal}{\kern0pt}{\isachargreater}{\kern0pt}\ v\isanewline
\ \ \ \ \ \ \ \ \ \ \ \ \ \ {\isacharbar}{\kern0pt}\ NONE\ {\isacharequal}{\kern0pt}{\isachargreater}{\kern0pt}\ {\isadigit{0}}\isanewline
\ \ end\isanewline
\isanewline
{\isacharparenleft}{\kern0pt}{\isacharasterisk}{\kern0pt}\ Fonction\ de\ vérification\ sémantique\ principale\ complète\ {\isacharasterisk}{\kern0pt}{\isacharparenright}{\kern0pt}\isanewline
fun\ semantic{\isacharunderscore}{\kern0pt}check\ {\isacharparenleft}{\kern0pt}absy{\isacharunderscore}{\kern0pt}data\ {\isacharcolon}{\kern0pt}\ absy{\isacharparenright}{\kern0pt}\ {\isacharparenleft}{\kern0pt}thy\ {\isacharcolon}{\kern0pt}\ theory{\isacharparenright}{\kern0pt}\ {\isacharcolon}{\kern0pt}\ theory\ {\isacharequal}{\kern0pt}\isanewline
\ \ let\isanewline
\ \ \ \ val\ {\isacharbraceleft}{\kern0pt}system{\isacharcomma}{\kern0pt}\ globals{\isacharunderscore}{\kern0pt}decls{\isacharcomma}{\kern0pt}\ varstab{\isacharcomma}{\kern0pt}\ locks{\isacharunderscore}{\kern0pt}decls{\isacharcomma}{\kern0pt}\ threads{\isacharunderscore}{\kern0pt}decls{\isacharbraceright}{\kern0pt}\ {\isacharequal}{\kern0pt}\ absy{\isacharunderscore}{\kern0pt}data\isanewline
\ \ \ \ \isanewline
\ \ \ \ val\ {\isacharunderscore}{\kern0pt}\ {\isacharequal}{\kern0pt}\ message\ false\ {\isacharparenleft}{\kern0pt}{\isachardoublequote}{\kern0pt}{\isacharequal}{\kern0pt}{\isacharequal}{\kern0pt}{\isacharequal}{\kern0pt}\ SEMANTIC\ ANALYSIS\ FOR\ SYSTEM{\isacharcolon}{\kern0pt}\ {\isachardoublequote}{\kern0pt}\ {\isacharcircum}{\kern0pt}\ \isanewline
\ \ \ \ \ \ \ \ \ \ \ \ \ \ \ \ \ \ \ \ \ \ \ \ \ \ Binding{\isachardot}{\kern0pt}name{\isacharunderscore}{\kern0pt}of\ system\ {\isacharcircum}{\kern0pt}\ {\isachardoublequote}{\kern0pt}\ {\isacharequal}{\kern0pt}{\isacharequal}{\kern0pt}{\isacharequal}{\kern0pt}{\isachardoublequote}{\kern0pt}{\isacharparenright}{\kern0pt}\isanewline
\ \ \ \ \isanewline
\ \ \ \ {\isacharparenleft}{\kern0pt}{\isacharasterisk}{\kern0pt}\ Vérification\ {\isadigit{1}}{\isacharcolon}{\kern0pt}\ Analyse\ des\ deadlocks\ {\isacharasterisk}{\kern0pt}{\isacharparenright}{\kern0pt}\isanewline
\ \ \ \ val\ deadlock{\isacharunderscore}{\kern0pt}warnings\ {\isacharequal}{\kern0pt}\ detect{\isacharunderscore}{\kern0pt}potential{\isacharunderscore}{\kern0pt}deadlock\ threads{\isacharunderscore}{\kern0pt}decls\isanewline
\ \ \ \ val\ {\isacharunderscore}{\kern0pt}\ {\isacharequal}{\kern0pt}\ if\ null\ deadlock{\isacharunderscore}{\kern0pt}warnings\ then\ \isanewline
\ \ \ \ \ \ \ \ \ \ \ \ \ \ message\ false\ {\isachardoublequote}{\kern0pt}{\isasymcheckmark}\ No\ potential\ deadlocks\ detected{\isachardoublequote}{\kern0pt}\isanewline
\ \ \ \ \ \ \ \ \ \ \ \ else\ \isanewline
\ \ \ \ \ \ \ \ \ \ \ \ \ \ {\isacharparenleft}{\kern0pt}message\ false\ {\isachardoublequote}{\kern0pt}DEADLOCK\ WARNINGS{\isacharcolon}{\kern0pt}{\isachardoublequote}{\kern0pt}{\isacharsemicolon}{\kern0pt}\isanewline
\ \ \ \ \ \ \ \ \ \ \ \ \ \ \ map\ {\isacharparenleft}{\kern0pt}message\ false{\isacharparenright}{\kern0pt}\ deadlock{\isacharunderscore}{\kern0pt}warnings{\isacharsemicolon}{\kern0pt}\ {\isacharparenleft}{\kern0pt}{\isacharparenright}{\kern0pt}{\isacharparenright}{\kern0pt}\isanewline
\ \ \ \ \isanewline
\ \ \ \ {\isacharparenleft}{\kern0pt}{\isacharasterisk}{\kern0pt}\ Vérification\ {\isadigit{2}}{\isacharcolon}{\kern0pt}\ Analyse\ des\ race\ conditions\ {\isacharasterisk}{\kern0pt}{\isacharparenright}{\kern0pt}\isanewline
\ \ \ \ val\ race{\isacharunderscore}{\kern0pt}warnings\ {\isacharequal}{\kern0pt}\ detect{\isacharunderscore}{\kern0pt}race{\isacharunderscore}{\kern0pt}conditions\ threads{\isacharunderscore}{\kern0pt}decls\ \ \isanewline
\ \ \ \ val\ {\isacharunderscore}{\kern0pt}\ {\isacharequal}{\kern0pt}\ if\ null\ race{\isacharunderscore}{\kern0pt}warnings\ then\isanewline
\ \ \ \ \ \ \ \ \ \ \ \ \ \ message\ false\ {\isachardoublequote}{\kern0pt}{\isasymcheckmark}\ No\ race\ conditions\ detected{\isachardoublequote}{\kern0pt}\isanewline
\ \ \ \ \ \ \ \ \ \ \ \ else\isanewline
\ \ \ \ \ \ \ \ \ \ \ \ \ \ {\isacharparenleft}{\kern0pt}message\ false\ {\isachardoublequote}{\kern0pt}RACE\ CONDITION\ WARNINGS{\isacharcolon}{\kern0pt}{\isachardoublequote}{\kern0pt}{\isacharsemicolon}{\kern0pt}\isanewline
\ \ \ \ \ \ \ \ \ \ \ \ \ \ \ map\ {\isacharparenleft}{\kern0pt}message\ false{\isacharparenright}{\kern0pt}\ race{\isacharunderscore}{\kern0pt}warnings{\isacharsemicolon}{\kern0pt}\ {\isacharparenleft}{\kern0pt}{\isacharparenright}{\kern0pt}{\isacharparenright}{\kern0pt}\isanewline
\ \ \ \ \isanewline
\ \ \ \ {\isacharparenleft}{\kern0pt}{\isacharasterisk}{\kern0pt}\ Vérification\ {\isadigit{3}}{\isacharcolon}{\kern0pt}\ Vérification\ des\ types\ pour\ chaque\ thread\ {\isacharasterisk}{\kern0pt}{\isacharparenright}{\kern0pt}\isanewline
\ \ \ \ fun\ check{\isacharunderscore}{\kern0pt}thread\ {\isacharparenleft}{\kern0pt}machine\ {\isacharcolon}{\kern0pt}absy{\isacharparenright}{\kern0pt}\ {\isacharparenleft}{\kern0pt}thread\ {\isacharcolon}{\kern0pt}\ thread{\isacharunderscore}{\kern0pt}absy{\isacharparenright}{\kern0pt}\ {\isacharcolon}{\kern0pt}\ unit\ {\isacharequal}{\kern0pt}\isanewline
\ \ \ \ \ \ let\isanewline
\ \ \ \ \ \ \ \ val{\isacharbraceleft}{\kern0pt}system{\isacharcomma}{\kern0pt}\ globals{\isacharunderscore}{\kern0pt}decls{\isacharcomma}{\kern0pt}\ varstab{\isacharcomma}{\kern0pt}\ locks{\isacharunderscore}{\kern0pt}decls{\isacharcomma}{\kern0pt}\ threads{\isacharunderscore}{\kern0pt}decls{\isacharbraceright}{\kern0pt}\ {\isacharequal}{\kern0pt}\ machine\isanewline
\ \ \ \ \ \ \ \ val\ {\isacharbraceleft}{\kern0pt}nom{\isacharunderscore}{\kern0pt}thread{\isacharcomma}{\kern0pt}\ locals{\isacharunderscore}{\kern0pt}decl{\isacharcomma}{\kern0pt}\ actions{\isacharcomma}{\kern0pt}\ locstab{\isacharbraceright}{\kern0pt}\ {\isacharequal}{\kern0pt}\ thread\isanewline
\ \ \ \ \ \ \ \ val\ thread{\isacharunderscore}{\kern0pt}name\ {\isacharequal}{\kern0pt}\ Binding{\isachardot}{\kern0pt}name{\isacharunderscore}{\kern0pt}of\ nom{\isacharunderscore}{\kern0pt}thread\isanewline
\ \ \ \ \ \ \ \ \isanewline
\ \ \ \ \ \ \ \ val\ {\isacharunderscore}{\kern0pt}\ {\isacharequal}{\kern0pt}\ message\ false\ {\isacharparenleft}{\kern0pt}{\isachardoublequote}{\kern0pt}Analyzing\ thread{\isacharcolon}{\kern0pt}\ {\isachardoublequote}{\kern0pt}\ {\isacharcircum}{\kern0pt}\ thread{\isacharunderscore}{\kern0pt}name{\isacharparenright}{\kern0pt}\isanewline
\ \ \ \ \ \ \ \ \isanewline
\ \ \ \ \ \ \ \ {\isacharparenleft}{\kern0pt}{\isacharasterisk}{\kern0pt}\ Conversion\ vers\ com\ pour\ analyse\ sémantique\ {\isacharasterisk}{\kern0pt}{\isacharparenright}{\kern0pt}\isanewline
\ \ \ \ \ \ \ \ val\ com{\isacharunderscore}{\kern0pt}program\ {\isacharequal}{\kern0pt}\ convert{\isacharunderscore}{\kern0pt}to{\isacharunderscore}{\kern0pt}com\ thy\ actions\ varstab\ locstab\isanewline
\isanewline
\ \ \ \ \ \ \ \ val\ initial{\isacharunderscore}{\kern0pt}state\ {\isacharequal}{\kern0pt}\ build{\isacharunderscore}{\kern0pt}initial{\isacharunderscore}{\kern0pt}state{\isacharunderscore}{\kern0pt}for{\isacharunderscore}{\kern0pt}thread\ globals{\isacharunderscore}{\kern0pt}decls\ locals{\isacharunderscore}{\kern0pt}decl\ thy\isanewline
\isanewline
\ \ \ \ \ \ \ \ val\ cps{\isacharunderscore}{\kern0pt}trace\ {\isacharequal}{\kern0pt}\ sem{\isacharunderscore}{\kern0pt}cps\ machine\ com{\isacharunderscore}{\kern0pt}program\ {\isacharparenleft}{\kern0pt}fn\ {\isacharunderscore}{\kern0pt}\ {\isacharequal}{\kern0pt}{\isachargreater}{\kern0pt}\ {\isacharbrackleft}{\kern0pt}{\isacharbrackright}{\kern0pt}{\isacharparenright}{\kern0pt}\ initial{\isacharunderscore}{\kern0pt}state\isanewline
\isanewline
\ \ \ \ \ \ \ \ val\ {\isacharunderscore}{\kern0pt}\ {\isacharequal}{\kern0pt}\ message\ false\ {\isacharparenleft}{\kern0pt}{\isachardoublequote}{\kern0pt}\ {\isasymcheckmark}\ CPS{\isacharminus}{\kern0pt}style\ trace\ for\ {\isachardoublequote}{\kern0pt}\ {\isacharcircum}{\kern0pt}\ thread{\isacharunderscore}{\kern0pt}name\ {\isacharcircum}{\kern0pt}\ {\isachardoublequote}{\kern0pt}{\isacharcolon}{\kern0pt}{\isachardoublequote}{\kern0pt}{\isacharparenright}{\kern0pt}\isanewline
\ \ \ \ \ \ \ \ val\ {\isacharunderscore}{\kern0pt}\ {\isacharequal}{\kern0pt}\ List{\isachardot}{\kern0pt}app\ {\isacharparenleft}{\kern0pt}fn\ ev\ {\isacharequal}{\kern0pt}{\isachargreater}{\kern0pt}\ message\ false\ {\isacharparenleft}{\kern0pt}{\isachardoublequote}{\kern0pt}\ \ \ \ {\isachardoublequote}{\kern0pt}\ {\isacharcircum}{\kern0pt}\ {\isacharat}{\kern0pt}{\isacharbraceleft}{\kern0pt}make{\isacharunderscore}{\kern0pt}string{\isacharbraceright}{\kern0pt}\ ev{\isacharparenright}{\kern0pt}{\isacharparenright}{\kern0pt}\ cps{\isacharunderscore}{\kern0pt}trace\isanewline
\ \ \ \ \ \ \ \ {\isacharparenleft}{\kern0pt}{\isacharasterisk}{\kern0pt}\ Vérification\ des\ types\ {\isacharasterisk}{\kern0pt}{\isacharparenright}{\kern0pt}\isanewline
\ \ \ \ \ \ \ \ val\ type{\isacharunderscore}{\kern0pt}errors\ {\isacharequal}{\kern0pt}\ type{\isacharunderscore}{\kern0pt}check{\isacharunderscore}{\kern0pt}thread\ thy\ thread\ globals{\isacharunderscore}{\kern0pt}decls\isanewline
\ \ \ \ \ \ \ \ val\ {\isacharunderscore}{\kern0pt}\ {\isacharequal}{\kern0pt}\ if\ null\ type{\isacharunderscore}{\kern0pt}errors\ then\isanewline
\ \ \ \ \ \ \ \ \ \ \ \ \ \ \ \ \ \ message\ false\ {\isacharparenleft}{\kern0pt}{\isachardoublequote}{\kern0pt}\ {\isasymcheckmark}\ Type\ checking\ passed\ for\ {\isachardoublequote}{\kern0pt}\ {\isacharcircum}{\kern0pt}\ thread{\isacharunderscore}{\kern0pt}name{\isacharparenright}{\kern0pt}\isanewline
\ \ \ \ \ \ \ \ \ \ \ \ \ \ \ \ else\isanewline
\ \ \ \ \ \ \ \ \ \ \ \ \ \ \ \ \ \ {\isacharparenleft}{\kern0pt}message\ false\ {\isacharparenleft}{\kern0pt}{\isachardoublequote}{\kern0pt}\ {\isasymcrossmark}\ Type\ errors\ in\ {\isachardoublequote}{\kern0pt}\ {\isacharcircum}{\kern0pt}\ thread{\isacharunderscore}{\kern0pt}name\ {\isacharcircum}{\kern0pt}\ {\isachardoublequote}{\kern0pt}{\isacharcolon}{\kern0pt}{\isachardoublequote}{\kern0pt}{\isacharparenright}{\kern0pt}{\isacharsemicolon}{\kern0pt}\isanewline
\ \ \ \ \ \ \ \ \ \ \ \ \ \ \ \ \ \ \ map\ {\isacharparenleft}{\kern0pt}fn\ err\ {\isacharequal}{\kern0pt}{\isachargreater}{\kern0pt}\ message\ false\ {\isacharparenleft}{\kern0pt}{\isachardoublequote}{\kern0pt}\ \ \ \ {\isachardoublequote}{\kern0pt}\ {\isacharcircum}{\kern0pt}\ err{\isacharparenright}{\kern0pt}{\isacharparenright}{\kern0pt}\ type{\isacharunderscore}{\kern0pt}errors{\isacharsemicolon}{\kern0pt}\ {\isacharparenleft}{\kern0pt}{\isacharparenright}{\kern0pt}{\isacharparenright}{\kern0pt}\isanewline
\ \ \ \ \ \ \ \ {\isacharparenleft}{\kern0pt}{\isacharasterisk}{\kern0pt}\ Vérification\ que\ lock\ et\ unlock\ respectent\ des\ regles\ d{\isacharprime}{\kern0pt}appartenance\ {\isacharasterisk}{\kern0pt}{\isacharparenright}{\kern0pt}\isanewline
\ \ \ \ \ \ \ \ val\ {\isacharunderscore}{\kern0pt}\ {\isacharequal}{\kern0pt}\ case\ check{\isacharunderscore}{\kern0pt}lock{\isacharunderscore}{\kern0pt}ownership\ thread{\isacharunderscore}{\kern0pt}name\ actions\ of\isanewline
\ \ \ \ \ \ \ \ \ \ NONE\ {\isacharequal}{\kern0pt}{\isachargreater}{\kern0pt}\ {\isacharparenleft}{\kern0pt}{\isacharparenright}{\kern0pt}\isanewline
\ \ \ \ \ \ \ \ {\isacharbar}{\kern0pt}\ SOME\ msg\ {\isacharequal}{\kern0pt}{\isachargreater}{\kern0pt}\ error\ msg\isanewline
\isanewline
\ \ \ \ \ \ \ \ {\isacharparenleft}{\kern0pt}{\isacharasterisk}{\kern0pt}\ Simulation\ d{\isacharprime}{\kern0pt}exécution\ pour\ détecter\ d{\isacharprime}{\kern0pt}autres\ problèmes\ {\isacharasterisk}{\kern0pt}{\isacharparenright}{\kern0pt}\isanewline
\ \ \ \ \ \ \ \ val\ initial{\isacharunderscore}{\kern0pt}state\ {\isacharequal}{\kern0pt}\ build{\isacharunderscore}{\kern0pt}initial{\isacharunderscore}{\kern0pt}state{\isacharunderscore}{\kern0pt}for{\isacharunderscore}{\kern0pt}thread\ globals{\isacharunderscore}{\kern0pt}decls\ locals{\isacharunderscore}{\kern0pt}decl\ thy\isanewline
\ \ \ \ \ \ \ \ val\ {\isacharparenleft}{\kern0pt}final{\isacharunderscore}{\kern0pt}state{\isacharcomma}{\kern0pt}\ execution{\isacharunderscore}{\kern0pt}trace{\isacharparenright}{\kern0pt}\ {\isacharequal}{\kern0pt}\ \isanewline
\ \ \ \ \ \ \ \ \ \ sem{\isacharunderscore}{\kern0pt}step\ thy\ machine\ \ com{\isacharunderscore}{\kern0pt}program\ {\isacharparenleft}{\kern0pt}fn\ {\isacharparenleft}{\kern0pt}s{\isacharcomma}{\kern0pt}\ t{\isacharparenright}{\kern0pt}\ {\isacharequal}{\kern0pt}{\isachargreater}{\kern0pt}\ {\isacharparenleft}{\kern0pt}s{\isacharcomma}{\kern0pt}\ t{\isacharparenright}{\kern0pt}{\isacharparenright}{\kern0pt}\ {\isacharparenleft}{\kern0pt}initial{\isacharunderscore}{\kern0pt}state{\isacharcomma}{\kern0pt}\ {\isacharbrackleft}{\kern0pt}{\isacharbrackright}{\kern0pt}{\isacharparenright}{\kern0pt}\isanewline
\ \ \ \ \ \ \ \ \isanewline
\ \ \ \ \ \ \ \ val\ {\isacharunderscore}{\kern0pt}\ {\isacharequal}{\kern0pt}\ message\ false\ {\isacharparenleft}{\kern0pt}{\isachardoublequote}{\kern0pt}\ \ {\isasymcheckmark}\ Execution\ simulation\ completed\ for\ {\isachardoublequote}{\kern0pt}\ {\isacharcircum}{\kern0pt}\ thread{\isacharunderscore}{\kern0pt}name{\isacharparenright}{\kern0pt}\isanewline
\ \ \ \ \ \ \ \ val\ {\isacharunderscore}{\kern0pt}\ {\isacharequal}{\kern0pt}\ message\ false\ {\isacharparenleft}{\kern0pt}{\isachardoublequote}{\kern0pt}\ Generated\ {\isachardoublequote}{\kern0pt}\ {\isacharcircum}{\kern0pt}\ \isanewline
\ \ \ \ \ \ \ \ \ \ \ \ \ \ \ \ \ \ \ \ \ \ \ \ \ \ \ \ \ Int{\isachardot}{\kern0pt}toString\ {\isacharparenleft}{\kern0pt}length\ execution{\isacharunderscore}{\kern0pt}trace{\isacharparenright}{\kern0pt}\ {\isacharcircum}{\kern0pt}\ {\isachardoublequote}{\kern0pt}\ CSP\ events{\isachardoublequote}{\kern0pt}{\isacharparenright}{\kern0pt}\isanewline
\ \ \ \ \ \ in\isanewline
\ \ \ \ \ \ \ \ {\isacharparenleft}{\kern0pt}{\isacharparenright}{\kern0pt}\isanewline
\ \ \ \ \ \ end\isanewline
\ \ \ \ \isanewline
\ \ \ \ val\ {\isacharunderscore}{\kern0pt}\ {\isacharequal}{\kern0pt}\ map\ {\isacharparenleft}{\kern0pt}check{\isacharunderscore}{\kern0pt}thread\ absy{\isacharunderscore}{\kern0pt}data{\isacharparenright}{\kern0pt}\ {\isacharparenleft}{\kern0pt}threads{\isacharunderscore}{\kern0pt}decls{\isacharparenright}{\kern0pt}\isanewline
\ \ \ \ \isanewline
\ \ \ \ {\isacharparenleft}{\kern0pt}{\isacharasterisk}{\kern0pt}\ Vérification\ {\isadigit{4}}{\isacharcolon}{\kern0pt}\ Cohérence\ globale\ du\ système\ {\isacharasterisk}{\kern0pt}{\isacharparenright}{\kern0pt}\isanewline
\ \ \ \ val\ {\isacharunderscore}{\kern0pt}\ {\isacharequal}{\kern0pt}\ message\ false\ {\isachardoublequote}{\kern0pt}{\isacharequal}{\kern0pt}{\isacharequal}{\kern0pt}{\isacharequal}{\kern0pt}\ GLOBAL\ SYSTEM\ CONSISTENCY\ {\isacharequal}{\kern0pt}{\isacharequal}{\kern0pt}{\isacharequal}{\kern0pt}{\isachardoublequote}{\kern0pt}\isanewline
\ \ \ \ \isanewline
\ \ \ \ {\isacharparenleft}{\kern0pt}{\isacharasterisk}{\kern0pt}\ Vérifier\ que\ tous\ les\ verrous\ utilisés\ sont\ déclarés\ {\isacharasterisk}{\kern0pt}{\isacharparenright}{\kern0pt}\isanewline
\ \ \ \ fun\ check{\isacharunderscore}{\kern0pt}lock{\isacharunderscore}{\kern0pt}declarations\ {\isacharparenleft}{\kern0pt}{\isacharparenright}{\kern0pt}\ {\isacharequal}{\kern0pt}\isanewline
\ \ \ \ \ \ let\isanewline
\ \ \ \ \ \ \ \ val\ declared{\isacharunderscore}{\kern0pt}locks\ {\isacharequal}{\kern0pt}\ map\ {\isacharparenleft}{\kern0pt}fn\ {\isacharparenleft}{\kern0pt}bdg{\isacharcomma}{\kern0pt}\ {\isacharunderscore}{\kern0pt}{\isacharcomma}{\kern0pt}\ {\isacharunderscore}{\kern0pt}{\isacharparenright}{\kern0pt}\ {\isacharequal}{\kern0pt}{\isachargreater}{\kern0pt}\ Binding{\isachardot}{\kern0pt}name{\isacharunderscore}{\kern0pt}of\ bdg{\isacharparenright}{\kern0pt}\ locks{\isacharunderscore}{\kern0pt}decls\isanewline
\ \ \ \ \ \ \ \ \isanewline
\ \ \ \ \ \ \ \ fun\ extract{\isacharunderscore}{\kern0pt}used{\isacharunderscore}{\kern0pt}locks{\isacharunderscore}{\kern0pt}from{\isacharunderscore}{\kern0pt}thread\ {\isacharparenleft}{\kern0pt}thread\ {\isacharcolon}{\kern0pt}\ thread{\isacharunderscore}{\kern0pt}absy{\isacharparenright}{\kern0pt}\ {\isacharequal}{\kern0pt}\isanewline
\ \ \ \ \ \ \ \ \ \ let\isanewline
\ \ \ \ \ \ \ \ \ \ \ \ fun\ extract{\isacharunderscore}{\kern0pt}from{\isacharunderscore}{\kern0pt}actions\ actions\ {\isacharequal}{\kern0pt}\isanewline
\ \ \ \ \ \ \ \ \ \ \ \ \ \ List{\isachardot}{\kern0pt}concat\ {\isacharparenleft}{\kern0pt}map\ {\isacharparenleft}{\kern0pt}fn\ action\ {\isacharequal}{\kern0pt}{\isachargreater}{\kern0pt}\isanewline
\ \ \ \ \ \ \ \ \ \ \ \ \ \ \ \ case\ action\ of\isanewline
\ \ \ \ \ \ \ \ \ \ \ \ \ \ \ \ \ \ LockA\ bdg\ {\isacharequal}{\kern0pt}{\isachargreater}{\kern0pt}\ {\isacharbrackleft}{\kern0pt}Binding{\isachardot}{\kern0pt}name{\isacharunderscore}{\kern0pt}of\ bdg{\isacharbrackright}{\kern0pt}\isanewline
\ \ \ \ \ \ \ \ \ \ \ \ \ \ \ \ {\isacharbar}{\kern0pt}\ UnlockA\ bdg\ {\isacharequal}{\kern0pt}{\isachargreater}{\kern0pt}\ {\isacharbrackleft}{\kern0pt}Binding{\isachardot}{\kern0pt}name{\isacharunderscore}{\kern0pt}of\ bdg{\isacharbrackright}{\kern0pt}\isanewline
\ \ \ \ \ \ \ \ \ \ \ \ \ \ \ \ {\isacharbar}{\kern0pt}\ IfelseA\ {\isacharparenleft}{\kern0pt}{\isacharparenleft}{\kern0pt}{\isacharunderscore}{\kern0pt}{\isacharcomma}{\kern0pt}\ then{\isacharunderscore}{\kern0pt}acts{\isacharparenright}{\kern0pt}{\isacharcomma}{\kern0pt}\ else{\isacharunderscore}{\kern0pt}acts{\isacharparenright}{\kern0pt}\ {\isacharequal}{\kern0pt}{\isachargreater}{\kern0pt}\isanewline
\ \ \ \ \ \ \ \ \ \ \ \ \ \ \ \ \ \ \ \ extract{\isacharunderscore}{\kern0pt}from{\isacharunderscore}{\kern0pt}actions\ then{\isacharunderscore}{\kern0pt}acts\ {\isacharat}{\kern0pt}\ extract{\isacharunderscore}{\kern0pt}from{\isacharunderscore}{\kern0pt}actions\ else{\isacharunderscore}{\kern0pt}acts\isanewline
\ \ \ \ \ \ \ \ \ \ \ \ \ \ \ \ {\isacharbar}{\kern0pt}\ WhileA\ {\isacharparenleft}{\kern0pt}{\isacharunderscore}{\kern0pt}{\isacharcomma}{\kern0pt}\ body{\isacharunderscore}{\kern0pt}acts{\isacharparenright}{\kern0pt}\ {\isacharequal}{\kern0pt}{\isachargreater}{\kern0pt}\ extract{\isacharunderscore}{\kern0pt}from{\isacharunderscore}{\kern0pt}actions\ body{\isacharunderscore}{\kern0pt}acts\isanewline
\ \ \ \ \ \ \ \ \ \ \ \ \ \ \ \ {\isacharbar}{\kern0pt}\ {\isacharunderscore}{\kern0pt}\ {\isacharequal}{\kern0pt}{\isachargreater}{\kern0pt}\ {\isacharbrackleft}{\kern0pt}{\isacharbrackright}{\kern0pt}{\isacharparenright}{\kern0pt}\ actions{\isacharparenright}{\kern0pt}\isanewline
\ \ \ \ \ \ \ \ \ \ in\isanewline
\ \ \ \ \ \ \ \ \ \ \ \ extract{\isacharunderscore}{\kern0pt}from{\isacharunderscore}{\kern0pt}actions\ {\isacharparenleft}{\kern0pt}{\isacharhash}{\kern0pt}actions\ thread{\isacharparenright}{\kern0pt}\isanewline
\ \ \ \ \ \ \ \ \ \ end\isanewline
\ \ \ \ \ \ \ \ \isanewline
\ \ \ \ \ \ \ \ val\ used{\isacharunderscore}{\kern0pt}locks\ {\isacharequal}{\kern0pt}\ List{\isachardot}{\kern0pt}concat\ {\isacharparenleft}{\kern0pt}map\ extract{\isacharunderscore}{\kern0pt}used{\isacharunderscore}{\kern0pt}locks{\isacharunderscore}{\kern0pt}from{\isacharunderscore}{\kern0pt}thread\ threads{\isacharunderscore}{\kern0pt}decls{\isacharparenright}{\kern0pt}\isanewline
\ \ \ \ \ \ \ \ val\ undeclared{\isacharunderscore}{\kern0pt}locks\ {\isacharequal}{\kern0pt}\ List{\isachardot}{\kern0pt}filter\ {\isacharparenleft}{\kern0pt}fn\ used\ {\isacharequal}{\kern0pt}{\isachargreater}{\kern0pt}\ \isanewline
\ \ \ \ \ \ \ \ \ \ not\ {\isacharparenleft}{\kern0pt}List{\isachardot}{\kern0pt}exists\ {\isacharparenleft}{\kern0pt}fn\ decl\ {\isacharequal}{\kern0pt}{\isachargreater}{\kern0pt}\ decl\ {\isacharequal}{\kern0pt}\ used{\isacharparenright}{\kern0pt}\ declared{\isacharunderscore}{\kern0pt}locks{\isacharparenright}{\kern0pt}{\isacharparenright}{\kern0pt}\ used{\isacharunderscore}{\kern0pt}locks\isanewline
\ \ \ \ \ \ in\isanewline
\ \ \ \ \ \ \ \ if\ null\ undeclared{\isacharunderscore}{\kern0pt}locks\ then\isanewline
\ \ \ \ \ \ \ \ \ \ message\ false\ {\isachardoublequote}{\kern0pt}{\isasymcheckmark}\ All\ used\ locks\ are\ properly\ declared{\isachardoublequote}{\kern0pt}\isanewline
\ \ \ \ \ \ \ \ else\isanewline
\ \ \ \ \ \ \ \ \ \ {\isacharparenleft}{\kern0pt}message\ false\ {\isachardoublequote}{\kern0pt}{\isasymcrossmark}\ Undeclared\ locks\ used{\isacharcolon}{\kern0pt}{\isachardoublequote}{\kern0pt}{\isacharsemicolon}{\kern0pt}\isanewline
\ \ \ \ \ \ \ \ \ \ \ map\ {\isacharparenleft}{\kern0pt}fn\ lock\ {\isacharequal}{\kern0pt}{\isachargreater}{\kern0pt}\ message\ false\ {\isacharparenleft}{\kern0pt}{\isachardoublequote}{\kern0pt}\ \ {\isachardoublequote}{\kern0pt}\ {\isacharcircum}{\kern0pt}\ lock{\isacharparenright}{\kern0pt}{\isacharparenright}{\kern0pt}\ undeclared{\isacharunderscore}{\kern0pt}locks{\isacharsemicolon}{\kern0pt}\ {\isacharparenleft}{\kern0pt}{\isacharparenright}{\kern0pt}{\isacharparenright}{\kern0pt}\isanewline
\ \ \ \ \ \ end\isanewline
\ \ \ \ \isanewline
\ \ \ \ val\ {\isacharunderscore}{\kern0pt}\ {\isacharequal}{\kern0pt}\ check{\isacharunderscore}{\kern0pt}lock{\isacharunderscore}{\kern0pt}declarations\ {\isacharparenleft}{\kern0pt}{\isacharparenright}{\kern0pt}\isanewline
\ \ \ \ \isanewline
\ \ \ \ fun\ check{\isacharunderscore}{\kern0pt}var{\isacharunderscore}{\kern0pt}declarations\ {\isacharparenleft}{\kern0pt}{\isacharparenright}{\kern0pt}\ {\isacharequal}{\kern0pt}\isanewline
\ \ \ \ \ \ let\isanewline
\ \ \ \ \ \ \ \ {\isacharparenleft}{\kern0pt}{\isacharasterisk}{\kern0pt}\ {\isacharminus}{\kern0pt}{\isacharminus}{\kern0pt}{\isacharminus}{\kern0pt}{\isacharminus}{\kern0pt}{\isacharminus}{\kern0pt}{\isacharminus}{\kern0pt}{\isacharminus}{\kern0pt}{\isacharminus}{\kern0pt}{\isacharminus}{\kern0pt}{\isacharminus}{\kern0pt}{\isacharminus}{\kern0pt}{\isacharminus}{\kern0pt}{\isacharminus}{\kern0pt}{\isacharminus}{\kern0pt}{\isacharminus}{\kern0pt}{\isacharminus}{\kern0pt}{\isacharminus}{\kern0pt}{\isacharminus}{\kern0pt}\ {\isadigit{1}}{\isachardot}{\kern0pt}\ Ensemble\ des\ noms\ globaux\ {\isacharminus}{\kern0pt}{\isacharminus}{\kern0pt}{\isacharminus}{\kern0pt}{\isacharminus}{\kern0pt}{\isacharminus}{\kern0pt}{\isacharminus}{\kern0pt}{\isacharminus}{\kern0pt}{\isacharminus}{\kern0pt}{\isacharminus}{\kern0pt}{\isacharminus}{\kern0pt}{\isacharminus}{\kern0pt}{\isacharminus}{\kern0pt}{\isacharminus}{\kern0pt}{\isacharminus}{\kern0pt}{\isacharminus}{\kern0pt}{\isacharminus}{\kern0pt}\ {\isacharasterisk}{\kern0pt}{\isacharparenright}{\kern0pt}\isanewline
\ \ \ \ \ \ \ \ val\ global{\isacharunderscore}{\kern0pt}names\ {\isacharequal}{\kern0pt}\isanewline
\ \ \ \ \ \ \ \ \ \ map\ {\isacharparenleft}{\kern0pt}fn\ {\isacharparenleft}{\kern0pt}bdg{\isacharcomma}{\kern0pt}\ {\isacharunderscore}{\kern0pt}{\isacharcomma}{\kern0pt}\ {\isacharunderscore}{\kern0pt}{\isacharcomma}{\kern0pt}\ {\isacharunderscore}{\kern0pt}{\isacharparenright}{\kern0pt}\ {\isacharequal}{\kern0pt}{\isachargreater}{\kern0pt}\ Binding{\isachardot}{\kern0pt}name{\isacharunderscore}{\kern0pt}of\ bdg{\isacharparenright}{\kern0pt}\ globals{\isacharunderscore}{\kern0pt}decls\isanewline
\ \ \ \ \isanewline
\ \ \ \ \ \ \ \ {\isacharparenleft}{\kern0pt}{\isacharasterisk}{\kern0pt}\ {\isacharminus}{\kern0pt}{\isacharminus}{\kern0pt}{\isacharminus}{\kern0pt}{\isacharminus}{\kern0pt}{\isacharminus}{\kern0pt}{\isacharminus}{\kern0pt}{\isacharminus}{\kern0pt}{\isacharminus}{\kern0pt}{\isacharminus}{\kern0pt}{\isacharminus}{\kern0pt}{\isacharminus}{\kern0pt}{\isacharminus}{\kern0pt}{\isacharminus}{\kern0pt}{\isacharminus}{\kern0pt}{\isacharminus}{\kern0pt}{\isacharminus}{\kern0pt}{\isacharminus}{\kern0pt}{\isacharminus}{\kern0pt}\ {\isadigit{2}}{\isachardot}{\kern0pt}\ Noms\ locaux\ d’un\ thread\ {\isacharminus}{\kern0pt}{\isacharminus}{\kern0pt}{\isacharminus}{\kern0pt}{\isacharminus}{\kern0pt}{\isacharminus}{\kern0pt}{\isacharminus}{\kern0pt}{\isacharminus}{\kern0pt}{\isacharminus}{\kern0pt}{\isacharminus}{\kern0pt}{\isacharminus}{\kern0pt}{\isacharminus}{\kern0pt}{\isacharminus}{\kern0pt}{\isacharminus}{\kern0pt}{\isacharminus}{\kern0pt}{\isacharminus}{\kern0pt}{\isacharminus}{\kern0pt}{\isacharminus}{\kern0pt}{\isacharminus}{\kern0pt}\ {\isacharasterisk}{\kern0pt}{\isacharparenright}{\kern0pt}\isanewline
\ \ \ \ \ \ \ \ fun\ local{\isacharunderscore}{\kern0pt}names{\isacharunderscore}{\kern0pt}of{\isacharunderscore}{\kern0pt}thread\ {\isacharparenleft}{\kern0pt}{\isacharbraceleft}{\kern0pt}locals{\isacharunderscore}{\kern0pt}decl{\isacharcomma}{\kern0pt}\ {\isachardot}{\kern0pt}{\isachardot}{\kern0pt}{\isachardot}{\kern0pt}{\isacharbraceright}{\kern0pt}\ {\isacharcolon}{\kern0pt}\ thread{\isacharunderscore}{\kern0pt}absy{\isacharparenright}{\kern0pt}\ {\isacharcolon}{\kern0pt}\ string\ list\ {\isacharequal}{\kern0pt}\isanewline
\ \ \ \ \ \ \ \ \ \ map\ {\isacharparenleft}{\kern0pt}fn\ {\isacharparenleft}{\kern0pt}{\isacharparenleft}{\kern0pt}bdg{\isacharcomma}{\kern0pt}\ {\isacharunderscore}{\kern0pt}{\isacharparenright}{\kern0pt}{\isacharcomma}{\kern0pt}\ {\isacharunderscore}{\kern0pt}{\isacharparenright}{\kern0pt}\ {\isacharequal}{\kern0pt}{\isachargreater}{\kern0pt}\ Binding{\isachardot}{\kern0pt}name{\isacharunderscore}{\kern0pt}of\ bdg{\isacharparenright}{\kern0pt}\ locals{\isacharunderscore}{\kern0pt}decl\isanewline
\ \ \ \ \isanewline
\ \ \ \ \ \ \ \ {\isacharparenleft}{\kern0pt}{\isacharasterisk}{\kern0pt}\ {\isacharminus}{\kern0pt}{\isacharminus}{\kern0pt}{\isacharminus}{\kern0pt}{\isacharminus}{\kern0pt}{\isacharminus}{\kern0pt}{\isacharminus}{\kern0pt}{\isacharminus}{\kern0pt}{\isacharminus}{\kern0pt}{\isacharminus}{\kern0pt}{\isacharminus}{\kern0pt}{\isacharminus}{\kern0pt}{\isacharminus}{\kern0pt}{\isacharminus}{\kern0pt}{\isacharminus}{\kern0pt}{\isacharminus}{\kern0pt}{\isacharminus}{\kern0pt}{\isacharminus}{\kern0pt}{\isacharminus}{\kern0pt}\ {\isadigit{3}}{\isachardot}{\kern0pt}\ Variables\ utilisées\ dans\ des\ actions\ {\isacharminus}{\kern0pt}{\isacharminus}{\kern0pt}{\isacharminus}{\kern0pt}{\isacharminus}{\kern0pt}{\isacharminus}{\kern0pt}\ {\isacharasterisk}{\kern0pt}{\isacharparenright}{\kern0pt}\isanewline
\ \ \ \ \ \ \ \ fun\ extract{\isacharunderscore}{\kern0pt}vars\ {\isacharparenleft}{\kern0pt}actions\ {\isacharcolon}{\kern0pt}\ a{\isacharunderscore}{\kern0pt}term\ list{\isacharparenright}{\kern0pt}\ {\isacharcolon}{\kern0pt}\ string\ list\ {\isacharequal}{\kern0pt}\isanewline
\ \ \ \ \ \ \ \ \ \ List{\isachardot}{\kern0pt}concat\ {\isacharparenleft}{\kern0pt}map\ {\isacharparenleft}{\kern0pt}fn\ action\ {\isacharequal}{\kern0pt}{\isachargreater}{\kern0pt}\isanewline
\ \ \ \ \ \ \ \ \ \ \ \ case\ action\ of\isanewline
\ \ \ \ \ \ \ \ \ \ \ \ \ \ \ \ AssignA\ \ {\isacharparenleft}{\kern0pt}bdg{\isacharcomma}{\kern0pt}\ {\isacharunderscore}{\kern0pt}{\isacharparenright}{\kern0pt}\ \ \ \ \ \ \ \ {\isacharequal}{\kern0pt}{\isachargreater}{\kern0pt}\ {\isacharbrackleft}{\kern0pt}Binding{\isachardot}{\kern0pt}name{\isacharunderscore}{\kern0pt}of\ bdg{\isacharbrackright}{\kern0pt}\isanewline
\ \ \ \ \ \ \ \ \ \ \ \ \ \ {\isacharbar}{\kern0pt}\ SendA\ \ \ \ {\isacharparenleft}{\kern0pt}bdg{\isacharcomma}{\kern0pt}\ {\isacharunderscore}{\kern0pt}{\isacharparenright}{\kern0pt}\ \ \ \ \ \ \ \ {\isacharequal}{\kern0pt}{\isachargreater}{\kern0pt}\ {\isacharbrackleft}{\kern0pt}Binding{\isachardot}{\kern0pt}name{\isacharunderscore}{\kern0pt}of\ bdg{\isacharbrackright}{\kern0pt}\isanewline
\ \ \ \ \ \ \ \ \ \ \ \ \ \ {\isacharbar}{\kern0pt}\ ReceiveA\ {\isacharparenleft}{\kern0pt}lbdg{\isacharcomma}{\kern0pt}\ sbdg{\isacharparenright}{\kern0pt}\ \ \ \ {\isacharequal}{\kern0pt}{\isachargreater}{\kern0pt}\ {\isacharbrackleft}{\kern0pt}Binding{\isachardot}{\kern0pt}name{\isacharunderscore}{\kern0pt}of\ lbdg{\isacharcomma}{\kern0pt}\isanewline
\ \ \ \ \ \ \ \ \ \ \ \ \ \ \ \ \ \ \ \ \ \ \ \ \ \ \ \ \ \ \ \ \ \ \ \ \ \ \ \ \ \ \ \ \ Binding{\isachardot}{\kern0pt}name{\isacharunderscore}{\kern0pt}of\ sbdg{\isacharbrackright}{\kern0pt}\isanewline
\ \ \ \ \ \ \ \ \ \ \ \ \ \ {\isacharbar}{\kern0pt}\ IfelseA\ \ {\isacharparenleft}{\kern0pt}{\isacharparenleft}{\kern0pt}{\isacharunderscore}{\kern0pt}{\isacharcomma}{\kern0pt}\ t{\isadigit{1}}{\isacharparenright}{\kern0pt}{\isacharcomma}{\kern0pt}\ t{\isadigit{2}}{\isacharparenright}{\kern0pt}\ \ \ {\isacharequal}{\kern0pt}{\isachargreater}{\kern0pt}\ extract{\isacharunderscore}{\kern0pt}vars\ t{\isadigit{1}}\ {\isacharat}{\kern0pt}\ extract{\isacharunderscore}{\kern0pt}vars\ t{\isadigit{2}}\isanewline
\ \ \ \ \ \ \ \ \ \ \ \ \ \ {\isacharbar}{\kern0pt}\ WhileA\ \ \ {\isacharparenleft}{\kern0pt}{\isacharunderscore}{\kern0pt}{\isacharcomma}{\kern0pt}\ body{\isacharparenright}{\kern0pt}\ \ \ \ \ \ \ {\isacharequal}{\kern0pt}{\isachargreater}{\kern0pt}\ extract{\isacharunderscore}{\kern0pt}vars\ body\isanewline
\ \ \ \ \ \ \ \ \ \ \ \ \ \ {\isacharbar}{\kern0pt}\ {\isacharunderscore}{\kern0pt}\ \ \ \ \ \ \ \ \ \ \ \ \ \ \ \ \ \ \ \ \ \ \ \ {\isacharequal}{\kern0pt}{\isachargreater}{\kern0pt}\ {\isacharbrackleft}{\kern0pt}{\isacharbrackright}{\kern0pt}{\isacharparenright}{\kern0pt}\ actions{\isacharparenright}{\kern0pt}\isanewline
\ \ \ \ \isanewline
\ \ \ \ \ \ \ \ {\isacharparenleft}{\kern0pt}{\isacharasterisk}{\kern0pt}\ {\isacharminus}{\kern0pt}{\isacharminus}{\kern0pt}{\isacharminus}{\kern0pt}{\isacharminus}{\kern0pt}{\isacharminus}{\kern0pt}{\isacharminus}{\kern0pt}{\isacharminus}{\kern0pt}{\isacharminus}{\kern0pt}{\isacharminus}{\kern0pt}{\isacharminus}{\kern0pt}{\isacharminus}{\kern0pt}{\isacharminus}{\kern0pt}{\isacharminus}{\kern0pt}{\isacharminus}{\kern0pt}{\isacharminus}{\kern0pt}{\isacharminus}{\kern0pt}{\isacharminus}{\kern0pt}{\isacharminus}{\kern0pt}\ {\isadigit{4}}{\isachardot}{\kern0pt}\ Manquantes\ dans\ un\ thread\ donné\ {\isacharminus}{\kern0pt}{\isacharminus}{\kern0pt}{\isacharminus}{\kern0pt}{\isacharminus}{\kern0pt}{\isacharminus}{\kern0pt}{\isacharminus}{\kern0pt}{\isacharminus}{\kern0pt}{\isacharminus}{\kern0pt}{\isacharminus}{\kern0pt}{\isacharminus}{\kern0pt}\ {\isacharasterisk}{\kern0pt}{\isacharparenright}{\kern0pt}\isanewline
\ \ \ \ \ \ \ \ fun\ undeclared{\isacharunderscore}{\kern0pt}in{\isacharunderscore}{\kern0pt}thread\ {\isacharparenleft}{\kern0pt}thr\ {\isacharcolon}{\kern0pt}\ thread{\isacharunderscore}{\kern0pt}absy{\isacharparenright}{\kern0pt}\ {\isacharcolon}{\kern0pt}\ string\ list\ {\isacharequal}{\kern0pt}\isanewline
\ \ \ \ \ \ \ \ \ \ let\isanewline
\ \ \ \ \ \ \ \ \ \ \ \ val\ declared\ {\isacharequal}{\kern0pt}\ global{\isacharunderscore}{\kern0pt}names\ {\isacharat}{\kern0pt}\ local{\isacharunderscore}{\kern0pt}names{\isacharunderscore}{\kern0pt}of{\isacharunderscore}{\kern0pt}thread\ thr\isanewline
\ \ \ \ \ \ \ \ \ \ \ \ val\ used\ \ \ \ \ {\isacharequal}{\kern0pt}\ extract{\isacharunderscore}{\kern0pt}vars\ {\isacharparenleft}{\kern0pt}{\isacharhash}{\kern0pt}actions\ thr{\isacharparenright}{\kern0pt}\isanewline
\ \ \ \ \ \ \ \ \ \ in\isanewline
\ \ \ \ \ \ \ \ \ \ \ \ List{\isachardot}{\kern0pt}filter\ {\isacharparenleft}{\kern0pt}fn\ v\ {\isacharequal}{\kern0pt}{\isachargreater}{\kern0pt}\ not\ {\isacharparenleft}{\kern0pt}List{\isachardot}{\kern0pt}exists\ {\isacharparenleft}{\kern0pt}fn\ d\ {\isacharequal}{\kern0pt}{\isachargreater}{\kern0pt}\ d\ {\isacharequal}{\kern0pt}\ v{\isacharparenright}{\kern0pt}\ declared{\isacharparenright}{\kern0pt}{\isacharparenright}{\kern0pt}\ used\isanewline
\ \ \ \ \ \ \ \ \ \ end\isanewline
\ \ \ \ \isanewline
\ \ \ \ \ \ \ \ {\isacharparenleft}{\kern0pt}{\isacharasterisk}{\kern0pt}\ {\isacharminus}{\kern0pt}{\isacharminus}{\kern0pt}{\isacharminus}{\kern0pt}{\isacharminus}{\kern0pt}{\isacharminus}{\kern0pt}{\isacharminus}{\kern0pt}{\isacharminus}{\kern0pt}{\isacharminus}{\kern0pt}{\isacharminus}{\kern0pt}{\isacharminus}{\kern0pt}{\isacharminus}{\kern0pt}{\isacharminus}{\kern0pt}{\isacharminus}{\kern0pt}{\isacharminus}{\kern0pt}{\isacharminus}{\kern0pt}{\isacharminus}{\kern0pt}{\isacharminus}{\kern0pt}{\isacharminus}{\kern0pt}\ {\isadigit{5}}{\isachardot}{\kern0pt}\ Agrégation\ globale\ {\isacharminus}{\kern0pt}{\isacharminus}{\kern0pt}{\isacharminus}{\kern0pt}{\isacharminus}{\kern0pt}{\isacharminus}{\kern0pt}{\isacharminus}{\kern0pt}{\isacharminus}{\kern0pt}{\isacharminus}{\kern0pt}{\isacharminus}{\kern0pt}{\isacharminus}{\kern0pt}{\isacharminus}{\kern0pt}{\isacharminus}{\kern0pt}{\isacharminus}{\kern0pt}{\isacharminus}{\kern0pt}{\isacharminus}{\kern0pt}{\isacharminus}{\kern0pt}{\isacharminus}{\kern0pt}{\isacharminus}{\kern0pt}{\isacharminus}{\kern0pt}{\isacharminus}{\kern0pt}{\isacharminus}{\kern0pt}{\isacharminus}{\kern0pt}{\isacharminus}{\kern0pt}\ {\isacharasterisk}{\kern0pt}{\isacharparenright}{\kern0pt}\isanewline
\ \ \ \ \ \ \ \ val\ missing\ {\isacharequal}{\kern0pt}\isanewline
\ \ \ \ \ \ \ \ \ \ List{\isachardot}{\kern0pt}concat\ {\isacharparenleft}{\kern0pt}map\ undeclared{\isacharunderscore}{\kern0pt}in{\isacharunderscore}{\kern0pt}thread\ threads{\isacharunderscore}{\kern0pt}decls{\isacharparenright}{\kern0pt}\isanewline
\ \ \ \ \ \ \ \ \ \ {\isacharbar}{\kern0pt}{\isachargreater}{\kern0pt}\ sort{\isacharunderscore}{\kern0pt}distinct\ string{\isacharunderscore}{\kern0pt}ord\isanewline
\ \ \ \ \ \ in\isanewline
\ \ \ \ \ \ \ \ if\ null\ missing\ then\isanewline
\ \ \ \ \ \ \ \ \ \ message\ false\ {\isachardoublequote}{\kern0pt}{\isasymcheckmark}\ All\ used\ variables\ are\ properly\ declared\ {\isacharparenleft}{\kern0pt}globals\ or\ locals{\isacharparenright}{\kern0pt}{\isachardoublequote}{\kern0pt}\isanewline
\ \ \ \ \ \ \ \ else\isanewline
\ \ \ \ \ \ \ \ \ \ {\isacharparenleft}{\kern0pt}message\ false\ {\isachardoublequote}{\kern0pt}{\isasymcrossmark}\ Undeclared\ variables\ {\isacharparenleft}{\kern0pt}neither\ global\ nor\ local{\isacharparenright}{\kern0pt}{\isacharcolon}{\kern0pt}{\isachardoublequote}{\kern0pt}{\isacharsemicolon}{\kern0pt}\isanewline
\ \ \ \ \ \ \ \ \ \ \ List{\isachardot}{\kern0pt}app\ {\isacharparenleft}{\kern0pt}fn\ v\ {\isacharequal}{\kern0pt}{\isachargreater}{\kern0pt}\ message\ false\ {\isacharparenleft}{\kern0pt}{\isachardoublequote}{\kern0pt}\ \ {\isachardoublequote}{\kern0pt}\ {\isacharcircum}{\kern0pt}\ v{\isacharparenright}{\kern0pt}{\isacharparenright}{\kern0pt}\ missing{\isacharparenright}{\kern0pt}\isanewline
\ \ \ \ \ \ end\isanewline
\ \ \ \ \isanewline
\ \ \ \ {\isacharparenleft}{\kern0pt}{\isacharasterisk}{\kern0pt}\ Appel\ de\ la\ vérification\ dans\ semantic{\isacharunderscore}{\kern0pt}check\ {\isacharminus}{\kern0pt}{\isacharminus}{\kern0pt}{\isacharminus}{\kern0pt}{\isacharminus}{\kern0pt}{\isacharminus}{\kern0pt}{\isacharminus}{\kern0pt}{\isacharminus}{\kern0pt}{\isacharminus}{\kern0pt}{\isacharminus}{\kern0pt}{\isacharminus}{\kern0pt}{\isacharminus}{\kern0pt}{\isacharminus}{\kern0pt}{\isacharminus}{\kern0pt}{\isacharminus}{\kern0pt}{\isacharminus}{\kern0pt}{\isacharminus}{\kern0pt}{\isacharminus}{\kern0pt}{\isacharminus}{\kern0pt}{\isacharminus}{\kern0pt}{\isacharminus}{\kern0pt}{\isacharminus}{\kern0pt}{\isacharminus}{\kern0pt}\ {\isacharasterisk}{\kern0pt}{\isacharparenright}{\kern0pt}\isanewline
\ \ \ \ val\ {\isacharunderscore}{\kern0pt}\ {\isacharequal}{\kern0pt}\ check{\isacharunderscore}{\kern0pt}var{\isacharunderscore}{\kern0pt}declarations\ {\isacharparenleft}{\kern0pt}{\isacharparenright}{\kern0pt}\isanewline
\ \ \ \ \isanewline
\ \ \ \ val\ {\isacharunderscore}{\kern0pt}\ {\isacharequal}{\kern0pt}\ message\ false\ {\isachardoublequote}{\kern0pt}{\isacharequal}{\kern0pt}{\isacharequal}{\kern0pt}{\isacharequal}{\kern0pt}\ SEMANTIC\ CHECK\ COMPLETED\ {\isacharequal}{\kern0pt}{\isacharequal}{\kern0pt}{\isacharequal}{\kern0pt}{\isachardoublequote}{\kern0pt}\isanewline
\ \ in\isanewline
\ \ \ \ thy\isanewline
\ \ end\isanewline
\isanewline
fun\ semantics\ {\isacharparenleft}{\kern0pt}absy{\isacharunderscore}{\kern0pt}cl{\isacharcolon}{\kern0pt}\ theory\ {\isacharminus}{\kern0pt}{\isachargreater}{\kern0pt}\ absy{\isacharparenright}{\kern0pt}\ {\isacharparenleft}{\kern0pt}thy\ {\isacharcolon}{\kern0pt}\ theory{\isacharparenright}{\kern0pt}\ {\isacharequal}{\kern0pt}\ \isanewline
\ \ let\isanewline
\ \ \ \ val\ absy{\isacharunderscore}{\kern0pt}data\ {\isacharequal}{\kern0pt}\ absy{\isacharunderscore}{\kern0pt}cl\ thy\isanewline
\ \ in\isanewline
\ \ \ \ semantic{\isacharunderscore}{\kern0pt}check\ absy{\isacharunderscore}{\kern0pt}data\ thy\isanewline
\ \ end\isanewline
\isanewline
fun\ sequence{\isacharunderscore}{\kern0pt}actions\ {\isacharbrackleft}{\kern0pt}{\isacharbrackright}{\kern0pt}\ {\isacharequal}{\kern0pt}\ SkipA\isanewline
\ \ {\isacharbar}{\kern0pt}\ sequence{\isacharunderscore}{\kern0pt}actions\ {\isacharbrackleft}{\kern0pt}a{\isacharbrackright}{\kern0pt}\ {\isacharequal}{\kern0pt}\ a\isanewline
\ \ {\isacharbar}{\kern0pt}\ sequence{\isacharunderscore}{\kern0pt}actions\ {\isacharparenleft}{\kern0pt}a{\isacharcolon}{\kern0pt}{\isacharcolon}{\kern0pt}rest{\isacharparenright}{\kern0pt}\ {\isacharequal}{\kern0pt}\ SeqA\ {\isacharparenleft}{\kern0pt}a{\isacharcomma}{\kern0pt}\ sequence{\isacharunderscore}{\kern0pt}actions\ rest{\isacharparenright}{\kern0pt}\isanewline
{\isacharparenleft}{\kern0pt}{\isacharasterisk}{\kern0pt}\isanewline
{\isacharparenleft}{\kern0pt}{\isacharasterisk}{\kern0pt}Fonction\ pour\ générer\ le\ code\ CSP\ complet\ {\isacharasterisk}{\kern0pt}{\isacharparenright}{\kern0pt}\isanewline
fun\ generate{\isacharunderscore}{\kern0pt}csp{\isacharunderscore}{\kern0pt}code\ {\isacharparenleft}{\kern0pt}absy{\isacharunderscore}{\kern0pt}data\ {\isacharcolon}{\kern0pt}\ absy{\isacharparenright}{\kern0pt}\ {\isacharparenleft}{\kern0pt}thy\ {\isacharcolon}{\kern0pt}\ theory{\isacharparenright}{\kern0pt}\ {\isacharcolon}{\kern0pt}\ string\ {\isacharequal}{\kern0pt}\isanewline
let\isanewline
\ \ \ \ val\ {\isacharbraceleft}{\kern0pt}system{\isacharcomma}{\kern0pt}\ globals{\isacharunderscore}{\kern0pt}decls{\isacharcomma}{\kern0pt}\ locks{\isacharunderscore}{\kern0pt}decls{\isacharcomma}{\kern0pt}\ threads{\isacharunderscore}{\kern0pt}decls{\isacharcomma}{\kern0pt}\ varstab{\isacharbraceright}{\kern0pt}\ {\isacharequal}{\kern0pt}\ absy{\isacharunderscore}{\kern0pt}data\isanewline
\ \ \ \ val\ system{\isacharunderscore}{\kern0pt}name\ {\isacharequal}{\kern0pt}\ Binding{\isachardot}{\kern0pt}name{\isacharunderscore}{\kern0pt}of\ system\isanewline
\isanewline
\ \ \ \ {\isacharparenleft}{\kern0pt}{\isacharasterisk}{\kern0pt}Génération\ des\ événements\ CSP\ {\isacharasterisk}{\kern0pt}{\isacharparenright}{\kern0pt}\isanewline
\ \ \ \ fun\ generate{\isacharunderscore}{\kern0pt}events\ {\isacharparenleft}{\kern0pt}{\isacharparenright}{\kern0pt}\ {\isacharequal}{\kern0pt}\isanewline
\ \ \ \ \ \ let\isanewline
\ \ \ \ \ \ \ \ val\ lock{\isacharunderscore}{\kern0pt}events\ {\isacharequal}{\kern0pt}\ map\ {\isacharparenleft}{\kern0pt}fn\ {\isacharparenleft}{\kern0pt}bdg{\isacharcomma}{\kern0pt}\ {\isacharunderscore}{\kern0pt}{\isacharcomma}{\kern0pt}\ {\isacharunderscore}{\kern0pt}{\isacharparenright}{\kern0pt}\ {\isacharequal}{\kern0pt}{\isachargreater}{\kern0pt}\ \isanewline
\ \ \ \ \ \ \ \ \ \ let\ val\ name\ {\isacharequal}{\kern0pt}\ Binding{\isachardot}{\kern0pt}name{\isacharunderscore}{\kern0pt}of\ bdg\isanewline
\ \ \ \ \ \ \ \ \ \ in\ {\isachardoublequote}{\kern0pt}channel\ lock{\isacharunderscore}{\kern0pt}{\isachardoublequote}{\kern0pt}\ {\isacharcircum}{\kern0pt}\ name\ {\isacharcircum}{\kern0pt}\ {\isachardoublequote}{\kern0pt}{\isacharcomma}{\kern0pt}\ unlock{\isacharunderscore}{\kern0pt}{\isachardoublequote}{\kern0pt}\ {\isacharcircum}{\kern0pt}\ name\ {\isacharcircum}{\kern0pt}\ {\isachardoublequote}{\kern0pt}\ {\isacharcolon}{\kern0pt}\ Int{\isachardoublequote}{\kern0pt}\isanewline
\ \ \ \ \ \ \ \ \ \ end{\isacharparenright}{\kern0pt}\ locks{\isacharunderscore}{\kern0pt}decls\isanewline
\isanewline
\ \ \ \ \ \ \ \ val\ global{\isacharunderscore}{\kern0pt}events\ {\isacharequal}{\kern0pt}\ map\ {\isacharparenleft}{\kern0pt}fn\ {\isacharparenleft}{\kern0pt}bdg{\isacharcomma}{\kern0pt}\ {\isacharunderscore}{\kern0pt}{\isacharcomma}{\kern0pt}\ {\isacharunderscore}{\kern0pt}{\isacharcomma}{\kern0pt}\ {\isacharunderscore}{\kern0pt}{\isacharparenright}{\kern0pt}\ {\isacharequal}{\kern0pt}{\isachargreater}{\kern0pt}\isanewline
\ \ \ \ \ \ \ \ \ \ let\ val\ name\ {\isacharequal}{\kern0pt}\ Binding{\isachardot}{\kern0pt}name{\isacharunderscore}{\kern0pt}of\ bdg\ \ \isanewline
\ \ \ \ \ \ \ \ \ \ in\ {\isachardoublequote}{\kern0pt}channel\ read{\isacharunderscore}{\kern0pt}{\isachardoublequote}{\kern0pt}\ {\isacharcircum}{\kern0pt}\ name\ {\isacharcircum}{\kern0pt}\ {\isachardoublequote}{\kern0pt}{\isacharcomma}{\kern0pt}\ update{\isacharunderscore}{\kern0pt}{\isachardoublequote}{\kern0pt}\ {\isacharcircum}{\kern0pt}\ name\ {\isacharcircum}{\kern0pt}\ {\isachardoublequote}{\kern0pt}\ {\isacharcolon}{\kern0pt}\ Int{\isachardoublequote}{\kern0pt}\isanewline
\ \ \ \ \ \ \ \ \ \ end{\isacharparenright}{\kern0pt}\ globals{\isacharunderscore}{\kern0pt}decls\isanewline
\ \ \ \ \ \ in\isanewline
\ \ \ \ \ \ \ \ String{\isachardot}{\kern0pt}concatWith\ {\isachardoublequote}{\kern0pt}{\isacharbackslash}{\kern0pt}n{\isachardoublequote}{\kern0pt}\ {\isacharparenleft}{\kern0pt}lock{\isacharunderscore}{\kern0pt}events\ {\isacharat}{\kern0pt}\ global{\isacharunderscore}{\kern0pt}events{\isacharparenright}{\kern0pt}\isanewline
\ \ \ \ \ \ end\isanewline
\isanewline
\ \ \ \ {\isacharparenleft}{\kern0pt}{\isacharasterisk}{\kern0pt}Conversion\ d{\isacharprime}{\kern0pt}une\ action\ en\ CSP\ {\isacharasterisk}{\kern0pt}{\isacharparenright}{\kern0pt}\isanewline
fun\ action{\isacharunderscore}{\kern0pt}list{\isacharunderscore}{\kern0pt}to{\isacharunderscore}{\kern0pt}csp\ {\isacharbrackleft}{\kern0pt}{\isacharbrackright}{\kern0pt}\ state\ thy\ {\isacharequal}{\kern0pt}\ {\isacharparenleft}{\kern0pt}{\isachardoublequote}{\kern0pt}SKIP{\isachardoublequote}{\kern0pt}{\isacharcomma}{\kern0pt}\ state{\isacharparenright}{\kern0pt}\isanewline
\ \ {\isacharbar}{\kern0pt}\ action{\isacharunderscore}{\kern0pt}list{\isacharunderscore}{\kern0pt}to{\isacharunderscore}{\kern0pt}csp\ {\isacharbrackleft}{\kern0pt}a{\isacharbrackright}{\kern0pt}\ state\ thy\ {\isacharequal}{\kern0pt}\ action{\isacharunderscore}{\kern0pt}to{\isacharunderscore}{\kern0pt}csp\ a\ state\ thy\isanewline
\ \ {\isacharbar}{\kern0pt}\ action{\isacharunderscore}{\kern0pt}list{\isacharunderscore}{\kern0pt}to{\isacharunderscore}{\kern0pt}csp\ {\isacharparenleft}{\kern0pt}a{\isacharcolon}{\kern0pt}{\isacharcolon}{\kern0pt}rest{\isacharparenright}{\kern0pt}\ state\ thy\ {\isacharequal}{\kern0pt}\isanewline
\ \ \ \ \ \ let\isanewline
\ \ \ \ \ \ \ \ val\ {\isacharparenleft}{\kern0pt}s{\isadigit{1}}{\isacharcomma}{\kern0pt}\ state{\isadigit{1}}{\isacharparenright}{\kern0pt}\ {\isacharequal}{\kern0pt}\ action{\isacharunderscore}{\kern0pt}to{\isacharunderscore}{\kern0pt}csp\ a\ state\ thy\isanewline
\ \ \ \ \ \ \ \ val\ {\isacharparenleft}{\kern0pt}s{\isadigit{2}}{\isacharcomma}{\kern0pt}\ state{\isadigit{2}}{\isacharparenright}{\kern0pt}\ {\isacharequal}{\kern0pt}\ action{\isacharunderscore}{\kern0pt}list{\isacharunderscore}{\kern0pt}to{\isacharunderscore}{\kern0pt}csp\ rest\ state{\isadigit{1}}\ thy\isanewline
\ \ \ \ \ \ in\isanewline
\ \ \ \ \ \ \ \ if\ {\isacharparenleft}{\kern0pt}a\ {\isacharequal}{\kern0pt}SkipA{\isacharparenright}{\kern0pt}\ then\ \isanewline
\ \ \ \ \ \ \ \ \ \ {\isacharparenleft}{\kern0pt}s{\isadigit{2}}{\isacharcomma}{\kern0pt}\ state{\isadigit{2}}{\isacharparenright}{\kern0pt}\isanewline
\ \ \ \ \ \ \ \ else\isanewline
\ \ \ \ \ \ \ \ \ \ {\isacharparenleft}{\kern0pt}s{\isadigit{1}}\ {\isacharcircum}{\kern0pt}\ {\isachardoublequote}{\kern0pt}{\isasymrightarrow}{\isachardoublequote}{\kern0pt}\ {\isacharcircum}{\kern0pt}\ s{\isadigit{2}}{\isacharcomma}{\kern0pt}\ state{\isadigit{2}}{\isacharparenright}{\kern0pt}\isanewline
\ \ \ \ \ \ end\isanewline
\isanewline
and\ action{\isacharunderscore}{\kern0pt}to{\isacharunderscore}{\kern0pt}csp\ action\ state\ thy\ {\isacharequal}{\kern0pt}\isanewline
\ \ case\ action\ of\isanewline
\ \ \ SkipA\ {\isacharequal}{\kern0pt}{\isachargreater}{\kern0pt}\ {\isacharparenleft}{\kern0pt}{\isachardoublequote}{\kern0pt}SKIP{\isachardoublequote}{\kern0pt}{\isacharcomma}{\kern0pt}\ state{\isacharparenright}{\kern0pt}\isanewline
\ \ {\isacharbar}{\kern0pt}\ LockA\ bdg\ {\isacharequal}{\kern0pt}{\isachargreater}{\kern0pt}\ {\isacharparenleft}{\kern0pt}{\isachardoublequote}{\kern0pt}lock{\isacharunderscore}{\kern0pt}{\isachardoublequote}{\kern0pt}\ {\isacharcircum}{\kern0pt}\ Binding{\isachardot}{\kern0pt}name{\isacharunderscore}{\kern0pt}of\ bdg\ {\isacharcircum}{\kern0pt}\ {\isachardoublequote}{\kern0pt}{\isacharbang}{\kern0pt}{\isadigit{0}}{\isachardoublequote}{\kern0pt}{\isacharcomma}{\kern0pt}\ state{\isacharparenright}{\kern0pt}\isanewline
\ \ {\isacharbar}{\kern0pt}\ UnlockA\ bdg\ {\isacharequal}{\kern0pt}{\isachargreater}{\kern0pt}\ {\isacharparenleft}{\kern0pt}{\isachardoublequote}{\kern0pt}unlock{\isacharunderscore}{\kern0pt}{\isachardoublequote}{\kern0pt}\ {\isacharcircum}{\kern0pt}\ Binding{\isachardot}{\kern0pt}name{\isacharunderscore}{\kern0pt}of\ bdg\ {\isacharcircum}{\kern0pt}\ {\isachardoublequote}{\kern0pt}{\isacharbang}{\kern0pt}{\isadigit{0}}{\isachardoublequote}{\kern0pt}{\isacharcomma}{\kern0pt}\ state{\isacharparenright}{\kern0pt}\isanewline
\ \ {\isacharbar}{\kern0pt}\ AssignA\ {\isacharparenleft}{\kern0pt}bdg{\isacharcomma}{\kern0pt}\ term{\isacharparenright}{\kern0pt}\ {\isacharequal}{\kern0pt}{\isachargreater}{\kern0pt}\isanewline
\ \ \ \ \ \ let\isanewline
\ \ \ \ \ \ \ \ val\ name\ {\isacharequal}{\kern0pt}\ Binding{\isachardot}{\kern0pt}name{\isacharunderscore}{\kern0pt}of\ bdg\isanewline
\ \ \ \ \ \ \ \ val\ typ\ {\isacharequal}{\kern0pt}\ fastype{\isacharunderscore}{\kern0pt}of\ term\isanewline
\ \ \ \ \ \ \ \ val\ v\ {\isacharequal}{\kern0pt}\ if\ typ\ {\isacharequal}{\kern0pt}\ {\isacharat}{\kern0pt}{\isacharbraceleft}{\kern0pt}typ\ bool{\isacharbraceright}{\kern0pt}\isanewline
\ \ \ \ \ \ \ \ \ \ \ \ \ \ \ \ then\ if\ eval{\isacharunderscore}{\kern0pt}term{\isacharunderscore}{\kern0pt}to{\isacharunderscore}{\kern0pt}bool\ thy\ term\ state\ then\ {\isadigit{1}}\ else\ {\isadigit{0}}\isanewline
\ \ \ \ \ \ \ \ \ \ \ \ \ \ \ \ else\ eval{\isacharunderscore}{\kern0pt}term{\isacharunderscore}{\kern0pt}to{\isacharunderscore}{\kern0pt}int\ thy\ term\ state\isanewline
\ \ \ \ \ \ \ \ val\ new{\isacharunderscore}{\kern0pt}state\ {\isacharequal}{\kern0pt}\ update{\isacharunderscore}{\kern0pt}state\ state\ name\ v\isanewline
\ \ \ \ \ \ in\isanewline
\ \ \ \ \ \ \ \ {\isacharparenleft}{\kern0pt}{\isachardoublequote}{\kern0pt}update{\isacharunderscore}{\kern0pt}{\isachardoublequote}{\kern0pt}\ {\isacharcircum}{\kern0pt}\ name\ {\isacharcircum}{\kern0pt}\ {\isachardoublequote}{\kern0pt}{\isacharbang}{\kern0pt}{\isachardoublequote}{\kern0pt}\ {\isacharcircum}{\kern0pt}\ Int{\isachardot}{\kern0pt}toString\ v{\isacharcomma}{\kern0pt}\ new{\isacharunderscore}{\kern0pt}state{\isacharparenright}{\kern0pt}\isanewline
\ \ \ \ \ \ end\isanewline
\ \ {\isacharbar}{\kern0pt}\ SendA\ {\isacharparenleft}{\kern0pt}bdg{\isacharcomma}{\kern0pt}\ term{\isacharparenright}{\kern0pt}\ {\isacharequal}{\kern0pt}{\isachargreater}{\kern0pt}\isanewline
\ \ \ \ \ \ let\isanewline
\ \ \ \ \ \ \ \ val\ name\ {\isacharequal}{\kern0pt}\ Binding{\isachardot}{\kern0pt}name{\isacharunderscore}{\kern0pt}of\ bdg\isanewline
\ \ \ \ \ \ \ \ val\ typ\ {\isacharequal}{\kern0pt}\ fastype{\isacharunderscore}{\kern0pt}of\ term\isanewline
\ \ \ \ \ \ \ \ val\ v\ {\isacharequal}{\kern0pt}\ if\ typ\ {\isacharequal}{\kern0pt}\ {\isacharat}{\kern0pt}{\isacharbraceleft}{\kern0pt}typ\ bool{\isacharbraceright}{\kern0pt}\isanewline
\ \ \ \ \ \ \ \ \ \ \ \ \ \ \ \ then\ if\ eval{\isacharunderscore}{\kern0pt}term{\isacharunderscore}{\kern0pt}to{\isacharunderscore}{\kern0pt}bool\ thy\ term\ state\ then\ {\isadigit{1}}\ else\ {\isadigit{0}}\isanewline
\ \ \ \ \ \ \ \ \ \ \ \ \ \ \ \ else\ eval{\isacharunderscore}{\kern0pt}term{\isacharunderscore}{\kern0pt}to{\isacharunderscore}{\kern0pt}int\ thy\ term\ state\isanewline
\ \ \ \ \ \ \ \ val\ new{\isacharunderscore}{\kern0pt}state\ {\isacharequal}{\kern0pt}\ update{\isacharunderscore}{\kern0pt}state\ state\ name\ v\isanewline
\ \ \ \ \ \ in\isanewline
\ \ \ \ \ \ \ \ {\isacharparenleft}{\kern0pt}{\isachardoublequote}{\kern0pt}update{\isacharunderscore}{\kern0pt}{\isachardoublequote}{\kern0pt}\ {\isacharcircum}{\kern0pt}\ name\ {\isacharcircum}{\kern0pt}\ {\isachardoublequote}{\kern0pt}{\isacharbang}{\kern0pt}{\isachardoublequote}{\kern0pt}\ {\isacharcircum}{\kern0pt}\ Int{\isachardot}{\kern0pt}toString\ v{\isacharcomma}{\kern0pt}\ new{\isacharunderscore}{\kern0pt}state{\isacharparenright}{\kern0pt}\isanewline
\ \ \ \ \ \ end\isanewline
\ \ {\isacharbar}{\kern0pt}\ ReceiveA\ {\isacharparenleft}{\kern0pt}local{\isacharunderscore}{\kern0pt}bdg{\isacharcomma}{\kern0pt}\ global{\isacharunderscore}{\kern0pt}bdg{\isacharparenright}{\kern0pt}\ {\isacharequal}{\kern0pt}{\isachargreater}{\kern0pt}\isanewline
\ \ \ \ \ \ let\isanewline
\ \ \ \ \ \ \ \ val\ global{\isacharunderscore}{\kern0pt}var{\isacharunderscore}{\kern0pt}name\ {\isacharequal}{\kern0pt}\ Binding{\isachardot}{\kern0pt}name{\isacharunderscore}{\kern0pt}of\ global{\isacharunderscore}{\kern0pt}bdg\isanewline
\ \ \ \ \ \ \ \ val\ v\ {\isacharequal}{\kern0pt}\ safe{\isacharunderscore}{\kern0pt}lookup\ global{\isacharunderscore}{\kern0pt}var{\isacharunderscore}{\kern0pt}name\ state\isanewline
\ \ \ \ \ \ \ \ val\ local{\isacharunderscore}{\kern0pt}var{\isacharunderscore}{\kern0pt}name\ {\isacharequal}{\kern0pt}\ Binding{\isachardot}{\kern0pt}name{\isacharunderscore}{\kern0pt}of\ local{\isacharunderscore}{\kern0pt}bdg\isanewline
\ \ \ \ \ \ \ \ val\ new{\isacharunderscore}{\kern0pt}state\ {\isacharequal}{\kern0pt}\ update{\isacharunderscore}{\kern0pt}state\ state\ local{\isacharunderscore}{\kern0pt}var{\isacharunderscore}{\kern0pt}name\ v\isanewline
\ \ \ \ \ \ in\isanewline
\ \ \ \ \ \ \ \ {\isacharparenleft}{\kern0pt}{\isachardoublequote}{\kern0pt}read{\isacharunderscore}{\kern0pt}{\isachardoublequote}{\kern0pt}\ {\isacharcircum}{\kern0pt}\ global{\isacharunderscore}{\kern0pt}var{\isacharunderscore}{\kern0pt}name\ {\isacharcircum}{\kern0pt}\ {\isachardoublequote}{\kern0pt}{\isacharbang}{\kern0pt}{\isachardoublequote}{\kern0pt}\ {\isacharcircum}{\kern0pt}\ Int{\isachardot}{\kern0pt}toString\ v{\isacharcomma}{\kern0pt}\ new{\isacharunderscore}{\kern0pt}state{\isacharparenright}{\kern0pt}\isanewline
\ \ \ \ \ \ end\isanewline
\ \ {\isacharbar}{\kern0pt}\ IfelseA\ {\isacharparenleft}{\kern0pt}{\isacharparenleft}{\kern0pt}cond{\isacharcomma}{\kern0pt}\ a{\isadigit{1}}{\isacharparenright}{\kern0pt}{\isacharcomma}{\kern0pt}\ a{\isadigit{2}}{\isacharparenright}{\kern0pt}\ {\isacharequal}{\kern0pt}{\isachargreater}{\kern0pt}\isanewline
\ \ \ \ \ \ let\isanewline
\ \ \ \ \ \ \ \ val\ b\ {\isacharequal}{\kern0pt}\ eval{\isacharunderscore}{\kern0pt}term{\isacharunderscore}{\kern0pt}to{\isacharunderscore}{\kern0pt}bool\ thy\ cond\ state\isanewline
\ \ \ \ \ \ \ \ val\ val{\isacharunderscore}{\kern0pt}str\ {\isacharequal}{\kern0pt}\ if\ b\ then\ {\isachardoublequote}{\kern0pt}true{\isachardoublequote}{\kern0pt}\ else\ {\isachardoublequote}{\kern0pt}false{\isachardoublequote}{\kern0pt}\isanewline
\ \ \ \ \ \ \ \ val\ {\isacharparenleft}{\kern0pt}then{\isacharunderscore}{\kern0pt}code{\isacharcomma}{\kern0pt}\ {\isacharunderscore}{\kern0pt}{\isacharparenright}{\kern0pt}\ {\isacharequal}{\kern0pt}\ action{\isacharunderscore}{\kern0pt}list{\isacharunderscore}{\kern0pt}to{\isacharunderscore}{\kern0pt}csp\ a{\isadigit{1}}\ state\ thy\isanewline
\ \ \ \ \ \ \ \ val\ {\isacharparenleft}{\kern0pt}else{\isacharunderscore}{\kern0pt}code{\isacharcomma}{\kern0pt}\ {\isacharunderscore}{\kern0pt}{\isacharparenright}{\kern0pt}\ {\isacharequal}{\kern0pt}\ action{\isacharunderscore}{\kern0pt}list{\isacharunderscore}{\kern0pt}to{\isacharunderscore}{\kern0pt}csp\ a{\isadigit{2}}\ state\ thy\isanewline
\ \ \ \ \ \ in\isanewline
\ \ \ \ \ \ \ \ {\isacharparenleft}{\kern0pt}{\isachardoublequote}{\kern0pt}{\isacharbraceleft}{\kern0pt}\ {\isacharparenleft}{\kern0pt}{\isachardoublequote}{\kern0pt}\ {\isacharcircum}{\kern0pt}\ val{\isacharunderscore}{\kern0pt}str\ {\isacharcircum}{\kern0pt}\ {\isachardoublequote}{\kern0pt}\ {\isacharampersand}{\kern0pt}\ {\isachardoublequote}{\kern0pt}\ {\isacharcircum}{\kern0pt}\ then{\isacharunderscore}{\kern0pt}code\ {\isacharcircum}{\kern0pt}\ {\isachardoublequote}{\kern0pt}{\isacharparenright}{\kern0pt}\ {\isacharbrackleft}{\kern0pt}{\isacharbrackright}{\kern0pt}\ {\isacharparenleft}{\kern0pt}{\isasymnot}{\isachardoublequote}{\kern0pt}\ {\isacharcircum}{\kern0pt}\ val{\isacharunderscore}{\kern0pt}str\ {\isacharcircum}{\kern0pt}\ {\isachardoublequote}{\kern0pt}\ {\isacharampersand}{\kern0pt}\ {\isachardoublequote}{\kern0pt}\ {\isacharcircum}{\kern0pt}\ else{\isacharunderscore}{\kern0pt}code\ {\isacharcircum}{\kern0pt}\ {\isachardoublequote}{\kern0pt}{\isacharparenright}{\kern0pt}\ {\isacharbraceright}{\kern0pt}{\isachardoublequote}{\kern0pt}{\isacharcomma}{\kern0pt}\ state{\isacharparenright}{\kern0pt}\isanewline
\ \ \ \ \ \ end\isanewline
\ \ {\isacharbar}{\kern0pt}\ WhileA\ {\isacharparenleft}{\kern0pt}cond{\isacharcomma}{\kern0pt}\ body{\isacharparenright}{\kern0pt}\ {\isacharequal}{\kern0pt}{\isachargreater}{\kern0pt}\isanewline
\ \ \ \ \ \ let\isanewline
\ \ \ \ \ \ \ \ val\ val{\isacharunderscore}{\kern0pt}str\ {\isacharequal}{\kern0pt}\ if\ eval{\isacharunderscore}{\kern0pt}term{\isacharunderscore}{\kern0pt}to{\isacharunderscore}{\kern0pt}bool\ thy\ cond\ state\ then\ {\isachardoublequote}{\kern0pt}true{\isachardoublequote}{\kern0pt}\ else\ {\isachardoublequote}{\kern0pt}false{\isachardoublequote}{\kern0pt}\isanewline
\ \ \ \ \ \ \ \ val\ {\isacharparenleft}{\kern0pt}body{\isacharunderscore}{\kern0pt}code{\isacharcomma}{\kern0pt}\ {\isacharunderscore}{\kern0pt}{\isacharparenright}{\kern0pt}\ {\isacharequal}{\kern0pt}\ action{\isacharunderscore}{\kern0pt}list{\isacharunderscore}{\kern0pt}to{\isacharunderscore}{\kern0pt}csp\ body\ state\ thy\isanewline
\ \ \ \ \ \ in\isanewline
\ \ \ \ \ \ \ \ {\isacharparenleft}{\kern0pt}{\isachardoublequote}{\kern0pt}{\isacharbraceleft}{\kern0pt}\ {\isacharparenleft}{\kern0pt}{\isachardoublequote}{\kern0pt}\ {\isacharcircum}{\kern0pt}\ val{\isacharunderscore}{\kern0pt}str\ {\isacharcircum}{\kern0pt}\ {\isachardoublequote}{\kern0pt}\ {\isacharampersand}{\kern0pt}\ \ µ\ X{\isachardot}{\kern0pt}{\isachardoublequote}{\kern0pt}\ {\isacharcircum}{\kern0pt}\ body{\isacharunderscore}{\kern0pt}code\ {\isacharcircum}{\kern0pt}\ {\isachardoublequote}{\kern0pt}\ {\isacharminus}{\kern0pt}{\isachargreater}{\kern0pt}\ X{\isacharparenright}{\kern0pt}\ {\isacharbrackleft}{\kern0pt}{\isacharbrackright}{\kern0pt}\ {\isacharparenleft}{\kern0pt}{\isasymnot}{\isachardoublequote}{\kern0pt}\ {\isacharcircum}{\kern0pt}\ val{\isacharunderscore}{\kern0pt}str\ {\isacharcircum}{\kern0pt}{\isachardoublequote}{\kern0pt}\ {\isacharampersand}{\kern0pt}\ {\isacharparenleft}{\kern0pt}SKIP{\isacharparenright}{\kern0pt}\ {\isacharparenright}{\kern0pt}\ {\isacharbraceright}{\kern0pt}{\isachardoublequote}{\kern0pt}{\isacharcomma}{\kern0pt}\ state{\isacharparenright}{\kern0pt}\isanewline
\ \ \ \ \ \ end\isanewline
\isanewline
\isanewline
\ \ \ \ {\isacharparenleft}{\kern0pt}{\isacharasterisk}{\kern0pt}\ Conversion\ d’un\ thread\ vers\ CSP\ {\isacharasterisk}{\kern0pt}{\isacharparenright}{\kern0pt}\isanewline
\ \ \ \ \ \ \ \ fun\ thread{\isacharunderscore}{\kern0pt}to{\isacharunderscore}{\kern0pt}csp\ globals{\isacharunderscore}{\kern0pt}decls\ {\isacharparenleft}{\kern0pt}thread\ {\isacharcolon}{\kern0pt}\ thread{\isacharunderscore}{\kern0pt}absy{\isacharparenright}{\kern0pt}\ thy\ {\isacharcolon}{\kern0pt}\ string\ {\isacharequal}{\kern0pt}\isanewline
\ \ \ \ \ \ \ \ \ \ let\isanewline
\ \ \ \ \ \ \ \ \ \ \ \ val\ {\isacharbraceleft}{\kern0pt}nom{\isacharunderscore}{\kern0pt}thread{\isacharcomma}{\kern0pt}\ actions{\isacharcomma}{\kern0pt}\ locals{\isacharunderscore}{\kern0pt}decl{\isacharcomma}{\kern0pt}\ {\isachardot}{\kern0pt}{\isachardot}{\kern0pt}{\isachardot}{\kern0pt}{\isacharbraceright}{\kern0pt}\ {\isacharequal}{\kern0pt}\ thread\isanewline
\ \ \ \ \ \ \ \ \ \ \ \ val\ thread{\isacharunderscore}{\kern0pt}name\ {\isacharequal}{\kern0pt}\ Binding{\isachardot}{\kern0pt}name{\isacharunderscore}{\kern0pt}of\ nom{\isacharunderscore}{\kern0pt}thread\isanewline
\ \ \ \ \ \ \ \ \isanewline
\ \ \ \ \ \ \ \ \ \ \ \ val\ initial{\isacharunderscore}{\kern0pt}state\ {\isacharequal}{\kern0pt}\ build{\isacharunderscore}{\kern0pt}initial{\isacharunderscore}{\kern0pt}state{\isacharunderscore}{\kern0pt}for{\isacharunderscore}{\kern0pt}thread\ globals{\isacharunderscore}{\kern0pt}decls\ locals{\isacharunderscore}{\kern0pt}decl\ thy\isanewline
\ \ \ \ \ \ \ \ \ \ \ \ val\ {\isacharparenleft}{\kern0pt}code{\isacharcomma}{\kern0pt}\ {\isacharunderscore}{\kern0pt}{\isacharparenright}{\kern0pt}\ {\isacharequal}{\kern0pt}\ action{\isacharunderscore}{\kern0pt}list{\isacharunderscore}{\kern0pt}to{\isacharunderscore}{\kern0pt}csp\ actions\ initial{\isacharunderscore}{\kern0pt}state\ thy\isanewline
\ \ \ \ \ \ \ \ \ \ in\isanewline
\ \ \ \ \ \ \ \ \ \ \ \ thread{\isacharunderscore}{\kern0pt}name\ {\isacharcircum}{\kern0pt}\ {\isachardoublequote}{\kern0pt}\ {\isacharequal}{\kern0pt}\ {\isachardoublequote}{\kern0pt}\ {\isacharcircum}{\kern0pt}\ code\ \isanewline
\ \ \ \ \ \ \ \ \ \ end\isanewline
\ \ \ \ \ \ \ \ \ \ \isanewline
\ \ \ \ fun\ generate{\isacharunderscore}{\kern0pt}threads\ globals{\isacharunderscore}{\kern0pt}decls\ threads{\isacharunderscore}{\kern0pt}decls\ thy\ {\isacharequal}{\kern0pt}\isanewline
\ \ \ \ \ \ \ \ String{\isachardot}{\kern0pt}concatWith\ {\isachardoublequote}{\kern0pt}{\isacharbackslash}{\kern0pt}n{\isachardoublequote}{\kern0pt}\ {\isacharparenleft}{\kern0pt}map\ {\isacharparenleft}{\kern0pt}fn\ t\ {\isacharequal}{\kern0pt}{\isachargreater}{\kern0pt}\ thread{\isacharunderscore}{\kern0pt}to{\isacharunderscore}{\kern0pt}csp\ globals{\isacharunderscore}{\kern0pt}decls\ t\ thy{\isacharparenright}{\kern0pt}\ threads{\isacharunderscore}{\kern0pt}decls{\isacharparenright}{\kern0pt}\isanewline
\isanewline
\ \ \ \ {\isacharparenleft}{\kern0pt}{\isacharasterisk}{\kern0pt}Génération\ du\ système\ global\ {\isacharasterisk}{\kern0pt}{\isacharparenright}{\kern0pt}\isanewline
\ \ \ \ fun\ generate{\isacharunderscore}{\kern0pt}system\ {\isacharparenleft}{\kern0pt}{\isacharparenright}{\kern0pt}\ {\isacharequal}{\kern0pt}\isanewline
\ \ \ \ \ \ let\isanewline
\ \ \ \ \ \ \ \ val\ thread{\isacharunderscore}{\kern0pt}names\ {\isacharequal}{\kern0pt}\ map\ {\isacharparenleft}{\kern0pt}fn\ t\ {\isacharequal}{\kern0pt}{\isachargreater}{\kern0pt}\ Binding{\isachardot}{\kern0pt}name{\isacharunderscore}{\kern0pt}of\ {\isacharparenleft}{\kern0pt}{\isacharhash}{\kern0pt}nom{\isacharunderscore}{\kern0pt}thread\ t{\isacharparenright}{\kern0pt}{\isacharparenright}{\kern0pt}\ threads{\isacharunderscore}{\kern0pt}decls\isanewline
\ \ \ \ \ \ in\isanewline
\ \ \ \ \ \ \ \ system{\isacharunderscore}{\kern0pt}name\ {\isacharcircum}{\kern0pt}\ {\isachardoublequote}{\kern0pt}\ {\isacharequal}{\kern0pt}\ {\isachardoublequote}{\kern0pt}\ {\isacharcircum}{\kern0pt}\ String{\isachardot}{\kern0pt}concatWith\ {\isachardoublequote}{\kern0pt}\ {\isacharbar}{\kern0pt}{\isacharbar}{\kern0pt}{\isacharbar}{\kern0pt}\ {\isachardoublequote}{\kern0pt}\ thread{\isacharunderscore}{\kern0pt}names\isanewline
\ \ \ \ \ \ end\isanewline
\isanewline
\ \ in\isanewline
\ \ \ \ String{\isachardot}{\kern0pt}concatWith\ {\isachardoublequote}{\kern0pt}{\isacharbackslash}{\kern0pt}n{\isacharbackslash}{\kern0pt}n{\isachardoublequote}{\kern0pt}\ {\isacharbrackleft}{\kern0pt}\isanewline
\ \ \ \ \ \ {\isachardoublequote}{\kern0pt}{\isacharminus}{\kern0pt}{\isacharminus}{\kern0pt}\ Channels{\isachardoublequote}{\kern0pt}{\isacharcomma}{\kern0pt}\isanewline
\ \ \ \ \ \ generate{\isacharunderscore}{\kern0pt}events\ {\isacharparenleft}{\kern0pt}{\isacharparenright}{\kern0pt}{\isacharcomma}{\kern0pt}\isanewline
\ \ \ \ \ \ {\isachardoublequote}{\kern0pt}{\isachardoublequote}{\kern0pt}{\isacharcomma}{\kern0pt}\isanewline
\ \ \ \ \ \ {\isachardoublequote}{\kern0pt}{\isacharminus}{\kern0pt}{\isacharminus}{\kern0pt}\ Threads{\isachardoublequote}{\kern0pt}{\isacharcomma}{\kern0pt}\isanewline
\ \ \ \ \ \ generate{\isacharunderscore}{\kern0pt}threads\ globals{\isacharunderscore}{\kern0pt}decls\ threads{\isacharunderscore}{\kern0pt}decls\ thy{\isacharcomma}{\kern0pt}\isanewline
\ \ \ \ \ \ {\isachardoublequote}{\kern0pt}{\isachardoublequote}{\kern0pt}{\isacharcomma}{\kern0pt}\isanewline
\ \ \ \ \ \ {\isachardoublequote}{\kern0pt}{\isacharminus}{\kern0pt}{\isacharminus}{\kern0pt}\ System{\isachardoublequote}{\kern0pt}{\isacharcomma}{\kern0pt}\isanewline
\ \ \ \ \ \ generate{\isacharunderscore}{\kern0pt}system\ {\isacharparenleft}{\kern0pt}{\isacharparenright}{\kern0pt}\isanewline
\ \ \ \ {\isacharbrackright}{\kern0pt}\isanewline
\ \ end\isanewline
{\isacharasterisk}{\kern0pt}{\isacharparenright}{\kern0pt}\isanewline
\isanewline
{\isacartoucheclose}%
\endisatagML
{\isafoldML}%
%
\isadelimML
%
\endisadelimML
%
\isadelimdocument
%
\endisadelimdocument
%
\isatagdocument
%
\isamarkupsection{Semantic Encoding%
}
\isamarkuptrue%
%
\endisatagdocument
{\isafolddocument}%
%
\isadelimdocument
%
\endisadelimdocument
\isakeywordONE{definition}\isamarkupfalse%
\ SKIP{\isacharunderscore}{\kern0pt}process\ {\isacharcolon}{\kern0pt}{\isacharcolon}{\kern0pt}\ {\isachardoublequoteopen}evs\ process{\isachardoublequoteclose}\isanewline
\ \ \isakeywordTWO{where}\ {\isachardoublequoteopen}SKIP{\isacharunderscore}{\kern0pt}process\ {\isasymequiv}\ Skip{\isachardoublequoteclose}\isanewline
%
\isadelimML
\isanewline
%
\endisadelimML
%
\isatagML
\isakeywordONE{ML}\isamarkupfalse%
{\isacartoucheopen}\isanewline
{\isacharparenleft}{\kern0pt}{\isacharasterisk}{\kern0pt}\ Convertir\ une\ liste\ de\ a{\isacharunderscore}{\kern0pt}term\ en\ une\ seule\ commande\ de\ type\ com\ {\isacharasterisk}{\kern0pt}{\isacharparenright}{\kern0pt}\isanewline
fun\ sequence{\isacharunderscore}{\kern0pt}actions\ {\isacharbrackleft}{\kern0pt}{\isacharbrackright}{\kern0pt}\ {\isacharequal}{\kern0pt}\ SkipA\isanewline
\ \ {\isacharbar}{\kern0pt}\ sequence{\isacharunderscore}{\kern0pt}actions\ {\isacharbrackleft}{\kern0pt}a{\isacharbrackright}{\kern0pt}\ {\isacharequal}{\kern0pt}\ a\isanewline
\ \ {\isacharbar}{\kern0pt}\ sequence{\isacharunderscore}{\kern0pt}actions\ {\isacharparenleft}{\kern0pt}a\ {\isacharcolon}{\kern0pt}{\isacharcolon}{\kern0pt}\ rest{\isacharparenright}{\kern0pt}\ {\isacharequal}{\kern0pt}\ SeqA\ {\isacharparenleft}{\kern0pt}a{\isacharcomma}{\kern0pt}\ sequence{\isacharunderscore}{\kern0pt}actions\ rest{\isacharparenright}{\kern0pt}\isanewline
{\isacharparenleft}{\kern0pt}{\isacharasterisk}{\kern0pt}\isanewline
fun\ subst{\isacharunderscore}{\kern0pt}with{\isacharunderscore}{\kern0pt}fresh{\isacharunderscore}{\kern0pt}env\ {\isacharparenleft}{\kern0pt}env{\isacharunderscore}{\kern0pt}type\ {\isacharcolon}{\kern0pt}\ typ{\isacharparenright}{\kern0pt}\ {\isacharparenleft}{\kern0pt}t\ {\isacharcolon}{\kern0pt}\ term{\isacharparenright}{\kern0pt}\ {\isacharcolon}{\kern0pt}\ term\ {\isacharasterisk}{\kern0pt}\ term\ {\isacharequal}{\kern0pt}\isanewline
\ \ let\isanewline
\ \ \ \ {\isacharparenleft}{\kern0pt}{\isacharasterisk}{\kern0pt}\ collect\ free\ variable\ names\ {\isacharasterisk}{\kern0pt}{\isacharparenright}{\kern0pt}\isanewline
\ \ \ \ val\ frees\ {\isacharequal}{\kern0pt}\ Term{\isachardot}{\kern0pt}add{\isacharunderscore}{\kern0pt}free{\isacharunderscore}{\kern0pt}names\ t\ {\isacharbrackleft}{\kern0pt}{\isacharbrackright}{\kern0pt}\isanewline
\isanewline
\ \ \ \ {\isacharparenleft}{\kern0pt}{\isacharasterisk}{\kern0pt}\ generate\ a\ fresh\ variable\ name\ {\isachardoublequote}{\kern0pt}{\isasymsigma}{\isachardoublequote}{\kern0pt}{\isacharcomma}{\kern0pt}\ adding\ primes\ if\ necessary\ {\isacharasterisk}{\kern0pt}{\isacharparenright}{\kern0pt}\isanewline
\ \ \ \ fun\ fresh{\isacharunderscore}{\kern0pt}name\ base\ {\isacharbrackleft}{\kern0pt}{\isacharbrackright}{\kern0pt}\ {\isacharequal}{\kern0pt}\ base\isanewline
\ \ \ \ \ \ {\isacharbar}{\kern0pt}\ fresh{\isacharunderscore}{\kern0pt}name\ base\ used\ {\isacharequal}{\kern0pt}\isanewline
\ \ \ \ \ \ \ \ \ \ let\isanewline
\ \ \ \ \ \ \ \ \ \ \ \ fun\ bump\ n\ {\isacharequal}{\kern0pt}\isanewline
\ \ \ \ \ \ \ \ \ \ \ \ \ \ let\ val\ candidate\ {\isacharequal}{\kern0pt}\ base\ {\isacharcircum}{\kern0pt}\ replicate{\isacharunderscore}{\kern0pt}string\ n\ {\isachardoublequote}{\kern0pt}{\isacharprime}{\kern0pt}{\isachardoublequote}{\kern0pt}\ in\isanewline
\ \ \ \ \ \ \ \ \ \ \ \ \ \ \ \ if\ member\ {\isacharparenleft}{\kern0pt}op\ {\isacharequal}{\kern0pt}{\isacharparenright}{\kern0pt}\ used\ candidate\ then\ bump\ {\isacharparenleft}{\kern0pt}n\ {\isacharplus}{\kern0pt}\ {\isadigit{1}}{\isacharparenright}{\kern0pt}\ else\ candidate\isanewline
\ \ \ \ \ \ \ \ \ \ \ \ \ \ end\isanewline
\ \ \ \ \ \ \ \ \ \ in\ bump\ {\isadigit{1}}\ end\isanewline
\isanewline
\ \ \ \ val\ sigma{\isacharunderscore}{\kern0pt}name\ {\isacharequal}{\kern0pt}\ fresh{\isacharunderscore}{\kern0pt}name\ {\isachardoublequote}{\kern0pt}{\isasymsigma}{\isachardoublequote}{\kern0pt}\ frees\isanewline
\ \ \ \ val\ sigma\ {\isacharequal}{\kern0pt}\ Free\ {\isacharparenleft}{\kern0pt}sigma{\isacharunderscore}{\kern0pt}name{\isacharcomma}{\kern0pt}\ env{\isacharunderscore}{\kern0pt}type{\isacharparenright}{\kern0pt}\isanewline
\isanewline
\ \ \ \ {\isacharparenleft}{\kern0pt}{\isacharasterisk}{\kern0pt}\ perform\ substitution{\isacharcolon}{\kern0pt}\ Free{\isacharparenleft}{\kern0pt}x{\isacharparenright}{\kern0pt}\ becomes\ sigma\ {\isachardollar}{\kern0pt}\ {\isachardoublequote}{\kern0pt}x{\isachardoublequote}{\kern0pt}\ {\isacharasterisk}{\kern0pt}{\isacharparenright}{\kern0pt}\isanewline
\ \ \ \ fun\ subst\ t\ {\isacharequal}{\kern0pt}\isanewline
\ \ \ \ \ \ case\ t\ of\isanewline
\ \ \ \ \ \ \ \ Free\ {\isacharparenleft}{\kern0pt}name{\isacharcomma}{\kern0pt}\ {\isacharunderscore}{\kern0pt}{\isacharparenright}{\kern0pt}\ {\isacharequal}{\kern0pt}{\isachargreater}{\kern0pt}\ sigma\ {\isachardollar}{\kern0pt}\ HOLogic{\isachardot}{\kern0pt}mk{\isacharunderscore}{\kern0pt}string\ name\isanewline
\ \ \ \ \ \ {\isacharbar}{\kern0pt}\ Const\ {\isacharunderscore}{\kern0pt}\ {\isacharequal}{\kern0pt}{\isachargreater}{\kern0pt}\ t\isanewline
\ \ \ \ \ \ {\isacharbar}{\kern0pt}\ Bound\ {\isacharunderscore}{\kern0pt}\ {\isacharequal}{\kern0pt}{\isachargreater}{\kern0pt}\ t\isanewline
\ \ \ \ \ \ {\isacharbar}{\kern0pt}\ Var\ {\isacharunderscore}{\kern0pt}\ {\isacharequal}{\kern0pt}{\isachargreater}{\kern0pt}\ t\isanewline
\ \ \ \ \ \ {\isacharbar}{\kern0pt}\ Abs\ {\isacharparenleft}{\kern0pt}n{\isacharcomma}{\kern0pt}\ T{\isacharcomma}{\kern0pt}\ body{\isacharparenright}{\kern0pt}\ {\isacharequal}{\kern0pt}{\isachargreater}{\kern0pt}\ Abs\ {\isacharparenleft}{\kern0pt}n{\isacharcomma}{\kern0pt}\ T{\isacharcomma}{\kern0pt}\ subst\ body{\isacharparenright}{\kern0pt}\isanewline
\ \ \ \ \ \ {\isacharbar}{\kern0pt}\ t{\isadigit{1}}\ {\isachardollar}{\kern0pt}\ t{\isadigit{2}}\ {\isacharequal}{\kern0pt}{\isachargreater}{\kern0pt}\ subst\ t{\isadigit{1}}\ {\isachardollar}{\kern0pt}\ subst\ t{\isadigit{2}}\isanewline
\isanewline
\ \ \ \ val\ t{\isacharunderscore}{\kern0pt}subst\ {\isacharequal}{\kern0pt}\ subst\ t\isanewline
\ \ in\isanewline
\ \ \ \ {\isacharparenleft}{\kern0pt}sigma{\isacharcomma}{\kern0pt}\ t{\isacharunderscore}{\kern0pt}subst{\isacharparenright}{\kern0pt}\isanewline
\ \ end\isanewline
{\isacharasterisk}{\kern0pt}{\isacharparenright}{\kern0pt}\isanewline
\isanewline
val\ sigma\ {\isacharcolon}{\kern0pt}\ term\ {\isacharequal}{\kern0pt}\ Free\ {\isacharparenleft}{\kern0pt}{\isachardoublequote}{\kern0pt}{\isasymsigma}{\isachardoublequote}{\kern0pt}{\isacharcomma}{\kern0pt}\ {\isacharat}{\kern0pt}{\isacharbraceleft}{\kern0pt}typ\ {\isachardoublequote}{\kern0pt}string\ {\isasymRightarrow}\ int{\isachardoublequote}{\kern0pt}{\isacharbraceright}{\kern0pt}{\isacharparenright}{\kern0pt}\ \ \ {\isacharparenleft}{\kern0pt}{\isacharasterisk}{\kern0pt}\ changer\ le\ type\ ici\ si\ besoin\ {\isacharasterisk}{\kern0pt}{\isacharparenright}{\kern0pt}\isanewline
\isanewline
\isanewline
fun\ subst{\isacharunderscore}{\kern0pt}with{\isacharunderscore}{\kern0pt}sigma\ {\isacharparenleft}{\kern0pt}t\ {\isacharcolon}{\kern0pt}\ term{\isacharparenright}{\kern0pt}\ {\isacharcolon}{\kern0pt}\ term\ {\isacharequal}{\kern0pt}\isanewline
\ \ let\isanewline
\ \ \ \ fun\ subst\ tm\ {\isacharequal}{\kern0pt}\isanewline
\ \ \ \ \ \ case\ tm\ of\isanewline
\ \ \ \ \ \ \ \ Free\ {\isacharparenleft}{\kern0pt}name{\isacharcomma}{\kern0pt}\ {\isacharunderscore}{\kern0pt}{\isacharparenright}{\kern0pt}\ \ \ \ \ \ \ \ \ \ \ {\isacharequal}{\kern0pt}{\isachargreater}{\kern0pt}\ sigma\ {\isachardollar}{\kern0pt}\ HOLogic{\isachardot}{\kern0pt}mk{\isacharunderscore}{\kern0pt}string\ name\isanewline
\ \ \ \ \ \ {\isacharbar}{\kern0pt}\ Const\ {\isacharunderscore}{\kern0pt}\ \ \ \ \ \ \ \ \ \ \ \ \ \ \ \ \ \ {\isacharequal}{\kern0pt}{\isachargreater}{\kern0pt}\ tm\isanewline
\ \ \ \ \ \ {\isacharbar}{\kern0pt}\ Bound\ {\isacharunderscore}{\kern0pt}\ \ \ \ \ \ \ \ \ \ \ \ \ \ \ \ \ \ {\isacharequal}{\kern0pt}{\isachargreater}{\kern0pt}\ tm\isanewline
\ \ \ \ \ \ {\isacharbar}{\kern0pt}\ Var\ {\isacharunderscore}{\kern0pt}\ \ \ \ \ \ \ \ \ \ \ \ \ \ \ \ \ \ \ \ {\isacharequal}{\kern0pt}{\isachargreater}{\kern0pt}\ tm\isanewline
\ \ \ \ \ \ {\isacharbar}{\kern0pt}\ Abs\ {\isacharparenleft}{\kern0pt}n{\isacharcomma}{\kern0pt}\ T{\isacharcomma}{\kern0pt}\ body{\isacharparenright}{\kern0pt}\ \ \ \ \ \ \ \ {\isacharequal}{\kern0pt}{\isachargreater}{\kern0pt}\ Abs\ {\isacharparenleft}{\kern0pt}n{\isacharcomma}{\kern0pt}\ T{\isacharcomma}{\kern0pt}\ subst\ body{\isacharparenright}{\kern0pt}\isanewline
\ \ \ \ \ \ {\isacharbar}{\kern0pt}\ tm{\isadigit{1}}\ {\isachardollar}{\kern0pt}\ tm{\isadigit{2}}\ \ \ \ \ \ \ \ \ \ \ \ \ \ \ {\isacharequal}{\kern0pt}{\isachargreater}{\kern0pt}\ subst\ tm{\isadigit{1}}\ {\isachardollar}{\kern0pt}\ subst\ tm{\isadigit{2}}\isanewline
\ \ in\isanewline
\ \ \ \ subst\ t\isanewline
\ \ end\isanewline
\isanewline
fun\ lambda{\isacharunderscore}{\kern0pt}sigma{\isacharunderscore}{\kern0pt}apply{\isacharunderscore}{\kern0pt}var\ {\isacharparenleft}{\kern0pt}env{\isacharunderscore}{\kern0pt}type\ {\isacharcolon}{\kern0pt}\ typ{\isacharparenright}{\kern0pt}\ {\isacharparenleft}{\kern0pt}var{\isacharunderscore}{\kern0pt}name\ {\isacharcolon}{\kern0pt}\ string{\isacharparenright}{\kern0pt}\ {\isacharcolon}{\kern0pt}\ term\ {\isacharequal}{\kern0pt}\isanewline
\ \ let\isanewline
\ \ \ \ val\ bound{\isacharunderscore}{\kern0pt}sigma\ {\isacharequal}{\kern0pt}\ Bound\ {\isadigit{0}}\isanewline
\ \ \ \ val\ body\ {\isacharequal}{\kern0pt}\ bound{\isacharunderscore}{\kern0pt}sigma\ {\isachardollar}{\kern0pt}\ HOLogic{\isachardot}{\kern0pt}mk{\isacharunderscore}{\kern0pt}string\ var{\isacharunderscore}{\kern0pt}name\isanewline
\ \ in\isanewline
\ \ \ \ Abs\ {\isacharparenleft}{\kern0pt}{\isachardoublequote}{\kern0pt}{\isasymsigma}{\isachardoublequote}{\kern0pt}{\isacharcomma}{\kern0pt}\ env{\isacharunderscore}{\kern0pt}type{\isacharcomma}{\kern0pt}\ body{\isacharparenright}{\kern0pt}\isanewline
\ \ end\isanewline
\isanewline
{\isacharparenleft}{\kern0pt}{\isacharasterisk}{\kern0pt}\ \isanewline
\ \ Crée\ une\ abstraction\ {\isasymlambda}{\isasymsigma}{\isacharcolon}{\kern0pt}typ{\isachardot}{\kern0pt}\ t\isanewline
\ \ en\ remplaçant\ les\ Free{\isacharparenleft}{\kern0pt}{\isachardoublequote}{\kern0pt}{\isasymsigma}{\isachardoublequote}{\kern0pt}{\isacharcomma}{\kern0pt}\ typ{\isacharparenright}{\kern0pt}\ par\ Bound\ {\isadigit{0}}\ dans\ t\isanewline
{\isacharasterisk}{\kern0pt}{\isacharparenright}{\kern0pt}\isanewline
fun\ make{\isacharunderscore}{\kern0pt}lambda{\isacharunderscore}{\kern0pt}sigma\ {\isacharparenleft}{\kern0pt}env{\isacharunderscore}{\kern0pt}type\ {\isacharcolon}{\kern0pt}\ typ{\isacharparenright}{\kern0pt}\ {\isacharparenleft}{\kern0pt}sigma{\isacharunderscore}{\kern0pt}name\ {\isacharcolon}{\kern0pt}\ string{\isacharparenright}{\kern0pt}\ {\isacharparenleft}{\kern0pt}t\ {\isacharcolon}{\kern0pt}\ term{\isacharparenright}{\kern0pt}\ {\isacharcolon}{\kern0pt}\ term\ {\isacharequal}{\kern0pt}\isanewline
\ \ let\isanewline
\ \ \ \ fun\ subst\ tm\ {\isacharequal}{\kern0pt}\isanewline
\ \ \ \ \ \ case\ tm\ of\isanewline
\ \ \ \ \ \ \ \ Free\ {\isacharparenleft}{\kern0pt}n{\isacharcomma}{\kern0pt}\ T{\isacharparenright}{\kern0pt}\ {\isacharequal}{\kern0pt}{\isachargreater}{\kern0pt}\ if\ n\ {\isacharequal}{\kern0pt}\ sigma{\isacharunderscore}{\kern0pt}name\ andalso\ T\ {\isacharequal}{\kern0pt}\ env{\isacharunderscore}{\kern0pt}type\ then\ Bound\ {\isadigit{0}}\ else\ tm\isanewline
\ \ \ \ \ \ {\isacharbar}{\kern0pt}\ Const\ {\isacharunderscore}{\kern0pt}\ {\isacharequal}{\kern0pt}{\isachargreater}{\kern0pt}\ tm\isanewline
\ \ \ \ \ \ {\isacharbar}{\kern0pt}\ Bound\ {\isacharunderscore}{\kern0pt}\ {\isacharequal}{\kern0pt}{\isachargreater}{\kern0pt}\ tm\isanewline
\ \ \ \ \ \ {\isacharbar}{\kern0pt}\ Var\ {\isacharunderscore}{\kern0pt}\ {\isacharequal}{\kern0pt}{\isachargreater}{\kern0pt}\ tm\isanewline
\ \ \ \ \ \ {\isacharbar}{\kern0pt}\ Abs\ {\isacharparenleft}{\kern0pt}n{\isacharcomma}{\kern0pt}\ T{\isacharcomma}{\kern0pt}\ body{\isacharparenright}{\kern0pt}\ {\isacharequal}{\kern0pt}{\isachargreater}{\kern0pt}\ Abs\ {\isacharparenleft}{\kern0pt}n{\isacharcomma}{\kern0pt}\ T{\isacharcomma}{\kern0pt}\ subst\ body{\isacharparenright}{\kern0pt}\isanewline
\ \ \ \ \ \ {\isacharbar}{\kern0pt}\ t{\isadigit{1}}\ {\isachardollar}{\kern0pt}\ t{\isadigit{2}}\ {\isacharequal}{\kern0pt}{\isachargreater}{\kern0pt}\ subst\ t{\isadigit{1}}\ {\isachardollar}{\kern0pt}\ subst\ t{\isadigit{2}}\isanewline
\ \ in\isanewline
\ \ \ \ Abs\ {\isacharparenleft}{\kern0pt}sigma{\isacharunderscore}{\kern0pt}name{\isacharcomma}{\kern0pt}\ env{\isacharunderscore}{\kern0pt}type{\isacharcomma}{\kern0pt}\ subst\ t{\isacharparenright}{\kern0pt}\isanewline
\ \ end\isanewline
{\isacharparenleft}{\kern0pt}{\isacharasterisk}{\kern0pt}\ Conversion\ d{\isacharprime}{\kern0pt}une\ action\ abstraite\ {\isacharparenleft}{\kern0pt}a{\isacharunderscore}{\kern0pt}term{\isacharparenright}{\kern0pt}\ vers\ le\ terme\ logique\ com\ {\isacharasterisk}{\kern0pt}{\isacharparenright}{\kern0pt}\isanewline
fun\ quote{\isacharunderscore}{\kern0pt}com\ {\isacharunderscore}{\kern0pt}\ SkipA\ {\isacharequal}{\kern0pt}\isanewline
\ \ \ \ \ \ {\isacharat}{\kern0pt}{\isacharbraceleft}{\kern0pt}term\ {\isachardoublequote}{\kern0pt}com{\isachardot}{\kern0pt}SKIP\ {\isacharcolon}{\kern0pt}{\isacharcolon}{\kern0pt}\ com{\isachardoublequote}{\kern0pt}{\isacharbraceright}{\kern0pt}\ \ \isanewline
\ \ {\isacharparenleft}{\kern0pt}{\isacharasterisk}{\kern0pt}{\isacharbar}{\kern0pt}\ quote{\isacharunderscore}{\kern0pt}com\ {\isacharparenleft}{\kern0pt}AssignA\ {\isacharparenleft}{\kern0pt}b{\isacharcomma}{\kern0pt}\ e{\isacharparenright}{\kern0pt}{\isacharparenright}{\kern0pt}\ {\isacharequal}{\kern0pt}\isanewline
\ \ \ \ let\isanewline
\ \ \ \ \ \ val\ var{\isacharunderscore}{\kern0pt}str\ {\isacharequal}{\kern0pt}\ HOLogic{\isachardot}{\kern0pt}mk{\isacharunderscore}{\kern0pt}string\ {\isacharparenleft}{\kern0pt}Binding{\isachardot}{\kern0pt}name{\isacharunderscore}{\kern0pt}of\ b{\isacharparenright}{\kern0pt}\isanewline
\isanewline
\ \ \ \ \ \ val\ {\isacharparenleft}{\kern0pt}sigma{\isacharcomma}{\kern0pt}\ e{\isacharprime}{\kern0pt}{\isacharparenright}{\kern0pt}\ {\isacharequal}{\kern0pt}\ subst{\isacharunderscore}{\kern0pt}with{\isacharunderscore}{\kern0pt}fresh{\isacharunderscore}{\kern0pt}env\ {\isacharat}{\kern0pt}{\isacharbraceleft}{\kern0pt}typ\ {\isachardoublequote}{\kern0pt}string\ {\isasymRightarrow}\ int{\isachardoublequote}{\kern0pt}{\isacharbraceright}{\kern0pt}\ e\isanewline
\ \ \ \ \ \ val\ lam\ {\isacharequal}{\kern0pt}\ Abs\ {\isacharparenleft}{\kern0pt}fst\ {\isacharparenleft}{\kern0pt}Term{\isachardot}{\kern0pt}dest{\isacharunderscore}{\kern0pt}Free\ sigma{\isacharparenright}{\kern0pt}{\isacharcomma}{\kern0pt}\ {\isacharat}{\kern0pt}{\isacharbraceleft}{\kern0pt}typ\ {\isachardoublequote}{\kern0pt}string\ {\isasymRightarrow}\ int{\isachardoublequote}{\kern0pt}{\isacharbraceright}{\kern0pt}{\isacharcomma}{\kern0pt}\ e{\isacharprime}{\kern0pt}{\isacharparenright}{\kern0pt}\isanewline
\ \ \ \ \ \ val\ lam\ {\isacharequal}{\kern0pt}\ Abs\ {\isacharparenleft}{\kern0pt}{\isachardoublequote}{\kern0pt}{\isasymsigma}{\isachardoublequote}{\kern0pt}{\isacharcomma}{\kern0pt}\ {\isacharat}{\kern0pt}{\isacharbraceleft}{\kern0pt}typ\ {\isachardoublequote}{\kern0pt}string\ {\isasymRightarrow}\ nat{\isachardoublequote}{\kern0pt}{\isacharbraceright}{\kern0pt}{\isacharcomma}{\kern0pt}\ e{\isacharprime}{\kern0pt}{\isacharparenright}{\kern0pt}\isanewline
\isanewline
\ \ \ \ in\isanewline
\ \ \ \ \ \ {\isacharat}{\kern0pt}{\isacharbraceleft}{\kern0pt}term\ assign{\isacharbraceright}{\kern0pt}\ {\isachardollar}{\kern0pt}\ var{\isacharunderscore}{\kern0pt}str\ {\isachardollar}{\kern0pt}\ lam\isanewline
\ \ \ \ end{\isacharasterisk}{\kern0pt}{\isacharparenright}{\kern0pt}\isanewline
\ \ \ {\isacharparenleft}{\kern0pt}{\isacharasterisk}{\kern0pt}\ {\isacharbar}{\kern0pt}\ quote{\isacharunderscore}{\kern0pt}com\ {\isacharparenleft}{\kern0pt}AssignA\ {\isacharparenleft}{\kern0pt}b{\isacharcomma}{\kern0pt}\ e{\isacharparenright}{\kern0pt}{\isacharparenright}{\kern0pt}\ {\isacharequal}{\kern0pt}\isanewline
\ \ \ \ \ \ \ \ let\isanewline
\ \ \ \ \ \ \ \ \ \ val\ var{\isacharunderscore}{\kern0pt}str\ {\isacharequal}{\kern0pt}\ HOLogic{\isachardot}{\kern0pt}mk{\isacharunderscore}{\kern0pt}string\ {\isacharparenleft}{\kern0pt}Binding{\isachardot}{\kern0pt}name{\isacharunderscore}{\kern0pt}of\ b{\isacharparenright}{\kern0pt}\isanewline
\ \ \ \ \ \ \ \ \ \ val\ e{\isacharprime}{\kern0pt}\ \ \ \ \ \ {\isacharequal}{\kern0pt}\ subst{\isacharunderscore}{\kern0pt}with{\isacharunderscore}{\kern0pt}sigma\ e\ \ \ \ \ \ \ \ \ \ \ \ \ \ \ {\isacharparenleft}{\kern0pt}{\isacharasterisk}{\kern0pt}\ {\isasymsigma}\ {\isacharprime}{\kern0pt}{\isacharprime}{\kern0pt}x{\isacharprime}{\kern0pt}{\isacharprime}{\kern0pt}\ {\isasymdots}\ {\isacharasterisk}{\kern0pt}{\isacharparenright}{\kern0pt}\isanewline
\ \ \ \ \ \ \ \ \ \ val\ lam\ \ \ \ \ {\isacharequal}{\kern0pt}\ Abs\ {\isacharparenleft}{\kern0pt}{\isachardoublequote}{\kern0pt}{\isasymsigma}{\isachardoublequote}{\kern0pt}{\isacharcomma}{\kern0pt}\ {\isacharat}{\kern0pt}{\isacharbraceleft}{\kern0pt}typ\ {\isachardoublequote}{\kern0pt}string\ {\isasymRightarrow}\ int{\isachardoublequote}{\kern0pt}{\isacharbraceright}{\kern0pt}{\isacharcomma}{\kern0pt}\ e{\isacharprime}{\kern0pt}{\isacharparenright}{\kern0pt}{\isacharasterisk}{\kern0pt}{\isacharparenright}{\kern0pt}\isanewline
{\isacharbar}{\kern0pt}\ quote{\isacharunderscore}{\kern0pt}com\ {\isacharunderscore}{\kern0pt}\ {\isacharparenleft}{\kern0pt}AssignA\ {\isacharparenleft}{\kern0pt}b{\isacharcomma}{\kern0pt}\ e{\isacharparenright}{\kern0pt}{\isacharparenright}{\kern0pt}\ {\isacharequal}{\kern0pt}\isanewline
\ \ \ \ let\isanewline
\ \ \ \ \ \ val\ var{\isacharunderscore}{\kern0pt}str\ {\isacharequal}{\kern0pt}\ HOLogic{\isachardot}{\kern0pt}mk{\isacharunderscore}{\kern0pt}string\ {\isacharparenleft}{\kern0pt}Binding{\isachardot}{\kern0pt}name{\isacharunderscore}{\kern0pt}of\ b{\isacharparenright}{\kern0pt}\isanewline
\ \ \ \ \ \ val\ expr\ {\isacharequal}{\kern0pt}\ subst{\isacharunderscore}{\kern0pt}with{\isacharunderscore}{\kern0pt}sigma\ e\ \ {\isacharparenleft}{\kern0pt}{\isacharasterisk}{\kern0pt}\ ou\ autre\ expression\ contenant\ {\isasymsigma}\ {\isacharprime}{\kern0pt}{\isacharprime}{\kern0pt}x{\isacharprime}{\kern0pt}{\isacharprime}{\kern0pt}\ {\isacharasterisk}{\kern0pt}{\isacharparenright}{\kern0pt}\isanewline
\ \ \ \ \ \ val\ lam\ {\isacharequal}{\kern0pt}\ make{\isacharunderscore}{\kern0pt}lambda{\isacharunderscore}{\kern0pt}sigma\ {\isacharat}{\kern0pt}{\isacharbraceleft}{\kern0pt}typ\ {\isachardoublequote}{\kern0pt}string\ {\isasymRightarrow}\ int{\isachardoublequote}{\kern0pt}{\isacharbraceright}{\kern0pt}\ {\isachardoublequote}{\kern0pt}{\isasymsigma}{\isachardoublequote}{\kern0pt}\ expr\isanewline
\ \ \ \ in\isanewline
\ \ \ \ \ \ {\isacharat}{\kern0pt}{\isacharbraceleft}{\kern0pt}term\ com{\isachardot}{\kern0pt}assign{\isacharbraceright}{\kern0pt}\ {\isachardollar}{\kern0pt}\ var{\isacharunderscore}{\kern0pt}str\ {\isachardollar}{\kern0pt}\ lam\isanewline
\ \ \ \ end\isanewline
\ \ \ \ \ \ \ \isanewline
\ \ {\isacharbar}{\kern0pt}\ quote{\isacharunderscore}{\kern0pt}com\ machine\ {\isacharparenleft}{\kern0pt}SeqA\ {\isacharparenleft}{\kern0pt}c{\isadigit{1}}{\isacharcomma}{\kern0pt}\ c{\isadigit{2}}{\isacharparenright}{\kern0pt}{\isacharparenright}{\kern0pt}\ {\isacharequal}{\kern0pt}\isanewline
\ \ \ \ \ \ {\isacharat}{\kern0pt}{\isacharbraceleft}{\kern0pt}term\ com{\isachardot}{\kern0pt}seq{\isacharbraceright}{\kern0pt}\ {\isachardollar}{\kern0pt}\ quote{\isacharunderscore}{\kern0pt}com\ machine\ c{\isadigit{1}}\ {\isachardollar}{\kern0pt}\ quote{\isacharunderscore}{\kern0pt}com\ machine\ c{\isadigit{2}}\isanewline
\ \ {\isacharbar}{\kern0pt}\ quote{\isacharunderscore}{\kern0pt}com\ machine\ {\isacharparenleft}{\kern0pt}LockA\ b{\isacharparenright}{\kern0pt}\ {\isacharequal}{\kern0pt}\isanewline
\ \ \ \ \ \ let\ \isanewline
\ \ \ \ \ \ \ \ val\ n{\isacharunderscore}{\kern0pt}str\ {\isacharequal}{\kern0pt}\ Binding{\isachardot}{\kern0pt}name{\isacharunderscore}{\kern0pt}of\ b\isanewline
\ \ \ \ \ \ \ \ val\ n\ {\isacharequal}{\kern0pt}\ HOLogic{\isachardot}{\kern0pt}mk{\isacharunderscore}{\kern0pt}number\ {\isacharat}{\kern0pt}{\isacharbraceleft}{\kern0pt}typ\ int{\isacharbraceright}{\kern0pt}\ {\isacharparenleft}{\kern0pt}get{\isacharunderscore}{\kern0pt}lock{\isacharunderscore}{\kern0pt}id\ n{\isacharunderscore}{\kern0pt}str\ machine{\isacharparenright}{\kern0pt}\isanewline
\ \ \ \ \ \ \ \ {\isacharparenleft}{\kern0pt}{\isacharasterisk}{\kern0pt}val\ n\ {\isacharequal}{\kern0pt}\ case\ Int{\isachardot}{\kern0pt}fromString\ n{\isacharunderscore}{\kern0pt}str\ of\isanewline
\ \ \ \ \ \ \ \ \ \ \ \ \ \ \ \ \ \ SOME\ i\ {\isacharequal}{\kern0pt}{\isachargreater}{\kern0pt}\ HOLogic{\isachardot}{\kern0pt}mk{\isacharunderscore}{\kern0pt}number\ {\isacharat}{\kern0pt}{\isacharbraceleft}{\kern0pt}typ\ nat{\isacharbraceright}{\kern0pt}\ i\isanewline
\ \ \ \ \ \ \ \ \ \ \ \ \ \ \ \ {\isacharbar}{\kern0pt}\ NONE\ {\isacharequal}{\kern0pt}{\isachargreater}{\kern0pt}\ error\ {\isacharparenleft}{\kern0pt}{\isachardoublequote}{\kern0pt}Invalid\ lock\ number{\isacharcolon}{\kern0pt}\ {\isachardoublequote}{\kern0pt}\ {\isacharcircum}{\kern0pt}\ n{\isacharunderscore}{\kern0pt}str{\isacharparenright}{\kern0pt}{\isacharasterisk}{\kern0pt}{\isacharparenright}{\kern0pt}\isanewline
\ \ \ \ \ \ in\ \isanewline
\ \ \ \ \ \ \ \ {\isacharat}{\kern0pt}{\isacharbraceleft}{\kern0pt}term\ {\isachardoublequote}{\kern0pt}com{\isachardot}{\kern0pt}lock\ {\isacharcolon}{\kern0pt}{\isacharcolon}{\kern0pt}\ int\ {\isasymRightarrow}\ com{\isachardoublequote}{\kern0pt}{\isacharbraceright}{\kern0pt}\ {\isachardollar}{\kern0pt}\ n\ \isanewline
\ \ \ \ \ \ end\isanewline
\ \ {\isacharbar}{\kern0pt}\ quote{\isacharunderscore}{\kern0pt}com\ machine\ {\isacharparenleft}{\kern0pt}UnlockA\ b{\isacharparenright}{\kern0pt}\ {\isacharequal}{\kern0pt}\isanewline
\ \ \ \ \ \ let\ \isanewline
\ \ \ \ \ \ \ \ val\ n{\isacharunderscore}{\kern0pt}str\ {\isacharequal}{\kern0pt}\ Binding{\isachardot}{\kern0pt}name{\isacharunderscore}{\kern0pt}of\ b\isanewline
\ \ \ \ \ \ \ val\ n\ {\isacharequal}{\kern0pt}\ HOLogic{\isachardot}{\kern0pt}mk{\isacharunderscore}{\kern0pt}number\ {\isacharat}{\kern0pt}{\isacharbraceleft}{\kern0pt}typ\ int{\isacharbraceright}{\kern0pt}\ {\isacharparenleft}{\kern0pt}get{\isacharunderscore}{\kern0pt}lock{\isacharunderscore}{\kern0pt}id\ n{\isacharunderscore}{\kern0pt}str\ machine{\isacharparenright}{\kern0pt}\isanewline
\ \ \ \ \ \ in\ \isanewline
\ \ \ \ \ \ \ \ {\isacharat}{\kern0pt}{\isacharbraceleft}{\kern0pt}term\ {\isachardoublequote}{\kern0pt}com{\isachardot}{\kern0pt}unlock\ {\isacharcolon}{\kern0pt}{\isacharcolon}{\kern0pt}\ int\ {\isasymRightarrow}\ com{\isachardoublequote}{\kern0pt}{\isacharbraceright}{\kern0pt}\ {\isachardollar}{\kern0pt}\ n\ \isanewline
\ \ \ \ \ \ end\isanewline
\ \ {\isacharbar}{\kern0pt}\ quote{\isacharunderscore}{\kern0pt}com\ {\isacharunderscore}{\kern0pt}\ {\isacharparenleft}{\kern0pt}SendA\ {\isacharparenleft}{\kern0pt}sv{\isacharcomma}{\kern0pt}\ e{\isacharparenright}{\kern0pt}{\isacharparenright}{\kern0pt}\ {\isacharequal}{\kern0pt}\isanewline
\ \ \ \ \ \ let\isanewline
\ \ \ \ \ \ \ \ val\ sv{\isacharunderscore}{\kern0pt}str\ {\isacharequal}{\kern0pt}\ HOLogic{\isachardot}{\kern0pt}mk{\isacharunderscore}{\kern0pt}string\ {\isacharparenleft}{\kern0pt}Binding{\isachardot}{\kern0pt}name{\isacharunderscore}{\kern0pt}of\ sv{\isacharparenright}{\kern0pt}\isanewline
\ \ \ \ \ \ \ \ val\ expr\ {\isacharequal}{\kern0pt}\ subst{\isacharunderscore}{\kern0pt}with{\isacharunderscore}{\kern0pt}sigma\ e\ \isanewline
\ \ \ \ \ \ \ \ val\ lam\ {\isacharequal}{\kern0pt}\ make{\isacharunderscore}{\kern0pt}lambda{\isacharunderscore}{\kern0pt}sigma\ {\isacharat}{\kern0pt}{\isacharbraceleft}{\kern0pt}typ\ {\isachardoublequote}{\kern0pt}string\ {\isasymRightarrow}\ int{\isachardoublequote}{\kern0pt}{\isacharbraceright}{\kern0pt}\ {\isachardoublequote}{\kern0pt}{\isasymsigma}{\isachardoublequote}{\kern0pt}\ expr\isanewline
\ \ \ \ \ \ in\ \isanewline
\ \ \ \ \ \ \ \ {\isacharat}{\kern0pt}{\isacharbraceleft}{\kern0pt}term\ {\isachardoublequote}{\kern0pt}com{\isachardot}{\kern0pt}send\ {\isacharcolon}{\kern0pt}{\isacharcolon}{\kern0pt}\ {\isacharparenleft}{\kern0pt}{\isasymsigma}\ {\isasymRightarrow}\ int{\isacharparenright}{\kern0pt}\ {\isasymRightarrow}\ string\ {\isasymRightarrow}\ com{\isachardoublequote}{\kern0pt}{\isacharbraceright}{\kern0pt}\ {\isachardollar}{\kern0pt}\ lam\ {\isachardollar}{\kern0pt}\ sv{\isacharunderscore}{\kern0pt}str\isanewline
\ \ \ \ \ \ \ end\isanewline
\isanewline
\ \ {\isacharbar}{\kern0pt}\ quote{\isacharunderscore}{\kern0pt}com\ {\isacharunderscore}{\kern0pt}\ {\isacharparenleft}{\kern0pt}ReceiveA\ {\isacharparenleft}{\kern0pt}v{\isacharcomma}{\kern0pt}\ sv{\isacharparenright}{\kern0pt}{\isacharparenright}{\kern0pt}\ {\isacharequal}{\kern0pt}\isanewline
\ \ \ \ \ \ let\isanewline
\ \ \ \ \ \ \ \ val\ v{\isacharunderscore}{\kern0pt}str\ {\isacharequal}{\kern0pt}\ HOLogic{\isachardot}{\kern0pt}mk{\isacharunderscore}{\kern0pt}string\ {\isacharparenleft}{\kern0pt}Binding{\isachardot}{\kern0pt}name{\isacharunderscore}{\kern0pt}of\ v{\isacharparenright}{\kern0pt}\isanewline
\ \ \ \ \ \ \ \ val\ sv{\isacharunderscore}{\kern0pt}str\ {\isacharequal}{\kern0pt}\ HOLogic{\isachardot}{\kern0pt}mk{\isacharunderscore}{\kern0pt}string\ {\isacharparenleft}{\kern0pt}Binding{\isachardot}{\kern0pt}name{\isacharunderscore}{\kern0pt}of\ sv{\isacharparenright}{\kern0pt}\isanewline
\ \ \ \ \ \ in\isanewline
\ \ \ \ \ \ \ \ {\isacharat}{\kern0pt}{\isacharbraceleft}{\kern0pt}term\ {\isachardoublequote}{\kern0pt}com{\isachardot}{\kern0pt}rec\ {\isacharcolon}{\kern0pt}{\isacharcolon}{\kern0pt}\ string\ {\isasymRightarrow}\ string\ {\isasymRightarrow}\ com{\isachardoublequote}{\kern0pt}{\isacharbraceright}{\kern0pt}\ {\isachardollar}{\kern0pt}\ v{\isacharunderscore}{\kern0pt}str\ {\isachardollar}{\kern0pt}\ sv{\isacharunderscore}{\kern0pt}str\isanewline
\ \ \ \ \ \ end\isanewline
\ \ {\isacharbar}{\kern0pt}\ quote{\isacharunderscore}{\kern0pt}com\ machine\ {\isacharparenleft}{\kern0pt}IfelseA\ {\isacharparenleft}{\kern0pt}{\isacharparenleft}{\kern0pt}cond{\isacharcomma}{\kern0pt}\ then{\isacharunderscore}{\kern0pt}branch{\isacharparenright}{\kern0pt}{\isacharcomma}{\kern0pt}\ else{\isacharunderscore}{\kern0pt}branch{\isacharparenright}{\kern0pt}{\isacharparenright}{\kern0pt}\ {\isacharequal}{\kern0pt}\isanewline
\ \ \ \ let\isanewline
\ \ \ \ \ {\isacharparenleft}{\kern0pt}{\isacharasterisk}{\kern0pt}\ val\ {\isacharparenleft}{\kern0pt}sigma{\isacharcomma}{\kern0pt}\ e{\isacharprime}{\kern0pt}{\isacharparenright}{\kern0pt}\ {\isacharequal}{\kern0pt}\ subst{\isacharunderscore}{\kern0pt}with{\isacharunderscore}{\kern0pt}fresh{\isacharunderscore}{\kern0pt}env\ {\isacharat}{\kern0pt}{\isacharbraceleft}{\kern0pt}typ\ {\isachardoublequote}{\kern0pt}string\ {\isasymRightarrow}\ int{\isachardoublequote}{\kern0pt}{\isacharbraceright}{\kern0pt}\ cond\isanewline
\ \ \ \ \ \ \ \ \ \isanewline
\ \ \ \ \ \ val\ lam\ {\isacharequal}{\kern0pt}\ Abs\ {\isacharparenleft}{\kern0pt}fst\ {\isacharparenleft}{\kern0pt}Term{\isachardot}{\kern0pt}dest{\isacharunderscore}{\kern0pt}Free\ sigma{\isacharparenright}{\kern0pt}{\isacharcomma}{\kern0pt}\ {\isacharat}{\kern0pt}{\isacharbraceleft}{\kern0pt}typ\ {\isachardoublequote}{\kern0pt}string\ {\isasymRightarrow}\ int{\isachardoublequote}{\kern0pt}{\isacharbraceright}{\kern0pt}{\isacharcomma}{\kern0pt}\ e{\isacharprime}{\kern0pt}{\isacharparenright}{\kern0pt}{\isacharasterisk}{\kern0pt}{\isacharparenright}{\kern0pt}\isanewline
\isanewline
\ \ \ \ \ \ val\ expr\ {\isacharequal}{\kern0pt}\ subst{\isacharunderscore}{\kern0pt}with{\isacharunderscore}{\kern0pt}sigma\ cond\isanewline
\ \ \ \ \ \ val\ lam\ {\isacharequal}{\kern0pt}\ make{\isacharunderscore}{\kern0pt}lambda{\isacharunderscore}{\kern0pt}sigma\ {\isacharat}{\kern0pt}{\isacharbraceleft}{\kern0pt}typ\ {\isachardoublequote}{\kern0pt}string\ {\isasymRightarrow}\ int{\isachardoublequote}{\kern0pt}{\isacharbraceright}{\kern0pt}\ {\isachardoublequote}{\kern0pt}{\isasymsigma}{\isachardoublequote}{\kern0pt}\ expr\isanewline
\ \ \ \ in\isanewline
\ \ \ \ \ \ {\isacharat}{\kern0pt}{\isacharbraceleft}{\kern0pt}term\ com{\isachardot}{\kern0pt}cond{\isacharbraceright}{\kern0pt}\ {\isachardollar}{\kern0pt}\ lam\ \isanewline
\ \ \ \ \ \ \ \ \ \ \ \ \ \ \ \ \ \ \ {\isachardollar}{\kern0pt}\ quote{\isacharunderscore}{\kern0pt}com\ machine\ {\isacharparenleft}{\kern0pt}sequence{\isacharunderscore}{\kern0pt}actions\ then{\isacharunderscore}{\kern0pt}branch{\isacharparenright}{\kern0pt}\isanewline
\ \ \ \ \ \ \ \ \ \ \ \ \ \ \ \ \ \ \ {\isachardollar}{\kern0pt}\ quote{\isacharunderscore}{\kern0pt}com\ machine\ {\isacharparenleft}{\kern0pt}sequence{\isacharunderscore}{\kern0pt}actions\ else{\isacharunderscore}{\kern0pt}branch{\isacharparenright}{\kern0pt}\isanewline
\ \ \ \ end\isanewline
\ \ \ \ \ \ \isanewline
\ \ {\isacharbar}{\kern0pt}\ quote{\isacharunderscore}{\kern0pt}com\ machine\ {\isacharparenleft}{\kern0pt}WhileA\ {\isacharparenleft}{\kern0pt}cond{\isacharcomma}{\kern0pt}\ body{\isacharparenright}{\kern0pt}{\isacharparenright}{\kern0pt}\ {\isacharequal}{\kern0pt}\isanewline
\ \ \ \ \ \ \ \ \ let\isanewline
\ \ \ \ \ \ val\ expr\ {\isacharequal}{\kern0pt}\ subst{\isacharunderscore}{\kern0pt}with{\isacharunderscore}{\kern0pt}sigma\ cond\ \isanewline
\ \ \ \ \ \ val\ lam\ {\isacharequal}{\kern0pt}\ make{\isacharunderscore}{\kern0pt}lambda{\isacharunderscore}{\kern0pt}sigma\ {\isacharat}{\kern0pt}{\isacharbraceleft}{\kern0pt}typ\ {\isachardoublequote}{\kern0pt}string\ {\isasymRightarrow}\ int{\isachardoublequote}{\kern0pt}{\isacharbraceright}{\kern0pt}\ {\isachardoublequote}{\kern0pt}{\isasymsigma}{\isachardoublequote}{\kern0pt}\ expr\isanewline
\ \ \ \ \ \ \ \ \ \ \ \isanewline
\ \ \ \ \ \ \ \ \ \ in\isanewline
\ \ \ \ \ \ \ \ \ \ \ \ {\isacharat}{\kern0pt}{\isacharbraceleft}{\kern0pt}term\ com{\isachardot}{\kern0pt}while{\isacharbraceright}{\kern0pt}\ {\isachardollar}{\kern0pt}\ lam\ \isanewline
\ \ \ \ \ \ \ \ \ \ \ \ \ \ \ \ \ \ \ \ \ \ \ \ \ {\isachardollar}{\kern0pt}\ quote{\isacharunderscore}{\kern0pt}com\ machine\ {\isacharparenleft}{\kern0pt}sequence{\isacharunderscore}{\kern0pt}actions\ body{\isacharparenright}{\kern0pt}\isanewline
\ \ \ \ \ \ \ \ \ \ end\isanewline
{\isacartoucheclose}\isanewline
\isanewline
\isakeywordONE{ML}\isamarkupfalse%
{\isacartoucheopen}\isanewline
{\isacharparenleft}{\kern0pt}{\isacharasterisk}{\kern0pt}\ Fonction\ pour\ déclarer\ les\ termes\ des\ threads\ {\isacharasterisk}{\kern0pt}{\isacharparenright}{\kern0pt}\isanewline
fun\ declare{\isacharunderscore}{\kern0pt}thread{\isacharunderscore}{\kern0pt}terms\ {\isacharparenleft}{\kern0pt}thy{\isacharcolon}{\kern0pt}\ theory{\isacharparenright}{\kern0pt}\ {\isacharparenleft}{\kern0pt}machine\ {\isacharcolon}{\kern0pt}\ absy{\isacharparenright}{\kern0pt}\ {\isacharparenleft}{\kern0pt}threads{\isacharcolon}{\kern0pt}\ thread{\isacharunderscore}{\kern0pt}absy\ list{\isacharparenright}{\kern0pt}\isanewline
\ \ \ \ \ \ \ \ \ \ \ \ \ \ \ \ \ \ \ \ \ \ \ \ \ {\isacharparenleft}{\kern0pt}varstab{\isacharcolon}{\kern0pt}\ {\isacharparenleft}{\kern0pt}term\ option{\isacharparenright}{\kern0pt}\ Symtab{\isachardot}{\kern0pt}table{\isacharparenright}{\kern0pt}\isanewline
\ \ \ \ \ \ \ \ \ \ \ \ \ \ \ \ \ \ \ \ \ \ \ \ \ {\isacharcolon}{\kern0pt}\ theory\ {\isacharequal}{\kern0pt}\isanewline
\ \ fold\ {\isacharparenleft}{\kern0pt}fn\ {\isacharbraceleft}{\kern0pt}nom{\isacharunderscore}{\kern0pt}thread{\isacharcomma}{\kern0pt}\ actions{\isacharcomma}{\kern0pt}\ {\isachardot}{\kern0pt}{\isachardot}{\kern0pt}{\isachardot}{\kern0pt}{\isacharbraceright}{\kern0pt}\ {\isacharequal}{\kern0pt}{\isachargreater}{\kern0pt}\ fn\ thy{\isacharprime}{\kern0pt}\ {\isacharequal}{\kern0pt}{\isachargreater}{\kern0pt}\isanewline
\ \ \ \ let\isanewline
\ \ \ \ \ \ val\ name\ {\isacharequal}{\kern0pt}\ Binding{\isachardot}{\kern0pt}name{\isacharunderscore}{\kern0pt}of\ nom{\isacharunderscore}{\kern0pt}thread\isanewline
\ \ \ \ \ \ val\ bname\ {\isacharequal}{\kern0pt}\ Binding{\isachardot}{\kern0pt}name\ {\isacharparenleft}{\kern0pt}name\ {\isacharcircum}{\kern0pt}\ {\isachardoublequote}{\kern0pt}{\isacharunderscore}{\kern0pt}term{\isachardoublequote}{\kern0pt}{\isacharparenright}{\kern0pt}\isanewline
\isanewline
\ \ \ \ \ \ {\isacharparenleft}{\kern0pt}{\isacharasterisk}{\kern0pt}\ Convert\ actions\ into\ com\ {\isacharasterisk}{\kern0pt}{\isacharparenright}{\kern0pt}\isanewline
\ \ \ \ \ \ val\ com{\isacharunderscore}{\kern0pt}ast\ {\isacharequal}{\kern0pt}\ sequence{\isacharunderscore}{\kern0pt}actions\ actions\isanewline
\ \ \ \ \ \ val\ com{\isacharunderscore}{\kern0pt}term\ {\isacharequal}{\kern0pt}\ quote{\isacharunderscore}{\kern0pt}com\ machine\ com{\isacharunderscore}{\kern0pt}ast\isanewline
\ \ \ \ \ \ val\ com{\isacharunderscore}{\kern0pt}term\ {\isacharequal}{\kern0pt}\ Type{\isachardot}{\kern0pt}constraint\ {\isacharat}{\kern0pt}{\isacharbraceleft}{\kern0pt}typ\ com{\isacharbraceright}{\kern0pt}\ com{\isacharunderscore}{\kern0pt}term\isanewline
\ \ \isanewline
\ \ \ \ \ \ {\isacharparenleft}{\kern0pt}{\isacharasterisk}{\kern0pt}\ Create\ a\ simple\ definition{\isacharcolon}{\kern0pt}\ thread{\isacharunderscore}{\kern0pt}name{\isacharunderscore}{\kern0pt}term\ {\isacharequal}{\kern0pt}\ com{\isacharunderscore}{\kern0pt}term\ {\isacharasterisk}{\kern0pt}{\isacharparenright}{\kern0pt}\isanewline
\ \ \ \ \ \ val\ lthy\ {\isacharequal}{\kern0pt}\ Named{\isacharunderscore}{\kern0pt}Target{\isachardot}{\kern0pt}theory{\isacharunderscore}{\kern0pt}init\ thy{\isacharprime}{\kern0pt}\isanewline
\ \ \ \ \ \ val\ {\isacharparenleft}{\kern0pt}{\isacharparenleft}{\kern0pt}{\isacharunderscore}{\kern0pt}{\isacharcomma}{\kern0pt}\ term{\isacharunderscore}{\kern0pt}def{\isacharparenright}{\kern0pt}{\isacharcomma}{\kern0pt}\ lthy{\isacharprime}{\kern0pt}{\isacharparenright}{\kern0pt}\ {\isacharequal}{\kern0pt}\isanewline
\ \ \ \ \ \ \ \ Local{\isacharunderscore}{\kern0pt}Theory{\isachardot}{\kern0pt}define\ {\isacharparenleft}{\kern0pt}{\isacharparenleft}{\kern0pt}bname{\isacharcomma}{\kern0pt}\ NoSyn{\isacharparenright}{\kern0pt}{\isacharcomma}{\kern0pt}\ {\isacharparenleft}{\kern0pt}{\isacharparenleft}{\kern0pt}Binding{\isachardot}{\kern0pt}empty{\isacharcomma}{\kern0pt}\ {\isacharbrackleft}{\kern0pt}{\isacharbrackright}{\kern0pt}{\isacharparenright}{\kern0pt}{\isacharcomma}{\kern0pt}\ com{\isacharunderscore}{\kern0pt}term{\isacharparenright}{\kern0pt}{\isacharparenright}{\kern0pt}\ lthy\isanewline
\ \ \ \ \ \ val\ thy{\isacharprime}{\kern0pt}{\isacharprime}{\kern0pt}\ {\isacharequal}{\kern0pt}\ Local{\isacharunderscore}{\kern0pt}Theory{\isachardot}{\kern0pt}exit{\isacharunderscore}{\kern0pt}global\ lthy{\isacharprime}{\kern0pt}\isanewline
\ \ \ \ \ \ \isanewline
\ \ \ \ \ \ val\ {\isacharunderscore}{\kern0pt}\ {\isacharequal}{\kern0pt}\ writeln\ {\isacharparenleft}{\kern0pt}{\isachardoublequote}{\kern0pt}{\isasymcheckmark}\ Defined\ command\ term{\isacharcolon}{\kern0pt}\ {\isachardoublequote}{\kern0pt}\ {\isacharcircum}{\kern0pt}\ Binding{\isachardot}{\kern0pt}name{\isacharunderscore}{\kern0pt}of\ bname{\isacharparenright}{\kern0pt}\isanewline
\isanewline
\ \ \ \ \ \ {\isacharparenleft}{\kern0pt}{\isacharasterisk}{\kern0pt}\ Optionally{\isacharcomma}{\kern0pt}\ create\ the\ semantic\ term\ separately\ {\isacharasterisk}{\kern0pt}{\isacharparenright}{\kern0pt}\isanewline
\ \ \ \ \ \ val\ sem{\isacharunderscore}{\kern0pt}bname\ {\isacharequal}{\kern0pt}\ Binding{\isachardot}{\kern0pt}name\ {\isacharparenleft}{\kern0pt}name\ {\isacharcircum}{\kern0pt}\ {\isachardoublequote}{\kern0pt}{\isacharunderscore}{\kern0pt}sem{\isachardoublequote}{\kern0pt}{\isacharparenright}{\kern0pt}\isanewline
\ \ \ \ \ \ val\ cont\ {\isacharequal}{\kern0pt}\ Abs\ {\isacharparenleft}{\kern0pt}{\isachardoublequote}{\kern0pt}s{\isachardoublequote}{\kern0pt}{\isacharcomma}{\kern0pt}\ {\isacharat}{\kern0pt}{\isacharbraceleft}{\kern0pt}typ\ {\isasymsigma}{\isacharbraceright}{\kern0pt}{\isacharcomma}{\kern0pt}\ {\isacharat}{\kern0pt}{\isacharbraceleft}{\kern0pt}term\ {\isachardoublequote}{\kern0pt}SKIP{\isacharunderscore}{\kern0pt}process{\isachardoublequote}{\kern0pt}{\isacharbraceright}{\kern0pt}{\isacharparenright}{\kern0pt}\isanewline
\ \ \ \ \ \ val\ sem{\isacharunderscore}{\kern0pt}term\ {\isacharequal}{\kern0pt}\ {\isacharat}{\kern0pt}{\isacharbraceleft}{\kern0pt}const\ Sem\isactrlsub {\isadigit{0}}{\isacharbraceright}{\kern0pt}\ {\isachardollar}{\kern0pt}\ com{\isacharunderscore}{\kern0pt}term\ {\isachardollar}{\kern0pt}\ cont\isanewline
\ \ \ \ \ \ \isanewline
\ \ \ \ \ \ val\ lthy{\isadigit{2}}\ {\isacharequal}{\kern0pt}\ Named{\isacharunderscore}{\kern0pt}Target{\isachardot}{\kern0pt}theory{\isacharunderscore}{\kern0pt}init\ thy{\isacharprime}{\kern0pt}{\isacharprime}{\kern0pt}\isanewline
\ \ \ \ \ \ val\ {\isacharparenleft}{\kern0pt}{\isacharparenleft}{\kern0pt}{\isacharunderscore}{\kern0pt}{\isacharcomma}{\kern0pt}\ sem{\isacharunderscore}{\kern0pt}def{\isacharparenright}{\kern0pt}{\isacharcomma}{\kern0pt}\ lthy{\isadigit{2}}{\isacharprime}{\kern0pt}{\isacharparenright}{\kern0pt}\ {\isacharequal}{\kern0pt}\isanewline
\ \ \ \ \ \ \ \ Local{\isacharunderscore}{\kern0pt}Theory{\isachardot}{\kern0pt}define\ {\isacharparenleft}{\kern0pt}{\isacharparenleft}{\kern0pt}sem{\isacharunderscore}{\kern0pt}bname{\isacharcomma}{\kern0pt}\ NoSyn{\isacharparenright}{\kern0pt}{\isacharcomma}{\kern0pt}\ {\isacharparenleft}{\kern0pt}{\isacharparenleft}{\kern0pt}Binding{\isachardot}{\kern0pt}empty{\isacharcomma}{\kern0pt}\ {\isacharbrackleft}{\kern0pt}{\isacharbrackright}{\kern0pt}{\isacharparenright}{\kern0pt}{\isacharcomma}{\kern0pt}\ sem{\isacharunderscore}{\kern0pt}term{\isacharparenright}{\kern0pt}{\isacharparenright}{\kern0pt}\ lthy{\isadigit{2}}\isanewline
\ \ \ \ \ \ val\ thy{\isacharprime}{\kern0pt}{\isacharprime}{\kern0pt}{\isacharprime}{\kern0pt}\ {\isacharequal}{\kern0pt}\ Local{\isacharunderscore}{\kern0pt}Theory{\isachardot}{\kern0pt}exit{\isacharunderscore}{\kern0pt}global\ lthy{\isadigit{2}}{\isacharprime}{\kern0pt}\isanewline
\ \ \ \ \ \ \isanewline
\ \ \ \ \ \ val\ {\isacharunderscore}{\kern0pt}\ {\isacharequal}{\kern0pt}\ writeln\ {\isacharparenleft}{\kern0pt}{\isachardoublequote}{\kern0pt}{\isasymcheckmark}\ Defined\ semantic\ term{\isacharcolon}{\kern0pt}\ {\isachardoublequote}{\kern0pt}\ {\isacharcircum}{\kern0pt}\ Binding{\isachardot}{\kern0pt}name{\isacharunderscore}{\kern0pt}of\ sem{\isacharunderscore}{\kern0pt}bname{\isacharparenright}{\kern0pt}\isanewline
\isanewline
\ \ \ \ in\isanewline
\ \ \ \ \ \ thy{\isacharprime}{\kern0pt}{\isacharprime}{\kern0pt}{\isacharprime}{\kern0pt}\isanewline
\ \ \ \ end{\isacharparenright}{\kern0pt}\isanewline
\ \ \ \ threads\ thy\isanewline
{\isacartoucheclose}\isanewline
\isanewline
\isakeywordONE{ML}\isamarkupfalse%
{\isacartoucheopen}\isanewline
{\isacharparenleft}{\kern0pt}{\isacharasterisk}{\kern0pt}\ Fonction\ pour\ générer\ les\ termes\ sémantiques\ {\isacharasterisk}{\kern0pt}{\isacharparenright}{\kern0pt}\isanewline
fun\ generate{\isacharunderscore}{\kern0pt}terms\ {\isacharparenleft}{\kern0pt}absy{\isacharunderscore}{\kern0pt}data\ {\isacharcolon}{\kern0pt}\ absy{\isacharparenright}{\kern0pt}\ {\isacharparenleft}{\kern0pt}thy\ {\isacharcolon}{\kern0pt}\ theory{\isacharparenright}{\kern0pt}\ {\isacharcolon}{\kern0pt}\ theory\ {\isacharequal}{\kern0pt}\isanewline
\ \ let\ \isanewline
\ \ \ \ val\ {\isacharbraceleft}{\kern0pt}system{\isacharcomma}{\kern0pt}\ globals{\isacharunderscore}{\kern0pt}decls{\isacharcomma}{\kern0pt}\ locks{\isacharunderscore}{\kern0pt}decls{\isacharcomma}{\kern0pt}\ threads{\isacharunderscore}{\kern0pt}decls{\isacharcomma}{\kern0pt}\ varstab{\isacharbraceright}{\kern0pt}\ {\isacharequal}{\kern0pt}\ absy{\isacharunderscore}{\kern0pt}data\isanewline
\ \ \ \ val\ system{\isacharunderscore}{\kern0pt}name\ {\isacharequal}{\kern0pt}\ Binding{\isachardot}{\kern0pt}name{\isacharunderscore}{\kern0pt}of\ system\isanewline
\isanewline
\ \ \ \ {\isacharparenleft}{\kern0pt}{\isacharasterisk}{\kern0pt}\ Fonction\ pour\ vérifier\ l{\isacharprime}{\kern0pt}égalité\ des\ termes\ dans\ une\ table\ {\isacharasterisk}{\kern0pt}{\isacharparenright}{\kern0pt}\isanewline
\ \ \ \ fun\ term{\isacharunderscore}{\kern0pt}equals{\isacharunderscore}{\kern0pt}in{\isacharunderscore}{\kern0pt}table\ {\isacharparenleft}{\kern0pt}tab\ {\isacharcolon}{\kern0pt}\ {\isacharparenleft}{\kern0pt}term\ option{\isacharparenright}{\kern0pt}\ Symtab{\isachardot}{\kern0pt}table{\isacharparenright}{\kern0pt}\ {\isacharparenleft}{\kern0pt}b\ {\isacharcolon}{\kern0pt}\ binding{\isacharparenright}{\kern0pt}\ {\isacharparenleft}{\kern0pt}t\ {\isacharcolon}{\kern0pt}\ term{\isacharparenright}{\kern0pt}\ {\isacharcolon}{\kern0pt}\ bool\ {\isacharequal}{\kern0pt}\isanewline
\ \ \ \ \ \ case\ Symtab{\isachardot}{\kern0pt}lookup\ tab\ {\isacharparenleft}{\kern0pt}Binding{\isachardot}{\kern0pt}name{\isacharunderscore}{\kern0pt}of\ b{\isacharparenright}{\kern0pt}\ of\isanewline
\ \ \ \ \ \ \ \ SOME\ {\isacharparenleft}{\kern0pt}SOME\ t{\isacharprime}{\kern0pt}{\isacharparenright}{\kern0pt}\ {\isacharequal}{\kern0pt}{\isachargreater}{\kern0pt}\ Term{\isachardot}{\kern0pt}aconv\ {\isacharparenleft}{\kern0pt}t{\isacharcomma}{\kern0pt}\ t{\isacharprime}{\kern0pt}{\isacharparenright}{\kern0pt}\isanewline
\ \ \ \ \ \ {\isacharbar}{\kern0pt}\ {\isacharunderscore}{\kern0pt}\ {\isacharequal}{\kern0pt}{\isachargreater}{\kern0pt}\ false\isanewline
\isanewline
\ \ \ \ {\isacharparenleft}{\kern0pt}{\isacharasterisk}{\kern0pt}\ Fonction\ pour\ mettre\ à\ jour\ l{\isacharprime}{\kern0pt}état\ {\isacharasterisk}{\kern0pt}{\isacharparenright}{\kern0pt}\isanewline
\ \ \ \ fun\ update{\isacharunderscore}{\kern0pt}state\ {\isacharparenleft}{\kern0pt}locstab\ {\isacharcolon}{\kern0pt}\ {\isacharparenleft}{\kern0pt}term\ option{\isacharparenright}{\kern0pt}\ Symtab{\isachardot}{\kern0pt}table{\isacharparenright}{\kern0pt}\isanewline
\ \ \ \ \ \ \ \ \ \ \ \ \ \ \ \ \ \ \ \ \ {\isacharparenleft}{\kern0pt}globtab\ {\isacharcolon}{\kern0pt}\ {\isacharparenleft}{\kern0pt}term\ option{\isacharparenright}{\kern0pt}\ Symtab{\isachardot}{\kern0pt}table{\isacharparenright}{\kern0pt}\isanewline
\ \ \ \ \ \ \ \ \ \ \ \ \ \ \ \ \ \ \ \ \ {\isacharparenleft}{\kern0pt}state\ {\isacharcolon}{\kern0pt}\ term\ {\isacharminus}{\kern0pt}{\isachargreater}{\kern0pt}\ int{\isacharparenright}{\kern0pt}\isanewline
\ \ \ \ \ \ \ \ \ \ \ \ \ \ \ \ \ \ \ \ \ {\isacharparenleft}{\kern0pt}x\ {\isacharcolon}{\kern0pt}\ binding{\isacharparenright}{\kern0pt}\isanewline
\ \ \ \ \ \ \ \ \ \ \ \ \ \ \ \ \ \ \ \ \ {\isacharparenleft}{\kern0pt}v\ {\isacharcolon}{\kern0pt}\ int{\isacharparenright}{\kern0pt}\ {\isacharcolon}{\kern0pt}\ term\ {\isacharminus}{\kern0pt}{\isachargreater}{\kern0pt}\ int\ {\isacharequal}{\kern0pt}\isanewline
\ \ \ \ \ \ let\isanewline
\ \ \ \ \ \ \ \ val\ x{\isacharunderscore}{\kern0pt}name\ {\isacharequal}{\kern0pt}\ Binding{\isachardot}{\kern0pt}name{\isacharunderscore}{\kern0pt}of\ x\isanewline
\ \ \ \ \ \ \ \ fun\ equals{\isacharunderscore}{\kern0pt}var\ y\ {\isacharequal}{\kern0pt}\isanewline
\ \ \ \ \ \ \ \ \ \ case\ Symtab{\isachardot}{\kern0pt}lookup\ locstab\ x{\isacharunderscore}{\kern0pt}name\ of\isanewline
\ \ \ \ \ \ \ \ \ \ \ \ SOME\ {\isacharparenleft}{\kern0pt}SOME\ t{\isacharprime}{\kern0pt}{\isacharparenright}{\kern0pt}\ {\isacharequal}{\kern0pt}{\isachargreater}{\kern0pt}\ Term{\isachardot}{\kern0pt}aconv\ {\isacharparenleft}{\kern0pt}y{\isacharcomma}{\kern0pt}\ t{\isacharprime}{\kern0pt}{\isacharparenright}{\kern0pt}\isanewline
\ \ \ \ \ \ \ \ \ \ {\isacharbar}{\kern0pt}\ {\isacharunderscore}{\kern0pt}\ {\isacharequal}{\kern0pt}{\isachargreater}{\kern0pt}\isanewline
\ \ \ \ \ \ \ \ \ \ \ \ \ \ {\isacharparenleft}{\kern0pt}case\ Symtab{\isachardot}{\kern0pt}lookup\ globtab\ x{\isacharunderscore}{\kern0pt}name\ of\isanewline
\ \ \ \ \ \ \ \ \ \ \ \ \ \ \ \ \ SOME\ {\isacharparenleft}{\kern0pt}SOME\ t{\isacharprime}{\kern0pt}{\isacharprime}{\kern0pt}{\isacharparenright}{\kern0pt}\ {\isacharequal}{\kern0pt}{\isachargreater}{\kern0pt}\ Term{\isachardot}{\kern0pt}aconv\ {\isacharparenleft}{\kern0pt}y{\isacharcomma}{\kern0pt}\ t{\isacharprime}{\kern0pt}{\isacharprime}{\kern0pt}{\isacharparenright}{\kern0pt}\isanewline
\ \ \ \ \ \ \ \ \ \ \ \ \ \ \ {\isacharbar}{\kern0pt}\ {\isacharunderscore}{\kern0pt}\ {\isacharequal}{\kern0pt}{\isachargreater}{\kern0pt}\ false{\isacharparenright}{\kern0pt}\isanewline
\ \ \ \ \ \ in\isanewline
\ \ \ \ \ \ \ \ fn\ y\ {\isacharequal}{\kern0pt}{\isachargreater}{\kern0pt}\ if\ equals{\isacharunderscore}{\kern0pt}var\ y\ then\ v\ else\ state\ y\isanewline
\ \ \ \ \ \ end\isanewline
\isanewline
\ \ \ \ {\isacharparenleft}{\kern0pt}{\isacharasterisk}{\kern0pt}\ Fonction\ d{\isacharprime}{\kern0pt}évaluation\ basique\ d{\isacharprime}{\kern0pt}un\ terme\ {\isacharasterisk}{\kern0pt}{\isacharparenright}{\kern0pt}\isanewline
\ \ \ \ fun\ eval{\isacharunderscore}{\kern0pt}term\ {\isacharparenleft}{\kern0pt}t\ {\isacharcolon}{\kern0pt}\ term{\isacharparenright}{\kern0pt}\ {\isacharparenleft}{\kern0pt}state\ {\isacharcolon}{\kern0pt}\ term\ {\isacharminus}{\kern0pt}{\isachargreater}{\kern0pt}\ int{\isacharparenright}{\kern0pt}\ {\isacharcolon}{\kern0pt}\ int\ {\isacharequal}{\kern0pt}\isanewline
\ \ \ \ \ \ case\ t\ of\isanewline
\ \ \ \ \ \ \ \ Const\ {\isacharparenleft}{\kern0pt}{\isacharat}{\kern0pt}{\isacharbraceleft}{\kern0pt}const{\isacharunderscore}{\kern0pt}name\ {\isachardoublequote}{\kern0pt}Groups{\isachardot}{\kern0pt}zero{\isacharunderscore}{\kern0pt}class{\isachardot}{\kern0pt}zero{\isachardoublequote}{\kern0pt}{\isacharbraceright}{\kern0pt}{\isacharcomma}{\kern0pt}\ {\isacharunderscore}{\kern0pt}{\isacharparenright}{\kern0pt}\ {\isacharequal}{\kern0pt}{\isachargreater}{\kern0pt}\ {\isadigit{0}}\isanewline
\ \ \ \ \ \ {\isacharbar}{\kern0pt}\ Const\ {\isacharparenleft}{\kern0pt}{\isacharat}{\kern0pt}{\isacharbraceleft}{\kern0pt}const{\isacharunderscore}{\kern0pt}name\ {\isachardoublequote}{\kern0pt}Groups{\isachardot}{\kern0pt}one{\isacharunderscore}{\kern0pt}class{\isachardot}{\kern0pt}one{\isachardoublequote}{\kern0pt}{\isacharbraceright}{\kern0pt}{\isacharcomma}{\kern0pt}\ {\isacharunderscore}{\kern0pt}{\isacharparenright}{\kern0pt}\ {\isacharequal}{\kern0pt}{\isachargreater}{\kern0pt}\ {\isadigit{1}}\isanewline
\ \ \ \ \ \ {\isacharbar}{\kern0pt}\ x\ {\isacharequal}{\kern0pt}{\isachargreater}{\kern0pt}\ state\ x\ \ {\isacharparenleft}{\kern0pt}{\isacharasterisk}{\kern0pt}\ Évaluation\ simplifiée{\isacharcomma}{\kern0pt}TODO\ {\isacharasterisk}{\kern0pt}{\isacharparenright}{\kern0pt}\isanewline
\isanewline
\ \ \ \ {\isacharparenleft}{\kern0pt}{\isacharasterisk}{\kern0pt}\ Fonction\ pour\ convertir\ une\ action\ en\ terme\ {\isacharasterisk}{\kern0pt}{\isacharparenright}{\kern0pt}\isanewline
\ \ \ \ fun\ action{\isacharunderscore}{\kern0pt}to{\isacharunderscore}{\kern0pt}term\ {\isacharparenleft}{\kern0pt}action\ {\isacharcolon}{\kern0pt}\ a{\isacharunderscore}{\kern0pt}term{\isacharparenright}{\kern0pt}\isanewline
\ \ \ \ \ \ \ \ \ \ \ \ \ \ \ \ \ \ \ \ \ \ \ {\isacharparenleft}{\kern0pt}state\ {\isacharcolon}{\kern0pt}\ term\ {\isacharminus}{\kern0pt}{\isachargreater}{\kern0pt}\ int{\isacharparenright}{\kern0pt}\isanewline
\ \ \ \ \ \ \ \ \ \ \ \ \ \ \ \ \ \ \ \ \ \ \ {\isacharparenleft}{\kern0pt}locstab\ {\isacharcolon}{\kern0pt}\ term\ option\ Symtab{\isachardot}{\kern0pt}table{\isacharparenright}{\kern0pt}\isanewline
\ \ \ \ \ \ \ \ \ \ \ \ \ \ \ \ \ \ \ \ \ \ \ {\isacharparenleft}{\kern0pt}globtab\ {\isacharcolon}{\kern0pt}\ term\ option\ Symtab{\isachardot}{\kern0pt}table{\isacharparenright}{\kern0pt}\isanewline
\ \ \ \ \ \ \ \ {\isacharcolon}{\kern0pt}\ term\ {\isacharasterisk}{\kern0pt}\ {\isacharparenleft}{\kern0pt}term\ {\isacharminus}{\kern0pt}{\isachargreater}{\kern0pt}\ int{\isacharparenright}{\kern0pt}\ {\isacharequal}{\kern0pt}\isanewline
\ \ \ \ \ \ case\ action\ of\isanewline
\ \ \ \ \ \ \ \ SkipA\ {\isacharequal}{\kern0pt}{\isachargreater}{\kern0pt}\ \isanewline
\ \ \ \ \ \ \ \ \ \ {\isacharparenleft}{\kern0pt}{\isacharat}{\kern0pt}{\isacharbraceleft}{\kern0pt}term\ {\isachardoublequote}{\kern0pt}com{\isachardot}{\kern0pt}SKIP\ {\isacharcolon}{\kern0pt}{\isacharcolon}{\kern0pt}\ com{\isachardoublequote}{\kern0pt}{\isacharbraceright}{\kern0pt}{\isacharcomma}{\kern0pt}\ state{\isacharparenright}{\kern0pt}\isanewline
\ \ \ \ \ \ \ \ \ \ \isanewline
\ \ \ \ \ \ {\isacharbar}{\kern0pt}\ AssignA\ {\isacharparenleft}{\kern0pt}name{\isacharcomma}{\kern0pt}\ assignation{\isacharparenright}{\kern0pt}\ {\isacharequal}{\kern0pt}{\isachargreater}{\kern0pt}\isanewline
\ \ \ \ \ \ \ \ \ \ let\isanewline
\ \ \ \ \ \ \ \ \ \ \ \ val\ v\ {\isacharequal}{\kern0pt}\ eval{\isacharunderscore}{\kern0pt}term\ assignation\ state\isanewline
\ \ \ \ \ \ \ \ \ \ \ \ val\ updated{\isacharunderscore}{\kern0pt}state\ {\isacharequal}{\kern0pt}\ update{\isacharunderscore}{\kern0pt}state\ locstab\ globtab\ state\ name\ v\isanewline
\ \ \ \ \ \ \ \ \ \ \ \ val\ name{\isacharunderscore}{\kern0pt}str\ {\isacharequal}{\kern0pt}\ HOLogic{\isachardot}{\kern0pt}mk{\isacharunderscore}{\kern0pt}string\ {\isacharparenleft}{\kern0pt}Binding{\isachardot}{\kern0pt}name{\isacharunderscore}{\kern0pt}of\ name{\isacharparenright}{\kern0pt}\isanewline
\ \ \ \ \ \ \ \ \ \ \ \ val\ assign{\isacharunderscore}{\kern0pt}term\ {\isacharequal}{\kern0pt}\ {\isacharat}{\kern0pt}{\isacharbraceleft}{\kern0pt}term\ assign{\isacharbraceright}{\kern0pt}\ {\isachardollar}{\kern0pt}\ name{\isacharunderscore}{\kern0pt}str\ {\isachardollar}{\kern0pt}\ assignation\isanewline
\ \ \ \ \ \ \ \ \ \ in\isanewline
\ \ \ \ \ \ \ \ \ \ \ \ {\isacharparenleft}{\kern0pt}assign{\isacharunderscore}{\kern0pt}term{\isacharcomma}{\kern0pt}\ updated{\isacharunderscore}{\kern0pt}state{\isacharparenright}{\kern0pt}\isanewline
\ \ \ \ \ \ \ \ \ \ end\isanewline
\isanewline
\ \ \ \ \ \ {\isacharbar}{\kern0pt}\ SeqA\ {\isacharparenleft}{\kern0pt}SkipA{\isacharcomma}{\kern0pt}\ C{\isacharparenright}{\kern0pt}\ {\isacharequal}{\kern0pt}{\isachargreater}{\kern0pt}\isanewline
\ \ \ \ \ \ \ \ \ \ action{\isacharunderscore}{\kern0pt}to{\isacharunderscore}{\kern0pt}term\ C\ state\ locstab\ globtab\isanewline
\isanewline
\ \ \ \ \ \ {\isacharbar}{\kern0pt}\ SeqA\ {\isacharparenleft}{\kern0pt}A{\isacharcomma}{\kern0pt}\ B{\isacharparenright}{\kern0pt}\ {\isacharequal}{\kern0pt}{\isachargreater}{\kern0pt}\isanewline
\ \ \ \ \ \ \ \ \ \ let\isanewline
\ \ \ \ \ \ \ \ \ \ \ \ val\ {\isacharparenleft}{\kern0pt}term{\isacharunderscore}{\kern0pt}A{\isacharcomma}{\kern0pt}\ stateA{\isacharparenright}{\kern0pt}\ {\isacharequal}{\kern0pt}\ action{\isacharunderscore}{\kern0pt}to{\isacharunderscore}{\kern0pt}term\ A\ state\ locstab\ globtab\isanewline
\ \ \ \ \ \ \ \ \ \ \ \ val\ {\isacharparenleft}{\kern0pt}term{\isacharunderscore}{\kern0pt}B{\isacharcomma}{\kern0pt}\ stateB{\isacharparenright}{\kern0pt}\ {\isacharequal}{\kern0pt}\ action{\isacharunderscore}{\kern0pt}to{\isacharunderscore}{\kern0pt}term\ B\ stateA\ locstab\ globtab\isanewline
\ \ \ \ \ \ \ \ \ \ \ \ val\ seq{\isacharunderscore}{\kern0pt}term\ {\isacharequal}{\kern0pt}\ {\isacharat}{\kern0pt}{\isacharbraceleft}{\kern0pt}term\ com{\isachardot}{\kern0pt}seq{\isacharbraceright}{\kern0pt}\ {\isachardollar}{\kern0pt}\ term{\isacharunderscore}{\kern0pt}A\ {\isachardollar}{\kern0pt}\ term{\isacharunderscore}{\kern0pt}B\isanewline
\ \ \ \ \ \ \ \ \ \ in\isanewline
\ \ \ \ \ \ \ \ \ \ \ \ {\isacharparenleft}{\kern0pt}seq{\isacharunderscore}{\kern0pt}term{\isacharcomma}{\kern0pt}\ stateB{\isacharparenright}{\kern0pt}\isanewline
\ \ \ \ \ \ \ \ \ \ end\isanewline
\isanewline
\ \ \ \ \ \ {\isacharbar}{\kern0pt}\ LockA\ name\ {\isacharequal}{\kern0pt}{\isachargreater}{\kern0pt}\isanewline
\ \ \ \ \ \ \ \ \ \ let\isanewline
\ \ \ \ \ \ \ \ \ \ \ \ val\ n{\isacharunderscore}{\kern0pt}str\ {\isacharequal}{\kern0pt}\ Binding{\isachardot}{\kern0pt}name{\isacharunderscore}{\kern0pt}of\ name\isanewline
\ \ \ \ \ \ \ \ \ \ \ \ val\ n\ {\isacharequal}{\kern0pt}\ HOLogic{\isachardot}{\kern0pt}mk{\isacharunderscore}{\kern0pt}number\ {\isacharat}{\kern0pt}{\isacharbraceleft}{\kern0pt}typ\ int{\isacharbraceright}{\kern0pt}\ {\isacharparenleft}{\kern0pt}get{\isacharunderscore}{\kern0pt}lock{\isacharunderscore}{\kern0pt}id\ n{\isacharunderscore}{\kern0pt}str\ absy{\isacharunderscore}{\kern0pt}data{\isacharparenright}{\kern0pt}\isanewline
\ \ \ \ \ \ \ \ \ \ \ \ val\ lock{\isacharunderscore}{\kern0pt}term\ {\isacharequal}{\kern0pt}\ {\isacharat}{\kern0pt}{\isacharbraceleft}{\kern0pt}term\ com{\isachardot}{\kern0pt}lock{\isacharbraceright}{\kern0pt}\ {\isachardollar}{\kern0pt}\ n\isanewline
\ \ \ \ \ \ \ \ \ \ in\isanewline
\ \ \ \ \ \ \ \ \ \ \ \ {\isacharparenleft}{\kern0pt}lock{\isacharunderscore}{\kern0pt}term{\isacharcomma}{\kern0pt}\ state{\isacharparenright}{\kern0pt}\isanewline
\ \ \ \ \ \ \ \ \ \ end\isanewline
\isanewline
\ \ \ \ \ \ {\isacharbar}{\kern0pt}\ UnlockA\ name\ {\isacharequal}{\kern0pt}{\isachargreater}{\kern0pt}\isanewline
\ \ \ \ \ \ \ \ \ \ let\isanewline
\ \ \ \ \ \ \ \ \ \ \ \ val\ n{\isacharunderscore}{\kern0pt}str\ {\isacharequal}{\kern0pt}\ Binding{\isachardot}{\kern0pt}name{\isacharunderscore}{\kern0pt}of\ name\isanewline
\ \ \ \ \ \ \ \ \ \ \ \ val\ n\ {\isacharequal}{\kern0pt}\ HOLogic{\isachardot}{\kern0pt}mk{\isacharunderscore}{\kern0pt}number\ {\isacharat}{\kern0pt}{\isacharbraceleft}{\kern0pt}typ\ int{\isacharbraceright}{\kern0pt}\ {\isacharparenleft}{\kern0pt}get{\isacharunderscore}{\kern0pt}lock{\isacharunderscore}{\kern0pt}id\ n{\isacharunderscore}{\kern0pt}str\ absy{\isacharunderscore}{\kern0pt}data{\isacharparenright}{\kern0pt}\isanewline
\ \ \ \ \ \ \ \ \ \ \ \ {\isacharparenleft}{\kern0pt}{\isacharasterisk}{\kern0pt}val\ n\ {\isacharequal}{\kern0pt}\ case\ Int{\isachardot}{\kern0pt}fromString\ n{\isacharunderscore}{\kern0pt}str\ of\isanewline
\ \ \ \ \ \ \ \ \ \ \ \ \ \ \ \ \ \ \ \ \ \ SOME\ i\ {\isacharequal}{\kern0pt}{\isachargreater}{\kern0pt}\ HOLogic{\isachardot}{\kern0pt}mk{\isacharunderscore}{\kern0pt}number\ {\isacharat}{\kern0pt}{\isacharbraceleft}{\kern0pt}typ\ nat{\isacharbraceright}{\kern0pt}\ i\isanewline
\ \ \ \ \ \ \ \ \ \ \ \ \ \ \ \ \ \ \ \ {\isacharbar}{\kern0pt}\ NONE\ {\isacharequal}{\kern0pt}{\isachargreater}{\kern0pt}\ error\ {\isacharparenleft}{\kern0pt}{\isachardoublequote}{\kern0pt}Invalid\ unlock\ number{\isacharcolon}{\kern0pt}\ {\isachardoublequote}{\kern0pt}\ {\isacharcircum}{\kern0pt}\ n{\isacharunderscore}{\kern0pt}str{\isacharparenright}{\kern0pt}{\isacharasterisk}{\kern0pt}{\isacharparenright}{\kern0pt}\isanewline
\ \ \ \ \ \ \ \ \ \ \ \ val\ unlock{\isacharunderscore}{\kern0pt}term\ {\isacharequal}{\kern0pt}\ {\isacharat}{\kern0pt}{\isacharbraceleft}{\kern0pt}term\ com{\isachardot}{\kern0pt}unlock{\isacharbraceright}{\kern0pt}\ {\isachardollar}{\kern0pt}\ n\isanewline
\ \ \ \ \ \ \ \ \ \ in\isanewline
\ \ \ \ \ \ \ \ \ \ \ \ {\isacharparenleft}{\kern0pt}unlock{\isacharunderscore}{\kern0pt}term{\isacharcomma}{\kern0pt}\ state{\isacharparenright}{\kern0pt}\isanewline
\ \ \ \ \ \ \ \ \ \ end\isanewline
\isanewline
\ \ \ \ \ \ {\isacharbar}{\kern0pt}\ {\isacharunderscore}{\kern0pt}\ {\isacharequal}{\kern0pt}{\isachargreater}{\kern0pt}\ \isanewline
\ \ \ \ \ \ \ \ \ \ error\ {\isacharparenleft}{\kern0pt}{\isachardoublequote}{\kern0pt}Unsupported\ action\ type\ in\ action{\isacharunderscore}{\kern0pt}to{\isacharunderscore}{\kern0pt}term{\isachardoublequote}{\kern0pt}{\isacharparenright}{\kern0pt}\isanewline
\isanewline
\ \ \ \ {\isacharparenleft}{\kern0pt}{\isacharasterisk}{\kern0pt}\ État\ par\ défaut\ {\isacharasterisk}{\kern0pt}{\isacharparenright}{\kern0pt}\isanewline
\ \ \ \ fun\ default{\isacharunderscore}{\kern0pt}state\ {\isacharparenleft}{\kern0pt}{\isacharunderscore}{\kern0pt}\ {\isacharcolon}{\kern0pt}\ term{\isacharparenright}{\kern0pt}\ {\isacharequal}{\kern0pt}\ {\isadigit{0}}\isanewline
\isanewline
\ \ \ \ {\isacharparenleft}{\kern0pt}{\isacharasterisk}{\kern0pt}\ Traiter\ tous\ les\ threads\ {\isacharasterisk}{\kern0pt}{\isacharparenright}{\kern0pt}\isanewline
\ \ \ \ val\ {\isacharunderscore}{\kern0pt}\ {\isacharequal}{\kern0pt}\ List{\isachardot}{\kern0pt}app\ {\isacharparenleft}{\kern0pt}fn\ {\isacharbraceleft}{\kern0pt}nom{\isacharunderscore}{\kern0pt}thread{\isacharcomma}{\kern0pt}\ actions{\isacharcomma}{\kern0pt}\ locstab{\isacharcomma}{\kern0pt}\ {\isachardot}{\kern0pt}{\isachardot}{\kern0pt}{\isachardot}{\kern0pt}{\isacharbraceright}{\kern0pt}\ {\isacharequal}{\kern0pt}{\isachargreater}{\kern0pt}\isanewline
\ \ \ \ \ \ case\ actions\ of\isanewline
\ \ \ \ \ \ \ \ {\isacharbrackleft}{\kern0pt}{\isacharbrackright}{\kern0pt}\ {\isacharequal}{\kern0pt}{\isachargreater}{\kern0pt}\ writeln\ {\isacharparenleft}{\kern0pt}{\isachardoublequote}{\kern0pt}Thread\ {\isachardoublequote}{\kern0pt}\ {\isacharcircum}{\kern0pt}\ Binding{\isachardot}{\kern0pt}name{\isacharunderscore}{\kern0pt}of\ nom{\isacharunderscore}{\kern0pt}thread\ {\isacharcircum}{\kern0pt}\ {\isachardoublequote}{\kern0pt}\ has\ no\ actions{\isachardoublequote}{\kern0pt}{\isacharparenright}{\kern0pt}\isanewline
\ \ \ \ \ \ {\isacharbar}{\kern0pt}\ {\isacharunderscore}{\kern0pt}\ {\isacharequal}{\kern0pt}{\isachargreater}{\kern0pt}\isanewline
\ \ \ \ \ \ \ \ \ \ let\isanewline
\ \ \ \ \ \ \ \ \ \ \ \ val\ {\isacharunderscore}{\kern0pt}\ {\isacharequal}{\kern0pt}\ writeln\ {\isacharparenleft}{\kern0pt}{\isachardoublequote}{\kern0pt}{\isacharequal}{\kern0pt}{\isacharequal}{\kern0pt}\ Processing\ thread\ {\isachardoublequote}{\kern0pt}\ {\isacharcircum}{\kern0pt}\ Binding{\isachardot}{\kern0pt}name{\isacharunderscore}{\kern0pt}of\ nom{\isacharunderscore}{\kern0pt}thread\ {\isacharcircum}{\kern0pt}\ {\isachardoublequote}{\kern0pt}\ {\isacharequal}{\kern0pt}{\isacharequal}{\kern0pt}{\isachardoublequote}{\kern0pt}{\isacharparenright}{\kern0pt}\isanewline
\ \ \ \ \ \ \ \ \ \ \ \ val\ com{\isacharunderscore}{\kern0pt}ast\ {\isacharequal}{\kern0pt}\ sequence{\isacharunderscore}{\kern0pt}actions\ actions\isanewline
\ \ \ \ \ \ \ \ \ \ \ \ val\ {\isacharparenleft}{\kern0pt}res{\isacharunderscore}{\kern0pt}term{\isacharcomma}{\kern0pt}\ final{\isacharunderscore}{\kern0pt}state{\isacharparenright}{\kern0pt}\ {\isacharequal}{\kern0pt}\ action{\isacharunderscore}{\kern0pt}to{\isacharunderscore}{\kern0pt}term\ com{\isacharunderscore}{\kern0pt}ast\ default{\isacharunderscore}{\kern0pt}state\ locstab\ varstab\isanewline
\ \ \ \ \ \ \ \ \ \ \ \ val\ {\isacharunderscore}{\kern0pt}\ {\isacharequal}{\kern0pt}\ writeln\ {\isacharparenleft}{\kern0pt}{\isachardoublequote}{\kern0pt}{\isasymcheckmark}\ Successfully\ processed\ thread\ {\isachardoublequote}{\kern0pt}\ {\isacharcircum}{\kern0pt}\ Binding{\isachardot}{\kern0pt}name{\isacharunderscore}{\kern0pt}of\ nom{\isacharunderscore}{\kern0pt}thread{\isacharparenright}{\kern0pt}\isanewline
\ \ \ \ \ \ \ \ \ \ in\isanewline
\ \ \ \ \ \ \ \ \ \ \ \ {\isacharparenleft}{\kern0pt}{\isacharparenright}{\kern0pt}\isanewline
\ \ \ \ \ \ \ \ \ \ end\isanewline
\ \ \ \ {\isacharparenright}{\kern0pt}\ threads{\isacharunderscore}{\kern0pt}decls\isanewline
\isanewline
\ \ \ \ {\isacharparenleft}{\kern0pt}{\isacharasterisk}{\kern0pt}\ Déclarer\ les\ termes\ des\ threads\ dans\ la\ théorie\ {\isacharasterisk}{\kern0pt}{\isacharparenright}{\kern0pt}\isanewline
\ \ \ \ val\ thy{\isacharprime}{\kern0pt}\ {\isacharequal}{\kern0pt}\ declare{\isacharunderscore}{\kern0pt}thread{\isacharunderscore}{\kern0pt}terms\ thy\ absy{\isacharunderscore}{\kern0pt}data\ threads{\isacharunderscore}{\kern0pt}decls\ varstab\isanewline
\ \ in\isanewline
\ \ \ \ thy{\isacharprime}{\kern0pt}\isanewline
\ \ end\isanewline
\isanewline
fun\ define{\isacharunderscore}{\kern0pt}cmd\ {\isacharparenleft}{\kern0pt}const{\isacharunderscore}{\kern0pt}name\ {\isacharcolon}{\kern0pt}\ string{\isacharparenright}{\kern0pt}\ {\isacharparenleft}{\kern0pt}rhs\ {\isacharcolon}{\kern0pt}\ term{\isacharparenright}{\kern0pt}\ {\isacharparenleft}{\kern0pt}thy\ {\isacharcolon}{\kern0pt}\ theory{\isacharparenright}{\kern0pt}\ {\isacharcolon}{\kern0pt}\ theory\ {\isacharequal}{\kern0pt}\isanewline
\ \ let\isanewline
\ \ \ \ val\ bname\ \ {\isacharequal}{\kern0pt}\ Binding{\isachardot}{\kern0pt}name\ const{\isacharunderscore}{\kern0pt}name\ \ \ \ \ \ \ \ \ \ \ {\isacharparenleft}{\kern0pt}{\isacharasterisk}{\kern0pt}\ nom\ de\ la\ constante\ {\isacharasterisk}{\kern0pt}{\isacharparenright}{\kern0pt}\isanewline
\ \ \ \ val\ lthy\ \ \ {\isacharequal}{\kern0pt}\ Named{\isacharunderscore}{\kern0pt}Target{\isachardot}{\kern0pt}theory{\isacharunderscore}{\kern0pt}init\ thy\ \ \ \ \ \ {\isacharparenleft}{\kern0pt}{\isacharasterisk}{\kern0pt}\ entrée\ en\ contexte\ local\ {\isacharasterisk}{\kern0pt}{\isacharparenright}{\kern0pt}\isanewline
\isanewline
\ \ \ \ {\isacharparenleft}{\kern0pt}{\isacharasterisk}{\kern0pt}\ \ \ definition\ const{\isacharunderscore}{\kern0pt}name\ where\ {\isachardoublequote}{\kern0pt}const{\isacharunderscore}{\kern0pt}name\ {\isacharequal}{\kern0pt}\ rhs{\isachardoublequote}{\kern0pt}\ \ \ {\isacharasterisk}{\kern0pt}{\isacharparenright}{\kern0pt}\isanewline
\ \ \ \ val\ {\isacharparenleft}{\kern0pt}{\isacharparenleft}{\kern0pt}{\isacharunderscore}{\kern0pt}{\isacharcomma}{\kern0pt}\ def{\isacharunderscore}{\kern0pt}thm{\isacharparenright}{\kern0pt}{\isacharcomma}{\kern0pt}\ lthy{\isacharprime}{\kern0pt}{\isacharparenright}{\kern0pt}\ {\isacharequal}{\kern0pt}\isanewline
\ \ \ \ \ \ Local{\isacharunderscore}{\kern0pt}Theory{\isachardot}{\kern0pt}define\ {\isacharparenleft}{\kern0pt}{\isacharparenleft}{\kern0pt}bname{\isacharcomma}{\kern0pt}\ NoSyn{\isacharparenright}{\kern0pt}{\isacharcomma}{\kern0pt}\ {\isacharparenleft}{\kern0pt}{\isacharparenleft}{\kern0pt}Binding{\isachardot}{\kern0pt}empty{\isacharcomma}{\kern0pt}\ {\isacharbrackleft}{\kern0pt}{\isacharbrackright}{\kern0pt}{\isacharparenright}{\kern0pt}{\isacharcomma}{\kern0pt}\ rhs{\isacharparenright}{\kern0pt}{\isacharparenright}{\kern0pt}\ lthy\isanewline
\isanewline
\ \ \ \ val\ thy{\isacharprime}{\kern0pt}\ \ \ {\isacharequal}{\kern0pt}\ Local{\isacharunderscore}{\kern0pt}Theory{\isachardot}{\kern0pt}exit{\isacharunderscore}{\kern0pt}global\ lthy{\isacharprime}{\kern0pt}\ \ \ \ {\isacharparenleft}{\kern0pt}{\isacharasterisk}{\kern0pt}\ on\ ressort\ en\ théorie\ {\isacharasterisk}{\kern0pt}{\isacharparenright}{\kern0pt}\isanewline
\isanewline
\ \ \ \ val\ {\isacharunderscore}{\kern0pt}\ {\isacharequal}{\kern0pt}\ writeln\ {\isacharparenleft}{\kern0pt}{\isachardoublequote}{\kern0pt}{\isasymcheckmark}\ defined\ “{\isachardoublequote}{\kern0pt}\ {\isacharcircum}{\kern0pt}\ const{\isacharunderscore}{\kern0pt}name\ {\isacharcircum}{\kern0pt}\ {\isachardoublequote}{\kern0pt}”\ {\isacharcolon}{\kern0pt}\ {\isachardoublequote}{\kern0pt}\ {\isacharcircum}{\kern0pt}\isanewline
\ \ \ \ \ \ \ \ \ \ \ \ \ \ \ \ \ \ \ \ \ Syntax{\isachardot}{\kern0pt}string{\isacharunderscore}{\kern0pt}of{\isacharunderscore}{\kern0pt}typ{\isacharunderscore}{\kern0pt}global\ thy{\isacharprime}{\kern0pt}\ {\isacharparenleft}{\kern0pt}fastype{\isacharunderscore}{\kern0pt}of\ rhs{\isacharparenright}{\kern0pt}{\isacharparenright}{\kern0pt}\isanewline
\ \ \ \ val\ {\isacharunderscore}{\kern0pt}\ {\isacharequal}{\kern0pt}\ writeln\ {\isacharparenleft}{\kern0pt}{\isachardoublequote}{\kern0pt}{\isachardoublequote}{\kern0pt}\ {\isacharcircum}{\kern0pt}\ Syntax{\isachardot}{\kern0pt}string{\isacharunderscore}{\kern0pt}of{\isacharunderscore}{\kern0pt}term{\isacharunderscore}{\kern0pt}global\ thy{\isacharprime}{\kern0pt}\ rhs{\isacharparenright}{\kern0pt}\isanewline
\ \ \ \ val\ {\isacharunderscore}{\kern0pt}\ {\isacharequal}{\kern0pt}\ writeln\ {\isacharparenleft}{\kern0pt}{\isachardoublequote}{\kern0pt}\ theorem{\isacharcolon}{\kern0pt}\ {\isachardoublequote}{\kern0pt}\ {\isacharcircum}{\kern0pt}\ Thm{\isachardot}{\kern0pt}string{\isacharunderscore}{\kern0pt}of{\isacharunderscore}{\kern0pt}thm{\isacharunderscore}{\kern0pt}global\ thy{\isacharprime}{\kern0pt}\ {\isacharparenleft}{\kern0pt}snd\ def{\isacharunderscore}{\kern0pt}thm{\isacharparenright}{\kern0pt}{\isacharparenright}{\kern0pt}\isanewline
\ \ in\isanewline
\ \ \ \ thy{\isacharprime}{\kern0pt}\isanewline
\ \ end\isanewline
{\isacartoucheclose}%
\endisatagML
{\isafoldML}%
%
\isadelimML
%
\endisadelimML
%
\isadelimdocument
%
\endisadelimdocument
%
\isatagdocument
%
\isamarkupsection{Impression%
}
\isamarkuptrue%
%
\endisatagdocument
{\isafolddocument}%
%
\isadelimdocument
%
\endisadelimdocument
%
\isadelimML
%
\endisadelimML
%
\isatagML
\isakeywordONE{ML}\isamarkupfalse%
{\isacartoucheopen}\ \ \ \isanewline
fun\ mk{\isacharunderscore}{\kern0pt}sumT\ {\isacharparenleft}{\kern0pt}T{\isadigit{1}}{\isacharcomma}{\kern0pt}\ T{\isadigit{2}}{\isacharparenright}{\kern0pt}\ {\isacharequal}{\kern0pt}\ Type\ {\isacharparenleft}{\kern0pt}{\isachardoublequote}{\kern0pt}Sum{\isacharunderscore}{\kern0pt}Type{\isachardot}{\kern0pt}sum{\isachardoublequote}{\kern0pt}{\isacharcomma}{\kern0pt}\ {\isacharbrackleft}{\kern0pt}T{\isadigit{1}}{\isacharcomma}{\kern0pt}\ T{\isadigit{2}}{\isacharbrackright}{\kern0pt}{\isacharparenright}{\kern0pt}\isanewline
\isanewline
val\ UNIV{\isacharunderscore}{\kern0pt}\ {\isacharequal}{\kern0pt}\ HOLogic{\isachardot}{\kern0pt}mk{\isacharunderscore}{\kern0pt}UNIV\ {\isacharat}{\kern0pt}{\isacharbraceleft}{\kern0pt}typ\ evs{\isacharbraceright}{\kern0pt}\isanewline
val\ empty{\isacharunderscore}{\kern0pt}evs{\isacharunderscore}{\kern0pt}set\ {\isacharequal}{\kern0pt}\ Const\ {\isacharparenleft}{\kern0pt}{\isacharat}{\kern0pt}{\isacharbraceleft}{\kern0pt}const{\isacharunderscore}{\kern0pt}name\ bot{\isacharbraceright}{\kern0pt}{\isacharcomma}{\kern0pt}\ {\isacharat}{\kern0pt}{\isacharbraceleft}{\kern0pt}typ\ {\isachardoublequote}{\kern0pt}evs\ set{\isachardoublequote}{\kern0pt}{\isacharbraceright}{\kern0pt}{\isacharparenright}{\kern0pt}\isanewline
\isanewline
\isanewline
fun\ mk{\isacharunderscore}{\kern0pt}fun{\isacharunderscore}{\kern0pt}upd\ {\isacharparenleft}{\kern0pt}f{\isacharcomma}{\kern0pt}\ x{\isacharcomma}{\kern0pt}\ y{\isacharparenright}{\kern0pt}\ {\isacharequal}{\kern0pt}\isanewline
\ \ let\isanewline
\ \ \ \ val\ T{\isacharunderscore}{\kern0pt}dom\ {\isacharequal}{\kern0pt}\ {\isacharat}{\kern0pt}{\isacharbraceleft}{\kern0pt}typ\ string{\isacharbraceright}{\kern0pt}\isanewline
\ \ \ \ val\ T{\isacharunderscore}{\kern0pt}cod\ {\isacharequal}{\kern0pt}\ {\isacharat}{\kern0pt}{\isacharbraceleft}{\kern0pt}typ\ int{\isacharbraceright}{\kern0pt}\isanewline
\ \ in\isanewline
\ \ \ \ Const\ {\isacharparenleft}{\kern0pt}{\isacharat}{\kern0pt}{\isacharbraceleft}{\kern0pt}const{\isacharunderscore}{\kern0pt}name\ fun{\isacharunderscore}{\kern0pt}upd{\isacharbraceright}{\kern0pt}{\isacharcomma}{\kern0pt}\isanewline
\ \ \ \ \ \ \ \ \ \ \ {\isacharparenleft}{\kern0pt}T{\isacharunderscore}{\kern0pt}dom\ {\isacharminus}{\kern0pt}{\isacharminus}{\kern0pt}{\isachargreater}{\kern0pt}\ T{\isacharunderscore}{\kern0pt}cod{\isacharparenright}{\kern0pt}\ {\isacharminus}{\kern0pt}{\isacharminus}{\kern0pt}{\isachargreater}{\kern0pt}\ T{\isacharunderscore}{\kern0pt}dom\ {\isacharminus}{\kern0pt}{\isacharminus}{\kern0pt}{\isachargreater}{\kern0pt}\ T{\isacharunderscore}{\kern0pt}cod\ {\isacharminus}{\kern0pt}{\isacharminus}{\kern0pt}{\isachargreater}{\kern0pt}\ {\isacharparenleft}{\kern0pt}T{\isacharunderscore}{\kern0pt}dom\ {\isacharminus}{\kern0pt}{\isacharminus}{\kern0pt}{\isachargreater}{\kern0pt}\ T{\isacharunderscore}{\kern0pt}cod{\isacharparenright}{\kern0pt}{\isacharparenright}{\kern0pt}\isanewline
\ \ \ \ {\isachardollar}{\kern0pt}\ f\ {\isachardollar}{\kern0pt}\ x\ {\isachardollar}{\kern0pt}\ y\isanewline
\ \ end\isanewline
\isanewline
fun\ make{\isacharunderscore}{\kern0pt}sigma{\isacharunderscore}{\kern0pt}term\ {\isacharparenleft}{\kern0pt}tab\ {\isacharcolon}{\kern0pt}\ term\ option\ Symtab{\isachardot}{\kern0pt}table{\isacharparenright}{\kern0pt}\ {\isacharcolon}{\kern0pt}\ term\ {\isacharequal}{\kern0pt}\isanewline
\ \ let\isanewline
\ \ \ \ val\ zero\ {\isacharequal}{\kern0pt}\ HOLogic{\isachardot}{\kern0pt}mk{\isacharunderscore}{\kern0pt}number\ {\isacharat}{\kern0pt}{\isacharbraceleft}{\kern0pt}typ\ int{\isacharbraceright}{\kern0pt}\ {\isadigit{0}}\isanewline
\ \ \ \ val\ base\ {\isacharequal}{\kern0pt}\ Abs\ {\isacharparenleft}{\kern0pt}{\isachardoublequote}{\kern0pt}{\isacharunderscore}{\kern0pt}{\isachardoublequote}{\kern0pt}{\isacharcomma}{\kern0pt}\ {\isacharat}{\kern0pt}{\isacharbraceleft}{\kern0pt}typ\ string{\isacharbraceright}{\kern0pt}{\isacharcomma}{\kern0pt}\ zero{\isacharparenright}{\kern0pt}\ \isanewline
\isanewline
\ \ \ \ fun\ apply{\isacharunderscore}{\kern0pt}update\ {\isacharparenleft}{\kern0pt}name{\isacharcomma}{\kern0pt}\ SOME\ value{\isacharparenright}{\kern0pt}\ acc\ {\isacharequal}{\kern0pt}\isanewline
\ \ \ \ \ \ \ \ \ \ mk{\isacharunderscore}{\kern0pt}fun{\isacharunderscore}{\kern0pt}upd\ {\isacharparenleft}{\kern0pt}acc{\isacharcomma}{\kern0pt}\ HOLogic{\isachardot}{\kern0pt}mk{\isacharunderscore}{\kern0pt}string\ name{\isacharcomma}{\kern0pt}\ value{\isacharparenright}{\kern0pt}\isanewline
\ \ \ \ \ \ {\isacharbar}{\kern0pt}\ apply{\isacharunderscore}{\kern0pt}update\ {\isacharparenleft}{\kern0pt}{\isacharunderscore}{\kern0pt}{\isacharcomma}{\kern0pt}\ NONE{\isacharparenright}{\kern0pt}\ acc\ {\isacharequal}{\kern0pt}\ acc\isanewline
\isanewline
\ \ \ \ val\ updated\ {\isacharequal}{\kern0pt}\ Symtab{\isachardot}{\kern0pt}fold\ apply{\isacharunderscore}{\kern0pt}update\ tab\ base\isanewline
\ \ in\isanewline
\ \ \ \ updated\ \ \isanewline
\ \ end\isanewline
\isanewline
val\ {\isacharunderscore}{\kern0pt}\ {\isacharequal}{\kern0pt}\isanewline
\ \ Outer{\isacharunderscore}{\kern0pt}Syntax{\isachardot}{\kern0pt}command\ \isactrlcommandUNDERSCOREkeyword {\isasymopen}SYSTEM{\isasymclose}\isanewline
\ \ \ \ \ \ {\isachardoublequote}{\kern0pt}defines\ System\ Specification{\isachardoublequote}{\kern0pt}\isanewline
\ \ \ \ \ \ {\isacharparenleft}{\kern0pt}parse{\isacharunderscore}{\kern0pt}system{\isacharunderscore}{\kern0pt}spec\ {\isachargreater}{\kern0pt}{\isachargreater}{\kern0pt}\isanewline
\ \ \ \ \ \ \ {\isacharparenleft}{\kern0pt}fn\ system{\isacharunderscore}{\kern0pt}data\ {\isacharequal}{\kern0pt}{\isachargreater}{\kern0pt}\ Toplevel{\isachardot}{\kern0pt}theory\ {\isacharparenleft}{\kern0pt}fn\ thy\ {\isacharequal}{\kern0pt}{\isachargreater}{\kern0pt}\isanewline
\ \ \ \ \ \ \ \ let\isanewline
\ \ \ \ \ \ \ \ \ \ {\isacharparenleft}{\kern0pt}{\isacharasterisk}{\kern0pt}\ {\isadigit{1}}{\isachardot}{\kern0pt}\ syntactic\ {\isacharampersand}{\kern0pt}\ semantic\ checks\ {\isacharminus}{\kern0pt}{\isacharminus}{\kern0pt}{\isacharminus}{\kern0pt}{\isacharminus}{\kern0pt}{\isacharminus}{\kern0pt}{\isacharminus}{\kern0pt}{\isacharminus}{\kern0pt}{\isacharminus}{\kern0pt}{\isacharminus}{\kern0pt}{\isacharminus}{\kern0pt}{\isacharminus}{\kern0pt}{\isacharminus}{\kern0pt}{\isacharminus}{\kern0pt}{\isacharminus}{\kern0pt}{\isacharminus}{\kern0pt}{\isacharminus}{\kern0pt}{\isacharminus}{\kern0pt}{\isacharminus}{\kern0pt}{\isacharminus}{\kern0pt}{\isacharminus}{\kern0pt}{\isacharminus}{\kern0pt}{\isacharminus}{\kern0pt}{\isacharminus}{\kern0pt}{\isacharminus}{\kern0pt}{\isacharminus}{\kern0pt}{\isacharminus}{\kern0pt}{\isacharminus}{\kern0pt}{\isacharminus}{\kern0pt}{\isacharminus}{\kern0pt}{\isacharminus}{\kern0pt}{\isacharminus}{\kern0pt}{\isacharminus}{\kern0pt}\ {\isacharasterisk}{\kern0pt}{\isacharparenright}{\kern0pt}\isanewline
\ \ \ \ \ \ \ \ \ \ val\ checked\ \ {\isacharequal}{\kern0pt}\ context{\isacharunderscore}{\kern0pt}check\ system{\isacharunderscore}{\kern0pt}data\ thy\isanewline
\ \ \ \ \ \ \ \ \ \ val\ thy{\isacharunderscore}{\kern0pt}sem\ \ {\isacharequal}{\kern0pt}\ semantic{\isacharunderscore}{\kern0pt}check\ checked\ thy\isanewline
\isanewline
\ \ \ \ \ \ \ \ \ \ val\ lthy{\isadigit{0}}\ {\isacharequal}{\kern0pt}\ Named{\isacharunderscore}{\kern0pt}Target{\isachardot}{\kern0pt}theory{\isacharunderscore}{\kern0pt}init\ thy{\isacharunderscore}{\kern0pt}sem\isanewline
\isanewline
\ \ \ \ \ \ \ \ \ \ fun\ define{\isacharunderscore}{\kern0pt}cmd\ {\isacharparenleft}{\kern0pt}{\isacharbraceleft}{\kern0pt}nom{\isacharunderscore}{\kern0pt}thread{\isacharcomma}{\kern0pt}\ actions{\isacharcomma}{\kern0pt}\ {\isachardot}{\kern0pt}{\isachardot}{\kern0pt}{\isachardot}{\kern0pt}{\isacharbraceright}{\kern0pt}\ {\isacharcolon}{\kern0pt}\ thread{\isacharunderscore}{\kern0pt}absy{\isacharparenright}{\kern0pt}\ lthy\ {\isacharequal}{\kern0pt}\isanewline
\ \ \ \ \ \ \ \ \ \ \ \ let\isanewline
\ \ \ \ \ \ \ \ \ \ \ \ \ \ val\ name\ \ \ \ \ \ {\isacharequal}{\kern0pt}\ Binding{\isachardot}{\kern0pt}name{\isacharunderscore}{\kern0pt}of\ nom{\isacharunderscore}{\kern0pt}thread\isanewline
\ \ \ \ \ \ \ \ \ \ \ \ \ \ val\ const{\isacharunderscore}{\kern0pt}bnd\ {\isacharequal}{\kern0pt}\ Binding{\isachardot}{\kern0pt}name\ {\isacharparenleft}{\kern0pt}name\ {\isacharcircum}{\kern0pt}\ {\isachardoublequote}{\kern0pt}{\isacharunderscore}{\kern0pt}cmd{\isachardoublequote}{\kern0pt}{\isacharparenright}{\kern0pt}\isanewline
\ \ \ \ \ \ \ \ \ \ \ \ \ \ val\ com{\isacharunderscore}{\kern0pt}term\ \ {\isacharequal}{\kern0pt}\ quote{\isacharunderscore}{\kern0pt}com\ checked\ {\isacharparenleft}{\kern0pt}sequence{\isacharunderscore}{\kern0pt}actions\ actions{\isacharparenright}{\kern0pt}\isanewline
\ \ \ \ \ \ \ \ \ \ \ \ \ \ val\ {\isacharparenleft}{\kern0pt}{\isacharparenleft}{\kern0pt}{\isacharunderscore}{\kern0pt}{\isacharcomma}{\kern0pt}\ {\isacharunderscore}{\kern0pt}{\isacharparenright}{\kern0pt}{\isacharcomma}{\kern0pt}\ lthy{\isacharprime}{\kern0pt}{\isacharparenright}{\kern0pt}\ {\isacharequal}{\kern0pt}\isanewline
\ \ \ \ \ \ \ \ \ \ \ \ \ \ \ \ \ \ \ \ Local{\isacharunderscore}{\kern0pt}Theory{\isachardot}{\kern0pt}define\ {\isacharparenleft}{\kern0pt}{\isacharparenleft}{\kern0pt}const{\isacharunderscore}{\kern0pt}bnd{\isacharcomma}{\kern0pt}\ NoSyn{\isacharparenright}{\kern0pt}{\isacharcomma}{\kern0pt}\isanewline
\ \ \ \ \ \ \ \ \ \ \ \ \ \ \ \ \ \ \ \ \ \ \ \ \ \ \ \ \ \ \ \ \ \ \ \ \ \ \ \ \ {\isacharparenleft}{\kern0pt}{\isacharparenleft}{\kern0pt}Binding{\isachardot}{\kern0pt}empty{\isacharcomma}{\kern0pt}\ {\isacharbrackleft}{\kern0pt}{\isacharbrackright}{\kern0pt}{\isacharparenright}{\kern0pt}{\isacharcomma}{\kern0pt}\ com{\isacharunderscore}{\kern0pt}term{\isacharparenright}{\kern0pt}{\isacharparenright}{\kern0pt}\ lthy\isanewline
\isanewline
\ \ \ \ \ \ {\isacharparenleft}{\kern0pt}{\isacharasterisk}{\kern0pt}Affichage{\isacharasterisk}{\kern0pt}{\isacharparenright}{\kern0pt}\isanewline
\ \ \ \ \ \ \ \ \ \ \ \ \ \ val\ {\isacharunderscore}{\kern0pt}\ {\isacharequal}{\kern0pt}\ writeln\ {\isacharparenleft}{\kern0pt}{\isachardoublequote}{\kern0pt}{\isasymcheckmark}\ Defined\ constant{\isacharcolon}{\kern0pt}\ {\isachardoublequote}{\kern0pt}\ {\isacharcircum}{\kern0pt}\ Binding{\isachardot}{\kern0pt}name{\isacharunderscore}{\kern0pt}of\ const{\isacharunderscore}{\kern0pt}bnd{\isacharparenright}{\kern0pt}\isanewline
\ \ \ \ \ \ \ \ \ \ \ \ \isanewline
\ \ \ \ \ \ \ \ \ \ \ \ \ \ val\ {\isacharunderscore}{\kern0pt}\ {\isacharequal}{\kern0pt}\ writeln\ {\isacharparenleft}{\kern0pt}{\isachardoublequote}{\kern0pt}\ \ {\isachardoublequote}{\kern0pt}\ {\isacharcircum}{\kern0pt}\ Syntax{\isachardot}{\kern0pt}string{\isacharunderscore}{\kern0pt}of{\isacharunderscore}{\kern0pt}term{\isacharunderscore}{\kern0pt}global\ thy{\isacharunderscore}{\kern0pt}sem\ com{\isacharunderscore}{\kern0pt}term{\isacharparenright}{\kern0pt}\isanewline
\isanewline
\ \ \ \ \ \ \ \ \ \ \ \ \ \ in\ lthy{\isacharprime}{\kern0pt}\ end\isanewline
\isanewline
\ \ \ \ \ \ \ \ \ \ val\ lthy{\isacharunderscore}{\kern0pt}final\ {\isacharequal}{\kern0pt}\ fold\ define{\isacharunderscore}{\kern0pt}cmd\ {\isacharparenleft}{\kern0pt}{\isacharhash}{\kern0pt}threads{\isacharunderscore}{\kern0pt}decls\ checked{\isacharparenright}{\kern0pt}\ lthy{\isadigit{0}}\isanewline
\ \ \ \ \ \ \ \ \ \ val\ thy{\isacharprime}{\kern0pt}\ \ \ \ \ \ \ {\isacharequal}{\kern0pt}\ Local{\isacharunderscore}{\kern0pt}Theory{\isachardot}{\kern0pt}exit{\isacharunderscore}{\kern0pt}global\ lthy{\isacharunderscore}{\kern0pt}final\isanewline
\isanewline
\ \ \ \ \ \ \ \ \ \ \ \ \isanewline
\ \ \ \ \ \ \ \ \ \ {\isacharparenleft}{\kern0pt}{\isacharasterisk}{\kern0pt}\ {\isacharminus}{\kern0pt}{\isacharminus}{\kern0pt}{\isacharminus}{\kern0pt}{\isacharminus}{\kern0pt}{\isacharminus}{\kern0pt}{\isacharminus}{\kern0pt}{\isacharminus}{\kern0pt}{\isacharminus}{\kern0pt}{\isacharminus}{\kern0pt}{\isacharminus}{\kern0pt}{\isacharminus}{\kern0pt}{\isacharminus}{\kern0pt}{\isacharminus}{\kern0pt}{\isacharminus}{\kern0pt}{\isacharminus}{\kern0pt}{\isacharminus}{\kern0pt}{\isacharminus}{\kern0pt}{\isacharminus}{\kern0pt}{\isacharminus}{\kern0pt}{\isacharminus}{\kern0pt}{\isacharminus}{\kern0pt}{\isacharminus}{\kern0pt}{\isacharminus}{\kern0pt}{\isacharminus}{\kern0pt}{\isacharminus}{\kern0pt}{\isacharminus}{\kern0pt}{\isacharminus}{\kern0pt}{\isacharminus}{\kern0pt}{\isacharminus}{\kern0pt}{\isacharminus}{\kern0pt}{\isacharminus}{\kern0pt}{\isacharminus}{\kern0pt}{\isacharminus}{\kern0pt}{\isacharminus}{\kern0pt}{\isacharminus}{\kern0pt}{\isacharminus}{\kern0pt}{\isacharminus}{\kern0pt}{\isacharminus}{\kern0pt}{\isacharminus}{\kern0pt}{\isacharminus}{\kern0pt}{\isacharminus}{\kern0pt}{\isacharminus}{\kern0pt}{\isacharminus}{\kern0pt}{\isacharminus}{\kern0pt}{\isacharminus}{\kern0pt}{\isacharminus}{\kern0pt}{\isacharminus}{\kern0pt}{\isacharminus}{\kern0pt}{\isacharminus}{\kern0pt}{\isacharminus}{\kern0pt}{\isacharminus}{\kern0pt}{\isacharminus}{\kern0pt}{\isacharminus}{\kern0pt}{\isacharminus}{\kern0pt}{\isacharminus}{\kern0pt}{\isacharminus}{\kern0pt}{\isacharminus}{\kern0pt}{\isacharminus}{\kern0pt}{\isacharminus}{\kern0pt}{\isacharminus}{\kern0pt}{\isacharminus}{\kern0pt}{\isacharminus}{\kern0pt}{\isacharminus}{\kern0pt}{\isacharminus}{\kern0pt}\ {\isacharasterisk}{\kern0pt}{\isacharparenright}{\kern0pt}\isanewline
\ \ \ \ \ \ \ \ \ \ {\isacharparenleft}{\kern0pt}{\isacharasterisk}{\kern0pt}\ {\isadigit{2}}{\isachardot}{\kern0pt}\ \ Construction\ du\ réseau\ GLOBALVARS\ {\isacharminus}{\kern0pt}{\isacharminus}{\kern0pt}{\isacharminus}{\kern0pt}{\isacharminus}{\kern0pt}{\isacharminus}{\kern0pt}{\isacharminus}{\kern0pt}{\isacharminus}{\kern0pt}{\isacharminus}{\kern0pt}{\isacharminus}{\kern0pt}{\isacharminus}{\kern0pt}{\isacharminus}{\kern0pt}{\isacharminus}{\kern0pt}{\isacharminus}{\kern0pt}{\isacharminus}{\kern0pt}{\isacharminus}{\kern0pt}{\isacharminus}{\kern0pt}{\isacharminus}{\kern0pt}{\isacharminus}{\kern0pt}{\isacharminus}{\kern0pt}{\isacharminus}{\kern0pt}{\isacharminus}{\kern0pt}{\isacharminus}{\kern0pt}{\isacharminus}{\kern0pt}{\isacharminus}{\kern0pt}{\isacharminus}{\kern0pt}\ {\isacharasterisk}{\kern0pt}{\isacharparenright}{\kern0pt}\isanewline
\ \ \ \ \ \ \ \ {\isacharparenleft}{\kern0pt}{\isacharasterisk}{\kern0pt}\ \ val\ globals{\isacharunderscore}{\kern0pt}names\ {\isacharcolon}{\kern0pt}\ term\ list\ {\isacharequal}{\kern0pt}\isanewline
\ \ \ \ \ \ \ \ \ \ \ \ map\ {\isacharparenleft}{\kern0pt}fn\ {\isacharparenleft}{\kern0pt}b{\isacharcomma}{\kern0pt}\ {\isacharunderscore}{\kern0pt}{\isacharcomma}{\kern0pt}\ {\isacharunderscore}{\kern0pt}{\isacharcomma}{\kern0pt}\ {\isacharunderscore}{\kern0pt}{\isacharparenright}{\kern0pt}\ {\isacharequal}{\kern0pt}{\isachargreater}{\kern0pt}\ HOLogic{\isachardot}{\kern0pt}mk{\isacharunderscore}{\kern0pt}string\ {\isacharparenleft}{\kern0pt}Binding{\isachardot}{\kern0pt}name{\isacharunderscore}{\kern0pt}of\ b{\isacharparenright}{\kern0pt}{\isacharparenright}{\kern0pt}\ {\isacharparenleft}{\kern0pt}{\isacharhash}{\kern0pt}globals{\isacharunderscore}{\kern0pt}decls\ checked{\isacharparenright}{\kern0pt}\isanewline
\ \ \ \ \ \ \ \ \ \ \isanewline
\ \ \ \ \ \ \ \ \ \ val\ globals{\isacharunderscore}{\kern0pt}list{\isacharunderscore}{\kern0pt}term\ {\isacharequal}{\kern0pt}\isanewline
\ \ \ \ \ \ \ \ \ \ \ \ HOLogic{\isachardot}{\kern0pt}mk{\isacharunderscore}{\kern0pt}list\ {\isacharat}{\kern0pt}{\isacharbraceleft}{\kern0pt}typ\ string{\isacharbraceright}{\kern0pt}\ globals{\isacharunderscore}{\kern0pt}names\isanewline
\ \ \ \ \ \ \ \ \ \ \isanewline
\ \ \ \ \ \ \ \ \ \ val\ globals{\isacharunderscore}{\kern0pt}mset{\isacharunderscore}{\kern0pt}term\ {\isacharequal}{\kern0pt}\isanewline
\ \ \ \ \ \ \ \ \ \ \ \ Cspm{\isacharunderscore}{\kern0pt}API{\isachardot}{\kern0pt}mk{\isacharunderscore}{\kern0pt}mset\ globals{\isacharunderscore}{\kern0pt}list{\isacharunderscore}{\kern0pt}term\isanewline
\ \ \ \ \ \ \ \ \ \ \isanewline
\ \ \ \ \ \ \ \ \ \ {\isacharparenleft}{\kern0pt}{\isacharasterisk}{\kern0pt}val\ sumT\ {\isacharequal}{\kern0pt}\ Type\ {\isacharparenleft}{\kern0pt}{\isachardoublequote}{\kern0pt}Sum{\isacharunderscore}{\kern0pt}Type{\isachardot}{\kern0pt}sum{\isachardoublequote}{\kern0pt}{\isacharcomma}{\kern0pt}\ {\isacharbrackleft}{\kern0pt}{\isacharat}{\kern0pt}{\isacharbraceleft}{\kern0pt}typ\ sema{\isacharunderscore}{\kern0pt}evs{\isacharbraceright}{\kern0pt}{\isacharcomma}{\kern0pt}\ {\isacharat}{\kern0pt}{\isacharbraceleft}{\kern0pt}typ\ vars{\isacharunderscore}{\kern0pt}ev{\isacharbraceright}{\kern0pt}{\isacharbrackright}{\kern0pt}{\isacharparenright}{\kern0pt}{\isacharasterisk}{\kern0pt}{\isacharparenright}{\kern0pt}\isanewline
\ \ \ \ \ \ \ \ \ \ val\ procT\ {\isacharequal}{\kern0pt}\ Cspm{\isacharunderscore}{\kern0pt}API{\isachardot}{\kern0pt}mk{\isacharunderscore}{\kern0pt}processT\ {\isacharat}{\kern0pt}{\isacharbraceleft}{\kern0pt}typ\ evs{\isacharbraceright}{\kern0pt}\isanewline
\ \ \ \ \ \ \ \ \ \ val\ sigma{\isadigit{0}}{\isacharunderscore}{\kern0pt}term\ {\isacharequal}{\kern0pt}\ make{\isacharunderscore}{\kern0pt}sigma{\isacharunderscore}{\kern0pt}term\ {\isacharparenleft}{\kern0pt}{\isacharhash}{\kern0pt}varstab\ checked{\isacharparenright}{\kern0pt}\ \ \ \ \isanewline
\ \ \ \ \ \ \isanewline
\ \ \ \ \ \ \ \ \ \ {\isacharparenleft}{\kern0pt}{\isacharasterisk}{\kern0pt}val\ GLOBALVARS{\isacharunderscore}{\kern0pt}const\ {\isacharequal}{\kern0pt}\isanewline
\ \ \ \ \ \ \ \ \ \ \ \ Const\ {\isacharparenleft}{\kern0pt}{\isacharat}{\kern0pt}{\isacharbraceleft}{\kern0pt}const{\isacharunderscore}{\kern0pt}name\ GLOBALVARS{\isacharbraceright}{\kern0pt}{\isacharcomma}{\kern0pt}\ {\isacharat}{\kern0pt}{\isacharbraceleft}{\kern0pt}typ\ string{\isacharbraceright}{\kern0pt}\ {\isacharminus}{\kern0pt}{\isacharminus}{\kern0pt}{\isachargreater}{\kern0pt}\ procT{\isacharparenright}{\kern0pt}\isanewline
\ \ \ \ \ \ \ \ \ \ \isanewline
\ \ \ \ \ \ \ \ \ \ val\ lam{\isacharunderscore}{\kern0pt}GLOBALVARS\ {\isacharequal}{\kern0pt}\isanewline
\ \ \ \ \ \ \ \ \ \ \ \ Abs\ {\isacharparenleft}{\kern0pt}{\isachardoublequote}{\kern0pt}idx{\isachardoublequote}{\kern0pt}{\isacharcomma}{\kern0pt}\ {\isacharat}{\kern0pt}{\isacharbraceleft}{\kern0pt}typ\ string{\isacharbraceright}{\kern0pt}{\isacharcomma}{\kern0pt}\ GLOBALVARS{\isacharunderscore}{\kern0pt}const\ {\isachardollar}{\kern0pt}\ Bound\ {\isadigit{0}}{\isacharparenright}{\kern0pt}{\isacharasterisk}{\kern0pt}{\isacharparenright}{\kern0pt}\isanewline
\isanewline
\ \ \ \ \ \ \ \ \ \ val\ GLOBALVARS{\isacharunderscore}{\kern0pt}partial\ {\isacharequal}{\kern0pt}\ Const\ {\isacharparenleft}{\kern0pt}{\isacharat}{\kern0pt}{\isacharbraceleft}{\kern0pt}const{\isacharunderscore}{\kern0pt}name\ GLOBALVARS{\isacharbraceright}{\kern0pt}{\isacharcomma}{\kern0pt}\ {\isacharat}{\kern0pt}{\isacharbraceleft}{\kern0pt}typ\ {\isasymsigma}{\isacharbraceright}{\kern0pt}\ {\isacharminus}{\kern0pt}{\isacharminus}{\kern0pt}{\isachargreater}{\kern0pt}\ {\isacharat}{\kern0pt}{\isacharbraceleft}{\kern0pt}typ\ string{\isacharbraceright}{\kern0pt}\ {\isacharminus}{\kern0pt}{\isacharminus}{\kern0pt}{\isachargreater}{\kern0pt}\ procT{\isacharparenright}{\kern0pt}\isanewline
\ \ \ \ \ \ \ \ \ \ val\ lam{\isacharunderscore}{\kern0pt}GLOBALVARS\ {\isacharequal}{\kern0pt}\ Abs\ {\isacharparenleft}{\kern0pt}{\isachardoublequote}{\kern0pt}idx{\isachardoublequote}{\kern0pt}{\isacharcomma}{\kern0pt}\ {\isacharat}{\kern0pt}{\isacharbraceleft}{\kern0pt}typ\ string{\isacharbraceright}{\kern0pt}{\isacharcomma}{\kern0pt}\ GLOBALVARS{\isacharunderscore}{\kern0pt}partial\ {\isachardollar}{\kern0pt}\ sigma{\isadigit{0}}{\isacharunderscore}{\kern0pt}term\ {\isachardollar}{\kern0pt}\ Bound\ {\isadigit{0}}{\isacharparenright}{\kern0pt}\isanewline
\isanewline
\ \ \ \ \ \ \ \ \ \ val\ temp{\isacharunderscore}{\kern0pt}test\ {\isacharequal}{\kern0pt}\ Cspm{\isacharunderscore}{\kern0pt}API{\isachardot}{\kern0pt}mk{\isacharunderscore}{\kern0pt}MultiInter\ globals{\isacharunderscore}{\kern0pt}mset{\isacharunderscore}{\kern0pt}term\ lam{\isacharunderscore}{\kern0pt}GLOBALVARS\isanewline
\ \ \ \ \ \ \ \ \ \ \ val\ {\isacharunderscore}{\kern0pt}\ {\isacharequal}{\kern0pt}\ writeln\ {\isacharparenleft}{\kern0pt}{\isachardoublequote}{\kern0pt}{\isacharbrackleft}{\kern0pt}SYSTEM{\isacharbrackright}{\kern0pt}\ \ global\ test\ {\isacharcolon}{\kern0pt}{\isacharbackslash}{\kern0pt}n\ \ {\isachardoublequote}{\kern0pt}\ {\isacharcircum}{\kern0pt}\isanewline
\ \ \ \ \ \ \ \ \ \ \ \ \ \ \ \ \ \ \ \ \ \ \ \ \ \ \ \ \ \ \ \ \ \ \ \ \ Syntax{\isachardot}{\kern0pt}string{\isacharunderscore}{\kern0pt}of{\isacharunderscore}{\kern0pt}term{\isacharunderscore}{\kern0pt}global\ thy{\isacharprime}{\kern0pt}\ temp{\isacharunderscore}{\kern0pt}test{\isacharparenright}{\kern0pt}\ \isanewline
val\ {\isacharunderscore}{\kern0pt}\ {\isacharequal}{\kern0pt}\ tracing\ {\isacharparenleft}{\kern0pt}{\isachardoublequote}{\kern0pt}{\isacharbackslash}{\kern0pt}n\ {\isacharbrackleft}{\kern0pt}DEBUG{\isacharbrackright}{\kern0pt}\ Type\ full{\isacharunderscore}{\kern0pt}term\ {\isacharequal}{\kern0pt}\ {\isachardoublequote}{\kern0pt}\ {\isacharcircum}{\kern0pt}\isanewline
\ \ \ \ \ \ \ \ \ \ \ \ \ \ \ \ \ \ \ \ \ \ \ \ \ \ \ \ \ \ \ \ \ Syntax{\isachardot}{\kern0pt}string{\isacharunderscore}{\kern0pt}of{\isacharunderscore}{\kern0pt}typ{\isacharunderscore}{\kern0pt}global\ {\isacharat}{\kern0pt}{\isacharbraceleft}{\kern0pt}theory{\isacharbraceright}{\kern0pt}\ {\isacharparenleft}{\kern0pt}fastype{\isacharunderscore}{\kern0pt}of\ {\isacharparenleft}{\kern0pt}temp{\isacharunderscore}{\kern0pt}test{\isacharparenright}{\kern0pt}\ {\isacharparenright}{\kern0pt}{\isacharparenright}{\kern0pt}\ \ \ \ \ \ \ \isanewline
\isanewline
\isanewline
\isanewline
\ \ \ \ \ \ \ \ \ \ val\ processes\ {\isacharcolon}{\kern0pt}\ term\ list\ {\isacharequal}{\kern0pt}\isanewline
\ \ \ \ \ \ \ \ \ \ \ \ map\ {\isacharparenleft}{\kern0pt}fn\ name\ {\isacharequal}{\kern0pt}{\isachargreater}{\kern0pt}\ lam{\isacharunderscore}{\kern0pt}GLOBALVARS\ {\isachardollar}{\kern0pt}\ name{\isacharparenright}{\kern0pt}\ globals{\isacharunderscore}{\kern0pt}names\isanewline
\ \ \ \ \ \ \ \ \ \ {\isacharparenleft}{\kern0pt}{\isacharasterisk}{\kern0pt}\ globals{\isacharunderscore}{\kern0pt}names\ {\isacharcolon}{\kern0pt}\ term\ list\ de\ strings\ {\isacharasterisk}{\kern0pt}{\isacharparenright}{\kern0pt}\isanewline
\ \ \ \ \ \ \ \ \ \ \isanewline
\ \ \ \ \ \ \ \ \ \ val\ globals{\isacharunderscore}{\kern0pt}net{\isacharunderscore}{\kern0pt}term\ {\isacharcolon}{\kern0pt}\ term\ option\ {\isacharequal}{\kern0pt}\isanewline
\ \ \ \ \ \ \ \ \ \ \ \ {\isacharparenleft}{\kern0pt}case\ processes\ of\isanewline
\ \ \ \ \ \ \ \ \ \ \ \ \ \ \ {\isacharbrackleft}{\kern0pt}{\isacharbrackright}{\kern0pt}\ {\isacharequal}{\kern0pt}{\isachargreater}{\kern0pt}\ NONE\isanewline
\ \ \ \ \ \ \ \ \ \ \ \ \ {\isacharbar}{\kern0pt}\ {\isacharbrackleft}{\kern0pt}t{\isacharbrackright}{\kern0pt}\ {\isacharequal}{\kern0pt}{\isachargreater}{\kern0pt}\ SOME\ t\isanewline
\ \ \ \ \ \ \ \ \ \ \ \ \ {\isacharbar}{\kern0pt}\ ts\ {\isacharequal}{\kern0pt}{\isachargreater}{\kern0pt}\ SOME\ {\isacharparenleft}{\kern0pt}foldr{\isadigit{1}}\ {\isacharparenleft}{\kern0pt}fn\ {\isacharparenleft}{\kern0pt}t{\isadigit{1}}{\isacharcomma}{\kern0pt}\ t{\isadigit{2}}{\isacharparenright}{\kern0pt}\ {\isacharequal}{\kern0pt}{\isachargreater}{\kern0pt}\ Cspm{\isacharunderscore}{\kern0pt}API{\isachardot}{\kern0pt}mk{\isacharunderscore}{\kern0pt}Inter\ t{\isadigit{1}}\ t{\isadigit{2}}{\isacharparenright}{\kern0pt}\ ts{\isacharparenright}{\kern0pt}{\isacharparenright}{\kern0pt}{\isacharasterisk}{\kern0pt}{\isacharparenright}{\kern0pt}\isanewline
\isanewline
\ \ \ \ \ \ \ \ \ \ {\isacharparenleft}{\kern0pt}{\isacharasterisk}{\kern0pt}\ {\isadigit{2}}{\isachardot}{\kern0pt}\ Construction\ du\ réseau\ GLOBALVARS\ compactifié\ {\isacharminus}{\kern0pt}{\isacharminus}{\kern0pt}{\isacharminus}{\kern0pt}{\isacharminus}{\kern0pt}{\isacharminus}{\kern0pt}{\isacharminus}{\kern0pt}{\isacharminus}{\kern0pt}{\isacharminus}{\kern0pt}{\isacharminus}{\kern0pt}{\isacharminus}{\kern0pt}{\isacharminus}{\kern0pt}{\isacharminus}{\kern0pt}{\isacharminus}{\kern0pt}{\isacharminus}{\kern0pt}{\isacharminus}{\kern0pt}{\isacharminus}{\kern0pt}{\isacharminus}{\kern0pt}{\isacharminus}{\kern0pt}{\isacharminus}{\kern0pt}{\isacharminus}{\kern0pt}{\isacharminus}{\kern0pt}{\isacharminus}{\kern0pt}{\isacharminus}{\kern0pt}{\isacharminus}{\kern0pt}{\isacharminus}{\kern0pt}\ {\isacharasterisk}{\kern0pt}{\isacharparenright}{\kern0pt}\isanewline
\ \ \ \ \ \ \ \ \ \ val\ globals{\isacharunderscore}{\kern0pt}names\ {\isacharcolon}{\kern0pt}\ term\ list\ {\isacharequal}{\kern0pt}\isanewline
\ \ \ \ \ \ \ \ \ \ \ \ map\ {\isacharparenleft}{\kern0pt}fn\ {\isacharparenleft}{\kern0pt}b{\isacharcomma}{\kern0pt}\ {\isacharunderscore}{\kern0pt}{\isacharcomma}{\kern0pt}\ {\isacharunderscore}{\kern0pt}{\isacharcomma}{\kern0pt}\ {\isacharunderscore}{\kern0pt}{\isacharparenright}{\kern0pt}\ {\isacharequal}{\kern0pt}{\isachargreater}{\kern0pt}\ HOLogic{\isachardot}{\kern0pt}mk{\isacharunderscore}{\kern0pt}string\ {\isacharparenleft}{\kern0pt}Binding{\isachardot}{\kern0pt}name{\isacharunderscore}{\kern0pt}of\ b{\isacharparenright}{\kern0pt}{\isacharparenright}{\kern0pt}\ {\isacharparenleft}{\kern0pt}{\isacharhash}{\kern0pt}globals{\isacharunderscore}{\kern0pt}decls\ checked{\isacharparenright}{\kern0pt}\isanewline
\ \ \ \ \ \ \ \ \ \ \isanewline
\ \ \ \ \ \ \ \ \ \ val\ globals{\isacharunderscore}{\kern0pt}list{\isacharunderscore}{\kern0pt}term\ {\isacharequal}{\kern0pt}\isanewline
\ \ \ \ \ \ \ \ \ \ \ \ HOLogic{\isachardot}{\kern0pt}mk{\isacharunderscore}{\kern0pt}list\ {\isacharat}{\kern0pt}{\isacharbraceleft}{\kern0pt}typ\ string{\isacharbraceright}{\kern0pt}\ globals{\isacharunderscore}{\kern0pt}names\isanewline
\ \ \ \ \ \ \ \ \ \ \isanewline
\ \ \ \ \ \ \ \ \ \ val\ globals{\isacharunderscore}{\kern0pt}mset{\isacharunderscore}{\kern0pt}term\ {\isacharequal}{\kern0pt}\isanewline
\ \ \ \ \ \ \ \ \ \ \ \ Cspm{\isacharunderscore}{\kern0pt}API{\isachardot}{\kern0pt}mk{\isacharunderscore}{\kern0pt}mset\ globals{\isacharunderscore}{\kern0pt}list{\isacharunderscore}{\kern0pt}term\isanewline
\ \ \ \ \ \ \ \ \ \ \isanewline
\ \ \ \ \ \ \ \ \ \ val\ procT\ {\isacharequal}{\kern0pt}\ Cspm{\isacharunderscore}{\kern0pt}API{\isachardot}{\kern0pt}mk{\isacharunderscore}{\kern0pt}processT\ {\isacharat}{\kern0pt}{\isacharbraceleft}{\kern0pt}typ\ evs{\isacharbraceright}{\kern0pt}\isanewline
\ \ \ \ \ \ \ \ \ \ val\ sigma{\isadigit{0}}{\isacharunderscore}{\kern0pt}term\ {\isacharequal}{\kern0pt}\ make{\isacharunderscore}{\kern0pt}sigma{\isacharunderscore}{\kern0pt}term\ {\isacharparenleft}{\kern0pt}{\isacharhash}{\kern0pt}varstab\ checked{\isacharparenright}{\kern0pt}\isanewline
\ \ \ \ \ \ \ \ \ \ \isanewline
\ \ \ \ \ \ \ \ \ \ val\ GLOBALVARS{\isacharunderscore}{\kern0pt}partial\ {\isacharequal}{\kern0pt}\ Const\ {\isacharparenleft}{\kern0pt}{\isacharat}{\kern0pt}{\isacharbraceleft}{\kern0pt}const{\isacharunderscore}{\kern0pt}name\ GLOBALVARS{\isacharbraceright}{\kern0pt}{\isacharcomma}{\kern0pt}\ {\isacharat}{\kern0pt}{\isacharbraceleft}{\kern0pt}typ\ {\isasymsigma}{\isacharbraceright}{\kern0pt}\ {\isacharminus}{\kern0pt}{\isacharminus}{\kern0pt}{\isachargreater}{\kern0pt}\ {\isacharat}{\kern0pt}{\isacharbraceleft}{\kern0pt}typ\ string{\isacharbraceright}{\kern0pt}\ {\isacharminus}{\kern0pt}{\isacharminus}{\kern0pt}{\isachargreater}{\kern0pt}\ procT{\isacharparenright}{\kern0pt}\isanewline
\ \ \ \ \ \ \ \ \ \ val\ lam{\isacharunderscore}{\kern0pt}GLOBALVARS\ {\isacharequal}{\kern0pt}\ Abs\ {\isacharparenleft}{\kern0pt}{\isachardoublequote}{\kern0pt}idx{\isachardoublequote}{\kern0pt}{\isacharcomma}{\kern0pt}\ {\isacharat}{\kern0pt}{\isacharbraceleft}{\kern0pt}typ\ string{\isacharbraceright}{\kern0pt}{\isacharcomma}{\kern0pt}\ GLOBALVARS{\isacharunderscore}{\kern0pt}partial\ {\isachardollar}{\kern0pt}\ sigma{\isadigit{0}}{\isacharunderscore}{\kern0pt}term\ {\isachardollar}{\kern0pt}\ Bound\ {\isadigit{0}}{\isacharparenright}{\kern0pt}\isanewline
\ \ \ \ \ \ \ \ \ \ \isanewline
\ \ \ \ \ \ \ \ \ \ {\isacharparenleft}{\kern0pt}{\isacharasterisk}{\kern0pt}\ Utilisation\ de\ MultiSync\ avec\ ensemble\ d{\isacharprime}{\kern0pt}événements\ vide\ pour\ compactifier\ {\isacharasterisk}{\kern0pt}{\isacharparenright}{\kern0pt}\isanewline
\ \ \ \ \ \ \ \ \ \ val\ globals{\isacharunderscore}{\kern0pt}net{\isacharunderscore}{\kern0pt}term\ {\isacharcolon}{\kern0pt}\ term\ option\ {\isacharequal}{\kern0pt}\isanewline
\ \ \ \ \ \ \ \ \ \ \ \ {\isacharparenleft}{\kern0pt}case\ globals{\isacharunderscore}{\kern0pt}names\ of\isanewline
\ \ \ \ \ \ \ \ \ \ \ \ \ \ \ {\isacharbrackleft}{\kern0pt}{\isacharbrackright}{\kern0pt}\ {\isacharequal}{\kern0pt}{\isachargreater}{\kern0pt}\ NONE\isanewline
\ \ \ \ \ \ \ \ \ \ \ \ \ {\isacharbar}{\kern0pt}\ {\isacharunderscore}{\kern0pt}\ {\isacharequal}{\kern0pt}{\isachargreater}{\kern0pt}\ \isanewline
\ \ \ \ \ \ \ \ \ \ \ \ \ \ \ let\isanewline
\ \ \ \ \ \ \ \ \ \ \ \ \ \ \ \ \ val\ multisync{\isacharunderscore}{\kern0pt}const\ {\isacharequal}{\kern0pt}\ Const\ {\isacharparenleft}{\kern0pt}{\isacharat}{\kern0pt}{\isacharbraceleft}{\kern0pt}const{\isacharunderscore}{\kern0pt}name\ MultiSync{\isacharbraceright}{\kern0pt}{\isacharcomma}{\kern0pt}\ \isanewline
\ \ \ \ \ \ \ \ \ \ \ \ \ \ \ \ \ \ \ \ \ \ \ \ \ \ \ \ \ \ \ \ \ \ \ \ \ \ \ \ \ \ \ \ \ {\isacharat}{\kern0pt}{\isacharbraceleft}{\kern0pt}typ\ {\isachardoublequote}{\kern0pt}evs\ set{\isachardoublequote}{\kern0pt}{\isacharbraceright}{\kern0pt}\ {\isacharminus}{\kern0pt}{\isacharminus}{\kern0pt}{\isachargreater}{\kern0pt}\ \isanewline
\ \ \ \ \ \ \ \ \ \ \ \ \ \ \ \ \ \ \ \ \ \ \ \ \ \ \ \ \ \ \ \ \ \ \ \ \ \ \ \ \ \ \ \ \ {\isacharat}{\kern0pt}{\isacharbraceleft}{\kern0pt}typ\ {\isachardoublequote}{\kern0pt}string\ multiset{\isachardoublequote}{\kern0pt}{\isacharbraceright}{\kern0pt}\ {\isacharminus}{\kern0pt}{\isacharminus}{\kern0pt}{\isachargreater}{\kern0pt}\ \isanewline
\ \ \ \ \ \ \ \ \ \ \ \ \ \ \ \ \ \ \ \ \ \ \ \ \ \ \ \ \ \ \ \ \ \ \ \ \ \ \ \ \ \ \ \ \ {\isacharat}{\kern0pt}{\isacharbraceleft}{\kern0pt}typ\ {\isachardoublequote}{\kern0pt}string\ {\isasymRightarrow}\ {\isacharparenleft}{\kern0pt}evs{\isacharparenright}{\kern0pt}\ process{\isachardoublequote}{\kern0pt}{\isacharbraceright}{\kern0pt}\ {\isacharminus}{\kern0pt}{\isacharminus}{\kern0pt}{\isachargreater}{\kern0pt}\ \isanewline
\ \ \ \ \ \ \ \ \ \ \ \ \ \ \ \ \ \ \ \ \ \ \ \ \ \ \ \ \ \ \ \ \ \ \ \ \ \ \ \ \ \ \ \ \ {\isacharat}{\kern0pt}{\isacharbraceleft}{\kern0pt}typ\ {\isachardoublequote}{\kern0pt}{\isacharparenleft}{\kern0pt}evs{\isacharparenright}{\kern0pt}\ process{\isachardoublequote}{\kern0pt}{\isacharbraceright}{\kern0pt}{\isacharparenright}{\kern0pt}\isanewline
\ \ \ \ \ \ \ \ \ \ \ \ \ \ \ \ \ val\ compact{\isacharunderscore}{\kern0pt}term\ {\isacharequal}{\kern0pt}\ multisync{\isacharunderscore}{\kern0pt}const\ {\isachardollar}{\kern0pt}\ empty{\isacharunderscore}{\kern0pt}evs{\isacharunderscore}{\kern0pt}set\ {\isachardollar}{\kern0pt}\ globals{\isacharunderscore}{\kern0pt}mset{\isacharunderscore}{\kern0pt}term\ {\isachardollar}{\kern0pt}\ lam{\isacharunderscore}{\kern0pt}GLOBALVARS\isanewline
\ \ \ \ \ \ \ \ \ \ \ \ \ \ \ in\isanewline
\ \ \ \ \ \ \ \ \ \ \ \ \ \ \ \ \ SOME\ compact{\isacharunderscore}{\kern0pt}term\isanewline
\ \ \ \ \ \ \ \ \ \ \ \ \ \ \ end{\isacharparenright}{\kern0pt}\isanewline
\ \ \ \ \ \ \ \ \ \ \isanewline
\ \ \ \ \ \ \ \ \ \ val\ {\isacharunderscore}{\kern0pt}\ {\isacharequal}{\kern0pt}\ {\isacharparenleft}{\kern0pt}case\ globals{\isacharunderscore}{\kern0pt}net{\isacharunderscore}{\kern0pt}term\ of\isanewline
\ \ \ \ \ \ \ \ \ \ \ \ \ \ \ \ \ \ \ \ SOME\ t\ {\isacharequal}{\kern0pt}{\isachargreater}{\kern0pt}\ writeln\ {\isacharparenleft}{\kern0pt}{\isachardoublequote}{\kern0pt}{\isacharbrackleft}{\kern0pt}SYSTEM{\isacharbrackright}{\kern0pt}\ réseau\ GLOBALVARS\ compactifié\ {\isacharcolon}{\kern0pt}{\isacharbackslash}{\kern0pt}n\ \ {\isachardoublequote}{\kern0pt}\ {\isacharcircum}{\kern0pt}\isanewline
\ \ \ \ \ \ \ \ \ \ \ \ \ \ \ \ \ \ \ \ \ \ \ \ \ \ \ \ \ \ \ \ \ \ \ \ \ \ Syntax{\isachardot}{\kern0pt}string{\isacharunderscore}{\kern0pt}of{\isacharunderscore}{\kern0pt}term{\isacharunderscore}{\kern0pt}global\ thy{\isacharprime}{\kern0pt}\ t{\isacharparenright}{\kern0pt}\isanewline
\ \ \ \ \ \ \ \ \ \ \ \ \ \ \ \ \ \ {\isacharbar}{\kern0pt}\ NONE\ {\isacharequal}{\kern0pt}{\isachargreater}{\kern0pt}\ writeln\ {\isacharparenleft}{\kern0pt}{\isachardoublequote}{\kern0pt}{\isacharbrackleft}{\kern0pt}SYSTEM{\isacharbrackright}{\kern0pt}\ Aucune\ variable\ globale\ déclarée{\isachardoublequote}{\kern0pt}{\isacharparenright}{\kern0pt}{\isacharparenright}{\kern0pt}\isanewline
\ \ \ \ \ \ \ \ \ \ \isanewline
\ \ \ \ \ \ \ \ \ \ \isanewline
\isanewline
\ \ \ \ \ \ \ \ \ \ {\isacharparenleft}{\kern0pt}{\isacharasterisk}{\kern0pt}\ {\isadigit{3}}{\isachardot}{\kern0pt}\ \ Construction\ du\ réseau\ SEMAPHORES\ {\isacharminus}{\kern0pt}{\isacharminus}{\kern0pt}{\isacharminus}{\kern0pt}{\isacharminus}{\kern0pt}{\isacharminus}{\kern0pt}{\isacharminus}{\kern0pt}{\isacharminus}{\kern0pt}{\isacharminus}{\kern0pt}{\isacharminus}{\kern0pt}{\isacharminus}{\kern0pt}{\isacharminus}{\kern0pt}{\isacharminus}{\kern0pt}{\isacharminus}{\kern0pt}{\isacharminus}{\kern0pt}{\isacharminus}{\kern0pt}{\isacharminus}{\kern0pt}{\isacharminus}{\kern0pt}{\isacharminus}{\kern0pt}{\isacharminus}{\kern0pt}{\isacharminus}{\kern0pt}{\isacharminus}{\kern0pt}{\isacharminus}{\kern0pt}{\isacharminus}{\kern0pt}{\isacharminus}{\kern0pt}{\isacharminus}{\kern0pt}{\isacharminus}{\kern0pt}{\isacharminus}{\kern0pt}{\isacharminus}{\kern0pt}{\isacharminus}{\kern0pt}{\isacharminus}{\kern0pt}{\isacharminus}{\kern0pt}{\isacharminus}{\kern0pt}\ {\isacharasterisk}{\kern0pt}{\isacharparenright}{\kern0pt}\isanewline
\ \ \ \ \ \ \ \ \ \ \isanewline
\ \ \ \ \ \ \ \ \ \ {\isacharparenleft}{\kern0pt}{\isacharasterisk}{\kern0pt}\ extraire\ les\ noms\ de\ verrous\ depuis\ checked\ {\isacharasterisk}{\kern0pt}{\isacharparenright}{\kern0pt}\isanewline
\ \ \ \ \ \ \ \ \ \ \isanewline
\ \ \ \ \ \ \ \ \ \ val\ lock{\isacharunderscore}{\kern0pt}ids{\isacharunderscore}{\kern0pt}terms\ {\isacharcolon}{\kern0pt}\ term\ list\ {\isacharequal}{\kern0pt}\isanewline
\ \ \ \ \ \ \ \ \ \ \ \ map\ {\isacharparenleft}{\kern0pt}fn\ {\isacharparenleft}{\kern0pt}b{\isacharcomma}{\kern0pt}\ {\isacharunderscore}{\kern0pt}{\isacharcomma}{\kern0pt}\ {\isacharunderscore}{\kern0pt}{\isacharparenright}{\kern0pt}\ {\isacharequal}{\kern0pt}{\isachargreater}{\kern0pt}\isanewline
\ \ \ \ \ \ \ \ \ \ \ \ \ \ \ \ \ \ \ HOLogic{\isachardot}{\kern0pt}mk{\isacharunderscore}{\kern0pt}number\ {\isacharat}{\kern0pt}{\isacharbraceleft}{\kern0pt}typ\ int{\isacharbraceright}{\kern0pt}\ {\isacharparenleft}{\kern0pt}get{\isacharunderscore}{\kern0pt}lock{\isacharunderscore}{\kern0pt}id\ {\isacharparenleft}{\kern0pt}Binding{\isachardot}{\kern0pt}name{\isacharunderscore}{\kern0pt}of\ b{\isacharparenright}{\kern0pt}\ checked{\isacharparenright}{\kern0pt}{\isacharparenright}{\kern0pt}\ \ \isanewline
\ \ \ \ \ \ \ \ \ \ \ \ \ \ \ \ \ \ \ {\isacharparenleft}{\kern0pt}{\isacharasterisk}{\kern0pt}\ get{\isacharunderscore}{\kern0pt}lock{\isacharunderscore}{\kern0pt}id\ {\isacharcolon}{\kern0pt}\ binding\ {\isacharminus}{\kern0pt}{\isachargreater}{\kern0pt}\ int\ {\isacharasterisk}{\kern0pt}{\isacharparenright}{\kern0pt}\isanewline
\ \ \ \ \ \ \ \ \ \ \ \ \ \ \ \ \ \ \ \ \ \ \ {\isacharparenleft}{\kern0pt}{\isacharhash}{\kern0pt}locks{\isacharunderscore}{\kern0pt}decls\ checked{\isacharparenright}{\kern0pt}\isanewline
\ \ \ \ \ \ \ \ \ \ \isanewline
\ \ \ \ \ \ \ \ \ \ val\ locks{\isacharunderscore}{\kern0pt}list{\isacharunderscore}{\kern0pt}term\ {\isacharequal}{\kern0pt}\isanewline
\ \ \ \ \ \ \ \ \ \ \ \ HOLogic{\isachardot}{\kern0pt}mk{\isacharunderscore}{\kern0pt}list\ {\isacharat}{\kern0pt}{\isacharbraceleft}{\kern0pt}typ\ int{\isacharbraceright}{\kern0pt}\ lock{\isacharunderscore}{\kern0pt}ids{\isacharunderscore}{\kern0pt}terms\isanewline
\ \ \ \ \ \ \ \ \ \ \isanewline
\ \ \ \ \ \ \ \ \ \ val\ locks{\isacharunderscore}{\kern0pt}mset{\isacharunderscore}{\kern0pt}term\ {\isacharequal}{\kern0pt}\isanewline
\ \ \ \ \ \ \ \ \ \ \ \ Cspm{\isacharunderscore}{\kern0pt}API{\isachardot}{\kern0pt}mk{\isacharunderscore}{\kern0pt}mset\ locks{\isacharunderscore}{\kern0pt}list{\isacharunderscore}{\kern0pt}term\ \ {\isacharparenleft}{\kern0pt}{\isacharasterisk}{\kern0pt}\ {\isacharbraceleft}{\kern0pt}{\isacharhash}{\kern0pt}\ {\isasymdots}\ {\isacharhash}{\kern0pt}{\isacharbraceright}{\kern0pt}\ {\isacharasterisk}{\kern0pt}{\isacharparenright}{\kern0pt}\isanewline
\ \ \ \ \ \ \ \ \ \ \isanewline
\ \ \ \ \ \ \ \ \ \ val\ SEMAPHORES{\isacharunderscore}{\kern0pt}const\ {\isacharequal}{\kern0pt}\isanewline
\ \ \ \ \ \ \ \ \ \ \ \ Const\ {\isacharparenleft}{\kern0pt}{\isacharat}{\kern0pt}{\isacharbraceleft}{\kern0pt}const{\isacharunderscore}{\kern0pt}name\ SEMAPHORES{\isacharbraceright}{\kern0pt}{\isacharcomma}{\kern0pt}\ {\isacharat}{\kern0pt}{\isacharbraceleft}{\kern0pt}typ\ int{\isacharbraceright}{\kern0pt}\ {\isacharminus}{\kern0pt}{\isacharminus}{\kern0pt}{\isachargreater}{\kern0pt}\ procT{\isacharparenright}{\kern0pt}\isanewline
\isanewline
\ \ \ \ \ \ \ \ \ val\ lam{\isacharunderscore}{\kern0pt}SEMAPHORES\ {\isacharequal}{\kern0pt}\isanewline
\ \ \ \ \ \ \ \ \ \ \ \ \ \ \ \ \ \ \ \ Abs\ {\isacharparenleft}{\kern0pt}{\isachardoublequote}{\kern0pt}idx{\isachardoublequote}{\kern0pt}{\isacharcomma}{\kern0pt}\ {\isacharat}{\kern0pt}{\isacharbraceleft}{\kern0pt}typ\ int{\isacharbraceright}{\kern0pt}{\isacharcomma}{\kern0pt}\ SEMAPHORES{\isacharunderscore}{\kern0pt}const\ {\isachardollar}{\kern0pt}\ Bound\ {\isadigit{0}}{\isacharparenright}{\kern0pt}\isanewline
\ \ \ \ \ \ \ \ val\ processes\ {\isacharcolon}{\kern0pt}\ term\ list\ {\isacharequal}{\kern0pt}\ map\ {\isacharparenleft}{\kern0pt}fn\ name\ {\isacharequal}{\kern0pt}{\isachargreater}{\kern0pt}\ lam{\isacharunderscore}{\kern0pt}SEMAPHORES\ {\isachardollar}{\kern0pt}\ name{\isacharparenright}{\kern0pt}\ lock{\isacharunderscore}{\kern0pt}ids{\isacharunderscore}{\kern0pt}terms\isanewline
\ \ \ \ \ \ \ \ \ \ {\isacharparenleft}{\kern0pt}{\isacharasterisk}{\kern0pt}\ globals{\isacharunderscore}{\kern0pt}names\ {\isacharcolon}{\kern0pt}\ term\ list\ de\ strings\ {\isacharasterisk}{\kern0pt}{\isacharparenright}{\kern0pt}\isanewline
\ \ \ \ \ \ \ \ \ \ \isanewline
\ \ \ \ \ \ \ {\isacharparenleft}{\kern0pt}{\isacharasterisk}{\kern0pt}\ \ val\ semaphores{\isacharunderscore}{\kern0pt}net{\isacharunderscore}{\kern0pt}term\ {\isacharcolon}{\kern0pt}\ term\ option\ {\isacharequal}{\kern0pt}\isanewline
\ \ \ \ \ \ \ \ \ \ \ \ \ \ case\ processes\ of\isanewline
\ \ \ \ \ \ \ \ \ \ \ \ \ \ \ \ {\isacharbrackleft}{\kern0pt}{\isacharbrackright}{\kern0pt}\ {\isacharequal}{\kern0pt}{\isachargreater}{\kern0pt}\ NONE\isanewline
\ \ \ \ \ \ \ \ \ \ \ \ \ \ {\isacharbar}{\kern0pt}\ {\isacharbrackleft}{\kern0pt}t{\isacharbrackright}{\kern0pt}\ {\isacharequal}{\kern0pt}{\isachargreater}{\kern0pt}\ SOME\ t\isanewline
\ \ \ \ \ \ \ \ \ \ \ \ \ \ {\isacharbar}{\kern0pt}\ ts\ {\isacharequal}{\kern0pt}{\isachargreater}{\kern0pt}\ SOME\ {\isacharparenleft}{\kern0pt}foldr{\isadigit{1}}\ {\isacharparenleft}{\kern0pt}fn\ {\isacharparenleft}{\kern0pt}t{\isadigit{1}}{\isacharcomma}{\kern0pt}\ t{\isadigit{2}}{\isacharparenright}{\kern0pt}\ {\isacharequal}{\kern0pt}{\isachargreater}{\kern0pt}\ Cspm{\isacharunderscore}{\kern0pt}API{\isachardot}{\kern0pt}mk{\isacharunderscore}{\kern0pt}Inter\ t{\isadigit{1}}\ t{\isadigit{2}}{\isacharparenright}{\kern0pt}\ ts{\isacharparenright}{\kern0pt}{\isacharasterisk}{\kern0pt}{\isacharparenright}{\kern0pt}\isanewline
\isanewline
\ \isanewline
\ \ \ \ \ \ \ \ \ \ val\ semaphores{\isacharunderscore}{\kern0pt}net{\isacharunderscore}{\kern0pt}term\ {\isacharcolon}{\kern0pt}\ term\ option\ {\isacharequal}{\kern0pt}\isanewline
\ \ \ \ \ \ \ \ \ \ \ \ {\isacharparenleft}{\kern0pt}case\ lock{\isacharunderscore}{\kern0pt}ids{\isacharunderscore}{\kern0pt}terms\ of\isanewline
\ \ \ \ \ \ \ \ \ \ \ \ \ \ \ {\isacharbrackleft}{\kern0pt}{\isacharbrackright}{\kern0pt}\ {\isacharequal}{\kern0pt}{\isachargreater}{\kern0pt}\ NONE\isanewline
\ \ \ \ \ \ \ \ \ \ \ \ \ {\isacharbar}{\kern0pt}\ {\isacharunderscore}{\kern0pt}\ {\isacharequal}{\kern0pt}{\isachargreater}{\kern0pt}\isanewline
\ \ \ \ \ \ \ \ \ \ \ \ \ \ \ let\isanewline
\ \ \ \ \ \ \ \ \ \ \ \ \ \ \ \ \ val\ empty{\isacharunderscore}{\kern0pt}evs{\isacharunderscore}{\kern0pt}set\ {\isacharequal}{\kern0pt}\ Const\ {\isacharparenleft}{\kern0pt}{\isacharat}{\kern0pt}{\isacharbraceleft}{\kern0pt}const{\isacharunderscore}{\kern0pt}name\ bot{\isacharbraceright}{\kern0pt}{\isacharcomma}{\kern0pt}\ {\isacharat}{\kern0pt}{\isacharbraceleft}{\kern0pt}typ\ {\isachardoublequote}{\kern0pt}evs\ set{\isachardoublequote}{\kern0pt}{\isacharbraceright}{\kern0pt}{\isacharparenright}{\kern0pt}\isanewline
\ \ \ \ \ \ \ \ \ \ \ \ \ \ \ \ \ val\ multisync{\isacharunderscore}{\kern0pt}const\ {\isacharequal}{\kern0pt}\ Const\ {\isacharparenleft}{\kern0pt}{\isacharat}{\kern0pt}{\isacharbraceleft}{\kern0pt}const{\isacharunderscore}{\kern0pt}name\ MultiSync{\isacharbraceright}{\kern0pt}{\isacharcomma}{\kern0pt}\ \isanewline
\ \ \ \ \ \ \ \ \ \ \ \ \ \ \ \ \ \ \ \ \ \ \ \ \ \ \ \ \ \ \ \ \ \ \ \ \ \ \ \ \ \ \ \ \ \ {\isacharat}{\kern0pt}{\isacharbraceleft}{\kern0pt}typ\ {\isachardoublequote}{\kern0pt}evs\ set{\isachardoublequote}{\kern0pt}{\isacharbraceright}{\kern0pt}\ {\isacharminus}{\kern0pt}{\isacharminus}{\kern0pt}{\isachargreater}{\kern0pt}\ \isanewline
\ \ \ \ \ \ \ \ \ \ \ \ \ \ \ \ \ \ \ \ \ \ \ \ \ \ \ \ \ \ \ \ \ \ \ \ \ \ \ \ \ \ \ \ \ \ {\isacharat}{\kern0pt}{\isacharbraceleft}{\kern0pt}typ\ {\isachardoublequote}{\kern0pt}int\ multiset{\isachardoublequote}{\kern0pt}{\isacharbraceright}{\kern0pt}\ {\isacharminus}{\kern0pt}{\isacharminus}{\kern0pt}{\isachargreater}{\kern0pt}\ \isanewline
\ \ \ \ \ \ \ \ \ \ \ \ \ \ \ \ \ \ \ \ \ \ \ \ \ \ \ \ \ \ \ \ \ \ \ \ \ \ \ \ \ \ \ \ \ \ {\isacharat}{\kern0pt}{\isacharbraceleft}{\kern0pt}typ\ {\isachardoublequote}{\kern0pt}int\ {\isasymRightarrow}\ {\isacharparenleft}{\kern0pt}evs{\isacharparenright}{\kern0pt}\ process{\isachardoublequote}{\kern0pt}{\isacharbraceright}{\kern0pt}\ {\isacharminus}{\kern0pt}{\isacharminus}{\kern0pt}{\isachargreater}{\kern0pt}\ \isanewline
\ \ \ \ \ \ \ \ \ \ \ \ \ \ \ \ \ \ \ \ \ \ \ \ \ \ \ \ \ \ \ \ \ \ \ \ \ \ \ \ \ \ \ \ \ \ {\isacharat}{\kern0pt}{\isacharbraceleft}{\kern0pt}typ\ {\isachardoublequote}{\kern0pt}{\isacharparenleft}{\kern0pt}evs{\isacharparenright}{\kern0pt}\ process{\isachardoublequote}{\kern0pt}{\isacharbraceright}{\kern0pt}{\isacharparenright}{\kern0pt}\isanewline
\ \ \ \ \ \ \ \ \ \ \ \ \ \ \ \ \ val\ compact{\isacharunderscore}{\kern0pt}term\ {\isacharequal}{\kern0pt}\ multisync{\isacharunderscore}{\kern0pt}const\ {\isachardollar}{\kern0pt}\ empty{\isacharunderscore}{\kern0pt}evs{\isacharunderscore}{\kern0pt}set\ {\isachardollar}{\kern0pt}\ locks{\isacharunderscore}{\kern0pt}mset{\isacharunderscore}{\kern0pt}term\ {\isachardollar}{\kern0pt}\ lam{\isacharunderscore}{\kern0pt}SEMAPHORES\isanewline
\ \ \ \ \ \ \ \ \ \ \ \ \ \ \ in\isanewline
\ \ \ \ \ \ \ \ \ \ \ \ \ \ \ \ \ SOME\ compact{\isacharunderscore}{\kern0pt}term\isanewline
\ \ \ \ \ \ \ \ \ \ \ \ \ \ \ end{\isacharparenright}{\kern0pt}\isanewline
\ \ \ \ \ \ \ \ \ \ \ \ \isanewline
\isanewline
\ \ \ \ \ \ \ \ \ \isanewline
\ \ \ \ \ \ \ \ \ \ \isanewline
\ \ \ \ \ \ \ \ \ \ {\isacharparenleft}{\kern0pt}{\isacharasterisk}{\kern0pt}val\ semaphores{\isacharunderscore}{\kern0pt}net{\isacharunderscore}{\kern0pt}term\ {\isacharequal}{\kern0pt}\ \isanewline
\ \ \ \ \ \ \ \ \ \ \ \ \ SOME{\isacharparenleft}{\kern0pt}\ Cspm{\isacharunderscore}{\kern0pt}API{\isachardot}{\kern0pt}mk{\isacharunderscore}{\kern0pt}MultiInter{\isacharunderscore}{\kern0pt}\ sumT\ locks{\isacharunderscore}{\kern0pt}mset{\isacharunderscore}{\kern0pt}term\ lam{\isacharunderscore}{\kern0pt}SEMAPHORES\ {\isacharparenright}{\kern0pt}\ {\isacharasterisk}{\kern0pt}{\isacharparenright}{\kern0pt}\ \ \ \ \ \ \ \ \ \ \isanewline
\isanewline
\ \ \ \ \ \ \ \ \ \ \ \ \ {\isacharparenleft}{\kern0pt}{\isacharasterisk}{\kern0pt}{\isacharbar}{\kern0pt}NONE\ {\isacharequal}{\kern0pt}{\isachargreater}{\kern0pt}\ writeln\ {\isacharparenleft}{\kern0pt}{\isachardoublequote}{\kern0pt}{\isacharbrackleft}{\kern0pt}SYSTEM{\isacharbrackright}{\kern0pt}\ \ réseau\ SEMAPHORES\ {\isacharcolon}{\kern0pt}{\isacharbackslash}{\kern0pt}n\ Vous\ n{\isacharprime}{\kern0pt}avez\ pas\ de\ Locks\ {\isacharbackslash}{\kern0pt}n{\isachardoublequote}{\kern0pt}{\isacharparenright}{\kern0pt}\ {\isacharasterisk}{\kern0pt}{\isacharparenright}{\kern0pt}\isanewline
\isanewline
\ \ \ \ \ \ \ \ \ \ {\isacharparenleft}{\kern0pt}{\isacharasterisk}{\kern0pt}\ {\isadigit{4}}{\isachardot}{\kern0pt}\ Construction\ du\ réseau\ des\ THREADS\ {\isacharminus}{\kern0pt}{\isacharminus}{\kern0pt}{\isacharminus}{\kern0pt}{\isacharminus}{\kern0pt}{\isacharminus}{\kern0pt}{\isacharminus}{\kern0pt}{\isacharminus}{\kern0pt}{\isacharminus}{\kern0pt}{\isacharminus}{\kern0pt}{\isacharminus}{\kern0pt}{\isacharminus}{\kern0pt}{\isacharminus}{\kern0pt}{\isacharminus}{\kern0pt}{\isacharminus}{\kern0pt}{\isacharminus}{\kern0pt}{\isacharminus}{\kern0pt}{\isacharminus}{\kern0pt}{\isacharminus}{\kern0pt}{\isacharminus}{\kern0pt}{\isacharminus}{\kern0pt}{\isacharminus}{\kern0pt}{\isacharminus}{\kern0pt}{\isacharminus}{\kern0pt}{\isacharminus}{\kern0pt}{\isacharminus}{\kern0pt}{\isacharminus}{\kern0pt}{\isacharminus}{\kern0pt}\ {\isacharasterisk}{\kern0pt}{\isacharparenright}{\kern0pt}\isanewline
\ \ \ \ \ \ \ \ \ {\isacharparenleft}{\kern0pt}{\isacharasterisk}{\kern0pt}\ val\ thread{\isacharunderscore}{\kern0pt}names\ {\isacharcolon}{\kern0pt}\ string\ list\ {\isacharequal}{\kern0pt}\isanewline
\ \ \ \ \ \ \ \ \ \ \ \ map\ {\isacharparenleft}{\kern0pt}fn\ {\isacharbraceleft}{\kern0pt}nom{\isacharunderscore}{\kern0pt}thread{\isacharcomma}{\kern0pt}\ {\isachardot}{\kern0pt}{\isachardot}{\kern0pt}{\isachardot}{\kern0pt}{\isacharbraceright}{\kern0pt}\ {\isacharequal}{\kern0pt}{\isachargreater}{\kern0pt}\ Binding{\isachardot}{\kern0pt}name{\isacharunderscore}{\kern0pt}of\ nom{\isacharunderscore}{\kern0pt}thread{\isacharparenright}{\kern0pt}\ {\isacharparenleft}{\kern0pt}{\isacharhash}{\kern0pt}threads{\isacharunderscore}{\kern0pt}decls\ checked{\isacharparenright}{\kern0pt}\isanewline
\isanewline
\ \ \ \ \ \ \ \ \ \ val\ thread{\isacharunderscore}{\kern0pt}consts\ {\isacharcolon}{\kern0pt}\ term\ list\ {\isacharequal}{\kern0pt}\isanewline
\ \ \ \ \ \ \ \ \ \ \ \ map\ {\isacharparenleft}{\kern0pt}fn\ name\ {\isacharequal}{\kern0pt}{\isachargreater}{\kern0pt}\isanewline
\ \ \ \ \ \ \ \ \ \ \ \ \ \ Const\ {\isacharparenleft}{\kern0pt}Sign{\isachardot}{\kern0pt}full{\isacharunderscore}{\kern0pt}name\ thy\ {\isacharparenleft}{\kern0pt}Binding{\isachardot}{\kern0pt}name\ {\isacharparenleft}{\kern0pt}name\ {\isacharcircum}{\kern0pt}\ {\isachardoublequote}{\kern0pt}{\isacharunderscore}{\kern0pt}cmd{\isachardoublequote}{\kern0pt}{\isacharparenright}{\kern0pt}{\isacharparenright}{\kern0pt}{\isacharcomma}{\kern0pt}\ {\isacharat}{\kern0pt}{\isacharbraceleft}{\kern0pt}typ\ com{\isacharbraceright}{\kern0pt}{\isacharparenright}{\kern0pt}{\isacharparenright}{\kern0pt}\isanewline
\ \ \ \ \ \ \ \ \ \ \ \ thread{\isacharunderscore}{\kern0pt}names\isanewline
\isanewline
\isanewline
\ \ \ \ \ \ \ \ \ \ val\ thread{\isacharunderscore}{\kern0pt}processes\ {\isacharcolon}{\kern0pt}\ term\ list\ {\isacharequal}{\kern0pt}\isanewline
\ \ \ \ \ \ \ \ \ \ \ \ map\ {\isacharparenleft}{\kern0pt}fn\ t\ {\isacharequal}{\kern0pt}{\isachargreater}{\kern0pt}\ sem{\isacharunderscore}{\kern0pt}const\ {\isachardollar}{\kern0pt}\ t{\isacharparenright}{\kern0pt}\ thread{\isacharunderscore}{\kern0pt}consts\isanewline
\isanewline
\ \ \ \ \ \ \ \ \ \ val\ thread{\isacharunderscore}{\kern0pt}net{\isacharunderscore}{\kern0pt}term\ {\isacharequal}{\kern0pt}\isanewline
\ \ \ \ \ \ \ \ \ \ \ \ {\isacharparenleft}{\kern0pt}case\ thread{\isacharunderscore}{\kern0pt}processes\ of\isanewline
\ \ \ \ \ \ \ \ \ \ \ \ \ \ \ \ {\isacharbrackleft}{\kern0pt}{\isacharbrackright}{\kern0pt}\ {\isacharequal}{\kern0pt}{\isachargreater}{\kern0pt}\ error\ {\isachardoublequote}{\kern0pt}No\ threads\ declared{\isachardoublequote}{\kern0pt}\isanewline
\ \ \ \ \ \ \ \ \ \ \ \ \ \ {\isacharbar}{\kern0pt}\ {\isacharbrackleft}{\kern0pt}t{\isacharbrackright}{\kern0pt}\ {\isacharequal}{\kern0pt}{\isachargreater}{\kern0pt}\ t\isanewline
\ \ \ \ \ \ \ \ \ \ \ \ \ \ {\isacharbar}{\kern0pt}\ ts\ {\isacharequal}{\kern0pt}{\isachargreater}{\kern0pt}\ foldr{\isadigit{1}}\ {\isacharparenleft}{\kern0pt}fn\ {\isacharparenleft}{\kern0pt}t{\isadigit{1}}{\isacharcomma}{\kern0pt}\ t{\isadigit{2}}{\isacharparenright}{\kern0pt}\ {\isacharequal}{\kern0pt}{\isachargreater}{\kern0pt}\ Cspm{\isacharunderscore}{\kern0pt}API{\isachardot}{\kern0pt}mk{\isacharunderscore}{\kern0pt}Inter\ t{\isadigit{1}}\ t{\isadigit{2}}{\isacharparenright}{\kern0pt}\ ts{\isacharparenright}{\kern0pt}{\isacharasterisk}{\kern0pt}{\isacharparenright}{\kern0pt}\isanewline
\ \ \ \ \ \ \ \ \ \ {\isacharparenleft}{\kern0pt}{\isacharasterisk}{\kern0pt}fun\ define{\isacharunderscore}{\kern0pt}cmd\ {\isacharparenleft}{\kern0pt}{\isacharbraceleft}{\kern0pt}nom{\isacharunderscore}{\kern0pt}thread{\isacharcomma}{\kern0pt}\ actions{\isacharcomma}{\kern0pt}\ {\isachardot}{\kern0pt}{\isachardot}{\kern0pt}{\isachardot}{\kern0pt}{\isacharbraceright}{\kern0pt}\ {\isacharcolon}{\kern0pt}\ thread{\isacharunderscore}{\kern0pt}absy{\isacharparenright}{\kern0pt}\ lthy\ {\isacharequal}{\kern0pt}\isanewline
\ \ \ \ \ \ \ \ \ \ \ \ let\isanewline
\ \ \ \ \ \ \ \ \ \ \ \ \ \ val\ name\ \ \ \ \ \ {\isacharequal}{\kern0pt}\ Binding{\isachardot}{\kern0pt}name{\isacharunderscore}{\kern0pt}of\ nom{\isacharunderscore}{\kern0pt}thread\isanewline
\ \ \ \ \ \ \ \ \ \ \ \ \ \ val\ const{\isacharunderscore}{\kern0pt}bnd\ {\isacharequal}{\kern0pt}\ Binding{\isachardot}{\kern0pt}name\ {\isacharparenleft}{\kern0pt}name\ {\isacharcircum}{\kern0pt}\ {\isachardoublequote}{\kern0pt}{\isacharunderscore}{\kern0pt}cmd{\isachardoublequote}{\kern0pt}{\isacharparenright}{\kern0pt}\isanewline
\ \ \ \ \ \ \ \ \ \ \ \ \ \ val\ com{\isacharunderscore}{\kern0pt}term\ \ {\isacharequal}{\kern0pt}\ quote{\isacharunderscore}{\kern0pt}com\ checked\ {\isacharparenleft}{\kern0pt}sequence{\isacharunderscore}{\kern0pt}actions\ actions{\isacharparenright}{\kern0pt}\isanewline
\ \ \ \ \ \ \ \ \ \ \ \ \ \ val\ {\isacharparenleft}{\kern0pt}{\isacharparenleft}{\kern0pt}lhs{\isacharcomma}{\kern0pt}\ {\isacharunderscore}{\kern0pt}{\isacharparenright}{\kern0pt}{\isacharcomma}{\kern0pt}\ lthy{\isacharprime}{\kern0pt}{\isacharparenright}{\kern0pt}\ {\isacharequal}{\kern0pt}\isanewline
\ \ \ \ \ \ \ \ \ \ \ \ \ \ \ \ \ \ \ \ Local{\isacharunderscore}{\kern0pt}Theory{\isachardot}{\kern0pt}define\ {\isacharparenleft}{\kern0pt}{\isacharparenleft}{\kern0pt}const{\isacharunderscore}{\kern0pt}bnd{\isacharcomma}{\kern0pt}\ NoSyn{\isacharparenright}{\kern0pt}{\isacharcomma}{\kern0pt}\isanewline
\ \ \ \ \ \ \ \ \ \ \ \ \ \ \ \ \ \ \ \ \ \ \ \ \ \ \ \ \ \ \ \ \ \ \ \ \ \ \ \ \ {\isacharparenleft}{\kern0pt}{\isacharparenleft}{\kern0pt}Binding{\isachardot}{\kern0pt}empty{\isacharcomma}{\kern0pt}\ {\isacharbrackleft}{\kern0pt}{\isacharbrackright}{\kern0pt}{\isacharparenright}{\kern0pt}{\isacharcomma}{\kern0pt}\ com{\isacharunderscore}{\kern0pt}term{\isacharparenright}{\kern0pt}{\isacharparenright}{\kern0pt}\ lthy\isanewline
\ \ \ \ \ \ \ \ \ \ \ \ \ \ val\ const{\isacharunderscore}{\kern0pt}name\ {\isacharequal}{\kern0pt}\ Local{\isacharunderscore}{\kern0pt}Theory{\isachardot}{\kern0pt}full{\isacharunderscore}{\kern0pt}name\ lthy{\isacharprime}{\kern0pt}\ const{\isacharunderscore}{\kern0pt}bnd\isanewline
\ \ \ \ \ \ \ \ \ \ \ \ \ \ val\ {\isacharunderscore}{\kern0pt}\ {\isacharequal}{\kern0pt}\ writeln\ {\isacharparenleft}{\kern0pt}{\isachardoublequote}{\kern0pt}{\isasymcheckmark}\ Defined\ constant{\isacharcolon}{\kern0pt}\ {\isachardoublequote}{\kern0pt}\ {\isacharcircum}{\kern0pt}\ const{\isacharunderscore}{\kern0pt}name{\isacharparenright}{\kern0pt}\isanewline
\ \ \ \ \ \ \ \ \ \ \ \ in\isanewline
\ \ \ \ \ \ \ \ \ \ \ \ \ \ {\isacharparenleft}{\kern0pt}const{\isacharunderscore}{\kern0pt}name{\isacharcomma}{\kern0pt}\ lthy{\isacharprime}{\kern0pt}{\isacharparenright}{\kern0pt}\ \ {\isacharparenleft}{\kern0pt}{\isacharasterisk}{\kern0pt}\ swapped{\isacharbang}{\kern0pt}\ {\isacharasterisk}{\kern0pt}{\isacharparenright}{\kern0pt}\isanewline
\ \ \ \ \ \ \ \ \ \ \ \ end{\isacharasterisk}{\kern0pt}{\isacharparenright}{\kern0pt}\isanewline
\ \ \ \ \ \ \ \ \ \ fun\ define{\isacharunderscore}{\kern0pt}cmd\ {\isacharparenleft}{\kern0pt}{\isacharbraceleft}{\kern0pt}nom{\isacharunderscore}{\kern0pt}thread{\isacharcomma}{\kern0pt}\ actions{\isacharcomma}{\kern0pt}\ {\isachardot}{\kern0pt}{\isachardot}{\kern0pt}{\isachardot}{\kern0pt}{\isacharbraceright}{\kern0pt}\ {\isacharcolon}{\kern0pt}\ thread{\isacharunderscore}{\kern0pt}absy{\isacharparenright}{\kern0pt}\ lthy\ {\isacharequal}{\kern0pt}\isanewline
\ \ \ \ \ \ \ \ \ \ \ \ let\isanewline
\ \ \ \ \ \ \ \ \ \ \ \ \ \ val\ name\ \ \ \ \ \ {\isacharequal}{\kern0pt}\ Binding{\isachardot}{\kern0pt}name{\isacharunderscore}{\kern0pt}of\ nom{\isacharunderscore}{\kern0pt}thread\isanewline
\ \ \ \ \ \ \ \ \ \ \ \ \ \ val\ const{\isacharunderscore}{\kern0pt}bnd\ {\isacharequal}{\kern0pt}\ Binding{\isachardot}{\kern0pt}name\ {\isacharparenleft}{\kern0pt}name\ {\isacharcircum}{\kern0pt}\ {\isachardoublequote}{\kern0pt}{\isacharunderscore}{\kern0pt}cmd{\isachardoublequote}{\kern0pt}{\isacharparenright}{\kern0pt}\isanewline
\ \ \ \ \ \ \ \ \ \ \ \ \ \ val\ com{\isacharunderscore}{\kern0pt}term\ \ {\isacharequal}{\kern0pt}\ quote{\isacharunderscore}{\kern0pt}com\ checked\ {\isacharparenleft}{\kern0pt}sequence{\isacharunderscore}{\kern0pt}actions\ actions{\isacharparenright}{\kern0pt}\isanewline
\ \ \ \ \ \ \ \ \ \ \ \ \ \ \isanewline
\ \ \ \ \ \ \ \ \ \ \ \ \ \ val\ {\isacharunderscore}{\kern0pt}\ {\isacharequal}{\kern0pt}\ Thm{\isachardot}{\kern0pt}cterm{\isacharunderscore}{\kern0pt}of\ {\isacharparenleft}{\kern0pt}Proof{\isacharunderscore}{\kern0pt}Context{\isachardot}{\kern0pt}init{\isacharunderscore}{\kern0pt}global\ {\isacharparenleft}{\kern0pt}Local{\isacharunderscore}{\kern0pt}Theory{\isachardot}{\kern0pt}exit{\isacharunderscore}{\kern0pt}global\ lthy{\isacharparenright}{\kern0pt}{\isacharparenright}{\kern0pt}\ com{\isacharunderscore}{\kern0pt}term\isanewline
\ \ \ \ \ \ \ \ \ \ \ \ \ \ \isanewline
\ \ \ \ \ \ \ \ \ \ \ \ \ \ val\ const{\isacharunderscore}{\kern0pt}type\ {\isacharequal}{\kern0pt}\ fastype{\isacharunderscore}{\kern0pt}of\ com{\isacharunderscore}{\kern0pt}term\isanewline
\ \ \ \ \ \ \ \ \ \ \ \ \ \ val\ lhs{\isacharunderscore}{\kern0pt}const\ {\isacharequal}{\kern0pt}\ Free\ {\isacharparenleft}{\kern0pt}Binding{\isachardot}{\kern0pt}name{\isacharunderscore}{\kern0pt}of\ const{\isacharunderscore}{\kern0pt}bnd{\isacharcomma}{\kern0pt}\ const{\isacharunderscore}{\kern0pt}type{\isacharparenright}{\kern0pt}\isanewline
\ \ \ \ \ \ \ \ \ \ \ \ \ \ val\ def{\isacharunderscore}{\kern0pt}eq\ {\isacharequal}{\kern0pt}\ Logic{\isachardot}{\kern0pt}mk{\isacharunderscore}{\kern0pt}equals\ {\isacharparenleft}{\kern0pt}lhs{\isacharunderscore}{\kern0pt}const{\isacharcomma}{\kern0pt}\ com{\isacharunderscore}{\kern0pt}term{\isacharparenright}{\kern0pt}\isanewline
\ \ \ \ \ \ \ \ \ \ \ \ \ \ \isanewline
\ \ \ \ \ \ \ \ \ \ \ \ \ \ val\ {\isacharparenleft}{\kern0pt}{\isacharparenleft}{\kern0pt}lhs{\isacharunderscore}{\kern0pt}term{\isacharcomma}{\kern0pt}\ {\isacharparenleft}{\kern0pt}{\isacharunderscore}{\kern0pt}{\isacharcomma}{\kern0pt}\ def{\isacharunderscore}{\kern0pt}thm{\isacharparenright}{\kern0pt}{\isacharparenright}{\kern0pt}{\isacharcomma}{\kern0pt}\ lthy{\isacharprime}{\kern0pt}{\isacharparenright}{\kern0pt}\ {\isacharequal}{\kern0pt}\ \isanewline
\ \ \ \ \ \ \ \ \ \ \ \ \ \ \ \ Specification{\isachardot}{\kern0pt}definition\isanewline
\ \ \ \ \ \ \ \ \ \ \ \ \ \ \ \ \ \ {\isacharbraceleft}{\kern0pt}verbose\ {\isacharequal}{\kern0pt}\ true{\isacharbraceright}{\kern0pt}\isanewline
\ \ \ \ \ \ \ \ \ \ \ \ \ \ \ \ \ \ {\isacharparenleft}{\kern0pt}SOME\ {\isacharparenleft}{\kern0pt}const{\isacharunderscore}{\kern0pt}bnd{\isacharcomma}{\kern0pt}\ NONE{\isacharcomma}{\kern0pt}\ NoSyn{\isacharparenright}{\kern0pt}{\isacharparenright}{\kern0pt}\ \ {\isacharparenleft}{\kern0pt}{\isacharasterisk}{\kern0pt}\ binding\ option\ {\isacharasterisk}{\kern0pt}{\isacharparenright}{\kern0pt}\isanewline
\ \ \ \ \ \ \ \ \ \ \ \ \ \ \ \ \ \ {\isacharbrackleft}{\kern0pt}{\isacharbrackright}{\kern0pt}\ \ \ \ \ \ \ \ \ \ \ \ \ \ \ \ \ \ \ \ \ \ \ \ \ \ \ \ \ \ \ {\isacharparenleft}{\kern0pt}{\isacharasterisk}{\kern0pt}\ parameters\ list\ {\isacharasterisk}{\kern0pt}{\isacharparenright}{\kern0pt}\isanewline
\ \ \ \ \ \ \ \ \ \ \ \ \ \ \ \ \ \ {\isacharbrackleft}{\kern0pt}{\isacharbrackright}{\kern0pt}\ \ \ \ \ \ \ \ \ \ \ \ \ \ \ \ \ \ \ \ \ \ \ \ \ \ \ \ \ \ \ {\isacharparenleft}{\kern0pt}{\isacharasterisk}{\kern0pt}\ premises\ list\ {\isacharasterisk}{\kern0pt}{\isacharparenright}{\kern0pt}\isanewline
\ \ \ \ \ \ \ \ \ \ \ \ \ \ \ \ \ \ {\isacharparenleft}{\kern0pt}{\isacharparenleft}{\kern0pt}Binding{\isachardot}{\kern0pt}empty{\isacharcomma}{\kern0pt}\ {\isacharbrackleft}{\kern0pt}{\isacharbrackright}{\kern0pt}{\isacharparenright}{\kern0pt}{\isacharcomma}{\kern0pt}\ def{\isacharunderscore}{\kern0pt}eq{\isacharparenright}{\kern0pt}\ \ \ \ {\isacharparenleft}{\kern0pt}{\isacharasterisk}{\kern0pt}\ attributes\ and\ equation\ {\isacharasterisk}{\kern0pt}{\isacharparenright}{\kern0pt}\isanewline
\ \ \ \ \ \ \ \ \ \ \ \ \ \ \ \ \ \ lthy\ \ \ \ \ \ \ \ \ \ \ \ \ \ \ \ \ \ \ \ \ \ \ \ \ \ \ \ \ {\isacharparenleft}{\kern0pt}{\isacharasterisk}{\kern0pt}\ local\ theory\ {\isacharasterisk}{\kern0pt}{\isacharparenright}{\kern0pt}\isanewline
\ \ \ \ \ \ \ \ \ \ \ \ \ \ \isanewline
\ \ \ \ \ \ \ \ \ \ \ \ \ \ val\ const{\isacharunderscore}{\kern0pt}name\ {\isacharequal}{\kern0pt}\ Local{\isacharunderscore}{\kern0pt}Theory{\isachardot}{\kern0pt}full{\isacharunderscore}{\kern0pt}name\ lthy{\isacharprime}{\kern0pt}\ const{\isacharunderscore}{\kern0pt}bnd\isanewline
\ \ \ \ \ \ \ \ \ \ \ \ \ \ val\ {\isacharunderscore}{\kern0pt}\ {\isacharequal}{\kern0pt}\ writeln\ {\isacharparenleft}{\kern0pt}{\isachardoublequote}{\kern0pt}{\isasymcheckmark}\ Defined\ constant{\isacharcolon}{\kern0pt}\ {\isachardoublequote}{\kern0pt}\ {\isacharcircum}{\kern0pt}\ const{\isacharunderscore}{\kern0pt}name{\isacharparenright}{\kern0pt}\isanewline
\ \ \ \ \ \ \ \ \ \ \ \ in\isanewline
\ \ \ \ \ \ \ \ \ \ \ \ \ \ {\isacharparenleft}{\kern0pt}const{\isacharunderscore}{\kern0pt}name{\isacharcomma}{\kern0pt}\ lthy{\isacharprime}{\kern0pt}{\isacharparenright}{\kern0pt}\ \ \isanewline
\ \ \ \ \ \ \ \ \ \ \ \ end\isanewline
\ \ \ \ \ \ \ \ \ \ {\isacharparenleft}{\kern0pt}{\isacharasterisk}{\kern0pt}\ {\isadigit{1}}{\isachardot}{\kern0pt}\ Définition\ des\ constantes\ de\ threads{\isacharcomma}{\kern0pt}\ et\ sortie\ du\ local{\isacharunderscore}{\kern0pt}theory\ {\isacharasterisk}{\kern0pt}{\isacharparenright}{\kern0pt}\isanewline
\ \ \ \ \ \ \ \ \ \ val\ {\isacharparenleft}{\kern0pt}cmd{\isacharunderscore}{\kern0pt}const{\isacharunderscore}{\kern0pt}names{\isacharcomma}{\kern0pt}\ lthy{\isacharunderscore}{\kern0pt}final{\isacharparenright}{\kern0pt}\ {\isacharequal}{\kern0pt}\isanewline
\ \ \ \ \ \ \ \ \ \ \ \ fold{\isacharunderscore}{\kern0pt}map\ define{\isacharunderscore}{\kern0pt}cmd\ {\isacharparenleft}{\kern0pt}{\isacharhash}{\kern0pt}threads{\isacharunderscore}{\kern0pt}decls\ checked{\isacharparenright}{\kern0pt}\ lthy{\isadigit{0}}\isanewline
\ \ \ \ \ \ \ \ \ \ \isanewline
\ \ \ \ \ \ \ \ \ \ val\ thy{\isacharprime}{\kern0pt}\ {\isacharequal}{\kern0pt}\ Local{\isacharunderscore}{\kern0pt}Theory{\isachardot}{\kern0pt}exit{\isacharunderscore}{\kern0pt}global\ lthy{\isacharunderscore}{\kern0pt}final\isanewline
\ \ \ \ \ \ \ \ \ \ \isanewline
\ \ \ \ \ \ \ \ \ \ {\isacharparenleft}{\kern0pt}{\isacharasterisk}{\kern0pt}\ {\isadigit{2}}{\isachardot}{\kern0pt}\ Maintenant\ on\ peut\ construire\ les\ termes\ Const{\isacharparenleft}{\kern0pt}{\isachardot}{\kern0pt}{\isachardot}{\kern0pt}{\isachardot}{\kern0pt}{\isacharparenright}{\kern0pt}\ depuis\ thy{\isacharprime}{\kern0pt}\ {\isacharasterisk}{\kern0pt}{\isacharparenright}{\kern0pt}\isanewline
\ \ \ \ \ \ \ \ \ \ val\ thread{\isacharunderscore}{\kern0pt}consts\ {\isacharcolon}{\kern0pt}\ term\ list\ {\isacharequal}{\kern0pt}\isanewline
\ \ \ \ \ \ \ \ \ \ \ \ map\ {\isacharparenleft}{\kern0pt}fn\ full{\isacharunderscore}{\kern0pt}name\ {\isacharequal}{\kern0pt}{\isachargreater}{\kern0pt}\ Const\ {\isacharparenleft}{\kern0pt}full{\isacharunderscore}{\kern0pt}name{\isacharcomma}{\kern0pt}\ {\isacharat}{\kern0pt}{\isacharbraceleft}{\kern0pt}typ\ com{\isacharbraceright}{\kern0pt}{\isacharparenright}{\kern0pt}{\isacharparenright}{\kern0pt}\ cmd{\isacharunderscore}{\kern0pt}const{\isacharunderscore}{\kern0pt}names\isanewline
\isanewline
\ \ \ \ \ \ \ \ \ \ {\isacharparenleft}{\kern0pt}{\isacharasterisk}{\kern0pt}\ A{\isachardot}{\kern0pt}\ \ noms\ des\ threads\ {\isacharminus}{\kern0pt}{\isacharminus}{\kern0pt}{\isacharminus}{\kern0pt}{\isacharminus}{\kern0pt}{\isacharminus}{\kern0pt}{\isacharminus}{\kern0pt}{\isacharminus}{\kern0pt}{\isacharminus}{\kern0pt}{\isacharminus}{\kern0pt}{\isacharminus}{\kern0pt}{\isacharminus}{\kern0pt}{\isacharminus}{\kern0pt}{\isacharminus}{\kern0pt}{\isacharminus}{\kern0pt}{\isacharminus}{\kern0pt}{\isacharminus}{\kern0pt}{\isacharminus}{\kern0pt}{\isacharminus}{\kern0pt}{\isacharminus}{\kern0pt}{\isacharminus}{\kern0pt}{\isacharminus}{\kern0pt}{\isacharminus}{\kern0pt}{\isacharminus}{\kern0pt}{\isacharminus}{\kern0pt}{\isacharminus}{\kern0pt}{\isacharminus}{\kern0pt}{\isacharminus}{\kern0pt}{\isacharminus}{\kern0pt}{\isacharminus}{\kern0pt}{\isacharminus}{\kern0pt}{\isacharminus}{\kern0pt}{\isacharminus}{\kern0pt}{\isacharminus}{\kern0pt}{\isacharminus}{\kern0pt}{\isacharminus}{\kern0pt}{\isacharminus}{\kern0pt}{\isacharminus}{\kern0pt}{\isacharminus}{\kern0pt}{\isacharminus}{\kern0pt}{\isacharminus}{\kern0pt}{\isacharminus}{\kern0pt}{\isacharminus}{\kern0pt}{\isacharminus}{\kern0pt}{\isacharminus}{\kern0pt}{\isacharminus}{\kern0pt}{\isacharminus}{\kern0pt}{\isacharminus}{\kern0pt}{\isacharminus}{\kern0pt}\ {\isacharasterisk}{\kern0pt}{\isacharparenright}{\kern0pt}\isanewline
\ \ \ \ \ \ \ \ \ \ val\ thread{\isacharunderscore}{\kern0pt}names\ {\isacharcolon}{\kern0pt}\ string\ list\ {\isacharequal}{\kern0pt}\isanewline
\ \ \ \ \ \ \ \ \ \ \ \ map\ {\isacharparenleft}{\kern0pt}fn\ {\isacharbraceleft}{\kern0pt}nom{\isacharunderscore}{\kern0pt}thread{\isacharcomma}{\kern0pt}\ {\isachardot}{\kern0pt}{\isachardot}{\kern0pt}{\isachardot}{\kern0pt}{\isacharbraceright}{\kern0pt}\ {\isacharequal}{\kern0pt}{\isachargreater}{\kern0pt}\ Binding{\isachardot}{\kern0pt}name{\isacharunderscore}{\kern0pt}of\ nom{\isacharunderscore}{\kern0pt}thread{\isacharparenright}{\kern0pt}\isanewline
\ \ \ \ \ \ \ \ \ \ \ \ \ \ \ \ {\isacharparenleft}{\kern0pt}{\isacharhash}{\kern0pt}threads{\isacharunderscore}{\kern0pt}decls\ checked{\isacharparenright}{\kern0pt}\isanewline
\ \ \ \ \ \ \ \ \ \ \isanewline
\isanewline
\ \ \ \ \ \ \ \ \ \ \isanewline
\ \ \ \ \ \ \ \ \ \ {\isacharparenleft}{\kern0pt}{\isacharasterisk}{\kern0pt}\ C{\isachardot}{\kern0pt}\ \ fabrique\ le\ réseau\ LOCALVARS\ pour\ un\ thread\ donné\ {\isacharminus}{\kern0pt}{\isacharminus}{\kern0pt}{\isacharminus}{\kern0pt}{\isacharminus}{\kern0pt}{\isacharminus}{\kern0pt}{\isacharminus}{\kern0pt}{\isacharminus}{\kern0pt}{\isacharminus}{\kern0pt}{\isacharminus}{\kern0pt}{\isacharminus}{\kern0pt}{\isacharminus}{\kern0pt}{\isacharminus}{\kern0pt}{\isacharminus}{\kern0pt}{\isacharminus}{\kern0pt}{\isacharminus}{\kern0pt}\ {\isacharasterisk}{\kern0pt}{\isacharparenright}{\kern0pt}\isanewline
\ \ \ \ \ \ \ \ \ \ fun\ localvars{\isacharunderscore}{\kern0pt}net\ {\isacharparenleft}{\kern0pt}{\isacharbraceleft}{\kern0pt}locals{\isacharunderscore}{\kern0pt}decl{\isacharcomma}{\kern0pt}\ locstab{\isacharcomma}{\kern0pt}\ {\isachardot}{\kern0pt}{\isachardot}{\kern0pt}{\isachardot}{\kern0pt}{\isacharbraceright}{\kern0pt}\ {\isacharcolon}{\kern0pt}\ thread{\isacharunderscore}{\kern0pt}absy{\isacharparenright}{\kern0pt}\ {\isacharcolon}{\kern0pt}\ term\ option\ {\isacharequal}{\kern0pt}\isanewline
\ \ \ \ \ \ \ \ \ \ \ \ \ \ if\ null\ locals{\isacharunderscore}{\kern0pt}decl\ then\ NONE\isanewline
\ \ \ \ \ \ \ \ \ \ \ \ \ \ else\ \ \ \ \isanewline
\ \ \ \ \ \ \ \ \ \ \ \ \ \ \ \ let\isanewline
\ \ \ \ \ \ \ \ \ \ \ \ \ \ \ \ \ \ \ \ val\ names\ {\isacharequal}{\kern0pt}\isanewline
\ \ \ \ \ \ \ \ \ \ \ \ \ \ \ \ \ \ \ \ \ \ map\ {\isacharparenleft}{\kern0pt}fn\ {\isacharparenleft}{\kern0pt}{\isacharparenleft}{\kern0pt}bdg{\isacharcomma}{\kern0pt}{\isacharunderscore}{\kern0pt}{\isacharparenright}{\kern0pt}{\isacharcomma}{\kern0pt}{\isacharunderscore}{\kern0pt}{\isacharparenright}{\kern0pt}\ {\isacharequal}{\kern0pt}{\isachargreater}{\kern0pt}\ HOLogic{\isachardot}{\kern0pt}mk{\isacharunderscore}{\kern0pt}string\ {\isacharparenleft}{\kern0pt}Binding{\isachardot}{\kern0pt}name{\isacharunderscore}{\kern0pt}of\ bdg{\isacharparenright}{\kern0pt}{\isacharparenright}{\kern0pt}\ locals{\isacharunderscore}{\kern0pt}decl\isanewline
\ \ \ \ \ \ \ \ \ \ \ \ \ \ \ \ \ \ \ \ val\ mset\ {\isacharequal}{\kern0pt}\ Cspm{\isacharunderscore}{\kern0pt}API{\isachardot}{\kern0pt}mk{\isacharunderscore}{\kern0pt}mset\ {\isacharparenleft}{\kern0pt}HOLogic{\isachardot}{\kern0pt}mk{\isacharunderscore}{\kern0pt}list\ {\isacharat}{\kern0pt}{\isacharbraceleft}{\kern0pt}typ\ string{\isacharbraceright}{\kern0pt}\ names{\isacharparenright}{\kern0pt}\isanewline
\ \ \ \ \ \ \ \ \ \ \ \ \ \ \isanewline
\ \ \ \ \ \ \ \ \ \ \ \ \ \ \ \ \ \ \ \ {\isacharparenleft}{\kern0pt}{\isacharasterisk}{\kern0pt}\ {\isacharminus}{\kern0pt}{\isacharminus}{\kern0pt}{\isacharminus}{\kern0pt}\ nouveau\ {\isasymsigma}\ pour\ CE\ thread\ {\isacharminus}{\kern0pt}{\isacharminus}{\kern0pt}{\isacharminus}{\kern0pt}\ {\isacharasterisk}{\kern0pt}{\isacharparenright}{\kern0pt}\isanewline
\ \ \ \ \ \ \ \ \ \ \ \ \ \ \ \ \ \ \ \ val\ sigma{\isacharunderscore}{\kern0pt}thread{\isacharunderscore}{\kern0pt}term\ {\isacharequal}{\kern0pt}\ make{\isacharunderscore}{\kern0pt}sigma{\isacharunderscore}{\kern0pt}term\ locstab\isanewline
\ \ \ \ \ \ \ \ \ \ \ \ \ \ \isanewline
\ \ \ \ \ \ \ \ \ \ \ \ \ \ \ \ \ \ \ \ val\ LOCALVARS{\isacharunderscore}{\kern0pt}partial\ {\isacharequal}{\kern0pt}\isanewline
\ \ \ \ \ \ \ \ \ \ \ \ \ \ \ \ \ \ \ \ \ \ \ \ \ \ Const\ {\isacharparenleft}{\kern0pt}{\isacharat}{\kern0pt}{\isacharbraceleft}{\kern0pt}const{\isacharunderscore}{\kern0pt}name\ LOCALVARS{\isacharbraceright}{\kern0pt}{\isacharcomma}{\kern0pt}\isanewline
\ \ \ \ \ \ \ \ \ \ \ \ \ \ \ \ \ \ \ \ \ \ \ \ \ \ \ \ \ \ \ \ \ \ {\isacharat}{\kern0pt}{\isacharbraceleft}{\kern0pt}typ\ {\isasymsigma}{\isacharbraceright}{\kern0pt}\ {\isacharminus}{\kern0pt}{\isacharminus}{\kern0pt}{\isachargreater}{\kern0pt}\ {\isacharat}{\kern0pt}{\isacharbraceleft}{\kern0pt}typ\ string{\isacharbraceright}{\kern0pt}\ {\isacharminus}{\kern0pt}{\isacharminus}{\kern0pt}{\isachargreater}{\kern0pt}\ procT{\isacharparenright}{\kern0pt}\isanewline
\ \ \ \ \ \ \ \ \ \ \ \ \ \ \isanewline
\ \ \ \ \ \ \ \ \ \ \ \ \ \ \ \ \ \ \ \ val\ lam\ {\isacharequal}{\kern0pt}\isanewline
\ \ \ \ \ \ \ \ \ \ \ \ \ \ \ \ \ \ \ \ \ \ \ \ \ \ Abs\ {\isacharparenleft}{\kern0pt}{\isachardoublequote}{\kern0pt}idx{\isachardoublequote}{\kern0pt}{\isacharcomma}{\kern0pt}\ {\isacharat}{\kern0pt}{\isacharbraceleft}{\kern0pt}typ\ string{\isacharbraceright}{\kern0pt}{\isacharcomma}{\kern0pt}\isanewline
\ \ \ \ \ \ \ \ \ \ \ \ \ \ \ \ \ \ \ \ \ \ \ \ \ \ \ \ \ \ \ LOCALVARS{\isacharunderscore}{\kern0pt}partial\ {\isachardollar}{\kern0pt}\ sigma{\isacharunderscore}{\kern0pt}thread{\isacharunderscore}{\kern0pt}term\ {\isachardollar}{\kern0pt}\ Bound\ {\isadigit{0}}{\isacharparenright}{\kern0pt}\isanewline
\ \ \isanewline
\ \ \ \ \ \ \ \ \ \ \ \ \ \ \ \ \ \ \ val\ processes\ {\isacharcolon}{\kern0pt}\ term\ list\ {\isacharequal}{\kern0pt}\isanewline
\ \ \ \ \ \ \ \ \ \ \ \ \ \ \ \ \ \ \ \ \ \ \ \ \ \ \ \ \ \ map\ {\isacharparenleft}{\kern0pt}fn\ name\ {\isacharequal}{\kern0pt}{\isachargreater}{\kern0pt}\ lam\ {\isachardollar}{\kern0pt}\ name{\isacharparenright}{\kern0pt}\ names\isanewline
\ \ \ \ \ \ \ \ \ \ \ \ \isanewline
\ \ \ \ \ \ \ \ \ \ \ \ \ \ \ \ \ \ \ \ val\ locals{\isacharunderscore}{\kern0pt}net{\isacharunderscore}{\kern0pt}term\ {\isacharcolon}{\kern0pt}\ term\ option\ {\isacharequal}{\kern0pt}\isanewline
\ \ \ \ \ \ \ \ \ \ \ \ \ \ \ \ \ \ \ \ {\isacharparenleft}{\kern0pt}case\ globals{\isacharunderscore}{\kern0pt}names\ of\isanewline
\ \ \ \ \ \ \ \ \ \ \ \ \ \ \ \ \ \ \ \ \ \ \ {\isacharbrackleft}{\kern0pt}{\isacharbrackright}{\kern0pt}\ {\isacharequal}{\kern0pt}{\isachargreater}{\kern0pt}\ NONE\isanewline
\ \ \ \ \ \ \ \ \ \ \ \ \ \ \ \ \ \ \ \ \ {\isacharbar}{\kern0pt}\ {\isacharunderscore}{\kern0pt}\ {\isacharequal}{\kern0pt}{\isachargreater}{\kern0pt}\ \isanewline
\ \ \ \ \ \ \ \ \ \ \ \ \ \ \ \ \ \ \ \ \ \ \ let\isanewline
\ \ \ \ \ \ \ \ \ \ \ \ \ \ \ \ \ \ \ \ \ \ \ \ \ val\ empty{\isacharunderscore}{\kern0pt}evs{\isacharunderscore}{\kern0pt}set\ {\isacharequal}{\kern0pt}\ Const\ {\isacharparenleft}{\kern0pt}{\isacharat}{\kern0pt}{\isacharbraceleft}{\kern0pt}const{\isacharunderscore}{\kern0pt}name\ bot{\isacharbraceright}{\kern0pt}{\isacharcomma}{\kern0pt}\ {\isacharat}{\kern0pt}{\isacharbraceleft}{\kern0pt}typ\ {\isachardoublequote}{\kern0pt}evs\ set{\isachardoublequote}{\kern0pt}{\isacharbraceright}{\kern0pt}{\isacharparenright}{\kern0pt}\isanewline
\ \ \ \ \ \ \ \ \ \ \ \ \ \ \ \ \ \ \ \ \ \ \ \ \ val\ multisync{\isacharunderscore}{\kern0pt}const\ {\isacharequal}{\kern0pt}\ Const\ {\isacharparenleft}{\kern0pt}{\isacharat}{\kern0pt}{\isacharbraceleft}{\kern0pt}const{\isacharunderscore}{\kern0pt}name\ MultiSync{\isacharbraceright}{\kern0pt}{\isacharcomma}{\kern0pt}\ \isanewline
\ \ \ \ \ \ \ \ \ \ \ \ \ \ \ \ \ \ \ \ \ \ \ \ \ \ \ \ \ \ \ \ \ \ \ \ \ \ \ \ \ \ \ \ \ \ \ \ \ \ \ \ \ {\isacharat}{\kern0pt}{\isacharbraceleft}{\kern0pt}typ\ {\isachardoublequote}{\kern0pt}evs\ set{\isachardoublequote}{\kern0pt}{\isacharbraceright}{\kern0pt}\ {\isacharminus}{\kern0pt}{\isacharminus}{\kern0pt}{\isachargreater}{\kern0pt}\ \isanewline
\ \ \ \ \ \ \ \ \ \ \ \ \ \ \ \ \ \ \ \ \ \ \ \ \ \ \ \ \ \ \ \ \ \ \ \ \ \ \ \ \ \ \ \ \ \ \ \ \ \ \ \ \ {\isacharat}{\kern0pt}{\isacharbraceleft}{\kern0pt}typ\ {\isachardoublequote}{\kern0pt}string\ multiset{\isachardoublequote}{\kern0pt}{\isacharbraceright}{\kern0pt}\ {\isacharminus}{\kern0pt}{\isacharminus}{\kern0pt}{\isachargreater}{\kern0pt}\ \isanewline
\ \ \ \ \ \ \ \ \ \ \ \ \ \ \ \ \ \ \ \ \ \ \ \ \ \ \ \ \ \ \ \ \ \ \ \ \ \ \ \ \ \ \ \ \ \ \ \ \ \ \ \ \ {\isacharat}{\kern0pt}{\isacharbraceleft}{\kern0pt}typ\ {\isachardoublequote}{\kern0pt}string\ {\isasymRightarrow}\ {\isacharparenleft}{\kern0pt}evs{\isacharparenright}{\kern0pt}\ process{\isachardoublequote}{\kern0pt}{\isacharbraceright}{\kern0pt}\ {\isacharminus}{\kern0pt}{\isacharminus}{\kern0pt}{\isachargreater}{\kern0pt}\ \isanewline
\ \ \ \ \ \ \ \ \ \ \ \ \ \ \ \ \ \ \ \ \ \ \ \ \ \ \ \ \ \ \ \ \ \ \ \ \ \ \ \ \ \ \ \ \ \ \ \ \ \ \ \ \ {\isacharat}{\kern0pt}{\isacharbraceleft}{\kern0pt}typ\ {\isachardoublequote}{\kern0pt}{\isacharparenleft}{\kern0pt}evs{\isacharparenright}{\kern0pt}\ process{\isachardoublequote}{\kern0pt}{\isacharbraceright}{\kern0pt}{\isacharparenright}{\kern0pt}\isanewline
\ \ \ \ \ \ \ \ \ \ \ \ \ \ \ \ \ \ \ \ \ \ \ \ \ val\ compact{\isacharunderscore}{\kern0pt}term\ {\isacharequal}{\kern0pt}\ multisync{\isacharunderscore}{\kern0pt}const\ {\isachardollar}{\kern0pt}\ empty{\isacharunderscore}{\kern0pt}evs{\isacharunderscore}{\kern0pt}set\ {\isachardollar}{\kern0pt}\ mset\ {\isachardollar}{\kern0pt}\ lam\isanewline
\ \ \ \ \ \ \ \ \ \ \ \ \ \ \ \ \ \ \ \ \ \ \ in\isanewline
\ \ \ \ \ \ \ \ \ \ \ \ \ \ \ \ \ \ \ \ \ \ \ \ \ SOME\ compact{\isacharunderscore}{\kern0pt}term\isanewline
\ \ \ \ \ \ \ \ \ \ \ \ \ \ \ \ \ \ \ \ \ \ \ end{\isacharparenright}{\kern0pt}\isanewline
\ \ \ \ \ \ \ \ \ \ \ \ \ \ \ \ in\isanewline
\ \ \ \ \ \ \ \ \ \ \ \ \ \ \ \ \ \ \ locals{\isacharunderscore}{\kern0pt}net{\isacharunderscore}{\kern0pt}term\isanewline
\ \ \ \ \ \ \ \ \ \ \ \ \ \ \ \ \ \ {\isacharparenleft}{\kern0pt}{\isacharasterisk}{\kern0pt}{\isacharparenleft}{\kern0pt}case\ processes\ of\isanewline
\ \ \ \ \ \ \ \ \ \ \ \ \ \ \ \ \ \ \ \ \ \ {\isacharbrackleft}{\kern0pt}{\isacharbrackright}{\kern0pt}\ {\isacharequal}{\kern0pt}{\isachargreater}{\kern0pt}\ NONE\isanewline
\ \ \ \ \ \ \ \ \ \ \ \ \ \ \ \ \ \ \ \ {\isacharbar}{\kern0pt}\ {\isacharbrackleft}{\kern0pt}t{\isacharbrackright}{\kern0pt}\ {\isacharequal}{\kern0pt}{\isachargreater}{\kern0pt}\ SOME\ t\isanewline
\ \ \ \ \ \ \ \ \ \ \ \ \ \ \ \ \ \ \ \ {\isacharbar}{\kern0pt}\ ts\ {\isacharequal}{\kern0pt}{\isachargreater}{\kern0pt}\ SOME\ {\isacharparenleft}{\kern0pt}foldr{\isadigit{1}}\ {\isacharparenleft}{\kern0pt}fn\ {\isacharparenleft}{\kern0pt}t{\isadigit{1}}{\isacharcomma}{\kern0pt}\ t{\isadigit{2}}{\isacharparenright}{\kern0pt}\ {\isacharequal}{\kern0pt}{\isachargreater}{\kern0pt}\ Cspm{\isacharunderscore}{\kern0pt}API{\isachardot}{\kern0pt}mk{\isacharunderscore}{\kern0pt}Inter\ t{\isadigit{1}}\ t{\isadigit{2}}{\isacharparenright}{\kern0pt}\ ts{\isacharparenright}{\kern0pt}{\isacharparenright}{\kern0pt}\isanewline
\ \ \ \ \ \ \ \ \ \ \ \ \ \ \ \ \ \ {\isacharasterisk}{\kern0pt}{\isacharparenright}{\kern0pt}\isanewline
\isanewline
\ \ \ \ \ \ \ \ \ \ \ \ \ \ \ \ end\isanewline
\ \ \ \ \ \ \ \ \ \ \ \ \isanewline
\ \ \ \ \ \ \ \ \ \ \isanewline
\ \ \ \ \ \ \ \ \ \ {\isacharparenleft}{\kern0pt}{\isacharasterisk}{\kern0pt}\ D{\isachardot}{\kern0pt}\ \ compose\ \ Sem\ ti{\isacharunderscore}{\kern0pt}cmd\ \ {\isacharbar}{\kern0pt}{\isacharbar}{\kern0pt}\ \ LOCALVARSᵢ\ \ {\isacharminus}{\kern0pt}{\isacharminus}{\kern0pt}{\isacharminus}{\kern0pt}{\isacharminus}{\kern0pt}{\isacharminus}{\kern0pt}{\isacharminus}{\kern0pt}{\isacharminus}{\kern0pt}{\isacharminus}{\kern0pt}{\isacharminus}{\kern0pt}{\isacharminus}{\kern0pt}{\isacharminus}{\kern0pt}{\isacharminus}{\kern0pt}{\isacharminus}{\kern0pt}{\isacharminus}{\kern0pt}{\isacharminus}{\kern0pt}{\isacharminus}{\kern0pt}{\isacharminus}{\kern0pt}{\isacharminus}{\kern0pt}{\isacharminus}{\kern0pt}{\isacharminus}{\kern0pt}{\isacharminus}{\kern0pt}{\isacharminus}{\kern0pt}{\isacharminus}{\kern0pt}{\isacharminus}{\kern0pt}{\isacharminus}{\kern0pt}{\isacharminus}{\kern0pt}{\isacharminus}{\kern0pt}{\isacharminus}{\kern0pt}\ {\isacharasterisk}{\kern0pt}{\isacharparenright}{\kern0pt}\isanewline
\ \ \ \ \ \ \ \ \ \isanewline
\ \ \ \ \ \ \ \ \ \ val\ sem{\isacharunderscore}{\kern0pt}const\ {\isacharequal}{\kern0pt}\ Const\ {\isacharparenleft}{\kern0pt}{\isacharat}{\kern0pt}{\isacharbraceleft}{\kern0pt}const{\isacharunderscore}{\kern0pt}name\ Sem{\isacharbraceright}{\kern0pt}{\isacharcomma}{\kern0pt}\ {\isacharat}{\kern0pt}{\isacharbraceleft}{\kern0pt}typ\ com{\isacharbraceright}{\kern0pt}\ {\isacharminus}{\kern0pt}{\isacharminus}{\kern0pt}{\isachargreater}{\kern0pt}\ procT{\isacharparenright}{\kern0pt}\isanewline
\ \ \ \ \ \ \ \ \ \ val\ thread{\isacharunderscore}{\kern0pt}processes\ {\isacharcolon}{\kern0pt}\ term\ list\ {\isacharequal}{\kern0pt}\isanewline
\ \ \ \ \ \ \ \ \ \ \ \ ListPair{\isachardot}{\kern0pt}map\ {\isacharparenleft}{\kern0pt}fn\ {\isacharparenleft}{\kern0pt}th{\isacharunderscore}{\kern0pt}absy{\isacharcomma}{\kern0pt}\ cmd{\isacharunderscore}{\kern0pt}term{\isacharparenright}{\kern0pt}\ {\isacharequal}{\kern0pt}{\isachargreater}{\kern0pt}\isanewline
\ \ \ \ \ \ \ \ \ \ \ \ \ \ \ \ let\isanewline
\ \ \ \ \ \ \ \ \ \ \ \ \ \ \ \ \ \ val\ sem{\isacharunderscore}{\kern0pt}proc\ {\isacharequal}{\kern0pt}\ sem{\isacharunderscore}{\kern0pt}const\ {\isachardollar}{\kern0pt}\ cmd{\isacharunderscore}{\kern0pt}term\isanewline
\ \ \ \ \ \ \ \ \ \ \ \ \ \ \ \ \ \ val\ lvars{\isacharunderscore}{\kern0pt}net\ {\isacharequal}{\kern0pt}\ localvars{\isacharunderscore}{\kern0pt}net\ th{\isacharunderscore}{\kern0pt}absy\ \isanewline
\ \ \ \ \ \ \ \ \ \ \ \ \ \ \ \ in\isanewline
\ \ \ \ \ \ \ \ \ \ \ \ \ \ \ \ \ \ case\ lvars{\isacharunderscore}{\kern0pt}net\ of\ \isanewline
\ \ \ \ \ \ \ \ \ \ \ \ \ \ \ \ \ \ NONE\ {\isacharequal}{\kern0pt}{\isachargreater}{\kern0pt}\ \ sem{\isacharunderscore}{\kern0pt}proc\ \isanewline
\ \ \ \ \ \ \ \ \ \ \ \ \ \ \ \ \ \ {\isacharbar}{\kern0pt}SOME\ t\ {\isacharequal}{\kern0pt}{\isachargreater}{\kern0pt}Cspm{\isacharunderscore}{\kern0pt}API{\isachardot}{\kern0pt}mk{\isacharunderscore}{\kern0pt}Sync\ sem{\isacharunderscore}{\kern0pt}proc\ UNIV{\isacharunderscore}{\kern0pt}\ t{\isacharparenleft}{\kern0pt}{\isacharasterisk}{\kern0pt}\ par{\isacharunderscore}{\kern0pt}const\ {\isachardollar}{\kern0pt}\ sem{\isacharunderscore}{\kern0pt}proc\ {\isachardollar}{\kern0pt}\ {\isacharat}{\kern0pt}{\isacharbraceleft}{\kern0pt}term\ UNIV{\isacharbraceright}{\kern0pt}\ {\isachardollar}{\kern0pt}\ t{\isacharasterisk}{\kern0pt}{\isacharparenright}{\kern0pt}\isanewline
\ \ \ \ \ \ \ \ \ \ \ \ \ \ \ \ end{\isacharparenright}{\kern0pt}\isanewline
\ \ \ \ \ \ \ \ \ \ \ \ \ \ {\isacharparenleft}{\kern0pt}{\isacharhash}{\kern0pt}threads{\isacharunderscore}{\kern0pt}decls\ checked{\isacharcomma}{\kern0pt}\ thread{\isacharunderscore}{\kern0pt}consts{\isacharparenright}{\kern0pt}\isanewline
\ \ \ \ \ \ \ \ \ \ \ \ \ \isanewline
\ \ \ \ \ \ \ \ \ \ {\isacharparenleft}{\kern0pt}{\isacharasterisk}{\kern0pt}\ E{\isachardot}{\kern0pt}\ \ réseau\ des\ threads\ {\isacharparenleft}{\kern0pt}{\isacharplus}{\kern0pt}\ locals{\isacharparenright}{\kern0pt}\ combiné\ avec\ \ {\isacharbar}{\kern0pt}{\isacharbar}{\kern0pt}{\isacharbar}{\kern0pt}\ {\isacharminus}{\kern0pt}{\isacharminus}{\kern0pt}{\isacharminus}{\kern0pt}{\isacharminus}{\kern0pt}{\isacharminus}{\kern0pt}{\isacharminus}{\kern0pt}{\isacharminus}{\kern0pt}{\isacharminus}{\kern0pt}{\isacharminus}{\kern0pt}{\isacharminus}{\kern0pt}{\isacharminus}{\kern0pt}{\isacharminus}{\kern0pt}{\isacharminus}{\kern0pt}{\isacharminus}{\kern0pt}{\isacharminus}{\kern0pt}{\isacharminus}{\kern0pt}{\isacharminus}{\kern0pt}\ {\isacharasterisk}{\kern0pt}{\isacharparenright}{\kern0pt}\isanewline
\ \ \ \ \ \ \ \ \ \ val\ thread{\isacharunderscore}{\kern0pt}net{\isacharunderscore}{\kern0pt}term\ {\isacharcolon}{\kern0pt}\ term\ {\isacharequal}{\kern0pt}\isanewline
\ \ \ \ \ \ \ \ \ \ \ \ {\isacharparenleft}{\kern0pt}case\ thread{\isacharunderscore}{\kern0pt}processes\ of\isanewline
\ \ \ \ \ \ \ \ \ \ \ \ \ \ \ \ {\isacharbrackleft}{\kern0pt}{\isacharbrackright}{\kern0pt}\ \ \ {\isacharequal}{\kern0pt}{\isachargreater}{\kern0pt}\ error\ {\isachardoublequote}{\kern0pt}No\ threads\ declared{\isachardoublequote}{\kern0pt}\isanewline
\ \ \ \ \ \ \ \ \ \ \ \ \ \ {\isacharbar}{\kern0pt}\ {\isacharbrackleft}{\kern0pt}t{\isacharbrackright}{\kern0pt}\ \ {\isacharequal}{\kern0pt}{\isachargreater}{\kern0pt}\ t\isanewline
\ \ \ \ \ \ \ \ \ \ \ \ \ \ {\isacharbar}{\kern0pt}\ ts\ \ \ {\isacharequal}{\kern0pt}{\isachargreater}{\kern0pt}\ foldr{\isadigit{1}}\ {\isacharparenleft}{\kern0pt}fn\ {\isacharparenleft}{\kern0pt}t{\isadigit{1}}{\isacharcomma}{\kern0pt}\ t{\isadigit{2}}{\isacharparenright}{\kern0pt}\ {\isacharequal}{\kern0pt}{\isachargreater}{\kern0pt}\ Cspm{\isacharunderscore}{\kern0pt}API{\isachardot}{\kern0pt}mk{\isacharunderscore}{\kern0pt}Inter\ t{\isadigit{1}}\ t{\isadigit{2}}{\isacharparenright}{\kern0pt}\ ts{\isacharparenright}{\kern0pt}\isanewline
\ \ \ \ \ \ \ \ \ \ \isanewline
\ \ \ \ \ \ \ \ \ \ \ \ \ \ \ \ \ \ \ \ val\ {\isacharunderscore}{\kern0pt}\ {\isacharequal}{\kern0pt}\ writeln\ {\isacharparenleft}{\kern0pt}{\isachardoublequote}{\kern0pt}{\isacharbrackleft}{\kern0pt}SYSTEM{\isacharbrackright}{\kern0pt}\ \ réseau\ THREADS\ {\isacharcolon}{\kern0pt}{\isacharbackslash}{\kern0pt}n\ \ {\isachardoublequote}{\kern0pt}\ {\isacharcircum}{\kern0pt}\isanewline
\ \ \ \ \ \ \ \ \ \ \ \ \ \ \ \ \ \ \ \ \ \ \ \ \ \ \ \ \ \ \ \ \ \ \ \ \ Syntax{\isachardot}{\kern0pt}string{\isacharunderscore}{\kern0pt}of{\isacharunderscore}{\kern0pt}term{\isacharunderscore}{\kern0pt}global\ thy{\isacharprime}{\kern0pt}\ thread{\isacharunderscore}{\kern0pt}net{\isacharunderscore}{\kern0pt}term{\isacharparenright}{\kern0pt}\isanewline
\ \ \ \ \ \ \ \ \ \ \ \ \ \ \ \ \ \ \ \ {\isacharparenleft}{\kern0pt}{\isacharasterisk}{\kern0pt}\ {\isadigit{5}}{\isachardot}{\kern0pt}\ Composition\ globale\ {\isacharcolon}{\kern0pt}\ \ {\isacharparenleft}{\kern0pt}SEMAPHORES\ {\isacharbar}{\kern0pt}{\isacharbar}{\kern0pt}{\isacharbar}{\kern0pt}\ GLOBALVARS{\isacharparenright}{\kern0pt}\ {\isacharbrackleft}{\kern0pt}{\isacharbar}{\kern0pt}UNIV{\isacharbar}{\kern0pt}{\isacharbrackright}{\kern0pt}\ \ THREADS\ \ {\isacharasterisk}{\kern0pt}{\isacharparenright}{\kern0pt}\isanewline
\isanewline
\ \ \ \ \ \ \ \ \ \ \ \ \ \ \ \ \ \ \ \ val\ net{\isadigit{1}}\ {\isacharequal}{\kern0pt}\ {\isacharparenleft}{\kern0pt}case\ globals{\isacharunderscore}{\kern0pt}net{\isacharunderscore}{\kern0pt}term\ of\ \isanewline
\ \ \ \ \ \ \ \ \ \ \ \ \ \ \ \ \ \ \ \ \ \ \ \ \ \ \ \ \ \ \ \ SOME\ t\ {\isacharequal}{\kern0pt}{\isachargreater}{\kern0pt}\ SOME{\isacharparenleft}{\kern0pt}\ Cspm{\isacharunderscore}{\kern0pt}API{\isachardot}{\kern0pt}mk{\isacharunderscore}{\kern0pt}Sync\ {\isacharparenleft}{\kern0pt}the\ semaphores{\isacharunderscore}{\kern0pt}net{\isacharunderscore}{\kern0pt}term{\isacharparenright}{\kern0pt}\ empty{\isacharunderscore}{\kern0pt}evs{\isacharunderscore}{\kern0pt}set\ t{\isacharparenright}{\kern0pt}\isanewline
\ \ \ \ \ \ \ \ \ \ \ \ \ \ \ \ \ \ \ \ \ \ \ \ \ \ \ \ \ \ \ \ \ \ \ \ \ \ \ \ \ \ {\isacharparenleft}{\kern0pt}{\isacharasterisk}{\kern0pt}\ par{\isacharunderscore}{\kern0pt}const\ \ {\isachardollar}{\kern0pt}\ thread{\isacharunderscore}{\kern0pt}net{\isacharunderscore}{\kern0pt}term\ \ {\isachardollar}{\kern0pt}\ {\isacharat}{\kern0pt}{\isacharbraceleft}{\kern0pt}term\ empty{\isacharbraceright}{\kern0pt}\ \ {\isachardollar}{\kern0pt}t\ {\isacharasterisk}{\kern0pt}{\isacharparenright}{\kern0pt}\isanewline
\ \ \ \ \ \ \ \ \ \ \ \ \ \ \ \ \ \ \ \ \ \ \ \ \ \ \ \ \ \ \ {\isacharbar}{\kern0pt}NONE\ {\isacharequal}{\kern0pt}{\isachargreater}{\kern0pt}\ \ semaphores{\isacharunderscore}{\kern0pt}net{\isacharunderscore}{\kern0pt}term{\isacharparenright}{\kern0pt}\isanewline
\ \ \isanewline
\ \ \ \ \ \ \ \ \ \ \ \ \ \ \ \ \ \ \ \ val\ full{\isacharunderscore}{\kern0pt}system{\isacharunderscore}{\kern0pt}term\ {\isacharequal}{\kern0pt}\ {\isacharparenleft}{\kern0pt}case\ net{\isadigit{1}}\ of\isanewline
\ \ \ \ \ \ \ \ \ \ \ \ \ \ \ \ \ \ \ \ \ \ \ \ \ \ \ \ \ \ \ \ \ \ \ \ \ \ \ \ \ \ SOME\ t\ {\isacharequal}{\kern0pt}{\isachargreater}{\kern0pt}\ Cspm{\isacharunderscore}{\kern0pt}API{\isachardot}{\kern0pt}mk{\isacharunderscore}{\kern0pt}Sync\ t\ UNIV{\isacharunderscore}{\kern0pt}\ thread{\isacharunderscore}{\kern0pt}net{\isacharunderscore}{\kern0pt}term\isanewline
\ \ \ \ \ \ \ \ \ \ \ \ \ \ \ \ \ \ \ \ \ \ \ \ \ \ \ \ \ \ \ \ \ \ \ \ \ \ \ \ \ \ {\isacharparenleft}{\kern0pt}{\isacharasterisk}{\kern0pt}\ par{\isacharunderscore}{\kern0pt}const\ \ {\isachardollar}{\kern0pt}\ net{\isadigit{1}}\ {\isachardollar}{\kern0pt}\ {\isacharat}{\kern0pt}{\isacharbraceleft}{\kern0pt}term\ UNIV{\isacharbraceright}{\kern0pt}\ {\isachardollar}{\kern0pt}\ {\isacharparenleft}{\kern0pt}t{\isacharparenright}{\kern0pt}\ {\isacharasterisk}{\kern0pt}{\isacharparenright}{\kern0pt}\isanewline
\ \ \ \ \ \ \ \ \ \ \ \ \ \ \ \ \ \ \ \ \ \ \ \ \ \ \ \ \ \ \ \ \ \ \ \ \ \ \ \ \ \ {\isacharbar}{\kern0pt}NONE\ {\isacharequal}{\kern0pt}{\isachargreater}{\kern0pt}\ thread{\isacharunderscore}{\kern0pt}net{\isacharunderscore}{\kern0pt}term{\isacharparenright}{\kern0pt}\isanewline
\ \ \ \ \ \ \ \ \ \ \isanewline
\ \ \ \ \ \ \ \ \ \ \ \ \ \ \ \ \ \ \ \ val\ {\isacharunderscore}{\kern0pt}\ {\isacharequal}{\kern0pt}\ writeln\ {\isacharparenleft}{\kern0pt}{\isachardoublequote}{\kern0pt}{\isacharbrackleft}{\kern0pt}SYSTEM{\isacharbrackright}{\kern0pt}\ \ réseau\ complet\ {\isacharcolon}{\kern0pt}{\isacharbackslash}{\kern0pt}n\ \ {\isachardoublequote}{\kern0pt}\ {\isacharcircum}{\kern0pt}\isanewline
\ \ \ \ \ \ \ \ \ \ \ \ \ \ \ \ \ \ \ \ \ \ \ \ \ \ \ \ \ \ \ \ \ \ \ \ \ Syntax{\isachardot}{\kern0pt}string{\isacharunderscore}{\kern0pt}of{\isacharunderscore}{\kern0pt}term{\isacharunderscore}{\kern0pt}global\ thy{\isacharprime}{\kern0pt}\ full{\isacharunderscore}{\kern0pt}system{\isacharunderscore}{\kern0pt}term{\isacharparenright}{\kern0pt}\isanewline
\ \ \ \ \ \ \ \ \ \ \ \ \ {\isacharparenleft}{\kern0pt}{\isacharasterisk}{\kern0pt}\ \ \ val\ {\isacharunderscore}{\kern0pt}\ {\isacharequal}{\kern0pt}\ tracing\ {\isacharparenleft}{\kern0pt}{\isachardoublequote}{\kern0pt}{\isacharbackslash}{\kern0pt}n\ {\isacharbrackleft}{\kern0pt}DEBUG{\isacharbrackright}{\kern0pt}\ Type\ threads\ {\isacharequal}{\kern0pt}\ {\isachardoublequote}{\kern0pt}\ {\isacharcircum}{\kern0pt}\isanewline
\ \ \ \ \ \ \ \ \ \ \ \ \ \ \ \ \ \ \ \ \ \ \ \ \ \ \ \ \ \ \ \ \ Syntax{\isachardot}{\kern0pt}string{\isacharunderscore}{\kern0pt}of{\isacharunderscore}{\kern0pt}typ{\isacharunderscore}{\kern0pt}global\ {\isacharat}{\kern0pt}{\isacharbraceleft}{\kern0pt}theory{\isacharbraceright}{\kern0pt}\ {\isacharparenleft}{\kern0pt}fastype{\isacharunderscore}{\kern0pt}of\ {\isacharparenleft}{\kern0pt}thread{\isacharunderscore}{\kern0pt}net{\isacharunderscore}{\kern0pt}term{\isacharparenright}{\kern0pt}\ {\isacharparenright}{\kern0pt}{\isacharparenright}{\kern0pt}\isanewline
\ \ \ \ \ \ \ \ \ \ \ \ \ \ \ \ \ \ val\ {\isacharunderscore}{\kern0pt}\ {\isacharequal}{\kern0pt}\ tracing\ {\isacharparenleft}{\kern0pt}{\isachardoublequote}{\kern0pt}{\isacharbackslash}{\kern0pt}n\ {\isacharbrackleft}{\kern0pt}DEBUG{\isacharbrackright}{\kern0pt}\ Type\ lam{\isacharunderscore}{\kern0pt}GLOBALVARS\ {\isacharequal}{\kern0pt}\ {\isachardoublequote}{\kern0pt}\ {\isacharcircum}{\kern0pt}\isanewline
\ \ \ \ \ \ \ \ \ \ \ \ \ \ \ \ \ \ \ \ \ \ \ \ \ \ \ \ \ \ \ \ \ Syntax{\isachardot}{\kern0pt}string{\isacharunderscore}{\kern0pt}of{\isacharunderscore}{\kern0pt}typ{\isacharunderscore}{\kern0pt}global\ {\isacharat}{\kern0pt}{\isacharbraceleft}{\kern0pt}theory{\isacharbraceright}{\kern0pt}\ {\isacharparenleft}{\kern0pt}fastype{\isacharunderscore}{\kern0pt}of\ {\isacharparenleft}{\kern0pt}the\ globals{\isacharunderscore}{\kern0pt}net{\isacharunderscore}{\kern0pt}term{\isacharparenright}{\kern0pt}\ {\isacharparenright}{\kern0pt}{\isacharparenright}{\kern0pt}\ \isanewline
\ \ \ \ \ \ \ \ \ \ \ \ \ \ \ \ \ \ val\ {\isacharunderscore}{\kern0pt}\ {\isacharequal}{\kern0pt}\ tracing\ {\isacharparenleft}{\kern0pt}{\isachardoublequote}{\kern0pt}{\isacharbackslash}{\kern0pt}n\ {\isacharbrackleft}{\kern0pt}DEBUG{\isacharbrackright}{\kern0pt}\ Type\ semaphores{\isacharunderscore}{\kern0pt}net{\isacharunderscore}{\kern0pt}term\ {\isacharequal}{\kern0pt}\ {\isachardoublequote}{\kern0pt}\ {\isacharcircum}{\kern0pt}\isanewline
\ \ \ \ \ \ \ \ \ \ \ \ \ \ \ \ \ \ \ \ \ \ \ \ \ \ \ \ \ \ \ \ \ Syntax{\isachardot}{\kern0pt}string{\isacharunderscore}{\kern0pt}of{\isacharunderscore}{\kern0pt}typ{\isacharunderscore}{\kern0pt}global\ {\isacharat}{\kern0pt}{\isacharbraceleft}{\kern0pt}theory{\isacharbraceright}{\kern0pt}\ {\isacharparenleft}{\kern0pt}fastype{\isacharunderscore}{\kern0pt}of\ {\isacharparenleft}{\kern0pt}the\ semaphores{\isacharunderscore}{\kern0pt}net{\isacharunderscore}{\kern0pt}term{\isacharparenright}{\kern0pt}\ {\isacharparenright}{\kern0pt}{\isacharparenright}{\kern0pt}\ \ \ \ \ \isanewline
\ \ \ \ \ \ \ \ \ \ \ \ \ \ \ \ \ \ val\ {\isacharunderscore}{\kern0pt}\ {\isacharequal}{\kern0pt}\ tracing\ {\isacharparenleft}{\kern0pt}{\isachardoublequote}{\kern0pt}{\isacharbackslash}{\kern0pt}n\ {\isacharbrackleft}{\kern0pt}DEBUG{\isacharbrackright}{\kern0pt}\ Type\ full{\isacharunderscore}{\kern0pt}term\ {\isacharequal}{\kern0pt}\ {\isachardoublequote}{\kern0pt}\ {\isacharcircum}{\kern0pt}\isanewline
\ \ \ \ \ \ \ \ \ \ \ \ \ \ \ \ \ \ \ \ \ \ \ \ \ \ \ \ \ \ \ \ \ Syntax{\isachardot}{\kern0pt}string{\isacharunderscore}{\kern0pt}of{\isacharunderscore}{\kern0pt}typ{\isacharunderscore}{\kern0pt}global\ {\isacharat}{\kern0pt}{\isacharbraceleft}{\kern0pt}theory{\isacharbraceright}{\kern0pt}\ {\isacharparenleft}{\kern0pt}fastype{\isacharunderscore}{\kern0pt}of\ {\isacharparenleft}{\kern0pt}full{\isacharunderscore}{\kern0pt}system{\isacharunderscore}{\kern0pt}term{\isacharparenright}{\kern0pt}\ {\isacharparenright}{\kern0pt}{\isacharparenright}{\kern0pt}{\isacharasterisk}{\kern0pt}{\isacharparenright}{\kern0pt}\ \ \ \ \ \isanewline
\ \ \ \ \ \ \ \ \ \ \ \ \ \ \ \ \ \ \ \ \ \isanewline
\ \ \ \ \ \ \ \ \ \ \ \ \ \ val\ {\isacharunderscore}{\kern0pt}\ {\isacharequal}{\kern0pt}Thm{\isachardot}{\kern0pt}cterm{\isacharunderscore}{\kern0pt}of\ {\isacharparenleft}{\kern0pt}Proof{\isacharunderscore}{\kern0pt}Context{\isachardot}{\kern0pt}init{\isacharunderscore}{\kern0pt}global\ thy{\isacharprime}{\kern0pt}{\isacharparenright}{\kern0pt}\ full{\isacharunderscore}{\kern0pt}system{\isacharunderscore}{\kern0pt}term\isanewline
\ \ \ \ \ \ \ \ \ \ \ \ \ \ val\ const{\isacharunderscore}{\kern0pt}bnd\ {\isacharequal}{\kern0pt}\ {\isacharhash}{\kern0pt}system\ checked\isanewline
\ \ \ \ \ \ \ \ \ \ \ \ \ \ val\ lthy{\isacharunderscore}{\kern0pt}sys\ {\isacharequal}{\kern0pt}\ Named{\isacharunderscore}{\kern0pt}Target{\isachardot}{\kern0pt}theory{\isacharunderscore}{\kern0pt}init\ thy{\isacharprime}{\kern0pt}\isanewline
\ \ \ \ \ \ \ \ \ \ \ \ \ \ \isanewline
\isanewline
\ \ \ \ \ \ \ \ \ \ \ \ \ \ val\ const{\isacharunderscore}{\kern0pt}type\ {\isacharequal}{\kern0pt}\ fastype{\isacharunderscore}{\kern0pt}of\ full{\isacharunderscore}{\kern0pt}system{\isacharunderscore}{\kern0pt}term\isanewline
\ \ \ \ \ \ \ \ \ \ \ \ \ \ val\ lhs{\isacharunderscore}{\kern0pt}const\ {\isacharequal}{\kern0pt}\ Free\ {\isacharparenleft}{\kern0pt}Binding{\isachardot}{\kern0pt}name{\isacharunderscore}{\kern0pt}of\ const{\isacharunderscore}{\kern0pt}bnd{\isacharcomma}{\kern0pt}\ const{\isacharunderscore}{\kern0pt}type{\isacharparenright}{\kern0pt}\isanewline
\ \ \ \ \ \ \ \ \ \ \ \ \ \ val\ def{\isacharunderscore}{\kern0pt}eq\ {\isacharequal}{\kern0pt}\ Logic{\isachardot}{\kern0pt}mk{\isacharunderscore}{\kern0pt}equals\ {\isacharparenleft}{\kern0pt}lhs{\isacharunderscore}{\kern0pt}const{\isacharcomma}{\kern0pt}\ full{\isacharunderscore}{\kern0pt}system{\isacharunderscore}{\kern0pt}term{\isacharparenright}{\kern0pt}\isanewline
\ \ \ \ \ \ \ \ \ \ \ \ \ \ \isanewline
\ \ \ \ \ \ \ \ \ \ \ \ \ \ val\ {\isacharparenleft}{\kern0pt}{\isacharparenleft}{\kern0pt}lhs{\isacharunderscore}{\kern0pt}term{\isacharcomma}{\kern0pt}\ {\isacharparenleft}{\kern0pt}{\isacharunderscore}{\kern0pt}{\isacharcomma}{\kern0pt}\ def{\isacharunderscore}{\kern0pt}thm{\isacharparenright}{\kern0pt}{\isacharparenright}{\kern0pt}{\isacharcomma}{\kern0pt}\ lthy{\isacharprime}{\kern0pt}{\isacharparenright}{\kern0pt}\ {\isacharequal}{\kern0pt}\ \isanewline
\ \ \ \ \ \ \ \ \ \ \ \ \ \ \ \ Specification{\isachardot}{\kern0pt}definition\isanewline
\ \ \ \ \ \ \ \ \ \ \ \ \ \ \ \ \ \ {\isacharbraceleft}{\kern0pt}verbose\ {\isacharequal}{\kern0pt}\ true{\isacharbraceright}{\kern0pt}\isanewline
\ \ \ \ \ \ \ \ \ \ \ \ \ \ \ \ \ \ {\isacharparenleft}{\kern0pt}SOME\ {\isacharparenleft}{\kern0pt}const{\isacharunderscore}{\kern0pt}bnd{\isacharcomma}{\kern0pt}\ NONE{\isacharcomma}{\kern0pt}\ NoSyn{\isacharparenright}{\kern0pt}{\isacharparenright}{\kern0pt}\ \ {\isacharparenleft}{\kern0pt}{\isacharasterisk}{\kern0pt}\ binding\ option\ {\isacharasterisk}{\kern0pt}{\isacharparenright}{\kern0pt}\isanewline
\ \ \ \ \ \ \ \ \ \ \ \ \ \ \ \ \ \ {\isacharbrackleft}{\kern0pt}{\isacharbrackright}{\kern0pt}\ \ \ \ \ \ \ \ \ \ \ \ \ \ \ \ \ \ \ \ \ \ \ \ \ \ \ \ \ \ \ {\isacharparenleft}{\kern0pt}{\isacharasterisk}{\kern0pt}\ parameters\ list\ {\isacharasterisk}{\kern0pt}{\isacharparenright}{\kern0pt}\isanewline
\ \ \ \ \ \ \ \ \ \ \ \ \ \ \ \ \ \ {\isacharbrackleft}{\kern0pt}{\isacharbrackright}{\kern0pt}\ \ \ \ \ \ \ \ \ \ \ \ \ \ \ \ \ \ \ \ \ \ \ \ \ \ \ \ \ \ \ {\isacharparenleft}{\kern0pt}{\isacharasterisk}{\kern0pt}\ premises\ list\ {\isacharasterisk}{\kern0pt}{\isacharparenright}{\kern0pt}\isanewline
\ \ \ \ \ \ \ \ \ \ \ \ \ \ \ \ \ \ {\isacharparenleft}{\kern0pt}{\isacharparenleft}{\kern0pt}Binding{\isachardot}{\kern0pt}empty{\isacharcomma}{\kern0pt}\ {\isacharbrackleft}{\kern0pt}{\isacharbrackright}{\kern0pt}{\isacharparenright}{\kern0pt}{\isacharcomma}{\kern0pt}\ def{\isacharunderscore}{\kern0pt}eq{\isacharparenright}{\kern0pt}\ \ \ {\isacharparenleft}{\kern0pt}{\isacharasterisk}{\kern0pt}\ attributes\ and\ equation\ {\isacharasterisk}{\kern0pt}{\isacharparenright}{\kern0pt}\isanewline
\ \ \ \ \ \ \ \ \ \ \ \ \ \ \ \ \ \ lthy{\isacharunderscore}{\kern0pt}sys\ \ \ \ \ \ \ \ \ \ \ \ \ \ \ \ \ \ \ \ \ \ \ \ \ {\isacharparenleft}{\kern0pt}{\isacharasterisk}{\kern0pt}\ local\ theory\ {\isacharasterisk}{\kern0pt}{\isacharparenright}{\kern0pt}\isanewline
\ \ \ \ \ \ \ \ \ \ \ \ \ \ \isanewline
\ \ \ \ \ \ \ \ \ \ \ \ \ \ {\isacharparenleft}{\kern0pt}{\isacharasterisk}{\kern0pt}\ Exit\ to\ global\ theory\ {\isacharasterisk}{\kern0pt}{\isacharparenright}{\kern0pt}\isanewline
\ \ \ \ \ \ \ \ \ \ \ \ \ \ val\ thy{\isacharprime}{\kern0pt}{\isacharprime}{\kern0pt}\ {\isacharequal}{\kern0pt}\ Local{\isacharunderscore}{\kern0pt}Theory{\isachardot}{\kern0pt}exit{\isacharunderscore}{\kern0pt}global\ lthy{\isacharprime}{\kern0pt}\isanewline
\ \ \ \ \ \ \ \ \ \ \ \ \ \ \ \ \ \isanewline
\ \ \ \ \ \ \ \ \ in\isanewline
\ \ \ \ \ \ \ \ \ \ \ \ \ \ thy{\isacharprime}{\kern0pt}{\isacharprime}{\kern0pt}\isanewline
\ \ \ \ \ \ \ \ \ \ \ \ end{\isacharparenright}{\kern0pt}{\isacharparenright}{\kern0pt}{\isacharparenright}{\kern0pt}\isanewline
{\isacartoucheclose}%
\endisatagML
{\isafoldML}%
%
\isadelimML
%
\endisadelimML
%
\isadelimdocument
%
\endisadelimdocument
%
\isatagdocument
%
\isamarkupsection{Tests%
}
\isamarkuptrue%
%
\endisatagdocument
{\isafolddocument}%
%
\isadelimdocument
%
\endisadelimdocument
\isakeywordONE{SYSTEM}\isamarkupfalse%
\ S\isanewline
\ \ \isakeywordTWO{globals}\ \isanewline
\ \ \isakeywordTWO{locks}\ \isanewline
\ \ \isakeywordTWO{thread}\ empty{\isacharcolon}{\kern0pt}\isanewline
\ \ \ \ \ \ \ \ \ \isakeywordTWO{actions}\ \isanewline
\ \ \isakeywordTWO{thread}\ empty{\isadigit{2}}{\isacharcolon}{\kern0pt}\isanewline
\ \ \ \ \ \ \ \ \ \isakeywordTWO{actions}\ \isanewline
end{\isacharsemicolon}{\kern0pt}\isanewline
\ \ \ \ \ \ \ \ \ \ \ \ \ \ \ \ \ \ \ \ \ \ \ \ \ \ \ \ \ \ \ \ \ \ \ \ \ \ \ \ \ \ \ \ \ \ \ \ \ \ \ \ \ \isanewline
\isanewline
\isakeywordONE{SYSTEM}\isamarkupfalse%
\ System{\isacharunderscore}{\kern0pt}A\isanewline
\ \ \isakeywordTWO{globals}\ x{\isacharcolon}{\kern0pt}{\isacartoucheopen}int{\isacartoucheclose}\ {\isacharequal}{\kern0pt}\ {\isacartoucheopen}{\isadigit{3}}\ {\isacharcolon}{\kern0pt}{\isacharcolon}{\kern0pt}\ int{\isacartoucheclose}\ var{\isadigit{1}}{\isacharcolon}{\kern0pt}{\isacartoucheopen}int{\isacartoucheclose}\ {\isacharequal}{\kern0pt}\ {\isacartoucheopen}{\isadigit{4}}{\isacharcolon}{\kern0pt}{\isacharcolon}{\kern0pt}int{\isacartoucheclose}\ \ var{\isadigit{4}}{\isacharcolon}{\kern0pt}{\isacartoucheopen}int{\isacartoucheclose}\ {\isacharequal}{\kern0pt}\ {\isacartoucheopen}{\isadigit{4}}{\isacharcolon}{\kern0pt}{\isacharcolon}{\kern0pt}int{\isacartoucheclose}\isanewline
\ \ \isakeywordTWO{locks}\ \ \ l{\isacharcolon}{\kern0pt}{\isacartoucheopen}{\isacharparenleft}{\kern0pt}{\isacharparenright}{\kern0pt}{\isacartoucheclose}\ \ \ \ l{\isadigit{2}}{\isacharcolon}{\kern0pt}{\isacartoucheopen}{\isacharparenleft}{\kern0pt}{\isacharparenright}{\kern0pt}{\isacartoucheclose}\ \ \ l{\isadigit{3}}{\isacharcolon}{\kern0pt}{\isacartoucheopen}{\isacharparenleft}{\kern0pt}{\isacharparenright}{\kern0pt}{\isacartoucheclose}\ \ \ \ \ \ \ \ \ \ \ \ \ \ \ \ \ \ \ \ \ \ \ \ \ \ \ \ \ \ \ \ \isanewline
\ \ \isakeywordTWO{thread}\ t{\isadigit{1}}\ {\isacharcolon}{\kern0pt}\ \isanewline
\ \ \ \ \ \ \ \ \ \isakeywordTWO{any}\ y\ {\isacharcolon}{\kern0pt}\ {\isacartoucheopen}int{\isacartoucheclose}\ {\isacharequal}{\kern0pt}\ {\isacartoucheopen}{\isadigit{3}}{\isadigit{6}}\ {\isacharcolon}{\kern0pt}{\isacharcolon}{\kern0pt}\ int{\isacartoucheclose}\ var{\isadigit{2}}\ {\isacharcolon}{\kern0pt}\ {\isacartoucheopen}int{\isacartoucheclose}\isanewline
\ \ \ \ \ \ \ \ \ \isakeywordTWO{actions}\ \isanewline
\ \ \ \ \ \ \ \ \ \isakeywordTWO{SKIP}{\isacharsemicolon}{\kern0pt}\isanewline
\ \ \ \ \ \ \ \ \ \isakeywordTWO{LOCK}\ l{\isacharsemicolon}{\kern0pt}\isanewline
\ \ \ \ \ \ \ \ \ \ y\ {\isacharless}{\kern0pt}{\isacharminus}{\kern0pt}\ var{\isadigit{1}}{\isacharsemicolon}{\kern0pt}\isanewline
\ \ \ \ \ \ \ \ \ \isakeywordTWO{UNLOCK}\ l{\isacharsemicolon}{\kern0pt}\isanewline
\ \ \ \ \ \ \ \ \ \isakeywordTWO{IF}\ {\isacartoucheopen}var{\isadigit{1}}{\isasymle}{\isacharparenleft}{\kern0pt}{\isadigit{3}}{\isacharcolon}{\kern0pt}{\isacharcolon}{\kern0pt}\ int{\isacharparenright}{\kern0pt}{\isacartoucheclose}\ \isakeywordTWO{THEN}\ \isanewline
\ \ \ \ \ \ \ \ \ \ \ \ \isakeywordTWO{WHILE}\ {\isacartoucheopen}var{\isadigit{1}}{\isasymle}{\isacharparenleft}{\kern0pt}{\isadigit{3}}{\isacharcolon}{\kern0pt}{\isacharcolon}{\kern0pt}\ int{\isacharparenright}{\kern0pt}{\isacartoucheclose}\ \isakeywordTWO{DO}\ \isanewline
\ \ \ \ \ \ \ \ \ \ \ \ \ \ \isakeywordTWO{LOCK}\ l{\isadigit{2}}{\isacharsemicolon}{\kern0pt}\isanewline
\ \ \ \ \ \ \ \ \ \ \ \ \ \ \isakeywordTWO{SKIP}{\isacharsemicolon}{\kern0pt}\isanewline
\ \ \ \ \ \ \ \ \ \ \ \ \ \ var{\isadigit{2}}\ {\isacharequal}{\kern0pt}\ {\isacartoucheopen}var{\isadigit{2}}\ {\isacharplus}{\kern0pt}\ {\isadigit{1}}\ {\isacharcolon}{\kern0pt}{\isacharcolon}{\kern0pt}int{\isacartoucheclose}{\isacharsemicolon}{\kern0pt}\isanewline
\ \ \ \ \ \ \ \ \ \ \ \ \ \ var{\isadigit{1}}\ {\isacharminus}{\kern0pt}{\isachargreater}{\kern0pt}\ {\isacartoucheopen}var{\isadigit{2}}\ {\isacharcolon}{\kern0pt}{\isacharcolon}{\kern0pt}\ int{\isacartoucheclose}{\isacharsemicolon}{\kern0pt}\isanewline
\ \ \ \ \ \ \ \ \ \ \ \ \ \ \isakeywordTWO{UNLOCK}\ l{\isadigit{2}}{\isacharsemicolon}{\kern0pt}\ \ \ \ \ \ \ \ \ \ \ \ \ \ \ \ \ \ \ \ \ \ \ \ \ \ \ \ \ \ \ \ \ \ \ \ \ \ \ \ \ \ \ \ \ \ \ \ \ \ \ \ \ \ \ \ \ \isanewline
\ \ \ \ \ \ \ \ \ \ \ \ \isakeywordTWO{DONE}\ \isanewline
\ \ \ \ \ \ \ \ \ \isakeywordTWO{ELSE}\isanewline
\ \ \ \ \ \ \ \ \ \ \ \ \isakeywordTWO{SKIP}{\isacharsemicolon}{\kern0pt}\isanewline
\ \ \ \ \ \ \ \ \ \isakeywordTWO{DONE}\isanewline
\ \ \ \ \ \ \ \ y\ \ {\isacharequal}{\kern0pt}\ {\isacartoucheopen}y{\isacharplus}{\kern0pt}{\isadigit{2}}\ {\isacharcolon}{\kern0pt}{\isacharcolon}{\kern0pt}\ int{\isacartoucheclose}{\isacharsemicolon}{\kern0pt}\isanewline
\ \ \ \ \ \ \ \ y\ \ {\isacharequal}{\kern0pt}\ {\isacartoucheopen}y{\isacharplus}{\kern0pt}{\isadigit{2}}\ {\isacharcolon}{\kern0pt}{\isacharcolon}{\kern0pt}\ int{\isacartoucheclose}{\isacharsemicolon}{\kern0pt}\isanewline
\ \ \ \ \ \ \ \ var{\isadigit{1}}\ {\isacharminus}{\kern0pt}{\isachargreater}{\kern0pt}\ {\isacartoucheopen}y\ {\isacharcolon}{\kern0pt}{\isacharcolon}{\kern0pt}\ int{\isacartoucheclose}{\isacharsemicolon}{\kern0pt}\isanewline
\isanewline
\ \ \isakeywordTWO{thread}\ t{\isadigit{2}}\ {\isacharcolon}{\kern0pt}\isanewline
\ \ \ \ \ \ \ \ \ \isakeywordTWO{any}\ var{\isadigit{2}}{\isacharcolon}{\kern0pt}{\isacartoucheopen}int{\isacartoucheclose}\ {\isacharequal}{\kern0pt}\ {\isacartoucheopen}{\isacharparenleft}{\kern0pt}{\isadigit{4}}{\isacharplus}{\kern0pt}{\isadigit{6}}{\isacharparenright}{\kern0pt}\ {\isacharcolon}{\kern0pt}{\isacharcolon}{\kern0pt}\ int{\isacartoucheclose}\ y\ {\isacharcolon}{\kern0pt}\ {\isacartoucheopen}int{\isacartoucheclose}\ {\isacharequal}{\kern0pt}\ {\isacartoucheopen}{\isadigit{2}}{\isadigit{8}}\ {\isacharcolon}{\kern0pt}{\isacharcolon}{\kern0pt}\ int{\isacartoucheclose}\isanewline
\ \ \ \ \ \ \ \ \ \isakeywordTWO{actions}\ \ \ \ \ \ \ \ \ \ \ \ \ \ \ \ \ \ \ \ \ \ \ \ \ \ \isanewline
\ \ \ \ \ \ \ \ \ \isakeywordTWO{SKIP}{\isacharsemicolon}{\kern0pt}\isanewline
\ \ \ \ \ \ \ \ \ \isakeywordTWO{LOCK}\ l{\isacharsemicolon}{\kern0pt}\isanewline
\isanewline
\ \ \ \ \ \ \ \ \ \isakeywordTWO{UNLOCK}\ l{\isacharsemicolon}{\kern0pt}\isanewline
\ \ \ \ \ \ \ \ \ \isakeywordTWO{IF}\ {\isacartoucheopen}{\isadigit{3}}\ {\isacharequal}{\kern0pt}\ {\isacharparenleft}{\kern0pt}{\isadigit{5}}{\isacharcolon}{\kern0pt}{\isacharcolon}{\kern0pt}\ int{\isacharparenright}{\kern0pt}{\isacartoucheclose}\ \isakeywordTWO{THEN}\ \isanewline
\ \ \ \ \ \ \ \ \ \ \ \ \isakeywordTWO{WHILE}\ {\isacartoucheopen}x\ {\isacharminus}{\kern0pt}{\isacharparenleft}{\kern0pt}{\isadigit{8}}\ {\isacharcolon}{\kern0pt}{\isacharcolon}{\kern0pt}int{\isacharparenright}{\kern0pt}\ {\isacharless}{\kern0pt}\ x\ {\isacartoucheclose}\ \isakeywordTWO{DO}\ \isanewline
\ \ \ \ \ \ \ \ \ \ \ \ \ \ \isakeywordTWO{LOCK}\ l{\isacharsemicolon}{\kern0pt}\ \isanewline
\ \ \ \ \ \ \ \ \ \ \ \ \ \ \isakeywordTWO{UNLOCK}\ l{\isacharsemicolon}{\kern0pt}\isanewline
\ \ \ \ \ \ \ \ \ \ \ \ \isakeywordTWO{DONE}\ \isanewline
\ \ \ \ \ \ \ \ \ \isakeywordTWO{ELSE}\isanewline
\ \ \ \ \ \ \ \ \ \ \ \ \isakeywordTWO{SKIP}{\isacharsemicolon}{\kern0pt}\isanewline
\ \ \ \ \ \ \ \ \ \isakeywordTWO{DONE}\isanewline
\ \ \ \ \ \ \ \ y\ {\isacharequal}{\kern0pt}\ {\isacartoucheopen}var{\isadigit{2}}\ {\isacharplus}{\kern0pt}\ {\isadigit{3}}\ {\isacharcolon}{\kern0pt}{\isacharcolon}{\kern0pt}\ int{\isacartoucheclose}{\isacharsemicolon}{\kern0pt}\ \isanewline
\ \ \isakeywordTWO{thread}\ t{\isadigit{3}}{\isacharcolon}{\kern0pt}\isanewline
\ \ \ \ \ \ \ \isakeywordTWO{any}\ test\ {\isacharcolon}{\kern0pt}int\ {\isacharequal}{\kern0pt}\ {\isacartoucheopen}{\isacharminus}{\kern0pt}{\isadigit{3}}\ {\isacharcolon}{\kern0pt}{\isacharcolon}{\kern0pt}int{\isacartoucheclose}\ \ var{\isacharunderscore}{\kern0pt}local\ {\isacharcolon}{\kern0pt}int\ {\isacharequal}{\kern0pt}\ {\isacartoucheopen}{\isadigit{4}}\ {\isacharcolon}{\kern0pt}{\isacharcolon}{\kern0pt}\ int{\isacartoucheclose}\ \ \ \isanewline
\ \ \ \ \ \ \ \isakeywordTWO{actions}\ \isanewline
\ \ \ \ \ \ \ \ \ \ var{\isacharunderscore}{\kern0pt}local\ {\isacharequal}{\kern0pt}\ {\isacartoucheopen}{\isacharparenleft}{\kern0pt}{\isadigit{4}}{\isacharplus}{\kern0pt}{\isadigit{4}}{\isadigit{2}}{\isacharparenright}{\kern0pt}\ {\isacharcolon}{\kern0pt}{\isacharcolon}{\kern0pt}\ int{\isacartoucheclose}{\isacharsemicolon}{\kern0pt}\isanewline
\ \ \ \ \ \ \ \ \ \ test\ {\isacharequal}{\kern0pt}\ {\isacartoucheopen}{\isadigit{2}}{\isacharasterisk}{\kern0pt}\ var{\isacharunderscore}{\kern0pt}local\ {\isacharcolon}{\kern0pt}{\isacharcolon}{\kern0pt}\ int{\isacartoucheclose}{\isacharsemicolon}{\kern0pt}\isanewline
\ \ \ \ \ \ \ \ \ \ test\ {\isacharequal}{\kern0pt}\ {\isacartoucheopen}test{\isacharplus}{\kern0pt}{\isadigit{3}}\ {\isacharcolon}{\kern0pt}{\isacharcolon}{\kern0pt}\ int{\isacartoucheclose}{\isacharsemicolon}{\kern0pt}\isanewline
\ \ \ \ \ \ \ \isakeywordTWO{IF}\ {\isacartoucheopen}x\ {\isachargreater}{\kern0pt}\ {\isacharparenleft}{\kern0pt}{\isadigit{6}}\ {\isacharcolon}{\kern0pt}{\isacharcolon}{\kern0pt}\ int{\isacharparenright}{\kern0pt}{\isacartoucheclose}\ \isanewline
\ \ \ \ \ \ \ \ \ \ \isakeywordTWO{THEN}\ \isanewline
\ \ \ \ \ \ \ \ \ \ \ \ \isakeywordTWO{WHILE}\ {\isacartoucheopen}x\ {\isachargreater}{\kern0pt}\ {\isacharparenleft}{\kern0pt}{\isadigit{2}}\ {\isacharcolon}{\kern0pt}{\isacharcolon}{\kern0pt}\ int{\isacharparenright}{\kern0pt}{\isacartoucheclose}\ \isakeywordTWO{DO}\ \isanewline
\ \ \ \ \ \ \ \ \ \ \ \ \ \ \ test\ {\isacharequal}{\kern0pt}\ {\isacartoucheopen}test{\isacharminus}{\kern0pt}{\isadigit{3}}\ {\isacharcolon}{\kern0pt}{\isacharcolon}{\kern0pt}\ int{\isacartoucheclose}{\isacharsemicolon}{\kern0pt}\isanewline
\ \ \ \ \ \ \ \ \ \ \ \ \isakeywordTWO{DONE}\ \isanewline
\ \ \ \ \ \ \ \ \ \ \isakeywordTWO{ELSE}\isanewline
\ \ \ \ \ \ \ \ \ \ \ \ test\ {\isacharequal}{\kern0pt}\ {\isacartoucheopen}test{\isacharplus}{\kern0pt}{\isadigit{3}}\ {\isacharcolon}{\kern0pt}{\isacharcolon}{\kern0pt}\ int{\isacartoucheclose}{\isacharsemicolon}{\kern0pt}\isanewline
\ \ \ \ \ \ \ \ \ \ \isakeywordTWO{DONE}\isanewline
end{\isacharsemicolon}{\kern0pt}\ \ \ \ \ \ \ \ \ \ \ \ \ \ \ \isanewline
\isanewline
\isakeywordONE{thm}\isamarkupfalse%
\ t{\isadigit{1}}{\isacharunderscore}{\kern0pt}cmd{\isacharunderscore}{\kern0pt}def\ empty{\isadigit{2}}{\isacharunderscore}{\kern0pt}cmd{\isacharunderscore}{\kern0pt}def\ \isanewline
\isakeywordONE{thm}\isamarkupfalse%
\ System{\isacharunderscore}{\kern0pt}A{\isacharunderscore}{\kern0pt}def\isanewline
%
\isadelimML
\isanewline
%
\endisadelimML
%
\isatagML
\isakeywordONE{ML}\isamarkupfalse%
{\isacartoucheopen}\isanewline
val\ temp\ \ {\isacharequal}{\kern0pt}\ \isactrlterm {\isasymopen}Sem{\isacharparenleft}{\kern0pt}t{\isadigit{1}}{\isacharunderscore}{\kern0pt}cmd{\isacharparenright}{\kern0pt}{\isasymclose}\isanewline
val\ temp{\isadigit{2}}\ {\isacharequal}{\kern0pt}\ \isactrlterm {\isasymopen}Sem{\isacharparenleft}{\kern0pt}t{\isadigit{1}}{\isacharunderscore}{\kern0pt}cmd{\isacharparenright}{\kern0pt}{\isasymclose}\isanewline
val\ temp{\isadigit{3}}\ {\isacharequal}{\kern0pt}\ \isactrlterm {\isasymopen}{\isacharparenleft}{\kern0pt}t{\isadigit{1}}{\isacharunderscore}{\kern0pt}cmd{\isacharparenright}{\kern0pt}{\isasymclose}\ \ \ \ \ \ \ \ \ \isanewline
val\ temp{\isadigit{4}}\ {\isacharequal}{\kern0pt}\ \isactrlterm {\isasymopen}empty{\isadigit{2}}{\isacharunderscore}{\kern0pt}cmd{\isasymclose}\isanewline
\isanewline
val\ tem{\isadigit{2}}{\isacharequal}{\kern0pt}\ Cspm{\isacharunderscore}{\kern0pt}API{\isachardot}{\kern0pt}mk{\isacharunderscore}{\kern0pt}Det\ temp\ temp{\isadigit{2}}\isanewline
{\isacartoucheclose}\isanewline
\isanewline
\isakeywordONE{ML}\isamarkupfalse%
{\isacartoucheopen}\isanewline
val\ temp\ {\isacharequal}{\kern0pt}\ \isactrlterm {\isasymopen}System{\isacharunderscore}{\kern0pt}A{\isasymclose}\isanewline
\isanewline
val\ c\ {\isacharequal}{\kern0pt}\ Thm{\isachardot}{\kern0pt}cterm{\isacharunderscore}{\kern0pt}of\ {\isacharat}{\kern0pt}{\isacharbraceleft}{\kern0pt}context{\isacharbraceright}{\kern0pt}\ temp\isanewline
\isanewline
{\isacartoucheclose}%
\endisatagML
{\isafoldML}%
%
\isadelimML
\isanewline
%
\endisadelimML
\isanewline
\isakeywordONE{find{\isacharunderscore}{\kern0pt}theorems}\isamarkupfalse%
\ {\isacharparenleft}{\kern0pt}{\isadigit{5}}{\isadigit{0}}{\isadigit{0}}{\isacharparenright}{\kern0pt}name\ {\isacharcolon}{\kern0pt}\ {\isachardoublequoteopen}CIL{\isachardoublequoteclose}\isanewline
\ \ \ \ \ \ \ \ \ \ \ \ \isanewline
\isakeywordONE{lemmas}\isamarkupfalse%
\ X\ {\isacharequal}{\kern0pt}\ System{\isacharunderscore}{\kern0pt}A{\isacharunderscore}{\kern0pt}def\ {\isacharbrackleft}{\kern0pt}simplified\ t{\isadigit{1}}{\isacharunderscore}{\kern0pt}cmd{\isacharunderscore}{\kern0pt}def\ t{\isadigit{2}}{\isacharunderscore}{\kern0pt}cmd{\isacharunderscore}{\kern0pt}def\ t{\isadigit{3}}{\isacharunderscore}{\kern0pt}cmd{\isacharunderscore}{\kern0pt}def\ Sem{\isacharunderscore}{\kern0pt}def{\isacharbrackright}{\kern0pt}\isanewline
\isanewline
%
\isadelimtheory
%
\endisadelimtheory
%
\isatagtheory
%
\endisatagtheory
{\isafoldtheory}%
%
\isadelimtheory
%
\endisadelimtheory
%
\end{isabellebody}%
\endinput
%:%file=$AFP/IMP_concur/CIL.thy%:%
%:%11=1%:%
%:%27=3%:%
%:%28=3%:%
%:%29=4%:%
%:%30=5%:%
%:%31=6%:%
%:%32=7%:%
%:%33=8%:%
%:%34=9%:%
%:%48=11%:%
%:%60=13%:%
%:%61=14%:%
%:%62=15%:%
%:%63=16%:%
%:%67=18%:%
%:%68=19%:%
%:%77=21%:%
%:%87=23%:%
%:%88=23%:%
%:%89=24%:%
%:%90=25%:%
%:%91=25%:%
%:%92=25%:%
%:%93=26%:%
%:%94=27%:%
%:%95=27%:%
%:%96=28%:%
%:%97=29%:%
%:%98=30%:%
%:%99=30%:%
%:%100=31%:%
%:%107=34%:%
%:%123=36%:%
%:%124=36%:%
%:%246=151%:%
%:%247=152%:%
%:%248=153%:%
%:%249=154%:%
%:%250=154%:%
%:%257=156%:%
%:%273=157%:%
%:%274=157%:%
%:%555=427%:%
%:%571=428%:%
%:%572=428%:%
%:%1490=1333%:%
%:%1500=1337%:%
%:%1501=1337%:%
%:%1502=1338%:%
%:%1505=1339%:%
%:%1510=1340%:%
%:%1511=1340%:%
%:%1679=1508%:%
%:%1680=1509%:%
%:%1681=1510%:%
%:%1682=1510%:%
%:%1721=1549%:%
%:%1722=1550%:%
%:%1723=1551%:%
%:%1724=1551%:%
%:%1875=1695%:%
%:%1891=1696%:%
%:%1892=1696%:%
%:%2299=2092%:%
%:%2309=2094%:%
%:%2310=2094%:%
%:%2311=2095%:%
%:%2312=2096%:%
%:%2313=2097%:%
%:%2314=2098%:%
%:%2315=2099%:%
%:%2316=2100%:%
%:%2317=2101%:%
%:%2318=2102%:%
%:%2319=2103%:%
%:%2320=2104%:%
%:%2321=2104%:%
%:%2322=2105%:%
%:%2323=2106%:%
%:%2324=2107%:%
%:%2325=2108%:%
%:%2326=2109%:%
%:%2327=2110%:%
%:%2328=2111%:%
%:%2329=2112%:%
%:%2330=2113%:%
%:%2331=2114%:%
%:%2332=2115%:%
%:%2333=2116%:%
%:%2334=2117%:%
%:%2335=2118%:%
%:%2336=2119%:%
%:%2337=2120%:%
%:%2338=2121%:%
%:%2339=2122%:%
%:%2340=2123%:%
%:%2341=2124%:%
%:%2342=2125%:%
%:%2343=2126%:%
%:%2344=2127%:%
%:%2345=2128%:%
%:%2346=2129%:%
%:%2347=2130%:%
%:%2348=2131%:%
%:%2349=2132%:%
%:%2350=2133%:%
%:%2351=2134%:%
%:%2352=2135%:%
%:%2353=2136%:%
%:%2354=2137%:%
%:%2355=2138%:%
%:%2356=2139%:%
%:%2357=2140%:%
%:%2358=2141%:%
%:%2359=2142%:%
%:%2360=2143%:%
%:%2361=2144%:%
%:%2362=2145%:%
%:%2363=2146%:%
%:%2364=2147%:%
%:%2365=2148%:%
%:%2366=2149%:%
%:%2367=2150%:%
%:%2368=2151%:%
%:%2369=2152%:%
%:%2370=2153%:%
%:%2371=2154%:%
%:%2372=2155%:%
%:%2373=2156%:%
%:%2374=2157%:%
%:%2375=2158%:%
%:%2376=2159%:%
%:%2377=2160%:%
%:%2378=2161%:%
%:%2379=2161%:%
%:%2380=2162%:%
%:%2381=2162%:%
%:%2384=2163%:%
%:%2389=2164%:%
%:%2390=2164%:%
%:%2397=2171%:%
%:%2398=2172%:%
%:%2399=2173%:%
%:%2400=2173%:%
%:%2410=2178%:%
%:%2413=2179%:%
%:%2414=2180%:%
%:%2415=2180%:%
%:%2416=2181%:%
%:%2417=2182%:%
%:%2418=2182%:%
%:%2419=2183%:%



% optional bibliography
\bibliographystyle{abbrv}
\bibliography{root}

\end{document}

%%% Local Variables:
%%% mode: latex
%%% TeX-master: t
%%% End:
\endinput
%:%file=$AFP/IMP_concur/document/root.tex%:%
